% 本文件由 LaTeX 排版 Agent 根据有效 translation.json 自动生成。
% author=Gregory the Great, Pope (c. 540-604); work_id=3602; title=圣大额我略书翰集

\chapter{\WorkTitle{圣大额我略书翰集}}

% structure: p0001 source=h1 target=omitted reason=page-title-already-rendered-by-parent

\Emph{第一卷：}\newline{}1 |\newline{}2 |\newline{}3 |\newline{}4 |\newline{}5 |\newline{}6 |\newline{}7 |\newline{}9 |\newline{}10 |\newline{}11 |\newline{}12 |\newline{}16 |\newline{}17 |\newline{}18 |\newline{}19 |\newline{}20 |\newline{}21 |\newline{}25 |\newline{}26 |\newline{}27 |\newline{}28 |\newline{}29 |\newline{}33 |\newline{}34 |\newline{}35 |\newline{}36 |\newline{}39 |\newline{}41 |\newline{}42 |\newline{}43 |\newline{}44 |\newline{}46 |\newline{}47 |\newline{}48 |\newline{}49 |\newline{}50 |\newline{}52 |\newline{}56 |\newline{}57 |\newline{}58 |\newline{}61 |\newline{}62 |\newline{}63 |\newline{}66 |\newline{}67 |\newline{}72 |\newline{}74 |\newline{}75 |\newline{}77 |\newline{}78 |\newline{}79 |\newline{}80\par

\Emph{第二卷：}\newline{}3 |\newline{}6 |\newline{}7 |\newline{}9 |\newline{}10 |\newline{}12 |\newline{}15 |\newline{}18 |\newline{}19 |\newline{}20 |\newline{}22 |\newline{}23 |\newline{}26 |\newline{}27 |\newline{}29 |\newline{}30 |\newline{}32 |\newline{}33 |\newline{}34 |\newline{}36 |\newline{}37 |\newline{}41 |\newline{}42 |\newline{}46 |\newline{}47 |\newline{}48 |\newline{}49 |\newline{}51 |\newline{}52 |\newline{}54\par

\Emph{第三卷：}\newline{}1 |\newline{}2 |\newline{}3 |\newline{}5 |\newline{}6 |\newline{}7 |\newline{}8 |\newline{}9 |\newline{}10 |\newline{}12 |\newline{}15 |\newline{}22 |\newline{}29 |\newline{}30 |\newline{}31 |\newline{}32 |\newline{}33 |\newline{}35 |\newline{}36 |\newline{}38 |\newline{}45 |\newline{}46 |\newline{}47 |\newline{}48 |\newline{}49 |\newline{}51 |\newline{}53 |\newline{}56 |\newline{}57 |\newline{}59 |\newline{}60 |\newline{}65 |\newline{}66 |\newline{}67\par

\Emph{第四卷：}\newline{}1 |\newline{}2 |\newline{}3 |\newline{}4 |\newline{}5 |\newline{}6 |\newline{}7 |\newline{}8 |\newline{}9 |\newline{}10 |\newline{}11 |\newline{}15 |\newline{}18 |\newline{}20 |\newline{}21 |\newline{}23 |\newline{}24 |\newline{}25 |\newline{}26 |\newline{}27 |\newline{}29 |\newline{}30 |\newline{}31 |\newline{}32 |\newline{}33 |\newline{}34 |\newline{}35 |\newline{}36 |\newline{}38 |\newline{}39 |\newline{}46 |\newline{}47\par

\Emph{第五卷：}\newline{}2 |\newline{}4 |\newline{}5 |\newline{}8 |\newline{}11 |\newline{}15 |\newline{}17 |\newline{}18 |\newline{}19 |\newline{}20 |\newline{}21 |\newline{}23 |\newline{}25 |\newline{}26 |\newline{}29 |\newline{}30 |\newline{}36 |\newline{}39 |\newline{}40 |\newline{}41 |\newline{}42 |\newline{}43 |\newline{}48 |\newline{}49 |\newline{}52 |\newline{}53 |\newline{}54 |\newline{}55 |\newline{}56 |\newline{}57 |\newline{}58\par

\Emph{第六卷：}\newline{}1 |\newline{}2 |\newline{}3 |\newline{}5 |\newline{}6 |\newline{}7 |\newline{}8 |\newline{}9 |\newline{}12 |\newline{}14 |\newline{}15 |\newline{}16 |\newline{}17 |\newline{}18 |\newline{}22 |\newline{}24 |\newline{}25 |\newline{}26 |\newline{}27 |\newline{}29 |\newline{}30 |\newline{}32 |\newline{}34 |\newline{}35 |\newline{}37 |\newline{}43 |\newline{}44 |\newline{}46 |\newline{}48 |\newline{}49 |\newline{}50 |\newline{}51 |\newline{}52 |\newline{}53 |\newline{}54 |\newline{}55 |\newline{}56 |\newline{}57 |\newline{}58 |\newline{}59 |\newline{}60 |\newline{}61 |\newline{}63 |\newline{}65 |\newline{}66\par

\Emph{第七卷：}\newline{}2 |\newline{}4 |\newline{}5 |\newline{}6 |\newline{}7 |\newline{}11 |\newline{}12 |\newline{}13 |\newline{}15 |\newline{}17 |\newline{}19 |\newline{}20 |\newline{}23 |\newline{}25 |\newline{}26 |\newline{}27 |\newline{}28 |\newline{}30 |\newline{}31 |\newline{}32 |\newline{}33 |\newline{}34 |\newline{}35 |\newline{}38 |\newline{}39 |\newline{}40 |\newline{}42 |\newline{}43\par

\Emph{第八卷：}\newline{}1 |\newline{}2 |\newline{}3 |\newline{}5 |\newline{}6 |\newline{}10 |\newline{}13 |\newline{}14 |\newline{}15 |\newline{}17 |\newline{}18 |\newline{}20 |\newline{}21 |\newline{}22 |\newline{}23 |\newline{}24 |\newline{}29 |\newline{}30 |\newline{}33 |\newline{}34 |\newline{}35\par

\Emph{第九卷：}\newline{}1 |\newline{}2 |\newline{}3 |\newline{}4 |\newline{}5 |\newline{}6 |\newline{}7 |\newline{}8 |\newline{}9 |\newline{}10 |\newline{}11 |\newline{}12 |\newline{}17 |\newline{}18 |\newline{}19 |\newline{}23 |\newline{}24 |\newline{}26 |\newline{}27 |\newline{}33 |\newline{}36 |\newline{}41 |\newline{}42 |\newline{}43 |\newline{}49 |\newline{}55 |\newline{}58 |\newline{}59 |\newline{}60 |\newline{}61 |\newline{}62 |\newline{}65 |\newline{}67 |\newline{}68 |\newline{}78 |\newline{}79 |\newline{}80 |\newline{}81 |\newline{}82 |\newline{}91 |\newline{}93 |\newline{}94 |\newline{}98 |\newline{}105 |\newline{}106 |\newline{}107 |\newline{}108 |\newline{}109 |\newline{}110 |\newline{}111 |\newline{}114 |\newline{}115 |\newline{}116 |\newline{}117 |\newline{}120 |\newline{}121 |\newline{}122 |\newline{}123 |\newline{}125 |\newline{}127\par

\Emph{第十卷：}\newline{}10 |\newline{}15 |\newline{}18 |\newline{}19 |\newline{}23 |\newline{}24 |\newline{}31 |\newline{}35 |\newline{}36 |\newline{}37 |\newline{}39 |\newline{}42 |\newline{}62 |\newline{}63\par

\Emph{第十一卷：}\newline{}1 |\newline{}12 |\newline{}13 |\newline{}25 |\newline{}28 |\newline{}29 |\newline{}30 |\newline{}32 |\newline{}33 |\newline{}35 |\newline{}36 |\newline{}37 |\newline{}38 |\newline{}40 |\newline{}44 |\newline{}45 |\newline{}46 |\newline{}47 |\newline{}50 |\newline{}54 |\newline{}55 |\newline{}56 |\newline{}57 |\newline{}58 |\newline{}59 |\newline{}60 |\newline{}61 |\newline{}62 |\newline{}63 |\newline{}64 |\newline{}65 |\newline{}66 |\newline{}67 |\newline{}68 |\newline{}69 |\newline{}76 |\newline{}77 |\newline{}78\par

\Emph{第十二卷：}\newline{}1 |\newline{}8 |\newline{}24 |\newline{}25 |\newline{}28 |\newline{}29 |\newline{}32 |\newline{}50\par

\Emph{第十三卷：}\newline{}1 |\newline{}5 |\newline{}6 |\newline{}7 |\newline{}8 |\newline{}9 |\newline{}10 |\newline{}12 |\newline{}18 |\newline{}22 |\newline{}26 |\newline{}27 |\newline{}31 |\newline{}34 |\newline{}38 |\newline{}39 |\newline{}40 |\newline{}41 |\newline{}42\par

\Emph{第十四卷：}\newline{}2 |\newline{}4 |\newline{}7 |\newline{}8 |\newline{}12 |\newline{}13 |\newline{}16 |\newline{}17\par

\SourceRecord{Translated by James Barmby.；Nicene and Post-Nicene Fathers, Second Series；Vols. 12-13.；Edited by Philip Schaff and Henry Wace.；Buffalo, NY: Christian Literature Publishing Co.,；1895.；<http://www.newadvent.org/fathers/3602.htm>.}

% structure: p0001 source=h1 target=section reason=page-title

\section{第一卷，第一封信}

\Emph{致西西里众主教。}\par

\Person{格列高利}，天主众仆之仆，致\Place{西西里}全境所有主教。\par

我们清楚地认识到，正如我们的前任认为宜当如此，将一切事务委托给同一人，实属非常必要；并且，在我们无法亲临之处，我们的权威应通过我们向其发送指示之人来代表。因此，藉天主的助佑，我们任命了我们的宗座副执事\Person{伯多禄}为我们在西西里省的代牧。对于他——我们已知藉天主的助佑，将我们教会全部产业的管辖委托给他——的品行，我们毫不怀疑。\par

我们也清楚地认识到，有一件应当实行之事：即每年一次，你们全体弟兄应在\Place{叙拉古}或\Place{卡塔纳}集会，按我们对他（指\Person{伯多禄}）的吩咐，接受应归于你们的尊敬；为的是，你们与前述我们的宗座副执事\Person{伯多禄}一起，得以以适当的明辨处理一切有关省区教会利益、或救济贫穷与受压迫者之急需、或劝诫众人以及处罚那些过失可能被证实之人的事务。愿敌对情绪——罪恶的养分——远离此会议，愿内心积怨消逝，愿那种极其可憎的灵魂不和消失。愿蒙天主悦纳的和谐与\OriginalTerm{爱德}{charity}，证明你们是祂的司铎。因此，要以如此的审慎与平静处理一切事务，使你们的集会最堪称为主教会议。\par

\SourceRecord{Translated by James Barmby.；Nicene and Post-Nicene Fathers, Second Series；Vol. 12.；Edited by Philip Schaff and Henry Wace.；Buffalo, NY: Christian Literature Publishing Co.,；1895.；<http://www.newadvent.org/fathers/360201001.htm>.}

% structure: p0001 source=h1 target=section reason=page-title

\section{第一卷，第二封信}

\Emph{致西西里总督犹斯底诺}\par

\Signature{额我略致西西里总督犹斯底诺}\par

我的口所言，我的良心所赞同；因为即使在你担任任何尊贵职务之前，我已深爱并深敬你。你举止的谦逊，甚至使一个本不情愿的人也对你产生了最初的爱慕之情。当我听说你前来管理西西里总督之职时，我极为欢喜。但自从我发现你与教会人士之间产生某种不睦，我便深感忧虑。如今你忙于民政管理，而我则负责教会治理，我们若彼此相爱，只要不影响公益，便可各得其宜。因此，我恳求你，因全能的天主，在他可畏的审判面前，我们必须为行为交账，愿你的荣耀常怀敬畏之心，决不容许任何事在我们之间哪怕引起轻微的不和。不要让利益引诱你偏离正义；不要让任何人的威胁或恩惠使你偏离正直的道路。看人生何其短暂：你行使司法权柄，要想一想你终有一天要面对哪位审判者。因此，当谨记：我们将所有的利益都留在此世，而有害的利益，我们只带与我们对立的控诉到审判台前。我们当追求那些死亡绝不能夺走、反而在今世终结时显为永恒的利益。\par

至于你所写关于粮食的事，尊贵的\Person{基托纳图斯}所述截然不同，他说所运送的仅是用于补充公共粮仓以清偿上一财务年所欠的数额。请关注此事，因为如果运送的粮食有任何短缺，那将不仅是某个人的死亡，而是全体百姓的死亡。\par

如今关于西西里产业的管理，我依天主的引导，派遣了这样一个人，你若如我所见的那样喜爱正义，必会与他完全同心。此外，关于你希望我善待你的请求，我如实相告：除非因古仇的诡计而生不义，我已深知你的荣耀如此谦逊，我必不耻于与你为友。\par

\SourceRecord{Translated by James Barmby.；Nicene and Post-Nicene Fathers, Second Series；Vol. 12.；Edited by Philip Schaff and Henry Wace.；Buffalo, NY: Christian Literature Publishing Co.,；1895.；<http://www.newadvent.org/fathers/360201002.htm>.}

% structure: p0001 source=h1 target=section reason=page-title

\section{卷一·第三书}

\Emph{致\Person{保禄}法学修士。}\par

\Person{额我略}致\Person{保禄}等。\par

尽管外人因我司铎职位的尊荣而向我微笑，我却不甚在意；但令我深感忧伤的是，你竟也为此而向我微笑，因为你明知我所渴慕的，却仍以为我获得了升迁。对我而言，若我之所愿得以实现——若我能完成你长久熟知的渴望，安享所渴盼的憩息——那才是最高的升迁。然而，如今我被困于罗马城，被此尊荣的锁链束缚，我甚至在致信阁下时仍有可庆幸之处：因为当最尊贵的主上前任执政官\Person{良}到来时，我猜想你不会留在西西里；而当你自己也因自身的尊荣被束缚，前来被困于罗马时，你将明了我的悲伤与苦涩。但显赫的\Person{莫伦提乌斯}主上、\OriginalTerm{机要秘书}{Chartularius}，到你那里时，我求你与他同心，应对罗马城当前的困境，因为城外我们不断被敌剑所刺。而城内因士兵的叛乱，我们更被危险重重压迫。此外，我们将在一切事上，将我们的副执事\Person{伯多禄}托付给阁下，他已受遣管理教会的产业。\par

\SourceRecord{Translated by James Barmby.；Nicene and Post-Nicene Fathers, Second Series；Vol. 12.；Edited by Philip Schaff and Henry Wace.；Buffalo, NY: Christian Literature Publishing Co.,；1895.；<http://www.newadvent.org/fathers/360201003.htm>.}

% structure: p0001 source=h1 target=section reason=page-title

\section{卷一 书信四}

\Emph{致君士坦丁堡主教若望}\par

额我略致君士坦丁堡主教若望。\par

如果爱德的德行在于爱近人，而我们也蒙命令要爱人如己，那么祢的福乐为何不爱我如同自己呢？因为我知道祢怀着何等热忱、何等焦虑，曾想逃避主教职务的重担；然而，当这同一重担被加在我身上时，祢却没有反对。显然，祢不爱我如同自己，因为祢宁愿我承担祢不愿被加于己身的重担。但我这不堪而软弱的人，掌管了一艘陈旧且严重破损的船（四面波浪涌入，船板因每日猛烈的风暴而吱嘎作响，预示着沉没），我恳求全能的天主，在我这危难中伸出祢祈祷的手援助我，因为祢远离我们在此地所受苦难的纷扰，更能热切地祈祷。\par

我的教务信函将尽快发送，因为我已在祝圣后立即派遣了我们的兄弟和同僚主教巴卡乌达作为此信的携带者，尽管有许多重要事务的催促。\par

\SourceRecord{Translated by James Barmby.；Nicene and Post-Nicene Fathers, Second Series；Vol. 12.；Edited by Philip Schaff and Henry Wace.；Buffalo, NY: Christian Literature Publishing Co.,；1895.；<http://www.newadvent.org/fathers/360201004.htm>.}

% structure: p0001 source=h1 target=section reason=page-title

\section{第一卷，第五封信}

\Emph{致皇帝之姊德奥克蒂斯塔}\par

格列高利致德奥克蒂斯塔等。\par

我的心灵以何等大的热诚匍匐于您的可敬之前，我无法完全用言语表达；但我也不费力去说出它，因为即使我沉默，您也在心中读到您对我热诚的感受。然而，我感到惊奇的是，您竟将此前赐予我的面容，从我最近担任的牧职中收回；在此牧职中，以主教之名，我被带回世俗；我被卷入如此巨大的世俗忧虑，甚至不记得在平信徒生活时曾遭受过这些。因为我失去了我安静的深层喜乐，似乎外表上升高而内里却下降。由此，我哀叹自己远离造物主的面容。因为我曾每日努力追求走出世界、走出肉身；驱散所有身体的幻象从灵魂的眼目，并无形地看见天上的喜乐；不仅用声音，而且在我心深处，我曾说：\BibleQuote{我的心曾对你说：我曾寻求你的面，你的面，主，我必寻求。}\BibleRef{咏 26:8}此外，在这世上无所欲求、无所畏惧，我似乎站在事物的某一巅峰，以至于我几乎相信通过先知所传主的应许已在我身上应验：\BibleQuote{我必将你举起，置于大地的崇高之处。}\BibleRef{依 58:14}因为那在心灵中藐视并践踏今世看似崇高荣耀的事物之人，便是被举起于大地的崇高之处。但是，因这试探的旋风突然从这巅峰摔落，我陷入恐惧和战栗，因为即使我对自己无所畏惧，我也极为害怕那些托付给我的人。四面我被俗务的波涛抛掷，被风暴淹没，以致我真可以说：\BibleQuote{我陷入深海，风暴将我淹没。}\BibleRef{咏 68:3}事务之后，我渴望回归我的心；但被虚妄思绪的喧嚣驱离，我无法回归。由此，我内在的自我远离了我，以致我无法听从先知的声音：\BibleQuote{背逆者啊，归向你的心吧。}\BibleRef{依 46:8}但被愚妄思绪所迫，我只得呼喊：\BibleQuote{我的心已弃我而去。}\BibleRef{咏 39:13}我曾爱慕默观生活之美，犹如辣黑耳，不生育却眼目清亮美丽\EditorialNote{创世纪第29章}，她虽在安静中较少结果，却更敏锐地看见光明。但不知因何审判，肋阿在夜间与我结合，即行动的生活；多产却目光柔弱；看得较少，却生育更多。我曾渴望同玛利亚坐在主脚前，聆听祂口中的言语；看哪，我却被迫同玛尔大在外务中服侍，为许多事操心忙碌\BibleRef{路 10:39}。我相信已从我身上逐出一群魔鬼，我愿忘记我所认识的人，安息于救主脚前；看哪，有人对我说话，迫使我违心而行：\BibleQuote{你回家去，宣告主为你做了何等大事。}\BibleRef{谷 5:19}但在如此多的世俗忧虑中，谁能宣讲天主奇妙的作为？对我而言，甚至回忆它们都已困难。因为在此尊位中，我被世俗事务的群集所迫，我看见自己属于那经上所说的人：\BibleQuote{当他们被高举时，你摔倒了他们。}\BibleRef{咏 72:18}祂不是说他们被高举之后你摔倒他们，而是当他们正在被高举时；因为所有恶人内里堕落，却因世俗尊荣的支撑似乎外表升高。因此他们的高举正是他们的堕落，因为他们倚靠虚假的光荣，却失去了真实的光荣。故此，祂又说：\BibleQuote{他们如烟消散，终归消灭。}\BibleRef{咏 36:20}因为烟在上升时消散，在扩展时消失。罪人生命伴随现世幸福时，确实如此，因为那显扬他的，也使他归于无有。故此，又经上记着：\BibleQuote{我的天主，求你使他们如车轮旋转。}\BibleRef{咏 82:14}因为车轮后部升起，前部落下。但我们身后的东西是今世的美物，我们遗留在后；面前的东西是永恒持久的，我们蒙召奔向它们，正如保禄见证说：\BibleQuote{忘记背后，努力面前。}\BibleRef{斐 3:13}因此，罪人今生被高举时，就如车轮，因为他在面前的事上堕落，却在身后的事上兴起。因为当他享受今生必留身后的荣耀时，他便从死后的事上坠落。确实有许多人知道如何控制外表的升高，以至于内里绝不因此堕落。经上记着：\BibleQuote{天主不弃绝有能者，因为祂自己也是大能的。}\BibleRef{约 36:5}并通过撒罗满说：\BibleQuote{明哲之人必握有权柄。}\BibleRef{箴 1:5}但对我来说，这些事是困难的，因为它们也极其沉重；心灵不乐意接受的，就不能适当掌控。看哪，我们最温和的皇帝陛下下令将猿猴变为狮子。确实，因其命令，它可以被称为狮子，但无法变成狮子。故此，他的仁爱必须亲自承担我所有过失和短缺的责任，因为他将权柄的职务交给了软弱的代理人。\par

\SourceRecord{Translated by James Barmby.；Nicene and Post-Nicene Fathers, Second Series；Vol. 12.；Edited by Philip Schaff and Henry Wace.；Buffalo, NY: Christian Literature Publishing Co.,；1895.；<http://www.newadvent.org/fathers/360201005.htm>.}

% structure: p0001 source=h1 target=section reason=page-title

\section{卷一，第六封书信}

\Emph{致贵族纳尔西斯}。\par

额我略致纳尔西斯等。\par

你崇高描述默观的甘饴，重燃了我堕落状态的哀叹，因为我听到自己在内里失去了什么，却同时在外表上，虽不配，攀上了治理的最高峰。须知我被如此巨大的忧伤打击，几乎不能说话；因为悲痛的阴暗遮蔽了我灵魂的眼目。凡所见者皆可悲，凡以为可喜者，在我看来却是可叹的。因为我反思，从安息的崇高巅峰坠落时，我攀登到了何等沮丧的世俗擢升高度。并且，由于我的过犯，我被从我主面前放逐到职务的流放中，我与先知一同，仿佛用被毁的耶路撒冷的话说：\BibleQuote{\Emph{安慰我的人远离了我}} \BibleRef{哀 1:16}。然而，当你寻求一个比喻来表达我的处境和名号，在信函中编织辞藻和慷慨陈词时，当然，最亲爱的弟兄，你是在称猿猴为狮子。在这里，我们看到你所做的正如我们常做的，即称长癣的幼崽为豹或老虎。因为我，我的好人，就如同失去了我的子女，因为我因世俗的挂虑而失去了义德之工。因此，\Quote{\Emph{不要叫我纳敖米，即美好的，而要叫我玛辣，因为我充满苦楚}} \BibleRef{卢 1:20}。至于你说我不该写信说\Quote{你该用\OriginalTerm{野牛}{bubali}在主田里耕耘}，因为在那显示给有福伯多禄的布帛中，\OriginalTerm{野牛}{bubali}与一切野兽都呈现于眼前；你自知附加一句：\BibleQuote{\Emph{宰杀吃吧}} \BibleRef{宗 10:13}。那么，你既尚未宰杀这些野兽，为何已因服从而想吃它们？或者，你不知道你所写的这头兽拒绝被你口中之剑宰杀吗？那么，你必须以那些你能刺伤并杀死（原文：因痛悔而杀死）之人，来满足你渴望的饥饿。\par

此外，关于我们弟兄的情况，我想，若天主助佑，将会如你所写。然而，我现在根本不该就此写信给我们最安详的主上，因为一开始不应以抱怨开头。但我已写信给我心爱的儿子、执事奥诺拉托，叫他在适当时机以恰当方式提及此事，并迅速告知我他们的答复。我请求代我向亚历山大勋爵、德奥多禄勋爵、我的儿子玛里诺、埃西西亚夫人、欧多基娅夫人和多米尼卡夫人致以问候。\par

\SourceRecord{Translated by James Barmby.；Nicene and Post-Nicene Fathers, Second Series；Vol. 12.；Edited by Philip Schaff and Henry Wace.；Buffalo, NY: Christian Literature Publishing Co.,；1895.；<http://www.newadvent.org/fathers/360201006.htm>.}

% structure: p0001 source=h1 target=section reason=page-title

\section{卷I，书信7}

\Emph{致安提约基雅宗主教阿纳斯塔西乌斯}。\par

额我略致阿纳斯塔西乌斯等。\par

我发现您的圣座所写的，对疲倦者是休息，对病者是健康，对渴者是泉源，对受暑者是荫凉。因为您的那些话语似乎甚至不是由血肉之舌表达的，因为您如此揭示了您对我怀有的属灵之爱，仿佛您的灵魂本身在说话。但非常艰难的是随后之事，因为您的爱命令我承担世俗的负担，并且，首先以属灵的方式爱了我之后，后来，如我所想，以现世的方式爱了我，您以您加给我的负担将我压倒在地；以至于我完全失去灵魂的正直，丧失默观的敏锐视力，可以这么说，不是出于预言的精神，而是从经验出发：\BibleQuote{我全然仆倒，极其卑微}\BibleRef{咏 118:107}。因为确实如此重大的事务负担压迫我，使我的心灵无法提升到天上的事物。我被诸多事务的波涛抛掷，在从前安闲的安逸之后，受动荡生活的风暴折磨，因此我可以真诚地说：\BibleQuote{我陷入深海，波涛淹没了我}\BibleRef{咏 68:3}。所以，您这站在德行岸上的人，请向我这处于危险中的人伸出您祈祷的手。至于您称我为主的嘴和灯，并声称我使许多人受益，这也增加了我的罪孽的负担，因为我的罪孽本应受惩罚，我却得到赞美而非惩罚。但我在这里被世俗事务的喧嚣所困扰，无法用言语表达；然而您可以从这封信的简短中看出，我对这位我超越所有人所爱之人说得如此之少。此外，我告知您，我已以最大努力恳求我们最安详的主上们，准许您在恢复您的尊严后，来到宗徒之长伯多禄的门限前，并在此与我同住，只要天主乐意；如此，在我有生之年堪当见到您时，我们可以通过谈论天上家乡来减轻我们旅途的疲惫。\par

\SourceRecord{Translated by James Barmby.；Nicene and Post-Nicene Fathers, Second Series；Vol. 12.；Edited by Philip Schaff and Henry Wace.；Buffalo, NY: Christian Literature Publishing Co.,；1895.；<http://www.newadvent.org/fathers/360201007.htm>.}

% structure: p0001 source=h1 target=section reason=page-title

\section{第一卷，第九封信}

\Emph{致副执事\Person{伯多禄}。}\par

\Person{额我略}致\Person{伯多禄}等。\par

\Person{额我略}，天主的仆人，司铎暨隐修院院长，该院位于西西里省\Place{帕诺尔莫斯}境内的圣\Person{德奥多禄}隐修院，告知我等：属于圣罗马教会的\Place{富洛尼阿库斯}农庄的人正试图侵犯与上述圣罗马教会农庄相邻的\Place{格尔迪尼亚}农庄的边界，而他们（\EditorialNote{即圣德奥多禄的隐修士们}）已无争议地拥有该农庄无数年。为此，我等希望你前往\Place{帕诺尔莫斯}城，以如下方式调查此事（以维护此前占有者的权利），即：如果你发现上述圣\Person{德奥多禄}隐修院无扰地拥有争议边界已达四十年，则不得允许其遭受任何损害，即使有利于圣罗马教会，也应以一切方式保障其不受干扰。但是，如果圣罗马教会的代理人能够证明该隐修院并非无争议地拥有其权利达四十年，而是在此期间曾就所述边界提出过任何问题，则应通过为此选定的仲裁人和平且合法地予以解决。因为我们不仅希望提出从未被争论的不当行为问题，也希望能迅速解决由他人而非我们提出的问题。因此，请你的经验使得一切调整得当，以致日后不再有与此事相关的问题提交给我们。另外，我们希望已故的招待所主管（\OriginalTerm{招待所主管}{Xenodochus}）\Person{巴卡乌达}的遗嘱，继续如同最初订立时那样有效。\par

十一月：第九\OriginalTerm{小纪}{Indiction}。\par

\SourceRecord{Translated by James Barmby.；Nicene and Post-Nicene Fathers, Second Series；Vol. 12.；Edited by Philip Schaff and Henry Wace.；Buffalo, NY: Christian Literature Publishing Co.,；1895.；<http://www.newadvent.org/fathers/360201009.htm>.}

% structure: p0001 source=h1 target=section reason=page-title

\section{第一卷，第十封}

\Emph{致\Person{巴高达}与\Person{阿格内鲁斯}主教}\par

\Signature{格里高利致\Person{巴高达}等}\par

住在\Place{泰拉奇纳}的希伯来人向我们请求许可，得以在我们的权下，持有他们迄今为止所拥有的会堂场地。但是，因为我们已获悉，同样的场地离教堂如此之近，以至于他们圣咏的声音都能传到教堂，我们已写信给我们同为主教的弟兄\Person{伯多禄}，告知他：如果情况是，从该处传来的声音在教堂中被听到，那么犹太人就必须停止在那里敬拜。因此，请阁下与上述我们的弟兄及同为主教者一起仔细视察该处，如果发现有任何对教堂的干扰，请在堡垒内提供另一个场所，使前述希伯来人可以聚集，以便他们能够无阻碍地举行他们的礼仪。但请阁下提供这样一个场所，以防他们被剥夺这一个，以致将来没有抱怨的理由。但我们禁止前述希伯来人受无理压迫或困扰；既然根据公义，他们获准在罗马法律的保护下生活，就让他们遵守他们所学的规诫，无人阻碍他们；然而不得允许他们拥有基督徒奴隶。\par

\SourceRecord{Translated by James Barmby.；Nicene and Post-Nicene Fathers, Second Series；Vol. 12.；Edited by Philip Schaff and Henry Wace.；Buffalo, NY: Christian Literature Publishing Co.,；1895.；<http://www.newadvent.org/fathers/360201010.htm>.}

% structure: p0001 source=h1 target=section reason=page-title

\section{卷一 书信十一}

\Emph{致贵族\Person{克莱门蒂娜}}。\par

\Person{额我略}致\Person{克莱门蒂娜}等。\par

接奉阁下惠函，提及已故、值得怀念的\Person{欧特里乌斯}之去世，兹告知：此一悲伤所扰动吾人之心，不亚于阁下。盖见备受赞誉之人渐次离世，而世界之倾颓已显于其因之果效。然吾人理当以皈依之明智慎虑，自世界抽身，免致同陷于其倾覆。诚然，丧失亲友之悲伤，当更可忍受，因吾人必死之处境本即需我们失去他们。尽管如此，对于肉体生活之援助的丧失，那位准许其离去的天主，有能力安慰，并亲自来填补空缺，成为安慰者。\par

吾人未能应允阁下请求，派遣执事\Person{阿纳托利乌斯}前往，乃因事理使然，非出于严格苛刻。盖吾人已任命其为管家，将主教府交其管理。\par

\SourceRecord{Translated by James Barmby.；Nicene and Post-Nicene Fathers, Second Series；Vol. 12.；Edited by Philip Schaff and Henry Wace.；Buffalo, NY: Christian Literature Publishing Co.,；1895.；<http://www.newadvent.org/fathers/360201011.htm>.}

% structure: p0001 source=h1 target=section reason=page-title

\section{第一卷，第十二封书信}

\Emph{致\Place{旧城（奥尔维耶托）}主教\Person{若望}}\par

\Person{额我略}致\Person{若望}等。\par

\Place{圣乔治隐修院}院长\Person{阿加皮托}告诉我们，他承受着来自阁下的许多委屈；不仅是在急需时可能对隐修院有所帮助的事情上，而且您甚至禁止在该隐修院举行弥撒，并阻止在那里埋葬死者。现在，如果情况如此，我们劝您停止这种不人道的行为，允许埋葬死者并在那里举行弥撒，不得再有反对，以免上述可敬的\Person{阿加皮托}不得不就所涉事项再次申诉。\par

\SourceRecord{Translated by James Barmby.；Nicene and Post-Nicene Fathers, Second Series；Vol. 12.；Edited by Philip Schaff and Henry Wace.；Buffalo, NY: Christian Literature Publishing Co.,；1895.；<http://www.newadvent.org/fathers/360201012.htm>.}

% structure: p0001 source=h1 target=section reason=page-title

\section{卷一，第十六封信}

\Emph{致\Place{阿奎莱亚}主教\Person{塞味禄}}\par

\Person{额我略}致\Person{塞味禄}等\par

当一个人行于歧途又重新走上正路时，主以全心拥抱他；那么后来，当这同一人离弃真理之路时，主为他忧伤，比当初他转离错谬时主因他而欢喜更为深切；因为不认识真理，其罪小于认识后不持守真理；出于错误的作为是一回事，明知故犯又是另一回事。我们以前曾因你融入了教会的合一而欢喜，现在更因你脱离了公教会团体而深感悲痛。因此，我们希望你，在持信者的敦促下，遵照最虔诚、最安详的皇帝之命，与你的追随者一同来到蒙福的宗徒\Person{伯多禄}的座前，以便藉天主的旨意召集一次主教会议，对你们中间存在的疑问作出判断。\par

\SourceRecord{Translated by James Barmby.；Nicene and Post-Nicene Fathers, Second Series；Vol. 12.；Edited by Philip Schaff and Henry Wace.；Buffalo, NY: Christian Literature Publishing Co.,；1895.；<http://www.newadvent.org/fathers/360201016.htm>.}

% structure: p0001 source=h1 target=section reason=page-title

\section{第一卷 第十七封信}

\Emph{致意大利全体主教。}\par

\Signature{格里高利致全体，等}\par

由于可憎的\Person{奥塔里特}在刚结束的复活节庆期中，禁止伦巴第人的子女在天主教信仰中领洗，为此罪恶，天主威严将他剪除，使他不得再见另一个复活节的庆期。因此，各位弟兄应当警告你们辖区内的所有伦巴第人：既然严重的死亡随处临近，他们应使这些已在亚略异端中领洗的子女与天主教信仰和好，从而平息悬在他们头上的全能上主的义怒。那么，警告你们所能及的人；用你们所有劝说的力量感化他们，引导他们归向正信；向他们宣讲无尽的永生；为使当你们来到严厉审判者面前时，因你们的关怀，能在自己身上显示牧者的收获。\par

\SourceRecord{Translated by James Barmby.；Nicene and Post-Nicene Fathers, Second Series；Vol. 12.；Edited by Philip Schaff and Henry Wace.；Buffalo, NY: Christian Literature Publishing Co.,；1895.；<http://www.newadvent.org/fathers/360201017.htm>.}

% structure: p0001 source=h1 target=section reason=page-title

\section{第一卷，第十八封书信}

\Emph{致\OriginalTerm{副执事}{Subdeacon} \Person{伯多禄}}\par

\Signature{额我略致伯多禄，等}\par

我们获悉，巴鲁塔尼亚教会之玛策禄，在帕诺尔莫斯城同一城中之圣亚德连隐修院内被指定补赎，不仅缺乏食物，且因衣物不足遭受不便。因此，我们认为有必要以此命令嘱尔阁下：为尔所见所需之食物、衣物及寝具，供其自身维持及仆役之需，尽行拨给；如此，其匮乏与赤身得以及时照料，尔所拨付于此人者，日后将归入尔之账目。故尔当如此行事：既成全我命，又因妥善安排此事，尔自身亦能分享其中利益。此外，尚有另一事嘱尔留意，不顾现已形成之旧习：即若西西里省中任何城市因其罪恶，因司铎失职而致无牧灵管理，尔当查看，在本教会圣职人员中，或在隐修院之外，有无堪当司铎职任者，先查其品行庄重，然后送交我等，以免各地方之羊群因牧者失职而长期困乏。但若尔发现任何空缺之地，在该教会中无一人适合此等尊位，则同样仔细查问后告知我等，以便提供天主视为堪当如此祝圣之人。因为一人之偏离，主的羊群在无牧人时处于悬崖间游荡之险，实为不当。如此，各地方之管理得以延续，亦不再有对失职者恢复原职之疑虑，且他们可更好悔改。\par

\SourceRecord{Translated by James Barmby.；Nicene and Post-Nicene Fathers, Second Series；Vol. 12.；Edited by Philip Schaff and Henry Wace.；Buffalo, NY: Christian Literature Publishing Co.,；1895.；<http://www.newadvent.org/fathers/360201018.htm>.}

% structure: p0001 source=h1 target=section reason=page-title

\section{第一卷，书信十九}

\Emph{致\Place{萨洛纳}主教\Person{纳塔利斯}}。\par

\Person{格里高利}致\Person{纳塔利斯}等。\par

您所寄来的教务会议纪要，其中总执事\Person{霍诺拉图斯}被定罪，我们看出其中充满纷争的种子，因为同一个人在同一时间既违背其意愿被擢升至司铎尊位，又被视作不配而免除执事职务。正如无人应被强迫接受不受欢迎的晋升是公正的，同样，我们认为，无人应无辜地被不公正地免去其圣职之服务。然而，既然可憎于天主的纷争为您在这一事件中的立场作了辩护，我们劝您恢复总执事\Person{霍诺拉图斯}的职位和职权，并同意为其神圣职务提供足够的辅助人员。如果你们之间仍存有嫌隙，请让上述总执事在接到提醒后，接受我们的听讼和调查；同时，请您派遣一位熟悉案情的人前来，以便在主助佑下，我们能在双方在场时，不徇情面地做出正义所认可的判决。\par

\SourceRecord{Translated by James Barmby.；Nicene and Post-Nicene Fathers, Second Series；Vol. 12.；Edited by Philip Schaff and Henry Wace.；Buffalo, NY: Christian Literature Publishing Co.,；1895.；<http://www.newadvent.org/fathers/360201019.htm>.}

% structure: p0001 source=h1 target=section reason=page-title

\section{第一卷，第二十封信}

\Emph{致\Place{萨洛纳}的执事\Person{奥诺拉图斯}}。\par

\Person{格里高利}致\Person{奥诺拉图斯}等。\par

我们读了您与您的主教彼此寄给我们的互相矛盾的信函，对你们之间如此缺乏爱德感到痛心。尽管如此，我们仍命令您继续履行您的职务；如果你们之间的纷争能在恩宠的能力下就地解决，我们相信这将对你们的灵魂大为有益。但若你们之间的不和已使您武装起来彼此对抗，以至您无意平息你们的不和之膨胀，那么就请毫不迟延地前来听取我们的裁决，并让您的主教为自己选派一位他中意的人，携带指示前来；这样，在仔细审查整个案件后，我们将裁定双方认为合适的事。但我们要让您知道，我们将就所有细节严格询问您，关于圣器——无论是您自己教堂的，还是从各教堂收集来的——是否正被以一切谨慎和忠诚保管。因为，若发现其中任何一件因疏忽或任何人的不诚实而丢失，您将因此罪责受牵连，因为您身为总执事，对该教堂的保管负有特别责任。\par

\SourceRecord{Translated by James Barmby.；Nicene and Post-Nicene Fathers, Second Series；Vol. 12.；Edited by Philip Schaff and Henry Wace.；Buffalo, NY: Christian Literature Publishing Co.,；1895.；<http://www.newadvent.org/fathers/360201020.htm>.}

% structure: p0001 source=h1 target=section reason=page-title

\section{第一书，第二十一函}

\Emph{致\Person{纳塔利斯}，\Place{撒罗纳}主教。}\par

\Person{额我略}致\Person{纳塔利斯}等。\par

我们已从你派来的执事\Person{斯德望}手中收到你尊座的书信，你在信中祝贺我们晋升。诚然，凡出于爱德的善意与热诚所奉献的，完全值得信任，因为理智促使你的主教尊位与我们一同喜乐。因此，我们为你的问候所鼓舞，凭良心宣告：我曾以一颗忧苦的心承担尊位的重担。然而，眼看我无法抗拒天主的谕令，我恢复了更为愉悦的心境。为此，我们写信恳求你尊座，使我们以及托付给我们照管的基督信徒的羊群，都能得到你祈祷的援助，以便在那保护的安全中，我们能有力量克服这个时代的飓风。\par

二月，\OriginalTerm{第九小纪}{ninth indiction}。\par

\SourceRecord{Translated by James Barmby.；Nicene and Post-Nicene Fathers, Second Series；Vol. 12.；Edited by Philip Schaff and Henry Wace.；Buffalo, NY: Christian Literature Publishing Co.,；1895.；<http://www.newadvent.org/fathers/360201021.htm>.}

% structure: p0001 source=h1 target=section reason=page-title

\section{第一卷，第二十五封书信}

《致君士坦丁堡主教若望及其他宗主教》\par

额我略致君士坦丁堡的若望、亚历山大的厄罗吉乌斯、安提约基雅的额我略、耶路撒冷的若望，以及安提约基雅的前宗主教阿纳斯塔西。 \OriginalTerm{同辈}{A paribus}。\par

当我想到，不配的我，虽全心抗拒，仍被迫背负牧职的重担时，悲伤的黑暗笼罩着我，我忧伤的心只见阴影，不见万物。因为主教被主拣选，除了作百姓的转求者外，还有何目的？那么，我若对自己所犯之罪在天主前尚无把握，又怎能大胆地作为他人罪过的转求者来到祂面前？若有人请我为他向一位对他发怒的伟人求情，而我又不认识那人，我会立刻回答：我不能为你转求，因为我与那人并无深交。那么，若人与人之间，我理当羞于为自己不熟悉的人转求，那么我为百姓在天主前获得转求者的地位，该是何等狂妄？因我生活的功绩并不能保证祂的友谊。而且，在这件事上，我发现了更严重的威胁：因为我们都知道，当一个不得宠的人被派去向一位发怒者求情时，后者的心会被激怒得更加严厉。我深恐信徒团体——主迄今容忍了他们的过犯——会因我的罪过加于他们而遭毁灭。但是，当我以某种方式压抑这恐惧，并以安慰的心将自己投入主教职务的关怀时，又因这职务的浩瀚而却步。\par

\Quote{因为我确实思量，需要何等警醒的关怀，才能使牧者思想纯洁、行为卓越、静默有节、言语有益、对众人同情相近、在默观中超越一切、因谦卑与善行者交友、因正义的热忱对恶者的恶行毫不妥协。} 当我以精细的探究试图查考这一切时，这考量的广阔本身反在细节上束缚了我。因为，如前所述，需要极大的关怀，以使\Quote{牧者思想纯洁、行为卓越，等等。} \EditorialNote{此处有一段长文，以此起头，以\Quote{超出秩序的范围}结尾，亦见于\WorkTitle{牧灵规则}卷二第二章。}\par

再者，当我思考牧者所需的工作时，我内心衡量需要何等精心的关怀，使他\Quote{行为卓越，好藉其生活为属下指明天路，等等。} \EditorialNote{见\WorkTitle{牧灵规则}卷二第三章至结尾。}\par

再者，当我思考牧者在言语和静默上的职责时，我战战兢兢地衡量多么必要使他\Quote{静默有节、言语有益，免得他或说出当隐藏的，或隐藏当说的，等等。} \EditorialNote{见\WorkTitle{牧灵规则}卷三第四章，直至\Quote{保持信仰的统一}。}\par

再者，当我思考牧者在同情与默观上当如何时，我内心衡量他\Quote{对众人同情相近，在默观中超越一切，好藉仁爱的怀抱，等等。} \EditorialNote{见\WorkTitle{牧灵规则}卷二第五章至结尾。}\par

再者，当我思考牧者在谦卑与严格上当如何时，我内心衡量多么必要使他\Quote{藉谦卑与善行者交友，藉正义的热忱对恶者的恶行毫不妥协，等等。} \EditorialNote{见\WorkTitle{牧灵规则}卷二第六章，直至\Quote{对待悖逆者}；结尾的措辞仅有细微变化，不影响含义。}因此，\Quote{伯多禄既从天主领受了，等等。} \EditorialNote{见\WorkTitle{牧灵规则}卷二第六章，直至\Quote{统治恶习而非统治弟兄}。}他良好地治理所领受的权柄，既懂得维护它，又懂得节制它。他良好地治理它，既藉它超越罪恶，又藉它使自己与他人平等。\par

再者，应如此持守谦卑之德，使治理之权不被松弛；免得一位主教过于降低自己，以致不能以纪律的锁链约束属下的生活；纪律的严厉也应持守，免得因热忱过烈而完全丧失温和。因为恶习常伪装成美德，以致吝啬愿显为节俭，挥霍显为慷慨，残忍显为正义的热忱，懈怠显为仁爱。因此，若两者不兼而行之，纪律与仁慈都远离其应有之道。反而应以极大的分辨技巧，保持公正审慎的仁慈与仁爱施罚的纪律。\Quote{因此，正如真理所教导的\BibleRef{路 10:34}，那人因撒玛黎雅人的照顾被带走，等等。} \EditorialNote{见\WorkTitle{牧灵规则}卷二第六章，直至\Quote{甜蜜的玛纳}。}\par

如此，既承担了牧职的重担，当我思考这一切以及许多类似的事时，我似乎成了我所不能成的人，尤其在此地，凡被称为牧者的人，都为外在的事务所重压；以致常不确定他是行使牧者之职还是属世贵人之职。的确，凡被置于弟兄之上治理他们者，不能完全免于外在事务；但仍需极大小心，免得被它们过分压垮。\Quote{因此，正确地告诉厄则克耳：司祭不可剃头，等等。} \EditorialNote{见\WorkTitle{牧灵规则}卷二第七章至结尾。}\par

但在此时，我看到无法采取如此审慎的管理方式，因为每日压在我身上的此类重大事务不仅耗尽心神，也摧残肉身生命。因此，至圣的弟兄，我恳求那将要来临的审判者、成千上万天使的集会、以及那些名字被写在诸天之上的首生者教会，求你以祈祷的代祷帮助在这牧职重担下疲惫的我，免得这重负超过我的力量压垮我。然而，我既记得经上所载：\BibleQuote{你们要彼此祈祷，为得痊愈}（\BibleRef{雅 5:16}），我也给予我所求的。但我将得到我所给予的。因为当我们借着祈祷的帮助与你相连时，仿佛在滑地上行走时彼此搀扶，由于爱德这伟大的安排，每个人的脚都会因靠着他人而更加稳固。\par

此外，既然心里相信能成义，口里承认能得救，我承认我领受并尊崇四福音书如同四部大公会议：即\Place{尼西亚}会议，其中\Person{亚略}的邪说被推翻；\Place{君士坦丁堡}会议，其中\Person{欧诺米}和\Person{马其敦}的谬误被驳斥；还有第一次\Place{厄弗所}会议，其中\Person{内斯多略}的不敬被谴责；以及\Place{加采东}会议，其中\Person{欧提基}和\Person{迪奥斯郭尔}的曲解被弃绝。我以全心拥护并极赞同地持守这些会议；因为圣教信仰的建筑如同立于四方之石上；凡不持守其坚固性的人，无论其生活行为如何，纵然看起来像石头，却仍在建筑之外。我也同样尊敬第五次大公会议，其中充满错误的所谓\Person{依巴斯}书信被弃绝；将天主与人之间的中保（\OriginalTerm{中保}{Mediator}）分为两个位格（\OriginalTerm{位格}{subsistences}）的\Person{戴奥道}被判定已陷入不敬的背信；而\Person{德奥道肋}的著作中攻击真福\Person{济利禄}信德的部分，被驳斥为狂妄地发表。凡上述可敬之会议所弃绝的人，我都弃绝；凡他们所尊崇的，我皆拥护；因为这些会议由普遍同意而建立，任何人若胆敢释放他们所束缚的，或束缚他们所释放的，他不是推翻会议，而是推翻自己。因此，凡持不同意见者，愿他受弃绝。但凡持守上述会议信德的人，愿来自天主圣父的平安，借着祂的圣子耶稣基督，与祂同体（\OriginalTerm{同体}{consubstantially}）为天主，在圣神的合一中生活并统治，至于无穷之世。阿们。\par

\SourceRecord{Translated by James Barmby.；Nicene and Post-Nicene Fathers, Second Series；Vol. 12.；Edited by Philip Schaff and Henry Wace.；Buffalo, NY: Christian Literature Publishing Co.,；1895.；<http://www.newadvent.org/fathers/360201025.htm>.}

% structure: p0001 source=h1 target=section reason=page-title

\section{第一卷，第二十六封信}

\Emph{致安提约基雅宗主教阿纳斯塔西乌斯}\par

\EditorialNote{本信开头与致同一阿纳斯塔西乌斯的第七封信相同，直至\Quote{站在德行之岸}一语；之后则接续如下。}\par

至于你称我为主的口舌与明灯，并断言我以言语裨益多人，且能光照多人，我承认你使我对自己的评估陷入极大疑惑。因为我省思自己，发现自己身上毫无这些美善的迹象；但我也省思你是何人，且不信你能说谎。于是，当我欲相信你所说时，我的软弱反驳我；当我欲争辩那些赞美之词时，你的圣德反驳我。但我恳求你，圣者，让我们在这场争论中达成某种一致：纵然事实并非如你所说，但因你所说，它或许可以如此。此外，我已将我的主教会议信函致予你，如同致予其他宗主教、你的弟兄一般；因为在我看来，你始终是全能天主恩赐所使你成为的，而不拘于人愿所认为你并非的那一面。我已将一些指示交付予持此函的\OriginalTerm{护卫官}{defensori}\Person{博尼法爵}，命他私下禀告尊驾。此外，我已将蒙福的宗徒伯多禄的钥匙送予你，他爱你；这些钥匙置于病人身上时，常藉许多奇迹发出光辉。\par

\SourceRecord{Translated by James Barmby.；Nicene and Post-Nicene Fathers, Second Series；Vol. 12.；Edited by Philip Schaff and Henry Wace.；Buffalo, NY: Christian Literature Publishing Co.,；1895.；<http://www.newadvent.org/fathers/360201026.htm>.}

% structure: p0001 source=h1 target=section reason=page-title

\section{\WorkTitle{第一辑，第廿七封书信}}

\Emph{致\Place{科林特}总主教\Person{亚纳斯大削}}\par

\Signature{\Person{额我略}致\Person{亚纳斯大削}等}\par

天主的判断何等难测，正因如此，当成为人类敬畏的对象；好使有限理性无法领悟之时，必然谦卑地俯首顺从，以心灵顺服的脚步跟随统治者的旨意。我自知软弱不能企及宗座之高度，宁愿推卸此重担，以免在牧职统治中因治理不当而失败。然而，既然我们不可违逆安排万有的主之意愿，我便顺从地跟随统治者的仁慈之手愿与我交往的方式。因为即使当前机会未曾出现，也必须使你的弟兄得知，主竟俯允我——尽管不配——主持宗座。既然理当如此，且通过我们派遣此信的携带者——即\Person{博尼法爵}\OriginalTerm{护卫者}{defensorem}——的机会出现，我们不但要谨以书信向你的弟兄献上爱德之问候，也要告知你我们的祝圣，如我们所信你愿我们如此行。因此，请你的爱德以回信使教会合一及你自身平安的佳音令我们欣喜；好使因地域相隔而承受的肉身分离，通过书信往来如同同在。我们也劝勉你，既然我们派遣上述携带此信者因某些必要事务前往最仁慈君王的脚下，且时局多变常在途中生出许多阻碍，愿你的司铎之情施予他所需的一切，无论是陆路行程的供应还是航行工具的获取，以便靠天主的仁慈，他能更快完成其预定行程。\par

\SourceRecord{Translated by James Barmby.；Nicene and Post-Nicene Fathers, Second Series；Vol. 12.；Edited by Philip Schaff and Henry Wace.；Buffalo, NY: Christian Literature Publishing Co.,；1895.；<http://www.newadvent.org/fathers/360201027.htm>.}

% structure: p0001 source=h1 target=section reason=page-title

\section{第一卷，第二十八封信}

\Emph{致里西农（在达尔马提亚）主教塞巴斯蒂安。}\par

\Signature{格列高利致塞巴斯蒂安等处。}\par

尽管我理应未收到尊驾的任何书信，然而我亦不忘自己的疏忽；我责备自己的懈怠，以爱的激励鞭策自己的怠惰，使那不愿自愿偿还所欠之人，甚至能在打击下学会偿还。此外，我通知您，我已准备好一份完整的陈情书，以恳切的祈祷呈于我们至虔敬的主上们，即他们应当派遣至福的安纳斯塔修斯宗主教，使用赐予他的\OriginalTerm{披肩}{pallium}，前往真福宗徒之长伯多禄的\OriginalTerm{门槛}{threshold}，与我一同举行\OriginalTerm{弥撒}{Mass}的隆重礼仪；以致尽管他未被允许返回自己的宗座，至少能在保留其尊位的情况下与我同住。但对于为何搁置我如此写就的内容，这封信的携带者将告知您。不过，请探明上述安纳斯塔修斯宗主教的心意，并以书信告知我他对此事希望如何处理。\par

\SourceRecord{Translated by James Barmby.；Nicene and Post-Nicene Fathers, Second Series；Vol. 12.；Edited by Philip Schaff and Henry Wace.；Buffalo, NY: Christian Literature Publishing Co.,；1895.；<http://www.newadvent.org/fathers/360201028.htm>.}

% structure: p0001 source=h1 target=section reason=page-title

\section{第一卷 第29封书信}

\Emph{致前长官和书记官亚里斯托布鲁}。\par

\Person{额我略}致\Person{亚里斯托布鲁}等。\par

我承认，我的口舌不足以充分表达我的情谊；但你自己的情谊更能向你述说我对你的一切感受。我听说你正遭受某些反对。然而，我并不为此感到十分忧伤，因为常会发生这样的情况：一艘本可凭借顺风抵达大洋深处的船，却在航程之初被逆风吹回，但正因为被吹回，才得以返回港口。此外，如果你偶然接到我的一封长信需要诠释，我求你，不要逐字翻译，而要译出大意；因为通常，当过分拘泥于字面时，思想的力量就会丧失。\par

\SourceRecord{Translated by James Barmby.；Nicene and Post-Nicene Fathers, Second Series；Vol. 12.；Edited by Philip Schaff and Henry Wace.；Buffalo, NY: Christian Literature Publishing Co.,；1895.；<http://www.newadvent.org/fathers/360201029.htm>.}

% structure: p0001 source=h1 target=section reason=page-title

\section{\WorkTitle{第一卷，第三十三封书信}}

\Emph{致意大利总督兼贵族\Person{Romanus}}\par

\Person{Gregory}致\Person{Romanus}等。\par

即使没有即刻的理由致书阁下，我们仍当关心您的健康与平安，以便通过频繁的互通音讯得知我们所愿听到的关于您的情况。此外，我们获悉\Place{Hortanum}城的主教\Person{Blandus}被阁下长期扣留在\Place{Ravenna}城。其结果是教会因无治理者而衰败，人民如无牧者；那里的婴儿因他们的罪未能领受圣洗而死去。再者，既然我们不相信阁下扣留他除了因某种可能的过犯外另有原因，那么理当召开一次会议，以揭露指控他的任何罪行。若在他身上发现足以导致他丧失司铎职位的过错，我们就有必要物色另一个人来接受祝圣，以免天主的教会无人照管，并且在那基督信仰不允许缺乏的事上陷于匮乏。但若阁下察觉他的情况与所述不符，恳请您准许他返回自己的教会，使他能对托付给他的灵魂尽忠职守。\par

三月；第九\OriginalTerm{小纪}{Indiction}。\par

\SourceRecord{Translated by James Barmby.；Nicene and Post-Nicene Fathers, Second Series；Vol. 12.；Edited by Philip Schaff and Henry Wace.；Buffalo, NY: Christian Literature Publishing Co.,；1895.；<http://www.newadvent.org/fathers/360201033.htm>.}

% structure: p0001 source=h1 target=section reason=page-title

\section{第一册，第三十四封书函}

\Emph{致\Person{文兰提乌斯}，前隐修士，\Place{叙拉古}贵族。}\par

\Person{额我略}致\Person{文兰提乌斯}等：\par

许多愚昧人曾以为，若我升迁至主教之职，便会拒绝致书于你，或与你保持书信往来。但事实并非如此；因为我职位自身的必要性迫使我不能缄默。因为经上记载：\BibleQuote{\Emph{你要大声喊叫，不可止息，扬起你的声音如号角}} \BibleRef{依 58:1}。又经上记载：\BibleQuote{\Emph{我立你作以色列家族的哨兵，你听了我口中的话，要代我警告他们}} \BibleRef{则 3:17}。随后又立即指出，若哨兵或听者因这宣告被保留或发出而有什么后果：\BibleQuote{\Emph{若我说恶人必死，你不警告他，也不对他讲话，使他离开恶道而生活，那恶人必因自己的罪而死；但我要由你手中追讨他的血债。但若你警告了恶人，他仍不离开他的邪恶和恶道，他固然要因自己的罪而死，但你却救了自己的灵魂。}}因此保禄也对厄弗所人说：\BibleQuote{\Emph{我今天手洁，不负你们任何人的血债。因为我从未回避将天主的全部计划传报给你们}} \BibleRef{宗 20:26}。那么，若他拒绝将天主的计划传报给他们，他便不会手洁于众人的血。因为当牧者拒绝责备犯罪者时，无疑他因缄默而杀害了他们。因此，被此事所迫，不论你愿意与否，我都要说话；因为我竭尽全力，要么你得以得救，要么我从你的死亡中获救。因为你记得你从前生活在何种状态，也知道你如今落到了何等境地，竟无视上主严正的责备。那么，趁还有时间，反省你的过错；趁你尚能畏惧将来审判者的严厉，现在就畏惧吧，免得那时你因为没有流泪避免它而觉得苦涩。思考经上所记载的：\BibleQuote{\Emph{你们要祈祷，不要叫你们的逃遁遇到冬天或安息日}} \BibleRef{玛 24:20}。因为冬天的严寒阻碍行路，而按法律的规定，安息日也不许行路。那么，企图在冬天或安息日逃跑的人，是想在不能行走的时候逃离严厉审判者的愤怒。所以，趁还有时间，趁还允许，逃离这极其可怕的责备：思考经上所记载的：\BibleQuote{\Emph{你手能做的事，要尽力去做，因为在你所要去的阴府里，没有工作，没有智谋，没有学问，没有智慧}} \BibleRef{训 9:10}。凭着福音的见证，你知道神圣的严重要求我们为闲谈负责，并为无益的话要求严格的交代。\BibleRef{玛 12:36}那么，想一想，如果它在审判中因一些人的闲谈而定罪，它会对邪恶的行为做什么呢？\Person{阿纳尼雅}曾许愿将钱献给天主\BibleRef{宗 5:2}，但后来受魔鬼的劝说而私自留下了。你知道他因此受了什么样的死亡惩罚。那么，如果那撤回已献给天主的钱的人应受死刑的惩罚，想想你在天主的审判中应受多么大的惩罚——你撤回了你自己，而不是钱，从全能的天主那里，你曾在隐修状态中把自己奉献给祂。所以，你若听从我的责备之言并遵行，最终你会明白这些话有多么和善甘甜。看，我承认，我说话是带着悲伤，因你的所作所为而忧心忡忡。我几乎说不出话来；但你那自知有愧的心，却几乎不能忍受它所听到的：脸红、窘迫、抗议。那么，若它不能忍受尘埃之言，在造物主的审判面前它又能做什么呢？然而，我承认天上恩宠的无限仁慈，因为它见你逃避生命，却仍然为你保存生命；它见你行为骄傲，却仍容忍你；它通过自己无用的仆人，向你传递责备和劝诫的话。此事如此重大，你当殷切思考保禄所说的：\BibleQuote{\Emph{我们劝告你们，弟兄们，不要白白地接受天主的恩宠。因为经上说：\InnerQuote{在悦纳的时候，我俯允了你；在救恩的日子，我帮助了你。}看！现在正是悦纳的时候；看！现在正是救恩的日子}} \BibleRef{格后 6:1}。\par

但我深知，当我的信被收到时，立刻有朋友聚集到你身边，你的文书食客被召来，关于生活目的的忠告竟向死亡推动者求取；这些人不是爱你，而是爱你所拥有的，他们只会告诉你当时令你喜欢的话。因为你亲自记得，你从前的顾问正是如此，他们引诱你犯下如此大罪。容我引用某位世俗作者的话给你：\Quote{凡事当与朋友商议，但应先考量朋友本身。}然而，你若想要一位顾问，我求你，以我为你的顾问。因为一个人若爱的是你本人，而不是你的财富，他给出的建议才最值得信赖。愿全能的天主让你心中明白，我的内心怀着何等关爱和爱德拥抱你——只要不冒犯天主的恩宠。因为我如此攻击你的过错，却如此爱你的本人；我如此爱你的本人，却不包容你过错的邪恶。因此，你若相信我爱你，请走到宗徒的门槛前，以我为顾问。但如果我因天主的缘故被认为过分尖锐，并因我的热忱而受怀疑，我将就这问题召集整个教会共同商议，凡众人认为为了善而应当做的，我绝不反对，反而欣然实行，并同意集体的决定。愿天主的恩宠在你完成我所告诫之事时护佑你。\par

\SourceRecord{Translated by James Barmby.；Nicene and Post-Nicene Fathers, Second Series；Vol. 12.；Edited by Philip Schaff and Henry Wace.；Buffalo, NY: Christian Literature Publishing Co.,；1895.；<http://www.newadvent.org/fathers/360201034.htm>.}

% structure: p0001 source=h1 target=section reason=page-title

\section{第一卷，第三十五封信}

\Emph{致特拉契纳主教伯多禄}\par

\Person{额我略}致\Person{伯多禄}等\par

犹太人\Person{若瑟}，即此信的持信者，告诉我们：住在\Place{特拉契纳}营地的犹太人，惯常在某些地方聚集庆祝他们的节庆，而您曾将他们从那里驱逐，他们便在您知情并同意下迁移到另一个地方，同样进行他们的节庆；如今他们抱怨又从同一地方被驱逐。但是，如果情况如此，我们希望您避免提供此类抱怨的理由，并允许他们按照惯例，在我们已经说过的地方聚集，即他们在您知情下获得作为聚会场所的地方。因为那些反对基督宗教的人，必须用温和、善意、劝告、说服来聚集到信德的合一，免得那些本可由讲道的甘饴和未来审判的可怕预期所吸引而信从的人，反被威胁和恐吓所驱退。因此，他们应该和善地聚集，从您这里听到天主的圣言，而不是害怕过度的严厉。\par

\SourceRecord{Translated by James Barmby.；Nicene and Post-Nicene Fathers, Second Series；Vol. 12.；Edited by Philip Schaff and Henry Wace.；Buffalo, NY: Christian Literature Publishing Co.,；1895.；<http://www.newadvent.org/fathers/360201035.htm>.}

% structure: p0001 source=h1 target=section reason=page-title

\section{第一卷 第三十六封信}

\Emph{致副执事伯多禄}\par

\Person{额我略}主教，天主的众仆之仆，致副执事\Person{伯多禄}。\par

我交给你的那册训令，在你前往西西里时必须仔细研读，以便对主教们格外留意，免得他们卷入世俗事务，除非为保护穷人所迫。但我认为，同一册训令中关于隐修士或圣职人员的规定，绝不可有任何变动。不过，愿你的阁下以极大的专注遵守这些事项，以达成我在此事上的愿望。此外，我听闻，从护民官安多尼诺的时代至今，在这最近十年间，有许多人遭受了来自罗马教会的暴行，以至于有人公开抱怨他们的地界被暴力侵犯，奴隶被夺走，动产被强行掠去，且未经任何司法程序。对于所有这类案件，我愿你的阁下保持警觉，凡在这十年间被发现以暴力夺走或借教会之名不义扣留之物，都凭我这道命令的权威，归还给发现其所属的人；免得那遭受暴力的人被迫前来我这里，承担如此长途跋涉的辛劳，如此一来，在我面前也无法查明他所说的是否属实。因此，仰望那将要来临的审判者的威严，归还所有被有罪地夺走之物，要知道，你若积聚[天上的]赏报而非财富，便为我带来极大的收益。但我们查明，大多数人抱怨的是丧失他们的奴隶，他们说，若有某个人的奴仆，或许是从其主人那里逃跑，自称属于教会，教会的管理者便立即将其作为教会的奴仆扣留，不经任何案件审理，反而蛮横地支持奴仆的话。这令我极为不悦，如同它违背真理的判断一样。因此，我愿你的阁下毫无延迟地纠正任何被如此行的事；并且，凡现在被教会产业持有的这类奴仆，既然未经审判而被夺走，也应在审判前予以归还；如此，若圣教会对他们有任何合法的权利，然后才可通过正常的法律程序剥夺其持有者的所有权。你要坚决地纠正所有这些事，因为你将真正是有福的宗徒伯多禄的战士，倘若在他的事上，即使他得不到任何东西，你也能守护真理。但是，如果你看到任何可以正义地主张为属于教会的东西，务必当心，绝不要试图以武力来主张这种权利；尤其是我已制定了一项处于绝罚之下的法令，即我们教会绝不可在任何城市或乡村田庄上设立\OriginalTerm{产权证书}{tituli}；但凡可基于理性为穷人主张之物，也应以理性来辩护；免得一件好事以不好的方式做，我们在全能天主面前连所正义寻求之事也被定为不义。此外，我请求你，要叫尊贵的平信徒和荣耀的\OriginalTerm{总督}{Prætor}因你的谦卑而爱你，而非因你的骄傲而畏惧你。然而，你若偶然知道他们正在对贫困者行不义，则立即把你的谦卑转为高昂，如此，当他们行善时你总是顺从，当他们行恶时你则反对。但你要如此行事，免得你的谦卑松懈，或你的权威僵化，为使正直调和谦卑，而谦卑使你的正直本身变得柔和。再者，既然主教们习惯为教宗的周年纪念在此聚集，禁止他们在我祝圣之日前来，因为愚昧且虚妄的铺张并不使我喜悦。但如果他们必须聚集，就让他们为宗徒之长伯多禄的周年纪念而来，感谢那位以其恩赐使他们成为牧者者。\Signature{祝好。写于四月朔日前第十七日，摩里修斯皇帝第九年。}\par

\SourceRecord{Translated by James Barmby.；Nicene and Post-Nicene Fathers, Second Series；Vol. 12.；Edited by Philip Schaff and Henry Wace.；Buffalo, NY: Christian Literature Publishing Co.,；1895.；<http://www.newadvent.org/fathers/360201036.htm>.}

% structure: p0001 source=h1 target=section reason=page-title

\section{第一卷 第三十九函}

\Emph{\Person{安特弥乌}副执事}。\par

\Person{额我略}致\Person{安特弥乌}，等。\par

我们曾嘱咐你，并记得后来也写信叮嘱你，要照顾穷人；若你在那些地区发现有人需要接济，应来信告知。你已费心办理，但涉及人数甚少。如今，我愿你一接到此命令，便立即向我父亲的姊妹\Person{帕特利亚}提供四十\OriginalTerm{索利多}{solidi}作为她孩子们的鞋钱，并四百\OriginalTerm{莫迪}{modii}小麦；向\Person{乌尔彼库斯}的遗孀\Person{帕拉蒂娜}夫人提供二十索利多和三百莫迪小麦；向\Person{费利克斯}的遗孀\Person{维维亚娜}夫人提供二十索利多和三百莫迪小麦。所有这些八十索利多应一并记入你的账目。但要速将你的收据总额带来，并蒙主助佑，在复活节前到达此地。\par

\SourceRecord{Translated by James Barmby.；Nicene and Post-Nicene Fathers, Second Series；Vol. 12.；Edited by Philip Schaff and Henry Wace.；Buffalo, NY: Christian Literature Publishing Co.,；1895.；<http://www.newadvent.org/fathers/360201039.htm>.}

% structure: p0001 source=h1 target=section reason=page-title

\section{第一卷，第四十一封信}

\Emph{致副执事伯多禄。}\par

额我略致伯多禄等。\par

托鲁姆城（\Emph{布鲁提亚的托里亚诺姆}）可敬的保利努斯主教告诉我们，由于蛮族的入侵，他的隐修士们已经四散，如今正在整个西西里游荡，并且因为没有领导者，他们既不关心自己的灵魂，也不留意修会的纪律。因此，我们命令你以一切谨慎和勤勉寻找并集合这些隐修士，将他们与他们的领导者——该主教——一起安置在位于墨西拿城的圣德奥多禄隐修院中，以便那些现在在那里、我们发现有需要领导者的人，以及那些你可能已经找到并带回的属于他团体的隐修士，都能在他的领导下共同事奉全能的主。此外，请知悉我们已将此事通知了同一城市可敬的斐理伯主教，以免在他所管理的教区中任何规定在未经他知晓的情况下被打乱。\par

\SourceRecord{Translated by James Barmby.；Nicene and Post-Nicene Fathers, Second Series；Vol. 12.；Edited by Philip Schaff and Henry Wace.；Buffalo, NY: Christian Literature Publishing Co.,；1895.；<http://www.newadvent.org/fathers/360201041.htm>.}

% structure: p0001 source=h1 target=section reason=page-title

\section{第一卷 第四十二封信}

\Emph{致副助祭\Person{安特弥乌斯}}。\par

\Person{额我略}致\Person{安特弥乌斯}等。\par

我们的弟兄及主教同仁\Person{若望}，通过其圣职人\Person{犹斯都}送交给我们的文书中，于众多事项之外告知我们如下：有些\Place{苏伦托}教区的隐修士随意从一个隐修院迁移到另一个隐修院，并因渴望世俗生活而脱离自己院长的会规；不仅如此（众所周知这是非法的），他们还各自企图拥有私人财物。因此，我们凭此命令你的经验：自此以后，不许任何隐修士从一个隐修院迁移到另一个，并且你不得允许他们中任何人拥有任何私人财物。但是，若有任何人竟敢如此妄为，应予以适当强制，将他送回他最初所居的隐修院，置于他原先逃脱的院长权下；免得我们若允许如此重大的恶行不加纠正而任其发展，那些失丧之人的灵魂将向其长上的灵魂追索。再者，若有任何圣职人员偶然成为隐修士，不得允许他们重新返回他们先前服务过的同一教会，或任何其他教会；除非某位隐修士的生活如此圣善，以至于他先前服务的主教认为他堪当司铎职，如此他可被该主教拣选，并由他祝圣于他认为适当的地方。此外，由于我们获悉，有些隐修士竟陷入如此大的邪恶，以致公开娶妻；你当以全部警醒搜寻他们，一旦找到，以应有的强制将他们送回他们曾为隐修士的隐修院。但是，对于宣发修道愿的圣职人员，也应当按照我们上所述及的予以处理，不可疏忽。因为如此你将在天主眼中蒙悦纳，并得以分享圆满的赏报。\par

\SourceRecord{Translated by James Barmby.；Nicene and Post-Nicene Fathers, Second Series；Vol. 12.；Edited by Philip Schaff and Henry Wace.；Buffalo, NY: Christian Literature Publishing Co.,；1895.；<http://www.newadvent.org/fathers/360201042.htm>.}

% structure: p0001 source=h1 target=section reason=page-title

\section{第一卷，书信第四十三}

\Emph{致\Person{莱安德尔}，\Place{西斯帕利斯}（\Place{塞维利亚}）主教}。\par

\Person{额我略}致\Person{莱安德尔}等。\par

我本愿以全副心神回复你的书信，若非我因牧职辛劳而疲惫不堪，更倾向于哭泣而非言语。这一点，你的尊位务必在阅读我的书信时予以理解并宽宥，因我对我深爱之人言语疏忽。诚然，我在此世被如此波涛抛掷，以至于完全无法将那艘陈旧腐朽的船驶入港口——此船，在天主隐秘的安排下，由我执掌舵柄。前方波涛涌来，侧面堆起泡沫海浪，后方风暴紧随。在这一切中，我不得安宁，有时被迫逆流而驶，有时转舵侧行以避开波涛的斜侧威胁。我呻吟，因我感觉到因我的疏忽，罪恶的舱底水增长，风暴猛烈冲击船只，腐朽的木板已发出船难的声响。我流泪回忆我如何失去了安息的宁静海岸，叹息着眺望那仍因事务之风逆吹而无法抵达的陆地。因此，最亲爱的兄弟，若你爱我，请在这波涛中向我伸出你祈祷的手；如此，因你在我的劳苦中帮助我，你将在益处本身中变得更加坚强。\par

然而，我无法用言语完全表达我得知我们共同之子、最光荣的国王\Person{雷卡雷多}已以最完全的虔诚皈依了天主教信仰时所得的喜乐。你在信中向我描述他的品性，使我虽未识其人却已爱他。但是，既然你知晓古老仇敌的诡计，他如何对征服者准备更猛烈的战争，愿你的圣德更谨慎地守护他，使他能完成他已经良好开端的事业，且不因完成善工而自高；使他能持守他所认识的信仰，也藉其生活的功德，并以他的行为表明他是永恒之国的公民，为的是在多年历程之后，他能从王国进入王国。\par

但关于\OriginalTerm{三次浸水}{trine immersion}，没有比你本人所感受到的正确更真实的回答了，即：信仰相同之处，用法的差异不会损害圣教会。如今我们三次浸水，象征三日埋葬的圣事，以致当婴孩第三次从水中被提起时，可表达三日后的复活。或者，若有人以为这是出于对至尊天主圣三的敬礼，那么将受圣洗者一次浸入水中亦无不可，因为三个\OriginalTerm{位格}{subsistences}中只有一个实体，所以无论三次还是一次在圣洗中浸洗婴孩，都无可指责，因为三次浸水可表示三位格，一次浸水可表示天主性的单一。但是，既然至今异端者惯常三次浸洗婴孩，我认为你们中间不应如此行，免得他们一边数算浸水次数，一边分裂天主性，并因继续照常行而夸口胜过我们的习惯。此外，我将下列书卷寄给你——对我而言最甘饴的弟兄——其目录附于下方。我在对真福约伯的阐释中所讲论的，即你在信中表示希望寄给你的，因其意义与言辞均显薄弱——正如我在讲道中所宣讲的那样——我尽力将其改写成论著形式，现正由抄写员誊写。若非因送信人匆忙，我本愿将全部内容毫无删减地传送给你，尤其是我为你的尊位而写下这部作品，为使我可被视为为我至爱之人而辛劳流汗。此外，若你从教会事务中得暇，你已知我的情况：即使身体缺席，我总视你与我同在，因为我在内心深处铭刻着你的面容。\Signature{五月颁发。}\par

\SourceRecord{Translated by James Barmby.；Nicene and Post-Nicene Fathers, Second Series；Vol. 12.；Edited by Philip Schaff and Henry Wace.；Buffalo, NY: Christian Literature Publishing Co.,；1895.；<http://www.newadvent.org/fathers/360201043.htm>.}

% structure: p0001 source=h1 target=section reason=page-title

\section{第一卷，第四十四函}

\Emph{致\Person{伯多禄}，\Place{西西里}副执事}\par

\Person{额我略}致\Person{伯多禄}，等。\par

关于我们迟迟未能遣回你的信使一事，我们因复活节庆典的诸项事务繁忙，所以不能让他更早出发。至于你所渴望得蒙指示的那些问题，你将在下文得知，我们在全面考虑之后，如何作出了决定。\par

我们确知，教会的农夫在谷物的价格上深受其苦，因为在丰收时，指定他们缴纳的金额并未按适当比例调整。我们愿意，在所有时节，无论谷物收成是多或少，比例皆按市场价格而定。我们也愿意，因海难损失的谷物应全额报账；但以你方在运送时毫不疏忽为条件，免得因错过恰当的运送时间而导致损失，此损失由你的过失造成。此外，我们认为，以六分之一税（\OriginalTerm{六分之一税}{sextariatics}）的方式向教会的农夫收取任何东西，或强迫他们缴纳比教会谷仓所用的更大的莫狄乌斯（\OriginalTerm{莫狄乌斯}{modius}），是极其错误和不义的。因此，我们藉此文告命令，不得以超过十八赛克斯塔里（\OriginalTerm{赛克斯塔里}{sextarii}）的莫狄乌斯（\OriginalTerm{莫狄乌斯}{modii}）向教会的农夫收取谷物；除非偶然有什么水手通常额外收取的东西，且他们自己证明在船上已消耗。\par

我们也确知，在教会的一些田产上实行着一种极不公正的征收，即管理人要求每七十个中缴纳三个半莫狄乌斯——这是言之可耻的事。即便如此仍不够；据说还按多年惯例征收了其他东西。我们对此做法深恶痛绝，并愿将其从教产中彻底铲除。但是，无论是在这项税或其它细微的杂税上，请你的经验判断每一镑多付了多少，以及什么是以不义方式从农夫那里收取的；并将一切减少为固定的缴纳额，在农夫的能力范围内，让他们缴纳总额达七十二：既不得征收超过一镑的谷物，也不得征收过重的磅数，也不得额外征收任何超过一镑的杂税；相反，通过你的估价，按其支付能力，让缴纳额凑至一定数目，这样便绝不会有任何可耻的征收。但是，唯恐在我死后，这些我们曾作为附加费不允、但允许增至固定缴纳额的税目，又再次以任何方式加征，以致缴纳总额增加，农夫被迫支付额外费用，我们希望你起草担保文书，由你签署，声明每人应缴纳如此数额，并排除谷物（\OriginalTerm{西里夸}{siliquæ}）、杂税或谷仓费用。此外，在这些项目中，凡过去应归教产管理员（\OriginalTerm{教产}{patrimonii}）的所有收益，我们愿意凭此命令从总支付额中归你。\par

万事先提醒你留意此事：在征收钱粮时不得使用不公正的衡器。若你发现有，要将其打破，并命人制作合乎标准的。因为我的儿子、天主的仆人迪亚科努斯（\OriginalTerm{迪亚科努斯}{Diaconus}）已经发现了不合意的衡器；但他没有权力更换。因此，我们愿意，除少量口粮（\OriginalTerm{口粮}{cibaria}）外，不得向教会的农夫征收超出公正衡器的任何东西。\par

此外，我们确知，首次布鲁达提奥（\OriginalTerm{布鲁达提奥}{burdatio}）的征收极大地削弱了我们的农夫，因为他们尚未能出售劳动果实时就须缴税；而且他们没有自己的钱支付，便向公共当铺借款，为此支付高昂代价；结果是他们因沉重的开支而疲弱。因此我们藉此训令命令，你的经验从公共基金中预付给他们本可能从陌生人那里借到的全部款项，并由教会的农夫逐渐偿还，按他们有能力支付时进行，以免他们暂时拮据时，以过低价格出售以后本可足够支付应缴数额的东西，但即便如此仍不足。\par

我们也得知，在农夫的婚姻上收取过高的费用：关于此，我们下令，婚姻费用不得超过一个索里都斯（\OriginalTerm{索里都斯}{solidus}）。若有贫穷者，他们应给得更少；若有富裕者，他们绝不可超过上述一个索里都斯的数额。而且我们不愿这些婚姻费用中的任何部分记入我们的账上，而是应归于管理人（\OriginalTerm{管理人}{conductorem}）的收益。\par

我们也确知，当某些农夫去世时，他们的亲属不得继承他们，而是他们的财物被收归教会使用：关于此事，我们规定，死者的亲属若住在教会的产业上，应作为继承人继承他们，且不得从死者的财产中抽走任何东西。但若有人留下年幼的子女，应选择谨慎的人掌管他们父母的财物，直到他们达到能够管理自己财产的年龄。\par

我们也查明了，若一个家族中有人犯了过错，他需赔偿，但不是亲自赔偿，而是以其财物赔偿。关于此事，我们命令：无论何人犯了过错，都应亲自承受应得的惩罚。此外，不得收受他的任何\OriginalTerm{馈金}{commodum}，除非可能是一些微不足道之物，可归为派去处理此事之官员的收益。我们也查明了，每当一个农夫不义地取自其佃农之物，确实会向农夫追讨，却未归还给原主。关于此事，我们命令：凡以暴力取自家族中任何人之物，都应归还给原主，而不是归入我们的收益，以免我们自己也成为暴力的帮凶。此外，我们愿意：若你的经验（Your Experience）在任何时候派遣你管辖之人处理超越庄园范围的事宜，他们可以从所派往之人那里接受小额谢礼，但仅限于他们自己得益；因为我们不愿教会的库银被不义之财玷污。我们也命令你的经验（Your Experience）注意此事：教会庄园上的农夫决不可因\OriginalTerm{馈金}{commodum}而被任命；否则，因为指望馈金，农夫常被更换；这种更换除了导致教会庄园永不被耕种之外，还有什么结果？但也不可让租约（即由教会租给农夫者）按应付租金总额来调整。我们希望你从教会庄园上收取的仓储费和存粮不超过惯例。但你自己仓库中我们命你准备的存粮，应从外人那里取得。\par

我们听说有三磅金子不义地从Subpatriana的农夫彼得那里取走了。关于此事，你要仔细审问\OriginalTerm{护民官}{defensorem}方提努斯（Fantinus）；如果它们显然是不义且不当取走的，就立即归还。我们也查明了，农民们第二次支付了特奥多西乌斯（Theodosius）曾向他们索取但未上交的\OriginalTerm{附捐}{burdation}，以致他们被双倍征税。这是因为他的财物不足以偿还教会债务。但由于我们通过我们的儿子、天主的仆人迪亚科努斯（Diaconus）获知，这笔亏空可用他的财产弥补，我们愿意：五十七个\OriginalTerm{索利多}{solidi}应原数归还给农民，不得有任何减免，以免他们被双倍征税。此外，如果他的财产在赔偿农民后还剩余四十个\OriginalTerm{索利多}{solidi}（这笔钱据说也在你手中），我们愿意将这些钱给他的女儿，以帮助她赎回她已典当的财物。我们也希望将她父亲的金杯（\OriginalTerm{酒杯}{batiolam}）归还给她。\par

光荣的\OriginalTerm{军长}{magister militum}康帕尼亚努斯（Campanianus）曾从他的瓦罗庄园（Varronian estate）每年拨出十二个\OriginalTerm{索利多}{solidi}给他的书记若望（John）。我们命令你每年毫无迟疑地支付给农夫欧普路斯（Euplus）的孙女，尽管她可能已收到了该欧普路斯的所有动产，或许除了他的现金。我们也希望你从他的现金中给她二十五个\OriginalTerm{索利多}{solidi}。据说有一个银碟子典当了一个\OriginalTerm{索利多}{solidus}，一个杯子典当了六个\OriginalTerm{索利多}{solidi}。在审问秘书多米尼库斯（Dominicus）或其他可能知情者之后，赎回这些典当物，并将上述小器皿归还原主。\par

我们感谢你的勤谨（Your Solicitude），因为在我吩咐你，在我兄弟的事务中，把他那笔钱还给他之后，你竟如此将这回事忘却，仿佛一个最末等的奴隶对你说了什么话一般。但现在，就叫你的疏忽（Your Negligence）——我不能说是你的经验（Your Experience）——去设法办成此事。凡你发现属于他的、在安托尼努斯（Antoninus）手中的任何财物，都尽快还给他。\par

关于犹太人萨尔平古斯（Salpingus）之事，已找到一封信，我们已将它转寄给你，以便你阅读后，完全了解他的案件以及据说与此案有关的一位寡妇的案件，就那五十一个\OriginalTerm{索利多}{solidi}——已知应归还之款——作出你认为公正的答复，使债权人无论如何不被不义地剥夺其应得之债。\par

他（遗嘱人）的遗赠的一半已给了安托尼努斯（Antoninus）；另一半将被赎回。我们愿意这一半从公共财物中补足给他；不仅给他，也给\OriginalTerm{护民官}{defensoribus}和外方人（\OriginalTerm{外方人}{pergrinis}），遗嘱人以遗赠的名义留给他们任何事物。我们也愿意将遗赠支付给\OriginalTerm{家族}{familiæ}；但这属于我们的关切。因此，你为我们这一方，即四分之三计算清楚后，就进行支付。\par

我们愿意你从卡努西乌姆（Canusium）教会款项中给该教会的圣职人员一些钱，以便他们如今遭受匮乏者能有些许生计；也愿若天主乐意祝圣一位主教，他也能有维持生活之资。\par

至于失足的司铎或任何其他圣职人员，我们希望你处理他们的财产时保持清洁，不受任何玷污。但要寻找那些懂得按天主而生活的最贫穷的正式隐修院，将失足者交付这些隐修院做补赎；失足者的财产应归于他们被交付做补赎的地方，以便那些负责纠正他们的人能藉其资财获得帮助。但如果他们拥有亲属，应将财产交给其合法亲属；但须为那些被交付做补赎的人充分提供赡养。若教会团体中的任何成员，无论是司铎、肋未人、隐修士、圣职者或其他人，若失足，我们愿意将他们交付补赎，但教会保留对其财产的权利。然而，应让他们在补赎期间获得足够维持生活的份额，免得因一无所有而对所交付的地方造成负担。若有人在教会领地内有亲属，应将财产交付给他们，由他们保管，但须受教会权利的约束。\par

三年前，西西里所有教会的副执事已按照罗马教会的惯例，被禁止与妻子有任何夫妻关系。但我认为，对那些不习惯如此节欲、且未曾事先许诺守贞的人，强迫他们与妻子分开，并因此（但愿不会）陷入更糟的境地，是严苛且不当的。因此，在我看来，从今以后应告知所有主教，不得擅自擢升任何未曾许诺贞洁生活的人为副执事；这样，以往并非刻意追求的事就不会被强制要求，而未来则可谨慎防范。但那些自三年前禁令以来已与妻子节欲相处的人，应受赞扬和奖赏，并应劝勉他们继续善行。至于那些自禁令以来不愿与妻子断绝关系的，我们不愿他们晋升至神圣品位；因为除非一个人在担任圣职之前已享有经认可的贞洁，否则不应接近祭台的职务。\par

关于商人\Person{李伯辣托}，他把自己托付给教会，居住在\Place{秦齐}庄园，我们希望每年为他提供一份供应；这份供应由你自己估价应为多少，并报告给我，记入你的账目。至于当前的征收期，我已从我们的儿子、天主的仆人\Person{迪亚科努斯}那里得到了消息。\par

有一位名叫\Person{若望}的隐修士去世了，他将\Person{范提努斯}\OriginalTerm{守护者}{defensorem}立为一半财产的继承人。将留给后者的部分交付给他，但嘱咐他不得再犯类似行为。但指定他因劳作应得的报酬，使他不会徒劳无益；并让他记住，靠教会薪俸生活的人不应贪求私利。但如果有任何东西，无罪的、不贪不义之财的，通过经办教会事务的人归于教会，那么这些人不应没有劳苦的果实。然而，如何酬报他们，应保留由我们裁断。\par

关于\Person{鲁斯提恰努斯}的款项，彻底查清此案，并执行你认为公正的决定。劝告尊贵的\Person{亚历山大}了结他与圣教会之间的诉讼；如果他或许忽略此事，你应在敬畏天主并保持尊重的情况下，尽力了结此案。我们也希望你在此事上花费一些；如果可行，只要他终止与我们之间的诉讼，就免去他必须付给他人的费用。\par

立即归还那位失足并被送入隐修院的天主婢女的捐赠，以便（如我上面所说）那个承担照管她的辛劳的地方，能因她拥有的财物而得供养。同时，收回她手中由他人持有的任何财物，并将其交给前述隐修院。\par

将\Place{维亚诺瓦}的客栈的付款按照你告诉我们数额寄给我们，因为你手头有这些款项。但根据你的裁量，给那个你委派在同一产业中的代理人一些东西。\par

关于那位曾与\Person{特奥多西乌斯}在一起的天主婢女，名唤\Person{埃克斯特拉尼亚}，我认为你应给她一份补贴，如果你认为有利的话，或者至少把她所作的捐赠归还给她。隐修院的那间房屋，\Person{安多尼努斯}曾以三十\OriginalTerm{苏勒德斯}{solidi}从隐修院购得，你应偿还银钱后立即归还。在彻底查清真相后，归还那些玛瑙瓶，我随此信的携带者送回给你。\par

如果\Person{撒图尼努斯}有空并且没有被您使用，就派他到我们这里来。\Person{坎帕纳}夫人名下的一位名叫\Person{费利克斯}的佃农——她曾给他自由并命令免除对他的审查——说副执事\Person{马克西穆斯}从他那里拿走了七十二个\OriginalTerm{索利多}{solidi}，为了支付这笔钱，他声称变卖或抵押了他在\Place{西西里}的所有财产。但律师们说，关于欺诈行为，他不能免于审查。然而，当他从\Place{坎帕尼亚}返回我们这里时，他在一场风暴中丧生。我们希望你找到他的妻子和儿女，赎回他抵押的一切，偿还他变卖所得，并且还给他们提供一些生活费；因为\Person{马克西穆斯}曾将此人派往\Place{西西里}，并在那里拿走了他所声称的钱。因此，查明他被拿走了什么，并立即归还给他的妻子和儿女。仔细阅读所有这些事情，并抛开你那种一贯的疏忽。我寄给农民们的信件使你在所有庄园被传阅，好让他们知道在我授权下，可以在哪些方面自卫，对抗不法行为；无论是原件还是副本，都交给他们。确保你毫不懈怠地遵守一切：因为，关于我写信给你要你遵守正义的事，我是无罪的；而如果你疏忽大意，你就有罪了。思考那将要来临的可怕审判者：让你的良心现在就以敬畏和战栗迎接祂的来临，免得那时它无故恐惧[?]——当天地在祂面前颤抖的时候。你已经听到我希望做什么：确保你做到。\par

\SourceRecord{Translated by James Barmby.；Nicene and Post-Nicene Fathers, Second Series；Vol. 12.；Edited by Philip Schaff and Henry Wace.；Buffalo, NY: Christian Literature Publishing Co.,；1895.；<http://www.newadvent.org/fathers/360201044.htm>.}

% structure: p0001 source=h1 target=section reason=page-title

\section{第一卷·书信四十六}

\Emph{致副助祭伯多禄。}\par

\Person{额我略}致\Person{伯多禄}等。\par

神圣的诫命劝勉我们要爱人如己；既然我们受命以这样的爱德爱他们，我们更应当供给他们肉身的需要，以减轻他们的困乏——纵使不能完全解决，至少也予以一些支持。因此，既然我们得知最可敬的\Person{戈迪斯卡尔库斯}之子不仅失明，而且缺乏食物，我们觉得必须尽己所能予以照顾。故此，我们命令你的经验，凭此命令每年供给他二十四\OriginalTerm{莫底乌斯}{modii}小麦，以及十二\OriginalTerm{莫底乌斯}{modii}豆类和二十\OriginalTerm{迪西马特}{decimates}酒，以维持生命；这些日后可在你的账目中列支。所以，你要这样行事，好使持本信者不致因耽误领受主的恩赐而抱怨，也使你因妥善施与恩惠而成为有份者。\par

\SourceRecord{Translated by James Barmby.；Nicene and Post-Nicene Fathers, Second Series；Vol. 12.；Edited by Philip Schaff and Henry Wace.；Buffalo, NY: Christian Literature Publishing Co.,；1895.；<http://www.newadvent.org/fathers/360201046.htm>.}

% structure: p0001 source=h1 target=section reason=page-title

\section{第一卷，第四十七封信}

\Emph{致\Place{阿尔勒}主教\Person{维吉利乌斯}与\Place{马赛}主教\Person{狄奥多鲁斯}。}\par

\Person{格里高利}致\Place{高卢}地区\Place{阿尔勒}主教\Person{维吉利乌斯}与\Place{马赛}主教\Person{狄奥多鲁斯}。\par

虽然迄今未能找到适当的时间与适当的人选致信于您二位，并适当回报你们的问候，但结果是我现在可以同时尽到爱德与兄弟之情的责任，并提及我们收到的某些人的抱怨，涉及如何拯救迷途者的灵魂。有许多人，虽确实属于犹太教，住在本省，并时常因各种商务事务前往马赛地区，已告知我们，那些地区定居的许多犹太人被带到圣洗之泉，更多的是出于强迫而非宣讲。如今，我认为此等情形下的意图值得称赞，并承认它出自对我主的爱。但我担心，除非有圣经的充分支持伴随这同样的意图，否则它要么没有有益的结果，要么甚至（天主不许）进一步导致我们希望拯救的灵魂丧亡。因为，当任何人被带到圣洗之泉，不是因宣讲的甘饴，而是因强迫，他就回到他旧日的迷信，并因重生而更坏地死去。因此，您二位要通过频繁的宣讲激励这些人，以便藉着他们老师的甘饴，他们更渴望改变他们的旧生活。因为这样我们的目的才能正确实现，而皈依者的心思不再回到他旧日的呕吐。因此，必须向他们讲论，能烧尽他们心中错误的荆棘，并以宣讲光照他们中的黑暗，这样您二位可通过你们频繁的劝诫为他们获得赏报，并带领他们，按天主所赐，达至新生命的重生。\par

\SourceRecord{Translated by James Barmby.；Nicene and Post-Nicene Fathers, Second Series；Vol. 12.；Edited by Philip Schaff and Henry Wace.；Buffalo, NY: Christian Literature Publishing Co.,；1895.；<http://www.newadvent.org/fathers/360201047.htm>.}

% structure: p0001 source=h1 target=section reason=page-title

\section{第一卷 第四十八封信}

\Emph{致撒丁总督德奥多禄。}\par

额我略致德奥多禄等。\par

你心中所怀的公义，理应彰显于你的行为。如今，圣维笃隐修院（该院由可敬者维杜拉所建）的女院长尤利亚纳向我们申诉，称你的官员多纳杜斯要求占有上述隐修院的产业。他自恃有你的庇护，不屑将案件交付司法审查。因此，请阁下责令该官员与上述天主的婢女共同将此事提交仲裁，以便仲裁者就争议所作的任何裁决均得以执行；如此，无论他发现自己应失去或保有何物，其所行之事便不视为德行之举，而归于法律的正义。\par

此外，一位虔诚的女士庞培雅娜，据悉在她自己的宅院内建立了一座隐修院，她投诉其已故女婿的母亲试图废除女婿的遗嘱，以便那位母亲对其儿子的财产的最后处置归于无效。为此，我们认为有必要以慈父般的爱德劝勉阁下，乐意地、公正地支持善举，并惠然下令确保这些人的合法权利得到保障。兹恳求主恩赐您的生活道路顺遂，并赐予您对尊贵职务的顺利管理。\par

\SourceRecord{Translated by James Barmby.；Nicene and Post-Nicene Fathers, Second Series；Vol. 12.；Edited by Philip Schaff and Henry Wace.；Buffalo, NY: Christian Literature Publishing Co.,；1895.；<http://www.newadvent.org/fathers/360201048.htm>.}

% structure: p0001 source=h1 target=section reason=page-title

\section{第一卷，第四十九封书信}

\Emph{致\Person{奥诺拉图斯}执事}。\par

\Person{额我略}致\Person{奥诺拉图斯}等。\par

既然我们承担了治理的职责（尽管不配），就有义务在我们力所能及的范围内援助有需要的弟兄。因此，我们的兄弟、\Place{卡利阿里}（\Place{卡利亚里}）都城的总主教\Person{雅努阿里乌斯}来到\Place{罗马}城，告诉我们：光荣的\OriginalTerm{军务长官}{magister militum}\Person{狄奥多鲁斯}，已知接受了\Place{撒丁岛}的公爵领地，正在那里做出许多违背我们最虔诚的君主命令的事——君主们以适当的仁慈和温和，消除了业主或帝国公民的许多困苦。因此，我们希望你适时向我们的最虔诚的君主陈述此事，以符合上述岛屿的居民公正合理的要求；鉴于先前也曾将神圣的帝国诏书发送给光荣的\OriginalTerm{军务长官}{Magister militum}\Person{埃丹丘斯}，他在第七次\OriginalTerm{小纪}{indiction}时任\Place{撒丁岛}公爵，命令纠正所有这些当前的冤屈，以便他们的命令——源于他们虔敬的慷慨——能被后来掌权的公爵们坚定不移地遵守，并且其益处不被管理者浪费；好让我们能在君主的仁慈统治下过安静的生活，并且对于他们以平静之心赐予臣民的规章，他们能在永恒审判者来临时得到加倍的补偿。\par

\SourceRecord{Translated by James Barmby.；Nicene and Post-Nicene Fathers, Second Series；Vol. 12.；Edited by Philip Schaff and Henry Wace.；Buffalo, NY: Christian Literature Publishing Co.,；1895.；<http://www.newadvent.org/fathers/360201049.htm>.}

% structure: p0001 source=h1 target=section reason=page-title

\section{卷书一，第五十封信}

\Emph{致副助祭安特米乌斯}。\par

\Person{额我略}致\Person{安特米乌斯}等。\par

既然，按天主的安排，蒙祂喜悦，我们获得了治理的职位，同样，我们也应当关心托付给我们的灵魂。如今我们发现，在\Place{欧摩菲安岛}上（众所周知，那里有真福宗徒之长\Person{伯多禄}的一座祈祷所），由于蛮族的凶猛逼迫，大批来自不同产业的男子及其妻子逃往那里避难。我们认为这并不适宜：因为附近另有避难之处，为何妇女应与隐修士同住？因此，我们以本命令训诫你的经验，从今以后，不得允许任何妇女（无论她是否属于教会管辖）在那里居住或逗留；而应让她们为自己准备避难之处（如上所述，附近有许多），任意选择；以便从此断绝与妇女的一切往来；否则，若我们停止尽一切可能的照顾并防范仇敌的陷阱，那么今后（愿天主不许）若发生任何不当之事，我们便难辞其咎。因此，不要迟延，将一千五百磅铅交给持此信的院长\Person{斐理克斯}，他在同一岛上需要这些铅；日后在结算账目时，当全部数量知晓后，可记入你的账中。你应如此进行，以便你自己也能取用一些，若可有效用于该岛的建筑。此外，由于岛上隐修士的团体生活艰苦，我们禁止接受十八岁以下的男孩进入这些隐修院。若现今已有，则让你的经验将他们迁出，送至\Place{罗马城}。我们希望你也在\Place{帕尔玛利亚岛}及其他岛上完全遵守此项规定。\par

\SourceRecord{Translated by James Barmby.；Nicene and Post-Nicene Fathers, Second Series；Vol. 12.；Edited by Philip Schaff and Henry Wace.；Buffalo, NY: Christian Literature Publishing Co.,；1895.；<http://www.newadvent.org/fathers/360201050.htm>.}

% structure: p0001 source=h1 target=section reason=page-title

\section{第一卷，第五十二封书信}

\Emph{致辩护者\Person{西玛库}}。\par

\Person{格列高利}致\Person{西玛库}等。\par

我的儿子执事\Person{波尼法爵}告诉我，阁下曾写信说，由虔敬的女士\Person{拉比纳}所建的一座隐修院现已准备好安置隐修士。我确实赞扬了您的殷切关怀；但我们希望另选一处，而不是原先指定用于此目的的地点；但考虑到时局不稳，须寻找一处临海高地，其位置天然有防御之便，或至少无需大动干戈即可设防。如此我们可以派遣隐修士前往，以期该岛本身至今尚无隐修院，因有此生活方式而得以改善。\par

为执行和料理此事，我们已指示持此命令的\Person{霍罗修斯}，阁下须与他一同巡视\Place{科西嘉}海岸；若在私人手中发现更合适的地点，我们准备支付合理价钱，以便我们得以做出稳妥安排。我们已命令上述\Person{霍罗修斯}前往\Place{戈尔戈尼亚}岛；请阁下陪同他前往，且对据我们所知已侵入那里的恶行施以惩罚，使你施加的惩罚能使该岛日后也保持改正。令同一院长\Person{霍罗修斯}整顿该岛的隐修院，然后速返我们这里。因此，请阁下如此行事，使你在两件事上，即在\Place{科西嘉}安置隐修院和整顿\Place{戈尔戈尼亚}隐修士之事上，尽快服从，不是服从我们的意愿，而是服从全能天主的旨意。\par

此外，我们希望在\Place{科西嘉}居住的司铎被禁止与妇女有任何交往，除非是母亲、姐妹或妻子，对这些妇女应保持贞洁。至于阁下写信给我的儿子上述执事\Person{波尼法爵}提及的三个人，请给他们你认为足够的一切，因为他们处于严重匮乏中；此事日后我们会在你的账目上准予报销。颁于七月。\par

\SourceRecord{Translated by James Barmby.；Nicene and Post-Nicene Fathers, Second Series；Vol. 12.；Edited by Philip Schaff and Henry Wace.；Buffalo, NY: Christian Literature Publishing Co.,；1895.；<http://www.newadvent.org/fathers/360201052.htm>.}

% structure: p0001 source=h1 target=section reason=page-title

\section{\WorkTitle{第一卷，第五十六封信}}

\Emph{致副执事伯多禄。}\par

额我略致伯多禄等。\par

我们极其渴望庆祝圣人的庆节，认为有必要将这封指导信函寄给您的经验，通知您：我们已安排，在主的助佑下，于八月隆重举行祝圣典礼，对象是最近在弟兄们住所内建造的荣福玛利亚祈祷所——该住所由院长马里尼亚努斯管理——以便我们所开始的工作，能藉主的运作而完成。然而，由于该住所的贫困，我们在那个庆日需要予以援助，因此，我们愿您为庆祝祝圣典礼而提供给穷人：十枚金币（\OriginalTerm{金币}{solidi}）、三十瓮（\OriginalTerm{瓮}{amphoræ}）葡萄酒、二百只羔羊、两桶（\OriginalTerm{桶}{orcæ}）油、十二只阉羊和一百只母鸡；这些日后可在您的账目中报销。请立即照办，不得有任何延误，好让我们的愿望，在天主恩准下，得以迅速实现。\par

\SourceRecord{Translated by James Barmby.；Nicene and Post-Nicene Fathers, Second Series；Vol. 12.；Edited by Philip Schaff and Henry Wace.；Buffalo, NY: Christian Literature Publishing Co.,；1895.；<http://www.newadvent.org/fathers/360201056.htm>.}

% structure: p0001 source=h1 target=section reason=page-title

\section{第一卷，第五十七封信}

\Emph{致\Person{塞维鲁}主教。}\par

额我略致\Person{塞维鲁}等。\par

我们从您的兄弟情谊的来函中得知，关于主教的选举，有人一致赞成\Person{奥克莱阿提努斯}；既然我们不认可他，他们不必再费心。但请通知该城居民：如果他们在自己的教会中找到任何适合此职务的人，全体将选择转到他身上。否则，持此信者将指明一个人，我已告知他，选举的通知应以此人为准。此外，您在对同一教会进行视察时，应谨慎细心，以使教会财产保持不受侵犯，其利益在您的管理下依惯例得到关注。\par

\SourceRecord{Translated by James Barmby.；Nicene and Post-Nicene Fathers, Second Series；Vol. 12.；Edited by Philip Schaff and Henry Wace.；Buffalo, NY: Christian Literature Publishing Co.,；1895.；<http://www.newadvent.org/fathers/360201057.htm>.}

% structure: p0001 source=h1 target=section reason=page-title

\section{第一卷，第五十八封信}

\Emph{致阿尔西辛努斯公爵、里米尼城的圣职人员、贵族与平民（\OriginalTerm{等级与平民}{ordini et plebi}）}\par

\Person{格列高利}致\Person{阿尔西辛努斯}等。\par

你们的爱德在期待一位主教时是多么热忱，你们呈递给我们的报告文本显示。但是，由于在类似情形中，祝圣者应当格外谨慎，我们正以应有的审慎关注此事。因此，我们以此函劝告你们的爱德：不必为\Person{奥克莱阿提努斯}的事麻烦向我们请求；但如果你们自己的城里发现某人能有益地承担此职，且我们不能对他提出异议，就请你们的意见赞成他。但如果找不到合适的人，我们已向持信者提及一人，你们同样可同意他。但你们要同心合意忠诚地祈祷，无论谁被祝圣，他既能有益于你们，又能显示相称于我们天主的司铎服务。\par

\SourceRecord{Translated by James Barmby.；Nicene and Post-Nicene Fathers, Second Series；Vol. 12.；Edited by Philip Schaff and Henry Wace.；Buffalo, NY: Christian Literature Publishing Co.,；1895.；<http://www.newadvent.org/fathers/360201058.htm>.}

% structure: p0001 source=h1 target=section reason=page-title

\section{卷一，第六十一封信}

\Emph{致\Person{盖纳狄乌斯}，\Place{非洲}贵族兼总督。}\par

\Person{格里高利}致\Person{盖纳狄乌斯}等。\par

你不断将天主敬畏置于眼前，追求正义，敌人被压伏的颈项可以作证；然而，为使基督的恩宠使你的荣耀常享同样的兴盛，请照你惯常的作法，迅速制止你发现的一切非法行为，如此，你以正义之武器武装，能以信德之力——一切德行之首——克胜敌对的攻击。现在，我们的兄弟、\Place{塔里斯}城的同僚主教\Person{马里尼亚努斯}含泪向我们陈诉：他城中的穷人到处受骚扰，因礼物或付款的开销而困苦；此外，他教会的神职人员受到军事长官\Person{特奥多鲁斯}手下人的严重侵扰，遭受身体伤害；此事竟发展到（说来惊人）他们被投入监狱，而他自己也在涉及他教会的事务上受到上述荣耀之人的严重阻碍。这些事情若属实，与共和国的纪律是何等相违，你们自己知道。既然阁下您宜纠正这一切，我向尊驾致意，要求您不再容许这些事发生；反之，严格命令他停止侵害教会，并不得有人因强加于他们的负担或超出理性所允许的付款而受委屈；如果有什么诉讼，应由法律秩序而非权力恐怖来裁决。那么，我恳求你，在主启示下，如此纠正这一切，藉你禁令的威吓，使荣耀的\Person{特奥多鲁斯}及其手下戒绝此类行为——即便不顾及正直，至少出于对你命令的恐惧；如此，为增加你的声誉与赏报，正义与自由可在你所托管的地区繁荣。\par

\SourceRecord{Translated by James Barmby.；Nicene and Post-Nicene Fathers, Second Series；Vol. 12.；Edited by Philip Schaff and Henry Wace.；Buffalo, NY: Christian Literature Publishing Co.,；1895.；<http://www.newadvent.org/fathers/360201061.htm>.}

% structure: p0001 source=h1 target=section reason=page-title

\section{第一卷·第六二封信}

\Emph{致撒丁岛\Place{卡腊利斯}（\Place{卡利亚里}）总主教\Person{雅努雅}。}\par

\Person{额我略}致\Person{雅努雅}等。\par

如果我们的主亲自藉圣经的证词宣称自己是寡妇的丈夫和孤儿的父亲，那么，我们作为祂身体的肢体，也应当以灵魂的最高情感努力效法元首，并在不损及正义的前提下，在必要时支持孤儿和寡妇。而且，我们得知一位虔诚的女子\Person{卡特辣}——她有一个儿子在蒙天主恩典由我们主持的神圣罗马教会服役——正受到某些人的勒索和骚扰，我们认为有必要通过此信劝勉您的兄弟，不要拒绝（在不损及正义的前提下）给予该女子保护，因为通过此类善行，您既使主成为您的债务人，也以爱德之链使我们更紧密相连。因为我们希望上述女子的事，无论现在还是将来，都由您裁决，使她得以摆脱诉讼的烦恼，却绝不免除服从公正判决的义务。现在我祈求主引领您的一生在祂面前昌盛，并亲自仁慈地将您领入那将要来临的荣耀王国。\par

\SourceRecord{Translated by James Barmby.；Nicene and Post-Nicene Fathers, Second Series；Vol. 12.；Edited by Philip Schaff and Henry Wace.；Buffalo, NY: Christian Literature Publishing Co.,；1895.；<http://www.newadvent.org/fathers/360201062.htm>.}

% structure: p0001 source=h1 target=section reason=page-title

\section{\WorkTitle{第一卷，第六十三封信}}

\Emph{致撒丁岛\Place{卡利亚里}（Caralis）主教\Person{雅纳略}}\par

\Person{额我略}致\Person{雅纳略}等。\par

尽管阁下以正义的热忱合宜地关注于保护各类人士，但我们相信，对于蒙我们书信所推荐的人，您将更加乐意予以援助。要知道，一位名叫\Person{庞培亚纳}的虔诚妇女，曾通过其家人向我们陈诉，她不断无理地遭受某些人的诸多委屈，并因此恳求我们致信向您推荐她。因此，我们以应有的爱德之情问候阁下，认为必须将上述妇女托付给您，使您本着正义，不让她受到任何违反公平的委屈，或轻率地承担任何费用。但若她有任何诉讼，应在选定的仲裁人面前辩论争议事项；无论有何裁决，均请通过您的协助平静地执行，以便您因此善工获得赏报，而蒙我们书信推荐的那位妇女也能因寻获正义而喜乐。\par

\SourceRecord{Translated by James Barmby.；Nicene and Post-Nicene Fathers, Second Series；Vol. 12.；Edited by Philip Schaff and Henry Wace.；Buffalo, NY: Christian Literature Publishing Co.,；1895.；<http://www.newadvent.org/fathers/360201063.htm>.}

% structure: p0001 source=h1 target=section reason=page-title

\section{第一卷·第六十六封信}

\Emph{致\Place{墨西拿}（\Place{梅塞内}）主教\Person{费利克斯}。}\par

\Signature{额我略致费利克斯等。}\par

那些我们发现给教会带来负担的习俗，我们理应出于关怀予以废止，以免任何人被迫向那些他们本应指望得到捐助的地方做出贡献。因此，您的职责是保持圣职人员及其他人的惯例完好无损，并每年向他们转交所惯常的：但今后我们禁止您向我们转交任何东西。并且，既然我们不喜爱\OriginalTerm{礼物}{xeniis}，我们已感谢地接受了您兄弟情谊送给我们的\OriginalTerm{帕尔马提亚特产}{Palmatianæ}，但已按合理价格将其出售，并将售价另行转交给您的兄弟情谊，以免您感到开销。此外，既然我们得知您的爱德渴望来到我们这里，我们通过本信劝告您不必劳烦前来：但请为我们祈祷，以便我们虽因路途遥远而分离，却因爱德在基督的帮助下在心灵上更加团结；为的是通过彼此代祷互相帮助，我们得以将未曾损坏的职务交还给那将要来临的审判者。\par

\SourceRecord{Translated by James Barmby.；Nicene and Post-Nicene Fathers, Second Series；Vol. 12.；Edited by Philip Schaff and Henry Wace.；Buffalo, NY: Christian Literature Publishing Co.,；1895.；<http://www.newadvent.org/fathers/360201066.htm>.}

% structure: p0001 source=h1 target=section reason=page-title

\section{第一卷，第六十七封书信}

\Emph{致副执事伯多禄}\par

\Person{额我略}致\Person{伯多禄}等。\par

若我们以仁慈的心怀，藉着显示怜悯来满足邻人的需要，我们无疑会发现主仁慈地垂允我们的祈求。如今我们获悉，\Person{帕斯托尔}视力极度衰弱，他有一妻二奴，且先前曾与光荣的\Person{约纳塔}女士在一起，正遭受极大的匮乏。因此，我们藉此命令函告\OriginalTerm{你的经验}{Your Experience}，切勿延迟给予他三百\OriginalTerm{斗}{modii}小麦及同样数量的豆子作为生计，这些日后可记入你的账目。因此，请如此行事，既使你本人因善工而获得赏报，又使我们的命令得以执行。八月。\par

\SourceRecord{Translated by James Barmby.；Nicene and Post-Nicene Fathers, Second Series；Vol. 12.；Edited by Philip Schaff and Henry Wace.；Buffalo, NY: Christian Literature Publishing Co.,；1895.；<http://www.newadvent.org/fathers/360201067.htm>.}

% structure: p0001 source=h1 target=section reason=page-title

\section{第一卷，第七十二封信}

\Emph{给\Person{伯多禄}副执事。}\par

额我略致\Person{伯多禄}等。\par

你已从先前的信函得知，我们曾希望居住在\Place{西西里岛}的弟兄和同为主教们，在此聚集，庆祝荣福宗徒\Person{伯多禄}的周年纪念。但因他们与前副执政官、尊贵的\Person{犹斯定}所涉诉讼，暂时耽搁了他们，且目前往返时间已不充裕，我们不愿他们在冬季前受扰。但\Place{阿格里琴托}的\Person{额我略}、\Place{卡塔纳}的\Person{良}和\Place{帕诺尔穆斯}的\Person{维克多}，我们无论如何都希望他们在冬季前来到我们这里。此外，从外人那里收集价值五十金镑的当年新粮，存放在\Place{西西里}不会腐烂的地方，以便我们在二月尽可能多派船只去那里运粮给我们。但若我们延迟派船，你应亲自提供船只，并靠主的助佑，在二月将同一批粮食运给我们；但须排除我们按惯例期待在九月或十月当前寄送的那批粮食。那么，让你的经验如此进行：在不烦扰教会任何\OriginalTerm{田庄隶属农民}{colonus}的情况下，收集粮食；因为这里收成如此微薄，除非靠天主的帮助从\Place{西西里}收集粮食，否则将面临严重的饥荒。但你务必全面守护历来被指定用于圣教会的船只，正如尊贵的前执政官\Person{良}所写给你的信函也一致指示你去做的那样。此外，许多人来此希望租用属于我们教会的各种土地或岛屿；有些我们确实拒绝了，但另一些我们已准许了他们的请求。但让你的经验留意圣教会的利益，要记住你曾在荣福宗徒\Person{伯多禄}至圣的遗体前接受管理他产业的权利。并且，即使有信函从我们这里送达你，也不可做任何有损产业的事，因为我们不记得曾给予过任何东西，也无意在没有充分理由的情况下放弃任何东西。\par

\SourceRecord{Translated by James Barmby.；Nicene and Post-Nicene Fathers, Second Series；Vol. 12.；Edited by Philip Schaff and Henry Wace.；Buffalo, NY: Christian Literature Publishing Co.,；1895.；<http://www.newadvent.org/fathers/360201072.htm>.}

% structure: p0001 source=h1 target=section reason=page-title

\section{\WorkTitle{第一卷，第七四封书信}}

\Quote{致根纳狄乌斯，阿非利加贵族与总督。}\par

\Person{额我略}致\Person{根纳狄乌斯}等。\par

主既使尊贵的阁下在此生军事战争中以胜利之光闪耀，您也当以全心全力之警觉，对抗教会的仇敌，好使您因两种凯旋而声名愈发彰显；当您在法庭争战中，亦为基督徒民众坚固地抵抗天主教会之敌，如主的战士勇敢地从事教会之战时，亦复如是。因为众所周知，宗教信仰异端之人，若得准许加害（愿主禁止），便悍然攻击天主信仰，企图将其异端之毒传播，以败坏基督身体之肢体。我们得知，他们竟悖主而昂首反对天主教会，妄想败坏基督徒宣信之信仰。但务请阁下抑制其企图，并将其骄傲之颈项屈服于正义之轭下。此外，请下令告诫天主教主教会议，不得以其资历而不顾其生活功德来任命首席主教；因为在主前，所嘉许的并非更显赫的品级，而是更善之生活的行为。至于首席主教本人，不应如惯常那样在乡间往来居留，而应住在他们所选的一个城市里，以便他能够更好地运用那归于他的尊威之影响力，来抵抗多纳徒派。再者，若努米底亚会议中有人愿来宗座，请允准之；并将那些意图指控其品格的人加以阻止。若通过您，分裂的教会得以合一，则阁下在造物主面前必得更大光荣之增加。因为当主看见祂所赐予的恩惠归于祂的荣耀时，祂愈益丰厚地赐予恩惠，正如同祂见到自己宗教的尊威因此而扩大。此外，我们本于应有之情，向您致上我们慈父般爱德的关怀，并恳求主使您的手臂强健，以克敌制胜，并以炽热的剑锋般的信仰热忱，磨砺您的心灵。\par

\SourceRecord{Translated by James Barmby.；Nicene and Post-Nicene Fathers, Second Series；Vol. 12.；Edited by Philip Schaff and Henry Wace.；Buffalo, NY: Christian Literature Publishing Co.,；1895.；<http://www.newadvent.org/fathers/360201074.htm>.}

% structure: p0001 source=h1 target=section reason=page-title

\section{第一卷，第七十五封信}

\Emph{致全非洲总督兼大元帅\Person{格纳迪乌斯}}\par

\Person{额我略}致\Person{格纳迪乌斯}总督兼大元帅等。\par

若贵上军事功绩如此卓著，并非源于您信德之功及基督宗教的恩宠，则其亦不足为奇，因我等知晓古代将领亦曾蒙受类似恩赐。然而，当天主恩准，您并非以血肉之筹划，而是以祈祷来预占未来的胜利，这便令人惊奇：您的荣耀并非来自世俗谋略，而是从天主那里由上倾注于您。因为何处不在传颂您功勋的美名？传闻称，您屡次征战并非出于嗜血之欲，而是为扩展我们所见天主受崇拜的共和国，旨在使基督之名通过信仰的宣讲在臣服的民族中传播。因为，正如您外在的英勇事迹使您今生显赫，同样，您内在品德的装饰——源于纯洁的心灵——也使您荣升，分享来日天堂的喜乐。因为我们得知，贵上曾为喂养蒙福的宗徒之长\Person{伯多禄}的羊群做了许多有益之事，以至于您通过提供达提提安移民，恢复了宗座产业中因缺乏适当耕作者而荒芜的不少部分。因此，无论您以基督徒的心怀赐予他什么，您都将在对未来审判的希望中获得回报。故而，我们认为应当将呈递此信的\Person{希拉鲁斯}推荐给阁下，以便您在他可能需要您帮助的事务中，赐予他一贯的关爱（但始终兼顾正义）。现在，我们以慈父之爱向您致候，恳求我们的天主和救主仁慈地保护阁下，以安慰神圣的共和国，并加强您手臂的力量，使祂的名通过附近民族日益广传。\par

\SourceRecord{Translated by James Barmby.；Nicene and Post-Nicene Fathers, Second Series；Vol. 12.；Edited by Philip Schaff and Henry Wace.；Buffalo, NY: Christian Literature Publishing Co.,；1895.；<http://www.newadvent.org/fathers/360201075.htm>.}

% structure: p0001 source=h1 target=section reason=page-title

\section{《第一书，第七十七封》}

\Emph{致努米底亚全体主教。}\par

额我略致努米底亚全体主教。\par

至亲的基督内弟兄们，若莠子令人烦恼地混杂在青绿的谷物中，农人的手就必须将它完全拔除，免得阻碍肥沃谷物的未来果实。因此，我们这些虽不配、却承担了主田园耕作的人，也要赶紧使谷物纯洁、不受莠子的一切阻碍，好让主的田园结出更丰盛的收成。你们曾通过我们的书记官希拉鲁斯向我们可敬的前任教宗请求，保留你们过去的一切习俗，这些习俗从真福的宗徒之长伯多禄所颁定的初规以来，久远的古人至今仍保存。我们呢，确实根据你们陈述的内容，允许你们的习俗（只要它显然不妄图损害天主教信仰）不受干扰地保留，无论是关于推举首席主教还是其他方面；但是，对于那些从多纳徒派中获得主教职位的人，我们绝对禁止他们被擢升到首席的尊严，即使他们的资历显示他们应当得到那职位。让他们满足于照顾托付给他们的民众，而不要企图在首席的至高位上胜过那些由天主教信仰在教会怀抱中教导并滋养的主教们。因此，至亲的弟兄们，你们要在主的爱德热忱中期待我们的劝诫，知道严厉的审判者将审察我们所做的一切，并会按我们工作的功绩，而非更高职位的特权，来认可我们每一个人。我恳求你们，彼此相爱，在基督内彼此和睦，同心一致反对异端者和教会的敌人。要关心邻人的灵魂：尽力以爱德的宣讲劝服众人来信从，同时也向他们摆出未来审判的恐怖；因为你们被立为牧者，羊群之主期待受托付的牧者带来加倍羊群的果实。如果祂预见因你们对祂羊群更勤勉的照顾而使之增长，祂必将以天上的王国多种恩赐装饰你们。再者，向你们表达兄弟之爱的问候，我祈求主使你们——祂已拣选为灵魂的牧者——在祂眼中堪当，并亲自如此安排我们在此世的行为，使祂能在来世中按它们所应得的方式接纳。\par

\SourceRecord{Translated by James Barmby.；Nicene and Post-Nicene Fathers, Second Series；Vol. 12.；Edited by Philip Schaff and Henry Wace.；Buffalo, NY: Christian Literature Publishing Co.,；1895.；<http://www.newadvent.org/fathers/360201077.htm>.}

% structure: p0001 source=h1 target=section reason=page-title

\section{第一卷，第七十八书}

\Emph{致科西嘉岛的主教\Person{良}}\par

\Person{额我略}致\Person{良}等。\par

我们的牧职责任催促我们以焦虑的关怀去援助一个没有司铎管理的教会。由于我们得知，\Place{萨奥纳}教会自其主教去世以来多年完全无人管理，我们认为有必要嘱咐你的兄弟情谊承担探访该教会的任务，以便通过你的安排促进其福祉。在该教会及其堂区内，我们授权你祝圣执事和司铎；然而，关于这些人，你应当谨慎查考，确保他们在个人方面没有任何被神圣法典所拒绝之处。但凡你的兄弟情谊认为堪当如此重大职务的人，查明其品行和行为适合接受圣职后，你可凭我们的权威许可，自由地擢升他们至上述职务。因此，我们希望你像对待自己的主教区一样，使用上述教会的一切财产，直到我们再次写信给你。因此，在所有这些事情上，你要如此勤勉和谨慎，以致通过你的安排，一切都能在天主的助佑下，为教会的益处而妥善安排。\par

\SourceRecord{Translated by James Barmby.；Nicene and Post-Nicene Fathers, Second Series；Vol. 12.；Edited by Philip Schaff and Henry Wace.；Buffalo, NY: Christian Literature Publishing Co.,；1895.；<http://www.newadvent.org/fathers/360201078.htm>.}

% structure: p0001 source=h1 target=section reason=page-title

\section{第一卷，第79封信}

\Emph{致\Place{科西嘉}主教\Person{玛尔定}。}\par

\Person{额我略}致\Person{玛尔定}等。\par

对于请求合理之事的人，我们理应善意垂听，以便请求者能获得他们所盼望的补救，而教会也不缺少牧者的悉心关怀。由于\Place{塔纳特斯}教会——你的兄弟曾在那里以司祭尊位获得荣耀——因罪恶已被敌意的野蛮占据并摧毁，以致你返回那里不再有任何希望，我们通过此信的权威，任命你在\Place{萨奥纳}教会为无可争议的枢机司铎，该教会长久以来缺乏主教的援助。因此，你要按照教会法典的训令，以对天主的爱，谨慎用心地安排和整理一切事，使你的兄弟因达成心愿而欢喜，并使天主教会因接纳你为枢机主教而充满相应的喜乐。\par

\SourceRecord{Translated by James Barmby.；Nicene and Post-Nicene Fathers, Second Series；Vol. 12.；Edited by Philip Schaff and Henry Wace.；Buffalo, NY: Christian Literature Publishing Co.,；1895.；<http://www.newadvent.org/fathers/360201079.htm>.}

% structure: p0001 source=h1 target=section reason=page-title

\section{第一卷·第八十书}

\Emph{致科西嘉岛的神职人员和贵族}\par

额我略致神职人员等…… \OriginalTerm{平等者之间}{A paribus}。\par

虽然长久以来，天主的教会没有主教并未使你们感到忧伤，但我们由于我们所担负的职务的职责，而且尤其被我们对你们的爱德所束缚，不得不为它的治理操心，知道在它的监督中同时有利于你们的灵魂。因为，如果羊群缺少牧人的照顾，它就容易落入埋伏者的陷阱。因此，既然\Place{萨奥纳}教会长期缺乏司铎的帮助，我们认为有必要设立我们的兄弟和同僚主教\Person{玛尔定}为其枢机司铎，但将视察的任务托付给我们的兄弟和同僚主教\Person{良}。我们也授予后者在其内及其堂区祝祭司铎和执事的权力，并允许他在逗留期间使用其财产，如同他是当地主教一样。因此，我们通过本函劝告你们，你们的爱德要全心接待上述视察员，并在一切合理的事上服从他，如同教会的子女所当行的，以便因你们的虔诚支持，他能完成一切有助于上述教会利益的事。\par

\SourceRecord{Translated by James Barmby.；Nicene and Post-Nicene Fathers, Second Series；Vol. 12.；Edited by Philip Schaff and Henry Wace.；Buffalo, NY: Christian Literature Publishing Co.,；1895.；<http://www.newadvent.org/fathers/360201080.htm>.}

% structure: p0001 source=h1 target=section reason=page-title

\section{第二卷·第三封书函}

\Emph{致\Person{Velox}，\OriginalTerm{军队长官}{Magister Militium}}\par

\Person{Gregory}致\Person{Velox}等\par

不久前，我们已通知阁下，有士兵准备好前往贵处；但因您的来信告知敌军已集结并向此处进发，我们因而将他们留在此地。不过，现在看来，派遣一定数量的士兵到您那里似乎是有利的，请阁下务必告诫并鼓励他们做好准备以应对劳累。当您找到机会时，请与我们的光荣之子\Person{Maurilius}和\Person{Vitalianus}商议，按照天主的助佑，为共和国的利益，他们指定您做什么，就去做什么。如果您查明那卑劣的\Person{Ariulph}正在向此地或拉文纳附近侵扰，请您像勇敢的人一样在他的后方尽力，使您的声望因天主的助佑而在共和国中因您的劳苦更加增长。不过，我们特别劝您做这件事：毫无延迟或借口地释放\Person{Maloin}、\Person{Adobin}、\Person{Vigild}和\Person{Grussing}一家，他们已知与光荣的\Emph{\OriginalTerm{军队长官}{Magister militum}} \Person{Maurilius}在一起，以便上述\Person{Maurilius}的属下到达贵处时，能毫无阻碍地与他们一同行进。\par

\EditorialNote{在\WorkTitle{Colbert}抄本和\WorkTitle{Paul. diac.}中，\OriginalTerm{10月朔日前的第五日}{Die. V. Kal. Oct.}，\OriginalTerm{第十小纪}{Indict. 10}}\par

\SourceRecord{Translated by James Barmby.；Nicene and Post-Nicene Fathers, Second Series；Vol. 12.；Edited by Philip Schaff and Henry Wace.；Buffalo, NY: Christian Literature Publishing Co.,；1895.；<http://www.newadvent.org/fathers/360202003.htm>.}

% structure: p0001 source=h1 target=section reason=page-title

\section{第二卷，第六封信}

\Emph{致那不勒斯人。}\par

额我略致居住在\Place{那不勒斯}的圣职界、贵族、士绅及平民。\par

虽然身为母亲之教会的属灵儿女所怀有的真诚敬意无需劝勉，但仍应以书信加以激励，以免其以为被忽视。为此，我以慈父般的爱德劝告你们的爱德：愿我们以许多眼泪同心合意地向我们的救赎主献上感谢，祂没有容许你们在如此悖逆的教师带领下行于无路之途，而是公开揭露了你们不称职的牧者的罪行。因为那德默特琉，即便此前已不配被称为主教，竟被发现卷入如此严重且恶劣的行径，若按其所行受毫不留情的审判，无疑会同时被神律和人间法律处以极刑。然而，因他被保留以行补赎，已被剥夺司铎职位的尊荣；我们不能容忍天主的教会长期没有教师，因为教规规定，牧者去世或撤职后，教会不可长久缺乏司铎。因此，我认为有必要通过本函劝告你们的爱德：既不要因延误，也不要因往常引发丑闻的不和，阻碍你们推选主教。你们当尽心寻觅这样一位人选：众人都能同心合意地以他为喜乐，且他绝不触犯神圣教规；为的是那曾因恶人之邪恶治理而被玷污的职位，无论被祝圣者是谁，愿他藉基督的恩宠并蒙祂悦纳，得以相称地担当并治理。\par

\SourceRecord{Translated by James Barmby.；Nicene and Post-Nicene Fathers, Second Series；Vol. 12.；Edited by Philip Schaff and Henry Wace.；Buffalo, NY: Christian Literature Publishing Co.,；1895.；<http://www.newadvent.org/fathers/360202006.htm>.}

% structure: p0001 source=h1 target=section reason=page-title

\section{第二卷·第七封}

\Emph{致\Person{马克西米安}，\Place{锡拉库萨}主教。}\par

\Person{额我略}致\Person{马克西米安}等。\par

当我们与弟兄们分担重担时，我们更有效地执行我们的属天使命。为此，我们任命您，我们最敬爱的弟兄和同为主教的弟兄，以宗座之名管理西西里所有教会，以便凡在那里被视为处于宗教状况的人，都应按我们的权柄服从您，为的是今后他们无需因轻微原因而长途航海来投奔我们。但若偶然有难以解决的问题，无论如何不能由您裁决，则唯有此类问题才应征求我们的裁决，以便我们因从最小的事务中解脱而能更有效地专注于更大的事务。且应明白，我们授予此权柄，不是给您的职位，而是给您本人，因为我们从您过去的生活中已能推断您未来的行为。\par

十二月，第十小纪。\par

\SourceRecord{Translated by James Barmby.；Nicene and Post-Nicene Fathers, Second Series；Vol. 12.；Edited by Philip Schaff and Henry Wace.；Buffalo, NY: Christian Literature Publishing Co.,；1895.；<http://www.newadvent.org/fathers/360202007.htm>.}

% structure: p0001 source=h1 target=section reason=page-title

\section{卷二 第九封书信}

\Emph{致那不勒斯人。}\par

额我略致居住在那不勒斯的\OriginalTerm{贵族和平民}{ordini et plebi}。\par

你们致我们的信函已表明你们对我们兄弟和同主教保禄的看法：我们祝贺你们，因你们与他短短几日的相处，竟如此渴望拥他为你们的\OriginalTerm{枢机主教}{cardinal bishop}。但是，在至极重要之事上不宜草率决定，因此，我们在基督助佑下，经过深思熟虑后将安排日后之事，同时，随着时间推移，他的品性将为你们所更熟悉。\par

所以，最亲爱的儿子们，如果你们真心爱戴上述之人，就应服从他，并以虔敬之心在和谐中满足他的意愿，好使你们彼此爱德的情感将你们如此紧密地联结，以致那在你们周围狂吠的仇敌，无法从你们任何人中间找到缝隙潜入而破坏你们的同心一致。再者，当我们看出上述主教已按我们所期望的、也为天主自己所赞许的，向天主奉献了灵魂的果实，此后，我们将遵循神圣的感召，就他的个人及你们的愿望，做应做的一切。\par

\SourceRecord{Translated by James Barmby.；Nicene and Post-Nicene Fathers, Second Series；Vol. 12.；Edited by Philip Schaff and Henry Wace.；Buffalo, NY: Christian Literature Publishing Co.,；1895.；<http://www.newadvent.org/fathers/360202009.htm>.}

% structure: p0001 source=h1 target=section reason=page-title

\section{第二卷 第十封信}

\Emph{致\Place{那不勒斯}主教\Person{保禄}}。\par

\Signature{额我略致保禄等。}\par

如果我们稳妥地履行所接受的司铎职务，那么天主的助佑和我们属灵儿女的爱戴无疑不会缺少。因此，请您的仁爱注意在一切事上表现自己，使该城圣职人员、贵族及全体人民共同为您所作的见证，因您的美善增长而更加坚固。那么，您应当如此致力于不断劝勉上述人民，使天主的农夫能将您因劳苦从他们那里收取的言语之果，储存在他的仓库中。但是，在我们能够（天主启示其旨意）就我们上述儿女所请求的事宜进行商议之前，我们准许从平信徒中任命圣职人员，也准许在您面前于该教堂内隆重举行释放奴隶的仪式。此外，我们希望您毫不迟疑地遵守上述教会中圣职界和长老们的习俗；您也要在教导他们方面如此勤谨看守，使他们远离一切不合宜或非法之事，在您的劝勉下，恒心持守，以应有的服从，在我们天主的服务中尽职。一月，第十小纪。\par

\SourceRecord{Translated by James Barmby.；Nicene and Post-Nicene Fathers, Second Series；Vol. 12.；Edited by Philip Schaff and Henry Wace.；Buffalo, NY: Christian Literature Publishing Co.,；1895.；<http://www.newadvent.org/fathers/360202010.htm>.}

% structure: p0001 source=h1 target=section reason=page-title

\section{第二卷·第十二封信}

\Emph{致\Place{亚里米伦}主教\Person{卡斯托里乌斯}}\par

\Person{额我略}致\Person{卡斯托里乌斯}等。\par

显贵女子\Person{蒂莫泰亚}通过一份请愿书（如下所述）告知我们，她在\Place{亚里米伦}城内属于她的一块土地上建立了一个祈祷所，她希望该祈祷所被祝圣以恭敬圣十字架。因此，最亲爱的弟兄，如果所述建筑在您城市的辖境内，并且已知没有遗体被埋葬在那里，那么，在首先接受了合法的捐赠基金之后，即她全部财产（奴隶除外）的三分之二，包括动产、固定装置和牲畜，为其终身保留用益权，并通过市政契约确保该捐赠基金之后，您应隆重祝圣上述祈祷所，但不举行公弥撒，条件是在将来不可在同一地点建造圣洗池，并且您不得任命一位枢机司铎。倘若她希望在那里举行弥撒，请告知她必须请求您的爱德派遣一位司铎，以免任何其他司铎擅自行事。此外，您将虔诚地存放她所准备的圣物。\par

\SourceRecord{Translated by James Barmby.；Nicene and Post-Nicene Fathers, Second Series；Vol. 12.；Edited by Philip Schaff and Henry Wace.；Buffalo, NY: Christian Literature Publishing Co.,；1895.；<http://www.newadvent.org/fathers/360202012.htm>.}

% structure: p0001 source=h1 target=section reason=page-title

\section{\WorkTitle{第二卷 第十五封信}}

\Emph{致保禄主教。}\par

国瑞致保禄等。\par

我曾委派贤弟暂时主持那不勒斯教会，为使你能通过善讲道理引领一切可能的人归向天主。正当你应当全心投入这项工作时，你却急于在为主结出果实之前返回，并请求我迅速处理同教会的事务，而我此刻对此事绝不能心不在焉。但怀着巩固该教会福祉的愿望，我认为需要经过长久慎重的考虑，以便能通过基督启示给我们的、可堪任命的人来安排其事务，从而巩固其良好状况。因此，贤弟目前应专心致志于拯救灵魂，使我对你的看法因你工作的成效而更加坚定。你关于执事伯多禄所写的一切，现已由前执政官德奥多禄向我们禀明。所以，既然我知道他对你忠心，并且如你所证，他关心教会的利益，他就无需畏惧任何人的反对或敌意，而应继续造福教会、服务天主，当他感到别人对他怀有怨恨时，反而要更加警醒；这样一来，他们就没有丝毫能力加害于他。此外，贤弟今后对他不应再有任何疑虑，因为任何隐秘的举动都不会对我产生影响。\par

\SourceRecord{Translated by James Barmby.；Nicene and Post-Nicene Fathers, Second Series；Vol. 12.；Edited by Philip Schaff and Henry Wace.；Buffalo, NY: Christian Literature Publishing Co.,；1895.；<http://www.newadvent.org/fathers/360202015.htm>.}

% structure: p0001 source=h1 target=section reason=page-title

\section{第二卷 第十八封信}

\Emph{致萨洛纳主教纳塔利斯}。\par

额我略致纳塔利斯等。\par

至爱的兄弟，我从许多由你城市而来的人获悉，你忽略了你的牧职责任，全心事奉宴饮。若非我亲身经历你的行事所证实，我本不该相信此一报告。因为你全然不致力于阅读，全然不注意劝勉，甚至对教会秩序的本身用途与目的一无所知，证据就在于你不知道如何对置于你之上的人保持崇敬。因为，当你已被我们可敬记忆中的前任书面禁止在心中保留对你总执事奥诺拉图斯长期不快的痛苦，而且此事也被我自己明确禁止时，你却置天主的命令于不顾，轻视我们的书信，试图以诡计，在提拔他至更高尊位的幌子下，贬黜上述你的总执事奥诺拉图斯。于是设计了这样一个计谋：他被撤去总执事职位后，你就可以召来另一个会顺应你生活方式的人；据我看来，上述那人之所以令你不悦，别无原因，只因他阻止你将圣器与圣衣赠予你的亲属。此事，我现在以及我们可敬记忆中的前任从前都希望接受一次准确的调查；但你因自知所为，一直拖延不派一位受训为案件审讯的代表前来。因此，你的弟兄之爱，即使在如此多次的劝诫之后，应当忏悔你错行的错误，并在收到我的信后立即恢复上述奥诺拉图斯的职位。倘若你拖延此事，须知此圣座授予你的白羊毛披肩的用法将被收回。但若你在失去白羊毛披肩后，仍坚持你的藐视（不服从），须知你将被剥夺主的体血的共融。在这之后，我们需要更全面地调查对你的指控，并以极其审慎和调查考虑你是否还能保留你的主教职。那违背正义规则而同意被提升到他人位置上的人，我们亦从上述总执事的尊位上予以罢免。并且，倘若他再敢在这同一职务上行使职务，让他知道他将被剥夺圣体的共融。因此，至爱的兄弟，你绝不要再激怒我们，免得我们在以爱德态度对待你时你藐视我们，届时你将在我们的严厉中感到我们极其强硬。因此，在恢复总执事奥诺拉图斯到他的职位后，速速派一位受训于案件的人到我们这里，他能通过他的陈述向我显示此事应如何公平地进行。因为我们已经命令上述总执事到我们这里来，以便在听取双方的陈述后，我们能做出任何公正且令全能天主喜悦的决定。因为我们不因私人感情而维护任何人，而是在天主的助佑下，持守正义的规则，搁置对任何人的情面。\par

\SourceRecord{Translated by James Barmby.；Nicene and Post-Nicene Fathers, Second Series；Vol. 12.；Edited by Philip Schaff and Henry Wace.；Buffalo, NY: Christian Literature Publishing Co.,；1895.；<http://www.newadvent.org/fathers/360202018.htm>.}

% structure: p0001 source=h1 target=section reason=page-title

\section{第二卷，第十九封信}

\Emph{致\Place{达尔马提亚}所有主教。}\par

\Person{额我略}致在整个\Place{达尔马提亚}的众主教。\par

虽然我们渴望通过书信的往来频繁地探望你们的兄弟之情，但当某些特殊情况需要我们关注时，我们愿意借此机会同时履行两项职责：既在探望中振作兄弟般的灵魂，又准确解释那些提请注意的事项，以免因无知而使心灵困惑。如今，我们的兄弟、\Place{萨洛纳}城的主教\Person{纳塔利斯}，希望将总执事\Person{奥诺拉图斯}擢升到司铎品位，而\Person{奥诺拉图斯}却拒绝被提升到更高的品位，他以所呈递的请愿书要求我那位神圣纪念的前任教宗，不应违背其意愿被提升。他声称此事并非为了促进他，而是出于对他的不满。因此，我那位神圣纪念的前任教宗写信给我们的兄弟和主教同事\Person{纳塔利斯}，禁止他违背\Person{奥诺拉图斯}的意愿提升他，或在他心中保留对他所怀的不满之痛。当我们也对这位\Person{纳塔利斯}施加同样的禁令时，他不仅无视天主的命令，而且藐视我们的信函，据说他狡猾地试图以降级的方式贬黜上述总执事，却借口提拔他到更高的尊位，这有违惯例。由此策划的是，在将他从总执事职位上撤除后，他可能召请另一个人来代替被撤职的总执事履行职务。现在，我们认为这位\Person{奥诺拉图斯}可能是因阻止主教将圣器给予他的亲属而招致主教的不满：我的前任教宗（神圣纪念者）从前和我现在都希望准确调查此事；但他（指主教）自知所为，却拖延派遣代表接受审判，唯恐其行为的真相暴露。因此，鉴于他已被多次书面劝诫，却始终顽固，我们已安排通过本函持信人再次送信劝诫他，以便本函持信人一到，他立即将总执事\Person{奥诺拉图斯}收回原位。如果他仍心怀硬化，乖戾地拖延不恢复其职位，我们命令，因他多次表现出的顽抗，剥夺他曾由本座授予的披带使用权。但是，如果在丧失披带后，他仍然坚持同样的顽固，我们命令禁止他分领主的身体和血。因为应当如此：他卑视那些以爱德接近他的人，现在却要面对公义中的严厉惩罚。因此，我们现在也不偏离那被上述主教所藐视的公义之道；但是，当那位在我们看来其罪责尚未明显的人被恢复原位时，我们责令\Person{纳塔利斯}主教派遣一位有训示的代表前来，能够通过其陈词向我们证明该主教的正直意图。因为我们也已命令上述总执事前来我们这里，以便在听取双方陈述后，我们决定任何公正的事，任何全能的天主所悦纳的事。因为我们不因私爱而偏袒任何人，而是依靠天主的助佑，不徇情面地持守公义之规。\par

\SourceRecord{Translated by James Barmby.；Nicene and Post-Nicene Fathers, Second Series；Vol. 12.；Edited by Philip Schaff and Henry Wace.；Buffalo, NY: Christian Literature Publishing Co.,；1895.；<http://www.newadvent.org/fathers/360202019.htm>.}

% structure: p0001 source=h1 target=section reason=page-title

\section{卷二，第二十封信}

\Emph{致副助祭\Person{安多尼诺}}。\par

\Person{额我略}致\Person{安多尼诺}等。\par

\Place{撒罗纳}教会总执事\Person{翁诺辣多}，在他所寄的请愿书中，从我蒙福记忆中的前任那里要求，他绝不应被其主教强迫违背自己意愿，以违背惯例的方式升任更高的品秩。\par

\EditorialNote{以下叙述后续过程，几乎与第十九封信逐字相同。}\par

因此我们以为宜以本命令之权威支持你，以便你前往\Place{撒罗纳}，至少尝试以劝诫劝导我们的兄弟及同主教\Person{纳塔利斯}（他已被如此多的信函劝诫），立即恢复上述\Person{翁诺辣多}之职位。但若他如往常一样，顽固地拖延不办，则以宗座权威禁止他使用由本座所赐之\OriginalTerm{披带}{pallium}。但若即使在失去披带后，你发现他仍坚持同样执拗，你应剥夺该主教参与圣体圣事之权利。此外，那违背正义准则、同意升任他人职位之人，我们命令免去其同一总执事之尊位。若他胆敢继续在同一地方履行职务，我们剥夺其参与圣体圣事之权利。因为，他以仁爱接近时，将正义视为无物，理应面对正义之严厉。因此，当总执事\Person{翁诺辣多}恢复其职位后，应由上述主教在你推动下，派人携带指示前来，此人能以其陈述向我们证明该主教之意愿是或曾是正当的。\par

\EditorialNote{以下内容与第十九封信结尾完全一致。}\par

至于我们的兄弟及同主教\Person{马尔科}，你应留意使他找到担保人，以便他尽快来到我们这里，为使他毫无迟延或耽搁，向我们报告其作为，并因此得以安全返回其住所。\par

\SourceRecord{Translated by James Barmby.；Nicene and Post-Nicene Fathers, Second Series；Vol. 12.；Edited by Philip Schaff and Henry Wace.；Buffalo, NY: Christian Literature Publishing Co.,；1895.；<http://www.newadvent.org/fathers/360202020.htm>.}

% structure: p0001 source=h1 target=section reason=page-title

\section{第二卷，第二十二封信}

\Emph{致伊利里库姆所有主教。}\par

额我略致所有主教等。\par

若维护古代惯例的秩序，便既为我们因你们的谨慎而喜乐，也使你们的兄弟情谊在自身的祝圣中稳妥。既然我们通过你们派遣的司铎\Person{马克西米安}和执事\Person{安德肋}所呈的信函得知，你们全体的一致同意和最宁静的君主的意愿，已集中于我们的兄弟和同为主教的\Person{若望}身上，我们感到极大的欢欣，因为在天主的指引下，这样一个被众人判断为称职的人被晋升到了司祭的职位。因此，依照你们的请求，我们以我们同意之权威，确认我们上述的兄弟和同为主教的\Person{若望}在他已被祝圣的司祭职中，并因授予他披肩而宣布我们认可他的祝圣。又因我们按惯例已将代牧职权交托给他，我们必须告诫你们的兄弟情谊，务须毫不迟疑地在属于教会秩序和正确纪律之事上，以及其它不被教规所禁止的事情上服从他；这样，你们选举他的正确判断便因你们所显示的服从而得以显明。\par

\SourceRecord{Translated by James Barmby.；Nicene and Post-Nicene Fathers, Second Series；Vol. 12.；Edited by Philip Schaff and Henry Wace.；Buffalo, NY: Christian Literature Publishing Co.,；1895.；<http://www.newadvent.org/fathers/360202022.htm>.}

% structure: p0001 source=h1 target=section reason=page-title

\section{卷二，信函二十三}

\Emph{致若望主教}\par

额我略致伊利里库姆的普里玛·犹斯提尼亚纳的主教若望。\par

很明显，一个人当选时众人意见一致，是品德的一个明证。既然我们从弟兄和同为主教者那里收到的报告宣称，整个会议的全体一致和最宁静的君主的意愿都召唤您登上司牧之位，我们便以极大的欢欣向全能的天主我们的造物主感恩，祂使您过去的生活和行动如此值得称道，以至于（这极其令人称赞）您赢得了所有人的认可。我们也与他们完全同意关于贤弟的人格。我们恳求全能的天主，既然祂的恩宠选择了您的爱德，就愿祂在一切事上保护您。我们按照惯例将\OriginalTerm{白羊毛披肩}{pallium}送给您，并更新我们的委托，任命您为\OriginalTerm{宗座代理}{vicar of the Apostolic See}，劝告您如此温和地对待您的属下，使他们被激发去爱您而非惧怕您。如果他们中有人因过失需要惩戒，您要注意纠正他们的过错，但绝不可从心中丢弃父爱。要警醒并勤勉地照顾托付给您的羊群，在纪律的热忱上严格，使潜伏的狼不能得逞扰乱主的羊圈，或有机可乘伤害羊群。要全心全意去赢得灵魂归于我们的天主；要知道我们接受牧者的名号不是为了安息，而是为了辛劳。那么，让我们在工作中显明我们的名号所意指的。如果我们正确权衡司祭职的特权，对于勤勉尽职的人它是尊荣，对于疏忽的人它必然是负担。因为，在天主眼中，此名使那些为拯救灵魂而辛劳勤奋的人通向永恒的荣耀，而对于闲散怠惰的人则导致惩罚。通过我们的口舌，托付给我们的百姓应得知另一生命的存在。贤弟的教导对于他们当是可喜的激励，而你的生活则是可效法的榜样。贤弟的讲道应向他们揭示该爱什么、该怕什么，你的效力应以此方式收获永赏的果实。但你的谨慎关怀尤其要约束你，决不可进行任何不合法的祝圣；每逢有人晋升到圣职，或可能到更高品级时，应按功绩而非贿赂或请求而祝圣。在任何祝圣中，绝不要让任何考虑以任何方式悄然影响贤弟，以免你陷入（愿天主阻止）\OriginalTerm{买卖圣职异端}{simoniacal heresy}的罗网。\Quote{人纵然赚得了全世界}，正如真理所说，\Quote{却赔上了自己的灵魂，为他有什么益处？}\BibleRef{谷 8:36}因此，我们必须在一切事上注目天主，轻视暂时和可朽之物，将我们心中的渴望导向永恒的美善。我完全不愿接受贤弟的礼物，因为我们从被掠夺和受苦的弟兄手中接受礼物，这似乎极不合适。但你的使者用另一个理由说服了我，将礼物献给贤弟的奉献不可被拒绝的人。为此你首先应当研究：如何准备不朽的礼物，献给那将要来临的灵魂的审判者，以便祂既垂顾你因有益的劳作，也同样垂顾我们因我们的劝勉。\par

\SourceRecord{Translated by James Barmby.；Nicene and Post-Nicene Fathers, Second Series；Vol. 12.；Edited by Philip Schaff and Henry Wace.；Buffalo, NY: Christian Literature Publishing Co.,；1895.；<http://www.newadvent.org/fathers/360202023.htm>.}

% structure: p0001 source=h1 target=section reason=page-title

\section{第二卷 第二十六封书信}

\Emph{致\Person{若望}主教。}\par

\Person{额我略}致\Person{若望}等。\par

既然我们已将视察\Place{那不勒斯}教会的工作委托给了我们的弟兄及同为主教的\Person{保禄}，因此，望您的仁爱不推辞承担\Place{内佩辛}教会的视察工作，为使根据逾越节庆典的要求，神圣礼仪所要求的一切，都能藉着您的操作，在各方面圆满达成。因此，在我们能考虑如何对待前述我们的弟兄及同为主教之前，望您的仁爱努力在所有事上显出如此熟练和警觉，以致上述主教的缺席完全不致被感觉到。\par

四月，第十\OriginalTerm{小纪}{Indiction}。\par

\SourceRecord{Translated by James Barmby.；Nicene and Post-Nicene Fathers, Second Series；Vol. 12.；Edited by Philip Schaff and Henry Wace.；Buffalo, NY: Christian Literature Publishing Co.,；1895.；<http://www.newadvent.org/fathers/360202026.htm>.}

% structure: p0001 source=h1 target=section reason=page-title

\section{卷二，书信二十七}

\Emph{致元老院议员\Person{鲁斯蒂齐亚娜}}\par

\Signature{额我略致\Person{鲁斯蒂齐亚娜}等}\par

收到您的来信后，得知您平安的消息令我欣慰，希望主以祂的仁慈保护并引导您的生命与作为。但我深感不解，为何您改变了您关于前往圣地朝圣之善工的心愿与誓愿；须知，当藉造物主的恩赐在心中孕育任何善念时，必须立即以热忱加以实行，以免那狡猾的阴谋者（指魔鬼）在设法诱捕灵魂之后，随后又设置障碍，使心灵因俗务缠身而衰弱，无法实现其渴望。因此，您有必要预见一切阻碍虔诚计划的因素，竭尽全力渴慕善工的果实，如此方能在现世安然生活，并在来世获得天国。至于您来信提及\Person{帕西武斯}曾试图散布一些对您不利的诽谤，另一方面请您思考，最虔诚的皇帝们不仅不愿听从这些诽谤，甚至严厉接待了诽谤的制造者；请您将灵魂的全部希望转向那一位，祂在此世强有力地阻止人造成他们所欲求的那么多伤害，以便祂能以自己臂膀的抵挡击退人的恶意，并仁慈地粉碎他们的企图，一如祂素常所行。我恳求您代我问候荣耀的\Person{阿皮奥}大人、\Person{欧瑟比亚}夫人、\Person{欧多克西乌斯}大人和\Person{格列戈利亚}夫人。\par

\SourceRecord{Translated by James Barmby.；Nicene and Post-Nicene Fathers, Second Series；Vol. 12.；Edited by Philip Schaff and Henry Wace.；Buffalo, NY: Christian Literature Publishing Co.,；1895.；<http://www.newadvent.org/fathers/360202027.htm>.}

% structure: p0001 source=h1 target=section reason=page-title

\section{第二卷 第二十九封信}

\Emph{致\Person{莫里利乌斯}和\Person{维塔利亚努斯}}。\par

\Person{格列高利}致\Person{莫里利乌斯}和\Person{维塔利亚努斯}，\OriginalTerm{军长}{magistris militum}。\par

收到阁下您的来信后，我们感谢天主，因为我们确信您平安无事；我们为您周详的准备感到极大的喜乐；您所提及的事也已立即办妥。但是，在您的人到达之后，尊贵的\Person{阿尔迪奥}写信给我们，说\Person{阿里乌尔夫}已经近在咫尺，我们担心派往您那里的士兵可能会落入他手中。然而在这里，只要天主赐予援助，我们的儿子，光荣的\OriginalTerm{军长}{magister militum}，也已准备好对抗他。但是，如果敌人亲自向这里推进，也请阁下您一如既往地，在他后方尽力而为。因为我们寄望于全能天主的德能，以及有福的\Person{伯多禄}本人——\OriginalTerm{宗徒之长}{Prince of the apostles}——的德能，他想要在他的瞻礼日流血，但愿他也会发现他毫不迟延地反对他。\par

\SourceRecord{Translated by James Barmby.；Nicene and Post-Nicene Fathers, Second Series；Vol. 12.；Edited by Philip Schaff and Henry Wace.；Buffalo, NY: Christian Literature Publishing Co.,；1895.；<http://www.newadvent.org/fathers/360202029.htm>.}

% structure: p0001 source=h1 target=section reason=page-title

\section{第二卷，第三十书}

\Emph{致\Person{莫里略}与\Person{维塔利亚努斯}。}\par

\Person{额我略}致军长\Person{莫里略}与\Person{维塔利亚努斯}。\par

我们曾通过我们的儿子\Person{维塔利亚努斯}以言词和信件恳求您的荣耀，嘱咐您与他沟通。但在正月十一日，\Person{阿里乌尔夫}送来了这封信，我们现转交给您。因此，当您阅读后，请查看\Place{苏阿纳}的人民是否坚守了他们向共和国承诺的忠诚，并从他们那里取得足够的、您可以信赖的人质；此外，再用誓言重新约束他们，归还您作为抵押从他们那里取走的东西，并通过您的谈话使他们恢复理智。但是，如果您明确查实他们已经与\Person{阿里乌尔夫}就投降事宜进行了交涉，或者至少已向他提供了人质，正如我们转交给您的\Person{阿里乌尔夫}的信件所暗示的那样，那么（在经过有益审慎的考虑后，以免您的灵魂或我的灵魂因誓言而负罪），请采取您认为对共和国有利的任何行动。但请您的荣耀如此行事，既不做出任何让我们的对手可以指责我们的事情，也不（愿主阻止）忽略任何共和国利益所要求的事项。此外，我荣耀的儿子们，请急切关注，因为据我所知，敌军已集结一支军队，据称驻扎在\Place{纳里纳}；如果（天主对他发怒）他决定转向此地，那么请你们在主可能协助的范围内劫掠他的据点，或者至少让您派遣的人仔细安排夜间警戒，以免任何不幸事件的消息传到我们这里。\par

\SourceRecord{Translated by James Barmby.；Nicene and Post-Nicene Fathers, Second Series；Vol. 12.；Edited by Philip Schaff and Henry Wace.；Buffalo, NY: Christian Literature Publishing Co.,；1895.；<http://www.newadvent.org/fathers/360202030.htm>.}

% structure: p0001 source=h1 target=section reason=page-title

\section{第二卷 信件三十二}

\Emph{致彼得，西西里的副执事。}\par

格列高利致彼得，等等。\par

据护守者罗玛诺斯报告，我得知天主婢女隐修院——它位于莫诺特乌斯农场——因属于维拉诺瓦教会的一个农场（据称已租赁给该隐修院）而受到本教会的侵害。若情况属实，请贵经验将农场归还，并退还你在两个小纪期间从那农场征收的款项。此外，既然有许多犹太人居住在教会的庄园上，我希望，若他们中有人愿意成为基督徒，可酌情减免其部分应缴款，以便其他人受此恩惠激励，也产生同样的愿望。\par

那些因年老不再生育的母牛，或似乎完全无用的公牛，应当卖掉，以便至少从其价格中获取一些收益。至于我们饲养的既不盈利的马群，我希望全部疏散，只保留四百匹较年轻的用于繁殖；这四百匹应当分给农民——每户若干匹——以便他们能在未来几年从这些马匹中向我们提供一些回报：因为我们为牧马人花费六十\OriginalTerm{金币}{solidi}，却从这些马群中得不到六十便士，这实在难以承受。那么，请贵经验如此进行：一部分分给所有农民，其余疏散并变现。但要在我们的所有产业中与牧马人这样安排，使他们能够通过耕种土地获得一些收益。所有可以在锡拉库萨或帕诺尔莫被教会索要的器具，必须在它们因年久而完全朽坏之前卖掉。\par

天主的仆人西里亚科斯兄弟到达罗马后，我仔细询问他是否曾就某位妇女案件中收受贿赂一事与你商议过。这位兄弟说，他是从你的讲述中得知案情的，因为你曾委托他查明受托支付贿赂的人是谁。我相信了此事，立刻亲切地接纳他进入眷顾，将他引荐给民众和神职人员，增加了他的\OriginalTerm{薪俸}{stipend}，将他提升为护守者中的较高等级，并在众人面前赞扬他的忠诚，因他在你的事奉中如此忠信；因此我已将他派回你处。然而，既然你十分匆忙，而我虽在病中却渴望见你，请你留下一位你已充分考验过的人接替你在锡拉库萨地区的职务，自己则赶快到我这里来，以便，如果全能的天主乐意，我们可以共同商议是否应由你本人返回那里，或另派他人接替你的位置。同时，我已派遣公证员贝内纳图斯前往帕诺尔莫地区接替你在教产中的职务，直到全能的天主命定祂所喜悦的事为止。\par

我严厉斥责了罗玛诺斯的轻率，因为在他管理的\OriginalTerm{客栈}{xenodochium}中，我现已发现，他更注重自己的利益而非天国的赏报。因此，如果你认为合适，就把他留在你的位置上。你要设法用警告和劝诫来加强他，使他能仁慈细心地对待\OriginalTerm{农民}{rusticos}，并在外乡人和本地人面前表现出改变和勤勉。然而，我这样说并非选定某人，而是将此留给你判断。我已选定帕诺尔莫地区接替你位置的人，这对我足够了；我希望你亲自为锡拉库萨地区提供一个人选。当你来时，请带上有关系或来自安托尼努斯财产的款项和\OriginalTerm{饰物}{ornamenta}。也带来你已征收的第九和第十小纪的款项，以及你所有的账目。请务必，若天主乐意，在圣西彼廉周年纪念日之前渡海到达本城，以免（天主禁止）在那个季节总是威胁海洋的星宿带来任何危险。\par

此外，我希望你知道，我心中不无内疚，因为我曾对天主的仆人普雷蒂奥苏斯严厉发怒（并非因他的严重过失），并使他悲伤痛苦地离开了我。我曾写信给主教大人，请求他，若愿意，就把那人送到我这里；但他完全不愿意。我不应使主教大人困扰，我也不能这样做；因为他既从事天主的事业，就应得到安慰的支持，而非被苦楚压垮。但那位普雷蒂奥苏斯，据我所闻，因不能回到我这里而极为痛苦。然而，如我所说，我不能使主教大人困扰，他不愿送他来，我在这两者之间犹豫不决。那么，若你小小的身躯中有更大的智慧，请设法处理此事，使我遂愿，且不使主教大人受到困扰。但若你见他有所困扰，就不要再提此事。然而，我对此感到不悦：他绝罚了尤塞比乌斯大人，一位如此高龄且身体很差的人。因此，你有必要私下对那位主教大人说，不要急于宣布判决，因为须由判决决定的案件，必须事先以谨慎和非常频繁的考量来权衡。\par

当招募官到来时——据我所闻，他们已在西西里招募新兵——请你命令你的代管人向他们赠送一些小礼物，以使他们对他怀有好感。但你在离开之前，也要按照古老习俗，向裁判官（\OriginalTerm{裁判官}{prætor}）的属员赠送一些东西；不过，要经由你留下代替你的人之手，以赢得他们对他的好感。此外，以免我们显得对他们有失礼貌，请指示你的代管人，在我们已向你阁下（\OriginalTerm{阁下}{Experience}）下达的、关于应给予任何个人或隐修院之物的命令上，完全遵行。等你到来时，我们将仰赖天主的助佑，共同考虑应如何安排这些事项。我派人送来、要经由你分施给穷人的三百索利多（\OriginalTerm{索利多}{solidi}），我认为不应托付他们自行处理。那么，让他们遵行我已提及的那些针对特定地点和人员的指示吧。\par

我现在记得此前曾写信说过，根据监护人（\OriginalTerm{监护人}{defensor}）安托宁的报告，我们欠隐修院或其他人的遗赠，应当按所指定的支付。我不知道为何你阁下（\OriginalTerm{阁下}{Experience}）迟迟未完成此事。因此，我们希望你用教会的款项，将我们应承担的这些遗赠份额全部付清，这样当你来我处时，穷人不会在你身后留下对你的哀叹。同时，也请你把已经找到的、与同一安托宁财产有关的契据带来。\par

据罗曼努斯报告，我得知雷登普图斯的妻子临终时口头吩咐，将一只银贝壳卖掉，所得金额分给她的被释奴，又将一只银盘留给了某所隐修院；关于这两项遗赠，我们希望她的愿望完全实现，以免因小事陷入更大的罪过。\par

此外，据隐修院院长马里尼亚努斯报告，我得知在普雷托里安隐修院（\OriginalTerm{普雷托里安隐修院}{Prætorian Monastery}）的建筑甚至尚未完工一半：既然如此，除你阁下（\OriginalTerm{阁下}{Experience}）的热忱外，我们还能称赞它什么呢？但即使现在，愿此劝诫激励你；并请尽力投入同一隐修院的建造。我曾说过不必为他们提供建筑费用，但我并未禁止他们建造隐修院。但请如此行事，务必命令你将在帕诺尔穆斯代你职位的人，以教会收入的费用建造同一隐修院，这样我便不会再收到隐修院院长的私下投诉。\par

再者，我得知你知道在农庄上有些事物——甚至相当数量——属于他人；但因某些人的恳求或出于胆怯，你不敢将它们归还给原主。但如果你真是基督徒，就应更惧怕天主的审判而非人的声音。请注意，我不停地就此事劝诫你；如果你忽视纠正，你也会听到我的声音作证反对你。如果你发现任何敬畏天主的平信徒，可以接受剃发并成为院长（\OriginalTerm{院长}{rector}）手下的代理人，我完全同意。也需向他们发送信件。\par

关于学者（\OriginalTerm{学者}{scholasticus}）康米苏斯之子的事，你已征询意见；看来他所要求的在法律上并不公正。我们不愿加重穷人的负担，使其受损；但既然他为此事费了心，我们希望给他五十索利多（\OriginalTerm{索利多}{solidi}），这当然要记在你的账目上。至于你在普罗基苏斯一案中为教会事务支出的费用，要么你从他的收入中自行补偿，要么若他的收入明显不足以偿还，你必须在这里从执事那里收取应得的款项。但不得妄言副执事格拉修斯之事，因为他的罪行要求直到生命末刻最严厉的补赎。\par

此外，你送给我一匹劣马和五匹良驮。那匹劣马我不能骑，它实在太差；那些良驮我不能骑，因为它们是驴。但若你愿意使我们满意，我们请求你给我们一些合适的。我们希望你给隐修院院长欧西比乌斯一百金索利多（\OriginalTerm{索利多}{solidi}），这当然要记在你的账目上。我们得知曾担任萨姆尼乌姆法官的西西尼乌斯在西西里正遭受严重贫困，我们希望你每年供给他二十十升量（\OriginalTerm{十升量}{decimates}）酒和四索利多。一位宗教人士（\OriginalTerm{宗教人士}{religiosus}）阿纳斯塔修斯据说住在帕诺尔穆斯城附近圣阿格娜的祈祷所，我们希望给他六金索利多。我们也希望六索利多——记在你的账目上——给予院长乌尔比库斯的母亲。关于天主的婢女奥诺拉塔之事，我的看法是这样：你应当带着她全部明显存在于劳里努姆主教若望任职之前的财产来见我。但同一婢女应与她的儿子同来，以便我们与她交谈，并做任何中悦天主的事。我们希望能将安托宁遗产中的《七书》（\OriginalTerm{七书}{Heptateuch}）卷册赠予普雷托里安隐修院，其余书籍由你带至此地。\par

\SourceRecord{Translated by James Barmby.；Nicene and Post-Nicene Fathers, Second Series；Vol. 12.；Edited by Philip Schaff and Henry Wace.；Buffalo, NY: Christian Literature Publishing Co.,；1895.；<http://www.newadvent.org/fathers/360202032.htm>.}

% structure: p0001 source=h1 target=section reason=page-title

\section{第二卷·第三十三篇书信}

\Emph{致 \Person{犹斯蒂诺} 总督}\par

\Person{额我略} 致 \Person{犹斯蒂诺} 等。\par

古代仇敌有一特色，就是对于那些因天主抗拒他而无法诱使其行恶的人，他暂时用虚假的报告损害他们的名声。既然关于我们的兄弟和同为主教的良有一种恶毒的谣言，散播了一些与他的司铎身份不符的事，我们就进行了严格和漫长的调查，看是否属实，并在所述之事上发现他并无过错。但为了使任何事似乎都没有遗漏，也为了不让任何可能的疑虑存留于我们心中，我们还在真福伯多禄最神圣的躯体前，令他额外作了严格的宣誓。他这样做之后，我们极其欢欣，因为通过这种证明，他的无辜显然彰显了出来。因此，请尊荣以全部爱德接纳前述之人，向他表示对司铎应有的敬重；也莫让任何疑虑存留于你心中，关于那些他现已澄清的指控。然而，你应在各方面紧紧依附上述主教，好让人看出你适当地、体面地在他身上尊敬天主，因为他是天主的服务者。\par

\SourceRecord{Translated by James Barmby.；Nicene and Post-Nicene Fathers, Second Series；Vol. 12.；Edited by Philip Schaff and Henry Wace.；Buffalo, NY: Christian Literature Publishing Co.,；1895.；<http://www.newadvent.org/fathers/360202033.htm>.}

% structure: p0001 source=h1 target=section reason=page-title

\section{卷二，第三十四封信}

\Emph{致\Person{Maximianus}，\Place{Syracuse}主教。}\par

\Person{Gregory}致\Person{Maximianus}等。\par

我记得曾屡次劝告你，决不可仓促判决。而如今，我得知你的尊驾在愤怒之中绝罚了最可敬的院长\Person{Eusebius}。我深感惊异的是，他往日的言行、年迈之龄以及久病不愈，竟都不能化解你的怒气。因为，无论他的过失如何，疾病的折磨本身也足以作为他的鞭笞了。一个已经受到天主管教的人，再加上人间的惩罚实属多余。不过，也许在这件事上你被允许做得过分，是为了让你在其他次要人物身上更为谨慎，并在你欲以判决打击任何人时深思熟虑。然而，现在请以与你先前激怒他时同等的温和来安慰此人，因为极其不公的是，那些最爱你的人却无故发现你对他们最为严厉。\par

\SourceRecord{Translated by James Barmby.；Nicene and Post-Nicene Fathers, Second Series；Vol. 12.；Edited by Philip Schaff and Henry Wace.；Buffalo, NY: Christian Literature Publishing Co.,；1895.；<http://www.newadvent.org/fathers/360202034.htm>.}

% structure: p0001 source=h1 target=section reason=page-title

\section{第二卷，第三十六封信}

\Emph{致院长欧瑟伯。}\par

额我略致欧瑟伯等。\par

愿你的爱德相信我，我为你的悲伤深感悲伤，仿佛我自己也在你身上受了委屈。但后来我得知，即使我们可敬的弟兄及同为主教的马克西米安已恢复对你的眷顾及共融，你的爱德仍不愿接受他的共融，那时我才知道先前所发生的事是公正的。天主的仆人在患难之时应当显出谦卑；但那些向长上自高自大的人，表明他们不屑作天主的仆人。的确，他先前所行之事本不应做；但你仍应谦卑地接受；再者，当他恢复对你的眷顾时，你本该心怀感激地回应。因为你没有这样做，我觉得我们无论如何都有流泪的理由。因为向尊重我们的人表示谦卑，算不得什么大事，连世俗之人也能做到；但我们尤其应当向那些使我们受苦的人表示谦卑。因为圣咏作者说：\BibleQuote{\Emph{求你垂视我在仇人前的谦卑}} \BibleRef{咏 9:14}。如果我们甚至不愿向我们的长辈谦卑，我们过的又是什么生活？所以，最亲爱的儿子，我恳求你，让一切苦毒从你心中消失，免得末日临近，古老的仇敌因不和的罪愆，拦阻我们进入永恒国度的道路。此外，我们已令副执事伯多禄将一百\OriginalTerm{索利多}{solidi}交给你的爱德，恳请你毫无不悦地收下。\par

\SourceRecord{Translated by James Barmby.；Nicene and Post-Nicene Fathers, Second Series；Vol. 12.；Edited by Philip Schaff and Henry Wace.；Buffalo, NY: Christian Literature Publishing Co.,；1895.；<http://www.newadvent.org/fathers/360202036.htm>.}

% structure: p0001 source=h1 target=section reason=page-title

\section{第二卷，书信三十七}

\Emph{致斯奎拉切（卡拉布里亚的斯奎拉切）主教若望。}\par

额我略致若望等\par

我们牧职的关怀提醒我们，为那些失去牧者的教会任命它们自己的主教，使他们能以牧灵的关切管理主的羊群。因此，我们认为必须任命你，若望，即已被敌人占领的\Place{里斯城}（\Place{里斯}，今\Place{阿莱西奥}？）的主教，为\Place{斯奎拉切}教会的枢机，使你得以继续你曾承担的灵魂牧养，顾及将来的报偿。尽管你因入侵的敌人被驱离自己的教会，必须管理另一个现在没有牧者的教会，但条件必须是：倘若前城从敌人手中解放，并在天主的保护下恢复原状，你应返回你最初被祝圣的教会。然而，若上述城市继续遭俘虏之灾，你必须留在我们使你入驻的这个教会。此外，我们命令你绝不可举行不合法的圣秩授予，也不可允许任何重婚者、或娶了非处女为妻者、或不识字者、或身体任何部分残缺者、或做过补赎者、或应服某种劳役者，领受圣秩。你若发现此类人士，绝不可为他们晋秩。一般非洲人，以及不明身份的陌生人，若申请教会品级，绝不可接受，因为有些非洲人是摩尼教徒，有些曾重领洗礼；而许多陌生人虽居小品位，却证明曾冒充更高的品位。我们也劝告你明智地看守托付给你的灵魂，更专注于赢得灵魂，而非现世的利益。殷勤地保管和处置教会的财产，使将来的审判者，当祂来审判时，能认可你在各方面配得地执行了你所承担的牧者职务。\par

\SourceRecord{Translated by James Barmby.；Nicene and Post-Nicene Fathers, Second Series；Vol. 12.；Edited by Philip Schaff and Henry Wace.；Buffalo, NY: Christian Literature Publishing Co.,；1895.；<http://www.newadvent.org/fathers/360202037.htm>.}

% structure: p0001 source=h1 target=section reason=page-title

\section{\WorkTitle{第二卷，第四十一封信}}

\Person{额我略}致\Person{阿里米努姆（里米尼）}主教\Person{卡斯托里乌斯}。\par

从附件呈文的内容可以看出，\Person{卢米诺苏斯}，\Person{圣安德肋}及\Person{圣多默}隐修院（位于\Place{阿里米努姆（里米尼）}）院长，向我们倾吐了何等可悲的恳求。关于此事，我们敦促阁下：该隐修院院长去世时，你的教会不得以任何借口干涉该隐修院现有或将来获得的财产的登记或保管。我们希望你只任命该团体全体会众共同同意的人为该隐修院院长，称其品德相称、适于修院纪律。此外，我们完全禁止主教在该处举办公开弥撒，以免在天主仆人的退修处引发群众集会，也免得妇女过于频繁出入（但愿天主禁止）成为丑闻的缘由，尤其对较单纯的灵魂。此外，我们规定，我们所写的这份文件应被仔细遵守，并在以后所有时间内由你和你的继任主教们保持有效且不受篡改；这样，在天主助佑下，你的教会既满足于自己的权利而不越界，同时该隐修院从此只受一般或教会法律权的管辖，免于一切烦恼和骚扰，能以最热忱的心灵完成其神圣工作。\par

\EditorialNote{在一些抄本中，以上信函被替换为以下内容，并附有关于隐修院特权的文件。}\par

\Emph{额我略致阿里米努姆（里米尼）主教卡斯托里乌斯。}\par

从附件呈文的内容可以看出，\Person{卢米诺苏斯}，\Person{圣安德肋}及\Person{圣多默}隐修院（位于\Place{阿里米努姆（里米尼）}）院长，向我们倾吐了何等可悲的恳求。因为从他的叙述我们得知，在非常多的隐修院中，隐修士们遭受了来自高级神职人员的许多偏见和骚扰。因此，阁下的责任在于通过一个健全的安排为他们将来的安宁作出准备，以便那些在天主服务中在此生活的人，在他的恩宠助佑下，能心无纷扰地坚持。但是，为了避免有人以某种理应改正而非延续的习俗为由，试图对隐修士造成任何形式的骚扰，有必要使我们下面所列的事项得到主教弟兄们如此仔细的遵守，以至于将来找不到任何可能引入不安的借口。\par

\Emph{论隐修院特权。}\par

因此，我们以我们的主\Person{耶稣基督}之名，并藉真福\Person{伯多禄}（宗徒之长，我们代替他治理此\Place{罗马教会}）的权威，禁止任何主教或世俗人士今后以任何方式图谋干预隐修院或其附属小室、庄园的收入、财产或文书，或诉诸任何诡计或勒索；但，如果偶然出现关于他们的教会与任何隐修院之间争议土地的案件，且无法友好解决，则应在所选隐修院院长及其他敬畏天主的父亲们面前，以最神圣的福音宣誓，不得故意拖延地终结。同样，在任何团体院长去世时，不得任命外人，只应任命同一团体中弟兄团体自愿一致选举、且选举中无欺骗或买卖圣职的人。但如果他们当中找不到合适人选，则应同样从其他隐修院选举一人予以任命。此外，院长一经任命，任何人不得以任何借口凌驾于他之上，除非（但愿天主禁止）有显而易见的罪行，证明应按神圣法典予以惩处。同样，应遵守以下规则：隐修士未经院长同意，不得为建立其他隐修院、或授予圣秩、或任何神职职务而从隐修院调离。我们还禁止主教为隐修院财产制作教会清册。但若情况需要，某地院长与其他院长就所获财产产生争议，此事也应由他们的商议或裁决终结。同样，在院长去世时，主教不得以任何借口干预隐修院现有、已赠予或将来获得之财产的登记或保管。我们还完全禁止主教在修院中举行公开弥撒，以免在天主仆人的退修处及避难所内制造群众集会的机会，或导致妇女不寻常的进入，这对她们的灵魂毫无益处。主教也不得在那里设置主教座，或拥有任何命令权、或举行任何圣秩授予（即使是最普通的），除非应地方院长的请求；这样隐修士们就能始终在院长的权力之下；任何主教都不得在未经院长证明及许可的情况下，将隐修士扣留在任何教会，或未经许可将其升迁至任何尊严。因此，我们规定，所有主教在将来任何时候都要严格遵守此书面文件，使之保持有效且不受篡改；这样，他们既满足于自己教会的权利而不越界，而隐修院也不受任何教会条件、强迫服务或对世俗政权的任何服从（惟保留教会法律权），而是免于一切烦恼和骚扰，能以最热忱的心灵完成其神圣工作。\par

\SourceRecord{Translated by James Barmby.；Nicene and Post-Nicene Fathers, Second Series；Vol. 12.；Edited by Philip Schaff and Henry Wace.；Buffalo, NY: Christian Literature Publishing Co.,；1895.；<http://www.newadvent.org/fathers/360202041.htm>.}

% structure: p0001 source=h1 target=section reason=page-title

\section{第二卷，第四十二封书信}

\Emph{致院长\Person{卢米诺苏斯}}\par

\Person{格列高利}致\Place{里米尼}的\Person{圣多玛斯}隐修院院长\Person{卢米诺苏斯}。\par

我们很高兴收到你和你的团体的请求，并按照圣教父的规章和法律形式同意你的请求。因为我们已经下令给我们的弟兄和同为主教\Person{卡斯特里奥}一封信，从而完全剥夺了他及其继任者伤害你隐修院的一切权力；使他再也不必来你们中间成为负担，也不得编制隐修院财产的清单，也不得在那里举行任何公开游行；仅保留给他这项司法权：他必须在一位已故院长的位置上，任命一位由团体共同同意选为合适的人。但是，既然这些事情已经完成，你们要勤于天主的工作，并专心致力于祈祷，免得你们似乎不是寻求内心的安宁以便祈祷，而是希望逃避主教对你们过失行为的严格监督。\par

\SourceRecord{Translated by James Barmby.；Nicene and Post-Nicene Fathers, Second Series；Vol. 12.；Edited by Philip Schaff and Henry Wace.；Buffalo, NY: Christian Literature Publishing Co.,；1895.；<http://www.newadvent.org/fathers/360202042.htm>.}

% structure: p0001 source=h1 target=section reason=page-title

\section{第二卷，第四十六封书信}

\Emph{致若望主教。}\par

\Person{额我略}致\Place{辣文纳}主教\Person{若望}。\par

我未能回复尊驾的多封来信，非因我怠惰，实因软弱。由于我的罪，当\Person{阿里乌尔弗}来到\Place{罗马}城，杀害一些人，又伤残另一些人时，我陷入了极大的忧伤，以致患上了绞痛之症。但我深感诧异，为何您至圣圣座对我素有的关怀，竟对此城及我的需要毫无裨益。然而，当您的信函送达我处时，我方知您确实在努力行事，却无人能受您之行动影响。因此，我将此归咎于我的罪：如今我们所关注的此人，既逃避与敌人作战，又禁止我们讲和；然而事实上，即便他此刻允许我们讲和，我们也完全无能为力，因为\Person{阿里乌尔弗}拥有\Person{奥塔尔}和\Person{诺杜尔夫}的军队，他要求先给予他们资助，才肯屈尊与我们谈论和平之事。\par

至于\Place{伊斯特利亚}主教们的事件，我从最虔诚的君王们下达给我的命令中，已得知您在信中所告之一切属实，他们命我暂缓强迫他们。我确实与您同心，并为您在来信中所表现的热忱与激情而大感喜乐，且承认自己在多方面欠您的情。然而，请知悉，我仍将就此同一事，以最大的热忱和自由，不断致函最庄严的主上们。此外，前述最卓越的\Person{罗马努斯·帕特里丘斯}的敌意不应令您动摇，因为既然我们在地位和等级上高于他，便更应以宽容和尊严容忍他的任何轻率行为。\par

然而，若有机会说服他，望您的兄弟之爱努力为之，使我们能与\Person{阿里乌尔弗}达成和议，哪怕仅稍有成效，因为兵丁已从\Place{罗马}城调离，此乃他亲知。但留在此地的\OriginalTerm{狄奥多西派人士}{Theodosiacs}，因未领到薪饷，难以被说服守卫城墙；如此城既被众人遗弃，若无和平，又将如何存续？\par

再者，关于那位从俘虏中被赎回的妇女，您曾写信请我们查问她的出身，我们愿您至圣圣座知悉，一位未知之人不易查寻。至于您所言有人已受祝圣又再受祝圣一事，这极为荒谬，且不符合像您这样秉性之人的考虑，除非能举出某些先例，在审判被指曾行此事之人时应予考量。但您的兄弟之爱切勿怀此见解。因为，正如已受洗者不应再受洗，已受祝圣者亦不能再就同一品级受祝圣。但若某人获得圣职时伴有轻微过失，则应判其为此过失作补赎，但仍保留其圣职。\par

关于\Place{那不勒斯}城，鉴于最卓越的\OriginalTerm{总督}{Exarch}之紧急坚持，我们告知您：据我们所知，\Person{阿里吉斯}已与\Person{阿里乌尔弗}联合，背弃对共和国的忠诚，并对此城图谋不轨；若不速派公爵前往，此城或已可视为沦陷。\par

至于您所言，应施舍给被焚毁的分裂者\Person{塞维鲁}之城，您的兄弟之爱持此看法，是因不知他向我们敌对者所送之贿赂。即便这些贿赂未送，我们也须考虑：怜悯应先施与信友，而后才及于教会的敌人。因为附近有\Place{法努姆}城，其中多人被俘；去年我本欲施舍于彼，但不敢穿越敌人之境。依我之见，您应派\Person{克劳狄}院长携一定数额的钱财前往，赎回他在那里发现的、被扣为奴以索取赎金的自由人，或任何仍被囚之人。至于所送之钱款数额，凡您决定，必合我意。此外，若您正与最卓越的\Person{罗马努斯·帕特里丘斯}商议，允许我们与\Person{阿里乌尔弗}讲和，我准备另遣一人至您处，以便更好地安排赎金事宜。\par

关于我们的兄弟和同袍主教\Person{纳塔利斯}，我一度深感忧虑，因我发现他在某些事上行事傲慢；但既已自改其行，他便胜过了我，也安慰了我的忧苦。与此相关，请告诫我们的兄弟和同袍主教\Person{马尔库斯}：在来见我们之前，先呈报帐目，然后若有必要，再往他处。若我们发现他行为良好，或许有必要将他曾管理的产业归还给他。\par

\SourceRecord{Translated by James Barmby.；Nicene and Post-Nicene Fathers, Second Series；Vol. 12.；Edited by Philip Schaff and Henry Wace.；Buffalo, NY: Christian Literature Publishing Co.,；1895.；<http://www.newadvent.org/fathers/360202046.htm>.}

% structure: p0001 source=h1 target=section reason=page-title

\section{第二卷，第47封信}

\Emph{致\Person{多米尼库斯}主教。}\par

\Person{额我略}致\Person{多米尼库斯}，\Place{迦太基}主教。\par

我们以最大的喜悦收到了您的兄弟情谊的来信，这些信由我们最可敬的弟兄和同为主教的\Person{多纳图斯}和\Person{克沃德乌尔图斯}，以及执事\Person{维克多}与公证员\Person{阿吉莱吉乌斯}之手，稍晚才送达我们。虽然我们曾以为因他们来得迟缓而有所损失，但我们从他们更为丰盛的爱德中得到了收益；因为从时间上的延迟来看，那爱德——我们凭天主的仁慈，通过您对司祭职务的默观、阅读的实践和成熟的年龄，知道它已牢牢植根于您心中——并未中断，反而增长。倘若它并非在您心中拥有极多丰盈的脉源，它就不会如此充沛地从您涌流而出。因此，最圣洁的兄弟，让我们以不可动摇的坚定紧握这众德之母与护卫。不要让诡诈之舌在我们中间削弱它，也不要让古敌的任何陷阱败坏它。因为这爱德联合分裂的，保持联合的。它高举卑微而不增浮夸；它贬抑高傲而不致沮丧。藉此，普世教会——即基督身体的联结——通过心灵对各部分的均衡，尽管因肢体多样而内含不同，仍在其各部分中喜乐。藉此，这些肢体既在他人喜乐中欢欣（虽自身受苦），也在他人忧伤中沮丧（虽自身喜乐）。因为，正如外邦人的导师所证：\BibleQuote{若一个肢体受苦，所有肢体都一同受苦；若一个肢体得荣耀，所有肢体都一同欢乐}，我毫不怀疑您为我等的困扰而叹息，正如我确信我们为您的平安而欢喜。\par

至于您的兄弟情谊因我们的祝圣而与我们一同欢喜，这向我显示了最真诚爱德的情感。但我承认，因默观这服务职分的重担，一股悲伤之力刺透了我的灵魂。因为司祭职的重担是沉重的；司铎必须首先生活，以成为他人的榜样，然后要警惕不因所显示的榜样而心高气傲。他应时常思想宣讲的职务，以极度的恐惧考虑主在即将离去为自己收取王国并将塔冷通交给仆人时所说：\BibleQuote{你们去做生意，直到我回来}（\BibleRef{路 19:13}）。这生意我们确实只有通过生活和言语赢得邻人的灵魂，才算是进行；通过宣讲天国的喜乐来坚固所有在神圣爱德中软弱的人；通过可怕地宣告地狱的惩罚来使倔强和胆怯者屈服；如果我们不因真理而姑息任何人；如果我们献身于天上的友谊，就不畏惧人的敌意。其实，圣咏作者正是这样表现自己，向天主献上了一种祭献，当他说：\BibleQuote{上主，我岂不是恨那恨你的人？我岂不是因你的仇敌而忧愁？我以极恨恨他们，他们成了我的仇敌}（\BibleRef{咏 138:21}）。但面对这重担，我为我的软弱而战栗，并仰望家主在收取王国后回来向我们算账。那时，我以何心面对他的到来，若我从所承担的生意中不给他带来收益，或几乎没有？因此，最亲爱的兄弟，请以你的祈祷帮助我；你见我为自己所恐惧的，你自己也要每日怀着焦虑的害怕思量。因为通过爱德的纽带，我关于自己所说的，也是你所关心的；而我愿你做的，也是我的事。\par

此外，关于您的兄弟情谊所写的教会特权，请毫无犹豫地持守，因为正如我们捍卫自己的权利，我们也尊重多个教会的权利。我既不以偏私给予任何教会超过其应得的，也不因野心之煽动而剥夺任何教会按权利属于它的；但我愿以各种方式尊敬我的弟兄，并相应地努力使每人得以在尊荣中提升，只要彼此之间没有权利上的反对。另外，我因您信使们的行为而大大与您同乐，从中显示您多么爱我，因为您向我派遣了被选的弟兄和儿子们。\par

\Signature{颁于八月第十日前之第十日，第十小纪。}\par

\SourceRecord{Translated by James Barmby.；Nicene and Post-Nicene Fathers, Second Series；Vol. 12.；Edited by Philip Schaff and Henry Wace.；Buffalo, NY: Christian Literature Publishing Co.,；1895.；<http://www.newadvent.org/fathers/360202047.htm>.}

% structure: p0001 source=h1 target=section reason=page-title

\section{第二卷，书信四十八}

\Emph{致科隆布主教}。\par

\Person{额我略}致\Person{科隆布}等。\par

众所周知，最亲爱的在基督内的弟兄，那古老的仇敌，曾以狡诈的劝说将第一个人从乐园的福乐贬谪到充满忧虑的此生，并在他身上将死亡的惩罚加诸于人类，如今以同样的狡诈，为了更容易攫取羊群，企图用注入的毒液来感染上主羊群的牧者，并已将他们视作自己的权利。然而我们，虽不配，却已肩负起宗徒之长\Person{伯多禄}代理下的宗座治理之责，我们的宗座牧职的本身迫使我们要尽己所能抵抗那共同的仇敌。如今，这些信件的携带者\Person{君士坦提乌斯}和\Person{穆斯图鲁斯}，在一份呈递给我们的请愿书中使我们得知，并且\Place{努米底亚}省\Place{普登提亚纳}教会的执事们断言，该教会的主教\Person{马克西米安努斯}因受\OriginalTerm{多纳特派}{Donatists}贿赂而腐化，竟以新的许可允许在他所居住的地方设立一位主教；此事虽先前惯例所允许，但天主教信仰禁止其继续存在。因此，我们认为有必要通过这封信函劝告您的弟兄之爱，当我们的\OriginalTerm{文书}{chartularius}\Person{希拉鲁斯}到达你处时，在同一案件中，应在一次聚集的主教总会议上，怀着对即将来临的审判者的畏惧，进行彻底而明智的调查。若该指控由这些信件的携带者以充分的证据对前述主教成立，则他必须被剥夺他所享有的尊位与职务，使他既能通过承认自己的过失而回到补赎的收益，也使他人不敢贸然尝试此类事情。\par

因为一个人若如所说那样，为收受金钱而将我们的主\Person{耶稣基督}出卖给异端者，他就该被解除处理祂至圣圣体圣血之奥迹的职务。此外，如果除了这项指控，他们之间还有任何争论，正如执事们自己的请愿书中所载，关于某些侵害或私人交易，您的弟兄之爱应连同我们前述的\OriginalTerm{文书}{chartularius}，在提出证据后充分调查，并公正地裁决各方。\par

此外，我们通过这些信件的携带者所提供的情报得知，\OriginalTerm{多纳特派}{Donatists}的异端因我们的罪过而日益蔓延，并且许多人因受贿而获得许可，被多纳特派重新受洗。弟兄，我们应当全心全意思考这事有多么严重。看，那豺狼不再于夜间潜行，而是在光天化日之下撕裂上主的羊群；我们见它屠戮羊群，却毫不焦虑，不用言语的标枪去抵抗它。那么，若我们连自己承担牧养之责的羊群都若无其事地任野兽撕碎，我们还能向上主呈上羊群繁增的什么果实呢？因此，让我们效法地上的牧者，激励自己的心，他们常在冬夜守望，受风雨霜冻之苦，惟恐哪怕一只羊，甚至可能不是有收益的羊，丧亡。若那偷袭者用贪婪的嘴咬伤了它，他们如何忙碌，心跳加速，气喘吁吁，以何等的呼喊跃向前去营救被掳的羊，被紧迫的需要所激动，惟恐因他们的疏忽而损失了什么，羊群的主会向他们追讨！让我们也警醒，免得有什么丧亡；若有什么偶然被夺取，让我们用神圣言语的呼喊将它带回上主的羊群，好使那牧者中的牧者怜悯地垂允，在祂的审判中认可我们守护了祂的羊栈。此外，你需要明智地注意：若同一位主教对这些信件的携带者提出任何正当的请愿，也应彻底调查；若他们自己因自己的过失而理应受到打击，我们宣布，绝不能因他们曾不辞辛劳来投奔我们而宽容他们。\par

八月，第十小纪。\par

\SourceRecord{Translated by James Barmby.；Nicene and Post-Nicene Fathers, Second Series；Vol. 12.；Edited by Philip Schaff and Henry Wace.；Buffalo, NY: Christian Literature Publishing Co.,；1895.；<http://www.newadvent.org/fathers/360202048.htm>.}

% structure: p0001 source=h1 target=section reason=page-title

\section{第二卷，第四十九封书信}

\Emph{致\Person{雅努阿里乌斯}总主教。}\par

\Person{额我略}致\Place{卡利亚里}总主教\Person{雅努阿里乌斯}。\par

如果我们怀有正直的心来审视我们所管理的司铎职务，那么个人爱德的和谐应当如此把我们与我们的子女结合起来，以致正如我们在名义上是父亲，也应在情感上证明我们确实是父亲。然而，我们应当像上面所说的那样，但为何却出现了如此众多的控告反对你的兄弟之德？我们仍然犹豫是否相信这些控告：但为了能够查明真相，我们派遣了我们的宗座公证人\Person{若望}前往你处，他受我们命令的支持，将使各方都服从所选仲裁人的裁决，并亲自执行这些裁决。因此，我们通过此信劝勉你的兄弟之德，要事先好好思考案件的实情；如果你发现你曾不义地夺取或持有任何东西，鉴于你的司铎身份，在审判之前要将其归还。\par

现在，在众多控告中，最杰出的\Person{依西多尔}控告说，他被你的兄弟之德以无效的理由开除教籍并予以弃绝。当我们想从当时在此地的一名你的神职人员那里了解这件事的原因时，他让我们明白，除了这个人冒犯了你之外，没有其他原因。这让我们非常难过；因为，如果是这样，你表明你没有想到天上的事，而是显示出你的交往是在地上的事中，竟然动用弃绝的诅咒来报复私人的冒犯；这是神圣规则所禁止的。因此，今后要十分谨慎小心，不要再为了维护你自己的冤屈而施加任何这样的惩罚。因为，如果你做了任何这类事，要知道将来它会报复在你身上。\par

\SourceRecord{Translated by James Barmby.；Nicene and Post-Nicene Fathers, Second Series；Vol. 12.；Edited by Philip Schaff and Henry Wace.；Buffalo, NY: Christian Literature Publishing Co.,；1895.；<http://www.newadvent.org/fathers/360202049.htm>.}

% structure: p0001 source=h1 target=section reason=page-title

\section{《第二卷，第五十一封书信》}

\Emph{致全体主教}\par

国瑞致全体主教，论 \OriginalTerm{三章}{Three Chapters} 事宜。\par

我极其欣慰地收到了你们的信函；然而，若我有幸能因你们从谬误中回转而欢喜，那将是更大的喜悦。你们的信函开头告知你们正遭受严重的迫害。但迫害若出于非理性，对救恩便毫无益处。因为任何人指望因罪恶而得赏报，实属不敬。你们应当知道，正如真福西彼廉所言，并非受苦使人成为殉道者，而是受苦的原因。既然如此，你们所提及的迫害，你们以此夸耀，是极不恰当的，因为你们并未因此更接近永恒的赏报。因此，愿信德的纯洁引导你们的爱德回归生养你们的母亲教会；愿你们的心意不因任何偏执而与和好的合一分离；愿没有任何劝说能阻止你们再次寻求正道。因为在处理三章的大公会议中，显然没有任何关乎信仰的事被推翻或丝毫改变；但是，正如你们所知，会议只涉及某些个人；其中一人，其著作明显偏离大公信仰的正直，被定罪并非不公。\par

此外，关于你们所写的自那时起意大利在诸省中尤其遭受鞭笞一事，你们不应将此曲解为一种责难，因为经上记载：\Emph{\Quote{主所爱的，祂必管教；凡所接纳的儿子，都必鞭打他}}（\BibleRef{希 12:6}）。那么，若果真如你们所言，意大利自那时起便更加为天主所爱，并得以完全蒙受悦纳，因为它堪当承受主的鞭打。然而，事实并非如你们所试图藉以侮辱她的那样，请听理性的声音。\par

自可敬的教宗维吉吕受命于皇城颁定罪罚判决反对当时的皇后狄奥多拉，或反对 \OriginalTerm{无首派}{Acephali} 之后，罗马城随即遭敌人攻击并攻陷。那么，这是否意味着无首派有理，或者他们被不公定罪，因为定罪后发生了这些事？绝无此理！因为无论是你们中的任何人，还是其他曾领受大公信仰奥迹的人，都不应该说或在任何方面承认这点。既然明此，现在也请你们最终从已作出的决定中退出。因此，为使你们心中充满完全的满足，并除去对三章的一切疑虑，我认为有益于将我先教宗、神圣记忆的教宗佩拉吉关于此主题所写的书寄给你们。若你们愿意反复阅读此书，摒弃固执己见的精神，我相信你们会完全遵循它，并无论如何回归与我们的合一。但若今后在阅读此书后，你们决定坚持目前的决定，无疑将表明你们是屈服于固执而非理性。因此，我再次以怜悯之心劝诫你们的爱德：既然在天主之下，我们的信仰纯洁在三章之事上仍保持完好，你们应除去一切心高气傲，回归期待并召唤其子女的母亲教会；并且越早越好，因为你们知道她每日都在期待你们。\par

\SourceRecord{Translated by James Barmby.；Nicene and Post-Nicene Fathers, Second Series；Vol. 12.；Edited by Philip Schaff and Henry Wace.；Buffalo, NY: Christian Literature Publishing Co.,；1895.；<http://www.newadvent.org/fathers/360202051.htm>.}

% structure: p0001 source=h1 target=section reason=page-title

\section{第二卷 第五十二封信}

\Emph{致纳塔利斯主教}\par

\Person{格列高利}致\Place{萨洛纳}主教\Person{纳塔利斯}。\par

仿佛忘记了先前信函的主旨，我曾决心对于你的福衔只说带有甘甜意味的话；但如今，你在你的书信中又以辩论的方式重提前信，我便再次被迫说一些也许我本不愿说的话。\par

因为，为了为宴席辩护，你的主教尊驾提到了\Person{亚巴郎}的宴席，据圣经的见证，据说他曾在其中款待了三位天使\EditorialNote{创世纪第18章}。鉴于这一榜样，如果我们得知你也款待天使，我们也不会责备你的福衔设宴。你又提到，\Person{依撒格}在饱足之时祝福了他的儿子\BibleRef{创 27:27}。关于旧约中的这两件事——它们既是按历史方式成就的，同时仍含有寓意——但愿我们能如此通读所行之事的事迹，以致能察觉并关注当做的事。因为确实，那一位（\Person{亚巴郎}）只问候了三位天使中的一位，宣明了天主圣三的各位是同一本体；另一位（\Person{依撒格}）在饱足时祝福了他的儿子，因为一个充满神圣筵席的人，其感官被扩展到预言的权能中。而圣经的话语正是神圣的筵席。所以，如果你勤勉地阅读——如果你从外表的事物汲取榜样而深入内在——你便会如同从田野中狩猎而得饱足，充满灵魂的胃腹，从而能够向你面前的儿子，即你所照管的子民，宣告将来的事。但那位讲论天主预言的人，已经对此世昏暗了；因为确实合理而相宜的是：那内在因明智而感官明亮的人，在此世因贪欲所见更少。\par

因此，将这些事拿到自己身上吧；如果你知道自己正如我所说，你丝毫无需怀疑我们的尊重。我也发现你的福衔欣然欢喜，如果你与世界的创造者一同承受\Quote{一个贪食的人}之名。对此我简要评论如下：如果你被不实地如此称呼，你确实与世界的创造者一同承受了此名；但如果是真实的，谁又能怀疑祂是虚假的呢？相似的名不能为你脱罪，如果原因不同。因为就连那被判死刑的盗贼也与祂一同背了十字架；但相似的钉十字架并不能使那被自己罪过捆绑的人脱罪。但现在，我以我所能献上的全部祈祷恳求天主，愿不仅那名号，而且其原因，将你最神圣的主教尊驾与我们的创造者联合起来。\par

此外，你的圣座在信中正确地赞扬了以施行爱德为意图而设的宴席。但你也应当知道，它们真正出于爱德，是在这样的情况下：在宴席上，不在背后诽谤不在场者的生活，无人被嘲笑指责，听到的不是关于世俗事务的无聊闲话，而是神圣读经的话语；身体也不被过分纵容，而只其软弱得到舒缓，以便保持健康以修德。所以，如果你在宴席中如此行事，我承认你是克己的大师。\par

至于你向我引述宗徒保禄的证言，他说：\Emph{那不吃的，不可判断那吃的}\BibleRef{罗 14:3}，我认为这完全不合理，因为一方面我不是那不吃的人，另一方面保禄在此的意思也不是说，基督的肢体，即那些在祂的身体（也就是祂的教会）中藉爱德的纽带彼此相连的人，应当对彼此毫不关心。当然，如果我与你毫无关系，你也与我无关，我本应被迫保持沉默，以免责备一个我无法纠正的人。那么，这条诫命只是针对那些试图判断没有交付给他们照管的人而给出的。但现在，由于天主的安排，我们是一体的，如果我们沉默不语地放过需要纠正的事，我们就有大错了。看，你的主教尊驾因我在宴席上责备了你而感到不快，而我，虽然地位超过你（尽管不在生活上），却准备受所有人指责，并受所有人改正。我只认为那人是我的朋友，通过他的口舌，我能在严厉的法官面前擦去我灵魂的污点。\par

至于你所提到的，最亲爱的弟兄，关于你因受磨难的重压而不能读书，我认为这几乎不能作为你的借口，因为保禄说：\Emph{凡是先前所写的，都是为教训我们而写的，使我们借着经书的忍耐和安慰，能怀有希望}\BibleRef{罗 15:4}。那么，如果圣经是为我们的安慰而准备的，我们越是感到在磨难的重担下疲惫，就越应当多读经书。但如果我们只依赖你在信中所引的那句话，即主说：\Emph{当他们交付你们时，不要思虑怎样或说什么；因为在那时刻，自会赐给你们应说什么；因为说话的不是你们，而是你们父的圣神在你们内说话}\BibleRef{玛 10:19}，那么我说，如果因充满圣神就不需要外在言语，圣经就是白白赐给我们了。但是，至亲的弟兄，在迫害时无疑信赖天主是一回事；当教会平安时我们应当做什么是另一回事。因为我们的责任是，借着同一位圣神，现在通过阅读学习，以便将来若有机会，也能在受苦中显明出来。\par

如今，我极其喜乐，因你在信中声明你正致力于劝勉。因为如此我便知道，你若费心也引导他人归向你的造物主，你便在明智地履行你职位的职责。但你在同一句中说你不如我，这使我在开始喜乐之后立刻忧伤，因我认为你是在嘲笑中给予我那些我实在不承认为应得的赞美。然而，我感谢全能的天主，因藉着你，异端者正被召回圣教会。但你也需要留意，那些在圣教会怀抱中的人也当如此生活，以免他们因邪恶的生活成为她的敌人。因为，他们若不投身于天上的渴望，而投身于世俗的欲望和快乐，则陌生人的儿子正在她的怀抱中被养育。\par

至于你声言你不可能不知道教会职位的等级，关于你我也完全知道；因此我深感困扰，因为，如果你知道事物的秩序，你却在关于我的事上未能知道，这对你更是过错的。因为，在我的前任和我本人为着总执事\Person{Honoratus}的事致信你的真福之后，随后，我俩的裁决被置若罔闻，该\Person{Honoratus}被剥夺了属于他的等级。此事若由四位宗主教中的任何一位所为，如此严重的忤逆绝不能未造成最严重的冒犯而被容许。尽管如此，既然你的兄弟之谊已回归你应有的地位，我不再计较对我或对我的前任的伤害。\par

至于你说我前任所传授并守护的当在我们时代也被遵守，我绝不敢在任何教会中侵犯我先祖关于我同僚司铎的规约，因为我若扰乱我弟兄的权利，便是伤害我自己。但是，当你所派遣的信使到达时，我将知晓你与前述总执事\Person{Honoratus}之间案情的权利；而我亲自审查此事将向你表明，你若拥有正义支持，你将不会从我这里受到任何伤害；正如你从未受过一样。但若正义支持常被提及的\Person{Honoratus}的辩护，我将通过宣告他无罪来表明，在审判中我甚至不认识我所认识的人。\par

关于绝罚的条款（若可如此说），那是必须附加在我们的信函上的（虽然第二次和第三次也插入了条件），你的真福无理由地抱怨，因为宗徒\Person{保禄}说：\BibleQuote{预备好了，要责罚一切不服从}\BibleRef{格后 10:6}。但是，让这些事过去吧：让我们回到我们现今所关心的事。因为，若主\Person{Natalis}按他所应做的行事，我不能不与他和好，因我知道我欠他的情谊多少。\par

\SourceRecord{Translated by James Barmby.；Nicene and Post-Nicene Fathers, Second Series；Vol. 12.；Edited by Philip Schaff and Henry Wace.；Buffalo, NY: Christian Literature Publishing Co.,；1895.；<http://www.newadvent.org/fathers/360202052.htm>.}

% structure: p0001 source=h1 target=section reason=page-title

\section{第二卷·第五十四封书信}

\Emph{以下是圣李奇尼安主教关于\WorkTitle{规则书}写给罗马城教宗圣额我略的书信。}\par

致最蒙福的主教宗额我略，李奇尼安主教。\par

\WorkTitle{规则书}是经由您的圣座颁布，并藉天主的恩宠传达给我们，我们怀着极大的喜悦阅读了它，因为我们发现其中包含的属灵规则。谁不会愉快地阅读这本书呢？在其中，通过不断的默想，他可以找到灵魂的药方；在其中，轻视这个变化无常的世界中瞬息万变的事物，他可以将灵魂的眼睛转向永生的稳固境地。您的这本书是一座所有德行的宫殿。在其中，智德辨明善恶的界限；义德给予各人应得的，使灵魂服从天主，身体服从灵魂。在其中，勇德在逆境和顺境中始终如一，既不被反对所摧毁，也不因成功而骄傲。在其中，节德克制情欲的狂暴，并明智地为快乐设定限度。在其中，您包含了所有与分享永生有关的事物：不仅为牧者制定了生活的规则，也为那些没有治理职务的人提供了生活的准则。因为牧者可以从您的四重划分中学习：他们在进入这一职务时应当如何；进入后应当过怎样的生活；应当如何教导以及教导什么；应当做什么以避免在如此崇高的司铎职位上骄傲自大。您的这番卓越教导得到了教会神圣的古代教父、圣师和护卫者的证明：希拉利、盎博罗削、奥古斯丁、额我略·纳齐盎——他们全都为您作证，如同先知们为宗徒们作证一样。圣希拉利在解释那位外邦人的导师——宗徒——的话时说：\Quote{因为他表示，属于纪律和道德的事物有助于司铎职的功德，如果教导和守卫信仰的科学所必需的事物也不缺少的话；因为仅仅无罪地行动，或仅仅博学地宣讲，都不能立刻构成一位良善而有益的司铎，因为，一个人虽然无罪，但若不博学，他只对自己有益；而一个人虽然博学，但若不无罪，他就没有教师的权威。}圣盎博罗削在他所写的关于\WorkTitle{论职务}的书中为您的这本书作了见证。圣奥古斯丁作证说：\Quote{在行动中，不应爱慕此世的尊严，也不应爱慕权力，因为太阳之下一切都是虚妄的。}但是，藉由这种尊严或权力所完成的工作本身，如果它被正确而有益地完成，那么它就对属下那合乎天主的益处有所贡献。因此宗徒说：\BibleQuote{谁若想要主教的职务，就是渴望一件善工。}他愿意解释什么是\OriginalTerm{监督}{episcopus}；这是一个表示工作而非尊严的称号。因为它是一个希腊词，源出于此——被置于他人之上的人，俯瞰他所管辖的人，意即照顾他们；因为\OriginalTerm{监督}{episcopacy}就是俯瞰。因此，如果我们愿意，可以用拉丁语说，履行主教的职务就是俯瞰；这样，一个喜欢凌驾于他人之上而不使他们受益的人可以明白他不是一位主教。因为如此，没有人被禁止渴望认识真理，为此目的，闲暇是值得称道的；但是，对于一种优越的地位——没有它人民就无法被治理——尽管它可以被体面地担任和管理，但觊觎它是不恰当的。因此，爱德寻求神圣的闲暇，以便有时间领悟和捍卫真理。但如果[治理的重担]被强加，那么它应当因爱德的责任而承担。但即使如此，也不应完全放弃对真理的喜爱，免得先前的甘甜被剥夺，而当前的责任成为重负（\WorkTitle{论天主圣三}卷八第一章）。\par

圣额我略·纳齐盎也为此作证，您的文风追随他，您也效法他的榜样，想要隐藏自己以避免司铎职的重担；这重担的性质如何，在您的全书中清楚地陈述了：然而您承担了您所害怕的。因为您的重担是向上托举，而非向下压制；不是将您沉入深渊，而是将您提升至星辰；藉着天主的恩宠、服从的功德和善工的功效，那因人性软弱而看似沉重的事物变得甘甜。因为您所讲说的与宗徒们及宗徒时代的人们一致。因为您本身是美好的，您说了美好的事物，并在其中显示了您的美好。我不愿您将自己比作一个画美好事物的丑陋画家，因为属灵的教导出自属灵的灵魂。大多数人评价人的画家高于无生命的画作。但不要将此视为奉承或谄媚，而是视为真理：因为撒谎既不适合我，也不适合您赞扬虚假的事物。那么，我虽坦诚直率，却已看出您及您的一切都是美好的，并看出我自己与您相比是足够丑陋的。因此，我恳求您，藉着您内丰盈的天主恩宠，不要拒绝我的祈祷，而是心甘情愿地教导我，我承认自己对此一无所知。因为我们是被迫必须做您所教导的事。\par

因为，当找不到一个精通司铎职的人时，除了像我这样不精通的人被祝圣外，还能做什么呢？您命令不可祝圣不精通的人。但让您的明智考虑一下：难道认识\Person{耶稣基督}，并认识祂被钉在十字架上，对一个人来说不足以算作精通吗？因为，如果这还不够，那么按照这本书，就没有人能被称为精通了。这样就没有人能做司铎，如果除非精通的人，没有人能做司铎。因为我们公然抵抗重婚者，以免圣事因此被玷污。如果一妻之夫在结婚前曾触摸过女人，那又怎样？如果他没有娶妻，却又并非没有触摸过女人，那又怎样？请您用您的笔安慰我们，使我们不至于因自己的罪或他人的罪而受罚。因为我们极其害怕被迫做我们不该做的事。看，必须服从您的训令，以便这样的人——权威所认可的人——被立为司铎；而所寻求的人却找不到。这样，由听道而来的信德将会消失；圣洗圣事将会消失，如果没有人施洗的话；那些至圣奥迹，通过司铎和辅礼人员而实现的，将会消失。无论哪种情况，危险都存在：要么必须祝圣一个不该祝圣的人，要么就没有人能举行或施行神圣奥迹。\par

几年前，\Place{西斯帕利斯}主教\Person{莱安德}从皇城返回时，途经我们这里，告诉我们说，您福乐曾发表过一些关于\WorkTitle{约伯传}的讲道集。但他匆匆经过，没有如我们所请求地把它们给我们看。但您后来写信给他谈及三重浸洗，并在信中说了，据我所知，您对那部作品不满意，并决定经过更成熟的考虑，将那些讲道集改为论著的形式。\par

我们确实有\Person{圣希拉里}，\Place{普瓦捷}主教的六卷书，他从\Person{奥利振}的希腊文翻译成拉丁文；但他没有按顺序解释整个\WorkTitle{圣约伯传}。我颇感惊讶，一位如此博学且如此神圣的人，竟会翻译\Person{奥利振}关于星辰的愚蠢故事。至圣神父，我无论如何也不能相信天上的光体是理性灵体，因为圣经并没有宣告它们是与天使一同造成的，或是与人类一同造成的。那么，愿您的福乐屈尊将这部作品，以及您在\WorkTitle{规则书}中提到已经写成的其他伦理著作，传送给我的卑微。因为我们是属于您的，并且乐于阅读您的作品。对我来说，正如您的\Person{额我略}所说，即使到极老之年仍要学习，是一件令人向往和光荣的事。愿\Person{天主圣三}保全您的冠冕不受伤害，以教导祂的教会，正如我们所希望的，至福的神父。\par

\SourceRecord{Translated by James Barmby.；Nicene and Post-Nicene Fathers, Second Series；Vol. 12.；Edited by Philip Schaff and Henry Wace.；Buffalo, NY: Christian Literature Publishing Co.,；1895.；<http://www.newadvent.org/fathers/360202054.htm>.}

% structure: p0001 source=h1 target=section reason=page-title

\section{第三卷，第一封信}

\Emph{致副执事伯多禄。}\par

\Person{额我略}致\Place{坎帕尼亚}的副执事\Person{伯多禄}。\par

送交我们的报告已经清楚表明，在\Place{卢库鲁斯堡垒}内，对我们的弟兄和同为主教的\Person{保禄}犯下了何等罪行。而且，由于\Place{坎帕尼亚}的可敬法官\Person{斯高拉提库斯}目前恰与我们同在，我们已特别委托他严格纠正如此大的悖逆之疯狂。但是，由于上述报告的送交者请求我们派遣某人代表我们，因此我们派遣副执事\Person{埃丕法尼乌斯}，他将与上述法官一同调查并查明骚乱是由谁煽动或调查的，并处以适当的惩罚。那么，愿你的经验全力以赴协助此案，以便查明真相，并对有罪者进行惩罚。因此，既然光荣的\Person{克莱门提娜}的奴隶据说与此同一罪行有关，并使用了煽动骚乱的言论，你要立即严格地惩罚他们，你的严厉不可因顾及她本人而松懈，因为他们既是高贵的妇女的仆人，因纯粹骄傲而犯罪，更应受到打击。但你也应当彻底调查，上述妇女是否知悉如此残暴的罪行，以及是否在她知晓下实施的，以便通过我们的巡视，所有人都能学习到：不仅对司铎动手，甚至在言语上冒犯司铎是多么危险。因为，如果此事处理松懈或有遗漏，要知道你将特别承担责备和风险；你在我们面前找不到任何借口。因为此事若得到最严格的调查和纠正，就会使你得到我们的认可；反之，若被轻描淡写，要知道我们的愤怒会加剧反对你。\par

此外，至于其他方面，如果有城内的奴隶逃往圣\Person{塞味利诺}隐修院或同一堡垒内的任何其他圣堂，一旦你得知此事，绝不可允许他们留在那里，而应将他们带到城内的圣堂；如果他们对自己的主人有正当的抱怨，必须在为他们做出适当安排后离开圣堂。但若他们犯了小罪，应立即使他们回到主人那里，主人应宣誓宽恕他们。\par

\SourceRecord{Translated by James Barmby.；Nicene and Post-Nicene Fathers, Second Series；Vol. 12.；Edited by Philip Schaff and Henry Wace.；Buffalo, NY: Christian Literature Publishing Co.,；1895.；<http://www.newadvent.org/fathers/360203001.htm>.}

% structure: p0001 source=h1 target=section reason=page-title

\section{第三卷，第二封信}

\Emph{致\Person{保禄}主教}。\par

\Person{额我略}致\Person{保禄}等。\par

虽然我们听说您所受的伤害，感到极大的痛心，但我们得知此事对您有利，仍感到安慰，因为据报给我们的事件经过揭示事实，您是为正直和公义而受害。所以，为使此事为您的弟兄带来更大的光荣，这次事件不应动摇您的坚定，也不应使您偏离真理之道。因为司铎若在受伤害后仍坚持真理之道，必获更大赏报。但是，为避免如此严重的邪恶的疯狂不受惩罚，且有害的违命行为变本加厉，我们已委托目前在座的尊贵的\Place{坎帕尼亚}法官\Person{斯柯拉斯提库斯}，应予以应得的压制，惩罚所行之事。但是，因您的人请求我们委派一位代表我们的人，故此，我们已为此派遣副执事\Person{厄丕法尼}前往\Place{那不勒斯}，他可与上述法官一同调查并查明真相，以便他能通过他的坚持，使那些被证实煽动或犯下如此重大罪行的人受到应有的惩处。\par

\SourceRecord{Translated by James Barmby.；Nicene and Post-Nicene Fathers, Second Series；Vol. 12.；Edited by Philip Schaff and Henry Wace.；Buffalo, NY: Christian Literature Publishing Co.,；1895.；<http://www.newadvent.org/fathers/360203002.htm>.}

% structure: p0001 source=h1 target=section reason=page-title

\section{第三卷 第三封信}

\Emph{致院长若望}\par

额我略致若望等\par

您的爱德曾请求我，波尼法爵弟兄可在您的隐修院被祝圣为院长（\OriginalTerm{院长}{praepositus}）；关于此请求，我甚为诧异为何之前未曾办妥。因为自我将此人交付于您之时起，您本应早已祝圣他。\par

关于圣若望的束腰外衣，您急切告知于我，我完全感到欣慰。但请您的爱德设法将此外衣寄送给我，或更好将此拥有外衣的主教本人，连同他的圣职人员和那件外衣一并送来，好使我们能享受此祝福，并能从这位主教及其圣职人员那里获益。我一向渴望了结与佛洛里亚努斯待决的诉讼，且已向他预付了多达八十\Emph{索利多}，我相信他打算将此作为隐修院债务的补偿；我极其希望此案得以了结，因为据说档案管理员斯德望促使上述佛洛里亚努斯将此案移交公审，而我们不愿卷入公开的诉讼。因此，我们必须作出某些让步，以便能就此案达成和解。此事完成后，我们将通知您的爱德。\par

但是你当全心关注兄弟们的灵魂。目前隐修院的声誉因你的疏忽而受损，这已足够。不要经常外出。为这些事务指定一位代理，而你当为自己留出阅读和祈祷的时间。\par

要热心款待；尽你所能，施舍给穷人；但须保留应当偿还给佛洛里亚努斯之物。此外，在你隐修院的弟兄中，我看不到他们对阅读的热忱。因此你须思考，这是何等大罪：天主从他人的奉献中赐予你生活费，而你却忽略学习天主的诫命。\par

此外，关于那六份，除非我们看到原始契约或副本，否则我们无能为力。但我已向天主之仆佛洛伦蒂努斯发出指令，若真相向他显明，他当将那六份归还于你；归还之后，我们将以租约方式让与其余六份，或改变其收益。\par

\SourceRecord{Translated by James Barmby.；Nicene and Post-Nicene Fathers, Second Series；Vol. 12.；Edited by Philip Schaff and Henry Wace.；Buffalo, NY: Christian Literature Publishing Co.,；1895.；<http://www.newadvent.org/fathers/360203003.htm>.}

% structure: p0001 source=h1 target=section reason=page-title

\section{第三卷，书信五}

\Emph{致副执事\Person{伯多禄}}\par

\Person{额我略}致\Place{坎帕尼亚}副执事\Person{伯多禄}：\par

我们无意干扰平信徒在审判中的特权，但当他们审判不公时，我们希望你以温和的权威予以抵制。因为约束暴力的平信徒并非违反法律，而是支持法律。既然据说\Place{奥尔蒂切卢姆}的\Person{费利克斯}的女婿\Person{德乌斯德迪特}对持此信之人施以暴力不公，并且至今仍非法扣押其财产，以至于她的寡妇身份的沮丧并未激发他的同情，反而加深了他的恶意，因此我们命令你的经验，针对上述之人，以及在上述妇女声称遭受不公的其他案件中，你要提供保护的援助，不允许任何人压迫她，以免你被发现忽视了在不损害公正的情况下所命令你的事，或者寡妇和其他穷人，在本地找不到帮助，因长途跋涉至此而承担费用。\par

\SourceRecord{Translated by James Barmby.；Nicene and Post-Nicene Fathers, Second Series；Vol. 12.；Edited by Philip Schaff and Henry Wace.；Buffalo, NY: Christian Literature Publishing Co.,；1895.；<http://www.newadvent.org/fathers/360203005.htm>.}

% structure: p0001 source=h1 target=section reason=page-title

\section{第三卷，第六封信}

\Emph{致若望主教}\par

\Person{额我略}致\Person{若望}，\Place{普里马·犹斯提尼亚那}主教。\par

在\Person{阿德里安}——\Place{特拜}城主教——如同面对仇敌一般从他的同司铎那里遭受了长久的苦难之后，他逃到了罗马城避难。虽然他最初控告的是\Person{若望}（\Place{拉里萨}主教），声称他在财务诉讼中不顾法律作出了判决，但此后他更严厉地控告的，却是你的兄弟情谊本人，指责你不公正地剥夺了他的司铎职位。然而，我们并不轻信未经调查的诉状，仔细审阅了案件的卷宗——无论是在我们的兄弟和同主教\Person{若望}面前，还是在你面前进行的。确实，关于上述\Person{若望}（\Place{拉里萨}主教）的判决（该判决因上诉而暂停），最虔诚的皇帝们已在给\Place{科林斯}主教的命令中充分裁定，而我们在基督的助佑下，也通过本信函携带者带给上述\Person{拉里萨的若望}的信件中作出了裁定。但是，在反复审查了皇帝命令委托你调查的那些相互冲突的判决，并检视了在\Person{若望}主教面前就受指控之人所进行的程序记录之后，我们发现你几乎没有调查任何与指定给你裁决的问题相关的事项，反而通过某种阴谋，找到了针对执事\Person{德默特琉}的证人，这些人为了定这位主教（即\Person{阿德里安}）的罪，竟声称他们曾听到\Person{德默特琉}就上述主教作证——这种事甚至听都不该听。而当\Person{德默特琉}本人否认做过此事时，看来你违背了司铎的惯例和纪律规范，将他作为一名被革除职位的执事交给了省总督。在遭受多次鞭打之后，他或许在酷刑的压力下说了些对主教不利的假话，但我们发现，直到事情结束，他对于被讯问的事情根本没有承认任何东西。我们也没有在卷宗中找到任何其他不利于\Person{阿德里安}的内容——无论是在证人证词中还是在\Person{阿德里安}的声明中。唯一的是，你的兄弟情谊——不知出于何种动机——藐视人法和天法，对他宣布了一项仓促的判决；即使该判决没有因上诉而中止，因为它违反了法律和教规，本身也不能合法成立。此外，在显而易见的上诉已提交给你之后，我们奇怪的是，你为何没有按照你的教会的代表向我们执事\Person{奥诺拉图斯}所作的承诺，派遣你的人来向我们报告你的判决。这一疏漏证明你要么是傲慢，要么是良心不安。因此，如果提交给我们的事实确凿，鉴于我们认为你滥用你在我之下的代理管辖权作出了不公正的裁决，我们将依靠基督的帮助，根据我们商议的结果，进一步就这些事项作出决定。\par

但就目前而言，凭借真福\Person{伯多禄}——宗徒之长——的权威，我们下令：首先撤销你的判决，使之完全无效，并在三十天内剥夺你的圣体共融，以便你以最大的忏悔和泪水为如此重大的过犯祈求我们的天主的宽恕。但是，如果我们得知你在执行我们这一判决时懈怠，要知道你的兄弟情谊不仅要因不公，还要因傲慢而受到更严厉的惩罚。至于我们上述兄弟和同主教\Person{阿德里安}，被你那既不符合教规也不符合法律的判决所定罪，我们命令：在基督与他同在的情况下，恢复他的职位和品级；这样，他既不会因你偏离正义之路的判决而受伤害，你的爱德也不会得不到纠正；如此，我们才能平息将来审判者的义怒。\par

\SourceRecord{Translated by James Barmby.；Nicene and Post-Nicene Fathers, Second Series；Vol. 12.；Edited by Philip Schaff and Henry Wace.；Buffalo, NY: Christian Literature Publishing Co.,；1895.；<http://www.newadvent.org/fathers/360203006.htm>.}

% structure: p0001 source=h1 target=section reason=page-title

\section{第三卷，第七封信}

\Emph{致若望主教。}\par

\Person{额我略}致\Place{拉里萨}主教\Person{若望}。\par

我们的弟兄，\Place{底比斯}城主教\Person{阿德里安}来到\Place{罗马}，痛诉他被你的弟兄及\Place{普里马·尤斯蒂尼亚纳}主教\Person{若望}既不合法律也不合教会法规地依据某些指控定罪。当我们久久未见对方任何代表前来此地反驳他的质疑时，我们便将你们面前进行的程序交付审阅，以便必要地查明真相。从这些程序中我们得知，被撤职的执事\Person{若望}与\Person{科斯马斯}——一人因身体的软弱，另一人因欺诈教会财产——曾就金钱事宜及刑事指控向我最虔诚的皇帝们提交了一份针对他的呈文。\par

他们下达给你们的谕令中，要求你们（即严格遵循法律与教会法规）审理此事，以便对金钱问题作出合法可靠的判决，但对于刑事指控，则须在彻底调查后向皇帝陛下禀报。如果你的弟兄以正确心态接受这些如此公正的谕令，就决不会接受那些因过犯已被撤职、且已怀有敌意之人对其主教提出的全面控告；尤其因他们向我最虔诚的主上们提交的呈文中暴露了不诚实，因他们宣称该呈文是在全体圣职人员同意下作出的。\par

此后，简略概括你们面前的某些程序：第一项指控涉及\Place{底比斯}执事\Person{斯德望}，据称主教\Person{阿德里安}虽已知晓其极可耻的生活，却未褫夺其职位的尊严。关于此项，除\Person{斯德望}一人——一个生活可耻、且因自身供认当受定罪的人——据称如此说外，并无证人证明\Person{阿德里安}主教对此知情。对他提出的第二项指控似乎是，他曾下令阻止婴孩领受圣洗圣事，致使他们带着未洗净的罪恶污秽死去。但出庭指证他的证人中，无人声明他们知道任何此类事情曾引起\Person{阿德里安}主教的注意，而只是说从婴孩的母亲们那里得知；据说这些母亲的丈夫因罪行已被逐出教会。但即便如此，他们也未声明那些婴孩在未受洗时死亡时刻来临，正如控方恶意呈文中所含的内容，因为显然他们在\Place{德米特里亚斯}城已受洗。刑事指控至此为止。\par

至于金钱事宜，你们如何裁决，可由受君王委派的人士依据最安详君主的至虔谕令所作的调查证明。因为，当屡次提到的\Person{阿德里安}对你们的判决提出上诉时，据我们从提交给\Place{普里马·尤斯蒂尼亚纳}主教\Person{若望}的四位证人证词中得知，他被投入最严厉的监禁，并被你的弟兄强迫出具一份文件，在其中他承认了对他的指控。诚然，在他出具的文件中，他承认了你们关于金钱问题的判决。但对于刑事指控，他以含糊可疑的方式触及，仿佛以某些迷雾阻挠你们的意图，且日后能更好地在混乱言词的晦暗中摆脱他的承认。当由他属下提交的上诉书及你们审理的其他程序已呈报给最虔诚的君主们，且我们的宗座执事\Person{奥诺拉图斯}与光荣的\OriginalTerm{文书官}{antigraphus}\Person{塞巴斯蒂安}受委派（如前所述）后，他被最安详的主上们免除所有进一步的命令。但我不晓得通过何等刻意设计的阴谋，又再次引出了一道帝国谕令，要求\Place{普里马·尤斯蒂尼亚纳}主教\Person{若望}就所有上述指控进行深入调查并作出判决。在此审讯中，\Person{阿德里安}主教的所有圣职人员以及执事\Person{德米特里乌斯}——后者在酷刑中——均宣称，所有这些对\Person{阿德里安}主教的诽谤均是经你的弟兄的诡计编造起来的。在你们面前对\Person{阿德里安}主教所作的任何刑事指控均未被证实。相反，违反教会法规与法律，另一次针对其执事\Person{德米特里乌斯}及其他人的残酷狡猾的调查又出现了，其中并未发现任何可使屡次提到的\Person{阿德里安}被合法定罪的理由，反而找到了释放他的依据。关于\Place{普里马·尤斯蒂尼亚纳}城主教\Person{若望}及其极不公正和可憎的判决，我们将采取进一步措施。至于\Person{阿德里安}主教，我们发现他既承受了与您司铎身份不相称的敌意，且因你的弟兄的判决而无正当理由地在金钱问题上被定罪。\par

既然他被上述\Person{若望}——\Place{普里马·尤斯蒂尼亚纳}的主教——违反法律和\OriginalTerm{教规}{canons}罢免，不能让他继续被剥夺其地位和荣誉，我们特此下令恢复他在其教会中的职位，并召回他至其应有尊严的秩序。尽管按道理您本应被剥夺主的身体的\OriginalTerm{共融}{communion}，因为您无视我可敬的前任的告诫——他藉此豁免了他及其教会不受您的权威管辖——您却再次僭越地对他们保留某种管辖权；然而我们，以更仁厚的判决，并仍允许您领受共融圣事，特此规定您的弟兄不得再行使过去您对他及其教会所拥有的任何管辖权。但根据我们前任的书面指示，若可能发生任何针对上述我们的司铎弟兄\Person{阿德里安}的案件，无论是关涉信仰、刑事还是金钱，若问题轻微，则由在皇城中的或将成为我们代表的人审理；若问题艰巨，则提交至宗座这里，以便在我们面前听审并裁决。但若您在任何时候，以任何借口或用任何隐秘手段试图违反我们的这些法令，须知我们判决您被剥夺圣体共融，除非在生命终结时，并得到罗马教宗的许可，否则不得领受。因为我们依据圣教父们的教导定下这条规则：凡不懂得服从神圣教规的人，也不配在神圣的祭台上举行或领受共融。此外，您的弟兄须立即将您据称至今仍保留的神圣财产或其他任何动产或不动产归还给他；我们已将一份明细清单附于本函，该清单已交付给我们。关于此事，若你们之间有任何疑问，我们希望由我们在皇城中的代表予以审理。\par

\SourceRecord{Translated by James Barmby.；Nicene and Post-Nicene Fathers, Second Series；Vol. 12.；Edited by Philip Schaff and Henry Wace.；Buffalo, NY: Christian Literature Publishing Co.,；1895.；<http://www.newadvent.org/fathers/360203007.htm>.}

% structure: p0001 source=h1 target=section reason=page-title

\section{第三卷 第八封信}

\SourceRecord{Translated by James Barmby.；Nicene and Post-Nicene Fathers, Second Series；Vol. 12.；Edited by Philip Schaff and Henry Wace.；Buffalo, NY: Christian Literature Publishing Co.,；1895.；<http://www.newadvent.org/fathers/360203008.htm>.}

% structure: p0001 source=h1 target=section reason=page-title

\section{第三卷 第九书}

\Emph{致副执事安东宁}\par

\Person{额我略}致\Person{安东宁}等。\par

我们听说，\Place{埃皮达鲁斯}城的主教\Person{弗洛伦提乌斯}，其财产先被没收，且未经司铎会议，因某些未证实的罪行而被判罪。而且，既然他不应受教会惩罚，因为并未宣判任何教会判决以定他的罪，我们命令你敦促我们的兄弟和同僚主教\Person{纳塔利斯}，使他将上述之人从据说被驱逐的流放地召回。并且，一旦召集了主教会议，如果针对他的指控能按教会法得到证实，我们愿意我们上述兄弟和同僚主教\Person{纳塔利斯}的判决对他生效。但是，如果他被全体审判宣告无罪，你不可允许他受到任何人的偏见，而且必须认真且严格地坚持将他的前述财产归还给他。因此，你越是感到此类交涉的负担沉重，便越要以更成熟和更警醒的执行，努力完成它们。\par

\SourceRecord{Translated by James Barmby.；Nicene and Post-Nicene Fathers, Second Series；Vol. 12.；Edited by Philip Schaff and Henry Wace.；Buffalo, NY: Christian Literature Publishing Co.,；1895.；<http://www.newadvent.org/fathers/360203009.htm>.}

% structure: p0001 source=h1 target=section reason=page-title

\section{卷三，第十封信}

\Emph{致副执事萨维努斯}。\par

\Person{额我略}致\Person{萨维努斯}等。\par

恶人已出去扰乱了你们的心，不明白他们所说的话，也不知道他们所坚持的是什么，假装在虔诚纪念的\Person{犹斯提尼安}时代，从\Place{加采东}神圣大公会议的信德中有所减损，而我们以全备的信德和全部的虔诚敬拜那次大公会议。同样，我们接受圣而公教会所有四次大公会议，如同我们接受四卷\WorkTitle{神圣福音}一样。但是，关于那些在大公会议结束后有所行动的人，在虔诚纪念的\Person{犹斯提尼安}时代曾有一些讨论；然而，信德在任何方面都没有被违反，而且关于这些人的任何事，也没有超出同一\Place{加采东}神圣大公会议所决定的。此外，我们诅咒任何胆敢从那次大公会议所颁布的信德定义中减损，或如同修订一般改变其含义的人；但正如它在那里被颁布，我们在各方面都持守它。因此，最亲爱的儿子，你应当回归到圣教会的合一中，以便你能在平安中度过你的日子；免得那恶神，虽然不能通过你的其他行为胜过你，却在你离世之日，因这事找到阻止你进入天上之国的途径。\par

\SourceRecord{Translated by James Barmby.；Nicene and Post-Nicene Fathers, Second Series；Vol. 12.；Edited by Philip Schaff and Henry Wace.；Buffalo, NY: Christian Literature Publishing Co.,；1895.；<http://www.newadvent.org/fathers/360203010.htm>.}

% structure: p0001 source=h1 target=section reason=page-title

\section{第三卷，第十二封信}

\Emph{致\Person{Maximianus}，主教。}\par

\Person{Gregory}致\Place{Syracuse}主教\Person{Maximianus}\par

我前些时候曾写信给你的兄弟，请你将那些对\Place{Agrigentum}城主教\Person{Gregory}提出控告的人送到\Place{罗马城}。现在我们借这封信敦促你立即办理此事。因此请尽快将当事人本人以及其余文件，即诉讼记录和已提交的请愿书送来。我们不允许任何拖延或借口；这样，当他们如我们所说的迅速被送到\Place{罗马城}后，我们就能在天主的助佑下知道如何最妥善地处理他。\par

\SourceRecord{Translated by James Barmby.；Nicene and Post-Nicene Fathers, Second Series；Vol. 12.；Edited by Philip Schaff and Henry Wace.；Buffalo, NY: Christian Literature Publishing Co.,；1895.；<http://www.newadvent.org/fathers/360203012.htm>.}

% structure: p0001 source=h1 target=section reason=page-title

\section{第三卷，第十五封信}

\Emph{致法官\Person{斯科拉斯蒂库}}\par

\Person{格雷戈里}致坎帕尼亚法官\Person{斯科拉斯蒂库}。\par

当我们因关怀\Place{那不勒斯}城而深感忧虑，该城失去了司铎的慰藉时，这些信件的携带者带着选举我们的副执事\Person{弗洛伦提乌}的法令到来，在如此沉重的思想负担下给了我们一些缓解。但是，当得知前述副执事逃离该城，流着泪恳求不要祝圣他时，我们的悲伤增加了，仿佛是因某种更沉重的天命。因此，谨此问候，我们敦促阁下召集城中的首领或民众，以便考虑选举另一位，他因\Person{基督}的安慰而堪当升任司铎职。然后，法令庄严通过并传至本城后，赖\Person{基督}的助佑，你们之间即可进行祝圣。但若你们找不到一位可共同同意的合适人选，至少选三位正直智慧之人，作为团体的代表派往本城，全体民众可同意他们的判断。也许，当他们来到此地，他们能寻得一位可无可指摘地祝圣为你们主教的人，为使你们失去牧者的城既不缺乏内在的行为监督者，又能在获得司铎照顾时，不给外来的敌意陷阱以可乘之机。\par

\SourceRecord{Translated by James Barmby.；Nicene and Post-Nicene Fathers, Second Series；Vol. 12.；Edited by Philip Schaff and Henry Wace.；Buffalo, NY: Christian Literature Publishing Co.,；1895.；<http://www.newadvent.org/fathers/360203015.htm>.}

% structure: p0001 source=h1 target=section reason=page-title

\section{第三卷·第二十二封信}

\Emph{致副执事\Person{安多尼诺}}\par

\Person{额我略}致副执事\Person{安多尼诺}，\Place{达尔马提亚}教产总管。\par

本地区普遍报道，我们的兄弟和同僚主教，萨洛纳教会的纳塔利斯已经去世。如果属实，你务必以最快速度和最大谨慎告诫该城的圣职人员和民众，使他们一致同意选举一位司铎接受祝圣；并且，当被选者的提名完成后，你要负责将其转达给我们，以便他经我们同意后祝圣，正如从古以来所行的。而你必须特别留意，在此次选举中，无论如何不可有任何形式的贿赂介入，也不可让任何人的庇护得势。因为若有人通过某些人的庇护当选，他在祝圣后出于对他们的尊重，就不得不迎合他们的意愿，从而导致该教会的财产减少，并且无法维持教会秩序。因此，在你的监督下，他们必须选举这样一个人：他不会不当屈从于任何人的意志，而是在其生活和言行的美德上堪当如此崇高品位的人。但是，关于同一教会的财产或装饰品，要令人在你面前忠实地编写一份清单。同时，为了不让任何财产丢失，要告诫执事雷斯佩克图斯和首席公证员（\OriginalTerm{首席公证员}{primicerium notariarum}）斯特凡努斯，让他们单独管理这些财产，并警告他们，若因疏忽致使财产减少，他们必须用自己的财物弥补。\par

此外，严令我们的兄弟和同僚主教\Person{马尔库斯}完全不得干预此事。因为，若我们得知他违背我们的意愿做了或试图做任何事，让他知道他将会招致不小的罪过和危险。但也须注意警告他，他必须仔细记录并完成他负责的我们教产的账目；为此，他应放下一切借口尽快从西西里地区来到我们这里。因此，他绝不可妄自干涉萨洛纳教会的事务，以免他对该教会承担更多责任，并可能被认定有罪。因为据说他拥有属于上述教会的许多物品；而且传言称他几乎是变卖其财产及其他非法行为的主谋。并且，若查明事实确实如传闻所说，他可以肯定，这绝不会不受到惩处。\par

让任何必要的开支由上述主教去世时在职的管家支付，这样他就可以按照他所知道的情况向未来的主教说明他的账目。我们命令你所做的一切事情，你确实必须与我们的儿子，卓越而最有口才的\Person{马尔塞鲁斯}商议后执行，以便你能够仔细而有效地完成这份指示文件中包含的所有内容，并且不会因疏忽而受到责备。\par

\SourceRecord{Translated by James Barmby.；Nicene and Post-Nicene Fathers, Second Series；Vol. 12.；Edited by Philip Schaff and Henry Wace.；Buffalo, NY: Christian Literature Publishing Co.,；1895.；<http://www.newadvent.org/fathers/360203022.htm>.}

% structure: p0001 source=h1 target=section reason=page-title

\section{第三卷，第二十九封书信}

\Emph{致米兰的长老与圣职人员}\par

额我略致米兰教会的长老、执事及圣职人员。\par

我们收到了你们爱德的书信，虽无署名，但由携带人——长老玛格努斯与圣职人员希玻律亲自担保。阅读后，我们发现你们一致赞同我们的儿子、你们教会的执事君士坦提乌斯，他与我相识已久。当我代表宗座在皇城时，他长期紧跟我；但我从未在他身上发现任何可指责之处。然而，长久以来我坚决决定不为任何人偏袒，以免其承担牧养之重担，我只能以祈祷陪伴你们的选举，愿全能的天主——祂永远预知我们未来的行为——赐予你们一位牧者，使在他的言语和举止中，你们能找到神圣劝勉的牧场；一位在其性格中谦卑与正直、严厉与慈爱同辉的人；一位不仅能在言语上，更能在生活中向你们展示生命之道的人；这样，你们的爱德能从他的榜样学会渴望永恒之乡而叹息。因此，最亲爱的儿子们，我们因着对我们职务监督职责的意识，敦促你们在为自己选择主教这件事上，不可有人只顾私利而不顾共同利益，以免若有人追求个人利益，被轻率的判断所欺骗：因为被贪欲束缚的心灵无法以自由的判断审查某人的优先资格。因此，考虑什么是为众人有益，你们当永远完全服从那被天主恩宠置于你们之上的人。因为一旦被置于你们之上，他不可再被你们审判；尽管现在他应被更彻底地审判，因为此后他可能不再被审判。但是，当天主许可一位牧者为你们祝圣后，你们当全心委身于他，并在他内服事那将祂置于你们之上的全能之主。\par

然而，既然上主的审判常按其功过为人民提供牧者，你们当寻求属灵之事，热爱天上之事，轻视暂时和易逝之事；并坚信，如果你们在行为中取悦天主，你们必将有一位取悦天主的牧者。看，这世上的一切，我们曾从圣经中听到注定要灭亡的，如今已见其毁灭。城市倾覆，营寨被拔，教堂被毁；我们土地上再无耕种之人。在我们这些所剩无几的人中，人的刀剑不断肆虐，同时伴随着从上而来的打击。因此，我们在眼前看见了许久前就听说要临到世界上的灾祸，大地本身已成为我们书卷的书页。既然万物都在消逝，我们应当思虑我们所爱的一切皆属虚无。因此，以忧心注视永恒审判者即将来临的日子，并以悔改预先面对其恐怖。以眼泪洗去你们一切过犯的污点。以暂时的哀哭平息那永远悬在你们头上的义怒。因为当我们的慈爱造物主在审判时来临，祂将以更大的恩宠安慰我们，正如祂现在看到我们为自己的过犯惩罚自己。\par

我们现在派遣我们的副执事若望，这些礼物的携带者，前往你们那里——愿天主恩宠相助——使他能协助你们当选的主教按照其前任的方式被祝圣。因为我们要求他人尊重我们的权利，同样我们也为所有人保存他们各自的权利。\par

\SourceRecord{Translated by James Barmby.；Nicene and Post-Nicene Fathers, Second Series；Vol. 12.；Edited by Philip Schaff and Henry Wace.；Buffalo, NY: Christian Literature Publishing Co.,；1895.；<http://www.newadvent.org/fathers/360203029.htm>.}

% structure: p0001 source=h1 target=section reason=page-title

\section{第三卷，第三十函}

\Emph{致副执事\Person{若望}。}\par

\Person{额我略}致\Person{若望}等。\par

既然宗座是因天主的安排而统管众教会，那么在我们多重的关怀中，当决定主教祝圣之事而等待我们的裁决时，更需我们特别留意。现今，在\Place{米兰}教会的主教\Person{劳伦提乌斯}去世后，该教会的圣职人员向我们报告说，他们一致同意选举我们的儿子、他们的执事\Person{君士坦提乌斯}。然而，他们的报告未经签署，故为避免任何疏忽，有必要让你持此命令的权柄前往\Place{热那亚}。并且，由于目前有许多\Place{米兰}人因野蛮人的凶暴而避难至彼处，你应召集他们，共同询问他们的意愿。若没有意见分歧破坏选举的共识——即若你确知众人仍一致希望并同意由上述我们的儿子\Person{君士坦提乌斯}继任——那么，你应依古代惯例，由他自己的主教们祝圣他，同时附上我们的权柄许可，并藉主的助佑；如此，既遵守了这一惯例，宗座得以保持其权力，同时也不削减它已授予他人的权利。\par

\SourceRecord{Translated by James Barmby.；Nicene and Post-Nicene Fathers, Second Series；Vol. 12.；Edited by Philip Schaff and Henry Wace.；Buffalo, NY: Christian Literature Publishing Co.,；1895.；<http://www.newadvent.org/fathers/360203030.htm>.}

% structure: p0001 source=h1 target=section reason=page-title

\section{第三卷 第三十一封信}

\Emph{致贵族\Person{罗马努斯}。}\par

额我略致贵族\Person{罗马努斯}，意大利总督。\par

我们相信阁下已知悉\Place{米兰}教会主教\Person{老楞佐}去世的消息。由于据我们从圣职团的报告所知，众人都一致同意选举我们的儿子、该教会执事\Person{君士坦提乌斯}，因此我们有必要，为遵守古老惯例，派遣我们教会的一名士兵，以促使他（根据他的发现，众人意愿与同意一致）由其本教区主教祝圣，正如古老惯例所要求，尽管仍须经我们同意。为此，我们以应尽的父爱致候阁下，恳请阁下，在正义认可下，无论上述\Person{君士坦提乌斯}当选与否，一旦有需要，赐予支持；以使这一服务既在此生敌人面前抬高您，又将来在生命之前向天主推荐您。因为他是我的属下，曾与我交往甚密。而您应将我们所知道属于我们的人视为属于您自己，并特别爱护他们。\par

\SourceRecord{Translated by James Barmby.；Nicene and Post-Nicene Fathers, Second Series；Vol. 12.；Edited by Philip Schaff and Henry Wace.；Buffalo, NY: Christian Literature Publishing Co.,；1895.；<http://www.newadvent.org/fathers/360203031.htm>.}

% structure: p0001 source=h1 target=section reason=page-title

\section{第三卷，第三十二封信}

\Emph{致总执事\Person{奥诺拉图斯}。}\par

\Person{额我略}致\Place{萨洛纳}总执事\Person{奥诺拉图斯}。\par

我们与前任的命令不久前已送达你的爱德，其中你被免除那些诬告你的指控；我们命令你毫无争议地复职于你的品级。然而，因为没过多久，你来到罗马城，抱怨你们中间有关圣器转让的不当行为，而且当我们这里有人可能回应你的反对意见时，你的主教\Person{纳塔利斯}却去世了，我们认为有必要通过本信进一步确认那些同样的命令，即我们前任和我们自己的（如已所述）不久前为你的无罪而发出的命令。因此，我们完全免除你所有对你的指控，我们希望你毫无争议地继续留在你的品级中，以至于上述那人所提出的问题不得以任何借口使你受到丝毫损害。此外，关于你的控诉要点，我们已严格命令\Person{安东尼努斯}副执事，他在你们地区担任我们所管理的圣教会产业的总管（在天主眷顾下），如果发现教会人员牵连其中，他应以最大的严厉和权威审理这些案件。但是，如果涉及的是教会司法权之效力无法触及之人，他应将每项要点下的证据存入公共档案，并毫不迟延地传送给我们，以便我们在准确知情后，知道如何藉基督的帮助处理此事。\par

\SourceRecord{Translated by James Barmby.；Nicene and Post-Nicene Fathers, Second Series；Vol. 12.；Edited by Philip Schaff and Henry Wace.；Buffalo, NY: Christian Literature Publishing Co.,；1895.；<http://www.newadvent.org/fathers/360203032.htm>.}

% structure: p0001 source=h1 target=section reason=page-title

\section{第三卷，第三十三封信}

\Emph{致\Person{Dynamius}，贵族。}\par

\Person{额我略}致高卢贵族\Person{Dynamius}。\par

忠实管理他人财物的人，表明他如何善于管理自己的财物。您的荣耀向我们显明这一点：您专注于每年的奉献，已将您收入的果实呈献给真福宗徒之长\Person{伯多禄}。您忠实地将他应得的给他，您就使这些礼物成为您自己的。因为确实，那些思想永恒光荣的世上显贵人物，应该如此行动，凭借他们优越的世俗权能，为自己获得非暂世性的赏报。因此，我们向您致以应尽的问候，恳求全能的天主以现世的美物充实您的生命，并将它延展至永恒崇高的喜乐。因为我们经由我们的儿子\Person{Hilarus}（另作\Person{Hilarius}）收到了来自上述我们教会收入的四百\Place{高卢}\OriginalTerm{索利多}{solidi}。我们现在作为真福宗徒伯多禄的祝福，送您一个小十字架，其中嵌有来自他锁链的圣物，那锁链曾一时束缚他的脖颈：但愿它们永远解除您的罪过。此外，在它四周的四个部分，含有来自圣\Person{老楞佐}的烤架的圣物，他曾在那上面被焚烧，那烤架，他的身体因真理而被火吞噬，愿它点燃您的灵魂，使之热爱上主。\par

\SourceRecord{Translated by James Barmby.；Nicene and Post-Nicene Fathers, Second Series；Vol. 12.；Edited by Philip Schaff and Henry Wace.；Buffalo, NY: Christian Literature Publishing Co.,；1895.；<http://www.newadvent.org/fathers/360203033.htm>.}

% structure: p0001 source=h1 target=section reason=page-title

\section{第三卷·第三十五封信}

\Emph{致副执事伯多禄。}\par

\Person{格列高利}致\Place{坎帕尼亚}副执事\Person{伯多禄}。\par

我们的兄弟同为主教的\Person{保禄}曾多次请求我们准许他返回自己的教会。我们见这请求合理，认为有必要同意他的请求。因此，你当召集那不勒斯教会的圣职人员，以便他们从他们中选出两三位，并务必派他们到这里来选举主教。但他们也应在致我们的信中说明，他们所派之人代表全体参与这次选举，以便他们的教会能有效祝圣自己的主教。因为我们不能再容许它长期没有自己的治理者。倘若他们试图以任何方式忽视你的劝诫，就以教会纪律的权威对待他们。因为凡不自愿同意此举者，将证明自己的乖戾。此外，请将一百\Emph{\OriginalTerm{索利多}{solidi}}和一个由他亲自挑选的孤儿赠予上述我们的兄弟同为主教的\Person{保禄}，作为他为同一教会劳苦的酬劳。此外，劝诫那些代表全体来此选举主教的人，要记住他们必须带来所有主教礼服，以及他们认为被选为主教者自己所需使用的足够金钱。但不可延误地派遣那些按我们所说选出的圣职人员，因为现时有那不勒斯城的诸位贵族在此，我们可以与他们商议选举主教之事，并在主的助佑下共同筹谋。\par

\SourceRecord{Translated by James Barmby.；Nicene and Post-Nicene Fathers, Second Series；Vol. 12.；Edited by Philip Schaff and Henry Wace.；Buffalo, NY: Christian Literature Publishing Co.,；1895.；<http://www.newadvent.org/fathers/360203035.htm>.}

% structure: p0001 source=h1 target=section reason=page-title

\section{第三卷，第三十六封信}

\Emph{致\Person{撒比诺}，\OriginalTerm{守护者}{Defensorem}。}\par

\Person{额我略}致\Person{撒比诺}，\Place{撒丁岛}的守护者。\par

某些严重事项传入我们耳中，需要按教规予以纠正；因此，我们责成你的经验，切勿疏忽，必须尽快促使我们的兄弟和同僚主教\Person{雅努雅留斯}，连同书记官\Person{若望}，排除一切借口，前来我们面前，以便在他面前，对所报告给我们的事情进行彻底调查。此外，如果虔信妇女\Person{蓬佩雅纳}和\Person{特奥多西亚}依其请求愿意前来此地，请你在各方面予以支持，使她们能藉你的帮助实现其愿望；但特别要尽一切努力，按他所请求的，带上最有口才的\Person{依西多尔}，以便在详尽审查他已知对\Place{卡拉利}教会提出的案件的是非曲直之后，他能依法了结此案。\par

此外，有人报告给我们关于司铎\Person{厄丕法尼}的一些个人不当行为，你必须认真调查一切，并尽快将据说与他一同犯罪的那些妇女，以及你认为可能了解此事内情的其他人，一并带来；以便真相大白，以符教会纪律之严正。\par

现在你要留心如此有效地完成这一切，以免因疏忽而招致任何责难，须知如果我们的这道命令在任何方面被懈怠执行，你必将自担风险。\par

\SourceRecord{Translated by James Barmby.；Nicene and Post-Nicene Fathers, Second Series；Vol. 12.；Edited by Philip Schaff and Henry Wace.；Buffalo, NY: Christian Literature Publishing Co.,；1895.；<http://www.newadvent.org/fathers/360203036.htm>.}

% structure: p0001 source=h1 target=section reason=page-title

\section{\WorkTitle{Book III, Letter 38}}

\Emph{致利伯廷努斯总督。}\par

\Person{额我略}致西西里总督\Person{利伯廷努斯}。\par

从你执政之初，天主就愿意你出去维护祂的事业，并出于祂的仁慈为你保留了这项报酬，伴随赞誉。因为据报，一个名叫\Person{纳萨斯}的极邪恶的犹太人，以应受惩罚的胆大妄为，在真福\Person{厄里亚}的名义下设立了一座祭台，并以亵圣的诱惑引诱许多基督徒在那里敬拜；而且，据说他还获得了基督徒奴隶，将他们用于自己的服务和利益。既然如此，他本该因如此重大的罪行受到最严厉的惩罚，但荣耀的\Person{尤斯蒂努斯}，如所写信告知我们的，被贪婪的诱惑所安抚，迟迟没有为对天主所犯的伤害复仇。然而，愿你的荣耀对所有这些事情进行严格调查，如果发现这些事情确已发生，请务必最严格和身体力行地惩罚这个邪恶的犹太人，这样你既能为自己赢得天主的恩宠，也能以这个榜样，为你自己的赏报，向后人展示典范。此外，根据法律的命令，应立即毫无含糊地释放他所获得的任何基督徒奴隶，以免（愿天主不许）基督徒的宗教因受制于犹太人而被玷污。因此，请以最快速度最严格地纠正这些事情，这样我们不仅会因你的纪律而感谢你，而且在需要时也能为你的良善作证。\par

\SourceRecord{Translated by James Barmby.；Nicene and Post-Nicene Fathers, Second Series；Vol. 12.；Edited by Philip Schaff and Henry Wace.；Buffalo, NY: Christian Literature Publishing Co.,；1895.；<http://www.newadvent.org/fathers/360203038.htm>.}

% structure: p0001 source=h1 target=section reason=page-title

\section{第三卷，第四十五封信}

\Emph{致安得烈主教}\par

\Person{额我略}致\Place{塔兰托}主教\Person{安得烈}[\Emph{塔兰托，在卡拉布里亚}]\par

一个人只有意识到自己的罪过，并努力以相称的补赎来平息永恒审判者的愤怒，才能无所畏惧地面对祂的审判台。我们现在确知你曾有一妾；关于此人，也引起了某些人的不利猜疑。但既然在可疑之事上不应作绝对判断，我们选择将此事交予你自己的良心。那么，如果你在领受圣秩后，记得曾被肉欲交合所玷污，你就必须辞去司铎之尊位，绝不可再妄自接近其职务的执行，因为你知道你若明知此罪却还隐瞒真相，执意留在现有圣秩中，你将冒着灵魂的危险举行圣事，并且无疑要在我们的天主面前交账。因此我们再次劝告你，如果你知道自己曾被古代仇敌的诡计所欺骗，就当趁你还来得及，以充分的补赎赶紧战胜他，免得，如我们所愿不发生的，你在审判之日与他同被定罪。然而，如果你并不自觉有罪，你就需要在现有圣秩中继续。\par

此外，由于你违反正当秩序，将一名在教会名册上的妇女判处用棍棒残酷殴打，虽然我们不认为她在八个月后因此死亡，但因为你并未尊重你的圣秩，我们判决你两个月内不得举行弥撒。同时，暂停你的职务，你应该为你的所作所为哀哭。因为理固宜然：既然值得称颂的司铎的榜样未能激励你采取适合你身份的宁静正直，那么至少矫正的药方应该迫使你这样做。\par

\SourceRecord{Translated by James Barmby.；Nicene and Post-Nicene Fathers, Second Series；Vol. 12.；Edited by Philip Schaff and Henry Wace.；Buffalo, NY: Christian Literature Publishing Co.,；1895.；<http://www.newadvent.org/fathers/360203045.htm>.}

% structure: p0001 source=h1 target=section reason=page-title

\section{第三卷，书信四十六}

\Emph{致若望主教。}\par

格里高利致加利波利（\Emph{即卡拉布里亚的加利波利}）主教若望。\par

从贵兄弟呈报我们的报告中得知，我们的兄弟和同为主教的安得烈无疑曾有一名姘妇。但是，由于不确定他在领受圣秩后是否与她有染，你务须以恳切的劝诫警告他：若他自知在领受圣秩后曾与她交合，就当辞去所担任的职务，不再履行圣职。而且，如果他明知自己做了这事，却隐瞒罪恶而擅自履行圣职，就当知道在神圣审判中他的灵魂将面临危险。\par

关于他在教会册籍上登记的那名妇女，他曾下令用棍棒责打她；尽管我们不相信她八个月后因此而死，然而，由于他如此惩罚她不符其神圣召叫，你要暂停他两个月举行弥撒圣祭，好让这羞辱至少教训他将来如何行事。\par

再者，前述主教的神职人员在提交给我们的请愿书（附于下方）中声称，他们遭受他的许多虐待。因此，请贵兄弟务必准确查明所有这些事，并以合理方式予以纠正和安排，使他们日后无需因此事再来此地。七月，小纪第十一年。\par

\SourceRecord{Translated by James Barmby.；Nicene and Post-Nicene Fathers, Second Series；Vol. 12.；Edited by Philip Schaff and Henry Wace.；Buffalo, NY: Christian Literature Publishing Co.,；1895.；<http://www.newadvent.org/fathers/360203046.htm>.}

% structure: p0001 source=h1 target=section reason=page-title

\section{第三卷，第四十七封书信}

\Emph{致\Place{萨洛纳}教会的圣职人员}\par

\Person{格列高利}致圣职人员等。\par

蒙爱的诸位，我们读了你们的书信，得知你们已选择了你们的\Person{Honoratus}总执事；我们知道，你们选择一位久经考验、生活庄重的人担任主教职，这完全令我们喜悦。我们也与你们一同认可他的个人品德，因他已为我们所知；我们也同样希望他按你们的意愿被祝圣为你们的司铎。为此，我们劝勉你们毫无歧义地坚持对他的选举。任何情况都不应使你们对他本人改变心意，因为这项值得赞扬的选举既已获认可，那么，若有任何人（愿天主阻止）引诱你们将爱从他身上转移，这便将给你们的心灵加诸重担，并在你们的声誉上留下不忠的污点。至于那些在此所愿的选举中不与你们同心的人，我们已令我们的副执事\Person{Antoninus}劝诫他们，使他们能与你们一致。我们也已就我们的弟兄及同为主教的\Person{Malchus}其人应如何处理，向他下达了训令。然而，因我们自己也已写信给他，我们相信他会立即保持安静，不再烦扰你们。万一他以任何方式疏忽而不服从，他的顽抗必将受到最严厉的教会惩罚之制裁。\par

\SourceRecord{Translated by James Barmby.；Nicene and Post-Nicene Fathers, Second Series；Vol. 12.；Edited by Philip Schaff and Henry Wace.；Buffalo, NY: Christian Literature Publishing Co.,；1895.；<http://www.newadvent.org/fathers/360203047.htm>.}

% structure: p0001 source=h1 target=section reason=page-title

\section{第三卷 第四十八封书信}

\Emph{致哥伦布主教}\par

\Person{额我略}致\Person{哥伦布}等人。\par

甚至在收到您的来信之前，我就因您应有的声誉而知道您是\Person{天主}的忠仆。如今收到来信，我更清楚地明白，早先传扬的名声确非虚传；我极为您的功绩而欣喜，因为您展现了值得称赞的生活与行为。既然我感受到这些恩赐是由至高威严的\Person{天主}所赐予，我向您祝贺；并称颂我们的施恩者\Person{天主}，祂从不拒绝将祂仁慈的恩赐给予谦卑的仆人。因此，我宣布这是真实的：您的兄弟之爱以爱德的火焰点燃我爱您之心，我的灵魂与您如此结合，以至于我虽然身在远方，却渴望见到您，心中也与您同在。您自己也能体会到这一点。因为事实上，在爱德中心灵的结合比肉体的临在更有力量。再者，您以整个心意、整个心灵、整个灵魂依附并献身于宗座，我如今已确知，正如您的信函见证之前，我已清楚知道。因此，首先以应有的爱德问候您，我劝勉您不要忘记您曾向真福\Person{伯多禄}、宗徒之长所许下的承诺。\par

因此，你当催促你教务会议的首席主教，绝不可让孩童领受圣秩，以免他们更快地攀登高位，反而更危险地跌倒。在圣职授任中不可有买卖圣职：任何人的权势或恳求都不应违背我们的这些禁令而获得任何东西。因为毫无疑问，如果有人不是因功绩，而是因偏袒（愿\Person{天主}禁止）或买卖圣职而被晋升至圣秩，\Person{天主}必然不悦。\par

因此，如果你知道这些事情正在发生，不要沉默，而要坚决反对；因为倘若你疏忽它们，或虽知情却隐瞒，罪的锁链不仅会束缚那些行事的人，而且在这件事上，\Person{天主}面前轻的罪责也会触及你。因此，如果任何类似的事情发生了，就应当以教会法的惩罚加以制止，免得如此大的恶行，伴随着别人的罪，因纵容而增强力量。\par

因此，我较早地允许了这位呈信者——你的执事\Person{维克托利努斯}——离开，我认为他是你的效仿者，我以爱德接待了他；通过他，我向你传送了真福\Person{伯多禄}的钥匙作为祝福，其中包含了他锁链的一部分。\par

最后，关于你在\Person{天主}之下努力召集的会议的合一与和平，请让你的爱德通过详细告知我一切来使我的心灵喜乐。\par

\SourceRecord{Translated by James Barmby.；Nicene and Post-Nicene Fathers, Second Series；Vol. 12.；Edited by Philip Schaff and Henry Wace.；Buffalo, NY: Christian Literature Publishing Co.,；1895.；<http://www.newadvent.org/fathers/360203048.htm>.}

% structure: p0001 source=h1 target=section reason=page-title

\section{第三卷 第四十九封书信}

\Emph{致\Person{阿德奥达图斯}主教。}\par

\Person{额我略}致\Place{努米底亚}省首席主教\Person{阿德奥达图斯}。\par

您的来信清楚表明了爱德之情如何使阁下与我们联系在一起；它们给我们带来了极大的喜乐，因为我们发现这些信函充满了仁爱精神，洋溢着悦乐天主的感情。正如我们简要提到的，您致我们的信函如此坦露了您的心迹，以致可以认为写信人从未曾离开我们。的确，那些情感不与彼此爱德相左的人不应被视为缺席。虽然如您在信中所说，您的体力和年龄都不允许您来到我们这里，使我们能因阁下的亲身到场而欣慰，但既然我们在情感上彼此合一，我们便完全彼此相临，因我们在因爱德而合一的心灵中彼此看见。此外，我们以合宜的爱德之情问候阁下，并劝勉您全心致力于在您于天主之下所担任的首席职位上智慧行事，使您既因达致这一品秩而灵魂获益，又能成为他人未来效法的良好榜样。\par

因此，要特别谨慎对待圣秩授予；绝不可让任何未年长且行为纯洁的人向往圣秩，以免他们因过早追求而永远失去应是的样貌。因为你们必须先检查那些要担任任何圣职者的生活和品行；为了使你们能够接纳配得上此职的人，不要让任何人的影响力或恳求诱骗你们。但首先必须谨慎，绝不能让买卖圣职在圣秩授予中有立足之地，以免（天主不许）更大的危险笼罩受祝圣者和祝圣者。因此，如果有时需要处理此类事务，要召集稳重的、有经验的人进入您的咨询会，并和他们共同商议。在所有其他人之前，您应当总是邀请我们的兄弟和同僚主教\Person{哥伦布}。因为我们相信，如果你按照他的建议行事，没有人会在任何方面找到您的过错；且要知道，这在我们看来就如同按照我们的建议行事一样可接纳；因为他的生活和品行在各方面都如此令我们满意，以致所有人都清楚知道，凡经他同意所行之事，必不会被任何瑕疵的污点所遮蔽。但本函的携带者，我们上述同僚主教的执事\Person{维克托里努斯}，如此传扬了您的功德，以致大大振奋了我们对您行为的精神。我们恳求全能的主使关于您的好名声在悦乐祂的行动中更加显扬。因此，当您正在筹备的会议在天主助佑下结束后，请以它的团结与和谐令我们欢喜，并告知我们所有情况。\par

\SourceRecord{Translated by James Barmby.；Nicene and Post-Nicene Fathers, Second Series；Vol. 12.；Edited by Philip Schaff and Henry Wace.；Buffalo, NY: Christian Literature Publishing Co.,；1895.；<http://www.newadvent.org/fathers/360203049.htm>.}

% structure: p0001 source=h1 target=section reason=page-title

\section{第三卷 第五十一封信}

\Emph{致\Person{马克西米安}主教。}\par

\Person{额我略}致\Place{锡拉库萨}主教\Person{马克西米安}。\par

那些与我朝夕相处的弟兄们极力催促我简要写一些我们在意大利所听说的关于教父们所行奇迹的事。为此我非常需要您的爱德之助，请简短地告诉我您所回忆起的以及您所知道的事。因为我记得您曾告诉我一些关于\Person{诺诺苏斯}院长大人的事，他当时与\Place{彭托米}的\Person{阿纳斯塔修斯}大人在一起，但我现在已忘记了。因此我恳求您立刻写信把这些或其他事情告诉我，如果您自己也不打算短期内来访的话。\par

\SourceRecord{Translated by James Barmby.；Nicene and Post-Nicene Fathers, Second Series；Vol. 12.；Edited by Philip Schaff and Henry Wace.；Buffalo, NY: Christian Literature Publishing Co.,；1895.；<http://www.newadvent.org/fathers/360203051.htm>.}

% structure: p0001 source=h1 target=section reason=page-title

\section{第三卷，第五十三封书信}

\Emph{致若望主教。}\par

额我略致君士坦丁堡主教若望。\par

虽然案情促使我思考，但爱德也催迫我写信，因为我曾一再致信给我最圣洁的兄弟主若望，却未收到他的回信。因为另有一个世俗之人冒用他的名字给我写信；倘若那些信确是他的，那我就太不警觉了，竟相信了他远非我后来所发现的那样。我曾就最可敬的司铎若望的案子以及依撒乌利亚的隐修士们的问题写信，其中一位已领受司铎职的隐修士在你们的教堂里遭到棍棒殴打；而你的最圣洁的弟兄之爱（从信上的署名来看）回信给我，声称对我所提及之事一无所知。对此答复我极为惊异，心中默默思量：若他说的是实话，那么还有什么比这更糟糕的事，就是如此对待天主的仆人，而近在咫尺的人竟毫无知觉？因为若狼吞吃羊而牧人不知道，牧人还有什么借口？但是，如果你的圣洁既知道我信中所提之事，也知道对司铎若望和依撒乌利亚的隐修士兼司铎阿塔纳修所做之事，却又写信给我说不知道，那我对此如何应答？因为真理通过其圣经说：\BibleQuote{说谎的口杀死灵魂}（参见\BibleRef{智 1:11}）。我恳求你，最圣洁的兄弟：难道你那如此伟大的节制竟到了这种地步，以致你想通过否认来向你的兄弟隐瞒你所知道已做的事？难道肉进入口中作食物，不是比谎言从口中出来欺骗邻人更好吗？尤其是当真理说：\BibleQuote{不是入口的污秽人，而是出口的污秽人}（参见\BibleRef{玛 15:11}）？但愿你最圣洁的心远离此类想法。那些信上署着你的名字，但我不认为它们是你的。我曾写信给最蒙福的主若望；但我相信是那个你的亲信回复的——那个年轻人，他迄今尚未学到任何关于天主的事；他不知道爱德的怀抱；他在恶行中受众人谴责；他每日用隐藏的遗嘱对许多人的死亡设下陷阱；他不敬畏天主，也不尊重人。请相信我，最圣洁的兄弟，你必须先纠正此人，好使那些亲近你的人因你的榜样而得改进。不要听从他的舌头：他应当遵照你的圣洁的劝告行事，而非你的圣洁受他言语的左右。因为，若你听他，我知道你不能与你的弟兄们保持平安。因为我，正如我的良心为我作证，不愿与任何人争吵；我尽一切努力避免争吵。尽管我极其渴望与全人类平安相处，但尤其与你，我极其爱你，只要你自己仍是我所认识的那个人。因为，若你不遵守教规，并企图撕毁教父的规条，我不知道你是谁。因此，最圣洁、最亲爱的兄弟，请如此行事，使我们彼此相认，免得古时的仇敌激起我们两人之间的冒犯，借其最残酷的胜利杀害许多人。至于我，为表明我不以骄傲的精神行事，若我之前提及的那个年轻人在你的弟兄之爱中并非行恶之首，我本可暂时对教规中随手可得的规定保持沉默，并放心地将最初来到我这里的人送回给你，因为我知道你的圣洁会以爱德接待他们。但即使现在，我仍要说：要么接纳这些人，恢复他们的职位，让他们安享宁静；要么，如果你不愿这样做，就按教父的规条和教规的定义来处理他们，不要与我争论。但若你两者都不做，我们固然不愿引发争吵，但如果你方挑起，我们也不回避。此外，你的弟兄之爱很清楚教规关于以鞭打恐吓人的主教的规定。因为我们被立为牧人，而非迫害者。杰出的宣讲者说：\BibleQuote{宣讲、劝勉、斥责，以百般的忍耐和教训}（参见\BibleRef{弟后 4:2}）。而以鞭打强求信德，这是前所未闻的宣讲。但这些事我不必在信中说太多，因为我已派我最亲爱的儿子、执事撒彼尼亚努斯作为我在教会事务中的代表，到我们主上的门槛前；他将更详细地与你就一切事交谈。除非你倾向与我们争论，你会发现在一切公义之事上他有所预备。我将他托付给你的蒙福之身，好让他至少能在我于皇城中所认识的主若望那里找到他。\par

\SourceRecord{Translated by James Barmby.；Nicene and Post-Nicene Fathers, Second Series；Vol. 12.；Edited by Philip Schaff and Henry Wace.；Buffalo, NY: Christian Literature Publishing Co.,；1895.；<http://www.newadvent.org/fathers/360203053.htm>.}

% structure: p0001 source=h1 target=section reason=page-title

\section{第三卷，第五十六封书信}

\Emph{致\Person{若望}主教。}\par

\Person{额我略}致\Place{拉文纳}主教\Person{若望}。\par

不久前，有人向我们提及有关您兄弟的一些事情，我们记得当\Person{卡斯托里乌斯}，我们主持的圣教会的公证员，前往您所在地区时，我们已经充分表明了自己的看法。因为我们听说，在您的教会中正发生一些违反习俗和谦逊之道的事情；而您很明白，唯有谦逊能高举司铎的职位。如今，若您的智慧曾友善地或以主教应有的严肃态度接受我们的劝诫，您就不该为此动怒，而应感谢我们，并改正这些行为。因为，即使是（但愿不会发生）不公正的指责，若无最耐心的容忍，也违背教会惯例。\par

但是，您兄弟过于激动了；当您内心膨胀，仿佛要为自己辩护时，您写道：\Quote{您只是在教会的子弟们被从圣器室遣散之后、在弥撒之时以及在庄严的祷文期间才使用\OriginalTerm{披肩}{pallium}；}您以极其明显的事实承认，自己僭用了某些违背普世教会惯例的东西。因为，在灰烬和麻衣的时期，在街市上、在人群的喧嚣中，您竟能合法地做您在穷人贵族的集会中、在教会的圣器室里声称是非法的行为，这怎么可能呢？然而，最亲爱的兄弟，我们认为这并非您所不知：在世界任何地方，几乎从未听说过任何一位总主教，除了在弥撒之时，为自己要求使用披肩的权利。而您对普世教会的这一惯例非常了解，您以您的书信最清楚地表明了这一点；在那些书信中，您附带寄来了我们可敬的前任教宗\Person{若望}的法令，其中规定我们前任所赐予您和您教会的所有特权性惯例都应当保留。那么，您承认普世教会的惯例是不同的，因为您声称凭特权有权做您所做的事。因此，我们认为，此事的疑虑已经消除。因为，要么所有总主教的惯例也应由您兄弟遵守，要么，若您说您的教会获得了特别许可，则应由您方出示罗马城前任教宗的法令，证明这些权利已被授予\Place{拉文纳}教会。但是，若未能出示，那么，既然您声称做这些事既非基于普遍惯例，亦非基于特权，则证明您所做的是僭越。而且，最亲爱的兄弟，我们将对未来的审判者说什么呢？如果我们为了（我不说教会性的，而是）某种世俗的尊荣而辩护那沉重的轭和链子挂在我们颈上的使用，认为自己若在短时间内没有如此重的负担就是降低了身份？我们渴望以披肩装饰自己，而我们的品行可能毫无装饰；然而，主教的颈上，没有比谦逊更光辉夺目的了。\par

因此，如果您兄弟坚决要用任何论据来维护您的荣誉，您有责任要么遵循普遍惯例（无需书面授权），要么凭书面展示的特权为自己辩护。或者，最终若两者皆不为，我们将不允许您为其他总主教树立此类僭越的榜样。但是，以免您或许认为我们在此信中忽略了兄弟爱德所要求之事，请知悉：我们已在我们档案中仔细查找了您教会的特权文件。确实，找到了一些东西，完全足以驳斥您兄弟的主张，但毫无有助于您教会在所涉问题上的争辩。因为，即使就您用以反对我们的您教会的那项惯例——我们之前写过，应由您方证明——我们要让您知道，我们已经充分考虑了此事，询问了我们的儿子们：执事\Person{伯多禄}、\OriginalTerm{首席书记官}{primicerius}\Person{高迪奥苏斯}，以及我们教座的\OriginalTerm{监护人}{defensorem}\Person{弥额尔}，还有其他人，他们曾受我们前任的派遣，因各种任务前往\Place{拉文纳}；他们非常肯定地否认，这些事曾在他们面前做过。因此，很显然，暗中进行之事必定是非法的僭越。因此，隐秘引入之事毫无坚实依据可使其持续。那么，您或您的前任擅自多做的事，请您本着爱德和兄弟般的善意，努力纠正。切勿试图——我不是说您主动，而是即使以他人甚至您的前任为榜样——偏离谦逊的规范。简言之，总结我上面所说的，我劝诫您：除非您能证明这曾由我的前任以特权方式许可给您，否则您不得再在街市上使用\OriginalTerm{披肩}{pallium}，以免您甚至失去在弥撒时使用它的权利——因为您竟敢在街市上僭用它。至于您在圣器室中坐着、戴着披肩接待教会的子弟们（这件事您兄弟既做过又否认过），我们现在暂且不予追究；因为遵循公会议的决定，我们拒绝惩罚被否认的小过失。然而，我们知道这事发生过一次两次，我们禁止再发生。但您兄弟要小心留意，以免开始时被宽恕的僭越，若继续下去，会受到更严厉的惩罚。\par

此外，你曾抱怨拉文纳城中某些司铎级人员涉及严重的刑事指控。我们希望你或就地调查他们的案件，或把他们送到这里来（除非因地点遥远而难以取证），以便在此审理。但是，如果他们倚仗权势人物的庇护——我们不信会如此——而藐视你的审判或拒绝前来，并且顽固地拒绝回应针对他们的指控，我们希望你第二次和第三次警告之后，禁止他们行使神圣职务的职权，并将他们的顽抗书面报告给我们，以便我们商议你该如何彻底查究他们的行为，并按照教规定义予以纠正。因此，让你的兄弟情谊知道：在这案件中我们完全免于责任，因为我们已经将彻底调查此事委托给你；如果他们的所有罪行未受惩罚，这次调查的全部重担就会危及你的灵魂。而且你要知道，亲爱的，如果你不用教规严厉的极致去纠正你神职人员的越轨行为，如果你让任何已被证实有这些越轨行为的人继续亵渎神圣的圣秩，你在未来的审判中将没有任何借口。\par

另外，你为支持你的神职人员使用餐巾一事所做的辩护，遭到我们这边神职人员的强烈反对，他们说这从未被准许给任何其他教会，而且拉文纳的神职人员，无论是在那里还是在罗马城，据他们所知，也从未以任何这样的方式擅自使用过；或者，如果曾以秘密篡夺的方式尝试过，那也不构成先例。但是，即使任何教会曾有过这样的擅自行为，他们断言也应该予以纠正，因为这并不是罗马教宗的恩准，而仅仅是一种偷偷摸摸的僭越。然而，为了保全你的兄弟情谊的荣誉，尽管我们上述神职人员不情愿，我们仍然允许你的首位执事（有人向我们证明他们以前使用过）在伺候你的时候使用餐巾。在其他任何时候或由任何其他人使用餐巾，我们严加禁止。\par

\SourceRecord{Translated by James Barmby.；Nicene and Post-Nicene Fathers, Second Series；Vol. 12.；Edited by Philip Schaff and Henry Wace.；Buffalo, NY: Christian Literature Publishing Co.,；1895.；<http://www.newadvent.org/fathers/360203056.htm>.}

% structure: p0001 source=h1 target=section reason=page-title

\section{第三卷，第五十七封书信}

\Emph{拉文纳主教若望致教宗额我略}\par

我最尊敬的同仆、你宗座的公证人\Person{卡斯托里乌}，将我主的信函交给了我，这信函由蜜与毒混合而成；然而它如此深植其刺，仍为医治措施留有余地。因为我主在责备骄傲并谈及随之而来的天主审判时，以某种方式宣称自己理应是温和与平静的。\par

那么，你断言我贪图新奇，僭用了超出前任所允许的披肩使用。愿我主的良心——它为天主右手所治理——绝不允许自己相信这一点；也愿他不要将最神圣的耳朵向常见传闻的不确定性敞开。首先，因为我虽是罪人，仍然知道违反教父们为我们设定的界限是多么严重的事，以及一切高傲只会导致跌倒。因为，若我们的祖先在君王面前也不能容忍骄傲，那么在司铎身上就更不可容忍了！其次，我记得我如何在你至圣的\Place{罗马}教会怀抱中受哺育，并藉天主的助佑在其内成长。我又怎敢如此胆大，竟反抗那向普世教会传递其法律的最神圣宗座呢？为维护其权威，正如天主所知，我已为自己招致了许多敌人的恶感。但我最可称颂的主，请勿以为我曾试图做任何违反古老习俗的事，正如许多甚至几乎全城居民都可作证，而上述最可敬的公证人，即使他未参与程序，也可能作证，因为，直到教会的子女从圣器室下来，执事们正进来立即前往祭台时，首位执事一直习惯于为\Place{拉文纳}教会的主教披上披肩，而主教也习惯于同样在庄严的祷文中使用它。\par

因此，愿无人试图在我主面前暗示任何反对我的话，因为若有人想这样做，他无法证明我曾引入任何新事。因为我如何服从你的命令并在需要时为你的利益服务，愿全能的天主向你至诚挚的心显明：我归咎于自己的罪过，致使在经历了内外如此多的劳苦与困难之后，竟配得上经历这样的转变。但这事又给予我安慰：至圣的教父们有时惩罚他们的儿子，目的只是为了使他们更加进步；而且，在这番忠诚与赔补之后，你不仅会为圣\Place{拉文纳}教会保全其古老特权，甚至会在你自己的时代赐予更多。\par

至于手帕——你的宗座声称我的司铎和执事使用手帕是一种僭越——我实言相告，提及此事令我烦恼，因为单凭真理本身（它唯独在我主面前有效）就已足够。因为既然这被允许给围绕城市设立的较小教会，那么我主的宗座也将能完全查明——若他俯允询问其首要宗座的尊敬神职人员——是否\Place{拉文纳}教会的司铎或肋未人每次为主教祝圣或事务来到\Place{罗马}时，他们都在你的至圣前任眼前带着手帕\Emph{前行}而未受任何指责。因此，当我这罪人由你的前任在那里祝圣时，我的所有司铎和执事在侍奉教宗大人时也使用手帕\Emph{前行}。既然我们的天主以其上智将万物交于你手中和最纯洁的良心，我以你所曾以品格装饰、如今以应得之尊严治理的宗座本身恳求你，不要因我的功德而减少与你密切相关的\Place{拉文纳}教会的特权；反而，如先知之言，愿这罪归在我和我父家身上，按其所应得。因此，为让你更加满意，我附上了你的前任们赐予圣\Place{拉文纳}教会的所有特权，尽管在你可敬的档案中，关于我前任祝圣时期的记录亦足为凭。但如今，在你查明真相后，无论你命令做什么，都在天主和你权下；因为我既愿服从我主宗座的命令，已留意不顾古老习俗，在接到进一步指示前停止使用。\par

\SourceRecord{Translated by James Barmby.；Nicene and Post-Nicene Fathers, Second Series；Vol. 12.；Edited by Philip Schaff and Henry Wace.；Buffalo, NY: Christian Literature Publishing Co.,；1895.；<http://www.newadvent.org/fathers/360203057.htm>.}

% structure: p0001 source=h1 target=section reason=page-title

\section{第三卷，书信59}

\Emph{致 \Person{塞昆迪努斯} 主教。}\par

\Person{格列高利} 致 \Person{塞昆迪努斯}，\Place{陶罗梅尼乌姆} 主教。[ \Emph{位于 \Place{西西里}}]。\par

我们之前曾下令，应将圣洗池从位于 \Place{马斯卡莱} 上方的 \Person{圣安德肋} 隐修院中移走，以免对隐修士造成不便，并应在现今圣洗池所在之处设立祭台。但这项命令的执行迄今被推迟。因此，我们劝诫你的兄弟之爱，在收到此信后不得再有任何拖延，而应将圣洗池本身填平，并立即在那里设立祭台，用以举行神圣奥迹；以便前述隐修士能更安全地自由举行天主的工程，并免得我们的心因疏忽而对你兄弟之爱动怒。\par

\SourceRecord{Translated by James Barmby.；Nicene and Post-Nicene Fathers, Second Series；Vol. 12.；Edited by Philip Schaff and Henry Wace.；Buffalo, NY: Christian Literature Publishing Co.,；1895.；<http://www.newadvent.org/fathers/360203059.htm>.}

% structure: p0001 source=h1 target=section reason=page-title

\section{第三卷，第六十封书信}

\Emph{致贵妇\Person{意大里加}}\par

\Person{额我略}致\Person{意大里加}等。\par

我们已经收到你的来信，其中充满甘饴，并欣喜地得知阁下安好。就我们自己的心意而言，由于父爱，我们不愿怀疑在那平静之下隐藏着任何潜在的恶意。愿全能天主成就此事，使我们如你所愿，以善意相待，而你亦以善意回应我们，并在你的行为中彰显你在言语中所表达的甜美。因为表面上最荣耀的健康与美丽，若内里有隐藏的疮伤，便毫无益处。而外在的和平为之提供护卫的不和，更应加以提防。至于阁下在上述书信中费心提醒我们的事情，请记得你已收到书面告知：我们不会就穷人的诉讼问题与你作任何可能引起冒犯或公众喧哗的解决。我们记得曾写信给你表明此意，并且，在天主助佑下，我们也知道如何以教会之节制，避免法律诉讼的争执，并按照那话，欣然忍受我们财产的剥夺。但我们想你应该知道：我们的沉默与忍耐，不会损害我之后未来教宗在穷人事务上的权益。因此，为履行我们上述的承诺，我们已决定对这些争议保持缄默；也不愿亲自介入这些交易，因为我们感到其中所显示的善意甚少。但是，荣耀的女儿，免得你因此以为我们完全放弃了关乎和睦之事，我们已训示即将前往\Place{西西里}的执事、我们的儿子\Person{西彼廉}：若你能以有益的方式、且不损害你灵魂地处理这些事务，他应凭我们的权威与你解决，而我们不再为这事务感到烦扰，因为它可借此友好地了结。如今，愿全能天主——祂深知如何使完全不可能之事成为可能——启迪你们两人，既以和平之心处理你们的事务，又为了你灵魂的益处，在涉及本教会穷人的事务上顾念他们的利益。\par

\SourceRecord{Translated by James Barmby.；Nicene and Post-Nicene Fathers, Second Series；Vol. 12.；Edited by Philip Schaff and Henry Wace.；Buffalo, NY: Christian Literature Publishing Co.,；1895.；<http://www.newadvent.org/fathers/360203060.htm>.}

% structure: p0001 source=h1 target=section reason=page-title

\section{第三卷，第六十五封信}

\Emph{致毛里修斯·奥古斯都}。\par

\Signature{格列高利致毛里修斯等}。\par

在全能天主面前，凡在对我们最安详的主上们的一切言行上未能保持纯洁无过者，就是有罪的。然而，我，陛下不配的仆人，在此陈情，既非以主教的身份，也非以共和国权利中的臣仆，而是以私人身份，因为，最安详的主上，从您尚未成为万主之主时，您便是属于我的。\par

当最杰出的郎吉努斯，御马官（\OriginalTerm{马夫}{stratore}），抵达此处时，我收到了主上们的法律。当时我因身体疾病疲惫不堪，未能作出任何回复。其中，主上们的虔敬规定，任何担任公职者都不得担任教会职务。对此我大为称赞，因为最明显的证据表明，急于脱离世俗身份而进入教会职务者，并非希望放弃世俗事务，而是变换之。但，同一法律又说不许他成为隐修士，这令我全然惊讶，因为他的账目可以通过隐修院呈报，其债务也可以从他入住的处所追索。因为不论一个人怀着多么虔诚的意愿想要成为隐修士，都应当先归还其不法所得，并且越是卸下负担，便越真实地为灵魂着想。同一法律又附加规定，凡手上受过烙印者不得成为隐修士。我向主上们坦白，这条法令令我大为惊恐，因为它堵住了许多人通往天堂之路，并且使迄今为止合法的事变为非法。因为有许多人即使在世俗状态也能度虔诚生活，但更有许多人除非放弃一切，否则无论如何都不能在天主内得救。但我这样对主上们说话，是什么？不过是尘土和蛆虫而已。然而，既然感到这条法令反对天主——万有的创造者，我便不能向主上们保持沉默。因为天上已经赐予主上们的虔敬统管全人类的权柄，目的就是为了帮助追求善的人，使通往天堂的道路更为宽广，使地上王国侍奉天上王国。看，明文规定说，一个人一旦被烙印为地上战士，除非他完成了服役或因身体孱弱被拒，否则不得充当我主耶稣基督的士兵。\par

对此，看，基督藉着我——祂和您的最末仆人——将回答，说：从书记升你为侍卫伯爵；从侍卫伯爵升你为\OriginalTerm{凯撒}{Cæsar}；从凯撒升你为皇帝；不仅如此，还使你成为众皇帝之父。我将我的司铎交在你手中；而你却要将你的士兵从我的服役中撤回吗？最虔诚的主上，我恳求你，回答你的仆人；当你的主来临并如此说话时，你将如何回答祂？\par

但或许有人相信，他们当中没有人是怀着纯洁的动机成为隐修士的。我，你不配的仆人，知道在我有生之年，有多少成为隐修士的士兵行了奇迹，施行了征兆和大事。但这条法律却禁止哪怕是这样的人中的一个成为隐修士。\par

我恳请我的主上查问，哪一位前任皇帝曾经颁布过这样的法律，并且更彻底地考虑，它是否应当被颁布。并且，这确实是一个极为严肃的考虑：在这个时代，任何人被禁止离开世俗；而此刻世界的终结正临近。因为看！不再迟延：天火焚烧，地火焚烧，诸元素烈焰，天使与总领天使、宝座与宰制者、率领者与掌权者，可畏的审判者将要显现。倘若祂赦免一切罪过，只说这条法律是反对祂本人而颁布的，请问还有什么借口呢？因此，凭着同一位可畏的审判者，我恳求你，不要让我主的一切眼泪、一切祈祷、一切守斋、一切施舍在全能天主的眼前失去光彩；反之，愿陛下通过解释或修改来减弱这条法律的效力，因为主上们对抗敌人的军队，当天主的军队因祈祷而增加时，就越发壮大。\par

我实是服从您的命令，已将此法律传至世界各地；然而，既然此法律本身全然不合全能天主之意，看，我已藉此陈情向我的最安详主上们声明了这一点。这样，我两方面都尽了我的职责，既顺从了皇帝，也没有为天主之故而对我所感保持沉默。\par

\SourceRecord{Translated by James Barmby.；Nicene and Post-Nicene Fathers, Second Series；Vol. 12.；Edited by Philip Schaff and Henry Wace.；Buffalo, NY: Christian Literature Publishing Co.,；1895.；<http://www.newadvent.org/fathers/360203065.htm>.}

% structure: p0001 source=h1 target=section reason=page-title

\section{第三卷，第六十六封书信}

\Emph{致医生\Person{德奥多若}}\par

\Signature{额我略致\Person{德奥多若}等}\par

全能天主和我的最慈悲的主上皇帝所赏赐给我的恩惠，我的口舌无法充分表达。为这些恩惠，我应当如何回报呢？岂不是该真诚地爱慕他们的足迹？但是，由于我的罪过，我不知道出于谁的暗示或怂恿，去年他在他的共和国里颁布了这样一条法律，以至凡真心爱慕他的人必须极其悲伤。我当时因病未能回复这条法律。但我刚才向我的主上提出了一些建议。因为他下令：任何人若曾担任任何公职，或曾是军需官，或曾在手上作过记号，或被编入士兵名册，除非其兵役已满，均不得成为隐修士。据熟悉旧法律的人说，这条法律最初是由\Person{尤利安}颁布的，我们都知道他多么反对天主。如果我们的最慈悲的主上这样做，也许是因为许多士兵成为隐修士，致使军队减少，那么，难道是凭士兵的勇力，全能天主才将\Place{波斯}帝国征服于他吗？岂不单是因他的眼泪被俯听，天主以他自己所不知的旨意，将\Place{波斯}帝国降服于他的帝国吗？\par

现在，我觉得极其严厉的是：他竟阻止他的士兵事奉那位将一切赐给他、且使他不仅统治士兵甚至统治司铎的天主。如果他的目的是防止财产损失，那么那些收留了士兵的隐修院为何不能偿还债务，只保留人员从事隐修生活呢？因这些事情使我非常忧伤，我已向我的主上陈述了此事。但请你的光荣趁有利时机私下将我的陈述呈递给他。因为我不愿由我的代表（\OriginalTerm{代表}{responsalis}）公开呈递，既然你亲近地事奉他，可以更自由、更公开地谈论有益于他灵魂的事，因为他事务繁忙，不容易找到心无挂虑的时候。所以，光荣之子，你为基督发言吧。若你被听从，那将有益于你前述主上和你自己的灵魂。若你不被听从，你只益了你自己的灵魂。\par

\SourceRecord{Translated by James Barmby.；Nicene and Post-Nicene Fathers, Second Series；Vol. 12.；Edited by Philip Schaff and Henry Wace.；Buffalo, NY: Christian Literature Publishing Co.,；1895.；<http://www.newadvent.org/fathers/360203066.htm>.}

% structure: p0001 source=h1 target=section reason=page-title

\section{第三卷，第六十七封书信}

\Emph{致\Person{多米提安}，都主教。}\par

\Person{额我略}致\Person{多米提安}等。\par

接获您至甘饴的福分之函，我极大喜乐，因函中对我多论及圣经。我在其中寻得我所爱的美味，便贪婪地吞食。其中也杂有许多关于外在及必要事务的内容。您行事如同为心神预备筵席，使所献的美味因多样性而更令人愉悦。诚然，外在事务如低劣平常的食物，滋味较逊，但经您如此巧妙处理，令人欣然接受，因为连卑微的食物也常因善厨者的调味而变得甜美。如今，既保持史事的真实，我先前关于其神圣意义所言者，不应被拒绝。因为，虽然因您之意，其意义或许不适合我的情形，但按其上下文，从中所引出之言可毫不迟疑地持守。因为她的侵犯者（即狄纳的）被称为那地的首领\BibleRef{创 34:2}，显然是指魔鬼，因我们的救主说：\BibleQuote{现在这世界的首领要被赶出去}\BibleRef{若 12:31}。他也寻求她为妻，因为恶灵急忙要合法地占有那已被他隐秘诱惑首先腐化的灵魂。因此雅各伯的儿子们极其忿怒，拿刀剑攻击舍根全家和他的地方\BibleRef{创 34:25}，因为凡有热忱者，也当攻击那些成为恶灵帮凶的人。他们先命他们行割损，然后在他们疼痛时杀死他们。因为严厉的师长，若不知如何节制热忱，虽藉宣讲割除腐败的偏向，但当犯过者已为自己所做的恶而哀恸时，仍常在纪律的严酷中粗暴，过于所当行。因为那些已割去包皮的人本不应死，因为凡哀恸邪淫之罪、将肉体的欢乐转为忧愁者，不应从师长那里经历纪律的严酷，以免人类的救主反被爱得少，若为祂的缘故灵魂受的苦超过适当。因此雅各伯也对这儿子们说：\BibleQuote{你们扰乱了我，使我在客纳罕人眼中成为可憎的}\BibleRef{创 34:30}。因为当师长仍残酷地攻击犯过者已哀恸之事时，软弱的心对救主的爱便会冷淡，因为它感觉自己在那本不宽恕自己的事上受了苦。\par

因此，我愿说这些，为表明我所陈明的含义与上下文并非不合。但由您圣善者从同一段落为我安慰所推论出者，我欣然接受，因为对圣经的理解，凡不违背纯正信仰者，不应拒绝。因为，正如从同一黄金，有人制项链，有人制戒指，有人制手镯为装饰，同样，从同一圣经知识，不同的诠释者通过无数理解方式，仿佛制成各样装饰品，却都为装饰天上的新娘。再者，我极其喜乐，您至甘饴的福分虽忙于世俗事务，仍警醒地将才智用于理解圣经。因为诚然如此，若前者不能完全避免，后者也不应完全放下。但我恳求您因全能天主，向在如此巨浪磨难中劳苦的我伸出祈祷之手，使我因您的转祷被提升至高处，而我被自己罪恶的重压压至深处。此外，虽我悲伤波斯皇帝未归化，却完全喜乐，因您向他宣讲了基督信仰；因为，虽他不配来到光明中，您的圣善者仍将获得宣讲的赏报。因为埃塞俄比亚人进浴池时是黑的，出来仍是黑的；但浴池看守人仍得到工钱。\par

此外，关于\Person{马乌里基}，您说得对：从阴影我可认识雕像；即，在小事上我可审度大事。此事上我们信任他，因誓言和人质将他的灵魂与我们联结。\par

\SourceRecord{Translated by James Barmby.；Nicene and Post-Nicene Fathers, Second Series；Vol. 12.；Edited by Philip Schaff and Henry Wace.；Buffalo, NY: Christian Literature Publishing Co.,；1895.；<http://www.newadvent.org/fathers/360203067.htm>.}

% structure: p0001 source=h1 target=section reason=page-title

\section{第四卷 第一封信}

\Emph{致\Person{君士坦提乌斯}主教。}\par

\Person{额我略}致\Person{君士坦提乌斯}，\Place{米兰}主教。\par

收到您的信后，我向全能天主献上极大的感恩，因为我得以因您的祝圣庆典而获慰藉。确实，众人因天主的恩赐，一致同意您当选，这一事实您应当深思熟虑，因为，在天主之后，您大大亏欠那些以如此谦逊的态度希望您被置于他们之上的人。\par

因此，你应以司铎的仁爱回应他们的行为，并以亲切的同情心关注他们的需要。如果在他们中有任何过失，用深思熟虑的责备予以斥责，好使你的司铎之怒掺和着甜蜜的味道，并且如此你即使在被大大畏惧时也能被下属所爱戴。这种行为也会在他们心中引发对你本人的极大敬畏；因为，正如仓促和习惯性的愤怒被鄙视，所以有分辨地对过失发怒，在大多数情况下，因其缓慢而变得可怕。\par

此外，我们的副执事\Person{若望}已返回，他报告了许多关于您的美德，我们恳求全能天主亲自成就祂所开始的事；为使祂显示您无论内外在善德上都已进步，既在此世人间，也在将来天使之中。\par

此外，我们按惯例将一件\OriginalTerm{披带}{pallium}寄给你，用于神圣的弥撒礼仪。但我恳求你，当你接受它时，要以谦逊维护它的尊贵和意义。\par

\SourceRecord{Translated by James Barmby.；Nicene and Post-Nicene Fathers, Second Series；Vol. 12.；Edited by Philip Schaff and Henry Wace.；Buffalo, NY: Christian Literature Publishing Co.,；1895.；<http://www.newadvent.org/fathers/360204001.htm>.}

% structure: p0001 source=h1 target=section reason=page-title

\section{第四卷 第二封信}

\Emph{致君士坦提乌斯主教。}\par

\Person{格列高利}致\Place{米兰}主教\Person{君士坦提乌斯}。\par

执事\Person{波尼法爵}，我最亲爱的儿子，通过您的信向我传达了一些私下消息：即三位主教，不是找到了机会而是寻隙，从您虔诚的共融中分离出来，声称您已同意谴责\OriginalTerm{三章}{Three Chapters}，并给出了一份担保书。此外，您的兄弟之爱清楚记得，是否在任何言语或文字中提到过\OriginalTerm{三章}{Three Chapters}；虽然您的兄弟之爱的前任\Person{劳伦提乌斯}曾向\OriginalTerm{宗座}{Apostolic See}送交了一份极其严格的担保书，有合法数量的最高贵人士在上面签字；其中我也在当时担任罗马城行政长官时同样签了字；因为既然这场分裂并非因任何事由而起，\OriginalTerm{宗座}{Apostolic See}理应审慎，旨在从各个方面维护司铎心中\OriginalTerm{普世教会}{Universal Church}的合一。但至于据说我们的女儿，王后\Person{泰奥德琳达}，在听到这个消息后，已退出与您的共融，这无论如何都是显而易见的：尽管她被坏人的言语稍稍诱惑，但等公证人\Person{希波利特斯}和院长\Person{若望}到达后，她必将千方百计寻求与您的兄弟之爱的共融。我也给她写了一封信，恳请您的兄弟之爱尽快转交给她。此外，关于那些似乎已分离的主教，我写了另一封信，当您让他们看过后，我相信他们必会在您的兄弟之爱面前悔改其骄傲的迷信。\par

此外，您准确而简要地告知了我所发生的事，不管是由\Person{阿戈}王还是由法兰克诸王所为。恳请您的兄弟之爱务必让我知道您迄今所查明的一切。但是，如果您看到伦巴第人的王\Person{阿戈}与\OriginalTerm{贵族}{Patrician}没有采取任何行动，请以我们的名义向他承诺，我准备关注他的事，如果他愿意与共和国有利地达成任何协议。\par

\SourceRecord{Translated by James Barmby.；Nicene and Post-Nicene Fathers, Second Series；Vol. 12.；Edited by Philip Schaff and Henry Wace.；Buffalo, NY: Christian Literature Publishing Co.,；1895.；<http://www.newadvent.org/fathers/360204002.htm>.}

% structure: p0001 source=h1 target=section reason=page-title

\section{第四卷 第三封信}

\Emph{致主教\Person{君士坦提乌斯}}\par

\Person{额我略}致\Place{米兰}主教\Person{君士坦提乌斯}。\par

现已获悉，你教区内某些主教，寻找而非找到了借口，企图脱离你手足之谊的合一，声称你曾在罗马城出具保证，表示你谴责了三章。而事实上，他们这样说是因为他们不知道我如何习惯于即使没有保证也信任你的手足之谊。因为如果有任何此类需要，单凭你口头之言便可信任。但我并不记得我们之间曾以言语或文字提及三章之事。至于他们，如果他们及早从错误中回头，就应宽恕他们，因为按宗徒\Person{保禄}的话：\Emph{\BibleQuote{他们不明白自己所说的，也不知道自己所肯定的。}} \BibleRef{弟前 1:7}。因为我们，在真理引导和良心见证下，声明我们完整无缺地持守神圣\Place{加采东}大公会议的信德，并且不敢在其定义上增加或删减任何东西。但是，如果有人胆敢自作主张，思虑任何或多或少的与之相反之事，并反对同一会议的信德，我们毫不犹豫地予以绝罚，并判定他脱离慈母教会的怀抱。因此，凡我这宣信不能使其归正者，不再爱\Place{加采东}会议，而是憎恨慈母教会的怀抱。所以，那些似乎如此胆大妄为之人，若曾因灵魂的热忱而如此发言，那么在他们接受这解释之后，就应返回你手足之谊的合一，不使自己与\Person{基督}的身体——即圣而公教会——分离。\par

\SourceRecord{Translated by James Barmby.；Nicene and Post-Nicene Fathers, Second Series；Vol. 12.；Edited by Philip Schaff and Henry Wace.；Buffalo, NY: Christian Literature Publishing Co.,；1895.；<http://www.newadvent.org/fathers/360204003.htm>.}

% structure: p0001 source=h1 target=section reason=page-title

\section{卷四，书信四}

\Emph{致泰奥德琳达王后。}\par

额我略致伦巴第人的王后泰奥德琳达。\par

据某些人的报告，我们得知，在几位主教的诱导下，阁下竟到了对圣教会犯下如此冒犯的地步，以致脱离了公教会同心合意的共融。如今，我们越是真诚地爱您，就越为您感到忧心忡忡：您竟然相信那些既不懂自己所谈论之事，又几乎无法理解所听到之事的无知愚人。\par

因为他们说，在虔诚纪念的犹斯提尼安时代，曾制定了一些违反加采东大公会议的规定；而他们自己既不阅读，也不相信那些阅读的人，便一直停留在他们自己就我们而臆想出来的同一错误中。因为我们——我们的良心作证——声明：关于这同一神圣加采东大公会议的信德，没有任何更改，没有任何侵犯；相反，在上述犹斯提尼安时代所做的一切，其做法都是为使加采东大公会议的信德不受任何干扰。此外，如果有人胆敢说或想任何事情违反该会议的信德，我们以绝罚介入，痛斥其见解。既然您在我方良心的见证下知道我们信德的纯正，那么，您绝不应使自己与公教会的共融分离，以免您所有的眼泪和一切善功归于乌有，如果发现它们与真信德相悖的话。因此，阁下应尽快致函我最可敬的弟兄和同僚主教龚思当休——他的信德及品行我早已深知——并以致他的信函表明您欣悦地接受了他的祝圣，并且您绝不与他的教会共融分离；虽然我认为我对此所说的已是多余：因为，尽管您心中曾有过某种程度的疑虑，但我相信，在我儿若望院长和依玻里多公证员到达时，这疑虑已从您心中消除了。\par

\SourceRecord{Translated by James Barmby.；Nicene and Post-Nicene Fathers, Second Series；Vol. 12.；Edited by Philip Schaff and Henry Wace.；Buffalo, NY: Christian Literature Publishing Co.,；1895.；<http://www.newadvent.org/fathers/360204004.htm>.}

% structure: p0001 source=h1 target=section reason=page-title

\section{第四卷 第五封信}

\Emph{致\Person{邦尼法爵}主教。}\par

\Person{额我略}致\Person{邦尼法爵}，\Place{雷吉乌姆}主教（\Emph{\Place{雷伊}}）。\par

司铎在有关神圣敬礼的事上受到劝诫，是件可耻的事。因为那时他们反被要求去做他们本应要求别人做的事。然而，为了免得阁下在有关天主的工作上有所疏忽（虽然我并不认为如此），我们认为特别就这一点劝诫你是合适的。因此，我们告诫你：绝不可因阁下的宽纵，让\Place{雷吉乌姆}城的圣职人员在他们的职务所要求的职责上有所松懈。反而，在有关天主的事上，要最立即、最热切地督促他们。我们也希望你留意上述圣职人员的声誉，以免传出任何不好的、违反教会纪律的事；因为他们的职务属于教会的装饰，而非污秽的行为。此外，我们规定：我们在西西里人所决定的事，也要由你的\OriginalTerm{副执事}{subdeacons}们遵守；也不可容许任何人因傲慢或鲁莽而违反我们的这项决定；这样，正如我们相信将会如此，上述一切被你最严格地执行，你就既不会成为我们劝诫的违犯者，也不会因对所托付的牧职规则有所懈怠而被指控有罪。\par

\SourceRecord{Translated by James Barmby.；Nicene and Post-Nicene Fathers, Second Series；Vol. 12.；Edited by Philip Schaff and Henry Wace.；Buffalo, NY: Christian Literature Publishing Co.,；1895.；<http://www.newadvent.org/fathers/360204005.htm>.}

% structure: p0001 source=h1 target=section reason=page-title

\section{第四卷 第六封信}

\Emph{致执事西彼廉}\par

\Person{额我略}致\Person{西彼廉}，执事及\Place{西西里}总代理。\par

有人向我们报告，\Place{卢卡尼亚}省有一名女子，名叫\Person{伯多禄尼拉}，通过\Person{阿格内鲁斯}主教的劝勉而皈依；她虽拥有自己财产的管理权，却仍通过特别的赠与契据，将全部财产赠与她所进入的隐修院。此外，上述主教去世时，将其一半财产留给了他的儿子\Person{阿格内鲁斯}，据称他是我们教会的公证员，另一半则给了上述隐修院。然而，因\Place{意大利}面临的灾祸，他们逃往\Place{西西里}避难，上述\Person{阿格内鲁斯}却据说败坏她的道德并玷污了她，发现她怀孕后，诱使她离开隐修院，并带走了她所有的财物，包括她自己的以及她可能从他父亲那里得到的一切；在犯下如此重大的罪行之后，他竟声称这些东西归他所有。因此，我们敦促你的仁爱，将上述男子和上述女子立即带到你面前，并对该案进行最彻底的调查。如果你发现情况正如向我们报告的那样，就用最严厉的净化手段裁决这桩被如此众多不义玷污的案件；以使严格的报应降临到上述男子身上，他既不念及自己的身份也不念及她的处境；同时，她先受惩罚并被送往隐修院补赎，然后所有从上述地方取走的财产，连同其所有果实和增益，都应该归还。\par

\SourceRecord{Translated by James Barmby.；Nicene and Post-Nicene Fathers, Second Series；Vol. 12.；Edited by Philip Schaff and Henry Wace.；Buffalo, NY: Christian Literature Publishing Co.,；1895.；<http://www.newadvent.org/fathers/360204006.htm>.}

% structure: p0001 source=h1 target=section reason=page-title

\section{第四卷·第七封书信}

\Emph{致\Person{格纳狄乌斯}，贵族。}\par

\Person{格列高利}致\Person{格纳狄乌斯}，贵族及\Place{非洲}总督。\par

我们确信，您宗教尊贵的心智因神圣之爱的热忱，对那些尤在教会内以不妥当方式所做之事倍感愤慨。因此，我们更加乐意将教会案件中过错的纠正交托给您，因为我们信赖您虔诚性情的倾向。那么，请尊驾知晓：据一些从\Place{非洲}地区来到我们这里的人报告，在\Place{努米底亚}议会中，许多事情正违反教父们的方式和教规的条例而行。由于无法再忍受屡次传到我们耳中关于此类事情的投诉，我们已将调查这些事委托给我们的弟兄和同为主教的\Person{哥伦布}——他的庄重名声，现已广为人知，不容我们怀疑。为此，以慈父般的情感问候您，我们劝告您的尊贵：在所有涉及教会纪律的事务中，您应向他提供您支持的援助，免得所做的不妥之事若不被调查和惩罚，就会借着长期存在的先例，以更大的放肆增长为未来的越轨行为。此外，最卓越的儿子，要知道，如果您寻求胜利，并为所托付给您行省的安全而努力，那么对于此目的，没有什么比竭尽可能约束司铎的生活和教会内部的战争更为有益。\par

\SourceRecord{Translated by James Barmby.；Nicene and Post-Nicene Fathers, Second Series；Vol. 12.；Edited by Philip Schaff and Henry Wace.；Buffalo, NY: Christian Literature Publishing Co.,；1895.；<http://www.newadvent.org/fathers/360204007.htm>.}

% structure: p0001 source=h1 target=section reason=page-title

\section{Book IV, Letter 8}

\Emph{致\Person{雅努阿里乌}主教。}\par

\Person{格列高利}致卡拉利斯（\Place{Cagliari}）主教\Person{雅努阿里乌}。\par

我们确实认为，您的地位本身已足以促使您奋力履行虔敬的职责。但是，为避免任何懈怠降低您的热忱，我们认为理应就这些事特别劝诫您。我们获知，您那里的\Person{斯德望}在离世时，通过遗嘱指示建立一座隐修院。但据说，由于继承人、可敬的\Person{特奥多西亚}女士的拖延，他的愿望至今未实现。因此，我们劝勉您的兄弟之爱对此事给予最大关注，并告诫上述女士，以便在一年内，她按指示建立隐修院，并毫无争议地依逝者的意愿完成一切。但若她因疏忽或狡诈而推迟实现这一计划（例如，若她无法在指定地点建立，而认为应另择他处，却因中间拖延而使此事被忽略），则我们希望由您的兄弟之爱尽力修建，并将遗留的财物和收益拨归这所可敬的处所。如此，你们既能在严厉的审判者面前免于懈怠的责罚，又能按我们最宗教的法律，以主教的热诚完成逝者被忽视的虔敬愿望。\par

\SourceRecord{Translated by James Barmby.；Nicene and Post-Nicene Fathers, Second Series；Vol. 12.；Edited by Philip Schaff and Henry Wace.；Buffalo, NY: Christian Literature Publishing Co.,；1895.；<http://www.newadvent.org/fathers/360204008.htm>.}

% structure: p0001 source=h1 target=section reason=page-title

\section{第四卷 第九封信}

\Emph{致\Person{雅奴雅里}主教}\par

额我略致\Place{卡辣利}（\Place{卡利亚里}）主教\Person{雅奴雅里}。\par

牧灵的热忱本身应当足以激励你，即使没有我们的帮助，也能有益而审慎地保护你所托管的羊群，并用勤勉的谨慎使之免于敌人的诡计。但是，因为我们发现你的爱德还需要我们权威的书面之言来增强你的坚定，我们必须借着兄弟之爱的劝勉，加强你向宗教热忱的懈怠倾向。\par

现在，我们得知你疏忽了对撒丁岛上那些天主的婢女们的隐修院的监护。尽管你的前任们曾明智地安排，由某些经认可的圣职人员负责照料她们的需要，但如今这已完全被忽视，以至于特别献身于天主的妇女被迫亲自前往公职人员那里处理赋税及其他义务，不得不奔波于村庄和农场之间以缴纳捐税，并不得体地混杂在属于男人的事务中。你的兄弟情谊应通过简单的纠正消除这一弊端：即谨慎地委派一位生活与品行经得起考验，且年龄与地位不会引起任何邪恶猜疑的人，他应怀着对天主的敬畏，协助这些隐修院的成员，使她们不再因任何私人或公共的理由，违反会规，走出她们可敬的院墙之外；而所有需要为她们办理的事，都由你所委派的人合理处理。至于修女们自身，应赞美天主，安守隐修院内，不再让信友们心中产生任何邪恶的猜疑。倘若她们中有人，或因先前的放纵，或因邪恶的免罚恶习，已被诱惑，或将来被引入奸淫失足的深渊，我们要求，在遭受相应刑罚的严厉之后，她被送往另一座更严格的童贞隐修院做补赎，在那里她可以致力于祈祷和斋戒，借着忏悔使自己受益，并成为更严格纪律的榜样，足以使他人畏惧。此外，任何与这类妇女一同被发现在不法行为中的人，若是平信徒，应被剥夺共融；若是圣职人员，也应撤职，并因他永远可悲的放纵而被投入隐修院。\par

我们也希望你每年召开两次主教会议，正如据称是你们教省的惯例，也为神圣教规的权威所命令；这样，若他们中有任何人的道德品质与其身份不符，他就能因弟兄们的友好规劝而被谴责，同时也能以父亲的谨慎采取措施，确保托付给他的羊群的安全和灵魂的得救。我们还获悉，犹太人的男奴和女奴，因信仰而逃到教会寻求庇护者，要么被归还给其不信的主人，要么按其价值被赎回以代替归还。因此，我们劝告你绝不可允许如此恶劣的习俗继续下去；而任何身为犹太人奴隶并逃到可敬之处寻求庇护的人，你不得使其遭受任何损害。相反，无论他之前已是基督徒，还是现在才受洗，都应借着教会怜悯的庇护，支持他对自由的诉求，而不让穷人受到任何损失。\par

主教不得擅自在已受洗的婴儿额上第二次傅圣化圣油；而司铎应在将受洗者的胸上傅油，以便之后由主教在他们额上傅油。\par

关于建立隐修院一事，既然不同的人已下令建造，若你发现任何被委以此责的人以不正当的借口拖延，我们希望你能按照法律的诫命精明地坚持，免得（愿天主不允许这样）因你的疏忽而使逝者的虔诚意愿受挫。此外，关于据称伯多禄先前下令在其房屋内建造的那座隐修院，我们已决定，你的兄弟情谊应准确调查那里的收入。若存在适当的供应，则在所有财产的减损及据称已被分散的部分追回后，应立即全力并无任何延迟地建立隐修院。但若资源不足或有害，我们希望你，在如所命令的那样仔细查考一切之后，向我们呈递一份报告，使我们得以知道，如何在天主的助佑下商议其建造事宜。那么，让你的兄弟情谊明智地注意上述各点，既不因违犯我们劝诫的要旨而被发现，也不因在牧职上太不热心而受天主的审判。\par

\SourceRecord{Translated by James Barmby.；Nicene and Post-Nicene Fathers, Second Series；Vol. 12.；Edited by Philip Schaff and Henry Wace.；Buffalo, NY: Christian Literature Publishing Co.,；1895.；<http://www.newadvent.org/fathers/360204009.htm>.}

% structure: p0001 source=h1 target=section reason=page-title

\section{第四卷·第十封书信}

\Emph{致达尔马提亚诸位主教}\par

\Person{额我略}致达尔马提亚诸位主教。\par

诸位兄弟本应出于对天主审判的敬畏而闭上肉性的眼睛，不应忽略任何有关天主和心灵正直的事，也不应因任何人的情面而偏离正义的直道。但如今你们的行为被世俗事务如此扭曲，以至于忘记了你们所拥有的司祭尊荣的全部路径和一切对天主的敬畏，只图谋取悦你们自己而不取悦天主，因此我们认为有必要向你们下达这些特别严格的书面命令，借着真福\Person{伯多禄}、宗徒之长的权威，我们命令：未经我们的同意和许可，不得在\Place{撒罗纳}城擅自为任何人覆手祝圣主教；也不得在该城祝圣任何人，除非按照我们所说的方式。\par

但是，如果你们出于自愿或因任何人的强迫，竟敢或试图做出任何违背此禁令的事，我们将判定你们被剥夺领受主圣体圣血的权分，从而使你们对事务的处理或违反我们命令的倾向，将你们与神圣奥迹隔绝，并且你们所祝圣的人不得被视为主教。因为我们不希望任何人被鲁莽地祝圣，如果其生活有可指责之处。因此，如果执事\Person{奥诺拉图斯}被证明是不配的，我们希望将可能当选者的生活和品行报告给我们，以便我们同意后可以妥善办理。\par

因为我们信赖全能天主，尽我们所能，绝不容许做任何损害我们灵魂的事；也绝不容许做任何损害你们的教会的事。但是，如果众人的自愿同意一致选定一个人，且因天主的恩宠他被证实堪当，并且无人反对其祝圣，我们希望你们在此\Place{撒罗纳}同一教会内，凭本信所授予的许可祝圣他；但排除\Person{马克西姆斯}这个人，关于他有许多恶报传到我们这里：除非他停止贪求更高的圣秩，依我看来，在充分调查后，他连现在所拥有的职位本身也应被剥夺。\par

\SourceRecord{Translated by James Barmby.；Nicene and Post-Nicene Fathers, Second Series；Vol. 12.；Edited by Philip Schaff and Henry Wace.；Buffalo, NY: Christian Literature Publishing Co.,；1895.；<http://www.newadvent.org/fathers/360204010.htm>.}

% structure: p0001 source=h1 target=section reason=page-title

\section{第四卷，第十一封信}

\Emph{致马克西米安主教。}\par

\Person{额我略}致\Place{锡拉库萨}主教\Person{马克西米安}。\par

很久以前，确实已经托付给你的贤弟，凭我们的权力，代为纠正西西里教会及其他可敬之处可能存在的任何过犯或不体面之事。但是，看到有申诉到达我们这里，说有些事情至今仍被忽视，我们认为应当再次特意激励你的贤弟去纠正它们。\par

因为我们得知，在教会新得的收益方面，关于其四份之一的规例并未得到遵行，反而各地方的主教只分发旧有收益的四份之一，将新获得的收益据为己用。因此，让你的贤弟积极纠正这种侵入的恶习，使无论是在以前的收益，还是现在或将来增加的收益中，四份之一都能按照规例分施。因为同一教会产业竟如同受两种不同法律——即强占法与圣规法——之评判，实属不当。\par

不可准许服务于教会的司铎、执事及其他任何品级的职员担任隐修院院长；而是让他们要么放弃圣职职务，晋升隐修会，要么，如他们决定继续担任院长之职，则绝不允许他们从事圣职工作。因为，如果一个人因其重要性而无法勤勉地履行其中任一职务，却被认为能胜任两者，从而使圣职秩序妨碍隐修生活，而隐修规则又反过来妨碍圣职效用，这是非常不合适的。此事我们业已考虑，提醒你的爱德：若某位主教离世，或因过犯而被撤职（愿天主免之），则应召集教长及全体圣职首要人员，在你的面前清点教会财产，准确登记所有发现之物，不得以实物或任何其他方式取走教会财产的任何部分，如据说以前曾以作为清点财产的酬劳为名发生过的那样。因为我们希望一切涉及保护穷人财产之事，都应如此执行，使在他们的事务中，为私利者之贪婪无可乘之机。\par

教会巡视员及其职员，当他们在不属于自己城市的堂区辛劳时，应按照你的指派，为他们的劳作领取一些补助。因为他们在提供服务的场合得到报酬是公正的。\par

我们最严厉地禁止年轻女子被立为女院长。因此，让你的贤弟不可允许任何主教授予贞女面纱，除非是年满六十岁的贞女，其年龄与品德当为此事所需；以便此事与上述各点，借主的助佑，因你严格要求的紧迫而得以纠正，你得以用圣规之纽带重新系住长久松弛的可敬之事，并且使神圣事务的安排，不是出自人不相称的意志，而是具有足够严格的规范。十月，小纪第十二年。\par

\SourceRecord{Translated by James Barmby.；Nicene and Post-Nicene Fathers, Second Series；Vol. 12.；Edited by Philip Schaff and Henry Wace.；Buffalo, NY: Christian Literature Publishing Co.,；1895.；<http://www.newadvent.org/fathers/360204011.htm>.}

% structure: p0001 source=h1 target=section reason=page-title

\section{第四卷，第十五封信}

\Emph{致雅努雅略主教。}\par

\Person{额我略}致\Place{卡利亚里}（\Emph{Cagliari}）主教\Person{雅努雅略}。\par

虔敬的妇人特奥多西亚，渴望实现其已故丈夫斯德望建造一座隐修院的意愿，恳求我们致函阁下，藉我们的推荐，她可更易蒙受您的援助。她声称，其丈夫曾指示应在名为皮斯切纳斯的田庄上建造隐修院，该田庄已归属已故主教多默的\OriginalTerm{旅舍}{Xenodochii}所有。尽管产业持有人允许她在非其自有土地上建院，但鉴于此举主有理由反对，我们以为宜同意其请求：即她愿赖主助佑，在其自称位于卡利亚里的自有宅院内，为天主的婢女建造一座隐修院。然而，因她述及该宅院已为宾客和访客所累，我们恳请阁下竭力以各种方式援助她，并赋予您保护的助力以成全其虔敬之心，致使您的协助与勤勉使您与她已故丈夫的热忱及她本人的奖赏有份。至于她请求供奉于该处的圣髑，我们愿由阁下以应有的敬意安置之。\par

\SourceRecord{Translated by James Barmby.；Nicene and Post-Nicene Fathers, Second Series；Vol. 12.；Edited by Philip Schaff and Henry Wace.；Buffalo, NY: Christian Literature Publishing Co.,；1895.；<http://www.newadvent.org/fathers/360204015.htm>.}

% structure: p0001 source=h1 target=section reason=page-title

\section{卷四 第十八封信}

\Emph{致院长玛吾尔}\par

额我略致玛吾尔等。\par

圣堂的关怀显然属于司铎职务，这促使我们必须如此用心，以免对圣堂出现任何疏忽之过。然而，我们得知曾托付给司铎的圣庞加爵圣堂屡遭疏忽，以致主日前来举行弥撒圣祭的人们因找不到司铎而抱怨离去。因此，经过成熟商议，我们决定撤换这些司铎，并赖天主的恩宠，在同一圣堂内为一隐修院建立一座隐修士团体，以使主持该隐修院的院长能对上述圣堂全面加以关注和照管。我们也认为适宜立你玛吾尔为这座隐修院的院长，并规定上述圣堂的土地以及凡已归其所有之物，或从其收入中所得之物，均归入你这隐修院，毫无减损地属于它；但条件是：上述圣堂内一切需施工或修缮之处，均由你负责完成，不得有误。\par

但为免在撤去原先托付于该圣堂的司铎后，圣堂似乎缺乏神圣礼仪的供应，我们特以此权令要你向其提供一位\OriginalTerm{外来的}{peregrine}司铎，以举行神圣弥撒圣祭；这位司铎必须同时居住在你的隐修院中，并从隐修院获得其生活供应。\par

但尤其要注意：在真福庞加爵至圣遗体之上，每日必行日课不可有缺。我们以本命令委托你实行的这些事，不仅你本人要执行，而且凡继承你职务地位者，亦应永远如此遵守并实现，以使今后上述圣堂再无疏忽之可能。\par

\SourceRecord{Translated by James Barmby.；Nicene and Post-Nicene Fathers, Second Series；Vol. 12.；Edited by Philip Schaff and Henry Wace.；Buffalo, NY: Christian Literature Publishing Co.,；1895.；<http://www.newadvent.org/fathers/360204018.htm>.}

% structure: p0001 source=h1 target=section reason=page-title

\section{第四卷，第二十封信}

\Emph{致\Person{马克西穆斯}，\OriginalTerm{篡位者}{Præsumptorem}}。\par

\Person{格列高利}致\Place{萨尔纳}的篡位者\Person{马克西穆斯}。\par

虽然任何人的生命功绩在其他方面无碍于其接受司铎圣职的祝圣，但谋求职位这一罪行本身便受到教规最严厉的谴责。如今我们得知，你或是以隐秘方式获取、或是假冒了最虔诚的君王的敕令，强行进入了本应人人敬重的司铎圣秩，而你的生活却不相称。我们毫不迟疑地相信了此事，因为你的生活和年龄并非我们所不知，更因为我们并非不了解我们最安详的主上皇帝的意向：他向来不习惯干预司铎的事务，以免因我们的罪过而受任何牵累。此外，还传出了一件前所未闻的恶行：即使在我们的禁令（在此禁令下，你及将祝圣你者均受绝罚）之后，据说你竟由军事力量扶持出场，并殴打长老、执事及其他神职人员。我们绝不可称此行为为祝圣，因为它是由被绝罚者举行的。因此，既然你毫无先例地侵犯了如此伟大的尊严——即司铎之职，我们责令：在我从我们的主上们或我们\OriginalTerm{代表}{responsalis}的来信中确知你是在真实而非隐秘的敕令下受祝圣之前，你及你的祝圣者不得擅自处理任何有关司铎职务的事，也不得接近神圣祭台的服务，直至你再次接到我们的通知。但若你胆敢违反此令，愿你受自天主及真福宗徒之长伯多禄的\OriginalTerm{绝罚}{anathema}，使你的惩罚成为其他公教会的鉴戒，因他们思量对你的判决。五月，小纪十二年。\par

\SourceRecord{Translated by James Barmby.；Nicene and Post-Nicene Fathers, Second Series；Vol. 12.；Edited by Philip Schaff and Henry Wace.；Buffalo, NY: Christian Literature Publishing Co.,；1895.；<http://www.newadvent.org/fathers/360204020.htm>.}

% structure: p0001 source=h1 target=section reason=page-title

\section{第四卷，第二十一封书信}

\Emph{致韦南修斯主教。}\par

格里高利致卢纳（在\Place{伊特鲁里亚}）主教\Person{韦南修斯}。\par

根据许多人的报告，我们得知居住在卢纳城中的犹太人拘留着基督徒奴隶。此事令我们愈感愤慨，因为你们的兄弟之谊对此容忍，令我们不安。依照您的职责及对基督宗教的敬意，您本应不留任何机会，让单纯的心灵（非经说服，而是以某种权威权利）服务于犹太人的迷信。因此，我们敦促您的兄弟之谊，依循最虔诚的法律所规定的途径，不许任何犹太人拥有基督徒奴隶。但若发现任何奴隶受其控制，应通过法律制裁的保护使其获得自由。至于犹太地产上的那些奴隶，尽管他们本身不受法律约束，但由于他们长期依附于土地，因租佃条件而有义务耕种，应允许他们继续耕种惯常的农田，向上述人士缴纳租金，并履行法律对佃农或本地居民所要求的一切，但不得再对他们施加额外负担。然而，若其中有人希望继续为奴，或有人愿意迁往别处，后者应自行考虑：他虽凭法律之力摆脱了奴役，却因自己的轻率而丧失了佃农的权利。在所有这些事情上，我们期望您如此明智地行事，既不让您成为一个失散羊群的有罪牧者，也不让您的懈怠在我们面前受到指责。\par

\SourceRecord{Translated by James Barmby.；Nicene and Post-Nicene Fathers, Second Series；Vol. 12.；Edited by Philip Schaff and Henry Wace.；Buffalo, NY: Christian Literature Publishing Co.,；1895.；<http://www.newadvent.org/fathers/360204021.htm>.}

% structure: p0001 source=h1 target=section reason=page-title

\section{第四卷 第二十三函}

\Emph{致巴尔巴里奇尼公爵奥斯皮托。}\par

额我略致奥斯皮托等。\par

既然你们种族中没有一个人是基督徒，我由此得知你比你们整个种族都优越，因为你在其中被发现是一名基督徒。因为，所有巴尔巴里奇尼都像无理性的禽兽般生活，不认识真天主，反而敬拜木石，而你敬拜真天主这一事实正表明你远胜过他们所有人。但你要以善行和言语实践你所接受的信德，并将你力所能及的奉献给你所信仰的基督，以便尽可能多的人被你带到祂面前，使他们受洗，并劝勉他们向往永生。如果你本人因忙于他事而无法亲自做到这些，我恳请你以我的问候，全力帮助我派往你们地区的那些人，即我的同僚主教斐理斯和我的儿子、天主的仆人基里亚库斯，借此在帮助他们的辛劳中，向全能的天主表达你的热忱，并使你所帮助其善工的祂的仆人在一切善行中成为你的帮助者。我们通过他们寄给你圣伯多禄宗徒的祝福，我恳请你如所应做的那样欣然接受。六月，第十二小纪。\par

\SourceRecord{Translated by James Barmby.；Nicene and Post-Nicene Fathers, Second Series；Vol. 12.；Edited by Philip Schaff and Henry Wace.；Buffalo, NY: Christian Literature Publishing Co.,；1895.；<http://www.newadvent.org/fathers/360204023.htm>.}

% structure: p0001 source=h1 target=section reason=page-title

\section{第四卷，第二十四封信}

\Emph{致\Place{撒丁岛}公爵\Person{扎巴达斯}。}\par

\Person{额我略}致\Person{扎巴达斯}等。\par

从我的兄弟及同主教\Person{斐理斯}，及天主的仆人\Person{基里亚库斯}的信函中，我们获悉了阁下的美德。我们向全能天主献上极大的感谢，因为\Place{撒丁岛}得到了这样一位公爵；他如此懂得如何在世间事务中向国家尽忠职守，也懂得如何向全能天主展示对天国的虔敬之情。因为他们曾写信给我说，你正在与\OriginalTerm{巴巴里奇尼}{Barbaricini}人商议和平条款，条件是要使这些巴巴里奇尼人归服于基督。为此我极其欢喜，而且，若蒙全能天主欣悦，我将迅速将你的贡献禀告我们最安详的君侯。因此，请完成你已开始的事业，向全能天主展示你心灵的虔诚，并竭尽全力帮助我们已派往你地区、为归化巴巴里奇尼人的人；须知这样的功业，无论在世俗君王面前，还是在天上君王眼中，都将大大有助于你。\par

\SourceRecord{Translated by James Barmby.；Nicene and Post-Nicene Fathers, Second Series；Vol. 12.；Edited by Philip Schaff and Henry Wace.；Buffalo, NY: Christian Literature Publishing Co.,；1895.；<http://www.newadvent.org/fathers/360204024.htm>.}

% structure: p0001 source=h1 target=section reason=page-title

\section{第四卷 第25封信}

\Emph{致\Place{撒丁岛}的贵族和业主。}\par

\Signature{\Person{额我略}致贵族等人。}\par

我从中获悉，我的兄弟和同为主教的\Person{费利克斯}以及我的儿子、天主的仆人\Person{齐亚库斯}的报告，得知你们几乎所有人都在自己的庄园里拥有崇拜偶像的农民（\OriginalTerm{农民}{rusticos}）。这使我非常难过，因为我知道属下的罪过连累他们上司的生命，并且，当属下的罪过未受纠正时，判决回落到那些管理他们的人身上。所以，卓越的儿子们，我劝你们，要全心全意地为你们的灵魂而热忱，并考虑你们将如何向全能天主为你们的属下交账。事实上，他们被托付给你们正是为此，既使他们为你们在世俗事务中服务，也使你们因对他们的照顾，为他们永恒中的灵魂提供所需。因此，如果他们向你们支付所欠的，你们为何不向他们支付你们所欠的呢？也就是说，阁下应殷勤地劝告他们，阻止他们脱离崇拜偶像的错误，以便通过使他们被吸引到信德，使全能天主对你们自己慈悲。因为，看，你们观察到这世界的结局近了；你们看到现在人的刀剑、现在天主的刀剑向我们发威；然而你们，真天主的崇拜者，却看到被托付给你们的那些人崇拜石头，而保持沉默。请问，在可畏的审判时，你们将说什么，当你们将天主的敌人掌握在手中，却蔑视征服他们使之归于天主、召回他们到祂那里？因此，向你们致以应有的问候，我恳求阁下热切地警醒，将自己奉献于对天主的热情，并赶快在你们的信中通知我，你们中何人已带领多少人归向\Person{基督}。那么，如果因任何原因你们无法这样做，请将此托付给我们上述的兄弟和同为主教的\Person{费利克斯}，或我的儿子\Person{齐亚库斯}，并为他们提供帮助以完成天主的工作，这样，在将来的报应中，你们将更能分享永生，因为你们现在为善工提供了帮助。\par

\SourceRecord{Translated by James Barmby.；Nicene and Post-Nicene Fathers, Second Series；Vol. 12.；Edited by Philip Schaff and Henry Wace.；Buffalo, NY: Christian Literature Publishing Co.,；1895.；<http://www.newadvent.org/fathers/360204025.htm>.}

% structure: p0001 source=h1 target=section reason=page-title

\section{第四卷·书信第二十六}

\Emph{致\Person{雅努亚里奥}主教。}\par

\Person{额我略}致\Place{卡拉利斯}（\Place{卡利亚里}）主教\Person{雅努亚里奥}。\par

我们从我们的同侪主教\Person{斐理斯}和院长\Person{奇里亚科}的报告得悉，在\Place{撒丁岛}上，司铎们受到世俗法官的压迫，而你的执事们蔑视你的兄弟之爱；而且，就目前所见，当你只追求纯朴时，纪律却被忽视了。因此我劝告你，排除一切借口，你要努力管理你所受托的教会，保持对圣职人员的纪律，不畏惧任何人的言语。但据我所闻，你已经禁止你的总执事与妇女同居，而到现在为止，你的这项禁令仍被藐视。除非他服从你的命令，我们的意愿是褫夺他的圣秩。\par

还有一件更令人痛惜的事，即：你的兄弟之爱的疏忽，竟容许属于圣教会的\OriginalTerm{农夫}{rusticos}至今仍留在不信中。如果我劝告你将那些不属于你的人引领到天主面前，而你自己却忽视拯救你自己的信徒脱离不信，那又有什么用呢？因此你必须在各方面警醒，促使他们悔改。因为，如果我在\Place{撒丁岛}上发现任何属于某位主教的异教农夫，我将严惩那位主教。\par

但现在，若有任何农夫被发现如此背信而顽固，竟拒绝前来归向主天主，他就必须以极重的税赋作负担，使他因被迫缴纳的痛苦而赶快走上正道。\par

我们亦得悉，有些领受圣秩者失足后（无论是在做完补赎之前或之后），又被召回他们的职务；这是我们完全禁止的，而且最神圣的教规也对此表示反对。因此，凡在领受任何圣秩后，陷入肉身之罪者，应丧失其圣秩，不得再近祭台从事职务。但是，为避免那些已受圣秩者永久沦丧，事先应谨慎考虑哪些人当受圣秩，即先应查看他们是否在多年生活中保持贞洁，是否关心诵读并喜爱施舍。也应询问此人是否曾两次结婚。还应查看他是否目不识丁，或对国家负有义务，以致在领受圣秩后被迫返回公职。因此，你的兄弟之爱应仔细查问这一切，使大家在经过认真审核后领受圣秩，无人易于在领受后遭免职。我们写给你的兄弟之爱的这些事，你要告知你下属的所有主教，因为我不愿亲自写信给他们，以免似乎削弱你的尊严。\par

我们还听说，有人因我们禁止司铎以圣油傅抹待洗者而感到不快。我们确实遵循了本教会的古老惯例：但如果确有人因此而感到不安，我们允许，在缺少主教的地方，司铎可以为待洗者傅抹圣油，甚至在额上傅抹。\par

\SourceRecord{Translated by James Barmby.；Nicene and Post-Nicene Fathers, Second Series；Vol. 12.；Edited by Philip Schaff and Henry Wace.；Buffalo, NY: Christian Literature Publishing Co.,；1895.；<http://www.newadvent.org/fathers/360204026.htm>.}

% structure: p0001 source=h1 target=section reason=page-title

\section{第四卷，第27封信}

\Emph{致\Person{雅努亚里}主教。}\par

\Person{格里高利}致\Person{雅努亚里}，\Place{卡拉利斯}（\Emph{卡利亚里}）主教。\par

你的兄弟情谊本当如此专注于虔敬的职责，以至于完全不需要我们的劝诫来促使你履行它们；然而，既然某些需要纠正的细节已为我们所知，那么另外有一封带有我们权威的信件寄给你，也并非不合宜之事。\par

因此，我们通知你，我们得知：在\Place{卡拉利斯}地区设立的\OriginalTerm{客舍}{Xenodochia}，向来有将账目详表定期呈报给该城主教的做法，也就是说，在主教的监护和照管下进行管理。现在，既然你的爱德被说成迄今忽略了此事，我们——如前所述——劝告那些在这些客舍中现在或曾经居住的人，应定期提交详细账目。并且应任命那些在生活、品行和勤勉方面最值得称许，且至少是\OriginalTerm{虔诚者}{religiosi}的人来管理它们，使法官无权烦扰他们，以免，若他们是能传唤到法庭的人，便可能造成浪费其所拥有的微薄资源的机会；关于这些资源，我们希望你极其留意，使之在无人知晓的情况下不被赠送给人，免得你的兄弟情谊的疏忽竟至于允许它们被掠夺。\par

此外，你知道，本信的携带者，长老\Person{埃皮法尼乌斯}，在某些撒丁人的信件中被刑事指控。我们按自己的意愿调查了他的案件，未发现对他指控的证据，便宣布他无罪，使他能恢复原位。因此，我们希望你追查对他指控的始作俑者；并且，除非发出那些信件的人准备用符合教规且最为严格的证据来支持他的指控，否则他绝不可接近圣体共融的奥秘。\par

另外，关于圣职人员\Person{保罗}，据说他多次被发现行为不端，并曾逃往\Place{非洲}，不顾圣职身份还俗；若是如此，我们已安排先对他进行体罚，然后交付补赎，以便根据宗徒的判词，通过肉身的折磨使灵魂得救，并且使他能用不断的眼泪洗去因邪恶行为而沾染的世俗罪恶污垢。\par

再者，根据教规的训令，任何\OriginalTerm{虔诚者}{religiosus}均不得与那些被暂停教会共融的人交往。\par

此外，对于圣职授任或圣职人员的婚姻，或蒙面贞女，任何人都不得擅取任何费用，除非他们自愿奉献某些东西。\par

至于那些离开隐修院还俗并嫁人的女子，应如何处理，我们已与你的兄弟情谊的上述长老详细交谈，你的圣德可从他的报告中获得更充分的信息。\par

再者，\OriginalTerm{虔诚的圣职人员}{religiosi clerici}应避免寻求或依赖在俗人士的庇护；而应在各方面按教规服从你的管辖权，以免因你的兄弟情谊的疏忽，导致你所治理的教会纪律涣散。\par

最后，关于那些与上述离开隐修院的女子犯罪、并据说现已被暂停共融的男子，若你的兄弟情谊察觉他们已为如此邪恶之事作了相称的忏悔，我们意愿你恢复他们与圣体共融。\par

\SourceRecord{Translated by James Barmby.；Nicene and Post-Nicene Fathers, Second Series；Vol. 12.；Edited by Philip Schaff and Henry Wace.；Buffalo, NY: Christian Literature Publishing Co.,；1895.；<http://www.newadvent.org/fathers/360204027.htm>.}

% structure: p0001 source=h1 target=section reason=page-title

\section{第四卷 第二十九封信}

\Emph{致雅努亚留主教。}\par

\Person{额我略}致\Place{卡腊利斯}（\Emph{卡利亚里}）主教\Person{雅努亚留}。\par

我们获悉，在撒丁省一个名叫福西亚纳的地方，据说曾有过祝圣主教的惯例；但因环境所迫，这惯例早已废弛。但如今，由于司铎短缺，我们得知有些外教人仍住在那里，过着野兽般的生活，完全不知敬拜天主，所以我们敦促你的友爱迅速按照古制在那里祝圣一位主教；也就是说，一位适合此项工作的人，并愿以牧灵热忱努力将迷途者引入主的羊群；如此，当他致力于拯救灵魂时，既不会让你发觉自己要求了多余的事，也不会让我们后悔自己枉然恢复了曾经中断的事。\par

\SourceRecord{Translated by James Barmby.；Nicene and Post-Nicene Fathers, Second Series；Vol. 12.；Edited by Philip Schaff and Henry Wace.；Buffalo, NY: Christian Literature Publishing Co.,；1895.；<http://www.newadvent.org/fathers/360204029.htm>.}

% structure: p0001 source=h1 target=section reason=page-title

\section{第四卷 第三十封书信}

\Emph{致君士坦丁娜・奥古斯塔}\par

格里高利致君士坦丁娜等\par

您虔诚之安详，以宗教热忱与圣德之爱著称，曾命令我向您送出圣保禄的头颅，或他遗体的其他部分，用于在皇宫中为恭敬同一圣保禄而建造的教堂。由于渴望从您那里接受命令，通过展现最迅速的服从，以此更激起您对我的眷顾，我更加痛苦，因为我既不能也不敢做您所吩咐的事。因为宗徒圣伯多禄和圣保禄的遗体在他们的教堂中闪耀着如此巨大的奇迹和可畏之事，以致一个人去那里祈祷都不可能不怀着极大的恐惧。简言之，当我有福记忆的前任教宗想要更换覆盖在有福宗徒伯多禄最神圣遗体上的银饰时，尽管距离同一遗体几乎十五英尺，一个不小的可怕征兆向他显现。不仅如此，我也以类似方式想要修改靠近圣保禄宗徒最神圣遗体不远处的某物；由于需要在他的坟墓附近挖到一定深度，那地方的监督发现了一些骨头，这些骨头确实不与同一坟墓相连；但是，由于他胆敢拿起它们并转移到另一个地方，一些可怕的征兆出现了，于是他突然死亡。\par

除此之外，当我有圣德记忆的前任教宗同样想要对殉道者圣老楞佐遗体不远处进行一些改进时，由于不知道可敬的遗体安放在哪里，在搜索过程中进行了挖掘，突然他的坟墓意外地被揭露；那些在场工作的隐修士和\OriginalTerm{堂宇管理员}{mansionarii}，他们看到了同一殉道者的遗体，虽然他们确实没有胆敢触碰，但都在十天内死亡，以致没有一人幸存，谁曾看见那义人的遗体。\par

此外，我最安详的夫人要知道，罗马人的习俗是，当他们赏赐圣徒的圣髑时，并不敢触碰身体的任何部分；而只是将一块布\OriginalTerm{圣布}{brandeum}放入一个盒子\OriginalTerm{圣盒}{pyxide}，并放置在最神圣的圣徒遗体附近；当它被取出时，它被以应有的恭敬存放在将要奉献的教堂中，于是那里产生的功效如此强大，就好像他们的遗体已被带到那特定地方一样。因此，在我有福记忆的教宗良的时代，据我们祖先传述，某些希腊人对这样的圣髑有所怀疑，上述教宗拿起剪刀剪了同一块布\OriginalTerm{圣布}{brandeum}，从切口流出了血。因为在罗马和整个西方地区，任何人若偶然想要触碰圣徒的遗体，都是不可容忍和亵渎神圣的；如果有人胆敢这样做，这鲁莽行为必定不会不受惩罚。为此我们大大惊奇于希腊人的习俗，他们自称取走圣徒的骨头；我们几乎不信。因为两年多前来到这里的某些希腊隐修士在夜晚的寂静中，在圣保禄教堂附近挖掘了露天躺着的死人尸体，并将他们的骨头保留在自己手中直到离开。当他们被抓获并受到仔细审讯为何这样做时，他们承认打算将这些骨头带到希腊冒充圣徒的圣髑。从这个例子中，正如已经说过的，我们心中产生了更大的怀疑，即他们是否真的如所声称的那样取走圣徒的骨头。\par

但我还能说什么关于有福宗徒们的遗体呢？众所周知，在他们殉难时，信徒们从东方来取回他们的遗体，视之为自己同胞的遗体。遗体被带到离城第二里程碑处，存放在一个叫作\OriginalTerm{卡特孔巴斯}{Catacumbas}的地方。但是，当全体群众聚在一起并试图从那里移走它们时，如此猛烈的雷电惊吓并驱散了他们，以致他们决不敢再尝试这样的事。然后，罗马人——因主的仁慈被认为配做此事——出去取走了他们的遗体，并将它们安放在现在存放的地方。\par

那么，最安详的夫人，有谁能如此鲁莽，既然知道这些事，还敢——我不说触碰他们的遗体——甚至连看它们一眼呢？因此，您给予我的这些命令，我本无法服从，在我看来，这并非出于您自己的意愿；而是某些人希望挑动您的虔诚反对我，以便夺去（天主禁止）您善意对我的眷顾，因此他们找出一个点，使我在那点上显得好像对您不服从。但我信赖全能的天主，您最仁慈的善意决不会从我这被偷走，并且您将永远享有圣宗徒们的德能，您全心全力爱他们，不是因他们的身体临在，而是因他们的保护。\par

此外，您同样命令送去的巾帕，与他的遗体在一起，因此不能被触碰，正如他的遗体不能被接近一样。但由于我最安详的夫人如此虔诚的愿望不应完全得不到满足，我将尽快送给您一些圣伯多禄宗徒亲自戴在颈上和手上的链条的部分，从这些链条中人们显示了众多奇迹；至少如果我成功用锉刀取下一些的话。因为，许多人经常来从这些链条上寻求祝福，希望得到一小部分锉屑，一位司铎带着锉刀等候，对于一些寻求者，链条上的一部分很快脱落，毫无延迟；但对于其他寻求者，锉刀在链条上长时间拉锉，却什么也得不到。六月，第十二小纪\par

\SourceRecord{Translated by James Barmby.；Nicene and Post-Nicene Fathers, Second Series；Vol. 12.；Edited by Philip Schaff and Henry Wace.；Buffalo, NY: Christian Literature Publishing Co.,；1895.；<http://www.newadvent.org/fathers/360204030.htm>.}

% structure: p0001 source=h1 target=section reason=page-title

\section{卷四，第三十一书}

\Emph{致御医德奥多禄}\par

格里高利致御医德奥多禄。\par

我亲自感谢全能天主，因为距离并不能分开真正彼此相爱者的心。因为，请看，最甘饴荣耀的儿子，我们身体相隔甚远，却因爱德而彼此相临。这是你的作为，你的书信所证实的，这是我与你同在时所体验的，这是你不在时我在你的荣耀中所承认的。愿这使你既为人所爱，又永远在全能天主面前相称。因为爱德是诸德之母，你结出善工之果，正是因为你心中存留着这果实之根。至于你所送来的、天主感动你为赎回俘虏而奉献的，我承认我既喜且忧地收下了。欢喜，是为了你，因我如此看出你正在天堂之乡预备居所；但极其忧伤是为了我自己，因为除了操心圣宗徒伯多禄的产业，如今还须为你——我最甘饴的儿子、德奥多禄大人——的产业交账，并要为是否妥善或疏忽地花费它而负责。但愿你心中倾注了祂自己仁慈心肠的全能天主，那赐你以卓越宣讲者论及我们救主所言——\BibleQuote{祂本是富有的，为了我们却成了贫穷的}\BibleRef{格后 8:9}——而忧虑挂虑的，愿祂在同一个救主来临时，使你显出在德行上是富有的，使你免于一切过犯，并赐你天国的喜乐代替地上的，永存的喜乐代替暂时的。\par

至于你所说的、你渴望在圣宗徒伯多禄最神圣的遗体旁为你所做的，请确信，即使你的舌头沉默，你的爱德也要求这样做。但愿我们配得为你祈祷：但我不配，我毫不怀疑。尽管如此，这里有许多配得的人，他们正因你的奉献而从仇敌手中被赎回，并忠实地事奉我们的造物主。关于他们，你做了所记载的事：\BibleQuote{把施舍存在穷人怀中，它会为你祈祷}\BibleRef{德 29:15}。\par

但是，既然爱得更多的人更为大胆，我对我最荣耀的儿子、德奥多禄大人最甘饴的性情有些怨言：即他从天主圣三接受了天才的恩赐、财富的恩赐、仁慈与爱德的恩赐，却仍不断纠缠于世俗的讼事，忙于不断的巡行，忽略每日阅读他救赎者的话语。因为圣经不就是全能天主给祂受造物的一封书信吗？确实，如果你的荣耀居于别处，并收到一位人间皇帝的来信，你绝不会拖延，不会休息，不会让睡眠合眼，直到你明白了那位人间皇帝所写的内容。\par

天国的皇帝，人及天使的主，为了你生命的裨益赐下了祂的书信；然而，荣耀的儿子，你竟然忽略了热切地阅读这些书信。那么，我恳求你，研究并每天默想你的造物主的话语。在天主的话语中学习天主的心，好使你更热切地向往永恒的事物，使你的灵魂燃起对天国喜乐的更大渴望。因为一个人在此世休息得越多，在他对造物主的爱中就越不安。但为使你这样行，愿全能天主将圣神，那安慰者，倾注于你：愿祂以自己的临在充满你的灵魂，并在充满的同时，安息它。\par

至于我，要知道我在此忍受许多、无数的苦涩。但我感谢全能天主，我所受的远少于我应得的。\par

我向你的荣耀推荐我的儿子、你的病人、纳尔塞斯大人。我确实知道你视他为完全蒙你推荐；但我请求你继续做着你在做的事，这样，通过请求我所见正在完成的，我可以因我的请求而分享你的赏报。此外，我已收到你的阁下的祝福，连同它所带着的爱德。而且我冒昧给你寄去一只鸭子及两只小鸭，作为对你爱德的承认，好使每当你的视线被引领看向它时，在事务的繁忙与喧嚣之中，也唤起你对我的纪念。\par

\SourceRecord{Translated by James Barmby.；Nicene and Post-Nicene Fathers, Second Series；Vol. 12.；Edited by Philip Schaff and Henry Wace.；Buffalo, NY: Christian Literature Publishing Co.,；1895.；<http://www.newadvent.org/fathers/360204031.htm>.}

% structure: p0001 source=h1 target=section reason=page-title

\section{第四卷，第三十二封信}

\Emph{致贵族\Person{纳尔塞斯}}\par

\Person{额我略}致\Person{纳尔塞斯}等。\par

您极甘甜的爱德在信中对我的善行说了许多赞美之词，对此我简答如下：\Emph{\BibleQuote{不要叫我\Person{纳敖米}，即美丽的；要叫我\Person{玛辣}，即苦涩的；因为我充满苦楚}}（\BibleRef{卢 1:20}）。\par

但是关于那些与我的兄弟和同为主教、最可敬的宗主教\Person{若望}悬而未决的司铎的案件，我认为，我们所面对的敌人正是你所断言渴望遵守法典的那个人。此外，我向您的爱德声明，我准备在全能天主的助佑下，尽我一切能力和影响来推进同一案件。并且，如果我看到其中宗座的法典未得到遵守，全能的天主将赐予我所能做的以对抗那些藐视法典之人。\par

至于您的爱德写信给我，要求我为您向我的儿子、首席医生兼前总督\Person{德敖多禄}表示感谢，我已经这样做了，并且从未停止尽可能的称赞您。此外，我恳请您原谅我简短回信；因为我被如此巨大的困苦所压，以至于不允许我读很多或写很多信。我只对您说这些：\Emph{\BibleQuote{因叹息的声音，我忘记了吃我的面包}}（\BibleRef{咏 101:5}）。所有与您在一起的人，我恳请您以我的名义问候。请代我问候\Person{多米尼加}夫人，我没有回复她的信，因为她虽是拉丁人，却用希腊文写信给我。\par

\SourceRecord{Translated by James Barmby.；Nicene and Post-Nicene Fathers, Second Series；Vol. 12.；Edited by Philip Schaff and Henry Wace.；Buffalo, NY: Christian Literature Publishing Co.,；1895.；<http://www.newadvent.org/fathers/360204032.htm>.}

% structure: p0001 source=h1 target=section reason=page-title

\section{第四卷，第三十三封信}

\Emph{致副执事\Person{安特米乌斯}。}\par

\Person{额我略}致\Person{安特米乌斯}等。\par

那些我们的救主恩赐他们从犹太人的灭亡中皈依到他那里的人，我们理当以合理的节制予以帮助；免得（但愿天主阻止这事）他们因缺乏食物而受苦。故我们以此命令的权威嘱咐你，不可疏忽每年向那希伯来人犹斯塔的子女支付金钱；即向\Person{犹利安}、\Person{雷登普图斯}和\Person{福尔图纳}，自即将到来的第十三小纪开始；并要知道，这笔付款务必记入你的账目。\par

\SourceRecord{Translated by James Barmby.；Nicene and Post-Nicene Fathers, Second Series；Vol. 12.；Edited by Philip Schaff and Henry Wace.；Buffalo, NY: Christian Literature Publishing Co.,；1895.；<http://www.newadvent.org/fathers/360204033.htm>.}

% structure: p0001 source=h1 target=section reason=page-title

\section{第四卷·第三十四封书信}

\Emph{致\Person{彭塔莱奥}总督。}\par

\Person{格里高利}致\Place{非洲}总督\Person{彭塔莱奥}。\par

法律如何严厉地惩处异端者那最可憎的败坏，阁下并非不知。因此，若这些人——我们信仰的正统性与法律的严格性均谴责他们——竟得到许可而在您的时代再次抬头，那决非轻罪。如今，据我们所知，在那些地区，多纳徒派的胆量已增长至如此地步：他们不仅以有害的僭越权柄，将天主教信仰的司铎逐出他们的教堂，而且竟还敢对那些已在真实宣认中为再生之水所洁净的人再次施洗。我们非常惊讶——若果真如此——当您被安置在那地区时，竟容许恶徒如此猖獗。首先请思量：若这些在他人时代曾被以正当理由镇压的人，在您的治理下却找到了猖獗的途径，那么您将留给人们什么样的判断呢？其次要知道：我们的天主将向您追讨丧亡者的灵魂，若您疏忽不去纠正——在可能的范围内——如此重大的可憎之事。请阁下勿怪我说这些话。因为我们爱您如同自己的子女，才向您指出我们确信对您有益的事。但请尽快将我们的弟兄与同为主教的\Person{保禄}送来，以免有人以任何借口获得机会阻碍他前来；这样，在更充分查明真相后，我们就能赖天主的助佑，以合理的处理方式解决此事，确定应如何着手惩罚如此严重的罪行。\par

\SourceRecord{Translated by James Barmby.；Nicene and Post-Nicene Fathers, Second Series；Vol. 12.；Edited by Philip Schaff and Henry Wace.；Buffalo, NY: Christian Literature Publishing Co.,；1895.；<http://www.newadvent.org/fathers/360204034.htm>.}

% structure: p0001 source=h1 target=section reason=page-title

\section{第四卷 第三十五封信}

\Emph{致 \Person{维克托} 与 \Person{哥伦布} 两位主教}.\par

\Person{额我略} 致非洲主教 \Person{维克托} 与 \Person{哥伦布}。\par

一种疾病若在初时被忽略，其后如何增强，我们凡在今生有份的，都从自身的急需中证明过了。如果它在初起时就得到高明医者的预见，我们知道它便会因被关注过晚而在造成大害之前消失。故此在灵魂疾病初起时，理智理当促使我们尽一切力量赶紧抵抗它们，免得我们忽略施用有益的良药，以致从我们手中夺去许多我们正力图为我们天主赢得的生命。因此，我们必须以警醒的谨慎看守我们自知被派作牧者的羊栈，使游荡的狼随处都有抵抗它的牧者，而不得其门而入。\par

我们发现，\OriginalTerm{多纳特派}{Donatists} 的毒刺在你们那里如此扰乱了主的羊群，仿佛它没有受到任何牧者的管束。而且有人报告给我们一件我们无法不沉重悲伤地谈起的事，即已有许多人被他们的毒牙撕碎。最后，他们说，为了以最邪恶的狂妄将公教司铎逐出他们的教堂，他们竟以最凶残的邪恶，甚至杀死了许多其他人，并对那些重生之水已赐予救恩的人施以重洗。这一切使我们忧心忡忡，因为，当你们身在那地时，竟容许可咒的狂妄犯下如此恶行。\par

在此事上，我们藉本函劝告你们，在与议并召集会议后，你们应热切并尽全力如此抵抗这仍初萌的疾病，免得它因忽略而获得力量，或将瘟疫的灾祸散布给托付你们照管的羊群。因为，若在任何方面你们（我们相信不会）忽略在初起时抵抗邪恶，他们必以自己错误的刀剑杀伤许多人。而实际上，最重要的事是让我们能够预先解救那些不被缠住的人，不让他们被魔鬼欺诈的圈套所诱捕。此外，预防一人受伤，胜过寻找如何治疗已受伤者。因此，思虑及此，请你们以恒切的祈祷和一切力量，赶快平息这亵渎的邪恶，以使后来的消息，藉基督恩宠的帮助，为我们带来对这些人的惩罚的喜悦多于对他们的过犯的悲哀。\par

此外，请尽可能设法速速将我们的弟兄和同道主教 \Person{保禄} 送至我们这里，以便从他那里更详细地了解如此大罪的原因后，我们能藉造物主的助佑，为这最凶残的邪恶施以合宜责备的良药。\par

\SourceRecord{Translated by James Barmby.；Nicene and Post-Nicene Fathers, Second Series；Vol. 12.；Edited by Philip Schaff and Henry Wace.；Buffalo, NY: Christian Literature Publishing Co.,；1895.；<http://www.newadvent.org/fathers/360204035.htm>.}

% structure: p0001 source=h1 target=section reason=page-title

\section{卷四，第三十六封书信}

\Emph{致利奥主教}\par

\Person{额我略}致\Place{卡塔尼亚}主教\Person{利奥}。\par

我们从许多人的报告中得悉，你们中间长久以来有一种习惯，允许副助祭与他们的妻子同居。天主的仆人，我们教座的执事，在我们前任的权威下，曾禁止任何人再行此事，规定当时已结婚的人应在两件事中选择其一：要么与妻子分居，要么无论如何不得再行施职。据报告，当时的副助祭\Person{斯佩齐奥苏斯}为此自行停职，直至去世，一直担任公证人之职，但停止了副助祭应尽的职务。他去世后，我们得知他的遗孀\Person{奥诺拉塔}因与一男子结合而被您的兄弟团体放逐到隐修院。因此，如果正如所说，她的丈夫自行停职，那么该妇人再婚不应受到损害，尤其如果她与那位副助祭结合时并无意放弃肉身之乐。\par

因此，如果您发现事实正如我们所知，您理应将上述妇人完全从隐修院释放，使她可以毫无顾虑地回到她的丈夫身边。\par

但今后，您的兄弟团体务必格外小心，对于任何可能升任此职的人，要极其谨慎地确保：如果他们已有妻子，不得享有与她们同居的自由；反而必须严格命令他们按照宗座的榜样遵守一切事项。\par

\SourceRecord{Translated by James Barmby.；Nicene and Post-Nicene Fathers, Second Series；Vol. 12.；Edited by Philip Schaff and Henry Wace.；Buffalo, NY: Christian Literature Publishing Co.,；1895.；<http://www.newadvent.org/fathers/360204036.htm>.}

% structure: p0001 source=h1 target=section reason=page-title

\section{第四卷，第三十八封信}

\Emph{致\Person{狄奥德琳达}王后}\par

\Person{额我略}致\Place{隆巴底}人王后\Person{狄奥德琳达}。\par

据某些人报告，我等得知，陛下被若干主教引领，甚至冒犯圣教会，暂停与公教一致的共融。如今我等愈真诚地爱您，就愈为您忧虑，因您相信无知的愚人，他们不但不知自己所谈何事，且几乎不能明白所听之事；他们自己既不阅读，也不相信阅读者，便固守他们自己对我等所妄想的同一错误。因为我等尊敬四次神圣大公会议：尼西亚会议，其中亚略被定罪；君士坦丁堡会议，其中马其顿尼被定罪；第一次厄弗所会议，其中聂斯多略被定罪；以及加采东会议，其中欧提克斯和狄奥斯库若被定罪；宣告凡思想异于这四次大公会议者，即外于真信仰。我等也谴责他们所谴责的，赦免他们所赦免的，并以绝罚击打任何胆敢对同样四次大公会议的信理加以增添或删减者，尤其对加采东会议，关于此会议，某些无知者的心中已生疑虑和迷信之端倪。\par

既然您从我明确的言论和宣认中得知我信仰的纯正，理应对真福伯多禄、宗徒之长之教会不再存有任何疑虑的顾忌；但要持守真信仰，使您的生命建立在教会的磐石上，即建立在真福伯多禄、宗徒之长的宣认上，免得您的一切眼泪和所有善工归于虚无，若它们被发现与真信仰无关。因为正如枝条离了根的活力便枯干，同样，善工无论看来多么美好，若与信仰的坚固分离，便一无是处。\par

因此，陛下理应从速致信我至可敬的兄弟和同僚主教君士坦提乌斯，关于他的信仰和生活，我久已深为确信，并通过致他的信函表明您如何欣喜地接受他的祝圣，且决不与他教会的共融分离；这样我等便可同欢共庆，恰如对待一位善良忠信的女儿。您也应知道，您和您的善工必蒙天主悦纳，倘在祂的审判来临之前，它们已获祂司铎之断案的认可。\par

\SourceRecord{Translated by James Barmby.；Nicene and Post-Nicene Fathers, Second Series；Vol. 12.；Edited by Philip Schaff and Henry Wace.；Buffalo, NY: Christian Literature Publishing Co.,；1895.；<http://www.newadvent.org/fathers/360204038.htm>.}

% structure: p0001 source=h1 target=section reason=page-title

\section{第四卷·第三十九书}

\Emph{致\Person{康斯坦提乌斯}主教}\par

\Person{额我略}致\Place{米兰}（\Emph{Milan}）主教\Person{康斯坦提乌斯}。\par

阅读了尊驾的信函，我们发觉你处于严重的困扰中，主要由于\Place{布雷西亚}（\Emph{Brescia}）的主教和市民，他们要求你寄给他们一封信，在信中你被要求发誓你没有谴责\OriginalTerm{三章}{Three Chapters}。现在，如果你的兄弟之爱的前任\Person{劳伦提乌斯}没有这样做，就不应当要求你这样做。但是，如果他做了，他就不是与普世教会同在，并且违背了他在保证书中所发誓的。然而，既然我们相信他遵守了他的誓言，并继续在\OriginalTerm{天主教会}{Catholic Church}的合一中，那么毫无疑问，他并未向他的任何主教发誓他没有谴责三章。因此，尊驾可以得出结论，你不应当被强迫做你的前任根本没有做过的事。但是，为了避免那些写信给你的人感到不快，请你给他们发一封信，在绝罚的介入下宣告，你既不从加采东大公会议的信仰中删除任何东西，也不接纳那些这样做的人，并且你谴责它所谴责的，宽恕它所宽恕的。我相信这样他们很快就会满意。\par

此外，关于你所写的，许多人因你在弥撒庆典中提及我们的兄弟、\Place{拉文纳}教会同为主教的\Person{若望}而感到不快，你应当查考古代的习惯；如果是习惯，就不应现在被愚人指责。但如果不是习惯，就不应当做可能引起一些人反感的事。然而我已尽力仔细查问，同一位\Person{若望}，我们的兄弟和同为主教，是否在祭台上提到你的名字；他们回答说没有这样做。如果他不提你的名字，我不知道有什么必要迫使你提他的名字。如果这样做可以不引起任何人的反感，你这样做是极值得称赞的，因为你显示出你对弟兄的爱德。\par

此外，关于你所写的你一直不愿意把我的信转交给王后\Person{特奥德林达}，因为信中提到了第五次大公会议，如果你认为她可能因此而不快，你不转交是做对了。因此我们现在按照你的建议行事，即我们只表达对四次大公会议的赞成。然而，关于后来在\Place{君士坦丁堡}举行的、被许多人称为第五次的大公会议，我愿你知道，它没有制定也没有持有任何违背那四次至圣大公会议的东西，因为其中关于信仰什么都没有做，只涉及人；而且是关于那些在加采东大公会议的法案中毫无记载的人，只是在教令颁布后，引起了讨论，最终对人采取了行动。然而我们仍然按你的愿望做了，没有提及这次大公会议。但我们也写信给我们女儿王后，告知你所写关于主教们的事。\Person{乌尔西奇努斯}曾写信给你反对我们的兄弟、同为主教的\Person{若望}，你应当通过写信给他，以温柔和理性制止他的意图。此外，关于\Person{福图纳图斯}，我们希望你兄弟之爱要小心，免得你以任何方式被恶人暗中影响。因为我听说他与你的前任\Person{劳伦提乌斯}一同在教会的餐桌前吃饭多年直至现在，他坐在显贵之中，签署文件，并且在我们兄弟知情的情况下在军队服役。现在，这么多年后，你的兄弟之爱认为他应当被赶出他现在所占的职位。这在我看来完全不合适。因此我通过他给你这个命令，但是私下。尽管如此，如果有什么合理的理由可以指控他，应当提交给我们裁决。但如果全能的天主愿意，我们将通过你的人送信给我们的儿子\Person{迪纳米乌斯}大人。\par

\SourceRecord{Translated by James Barmby.；Nicene and Post-Nicene Fathers, Second Series；Vol. 12.；Edited by Philip Schaff and Henry Wace.；Buffalo, NY: Christian Literature Publishing Co.,；1895.；<http://www.newadvent.org/fathers/360204039.htm>.}

% structure: p0001 source=h1 target=section reason=page-title

\section{第四卷，第四十六封书信}

\Emph{致贵族妇女鲁斯蒂恰纳。}\par

额我略致鲁斯蒂恰纳等。\par

收到尊函，得知您已抵达西乃山，我甚为欣喜。但请相信我，我本也愿与您同行，却绝不情愿与您一同返回。然而，我很难相信您曾到过圣地并见过许多教父。因为我相信，倘若您见过他们，您绝不会如此迅速地返回君士坦丁堡城。而今，既然对这样一座城市的眷爱丝毫未从您心中离去，我怀疑尊驾并未从心底投身于那些用肉眼看见的圣物。愿全能的天主以祂仁慈的恩宠光照您的心灵，赐予您智慧，使您思想一切暂时之事物是何等虚幻，因为正当我们如此交谈时，时光飞驰，审判者临近，而那时刻甚至已近在眼前，那时我们将违背己意交出这世界，而我们本不愿自愿交出。我恳请代我问候阿皮奥大人、尤塞比亚夫人及其女儿们。至于您信中所托付的那位乳母，我极为敬重她，希望她不受任何侵扰。然而，我们正遭受如此巨大的困窘，以至于在目前这个时刻，我们连自己也无法免除征敛（\OriginalTerm{征敛}{angariis}）与负担。\par

\SourceRecord{Translated by James Barmby.；Nicene and Post-Nicene Fathers, Second Series；Vol. 12.；Edited by Philip Schaff and Henry Wace.；Buffalo, NY: Christian Literature Publishing Co.,；1895.；<http://www.newadvent.org/fathers/360204046.htm>.}

% structure: p0001 source=h1 target=section reason=page-title

\section{第四卷，第四十七封信}

\Emph{致执事\Person{萨比尼亚努斯}}。\par

\Person{格里高利}致\Person{萨比尼亚努斯}等。\par

你知道背信者\Person{马克西姆斯}一案中发生了什么事。因为至尊的皇帝陛下曾下令不应祝圣他，之后他却愈发骄傲。原来光荣的贵族\Person{罗曼努斯}属下的人收受了他的贿赂，致使他被祝圣，其方式险些杀死副执事兼产业总管\Person{安多尼努斯}，若非他逃走的话。但在我得知他违背理性和惯例被祝圣后，我曾致信给他，告知他不得擅自主持弥撒大礼，除非我先从我们至尊的陛下们那里了解他们对他有何命令。而这些我的信件，已在城内公开宣读或张贴，他却令人当众撕毁，从而更加公开地跳出来蔑视宗座。你知道我对此会如何容忍，因为我之前就已准备好宁愿死去，也不愿蒙福宗徒\Person{伯多禄}的教会在我的时代堕落。此外，你深知我的作风：我忍耐长久；但一旦我决定不再忍耐，我会欣然面对一切危险。因此，必须依靠天主的帮助迎击危险，以免他被驱使犯下过分的罪。注意我所说的话，并思量是何等巨大的悲痛促成了它。\par

但我听闻他派了一位不知名的神职人员到君士坦丁堡，说\Person{马尔库斯}主教在狱中被处死以换取金钱。关于此事，有一点你可以向陛下们简要提及：如果作为他们仆人的我愿意与伦巴第人的死亡有任何关系，那么今日伦巴第民族就既无王、公爵、也无伯爵，且会陷入极度的混乱。然而，由于我敬畏天主，我不愿与任何人的死亡有牵连。至于\Person{马尔库斯}主教，他既没有被关押，也未遭任何磨难；在他辩护并被判刑的那天，他在我不知情的情况下被公证人\Person{波尼法爵}带到他的家中，那里为他准备了晚餐，他在那里用餐，受到该\Person{波尼法爵}的礼遇，并在夜间突然去世，我想你已知悉。此外，我本打算派我们的\Person{埃克西拉腊图斯}到你那里处理那件事；但因我认为此事已经了结，故停止这样做。\par

\SourceRecord{Translated by James Barmby.；Nicene and Post-Nicene Fathers, Second Series；Vol. 12.；Edited by Philip Schaff and Henry Wace.；Buffalo, NY: Christian Literature Publishing Co.,；1895.；<http://www.newadvent.org/fathers/360204047.htm>.}

% structure: p0001 source=h1 target=section reason=page-title

\section{第五卷，第二封信}

\Emph{致\Person{费利克斯}主教与\Person{西里亚库斯}院长}。\par

\Person{额我略}致\Person{费利克斯}等。\par

呈交给你们的报告的主旨充分说明了虔诚贵妇\Person{特奥多西娅}的投诉；我们在阅读时发现，其中有许多指控，不符合司铎的温和，是针对我们的兄弟和同僚主教\Person{雅努阿里乌斯}的；以至于在她为天主的仆人建立了一座隐修院之后，据说在圣堂奉献礼的那一刻展现了所有涉及贪婪、骚乱和错误的事情。因此，如果事实正如我们在她上述陈述中所发现的，并且如果你们还知道有任何其他不妥之事发生，我们劝告你们，在首先纠正所有错事之后，你们要督促\Place{阿吉利塔努斯}隐修院的院长\Person{穆西库斯}，让他立即特别关注那些他已开始安置在那里的隐修士们；为使这神圣之地在主的助佑下由你们以体面和规范的方式整顿好，既不让我们因前述虔诚贵妇频繁的抱怨——她的善良愿望未得实现——而受到困扰，也不让如此虔诚的计划因你们的疏忽而荒废（我们相信不会如此），以致损害你们的灵魂。\par

\SourceRecord{Translated by James Barmby.；Nicene and Post-Nicene Fathers, Second Series；Vol. 12.；Edited by Philip Schaff and Henry Wace.；Buffalo, NY: Christian Literature Publishing Co.,；1895.；<http://www.newadvent.org/fathers/360205002.htm>.}

% structure: p0001 source=h1 target=section reason=page-title

\section{第五卷·第四封书信}

\Emph{致\Person{康斯坦提乌斯}主教。}\par

\Person{格列高利}致梅迪奥拉努姆（\Place{米兰}）主教\Person{康斯坦提乌斯}。\par

如果允许失职者恢复其品级，教会纪律的力量无疑会被削弱，因为各人在复职的希望中，便不再畏惧放纵其邪恶的倾向。例如，你的弟兄已咨询我们：\Person{阿曼迪努斯}，前司铎兼前院长，因过失被你的前任免职，是否应恢复其品级？此事不允许；我们判定绝不能这么做。然而，倘若他的生活方式堪得，鉴于他已被完全剥夺圣职，可按你的酌定，在隐修院中给他安排一个位于其他隐修士之前的位置。因此，最重要的是，务必小心，勿让任何人的恳求说服你以任何方式恢复失职者的圣秩，免得这种惩罚被认为不是为他们永久规定的，而只是权宜之计。\par

至于前司铎\Person{维塔利安}，关于此人你写道应严加看管，我们将把他送往\Place{西西里}，使他失去一切离开的希望，这样他或许至少能约束自己，作补赎之哀哭。至于\Place{维纳斯港}的\Person{约比努斯}，曾任执事兼院长，我们已判定剥夺其职务，并写信令他人受祝圣接替其位。同样，我们也判定那三位副执事（你的弟兄已通知我们他们失职）应永远停止并剥夺其职务，不得准许他们享有任何超出平信徒共融的待遇。此外，我们已判定前司铎\Person{撒杜尔尼努斯}提供保证，绝不再擅自履行其圣秩的职务。我们愿他保持被剥夺圣秩的状态，留在他原先所在的岛上，准许他对隐修院负有并行使照料与关怀；因为我们相信，他的失足使他更加谨慎，如今他将更仔细地守护那些托付给他的人。\par

此外，关于你教会的公证人\Person{若望}，我们藉以爱你且长久爱你的爱德提醒我们写信，免得你在他过失激怒你时对他下令。因此，提防这一点，务必尽你所能全面调查你教会的财产；这样，你既不会得罪天主，他也无法在人前找到指控你的理由。因为我们写信，并非无理由地为\Person{若望}辩护或夸奖他本人，而是免得你的灵魂在愤怒的刺激下背负任何罪过。所以，如前所述，你必须绝不可忽略在敬畏天主中，对教会财产进行彻底调查。\par

此外，你至亲爱之弟兄的书信使我们对\Person{福图纳图斯}其人感到十分惊讶。但那封信要么不是你口述的，要么，如果它确实是你的，我们在其中绝认不出我们的弟兄，主内\Person{康斯坦提乌斯}。因为你应当已经并仍然应当注意，我们是为你名誉的缘故而写信。因为，当他声称在你们中受冤屈，且未能获得\Emph{监护者（\OriginalTerm{监护者}{defensor}）}的帮助时，他暗示的除了你方的恶意还能是什么？因此，为使此事不致在某种程度上玷污你的名誉，也不致有理由让你的教会遭受任何损害，你应当派遣一位你指定的人来这里，以便审查案情，并结束此事，不带有你的恶意。尤其为此原因：如果在他的申诉之后，在你们中间作出的判决对你有利，他就会被认为不是合理地，而仅仅是通过权力被击败。但我们，出于与你相连的爱德，不停地劝告你做有益于你名誉的事，因为我们知道，尽管这劝告暂时使你忧伤，但当争端的敌意过去后，它将带给你喜乐。在九月，第十三小纪。（梵蒂冈抄本：\Emph{十二月，第十三小纪。}）\par

\SourceRecord{Translated by James Barmby.；Nicene and Post-Nicene Fathers, Second Series；Vol. 12.；Edited by Philip Schaff and Henry Wace.；Buffalo, NY: Christian Literature Publishing Co.,；1895.；<http://www.newadvent.org/fathers/360205004.htm>.}

% structure: p0001 source=h1 target=section reason=page-title

\section{\WorkTitle{第五卷 第五封书信}}

\Emph{致\Person{多米尼库斯}主教}\par

\Person{格列高利}致\Place{迦太基}主教\Person{多米尼库斯}。\par

您的代表（\OriginalTerm{代表}{responsalis}）普洛斯珀，此信的携带者，已与我们在一起，并在表达了您的爱德的其他事情之后，将您的第二封信连同皇帝的谕令的陈述以及一份关于在你们当中举行的教务会议的记录交给了我们。读了这一切，我们为您的牧灵热忱感到喜悦，也为我们的最虔诚的君主们没有听从受贿之人以宗教为借口对您提出的诬告而高兴；但尤其令我们高兴的是，您的弟兄之德已如此致力于保护非洲行省，以致丝毫没有疏忽以司铎的热忱去抑制异端的邪僻派别。关于平息这些派别，我们记得在参考您的爱德的信件之前就已如此充分地定下了法令，以至于我们认为不必再向你们回复任何关于它们的事情。虽然如此，并且我们渴望所有异端者总是被公教司铎以活力和理性所压制，但在彻底审视了你们所做的之后，我们实际上担心，这可能会（愿主避免）引起对其他教务会议的首席主教的冒犯。因为在你们的决议结束时，你们发布了一项判决，其中在命令查究那些异端者时，你们又提出那些疏忽职责者将遭受剥夺财产和地位的处罚。因此，最亲爱的弟兄，在处理涉及我们自身之外需要纠正的事务时，最好首先保持我们彼此间的爱德，并且我们应在思想上服从（我认为这对您的庄重尤为合适）甚至地位低于我们的人。因为那时，当你们按照司铎之职所要求的，努力保持教会内的和谐时，你们将以全部联合的力量更有力地应对异端的谬误。\par

\SourceRecord{Translated by James Barmby.；Nicene and Post-Nicene Fathers, Second Series；Vol. 12.；Edited by Philip Schaff and Henry Wace.；Buffalo, NY: Christian Literature Publishing Co.,；1895.；<http://www.newadvent.org/fathers/360205005.htm>.}

% structure: p0001 source=h1 target=section reason=page-title

\section{第五卷 第八封信}

\\Emph{致\\Person{西彼廉}，执事。}\par

\\Person{格里高利}致执事兼\\Place{西西里}产业总管\\Person{西彼廉}。\par

关于我们产业上的摩尼教徒，我曾屡次敦促阁下以最大的勤奋督促他们，并使他们回归大公信仰。若时机需要，则亲自查询，若事务繁忙，则通过他人进行。此外，我听说在我们产业上有一些希伯来人，他们无论如何都不愿归向天主。但在我看来，你应致函所有已知有希伯来人居住的我们产业，以我之名特别应许他们：无论其中何人归向我们的真主天主耶稣基督，他所持份地的负担将得以减轻。我希望如此办理：若有一人应缴一枚\\OriginalTerm{索利多}{solidus}，则豁免其三分之一；若应缴三或四枚，则豁免一枚\\OriginalTerm{索利多}{solidus}；若更多，则仍按相同比例豁免，或至少按阁下认为合适的方式办理，使归正者负担略减，而教会亦不必承担过重费用。我们如此行并非无益，若借减轻他们缴付的负担而引领他们归向基督的恩宠，因为即便他们自己以微弱的信德而来，但由他们所生之人将来会以更大的信德受洗：如此，我们或得着他们，或得着他们的子孙。而且，为基督的缘故，我们免除他们任何数额的付款都算不上什么大事。再者，前些时候，当执事若望来时，阁下曾写信给我，我当即全部阅读，但隔了许多天才回复；随后，在如此延迟之后，我凭记忆回复了一切细节。但现在我认为有一要点被我遗忘，我怀疑我未作答复。因为你曾写道，正通过某些承办人向农民预支贷款以偿还债务，以免他们因向他人借贷而遭受苛求或物价上涨的负担。此点为我最所乐意；若我确已就此写过，则请遵循我所写的内容。但若如我所疑，我在回复中未就此给予明确指示，则你务须毫不迟疑地为农民的利益预支款项，因为教会的产业不会因此受损，而农民将从中获益。此外，若有其他你认为有利之事，你须毫无犹豫地执行。\par

\SourceRecord{Translated by James Barmby.；Nicene and Post-Nicene Fathers, Second Series；Vol. 12.；Edited by Philip Schaff and Henry Wace.；Buffalo, NY: Christian Literature Publishing Co.,；1895.；<http://www.newadvent.org/fathers/360205008.htm>.}

% structure: p0001 source=h1 target=section reason=page-title

\section{第五卷，第十一封}

\Emph{致\Person{若望}主教。}\par

\Person{额我略}致\Place{拉文纳}主教\Person{若望}。\par

我得知你的弟兄因被理性之责禁止在\OriginalTerm{连祷}{litanies}中佩戴\OriginalTerm{披肩}{pallium}而深感不安。但你通过最杰出的贵族、最显赫的行政长官以及你城中的其他高贵人士，恳切请求获准此事。如今，我们经仔细询问曾担任你弟兄的执事的\Person{阿德奥达图斯}后，查实你的前任们从未在连祷中使用披肩，除非在圣\Person{若翰洗者}、圣宗徒\Person{伯多禄}以及圣殉道者\Person{亚坡里纳里}的庆节上。但我们绝无理由相信他，因为我们许多代表多次到过你弟兄的城，他们声称从未见过此类情形。此事应更相信多数人而非一人，且此人是为自己的教会作证。然而，我们不愿你的弟兄受困扰，也不愿我们的子民的请求在我们面前无效，我们准许在圣\Person{若翰洗者}圣诞日、圣宗徒\Person{伯多禄}以及圣殉道者\Person{亚坡里纳里}的庆节，以及你祝圣典礼的庆典日使用披肩，直至我们获得更确切的认识。但在更衣室内，依照旧例，在迎接并遣散教会之子们后，你的弟兄可佩戴披肩，然后前往举行弥撒礼仪，切不可妄加鲁莽僭越之胆；否则，在外部服饰上越序夺取，可能忽略本该按序行之之事。十月颁赐，第13次小纪。\par

\SourceRecord{Translated by James Barmby.；Nicene and Post-Nicene Fathers, Second Series；Vol. 12.；Edited by Philip Schaff and Henry Wace.；Buffalo, NY: Christian Literature Publishing Co.,；1895.；<http://www.newadvent.org/fathers/360205011.htm>.}

% structure: p0001 source=h1 target=section reason=page-title

\section{第五卷，第十五封信}

\Emph{致若望主教}\par

格列高利致拉文纳主教若望。\par

首先，令我忧伤的是：你的兄弟以双心写信给我，在信中显露一种谄媚，而在世俗交往中却口吐另一种言语。其次，令我痛心的是，我兄弟若望至今仍保留着公证人少年时常有的讥嘲之舌。他出言尖刻，似乎以此为乐。他对朋友当面奉承，背后诋毁。第三，令我悲伤且深恶痛绝的是，他竟将可耻的罪行加于仆人之身，无论何时，称他们为\Quote{柔弱}；更可悲的是，此等事竟公开行之。此外，既无纪律约束司铎的生活，他只显露出作为他们的主人。最后一事，在傲慢的迹象中最为首要，即他在教堂外使用\OriginalTerm{披带}{pallium}——此事在我前任之时他从未敢行，且据我代表所证，他任何前任也未曾敢行（或许在安放圣髑时除外，但关于圣髑，仅有一人声称如此）；然而，在我任内，他竟对我蔑视，以极度狂妄不仅行了此事，甚至习以为常。\par

从这一切我看出，对于他，\OriginalTerm{主教职}{Episcopacy}的尊荣全在于外表，而非内心。我实在感谢全能的天主，当此事（我前任们从未听闻）传到我耳中时，\Person{伦巴第人}正驻扎在我与拉文纳城之间。或许我心中曾想向人显示我能何等严厉。\par

然而，唯恐你以为我欲贬抑或减损你教会的尊荣，请回想拉文纳的执事在罗马庄严举行弥撒时曾站立何处，再查问他如今站立何处；你就会明白我实愿尊荣拉文纳教会。但任何人出于骄傲而攫取任何事物，这是我不能容忍的。\par

尽管如此，我已就此事致信我在君士坦丁堡的执事，要他向所有统辖三十甚至四十位主教的人查询。若何处有此习俗，即游行\OriginalTerm{祷文}{litanies}时佩戴\OriginalTerm{披带}{pallium}，但愿在我这里，拉文纳教会的尊荣丝毫不被减损。\par

因此，最亲爱的兄弟，请反思我上述一切：思念你蒙召之日；思考你将为何背负主教职之重担而交账。改正那些公证人的品性。看看一个主教在言语和行为上应当如何。对你的弟兄完全真诚。不要口是心非。不要想要显得超过你实际所是，这样你才能超过你所显出的。请相信我，当我到达现职时，我对你怀有如此的敬重与爱德，以致你若曾愿持守我这爱德，你必永不会找到像我这样的兄弟，如此真心爱你，在一切虔敬中与你同心；但当我得知你的言语和品性时，我承认我退缩了。我因此恳请你，因全能的天主，改正我所提及的一切，尤其是诡诈之恶。请允许我爱你；无论今生还是来世，被你弟兄所爱或于你有益。但请以行为而非言语回应这一切。\par

\SourceRecord{Translated by James Barmby.；Nicene and Post-Nicene Fathers, Second Series；Vol. 12.；Edited by Philip Schaff and Henry Wace.；Buffalo, NY: Christian Literature Publishing Co.,；1895.；<http://www.newadvent.org/fathers/360205015.htm>.}

% structure: p0001 source=h1 target=section reason=page-title

\section{第五卷 · 第十七封信}

\Emph{致执事西彼廉}。\par

额我略致西彼廉等。\par

我收到了你关于主玛克西米安于十一月去世的极为苦涩的书信。他确实已获得他所渴望的赏报，但不幸的叙拉古城人民值得同情，因为他们不配长久拥有这样一位牧者。因此，务必请你的爱德审慎留意，务要在这同一教会中选立一位接替者，使他在主玛克西米安之后，不致不配地获得同样的治理职位。我确实相信大多数人会推选司铎特拉扬，据说他品性善良，但我怀疑他不适合治理那个地方。然而，如果找不到更好的人选，且对他没有指控，在万分必要的情况下，可以俯就他。但若在这件事上征询我的意愿，我私下告知你我的愿望：因为在这同一教会中，我看没有人比卡塔尼亚教会的总执事若望更适合在主玛克西米安之后继任。如果他的选举能够实现，我相信他将是一个非常合适的人选。但你也必须先私下查问他，看他是否有任何可能妨碍的指控。如果他发现毫无过失，就可以被正当地选立。如果这样做，我们的兄弟及同僚主教莱奥也必须允许他离开，以便他可以自由地被祝圣。因此，我已注意将这些事告知你的爱德；现在你应力求周详地观察各方，并安排那中悦天主的事。\par

\SourceRecord{Translated by James Barmby.；Nicene and Post-Nicene Fathers, Second Series；Vol. 12.；Edited by Philip Schaff and Henry Wace.；Buffalo, NY: Christian Literature Publishing Co.,；1895.；<http://www.newadvent.org/fathers/360205017.htm>.}

% structure: p0001 source=h1 target=section reason=page-title

\section{第五卷，第十八封书信}

\Emph{致若望主教}\par

额我略致君士坦丁堡主教若望。\par

当贤弟被擢升到司祭尊位时，你记得你发现教会的和平与和谐是怎样的。但是，我不知是出于何种大胆或何种骄矜膨胀，你竟试图夺取一个新名号，致使所有弟兄的心可能因而感到冒犯。对此我极为惊异，因为我记得你是多么宁愿逃避主教职务，而不是得到它。然而，如今你得到了它，却如此行使职权，仿佛你曾以野心奔向前去。因为，你曾承认自己不配被称为主教，而今却已堕落到如此地步，竟轻视众兄弟，贪求被称为唯一的主教。\newline{}关于此事，我可敬的前任培拉奇曾向圣座致送重要的信函；他在信中废除了你们为处理我们昔日的兄弟同主教额我略的案件而召集的会议法案，因为那可恶的骄矜头衔，并且禁止他按惯例派往我们主上的朝门的总执事与你共祭。但他去世后，我这不堪者继承了教会的治理，我曾通过其他代表，也通过我们共同的孩子执事撒比尼雅诺，用心地口头而非书面告诉贤弟，希望你克制自己，不要如此僭越。并且，若你拒绝改正，我禁止他与你共祭；这样我先凭一种羞耻感恳请圣座，以便如果这可恶亵渎的主张不能通过羞耻改正，可随后采取严格的教会法规措施。\newline{}而且，既然要切除的疮应先以轻柔的手抚摸，我恳求你，我哀求求你，并用我能力范围内的一切温柔要求你，使贤弟驳斥所有阿谀奉承并献给你这错误名号的人，也不要愚昧地同意被以这骄傲头衔称呼。因为我实在含泪而言，且出自内心深处将这归咎于我的罪过：这位被置于主教等级中、目的正是为了引导他人灵魂归于谦逊的我的兄弟，至今竟不能被引导归于谦逊；这位向他人教授真理者，竟不肯教导自己，即使我恳求他。\par

我请你考虑，在这莽撞的僭越中，整个教会的和平受到扰乱，并且与倾注于众人公有的恩宠相矛盾；在这恩宠中，你自己无疑将有能力成长，只要你下定决心这样做。而你越克制自己不篡夺那个骄傲愚昧的头衔，你就变得越大；你越不因贬抑兄弟而专心于篡夺，你就越有进步。\newline{}因此，最亲爱的兄弟，请你全心爱慕谦逊，借着谦逊，所有弟兄的和睦与神圣普世教会的合一得以保持。\newline{}诚然，宗徒保禄听见有人说：\BibleQuote{我属保禄，我属阿颇罗，但我属基督}\BibleRef{格前 1:13}，极其恐惧地看待主身体的这种撕裂，他们仿佛将自己连接到其他头上，于是喊道：\BibleQuote{保禄为你们被钉十字架了吗？或者你们受洗是归于保禄名下吗？}\BibleRef{格前 1:13}。如果他尚且避免将基督的肢体部分地置于某些头下，仿佛在基督之外，即使这些头是宗徒本人，那么，在最后的审判时，你将如何面对作为普世教会之头的基督，因为你企图借着\OriginalTerm{普世}{Universal}的称号，将他所有的肢体都置于你自己之下？\newline{}我问，在这错误的名号中，谁被提出作为效法的榜样，岂不是那个蔑视与他一同被造的天使军团，试图跃升到独一无二的巅峰，为显得不受任何管辖、独自凌驾于众人之上的那位？他甚至说：\BibleQuote{我要升到天上，我要高举我的宝座高于天上的星辰：我要坐在会众的山上，在极北之处：我要升到云的高处；我要与至上者相似。}\BibleRef{依 14:13}\par

因为你的众弟兄，普世教会的主教们，若不是天上的星辰，他们的生活和宣讲在人的罪恶和错误中，如同在黑夜的阴影中一同闪耀，又是什么呢？当你渴望借这骄傲的称号将自己置于他们之上，并践踏他们的名号与你相比时，你所说的，岂不就是\Quote{我要升到天上；我要把我的宝座高举在天空的星辰之上}吗？所有主教不都是云彩，既以讲道的话语降雨，又以善行的光芒闪耀吗？当你的兄弟之爱轻视他们，并试图将他们压在自己之下时，你所说的，岂不就是古仇所说的\Quote{我要升到云层高处}吗？当我含泪看到这一切，并对天主隐藏的审判战栗时，我的恐惧增加，我的心无法容纳叹息，因为这位极其圣洁的人，\Person{若望}主，有如此伟大的节制和谦卑，竟因熟悉之舌的诱惑，陷入如此极端的骄傲，以致在贪求那不当的名称时，试图像那一位，他因骄傲而想如同天主，却失去了赐予他的肖似的恩宠，并且因为他寻求虚假的荣耀，从而丧失了真正的福乐。诚然，\Person{伯多禄}，宗徒之长，他本人是圣而普世教会的一员，\Person{保禄}、\Person{安德肋}、\Person{若望}——他们若不是个别团体的首领，又是什么呢？然而，所有人都在一个元首之下成为肢体。并且（用简短的话语概括一切），法律以前的圣人，法律下的圣人，恩宠下的圣人，所有这些组成主身体的，都被设立为教会的肢体，他们中没有一个人愿意被称为普世的。现在，愿你的圣德承认，你在内心膨胀到何等程度，竟渴望被称为那个名称，而任何真正圣洁的人都不敢自称此名。\par

正如你的兄弟之爱所知，岂不是曾有这样的情况：这宗座——我因天主的眷顾而服务——的首长们，被可敬的\Place{加采东}大公会议赐予被称为\Quote{普世的}的荣誉吗？然而，他们中无一人曾愿意被称为这样的称号，或抓住这个不明智的名称，以免因教宗职的尊位而将特殊的荣耀归于自己，从而似乎否定了所有弟兄。\par

但我知道，这一切都源于那些以欺骗性的亲密服侍你的圣德的人；我恳求你的兄弟之爱明智地提防他们，不要让自己被他们的话语欺骗。因为他们越是赞美你，就越应被视为更大的敌人。离弃这样的人吧；如果他们必须欺骗，至少让他们欺骗世俗之人的心，而不是司铎的心。\BibleQuote{让死人埋葬他们的死人}\BibleRef{路 9:60}。但要与先知同说：\BibleQuote{愿那些向我说\InnerQuote{啊哈，啊哈}的人，转身受辱}\BibleRef{咏 69:4}。又说：\BibleQuote{但不要让罪人的油傅抹我的头}\BibleRef{咏 140:5}。\par

因此，智慧人也善加劝诫：\BibleQuote{与多人和平相处，但只从千人中选一人作你的谋士}\BibleRef{德 6:6}。因为\BibleQuote{恶交败坏善俗}\BibleRef{格前 15:33}。古仇无法侵入坚强的心时，就寻找与他们相关的软弱之人，并如同以梯子靠着高墙攀登。这样，他藉与\Person{亚当}相关的女人欺骗了他。同样，当他杀害有福的\Person{约伯}的儿女时，他留下了软弱的女人，以便他自己无法渗入他的心，却至少能通过女人的话做到。因此，无论什么样的软弱而世俗的人在你身边，愿他们被自己那劝诱的话语和谄媚所粉碎，因为他们因那假装爱慕者的乖僻，为自己招致了天主的永恒敌意。\par

确实，从前藉\Person{若望宗徒}宣告说：\BibleQuote{孩子们，现在是末时了}\BibleRef{若一 2:18}，正如真理所预言的。如今，瘟疫和刀剑肆虐世界，民族攻击民族，地球的球体震动，裂开的大地与居民一同消溶。因为一切预言的事都已应验。骄傲之王临近了，（可怖的是！）有一队司祭正在为他预备，因为那些被设立为谦卑领袖的人，却将自己投身于骄傲的轭下。但在这事上，即使我们的口舌完全不抗议，那以自己的位格特别反对骄傲恶习者的权能，已举起为报复高抬自己的人。因此经上记载：\BibleQuote{天主阻挡骄傲的人，赐恩宠给谦卑的人}\BibleRef{雅 4:6}。又说：\BibleQuote{凡心高气傲的，在天主面前为不洁}\BibleRef{箴 16:5}。因此，论到骄傲的人经上写着：\BibleQuote{尘土和灰烬为什么骄傲？}\BibleRef{德 10:9}。因此，真理亲自说：\BibleQuote{凡自高自大的，必被贬抑}\BibleRef{路 14:11}。并且，为藉谦卑将我们引回生命之路，祂屈尊在自己身上展示祂所教导我们的，说：\BibleQuote{跟我学习，因为我是良善心谦的}\BibleRef{玛 11:29}。为此，天主的独生子取了我们的软弱形像；为此，那不可见的不仅显现为可见的，甚至被轻看；为此，祂忍受了侮辱的嘲笑，讥讽的辱骂，痛苦的折磨，好使天主在祂的谦卑中教导人不要骄傲。那么，谦卑的德行何其伟大！为了教导这德行，那位无比伟大的竟变得微小，甚至至死受苦！因为魔鬼的骄傲是我们丧亡的根源，而天主的谦卑却被发现是我们救赎的方法。也就是说，我们的仇敌，受造于万有之中，却渴望显得超乎万有之上；而我们的救赎主，仍然超乎万有之上，却屈尊在万有中成为微小。\par

我们这些主教，从我们救赎主的谦卑中领受了尊位，却效仿那仇敌本身的骄傲，对此我们又能说什么呢？看，我们知道我们的造物主从祂至高之巅降下，为的是赐荣耀给人类；而我们这些受造于尘土中最卑微者，却以贬抑我们的弟兄为荣。天主甚至谦卑自下到我们的尘土；而尘土之人却昂首向天，用舌头在地上横行，毫不羞愧，也不畏惧被高举；人本是腐朽，人子本是虫豸。\par

最亲爱的弟兄，让我们回想最智慧的撒罗满所说的这话：\BibleQuote{闪电先于雷鸣，心骄必先于衰败}\BibleRef{德 32:10}；另一方面又续言：\Quote{谦卑先于光荣。}因此，我们若力求达到真正的崇高，就当心谦。绝不要让我们的心眼被得意之烟所蒙蔽——这烟越是升起，便越快消散。让我们思考我们救赎主的诫命是如何劝诫我们的，祂说：\BibleQuote{神贫的人是有福的，因为天国是他们的}\BibleRef{玛 5:3}。因此，祂也藉先知说：\BibleQuote{但我要垂顾谁呢？惟有那谦卑、安静、敬畏我言语的人}\BibleRef{依 66:2}。的确，当主想要引导祂那仍被软弱困扰的门徒之心回归谦卑之路时，祂说：\BibleQuote{谁若愿意在你们中为首，就当作你们的奴仆}\BibleRef{玛 20:27}。由此清楚可见，谁在心中谦卑，谁才真正被高举。因此，我们当害怕被列于那些寻求会堂中首位、市场上问安、并被人称为辣彼的人中。因为相反，主对门徒说：\BibleQuote{但你们不要被称为\InnerQuote{辣彼}，因为你们的导师只有一位，你们众人都是兄弟。也不要称地上的人为父，因为你们的父只有一位}\BibleRef{玛 23:7}\par

那么，最亲爱的弟兄，在将来审判那可怕的审查中，你若渴望在世上不仅被称为父，甚至被称为总父，你又将如何作答呢？因此，务必提防恶人的恶念，远离一切引人跌倒的煽动。\BibleQuote{（固然）免不了有绊脚石；但绊倒人的那个人有祸了！}\BibleRef{玛 18:7}。看，因这可憎的傲慢头衔，教会撕裂，众弟兄的心都被触怒。难道你忘了真理所说的话？\BibleQuote{无论谁使这些信我的小子中的一个跌倒，倒不如拿一块驴拉的磨石拴在他的颈上，沉在深海里。}\BibleRef{玛 18:6}（但经上记载：\BibleQuote{爱不求自己的益处}\BibleRef{格前 13:4}。看，你的兄弟之情甚至把不属于自己的东西据为己有。又记载：\BibleQuote{论尊敬，彼此争先}\BibleRef{罗 12:10}。而你却企图夺去所有人的尊荣，单为自己非法篡夺。最亲爱的弟兄，经上记载的\BibleQuote{与众人和平相处，追求圣德；若无圣德，谁也见不着主}\BibleRef{希 12:14}又在何处？经上记载的\BibleQuote{缔造和平的人是有福的，因为他们要称为天主的子女}\BibleRef{玛 5:9}又在何处？\par

你应当审度，免得有苦根的苗芽生长起来搅扰你，因而使许多人玷污。但即使我们疏于审度，上天的审判也会警惕如此巨大的骄傲膨胀。而我们——面对这邪恶企图所犯下的如此重大的罪过——我们，我说，在观察真理所吩咐的：\BibleQuote{如果你的弟兄得罪了你，去，要在你和他两人之间规劝他。如果他听从了你，你便赢得了你的兄弟。但他如果不听从，你就另带一个人或两个人去，为叫任何事实，凭两三个证人的口供得以确定。若是他仍不听从他们，你要告诉教会；如果他连教会也不听从，你就把他看作外邦人或税吏}\BibleRef{玛 18:15}。因此，我已一次又一次地通过我的代表，以谦卑的话语小心地斥责这侵犯整个教会的罪过；如今我亲自写信。我在谦卑方面所应尽的一切义务，我都没有忽略。但是，倘若在我斥责时被轻视，那么我就必须诉诸教会了。\par

因此，愿全能的天主向你的兄弟之情显示，我对你怀有多么大的爱德，促使我说这样的话；并且我在此事上，不是反对你，而是为你忧伤。但情况如此，我在此必须将福音的训导、教会的法令和弟兄们的益处置于我所深爱之人的情面之上。\par

关于若望与阿塔纳削两位司铎的案件，我收到了你圣座最甘饴可亲的信函。赖主助佑，我将在另一封信中答复你。如今，我被蛮族的刀剑四面围困，承受着如此巨大的苦难，以致不容我——不必说谈论许多事——甚至连呼吸都几乎不能。写于一月朔日；第十三小纪。\par

\SourceRecord{Translated by James Barmby.；Nicene and Post-Nicene Fathers, Second Series；Vol. 12.；Edited by Philip Schaff and Henry Wace.；Buffalo, NY: Christian Literature Publishing Co.,；1895.；<http://www.newadvent.org/fathers/360205018.htm>.}

% structure: p0001 source=h1 target=section reason=page-title

\section{第五卷·第十九封}

\Emph{致执事萨比尼亚努斯}。\par

格里高利致萨比尼亚努斯等。\par

关于我们最可敬的弟兄——君士坦丁堡主教若望的案件，我不愿写两封信。但我已简要起草了一封，似乎能兼顾两方面的要求：即既诚实又友善。\par

因此，请你的爱德务必将我现在遵皇帝意愿寄给他的这封信交给他。因为随后将另有一封信寄给他，那封信不会令他的傲慢欢喜。他竟至于此：以司铎若望的案件为借口，将案卷转寄至此，其中几乎每一行都称自己为\OriginalTerm{普世}{\GreekText{οἰκουμενικὸν}}牧首。但我寄望于全能天主，至高威严将挫败他的伪善。我却奇怪他怎能如此欺骗你的爱德，以致你竟让皇帝陛下被说服，亲自就此事写信给我，告诫我与他保持和平。因为，若皇帝陛下愿秉持正义，他本应告诫他戒除那傲慢的称号，然后我们之间便立即有了和平。然而，我担心你们并未完全考虑，我们前述的弟兄若望此举是何等狡猾。因为他的目的正在于此：一方面使皇帝陛下得到服从，从而他自己似乎在其虚妄中得到确认；另一方面若我不服从，那么皇帝的心思便可能被我激怒。但我们必坚守正道，在此事中除了全能之主外无所畏惧。因此，你的爱德在任何事上都不要害怕。你看到世上一切反对真理的高傲之事，为真理的缘故都要轻看。要信赖全能天主的恩宠和蒙福的宗徒伯多禄的助佑。要记住真理的声音说：\BibleQuote{那在你们内的，比那在世界上的更大}（\BibleRef{若一 4:4}）。在此事中，无论必须做什么，都要以最大的权威去做。因为如今我们既已无法以任何方式免受敌人刀剑的伤害，为了对共和国的爱，我们已失去了金银、奴隶和衣物，那么若再因那些人而失去信德，就太耻辱了。因为赞同那可怕的称号，无异于失去信德。因此，正如我已在先前信中所写的，你切不可擅自与他继续交往。\par

\SourceRecord{Translated by James Barmby.；Nicene and Post-Nicene Fathers, Second Series；Vol. 12.；Edited by Philip Schaff and Henry Wace.；Buffalo, NY: Christian Literature Publishing Co.,；1895.；<http://www.newadvent.org/fathers/360205019.htm>.}

% structure: p0001 source=h1 target=section reason=page-title

\section{第五卷，第二十书}

\Emph{致奥古斯都·莫里斯}\par

额我略致莫里斯等。\par

我们最虔敬、由天主所指定的主上，在他其他庄严的关怀与负担中，也以属灵的热忱正直地守护着司铎之间的和平，因为他虔诚而真实地认为，没有人能妥善治理地上之事，除非他懂得如何处理神圣之事；并且共和国的和平系于普世教会的和平。因为，最宁静的主上，如果司铎们同心合意，致力于以口舌为你们向救赎主祈祷，并且按他们应当做的，也以他们的功德祈祷，那么有什么人的权势、什么血肉之臂的力量敢对你们最基督化的帝国的崇高地位举起不敬之手呢？或者，有什么最野蛮民族的刀剑会如此残酷地前进屠杀信友，除非我们这些被称为司铎却并非司铎之人的生活被最邪恶的行为所压垮？但是，当我们忽略与我们相关的事，而思想那些与我们无关的事时，我们将自己的罪恶与蛮族的力量相结合，而我们的过错压垮了共和国的力量，磨利了敌人的刀剑。然而，我们要为自己说什么呢？我们这些不配治理天主之民的人，用我们罪恶的重担压迫他们；我们以榜样摧毁我们以口舌所宣讲的；我们以行为教导不义，却仅以声音陈述正义之事。我们的骨头因斋戒而磨损，内心却膨胀；身体披着粗劣的衣服，心中却高抬自己超过紫袍；我们卧于灰烬，却藐视崇高；我们身为谦逊之师，却是骄傲之首；在羊的面孔背后隐藏着狼的牙齿。这一切的结局，除了我们欺骗人，却在天主面前显明之外，还能是什么呢？因此，最虔敬的主上极其明智地，为遏制战争行动而寻求教会的和平，并为了巩固和平，屈尊使司铎们的心回归和谐。这确实是我所愿望的；就我而言，我顺从于他最宁静的命令。但是，由于这并非我的事，而是天主的事；由于虔诚的法律，由于可敬的会议，由于我们主耶稣基督的命令本身，被某种傲慢而浮夸的措辞的发明所扰乱，愿最虔敬的主上切开伤处，并以至尊权威的锁链捆绑这抗拒的病人。因为，紧紧包扎这些事，你们就解救了共和国；而剪除这些事，你们就为你们统治的延续作了预备。\par

因为凡认识福音的人都清楚，藉主的声音，整个教会的关怀被托付给了圣宗徒、所有宗徒之长伯多禄。因为经上对他说：\BibleQuote{伯多禄，你爱我吗？你喂养我的羊群}\BibleRef{若 21:17}。经上又对他说：\BibleQuote{看，撒殚想要筛你们像筛麦子一样；但是，我已为你祈求，伯多禄，使你的信德不致失落。当你回头以后，要坚固你的弟兄}\BibleRef{路 22:31}。经上又对他说：\BibleQuote{你是伯多禄，在这磐石上我要建立我的教会，阴间的门不能胜过她。我要把天国的钥匙交给你；凡你在地上所束缚的，在天上也要被束缚；凡你在地上所释放的，在天上也要被释放}\BibleRef{玛 16:18}。\par

看，他接受了天国的钥匙，赐给了他捆绑和释放的权柄，整个教会的关怀和元首之位交给了他，然而他并没有被称为普世宗徒；而我最神圣的同伴，司铎若望，却试图被称为普世主教。我不得不呼喊说：\Emph{O tempora, O mores!}\par

看，欧洲各地区的一切都落入了蛮族的势力，城市被摧毁，营垒被推翻，省份人烟稀少，没有农夫居住土地，偶像崇拜者每日肆虐并统治以屠杀信徒，然而司铎们——本应痛哭流涕卧于尘土和灰烬中——却为自己寻求虚名的称号，并以新奇而亵渎的头衔自夸。\par

最虔敬的主上，我在这件事上是在维护自己的事业吗？我是在怨恨自己所受的特殊委屈吗？不，这是全能天主的事业，是普世教会的事业。\par

这是谁，竟敢违抗福音的规诫，违反教规的法令，擅自为自己篡夺一个新名？但愿他独自一人，若能在不损及他人的情况下——他渴望成为普世者。\par

而且我们确实知道，君士坦丁堡教会的许多司铎曾陷入异端的漩涡，不仅成为异端者，甚至成为异端首领。因为从这里出了聂斯多略，他以为耶稣基督——天主与人的中保——是两个位格，因他不相信天主能成为人，竟至于跌入犹太式的背信。从这里出了马其顿尼，他否认天主圣神与圣父、圣子同一本体。那么，如果在那个教会中有人取用了那个名称，使自己成为一切善人的首领，那么随之而来的便是普世教会从它的地位上坠落（愿天主阻止），因为那被称为普世者跌倒了。但是，愿那亵渎之名远离基督徒的心，它夺去了所有司铎的尊荣，而由一人疯狂地妄自夺取。\par

的确，为尊敬宗徒之长伯多禄，可敬的加采东大公会议曾将此名献给罗马教宗。但是，他们之中从未有人同意使用这个独特的名号，以免因将某物单独给予一人而剥夺司铎们普遍应得的尊荣。那么，为何我们不愿寻求这头衔的荣耀（即使被献上时），而另一人却未受献上而擅自夺取呢？\par

因此，应当由我们最虔敬的诸位主上的命令来折服那不屑于服从教规训令的人；应当约束那损害神圣普世教会、心高气傲、贪恋以独特之名自夸、并且通过一个专属自己的头衔将自己置于你们帝国尊严之上的人。\par

看，我们都因这事而受冒犯。那么，让冒犯的制造者回归正途；所有司铎间的纷争就会止息。就我而言，我是所有司铎的仆人，只要他们生活得当，合乎司铎身份。因为凡因虚浮荣耀的膨胀，昂首对抗全能天主和教父规条的人，我信赖全能天主，他必不会使我向他低头，即使用刀剑也不行。\par

此外，关于我们在此城听闻此头衔时所发生的一切，我已详尽告知我的执事兼\OriginalTerm{代表}{responsalis}\Person{Sabinianus}。愿我的主上们的慈爱顾念我，视我为你们的人——你们向来爱护、抬举我，超过他人；而我切愿服从你们，却又惧怕在至高可畏的审判中获罪。因此，依照上述执事\Person{Sabinianus}的请求，愿我至为慈仁的主上或者俯允审理此事，或者促使那屡次提及的人最终放弃这一企图。若因陛下最公正的判决，他服从你们的命令，即便是温和的命令，我们将感谢全能天主，并为因你们而赐予整个教会的平安而喜乐。但若他仍坚持当前的争执，我们确信真理的这句话已经应验：\Quote{\BibleQuote{凡高举自己的，必被贬抑}}（\BibleRef{路 14:11}）。经上又记载：\Quote{\BibleQuote{心骄必败}}（\BibleRef{箴 16:18}）。我，顺从主上们的命令，既已温和地写信给我上述的同侪司铎，并谦卑地劝告他改除贪图虚浮荣耀的毛病。因此，若他愿意听从我，他便有一位忠心的弟兄。但若他坚持骄傲，我已看到后果：他将发现那一位与他为敌，经上记载：\Quote{\BibleQuote{天主拒绝骄傲的人，赐恩宠于谦逊的人}}（\BibleRef{雅 4:6}）。\par

\SourceRecord{Translated by James Barmby.；Nicene and Post-Nicene Fathers, Second Series；Vol. 12.；Edited by Philip Schaff and Henry Wace.；Buffalo, NY: Christian Literature Publishing Co.,；1895.；<http://www.newadvent.org/fathers/360205020.htm>.}

% structure: p0001 source=h1 target=section reason=page-title

\section{第五卷，第二十一封信}

\Emph{致\Person{君士坦丁娜}\OriginalTerm{奥古斯塔}{Augusta}}。\par

\Person{额我略}致\Person{君士坦丁娜}等。\par

全能的天主将您虔敬之心握在祂右手中，祂既藉着您保护我们，也为您暂时的作为预备永恒的赏报。因为我从我的代理人六品副祭\Person{萨比尼亚努斯}的信中得知，陛下以何等公义关心着蒙福的宗徒之长\Person{伯多禄}的事业，反对某些骄傲谦卑、假意仁慈的人。我信赖我们救赎主的恩赐，为这些您向最安详的主上及其最虔诚的儿子们所行的善事，您也必将在天国得到报偿。毫无疑问，您将获得永恒的恩惠，从罪恶的锁链中解脱，因为若您在他的教会的事业上使他成为您的债务人，那束缚与释放的权柄已交给了他。因此我仍恳求您，不容许任何人的伪善胜过真理，因为有的人如那卓越的宣讲者所言，以甜言蜜语欺骗无知者的心——他们衣着卑贱，内心却骄傲自大。他们假装鄙视世上的一切，却企图为自己谋取世上所有的一切。他们在众人前自认不配，却不能满足于私下的名号，因为他们贪图那些能使自己显得比众人更配的事物。因此，愿您的虔敬，全能的天主已任命您与我等最安详的主上治理普世，藉着您对正义的偏爱，事奉那从祂领受如此伟大权柄的祂，好使您能更稳妥地治理托付给您的世界，因为您在实行真理时更真实地事奉了万有的创造者。\par

此外，我告诉您，我收到了最虔诚的主上的来信，要我同我的兄弟及同司铎\Person{若望}和好。的确，宗教主上如此命令司铎是合宜的。但是，当这位我的兄弟以新的僭越和骄傲自称普世主教，并在我们可敬的前任教宗时代，使自己在主教会议中被称为如此骄傲的称号时——尽管那次会议的所有法令因被宗座否定而废除——最安详的主上给了我一个颇令人苦恼的暗示，即他没有斥责那行事骄傲的人，却试图说服我放弃我的主张，而我在这一事业中是在维护福音、教规、谦逊和正直的真理；反观我前述的兄弟及同司铎，却在违背福音原则、蒙福的宗徒\Person{伯多禄}、所有教会以及教规的规定。但主是全能者，万物都在祂手中；经上记载：\BibleQuote{任何智慧、明哲或谋略，都不能对抗主}\BibleRef{箴 21:30}。确实，我多次提及的最圣洁的兄弟试图说服我最安详的主上许多事；但我深知，他所有的祈祷和眼泪都不会允许我的主上在任何事上被任何人以违反理性或损害其灵魂的方式所欺骗。\par

然而，令人非常苦恼且难以耐心忍受的是，我前述的兄弟及同主教竟藐视所有其他人，企图被称为唯一的（主教）。但在他的这种骄傲中，除了说明假基督的时代已经临近之外，还能表示什么呢？因为他确实是在效法那一位，他鄙视与天使军团共享社交之乐，企图跃上独一无二的高峰，说：\BibleQuote{我要将我的宝座高举在星宿之上；我要坐在聚会之山上，在北方的极处；我要升到高云之上，我要与至高者相似}\BibleRef{依 14:13}。因此，我恳求您因全能的天主，不容许您虔敬的时代被一个人的自高所玷污，也绝不对此等乖谬的称号给予任何赞同；且在此事上，愿您的虔敬绝不轻视我；因为尽管\Person{额我略}的罪过如此重大，理应遭受这些事，但宗徒\Person{伯多禄}却没有任何罪过，致使他在你们的时代应当如此受苦。因此，我一再因全能的天主恳求您，正如你们的先祖君主寻求圣宗徒\Person{伯多禄}的恩宠，你们也应同样既为自己寻求，又加以保持；并且，他在你们中的尊荣绝不因我们这些不配事奉他的人的罪过而有所减损，因为他既能现在在一切事上作你们的助佑，将来也能赦免你们的罪过。\par

此外，至今已整整七年，我们在这城中生活在伦巴第人的刀剑之下。这座教会每日为他们在他们中间生存而花费多少，不可胜数。但我简要指出，正如在拉文纳地区，我的主上们的虔敬为意大利的第一支军队设立了一位\OriginalTerm{司库}{sacellarium}，以应付每日的经常性开支，同样，我在这城中也是他们为此等目的而设的司库。然而，这座教会，一面同时不断地为圣职人员、隐修院、穷人、民众，此外还为伦巴第人花费如此之多，看哪，它仍被所有教会的苦难所压，这些教会因这一个人的骄傲而大大叹息，尽管她们不敢说什么。\par

此外，在未通知我和我的\OriginalTerm{代理}{responsalis}的情况下，萨洛纳城的一位主教被祝圣；此事在先前的任何君主治下都从未发生过。我听闻后，立即传信给那违规的非法祝圣者，命他绝不可擅自举行弥撒的隆重礼仪，除非我们先从我们最安详的君主们那里得知他们已下令如此行事；并且我以绝罚之令禁止他。但他却藐视我、轻看我，倚仗某些世俗人物的胆大妄为（据说他给了他们许多贿赂，以致穷尽其教会），迄今竟擅自举行弥撒，而且拒绝按我主上们的命令前来见我。如今，我顺从他们虔诚的谕令，从心底已宽恕了这位未经我同意而被祝圣的马克西穆斯，他在祝圣中越过我和我的\OriginalTerm{代理}{responsalis}的僭越，如同他是在我的权下被祝圣一样。但他其他的过犯——即我所听闻的肉身上的罪行、他通过贿赂当选、以及在受绝罚期间擅自举行弥撒——这些事，为了天主的缘故，我不能不加调查就放过。但我希望并恳求主，莫让他在所传闻的这些事情上被查出过犯，使他的案件得以解决，而无损于我的灵魂。然而，在查明之前，我最安详的君主已在发出的命令中吩咐我，当他到来时，要以尊敬接待他。而一个被传有如此多且如此性质之事的人，在尚未经调查和甄别（这本当首先进行）之前就受尊敬，这是非常严重的事。而且，如果托付给我的诸位主教的事，在我最虔诚的君主们面前，竟在别人的庇护下被裁决，那我这个不幸之人在此教会中又能何为？但是，我的主教们藐视我，并求助于世俗审判官来反对我，我感谢全能的天主，将这一切归咎于我的罪过。我仅此略述，因我在稍等片刻；若他长久拖延不来见我，我将毫不犹豫地对他施以严格的教会纪律。但我信赖全能的天主，祂会赐予我们最虔诚的君主们长寿，并借着您的手为我们安排一切，不是按我们的罪过，而是按祂恩宠的恩赐。因此，我将这些事禀告我最安详的夫人，因为我不是不知道，殿下最纯洁的良心以何等大的热忱为正义所感动。\par

\SourceRecord{Translated by James Barmby.；Nicene and Post-Nicene Fathers, Second Series；Vol. 12.；Edited by Philip Schaff and Henry Wace.；Buffalo, NY: Christian Literature Publishing Co.,；1895.；<http://www.newadvent.org/fathers/360205021.htm>.}

% structure: p0001 source=h1 target=section reason=page-title

\section{第五卷 第二十三封信}

\Emph{致公证人\Person{卡斯托里乌}。}\par

\Person{额我略}致\Person{卡斯托里乌}等。\par

我们听到我们的兄弟及同为主教\Person{若望}去世的消息，深感悲伤，尤其因为该城此时失去了牧灵照顾的慰藉。因此，由于教会自身的许多益处要求在\Person{基督}的指引下毫不迟延地祝圣一位司铎，我们因而命令你阁下以全部紧迫劝勉圣职人员和民众，切勿延迟选举一位司铎加以祝圣。然而，在一切事之前，我们希望向你强调，他们在公共事业中不要关注自己的私利。因此，在这次选举中不可有贪婪，以免他们贪图报酬而失去辨别选择的能力，并且认为那适合此职的人，可能是取悦了他们，不是因他的功德，而是因他的礼物。因为他们应当特别且绝对知道这一点：那凡敢将天主的恩赐当作商品，以为可以用价钱购买的人，不但不配司铎职，而且肯定会罪上加罪。因此，不要选那慷慨行贿的，而要选那因功德而配的。因为刑罚将波及被选者和选举者，如果他们以亵渎的心意图玷污司铎职的纯洁。此外，不论被选者是一人或两人，务必告诫五位年长长老和五位首要人士一同前来见我们。至于圣职人员，若除决定前来的人之外，你认为有其他人必须出席，就立即派他们来我们这里，以便在整顿教会时，没有任何借口或拖延。\par

\SourceRecord{Translated by James Barmby.；Nicene and Post-Nicene Fathers, Second Series；Vol. 12.；Edited by Philip Schaff and Henry Wace.；Buffalo, NY: Christian Literature Publishing Co.,；1895.；<http://www.newadvent.org/fathers/360205023.htm>.}

% structure: p0001 source=h1 target=section reason=page-title

\section{第五卷，第二十五封信}

\Emph{致\Person{塞维鲁}，主教。}\par

\Person{额我略}致\Person{塞维鲁}，\Place{菲库卢姆}主教。\par

送交我们的报告告知我们\Person{若望}主教的去世。因此，我们郑重委托您的兄弟情谊担任丧主之教会的视察工作：这工作您应如此执行，以致无人敢干涉有关圣职人员的晋升、收入、饰物、服务，或属于同一教会产业的一切其他事项。\Emph{依惯例}。\par

\SourceRecord{Translated by James Barmby.；Nicene and Post-Nicene Fathers, Second Series；Vol. 12.；Edited by Philip Schaff and Henry Wace.；Buffalo, NY: Christian Literature Publishing Co.,；1895.；<http://www.newadvent.org/fathers/360205025.htm>.}

% structure: p0001 source=h1 target=section reason=page-title

\section{卷五，第26封信}

\Emph{致\Place{拉文纳}人民}\par

\Person{额我略}致\Place{拉文纳}的圣职人员、贵族和平民\par

获悉你们主教去世的消息，我们已托付我们的弟兄和同为主教的\Place{菲库卢姆}的\Person{塞维鲁斯}负责探望这一失去牧者的教会，并责成他不允许任何人在圣职人员晋升、教会收益、圣器装饰和礼仪服务方面擅自行事。你们应当服从他勤勉的劝勉。\Emph{按惯例。}\par

\SourceRecord{Translated by James Barmby.；Nicene and Post-Nicene Fathers, Second Series；Vol. 12.；Edited by Philip Schaff and Henry Wace.；Buffalo, NY: Christian Literature Publishing Co.,；1895.；<http://www.newadvent.org/fathers/360205026.htm>.}

% structure: p0001 source=h1 target=section reason=page-title

\section{第五卷，第29封信}

\Emph{致\Person{文科马鲁斯}，守护者（\OriginalTerm{守护者}{Defensor}）}。\par

\Person{额我略}致\Person{文科马鲁斯}等。\par

为了教会的利益，我们愿意并乐于：倘若你不受任何身体服役条件的约束，也未曾担任其他城市的圣职人员，且对你没有教规上的反对，你便要接受教会守护者的职位，以便你廉洁而乐意地执行我们为穷人利益所命令你的一切，使用经过慎重考虑后授予你的这一特权，忠诚地尽力完成我们命令你的一切，因为你要在我们的天主的审判下为你的行为交账。这封信我们口授，由我们教会的公证人\Person{帕特里乌斯}书记成文；三月，\OriginalTerm{小纪}{Indiction}第十三年。\par

\SourceRecord{Translated by James Barmby.；Nicene and Post-Nicene Fathers, Second Series；Vol. 12.；Edited by Philip Schaff and Henry Wace.；Buffalo, NY: Christian Literature Publishing Co.,；1895.；<http://www.newadvent.org/fathers/360205029.htm>.}

% structure: p0001 source=h1 target=section reason=page-title

\section{\WorkTitle{Book V, Letter 30}}

\Emph{To Mauricius Augustus}\par

格里高利致皇帝默里修斯等：\par

我主之仁爱，素来仁慈地扶持尔等臣仆，今在此处彰显于如此仁慈的供应中，以致所有软弱之人的匮乏，皆因尔等慷慨之援助而获缓解。为此，我们众人都以祈祷和泪水恳求全能天主，祂感动了陛下仁慈之心以行此事，愿祂以永恒的爱护保守我主之帝国平安，并藉其至高威能之助，在万民中拓展他们的胜利。我的同仆布萨带来的三十镑黄金，已由斯克瑞波忠实地分给了司铎、贫困者及其他人员。又因有某些献身于宗教生活的女性（\OriginalTerm{献身于宗教生活的女性}{sanctimoniales fœminæ}）从各省来到此城，逃避被掳掠之后果，凡有空位者已安置于隐修院，但其余未能容留者，则过着极端贫困的生活。故视为合宜，将从救济盲人、残障者及弱者之款项中节省下来的部分分给她们，使不仅本地的贫穷者，连到此的异乡人也都能领受我主的怜悯。因此，众人都一致为我主之生命祈祷，愿全能天主赐予你们长久而安宁的生活，并赐福尔等圣主最幸福的后裔，使之在罗马共和国中长期昌盛。士兵的薪饷也已由上述我的同仆斯克瑞波，在著名的军团长（\OriginalTerm{军团长}{magister militum}）卡斯图斯在场下分配，使所有人都怀着感谢，以应有的纪律领受我主的赏赐，并杜绝了从前在他们中间曾盛行的所有怨言。\par

\SourceRecord{Translated by James Barmby.；Nicene and Post-Nicene Fathers, Second Series；Vol. 12.；Edited by Philip Schaff and Henry Wace.；Buffalo, NY: Christian Literature Publishing Co.,；1895.；<http://www.newadvent.org/fathers/360205030.htm>.}

% structure: p0001 source=h1 target=section reason=page-title

\section{第五卷，第三十六封}

\Emph{致塞维鲁斯}, \Emph{法学士}\par

格里高利致塞维鲁斯，\Emph{法学士}致总督。\par

那些协助审判官并因真诚的依恋而与他们相连的人，应当向他们提出建议，并提示那些既能拯救他们的灵魂又不减损他们名声的事。既然如此，既然我们知道您以何等真诚的忠诚爱戴那最卓越的总督，我们谨告知您的尊驾所发生的事，以便您在了解后，能促使他合理地同意这些事。\par

那么，请知悉，伦巴第人的王阿基乌尔夫并不拒绝缔结全面和平，只要主执政官同意仲裁。因为他抱怨说，在和平期间，他的地区发生了许多暴力行为。既然若能找到合理的仲裁理由，他希望能得到满足，他自己也承诺，如果在和平期间他方有不当行为，他将在各方面给予满足。既然如此，他所要求的显然是合理的，应当予以同意，因此应有仲裁，以便若任何一方有不当行为，可予以调整；如此，在天主的眷顾下，方能建立全面和平；因为这对我们所有人何等必要，您深知。因此，请像您一贯明智地行事，使最卓越的总督毫不拖延地同意此事，以免和平似乎被他拒绝，这是不应有的。因为，若他不愿同意，他（阿基乌尔夫）确实再次承诺与我们缔结单独和平；但我们知道，各种岛屿和其他地方无疑会在那种情况下遭毁。然而，请他（总督）考虑这些事，并尽快缔结和平，以便至少在此休战期间我们能有些许安宁，而共和国的兵力也能在天主的助佑下更好地恢复以抵抗。\par

\SourceRecord{Translated by James Barmby.；Nicene and Post-Nicene Fathers, Second Series；Vol. 12.；Edited by Philip Schaff and Henry Wace.；Buffalo, NY: Christian Literature Publishing Co.,；1895.；<http://www.newadvent.org/fathers/360205036.htm>.}

% structure: p0001 source=h1 target=section reason=page-title

\section{第五卷·第三十九封信}

\Emph{致阿纳斯塔西乌斯主教}。\par

\Person{额我略}致\Place{安提约基亚}主教\Person{阿纳斯塔西乌斯}。\par

\BibleQuote{天主受享光荣于高天，主爱的人在世享平安}\BibleRef{路 2:14}，因为那条曾经使\Place{安提约基亚}的岩石干涸的大河，如今终于复归其正道，浇灌着附近低处的山谷，使之结出果实，有三十倍的，有六十倍的，有一百倍的。如今无疑有许多灵魂的花朵正在其山谷中生长，并因你舌头的溪流而结出成熟的果实。为此，我们以心灵和口舌的声音，从内心深处向全能的天主献上应有的赞美，并为你的真福而喜乐，不仅与你一同喜乐，也与所有属于你管辖的人一同喜乐。我收到了你圣座的来信，对我而言极为甘甜愉悦，而我们自己，若可如此说，正与你一同在同劳苦中流汗。我确实知道，在经历过那种安息的高峰、以心灵之手触摸天上奥秘之后，外在事务的负担对你而言必定是何等沉重。但请记得你治理的是一个宗座，并且因你成了众人所成的一切，而更加容易抚慰忧伤。在列王纪中，正如你娴熟的圣座所知，描述了一个人，他能左右手并用\BibleRef{编上 12:2}。关于此事，我毫不怀疑我往昔最甘甜、最神圣的主保、主\Person{阿纳斯塔西乌斯}，在他将世间事业引向天上的益处时，将左手当作右手使用；这样，他天上的专注可以借右手完成其工作，同时，当他因关心现世事物而被引向正义的利益时，左手也能获得右手的力量。\par

的确，这些事不能没有重大的劳苦和烦恼。但让我们记住先我们之前的人的劳苦；我们所忍受的就不会是沉重的。因为\BibleQuote{我们必须经过许多苦难，才能进入天主的国。}\BibleRef{宗 14:22} 并且，\BibleQuote{我们被压得超出力量，甚至超过了限度，以致我们厌倦了生命。但我们自己也得了死亡的判决，使我们不依靠自己。}\BibleRef{格后 1:8} 然而，\BibleQuote{我认为现时的苦难，与将来要在我们身上显示的光荣，是不能相比的。}\BibleRef{罗 8:18} 那么，我们这些软弱的羊，如何能不经劳苦而通过这世界的炎热，要知道连公羊都曾承受过沉重的劳苦呢？\par

此外，我在这片土地上所遭受的种种磨难，无论是来自伦巴第人的刀剑、法官的不义、事务的繁忙、对属民的操心，还是身体的病痛，我都无法用笔或口来表达。关于这些事，即使我能简要地说些什么，我也犹豫，免得在你为你自己的磨难所苦之际，再给你至圣的爱德加上我的磨难。愿全能的天主以其丰厚的慈爱，用一切安慰充满你至圣真福的心灵，并因你的转求，在某时赐予我这不配的人从我遭受的这些邪恶中安息。阿们。恩宠。这些词句，正如你所见，取自你曾写的话，我插入我的书信中，为使你的真福明白，对圣\Person{依纳爵}而言，他不仅属于你，也属于我们。因为，既然我们共同拥有他的导师、宗徒之长，那么同样，我们中谁也不应独自占有这位宗徒之长的门徒。此外，我们已领受了你的祝福，它馨香而甘饴，以应有的情感领受。我们感谢全能的天主，你所做的、所说的、所给予的，都是芬芳甘饴的。因此，为了你的生命，让我们一同说，让我们大家说：\BibleQuote{天主受享光荣于高天，主爱的人在世享平安。}\par

\SourceRecord{Translated by James Barmby.；Nicene and Post-Nicene Fathers, Second Series；Vol. 12.；Edited by Philip Schaff and Henry Wace.；Buffalo, NY: Christian Literature Publishing Co.,；1895.；<http://www.newadvent.org/fathers/360205039.htm>.}

% structure: p0001 source=h1 target=section reason=page-title

\section{卷五，书信四十}

\Emph{致奥古斯都·摩里修斯}\par

\Signature{格里高利致摩里修斯等}\par

朕主的虔诚在其至为安和的诏令中，虽意在反驳我某些事，却因宽宥我而绝未宽宥我。因其中所用\Quote{纯朴}一词，实有礼貌地称我为愚笨。诚然，在圣经中，当纯朴以善意被提及，常与明智及正直相联。故经上论及圣约伯写道：\Quote{那人纯朴而正直}(\BibleRef{约 1:1})。圣保禄宗徒亦劝勉说：\Quote{你们在恶事上要纯朴，在善事上要明智}(\BibleRef{罗 16:19})。真理亲自教训说：\Quote{你们要明智如蛇，纯朴如鸽}(\BibleRef{玛 10:16})；由此表明，若明智缺了纯朴，或纯朴缺了明智，便毫无益处。为使祂的仆人凡事练达，祂愿意他们既纯朴如鸽，又明智如蛇，好让蛇的狡黠磨利鸽的纯朴，鸽的纯朴调和蛇的狡黠。\par

因此，我——在朕主最安和的诏令中被指为无明智之纯朴，受阿里乌尔夫诡计所欺——显然且无疑地被称作愚笨；连我自己也承认如此。因即便陛下缄默，事实亦在呼喊。倘若我不愚笨，断不会在此伦巴第人刀剑中忍受今日之苦。此外，关于我所述阿里乌尔夫诚心预备与共和国议和一事，因我不被相信，又被责为说谎。然我虽非司铎，却知对司铎而言——作为真理的仆人而被视为欺诈——是严重的伤害。我已有时察觉，诺尔杜尔夫在我之前被信，莱奥在我之前被信，如今那些看似在陛下左右者，其言比我更易得信任。\par

诚然，若我地的俘虏境况非日益加剧，我甘愿默受对自己的轻蔑与讥讽。但这令我痛苦至极：因我背负谎言之名，致使意大利日日被伦巴第人轭下掳掠。我的陈述丝毫未被采信，而敌人的力量大增。然我向至虔之主进言：望他对我不妨作任何恶评，但关乎共和国之利及解救意大利之事，切莫轻易信任何人，宁信事实而非言语。再者，愿我主勿因其属世权柄而遽然轻慢司铎，却因祂的仆人之故，以卓越的考量治理他们，同时亦当给予应得的敬重。因圣经中司铎有时被称为神，有时被称为天使。经梅瑟论及被要求宣誓者云：\Quote{带他到神明前}(\BibleRef{出 22:8})，即到司铎前。又写道：\Quote{不可咒骂神明}(\BibleRef{出 22:28})，即司铎。先知说：\Quote{司铎的唇该保存知识，人应从他的口中寻求训诲，因为他是万军上主的天使}(\BibleRef{拉 2:7})。那么，若陛下屈尊尊敬那些连天主在祂圣言中都尊称为\Quote{天使}或\Quote{神}的人，又有何奇怪？\par

教会史亦见证：当指控主教的书面诉状呈递至可敬的君士坦丁皇帝面前时，他确收了诉状，但随即召集被控的主教，当着他们的面烧毁了所收的诉状，说道：\Quote{你们是由真神所立的神。去吧，在你们中间解决你们的事，因为我们不宜审判神。}然而，至虔之主，在这句话中，他因谦卑所赋予自己的，比因尊敬所赋予他们的更多。在他之前，共和国中曾有异教君王，他们不认识真神，却崇拜木石之神，且仍对其司铎致以最高敬意。那么，若基督徒皇帝屈尊尊敬真神的司铎，岂非当？而异教君王，如我们已言，尚知尊敬侍奉木石之神的司铎。我向朕主之虔诚进此言，非为自己，乃为所有司铎。因我是一个罪人。我每日不停地得罪全能天主，故猜想，在可怕的审判时，这每日不停的鞭笞或可作些补赎。我相信，你们越严厉地折磨我——这侍奉祂不善者——便越能平息同一位全能天主的义怒。因我已受许多鞭笞，而当主人的诏令再加于身时，我反倒找到了未期盼的安慰。若我能，将简略列举这些鞭笞。\par

首先，我与托斯卡纳的伦巴第人所缔结的和平，未曾耗费共和国丝毫，竟从我手中被剥夺。随后，和平被破坏，士兵从罗马城被调离。有些人被敌人杀害，另一些人被安置在\Place{纳尔尼}和\OriginalTerm{佩鲁贾}{Perugia}；而罗马被遗弃，以便守住\Place{佩鲁贾}。此后，更沉重的打击是\Person{阿吉卢尔夫}的到来，以至我亲眼目睹罗马人像狗一样被绳子拴着脖子，准备被带去\Place{法兰克}贩卖。而且，因为我们这些在天主保护下留在城内的人逃脱了他的魔掌，便由此寻得一个借口，使我们显得有罪；即因为粮食短缺——这在这座城市里无论如何也无法大量长期储存——正如我在另一份陈情中更详细说明的那样。就我个人而言，我丝毫未感不安，因为我声明，我的良心为我作证，我准备承受任何逆境，只要我能以灵魂的安全从中脱身。但对于卓越的官员\Person{格列高利}总督和\Person{卡斯托里乌斯}\OriginalTerm{军事指挥官}{magistro militum}，我深感忧心，因为他们毫不懈怠地尽力而为，在围城期间不眠不休地警戒守卫，而在此之后，却遭到我主上们的严厉震怒。关于他们，我清楚地明白，不利于他们的不是他们的行为，而是我这个人。因为，他们既与我一同在患难中劳苦，也在劳苦后一同受困扰。\par

至于我主上们的敬虔向我提出全能天主可畏而可怕的审判，我恳求你们，因同一位全能天主之名，不要再这样做。因为我们还不知道届时我们中间谁将如何站立。卓越的宣讲者\Person{保禄}说：\BibleQuote{什么都不要判断，直到主来到，祂要照亮黑暗中的隐密事，并显明人心中的计谋}\BibleRef{格前 4:5}。但我简短地说：作为不堪当的罪人，我倚靠耶稣来临时祂的慈悲胜于倚靠你们敬虔的正义。关于这审判，有许多事是人所不知的；因为或许祂会责备你们所赞美的，赞美你们所责备的。因此，在这一切不确定之中，我唯有转向眼泪，祈祷同一位全能天主以祂的手引导我们最敬虔的主上，并在那可怕的审判中使他免于一切过失。愿祂使我这样取悦于众人——若有必要——而不冒犯祂永恒的恩宠。\par

\SourceRecord{Translated by James Barmby.；Nicene and Post-Nicene Fathers, Second Series；Vol. 12.；Edited by Philip Schaff and Henry Wace.；Buffalo, NY: Christian Literature Publishing Co.,；1895.；<http://www.newadvent.org/fathers/360205040.htm>.}

% structure: p0001 source=h1 target=section reason=page-title

\section{第五卷，第四十一封信}

\Emph{致君士坦丁娜·奥古斯塔。}\par

国瑞致君士坦丁娜等。\par

深知我最宁静的夫人如何思想天上的家乡和她灵魂的生命，我认为，若我对于因敬畏天主而应向她陈述之事保持缄默，我将负大过。\par

得知撒丁岛上的许多土著仍然依照其种族的邪恶习俗向偶像献祭，而该岛上的司铎在宣讲我们的救赎主方面懈怠，我便派遣了意大利的一位主教前往，他借主的合作引领了许多土著归信。但他向我报告了一件渎神之事，即该岛那些向偶像献祭的人向法官行贿以获准这样做。而且，当他们中一些人受洗并停止向偶像献祭后，该岛的同一位法官甚至在他们受洗后仍索取他们先前为获得献祭许可而惯常支付的同一笔款项。当上述主教责备他时，他回答说，他曾允诺如此大的一笔\Emph{\OriginalTerm{贿赂}{suffragium}}，以至于他除了依靠此类案件所得外无法凑足。但科西嘉岛受到如此众多的敲诈者和如此沉重的勒索负担的压迫，以致岛内居民除了卖掉子女外几乎无法凑足被勒索的款项。由此导致该岛的地主们抛弃虔诚的共和国，被迫投靠最邪恶的伦巴底民族。因为还有什么样的苦难比被如此束缚和压榨以至于被迫卖掉子女更严酷或更残忍的呢？此外，在西西里岛，有一位司提反，海事部的\Emph{\OriginalTerm{书记}{chartularius}}，据说从事如此不法行为和压迫，侵占属于不同人的地方，未经任何法律程序就在产业和房屋上张贴告示，以致如果我想细数我所听说的他的每一件恶行，我无法用一卷书完成此事。\par

愿我最宁静的夫人明智地关注这一切，并减轻被压迫者的哀叹。因为我猜疑这些事尚未到达您最虔诚的耳中。因为若它们能到达，它们绝不会持续至今。但现在应适时向我最虔诚的主上陈述，以使他从自己的灵魂、从帝国、从他的儿子们身上除去如此重大的罪担。我知道他会说，从上述诸岛征收的一切都转交给我们用于意大利的开支。但对此我建议，即使拨给意大利的经费较少，他仍应使他的帝国摆脱被压迫者的眼泪。因为也许，在这片土地上如此巨大的开支之所以收益较少，是由于征收金钱时掺杂了罪恶。因此，愿我们最宁静的主上们下令，任何东西都不得以罪恶的方式征收。我知道，尽管为共和国利益付出的较少，但共和国反而得到更多的帮助。而且，虽然也许因经费较少而帮助较少，但与其您在天主永生之路上遇到任何阻碍，不如我们不暂时存活。因为当父母为了孩子不受折磨而与他们分离时，他们的感受如何、心情如何，可想而知。但如何为别人的孩子着想，有自己孩子的人是清楚的。因此，让我简要陈述这些事就够了，免得若您的虔敬不知晓这些地区所发生的事，我将在严厉的审判者面前为我缄默的罪过而受罚。\par

\SourceRecord{Translated by James Barmby.；Nicene and Post-Nicene Fathers, Second Series；Vol. 12.；Edited by Philip Schaff and Henry Wace.；Buffalo, NY: Christian Literature Publishing Co.,；1895.；<http://www.newadvent.org/fathers/360205041.htm>.}

% structure: p0001 source=h1 target=section reason=page-title

\section{第五卷，书信四十二}

\Emph{致塞巴斯蒂安，主教}\par

\Person{额我略}致\Person{塞巴斯蒂安}，锡尔米乌姆主教。\par

我已收到贤弟最甜美、最愉悦的书信，这信虽使你从未离开我心，却仍令你的圣座仿佛亲身与我同在。但我恳求全能的天主以祂的右手保护你，赐予你在此世平安的生活，并在合祂心意之时，赐予永恒的赏报。我恳求你，若你以我们在一起时你常常爱我的那份爱来爱我，就请更热切地为我祈祷，使全能的天主解开我罪恶的束缚，使我摆脱这腐朽的负担，在祂面前得以自由站立。因为，不论天乡的甜蜜多么无法估量地吸引人，此生仍有无数忧苦，每日催迫我们热爱天上的事物。而正是这些忧苦使我格外喜悦，因为它们不让我在此世有任何事物可喜悦。\par

至圣的兄弟，我们无法描述在这片土地上因你的朋友罗马诺斯大人所受的苦。然而我可以简要地说，他对我们的恶意已超过伦巴第人的刀剑；以至于杀戮我们的敌人似乎比共和国的法官更仁慈，后者以他们的恶意、掠夺和欺骗令我们焦虑不堪。而且，同时承担主教和神职人员的牧职，以及隐修院和民众的照顾，还要时刻警惕敌人的阴谋，并总是怀疑公爵们的欺诈和恶意；这一切包含多少劳苦和忧苦，贤弟你既爱遭受这些痛苦的我越纯粹，就越能真实地评估。\par

此外，在向你致我当有的问候时，我告知你，从\OriginalTerm{辩护官}{defensor}博尼法斯的报告中得知，我们的兄弟、至圣的宗主教阿纳斯塔修斯大人曾希望将一座城市的教会治理托付于你，而你拒绝同意。我极为欣悦地赞同并极力称赞你的这种情感和智慧；我认为你是幸福的，而我是不幸的，竟然在这种时候同意承担教会治理之责。然而，如果出于对兄弟们的谦让，并专注慈悲善工，你有朝一日决定同意这样的提议，我恳求你绝不可将别人的爱置于我的爱之前。因为在\Place{西西里岛}上有一些没有主教的教会。如果因天主的指引你乐意承担一个教会的治理，你能够在有福的宗徒伯多禄的门槛旁，借助他的帮助做得更好。但若你并不乐意，就安然保持现状，愿此决心留在你内；并为我们这些不幸的人祈祷。愿全能的天主在任何祂愿意你所在之处，保护你于祂的眷顾之下，并引领你获得天上的赏报。\par

\SourceRecord{Translated by James Barmby.；Nicene and Post-Nicene Fathers, Second Series；Vol. 12.；Edited by Philip Schaff and Henry Wace.；Buffalo, NY: Christian Literature Publishing Co.,；1895.；<http://www.newadvent.org/fathers/360205042.htm>.}

% structure: p0001 source=h1 target=section reason=page-title

\section{第五卷，第四十三封书信}

\Emph{致欧洛吉乌斯与阿纳斯塔修斯主教。}\par

额我略致亚历山大里亚主教欧洛吉乌斯与安提约基雅主教阿纳斯塔修斯。\par

那位卓越的宣道者说：\BibleQuote{我既然是外邦人的宗徒，我就光荣我的职务}\BibleRef{罗 11:13}；又在另一处说：\BibleQuote{我们在你们中间如同小孩一般}\BibleRef{得前 2:7}。他无疑为后世的我等树立了榜样：我们当心存谦卑，同时尊重我们品秩的尊严，既不因谦卑而畏缩，也不因高位而骄傲。\newline{} 八年前，在我的前任、圣德可纪念的佩拉吉乌斯任上，我们身处君士坦丁堡的兄弟与同侪主教若望，借另一事端召集了一次主教会议，在会中妄图自称\OriginalTerm{普世主教}{Universal Bishop}。我的前任一得知此事，便发出信函，借圣宗徒伯多禄的权威撤销了该会议的决议；我已将信函副本寄呈贵圣德。此外，他禁止随侍最虔诚的君上处理教会事务的执事，与上述我们的同侪司铎一同举行弥撒圣祭。\newline{} 我与前任同心，也向前述我们的同侪司铎发出类似信函，窃以为应将副本送交贵真福，特为此目的：我们可先以适度力量攻击我们前述兄弟在这事上的心意——他以一项新的骄傲之举，扰乱了普世教会的全部心肠。然而，若他全然不肯从高涨的骄傲中弯曲，那么，在全能天主的助佑下，我们将更具体地思考应当如何行事。\par

因为，按贵可敬的圣德所知，这\OriginalTerm{普世}{Universality}的名号，是由迦克敦神圣会议献给那藉天主眷顾我所服务的宗座之教宗的。但我的任何一位前任从未同意使用这如此亵渎的称号；因为，倘若一位宗主教被称为\OriginalTerm{普世}{Universal}，则其余人的宗主教之名便受贬损。但远非如此，远非如此，基督徒之心岂可愿有人为自己攫取那似乎会多少减损弟兄荣耀的称号？既然我们尚且不愿接受他人献上的这尊荣，那么任何人竟想强行篡夺它，是何等可耻啊。\par

因此，贵圣德在信函中切勿称任何人为\OriginalTerm{普世}{Universal}，以免您将不应给与的尊荣献给他人，从而减损了您应有的尊荣。也莫让任何不良的疑虑因我们最安宁的君主而使您不安，因为他是敬畏全能天主的，决不会同意做出任何违反福音规诫、违反最神圣教规的事。至于我，尽管与您相隔漫长的陆地和海洋，但在心中却与您完全相连。我相信贵圣德对我也是如此；因为，当您以爱还爱我时，您并不远离我。因此，我们更加感谢那芥菜籽\BibleRef{玛 13:31}，它从看似微小而可鄙的种子，竟被如此蔓延各处，枝干从同一根上升并伸展，以致天上的飞鸟都能在它们里面搭窝。也要感谢那酵母，它放在三斗面中，使全人类的团在合一中发起来\BibleRef{玛 13:33}；还要感谢那非人手从山上凿出的小石头，它已布满全地\BibleRef{达 2:35}，并且为此目的处处扩展，好使从归为一体的人类中，整个教会的身体得以完成，于是各肢体的这一区分便有助于整体相系的益处。\par

因此，我们也不远离你们，因为在无所不在的那一位里，我们是一体。让我们感谢祂，祂除灭了仇恨，使在祂的肉体内，整个世界只有一个羊群，一个羊栈，祂自己是唯一的牧者；让我们常记真理的宣道者的劝诫：\BibleQuote{尽力以和平的联系，保持心神的合一}\BibleRef{弗 4:3}，以及\BibleQuote{与众人追求和平与圣德，没有圣德，谁也见不到天主}\BibleRef{希 12:14}。他还对其他门徒说：\BibleQuote{若是可能，就你们方面，要与众人和平相处}\BibleRef{罗 12:18}。\par

但是，因为和平只能在双方之间建立，当恶人逃避和平时，善人却应将和平保存在内心深处。因此，经上也说得极好：\BibleQuote{就你们方面}，意思是即使和平被恶人的心拒绝，它仍应留在我们内。我们真正保持这样的和平，是在我们以爱德和持久的正义对待骄傲者的过错时，我们爱他们而憎恨他们的恶习。因为人是天主的工程，但恶习是人的工程。让我们分辨什么是天主所造的，什么是人所造的，既不要因人的错误而恨人，也不要因人的缘故而爱错误。\par

因此，让我们同心一致地攻击那人身上的骄傲之恶，好使那人自己先从他的敌人——即他的错误——中获得自由。我们的全能救赎主必会供给爱德与正义以力量；祂必会赐给我们，虽远隔两地，仍得祂圣神的合一；正是祂的作为，使教会如同按照世界的四面建造成的方舟，以不朽坏的木板和爱德的沥青连接在一起，不为任何反对的风，也不为从外面涌来的任何波涛所扰乱。\par

但既然藉着祂的恩宠引导我们，我们应当寻求没有外来的波浪使我们困惑，同样我们应当全心祈祷，最亲爱的弟兄们，愿祂眷顾的右手能抽出我们内在积存的舱底水。因为确实我们的仇敌魔鬼，它向谦卑者发怒，如同咆哮的狮子巡行，寻找可吞食的人\BibleRef{伯前5:8}，不再如我们所见的巡行于羊圈，而是如此坚决地咬住教会某些必要成员，除非藉着主的恩宠，谨慎的牧者群齐心救援，无人能怀疑它很快会撕毁整个羊圈——愿天主禁止。考虑，最亲爱的弟兄们，是谁紧随其后，它的临近甚至已在司铎中爆发如此乖张的开始。因为正是因为他临近，经上记着说：\BibleQuote{他是骄傲一切子民的王}\BibleRef{约41:25}——我不得不在极大悲痛中说——我们的兄弟和同为主教的若望，蔑视主的命令、规诫及教父的规矩，试图通过傲慢在名义上作它的先驱。\par

但愿全能天主向你的真福显示，我是以何等剧痛的呻吟被这种思虑折磨；他，那个曾对我而言最谦和的人，那个曾为众人所爱的人，那个看起来专注于施舍、善行、祈祷和斋戒的人，从他所坐的灰烬中，从他所宣讲的谦卑中，变得如此自夸，竟试图将一切归于自己，并借着一种浮夸表达的傲慢，企图将基督的所有肢体——它们只与一个头即基督结合——都屈服于自己。也不奇怪，那同一诱惑者知道骄傲是一切罪的开端，先前在第一个人的事上就已使用它，现在也在一些人的德行终点前放置它，为那些似乎通过他们生活的善目标已多少逃脱了他最残酷之手的人，在善工的目标处，即在完美的终点，设下陷阱。\par

因此我们应当热切祈祷，以不断的恳求哀求全能天主，愿祂从那个人的灵魂上除去这个错误，并从教会的合一与谦卑中除去这骄傲与混乱的祸害。藉着主的恩宠，我们应当同心合作，竭尽全力预备，免得在一种表达方式的毒害中，基督身体内的活肢体死去。因为，若容许这种表达方式被随意使用，所有宗主教的名誉就被否认；而当那被称为普世的人若在他的错误中灭亡，将找不到任何主教仍处于真理的状态中。\par

那么，你们应当坚定而无偏见地持守教会，如同你们所领受的，不要让这魔鬼篡夺的企图得到你们的支持。站稳吧，稳固地站稳；切勿发出或接受任何带有\Quote{普世}这一虚假称号的文件。命令你们管辖的所有主教远离这傲慢的玷污，以使普世教会承认你们为宗主教，不仅在于善工，也在于真理的权威。但是，若偶然恶果随之而来，我们应当同心坚持，甚至以死表明：在损害普遍利益时，我们不特别喜爱自己的任何事物。让我们同保禄一起说：\BibleQuote{为我而言，生活就是基督，死亡就是利益}\BibleRef{斐1:21}。让我们听听所有牧者之首所说的话：\BibleQuote{若你们为正义而受苦，便是有福的}\BibleRef{伯前3:14}。因为请相信我，我们所领受的为宣讲真理的尊严，若必要情况要求，我们将更安全地放弃它，而非为保持同一真理而保留它。最后，请为我祈祷，正如你最亲爱的真福所当行的，使我能在行为上显明我如此大胆对你们所说的。\par

\SourceRecord{Translated by James Barmby.；Nicene and Post-Nicene Fathers, Second Series；Vol. 12.；Edited by Philip Schaff and Henry Wace.；Buffalo, NY: Christian Literature Publishing Co.,；1895.；<http://www.newadvent.org/fathers/360205043.htm>.}

% structure: p0001 source=h1 target=section reason=page-title

\section{第五卷，第四十八封信}

\Emph{致\Person{安德烈}，\OriginalTerm{学者}{Scholasticus}}。\par

\Person{额我略}致\Person{安德烈}等。\par

我们一直渴望满足最杰出的显贵大人关于总执事\Person{多纳图斯}之位的愿望；但是，鉴于鲁莽地为任何人覆手极其危险灵魂，我们谨慎地通过彻底调查审视了他的生活与行为。由于我们发现许多事情——正如我们已写信给上述显贵大人——使他远离主教职，我们畏惧天主的审判，认为不宜同意他的祝圣。我们也没有胆量祝圣那位不熟悉圣咏的司铎\Person{若望}，因为这情况显然表明他对自己不够认真。这些候选人被排除后，我们敦促各方从他们自己的人中选出一位，他们宣称无人适合此职，我们与他们一同更为忧虑，最终他们异口同声、一致恳求我们可敬的弟兄、司铎\Person{马里尼亚诺}——他们得知他曾与我一同在隐修院生活很长时间。他推辞此职，最后通过各种方式勉强被说服同意他们的请求。由于我们熟知他的生活，知道他在救灵上谨慎，我们没有延迟他的祝圣。因此，愿阁下恰当地接待他，并在他初任时伸出援助之手。因为如您所知，任何职务的新手对所有人来说都是非常艰难的。但我深信全能天主既已恩准他治理自己的羊群，必会激励他关注内在之事，并以恩宠的慈爱安慰他处理外在事务。而且，由于他长久享受宁静后，正如我们先前所说，他的新手身份无疑会让他受到困扰，我恳求，当他逃离世俗风暴的旋风来到您身边时，愿他总能在您心中找到避风的港湾，并因您的爱德恩惠而欢欣。不久您就会知道自己与他们能达成多大一致；因为他是不情愿地接受主教职的。\par

\SourceRecord{Translated by James Barmby.；Nicene and Post-Nicene Fathers, Second Series；Vol. 12.；Edited by Philip Schaff and Henry Wace.；Buffalo, NY: Christian Literature Publishing Co.,；1895.；<http://www.newadvent.org/fathers/360205048.htm>.}

% structure: p0001 source=h1 target=section reason=page-title

\section{第五卷，书信四十九}

\Emph{致莱安德，主教。}\par

\Person{额我略}致\Person{莱安德}，\Person{伊斯帕利斯}（\OriginalTerm{塞维利亚}{Seville}）主教。\par

怀着何等热切渴望见您，您可在我心版上读到，因您深深爱着我。但由于山川阻隔，无法相见，我遵从爱德对您的指示，在我们的共同儿子\Person{普罗比努斯}司铎抵达此处时，将我在就任主教之初所著的\WorkTitle{牧灵规则}一书，以及您已知我早已完成的对真福约伯的诠释诸书，呈递给尊驾。确实，那部作品的第三、第四部分的一些篇章，我尚未寄给仁爱阁下，因为我已经将那些部分的篇章仅给了隐修院。因此，我恳请尊驾认真翻阅这些寄去的书卷，并更加深切地哀叹我的罪过，免得我因似乎明知当行却未行而遭受更严厉的谴责。但在这教会中，我承受着何等巨大的事务喧嚣，这封信的简短本身也足以向您的仁爱表明，因为我向那个我最爱的人说如此少的话。\par

\SourceRecord{Translated by James Barmby.；Nicene and Post-Nicene Fathers, Second Series；Vol. 12.；Edited by Philip Schaff and Henry Wace.；Buffalo, NY: Christian Literature Publishing Co.,；1895.；<http://www.newadvent.org/fathers/360205049.htm>.}

% structure: p0001 source=h1 target=section reason=page-title

\section{第五卷，书信五十二}

\Emph{致\Person{若望}总主教}\par

\Person{额我略}致\Person{若望}\Place{科林托}总主教。\par

我们的兄弟及同为主教者\Person{塞孔迪诺}的公义与关怀，早已为我们所熟知，如今更由你书信的内容得以显明。在此事上，他极大地取悦了我们，并使我们欣喜，因在调查前任主教\Person{阿纳斯塔修斯}一案中，他既勤勉地履行了审慎之责，又按正义及应然之理，对所发现的罪行做出了判决。然而在一切事上，我们感谢全能天主，因为当某些控告者退缩时，祂使真相被他知晓，以免如此重大罪行的肇始者逃脱发觉。但是，鉴于在上述\Person{阿纳斯塔修斯}被公正定罪及罢免的判决中，我们上述的兄弟及同为主教者如此处理某些人的过犯，即将其留待我们审判，我们因而宜借此书信表明，关于他们应持守及遵行何事。\par

关于持信人\Person{保禄}执事，尽管他的过错极其可耻——即受许诺迷惑，未控告其近日被罢免的前任主教，且因贪心之急切，违背自己的灵魂，同意缄默而不愿声明真相——然而，因我们宜以仁慈而非严厉，我们宽恕此过，并决定他应恢复其品级及职位。因为我们认为，自判决宣布以来他所承受的忧苦，足以惩罚此过。至于\Person{尤菲米乌斯}与\Person{托马斯}，他们因放弃控告而接受了圣秩，我们决定他们应被剥夺这些圣秩，罢免后保持如此；并规定他们永不可借任何托词或理由恢复圣秩。因为，他们接受所获之尊位，并非因其功绩，而是作为恶行之酬报，这极不合宜，且违背教会纪律之规则。然而，因我们宜倾向于仁慈而非严格正义，我们决定上述\Person{尤菲米乌斯}与\Person{托马斯}应恢复其原品级及职位（但仅限此原品级，即他们被擢升圣秩前所处的品级），并终生领取这些职位之俸禄，如前所惯。此外，关于读经员\Person{克莱马修斯}，我出于同样仁爱之心，指定他应恢复其品级及职位。对于以上所有人，即\Person{保禄}执事、\Person{尤菲米乌斯}、\Person{托马斯}及\Person{克莱马修斯}，请阁下按其各人所处的品级及职位，依其向来领取之方式，从本第十三小纪起，毫无减损地供给其报酬。既然上述\Person{保禄}执事声称为贵教会的益处花费甚多，并希望借阁下之助追回，我们敦促，若果如此，你应以一切可能方式协助他，并支持他追回所赠之物，因为无理由容许他为大众之利所花费者遭受不公损失。此外，请阁下立即归还上述三磅黄金，即应我们上述兄弟及同为主教者\Person{塞孔迪诺}之请，该\Person{保禄}执事为贵教会的益处所给予的，以免（愿天主不许）你似乎非理性地，而仅凭任性加重他的负担。\par

\SourceRecord{Translated by James Barmby.；Nicene and Post-Nicene Fathers, Second Series；Vol. 12.；Edited by Philip Schaff and Henry Wace.；Buffalo, NY: Christian Literature Publishing Co.,；1895.；<http://www.newadvent.org/fathers/360205052.htm>.}

% structure: p0001 source=h1 target=section reason=page-title

\section{第五卷，书信第五十三}

\Emph{致\Person{维吉里}主教。}\par

\Person{额我略}致\Person{维吉里}，\Place{Arelate}（\Place{阿尔勒}）主教。\par

啊，爱德何其美好！它通过心中的图像，将不在场者如同在场者一样呈现给我们；藉着爱德，分裂的得以合一，混乱的得以安定，不平等的得以相连，不完美的得以成全。那位杰出的宣讲者正确地称它为全德的联系，因为虽然其他德行确实产生完美，但爱德将它们结合起来，使它们不再能从爱者的心中松开。亲爱的兄弟，我发觉你充满了这一德行，因为无论是从高卢地区来的人，还是你写给我的信函，都向我证实了这一点。\par

至于你在来信中，按照古老的惯例，请求使用祭披和宗座代牧权，我绝不会怀疑你追求的是短暂权力的崇高或外在敬礼的装饰。但既然众所周知，神圣信仰从宗座传到高卢地区，当你的兄弟情谊请求恢复宗座的古老惯例时，这岂不是善良的后代回归母亲的怀抱吗？因此，我们欣然同意你所请求的，以免我们似乎要么撤回了对你应有的任何荣誉，要么轻视了我们最杰出的儿子希尔德贝尔特国王的请求。然而，目前的情况要求更大的热忱，以便随着尊严的增加，关注也得以增长，在守护他人方面的警惕也得以增长，你生活的功德成为你下属的榜样，并且你的兄弟情谊永远不要通过赋予你的尊严寻求你自己的利益，而是寻求天乡的利益。因为你蒙福的宗徒叹息着说：\BibleQuote{众人只顾自己的事，不顾耶稣基督的事} \BibleRef{斐 2:21}。\par

因为我们从某些人的报告中得知，在高卢和日耳曼地区，没有人能不经送礼而获得圣职。如果是这样，我含泪而言，我叹息而说：当司祭品秩内在堕落时，它也无法长久外在站立。因为我们从福音中知道我们的救赎主亲自做了什么：他进入圣殿，推倒了卖鸽子者的座位（\BibleRef{玛 21:12}）。因为卖鸽子就是为圣神收取暂时的代价，圣神与天主同性同体，全能的天主通过覆手赐予人。从这种恶行，已经预示着将有什么后果。因为那些胆敢在天主圣殿里卖鸽子的人，由于天主的审判，他们的座位倒塌了。\par

事实上，这种过犯在下属中蔓延加剧。因为一个人通过金钱被提升到任何神圣品级，在他晋升的根源上已经腐化，就更容易把他所买的卖给别人。那么经文所写的在哪里：\BibleQuote{你们白白得来的，也要白白分施} \BibleRef{玛 10:8}？\par

而且，既然西蒙尼异端是第一个兴起攻击圣教会的，为何不思考，为何不看，凡用金钱祝圣任何人的人，在提升他时，就使他成为异端者？\par

另一件非常可憎的事情也被报告给我们：有些人，身为平信徒，因贪图暂时的荣耀，在主教去世后剃发，并且一下子成为司铎。在这种情况下，已经知道他是怎样的人：突然从平信徒状态晋升到神圣领导地位。一个从未当过兵的人，敢于成为修会的领袖。从未听过别人讲道的人，怎么能讲道？或者，尚未哀悼过自己罪过的人，怎么能纠正别人的过错？而保禄宗徒禁止\OriginalTerm{新教友}{neophyte}进入圣秩，我们应当理解，正如那时新教友称为刚在信仰中栽种的人，同样我们现在把那些在圣洁生活中还是新手的人也视为新教友。\par

此外，我们知道，墙壁被建造之后，不会立刻承受木料的重压，而是要等到它们新造的潮气干燥，否则如果在它们尚未稳固之前就压上重物，就会使整个建筑一起倒塌到地上。当我们砍伐树木用于建筑时，我们等待它们青翠的潮气先干燥，否则如果还在新鲜时就压上建筑的重压，它们会因过于新鲜而弯曲，并因过早被抬高而更快折断倒下。那么，为什么在人类身上不谨慎地关注这一点，而即使在木材和石头上也如此仔细地考虑呢？\par

因此，你的兄弟情谊必须注意劝诫我们最杰出的儿子希尔德贝尔特国王，使他完全从他的王国中除去这种罪的污点，以便全能的天主看到他所爱的是天主所爱的，他所避免的是天主所恨的，从而赐予他更大的赏报。\par

因此，我们按照古老惯例，在天主之下，将我们在我们最杰出的儿子希尔德贝尔特国王统治下的各教会的代牧权托付给你的兄弟情谊，但各个总主教按照古老惯例保留他们应有的尊严。我们还送给你一件祭披，供你在教会内举行弥撒时使用。此外，若有任何主教因故想要远行，没有你圣座的授权，他不得前往别处。如果主教之间产生了信仰问题或其他问题，且不易解决，应在十二位主教的集会中讨论和决定。但如果在调查真相后仍无法决定，应提交我们的判断。\par

愿全能的天主将你置于祂的保护之下，并恩赐你以你的行为保持你所获得的尊严。\newline{} \Signature{公元八月十二日，第13期小纪。}\par

\SourceRecord{Translated by James Barmby.；Nicene and Post-Nicene Fathers, Second Series；Vol. 12.；Edited by Philip Schaff and Henry Wace.；Buffalo, NY: Christian Literature Publishing Co.,；1895.；<http://www.newadvent.org/fathers/360205053.htm>.}

% structure: p0001 source=h1 target=section reason=page-title

\section{第五卷，第五十四封书函}

\Emph{致\Person{希尔代贝特}王国全体主教。}\par

\Person{额我略}致\Place{高卢}隶属于\Person{希尔代贝特}王国的全体主教。\par

天主上智的安排为此规定了应有不同的等级和分明的秩序，为的是当下位者对上位者显示敬意，而上位者施爱于下位者时，从多样性中便产生了和谐的连结，各项职务的管理也能得以恰当承担。的确，整体无法以其他方式存续，除非这种伟大的差异秩序维系着它。再者，受造物不能以绝对平等的方式被治理或生存，这一教训来自天军的榜样，因为既有天使，也有总领天使，显然它们并非平等；而是在能力和等级上，如你们所知，彼此不同。如果连这些无罪者之中都显然有这种区别，那么有谁能在明知甚至天使都服从这种秩序的情况下，拒绝自愿顺服呢？因为从此和平与爱德互相拥抱，和谐的真诚在彼此相爱中坚定不移，而这爱是悦乐天主的。\par

既然每一项职责唯有在有一位可向其请示的领导人时才能得到有益地履行，因此我们认为，在我们最杰出的儿子希尔代贝特国王统治下的各教会，根据古老习惯，将我们的代牧权授予我们的兄弟、阿尔勒城的维吉利乌斯主教，是合宜的，以便天主教的信仰，即四大神圣会议的完整，在天主的保护下以殷勤的虔诚得以保存，并且，如果在我们的兄弟和同僚司铎之间偶然发生任何争执，他能作为宗座代表，以其权威的严谨，以谨慎的节制平息争执。我们也委托他，如果在某些案件中出现需要他人出席的争议，他应召集适当数量的兄弟和同僚主教，在适当考虑公平的情况下有益地讨论此事，并以教规的完整性作出裁决。但如果（愿天主的大能避免）在信仰问题上发生争执，或出现任何可能引起严重疑问的事务，而他需要宗座的判断以取代自己的权柄，我们指示他，在勤勉查究真相后，应自行报告将此问题提交我们认知，以便通过适当的判决终结疑问。\par

并且，由于主教们有必要在他（我们授予代牧权者）面前于适当时机集合开会，只要他认为合适，我们劝告你们中任何人不得擅自行事，违抗他的命令，或推迟参加全体会议，除非或许因身体虚弱阻碍某人，或有正当理由允许其缺席。然而，凡因故无法出席主教会议者，应派一位司铎或执事代替自己，以便凭天主的助佑由我们的代牧所决定的事项，能通过他所派遣者的忠实报告传达给缺席者，并以坚定不移的恒心予以遵守，从而没有借口胆敢违反。\par

关于此事，我们也提醒你们，在未经我们前述兄弟和同僚主教维吉利乌斯的授权时，你们中任何人都不得以任何方式远行到远处，因为知道我们的前任们授予其前任代牧权的命令无疑规定了这一点。\par

此外，我们劝告你们每一位要留心自己的职责，以便那渴望获得牧羊报酬的人，以谨慎和祈祷守护托付给他的羊群，以免游走的狼侵入并撕裂托付给他的羊，而在报应时得到惩罚而非奖赏。因此，最亲爱的弟兄们，我们希望并恳求全能的天主以我们全部的祈祷，使你们在祂爱德的恒心中越来越热切，并特别恩赐你们在教会的和平与彼此合一中被保守。\par

有人向我们报告，有人通过西满尼异端被晋升至神圣神品；我们已经命令我们上述兄弟和同僚主教维吉利乌斯，此事必须完全禁止；为使你们的兄弟团体知晓并认真遵守，我们给他的信函将在你们面前宣读。颁于八月十二日，第十三小纪。\par

\SourceRecord{Translated by James Barmby.；Nicene and Post-Nicene Fathers, Second Series；Vol. 12.；Edited by Philip Schaff and Henry Wace.；Buffalo, NY: Christian Literature Publishing Co.,；1895.；<http://www.newadvent.org/fathers/360205054.htm>.}

% structure: p0001 source=h1 target=section reason=page-title

\section{\WorkTitle{第五卷，书信五十五}}

\Emph{致希尔德贝尔特王}\par

\Person{额我略}致法兰克王\Person{希尔德贝尔特}。\par

尊驾的来函令我们无比欣喜，因它证明您以虔诚的孝爱，关心司铎应得的尊荣与敬意。您藉此向众人表明，您是忠信的天主崇拜者，因您以应有的、悦纳人的敬礼爱戴祂的司铎，并以基督徒的热忱，竭力促成任何能提升他们地位的事。因此，我们欣然接受您所写的，并乐意恩准您所渴望的；为此，我们蒙天主的恩宠，依照古老惯例及尊驾的意愿，将我们的代牧管辖权托付给我们的兄弟、\Place{阿尔勒}城的主教\Person{维尔吉里}，并依古制赐予他使用披带的权利。\par

然而，由于有些事已被禀报给我们，它们极其冒犯全能的天主，并玷污司铎应得的尊荣与敬意，我们恳求这些事能藉您权力的惩治而得到全面改正，以免桀骜不驯、乖谬的行为与您的虔诚相悖，使您的王国或您的灵魂（愿天主禁止）因他人的罪过而受重负。\par

此外，我们得知：在主教去世时，有些人身为平信徒，剃发之后，一跃而升任主教。如此，未先作门徒者，竟因轻率的野心成为师傅。由于他未曾学习应教导之事，便空有司铎职衔；因为他在言语和行为上仍如前是个平信徒。那么，他怎能为他人的罪过代祷，他既未先为自己的罪过哀哭？这样的牧者不是保护，而是欺骗羊群；因为他因羞耻而不能劝诫他人做他自己也不做的事，这岂不是主的人民成为强盗的猎物，从本应获得强大健康保护之源，反而招致毁灭？此事何其恶劣、何其乖谬，愿尊驾的高见甚至从您自己的行政事务中也可想见。因为，您肯定不会将军队交给一个长官，除非他的行事和忠诚已先显明；除非他先前生活中的德行与勤勉已证明他合适。但，如果军队的指挥权只交给这样的人，那么由此类比，灵魂的牧者应当是怎样的，便易于推断。然而，这对我们是一种耻辱——我们羞于说出——司铎们攫取领导权，却连宗教战争的开端都未曾看见。\par

此外，还有一件最可憎之事，也已被禀报给我们：圣秩因\OriginalTerm{西满异端}{simoniacal heresy}而授予，即为收受贿赂。既然以价格而非功绩擢升某人至任何圣秩，乃是极具毒害、且与普世教会相悖之事，我们恳请尊驾下令将这如此可憎的恶行从您的王国中驱逐。因为，那人不害怕用金钱购买天主的恩赐，并妄图以代价获取他本不配因恩宠而拥有之物，便证明自己完全不配此职。\par

最杰出的儿子，我劝诫您这些事，正因为如此：我渴望您的灵魂得救。我原本早该写信论及这些，若非无数事务阻碍了我的意愿。如今，既然回复您书信的适当时机已至，我并未省略我应尽之责。为此，我们以父爱的亲情问候尊驾，恳求我们吩咐上述兄弟及同僚主教所行、所守的一切，都能在您眷顾的保护下得以实行，并且您不容许任何人因傲慢或骄矜而加以扰乱。反而，正如在他的前任、在您光荣之父的治世下所遵守的，如今也因您的协助，以热忱的虔敬得以遵守。因此，我们理当得到回报；我们既未延迟成就您的意愿，您也当为天主及宗徒之长圣伯多禄的缘故，使我们所定的条例在一切方面得到遵守；如此，尊驾那值得赞颂、中悦天主的声誉方能广传四方。八月十二日，\OriginalTerm{小纪}{Indiction}第十三年。\par

\SourceRecord{Translated by James Barmby.；Nicene and Post-Nicene Fathers, Second Series；Vol. 12.；Edited by Philip Schaff and Henry Wace.；Buffalo, NY: Christian Literature Publishing Co.,；1895.；<http://www.newadvent.org/fathers/360205055.htm>.}

% structure: p0001 source=h1 target=section reason=page-title

\section{第五卷，第五十六封书信}

\Emph{致马里尼亚努斯主教}\par

额我略致\Place{拉文纳}主教马里尼亚努斯。\par

蒙宗座的仁爱与古代惯例的规约，我们以为宜将\OriginalTerm{肩带}{pallium}的使用授予你的弟兄，你已知奉命承担了拉文纳教会的治理职务。且要记住，你只能在你们本城的教堂内，在众子（即\Emph{平信徒}）散去之后，从接见厅前往举行神圣弥撒庆典的路上使用它；但当弥撒结束后，你要注意再次将它收存在接见厅中。但在教堂外，我们不再允许你使用它，除非一年四次，在我们向你前任若望提及的那些连祷中；同时给你这项劝诫：就如藉着主的恩赐，你从我们这里获得了此类饰物以尊崇司祭职，同样你也要努力以品行与行为的正直来装饰你所承担的职务，以光荣基督。因为，如果你身体的这件服饰与你灵魂的美善相合，你将因这两种相互呼应的装饰而显耀。此外，所有过去显然已经授予你们教会的特权，我们也以我们的权威予以确认，并规定它们继续保持不受侵犯。\par

\SourceRecord{Translated by James Barmby.；Nicene and Post-Nicene Fathers, Second Series；Vol. 12.；Edited by Philip Schaff and Henry Wace.；Buffalo, NY: Christian Literature Publishing Co.,；1895.；<http://www.newadvent.org/fathers/360205056.htm>.}

% structure: p0001 source=h1 target=section reason=page-title

\section{第五卷，第五十七封书信}

\Emph{致若望主教。}\par

额我略致格林多主教若望\par

如今，我们无隐不察的天主，既已从祂教会的治理中驱逐了可憎的污秽瘟疫，且乐意擢升你掌管教会，故你务必殷切谨慎，好使主的羊群在承受前任牧者带来的创伤和种种祸害之后，能由你的仁爱获得安慰和良药。因此，让你行动的双手拭去先前传染的污点，不容那可恶邪恶的丝毫痕迹存留。\par

因此，让你对属下的关怀值得赞扬。温和地展现纪律，明智地施行责备。让仁慈缓和愤怒，让热忱激励仁慈：使二者互相调和，以免过度的惩罚造成不应有的痛苦，或松弛懈怠损害纪律的公正。让你的仁爱的行为成为托付给你之民众的教训。让他们在你身上看到所应爱的，并觉察所应效法的。让他们通过你的榜样学习如何生活。让他们不因你的引领而偏离正道；让他们跟随你找到通往天主的道路；这样，你从人类救主那里所得的赏报，将与你为祂赢得的灵魂一样多。所以，劳苦吧，最亲爱的弟兄，将你全部的心神活动导向这样的目标，使你有朝一日堪当听到：\BibleQuote{好啊！善良忠信的仆人，进入你主的福乐吧！}\BibleRef{玛 25:21}\par

正如你通过我们的弟兄和同僚主教安德肋转交的信件中所求，我们已将\OriginalTerm{白羊毛披肩}{pallium}寄给你，你需按照你的前任们在我等前任的许可下被证明曾使用的方式使用它。\par

此外，我们听闻在那些地区，无人能不经贿赂而获得任何圣职。若果真如此，我含泪而言，悲叹地说：当司铎的品级内在堕落时，外在也无法长久站立。因为我们从福音得知我们的救主亲自做了什么：祂怎样进了圣殿，推翻了卖鸽子者的座位\BibleRef{玛 21:12}。因为卖鸽子就是为圣神收取现世的代价，而圣神与全能的天主同性同体，是通过覆手赐予人的。正如我先前所说，这邪恶的后果已被预示：那些胆敢在天主圣殿中卖鸽子者的座位，因天主的审判而倒塌。事实上，这过犯在属下中越发蔓延。因为一个人若在晋升的最根本处就已沾染污点，他便会更愿意将他所买来的卖给他人。既然西门异端是首先兴起反对圣教会的，为何不思考，为何看不见，任何人用价钱晋升他人，就是使他成为异端者？那时，又哪里还有\BibleQuote{你们白白得来的，也要白白分施}\BibleRef{玛 10:8}这句话呢？既然圣教会彻底谴责这最凶残的邪恶，我们劝勉你的仁爱，务必以你全部的热忱，在所有你辖下的地方，抑制这如此可憎且重大的罪过。因为，若我们今后察觉此类事情发生，我们将不再以言语，而是以教会法律来纠正；并且我们将开始对你有不同的看法，这不应发生。\par

再者，你的仁爱知道，从前\OriginalTerm{白羊毛披肩}{pallium}若不收取代价是不予赐予的。但因此事不合宜，我们在真福宗徒之长伯多禄的遗体前召开了会议，严令禁止为此事或为授圣职收取任何财物。\par

你的职责便是：既不可因代价，也不可因某些人的恩泽或请求，而同意任何人晋升圣职。因为，如我们所说，这是大罪，我们绝不能容忍其继续而不加责斥。\par

我延迟接纳上述我们的弟兄和同僚主教安德肋，因为通过我们的弟兄和同僚主教塞昆狄努斯的报告，我们得知他在针对拉里萨的若望的诉讼中伪造了信函，自称来自我们。若非你的善心曾促动我们，我们绝不会接纳他。八月十五日，小纪十三。\par

\SourceRecord{Translated by James Barmby.；Nicene and Post-Nicene Fathers, Second Series；Vol. 12.；Edited by Philip Schaff and Henry Wace.；Buffalo, NY: Christian Literature Publishing Co.,；1895.；<http://www.newadvent.org/fathers/360205057.htm>.}

% structure: p0001 source=h1 target=section reason=page-title

\section{第五卷，第五十八封书信}

\Emph{致全赫拉迪亚的主教们}\par

额我略致驻守赫拉迪亚省的所有主教。\par

最亲爱的弟兄们，我与你们一起感谢全能的天主，祂使远古仇敌所引入的隐藏疮伤为众人所知，并以有益的切口将其从祂教会的身体上切除。我们既有理由喜乐，也有理由悲伤；喜乐是因为惩处了罪行，悲伤则是因为一位兄弟的跌倒。然而，既然一人的跌倒多半会成为另一人的保障，凡害怕跌倒的人，就当留意：不给仇敌以可乘之机，也不以为所行之事能隐藏不露。因为真理宣告说：\BibleQuote{没有隐藏的事不被揭露的}，\BibleRef{玛 10:26}。这声音已是吾人行为的先驱，祂亲自作证，将秘密所做的一切都带到众人面前。在既是见证者又是审判者的祂面前，谁能试图隐藏自己的行为呢？然而，有时当我们专注一事时，却疏忽了另一事，因此人人都当警惕仇敌的一切罗网，免得在一处获胜，却在另一处被击败。地上的仇敌想要攻占设防之地时，也运用同样的战术。他暗中设下伏兵，却装作完全专注于攻打一处，以致当众人为危险迫近之地而奔往防守时，其他未曾被怀疑的地方就被夺取了。结果，那本已被察觉并被对手的英勇击退者，却以诡诈获得了无法通过战斗得到的东西。既然在这一切事上都需要神圣保护的援助，我们每个人都当以心灵之声向主呼求，说：\BibleQuote{上主，求你不要远离我，我的勇力，快来救扶我！}\BibleRef{咏 21:20}。因为显然，除非祂亲自帮助并保护那些呼求祂的人，我们的仇敌就不能被战胜。\par

此外，要知道，我们通过我们的兄弟和同为主教的安德肋收到了您的爱德的来信，已将披肩\OriginalTerm{披肩}{pallium}送交给了我们的兄弟、格林多主教若望；你们理应服从他，尤其因为古代习俗的秩序要求如此，而你们亲自见证的他的优点也促使如此。因为从某些人给我的报告中得知，在那些地区，没有人不支付代价而获得任何圣职。如果真是这样，我要含泪宣告，我要叹息声明：当司祭职内在堕落后，它也无法在外在长久站立。因为从福音中我们知道我们的救赎者曾亲自做了什么：祂进入圣殿，推倒了卖鸽子者的凳子。因为卖鸽子实际上就是为圣神接受现世的报酬，而全能的天主通过覆手将与自己同体的圣神赐予人。正如我之前所说，这邪恶的后果已被暗示：那些胆敢在天主圣殿中卖鸽子者的凳子，因天主的审判而倒塌。实际上，这过犯在属下中蔓延加剧。因为那位在晋铎的根源上已被玷污的人，自己更准备将所买来的转卖给他人。经上记载的在哪里呢？\BibleQuote{你们白白得来的，也要白白分施。}\BibleRef{玛 10:8}。既然买卖圣职异端\OriginalTerm{买卖圣职异端}{simoniacal heresy}是首先兴起反对圣教会的，为什么我们不考虑、看不见：凡因价而擢升某人，就使他成了异端者呢？因此，我们劝诫你们，任何人都不许再容忍此事发生，或敢因任何人的情面或请求而擢升任何人至圣秩，除非其生活与行为的品格已证明他堪当。因为如果我们今后发现相反的情况，要知道它将受到严厉而合乎教规的惩罚。\Signature{颁于八月十五日，第十三小纪\OriginalTerm{小纪}{Indiction}。}\par

\SourceRecord{Translated by James Barmby.；Nicene and Post-Nicene Fathers, Second Series；Vol. 12.；Edited by Philip Schaff and Henry Wace.；Buffalo, NY: Christian Literature Publishing Co.,；1895.；<http://www.newadvent.org/fathers/360205058.htm>.}

% structure: p0001 source=h1 target=section reason=page-title

\section{第六卷 第一封信}

\Emph{致马里尼亚诺主教}\par

额我略致拉文纳主教马里尼亚诺。\par

不合理的请求不应允准，同样，凡请求合法之事也不应置诸不理。现今贤弟的司铎、执事及圣职人员向我们提交陈情书，申诉已故的前任若望立下遗嘱，以各种遗赠加重其教会负担。他们请求，凡有损其教会之遗赠，因法律所禁，无论以何借口均不得支付。虽然，由于他已放弃继承与承继，没有任何理由迫使贤弟满足此类要求，但我们仍特此劝勉贤弟：对于他违反法律规定、就属于其教会的财产或其在主教任内所取得的财产所作的遗赠，贤弟既不要予以支持，也不要无论如何同意。但是，若他对自己在主教任前拥有且尚未赠予其教会的私人财产希望或指示有所处置，则此处置必须在各方面视为有效，任何教会人员不得以任何借口无理企图撤销。\par

然而，由于他在生时曾多次恳求我们，以我们的权威确认他授予自己在圣亚坡理纳教堂旁所建隐修院之物，我们已应允此事。因此，我们认为有必要劝勉贤弟：凡他在该隐修院中授予并设立之物，不可使之减少，而应确保一切得以保存并稳固确立。既然他已知在遗嘱中提及此隐修院及授予该院之财产，你们必须明白，我们确认此部分并非因我们遵循他的遗愿，而是因为正如我们所说，我们在他在生时已应允他。因此，贤弟当速速妥善完成这一切，以使他在上述隐修院所设立并经我们确认之事得以维持，而他为损害其教会而遗嘱指示给予或执行之事不得有效，因为法律禁止之。\par

\SourceRecord{Translated by James Barmby.；Nicene and Post-Nicene Fathers, Second Series；Vol. 12.；Edited by Philip Schaff and Henry Wace.；Buffalo, NY: Christian Literature Publishing Co.,；1895.；<http://www.newadvent.org/fathers/360206001.htm>.}

% structure: p0001 source=h1 target=section reason=page-title

\section{第六卷，第二封信}

\Emph{致\Place{拉文纳}的圣职人员与教友}\par

\Person{额我略}致\Place{拉文纳}教会的圣职人员与教友。\par

我们得知，某些人受恶神唆使，妄图用虚假言论败坏你们的思想，涉及我们的弟兄及同主教\Person{马里尼亚诺斯}的名誉；他们声称，这位我们的弟兄对神圣的加采东大公会议的尊崇不及他所应尽的。关于此事，他本人将亲自向你们所有人证实他信仰的正统，而我们充分作证：他自幼在神圣的公教会怀抱中成长，一生持守信仰的正确宣讲，并以其生命为证。因为，他尊崇神圣的尼西亚大公会议（在该会议上\Person{亚略}受谴责）、君士坦丁堡大公会议（在该会议上\Person{马其顿尼}受谴责）、第一次厄弗所大公会议（在该会议上\Person{聂斯多略}受谴责），以及神圣的加采东大公会议（在该会议上\Person{狄奥斯哥禄}与\Person{欧提基}受谴责）。任何人若胆敢说任何反对这四次大公会议之信仰、反对神圣记忆的教宗\Person{良}之书函与定义的话，应受绝罚。因此，你们既已获得最充分的满足，就当以纯洁之心、全备的爱德热爱你们的牧者，好使你们的这位牧者那在天主前纯净倾流的转祷，能有利于你们。\par

\SourceRecord{Translated by James Barmby.；Nicene and Post-Nicene Fathers, Second Series；Vol. 12.；Edited by Philip Schaff and Henry Wace.；Buffalo, NY: Christian Literature Publishing Co.,；1895.；<http://www.newadvent.org/fathers/360206002.htm>.}

% structure: p0001 source=h1 target=section reason=page-title

\section{第六卷，第三封信}

\Emph{致萨洛纳的玛克西穆}\par

\Person{额我略}致\Person{玛克西穆}，萨洛纳教会的觊觎者。\par

每当听说有违反教会纪律之事，我们决不敢不加审查，免得因纵容而在天主面前获罪。如今我们有耳闻，您是通过买卖圣物的异端而受祝圣的。不仅如此，我们在此还听到关于您的许多其他事情，其中有一件尤其促使我们认为有必要通过信件紧急禁止您举行弥撒庆典，直到我们能更确切地查明情况。为此，为避免教会子女长期无牧，也为了避免若这些传言不经审查，此类恶习会蔓延至许多人，我们劝您放下一切托词，尽快前来我们这里，以便我们能在秉公的前提下查知这些事，并按照教会规章加以终结，基督指引我们路径。但您要如此行事，以免您一再拖延前来，以致您的缺席反而使您更受这些指控的牵连，也免得我们被迫在会议上对您作出更严厉的判决——不仅针对您逃避洗脱的指控罪行，也针对您不顺服的过失，即视您为顽抗者。\par

\SourceRecord{Translated by James Barmby.；Nicene and Post-Nicene Fathers, Second Series；Vol. 12.；Edited by Philip Schaff and Henry Wace.；Buffalo, NY: Christian Literature Publishing Co.,；1895.；<http://www.newadvent.org/fathers/360206003.htm>.}

% structure: p0001 source=h1 target=section reason=page-title

\section{第六卷 第五封信}

\Emph{致\Person{布伦希尔德}王后。}\par

\Person{额我略}致\Place{法兰克}王后\Person{布伦希尔德}。\par

你的贤明的良善与可悦天主之处，既彰显于你治理王国，亦彰显于你教诲儿子。你不但以睿智的关切，为他保全了现世之荣耀无虞，并且以合宜的母爱、可赞的教养之顾，将他的心思扎根于真信德中，从而也顾及了永生的赏报。因此，他当超越万国之邦，非无端由，因为他既纯然地崇拜，又真诚地宣认万民的造主。但为使信德在他身上藉行为更可赞地闪耀，愿你的劝勉之言激励他，从而，正如王权使他高居人上，亦使善行在天主面前成为大。\par

既然以往在许多实例中，我们对你的贤明的基督信仰怀有信心，我们因你全心所爱的宗徒之长圣伯多禄之爱，恳求你：愿你的庇护之助，珍视我们至爱的子、携此信之司铎\Person{坎迪杜斯}，以及他所治理的小产业，使他得着你扶持之恩，既能忠实管理这小产业——它显然有益于穷人之费用——又能收回任何曾被夺去之物，归于这小产业之所有。因为此产业长久以来未有教会人士被派管理，这并非不增加你的赞美。愿你的贤明，如此情愿关注我们所求之事，以致蒙福的宗徒之长圣伯多禄——主耶稣基督赐予他束缚与释放之权——既使你的贤明在此世因你的后裔欢欣，亦在多年之后，使你得脱一切邪恶，立于永恒审判者面前。\par

\SourceRecord{Translated by James Barmby.；Nicene and Post-Nicene Fathers, Second Series；Vol. 12.；Edited by Philip Schaff and Henry Wace.；Buffalo, NY: Christian Literature Publishing Co.,；1895.；<http://www.newadvent.org/fathers/360206005.htm>.}

% structure: p0001 source=h1 target=section reason=page-title

\section{第六卷，第六封信}

\Emph{致希尔德贝尔特王。}\par

\Person{额我略}致\Person{法兰克王希尔德贝尔特}。\par

君王的高贵超出常人多少，事实上，您的王国的崇高地位也超出其他民族的王国多少。然而，做王并不稀奇，因为还有其他王；但做一位\OriginalTerm{公教信徒}{Catholic}，却是别人不配做的，这已足够。因为正如一盏大灯的光辉在尘世之夜的黑暗中以其光明闪耀，同样，您信德的明光也在其他民族的黑暗背信之中闪烁发光。其他君王所夸耀的一切，您都拥有。但他们在这方面远远不及，因为他们没有您所拥有的首要善物。那么，为使他们在行为上也如在信德上一样被超越，但愿您的尊威常向您的臣民显示仁慈。若有任何触怒您的事情，未经调查不要惩罚。因为您越是克制自己的权力，感到自己不可为所欲为，就越能中悦万王之王，即全能的主。\par

如今，您对真福宗徒之长\Person{伯多禄}的爱德显然表明您在心志和行为上持守了信德的纯洁，因为在他的产业上，在您至高权柄的治理之下，至今得到良好的管理并保存。但由于贵族\Person{迪纳缪斯}——我们曾推荐他照管这份产业——据我们所知，如今无法管理它，唯恐在您地区的这份小产业因疏忽而毁坏，因此我们派遣了本函的呈递人，我们至爱的儿子司铎\Person{坎迪杜斯}前去管理。我们恳请您以慈父之爱首先向他致意，并在各方面举荐他，请求：若该处有任何不公正之事，或同一份小产业的任何财物被某人扣留，愿您加以纠正，并将被转移之物归还原主；如此，您的公正与信德将在万民面前闪耀，这将是极其光荣和可赞美的。\par

此外，我们向您的尊威遣送了\Person{圣伯多禄}的钥匙，其中含有他锁链的一部分，以便挂在您的颈上时，保护您免受一切邪恶。\par

\SourceRecord{Translated by James Barmby.；Nicene and Post-Nicene Fathers, Second Series；Vol. 12.；Edited by Philip Schaff and Henry Wace.；Buffalo, NY: Christian Literature Publishing Co.,；1895.；<http://www.newadvent.org/fathers/360206006.htm>.}

% structure: p0001 source=h1 target=section reason=page-title

\section{第六卷 第七封信}

\Emph{致\Person{坎迪杜斯}司铎。}\par

\Person{格列高利}致前往\Place{高卢}教产的\Person{坎迪杜斯}司铎。\par

现今你蒙我们的主天主耶稣基督的助佑，前往管理在\Place{高卢}的教产，我们希望你爱德用你收到的钱为穷人购买衣服，或购买十七八岁的\Place{英格兰}男孩，他们若能舍身于隐修院中归奉天主，必能获益。如此，\Place{高卢}的钱财（不能在我国花用）即可在当地善加使用。此外，你若能从称为\Emph{ablatae}的税赋收入中取得任何款项，我们也希望你用这钱为穷人购买衣服，或如我们之前所说，购买能在全能的天主的工作中获益的男孩。但因为在那里能找到的男孩都是外教人，我希望派遣一位司铎与他们一同前来，以防路上有疾病发生，可以给那些看来将要死去的人施洗。因此，愿你爱德如此行事，不要耽误时间，勤勉完成这些事。\par

\SourceRecord{Translated by James Barmby.；Nicene and Post-Nicene Fathers, Second Series；Vol. 12.；Edited by Philip Schaff and Henry Wace.；Buffalo, NY: Christian Literature Publishing Co.,；1895.；<http://www.newadvent.org/fathers/360206007.htm>.}

% structure: p0001 source=h1 target=section reason=page-title

\section{第六卷·第八封信}

\Emph{致\Place{埃皮鲁斯}诸位主教。}\par

\Person{额我略}致\Place{埃皮鲁斯}诸位主教：\Person{德奥多禄}、\Person{德默特琉}、\Person{斐理伯}、\Person{则诺}及\Person{阿尔基松}。\par

最亲爱的弟兄们，你们的来信告知我们，由于天主的恩宠，我们的兄弟安德肋已庄严地被祝圣为\Place{尼科波利斯}城的主教。既然你们表明，他的祝圣是在圣职人员和省内人士的赞同下进行的，我们感到喜乐；并祈求你们为他所见证的美善能常在他内，并藉天主恩宠的合作而日益增长，因为牧者的良善便是其属下的得救。你们的责任，便是殷勤地效法那在你们称颂中所显示为令你们喜悦的在他身上的一切。因为任何人若不愿效法那令他喜悦的美善，在人间便是过错，在天主面前便是可罚的。因此，愿你们的服从为你们的见证增添信誉。愿无人违抗他为了教会的共同利益、在保持正直的前提下所命令的。愿你们每人乐于展现他的热忱，以便在你们中间存有天主所喜悦的、持久的司祭和谐时，没有恶意能使你们脱离彼此爱德的束缚，也没有分歧能扰乱你们。因为狡猾的仇敌将无法进入你们的心，因为他知道，在真诚的爱德占有一席之地的地方，他丝毫不能被接纳或收留。\par

此外，最亲爱的弟兄们，请当心，并要对托付给你们的羊群付出你们所承担的、也是你们所亏欠的警醒；藉着警醒和祈祷，抵抗仇敌的诡计。以无玷的信德，将你们所治理的人民交托给我们的天主，好使你们的司祭职分在永恒的审判者面前，为你们带来不应是惩罚，而是冠冕。\par

你们要知道，我们已给上述我们的兄弟和同为主教的安德肋送去了一条\OriginalTerm{披肩}{pallium}，并授予他所有我们的前任授予他前任的特权。\par

此外，我们听说在你们地区，神圣的圣秩是因着给予报酬而授予的。如果真是这样，我含泪说道，我叹息着宣告等等。\EditorialNote{参见第五卷·第五十三封信，至\Quote{成为异端者}处。}因此，我劝勉并恳求你们务必留心这一点：任何给予的报酬、任何情面、任何人的恳求，都不得对神圣的圣秩提出要求，而应提升那品行端正、生活堪受称赞的人担任此职。因为倘若我们，正如我们不相信会发生的那样，察觉到任何此类事情，我们将按合宜的方式，以教规的严厉加以纠正。愿全能的天主，祂以智慧之大能奇妙地安排万有，并守护祂所安排的一切，赐予你们意愿和践行祂所命令的能力。\par

\SourceRecord{Translated by James Barmby.；Nicene and Post-Nicene Fathers, Second Series；Vol. 12.；Edited by Philip Schaff and Henry Wace.；Buffalo, NY: Christian Literature Publishing Co.,；1895.；<http://www.newadvent.org/fathers/360206008.htm>.}

% structure: p0001 source=h1 target=section reason=page-title

\section{第六卷，第九封信}

致\Person{多努斯}主教。\par

\Person{格列高利}致\Place{默西拿}（\Emph{Messene}）主教\Person{多努斯}。\par

宗座之仁爱及古代惯例之秩序，促使我等认为适宜赐予尔\OriginalTerm{披带}{pallium}，尔已知在\Place{默西拿}教会承担治理之职；即在尔前任使用披带之时及其方式，我等无争议；同时告诫尔，正如尔因从我等获得此类装饰以荣耀司铎职务而欢喜，亦当努力以品行与事工之正直，在\Person{基督}的光荣中装饰尔在我等权下所承担的职务。如此，若此身体之服饰与尔灵魂之各种美德相符，尔将因相互呼应的装饰而显赫。凡已知古时授予尔教会之一切特权，我等以权柄予以确认，并裁定其应继续有效，不受侵犯。\par

\SourceRecord{Translated by James Barmby.；Nicene and Post-Nicene Fathers, Second Series；Vol. 12.；Edited by Philip Schaff and Henry Wace.；Buffalo, NY: Christian Literature Publishing Co.,；1895.；<http://www.newadvent.org/fathers/360206009.htm>.}

% structure: p0001 source=h1 target=section reason=page-title

\section{第六卷，第十二封信}

\Emph{致\Person{蒙塔纳}与\Person{托马斯}。}\par

\Person{格列高利}致\Person{蒙塔纳}等。\par

由于我们的救赎者，万有的造物主，为了这个目的屈尊就卑，取了人性，即藉祂天主性的恩宠，打破那束缚我们的奴役锁链，将我们恢复至原初的自由；因此，若有人将本性原本生而自由、却因万民法而受制于奴役之轭的人，藉获释的恩惠，恢复他们生而享有的自由，这便是一项有益的行为。因此，受仁爱及对本案例的顾念所感动，我们使你们，\Person{蒙塔纳}与\Person{托马斯}，圣罗马教会（我们赖天主助佑所事奉的教会）的仆人，从今日起成为自由人及\Place{罗马}公民，并将你们的所有私产全部交还给你们。\par

你，\Person{蒙塔纳}，既声明立志度隐修生活，我等今日便赠予并赐给你两\OriginalTerm{翁西亚}{unciæ}，此乃司铎\Person{高迪奥索}依其遗嘱安排而遗赠予你的，但有一条件，即一切均须完全归于\Place{圣老楞佐隐修院}的利益，该院由女院长\Person{康斯坦蒂娜}管理，而你也蒙天主的仁慈，即将在那里宣发圣愿。然而，若发现你以任何方式隐瞒了上述\Person{高迪奥索}所遗留财产的任何部分，则毫无疑问，全部财产必须转移至我教会的名下。\par

再者，对于你，上述的\Person{托马斯}，我等为增进你的自由，也愿你为书记服务，我等亦于今日依此释奴文书赐予并赠给你五\OriginalTerm{翁西亚}{unciæ}，此乃前述司铎\Person{高迪奥索}依其遗嘱以继承权形式留给你的，连同他曾赐予你母亲的嫁妆；但须附以下法律及条件：若你去世时无合法子女（即合法婚姻所出之子女），则我等所赐予你的一切，须完整无损地归还至圣罗马教会名下。但若你有婚生子女（如上所述，且为法律所承认），且你离世时他们仍健在，则我等指定你为同一财产的无条件主人，并赐你全权就此财产订立遗嘱。因此，凡我等藉此释奴证书所指定并赐予之事，你须知我等及我等继任者必毫无异议地遵守。因为正义与理性的法则提示，凡渴望其命令为继任者所遵守之人，自己无疑也应遵守前任的旨意与规定。我等已口述此释奴证书由书记\Person{帕泰里乌斯}书写，并与三位首席司铎及三位执事一同亲手签署，以臻完全稳妥，并已交付于你。\par

订于\Place{罗马}城。\par

\SourceRecord{Translated by James Barmby.；Nicene and Post-Nicene Fathers, Second Series；Vol. 12.；Edited by Philip Schaff and Henry Wace.；Buffalo, NY: Christian Literature Publishing Co.,；1895.；<http://www.newadvent.org/fathers/360206012.htm>.}

% structure: p0001 source=h1 target=section reason=page-title

\section{第六卷 第十四封书信}

\Emph{致\Person{纳尔塞斯}伯爵}。\par

\Person{额我略}致\Person{纳尔塞斯}等。\par

你的爱德，渴望了解我们的意见，不辞辛劳写信来询问我们对那本反对司铎\Person{阿塔纳修斯}的书有何看法，这本书已经寄给了我们。在仔细阅读了其中一些部分后，我们发现他陷入了摩尼教的错误。但是，那个用标记指出某些地方为异端的人，他自己也陷入了白拉奇异端；因为他把某些表达得完全符合公教信仰且完全正统的观点标为异端。因为当这样写着：亚当犯罪时，他的灵魂死了，作者后来表明它是如何死的，即它失去了其状况的幸福。谁否认这一点，就不是公教徒。因为天主曾说过：\BibleQuote{在你们吃的那一天，你们必定要死}\BibleRef{创 2:17}。因此，当亚当吃了禁树的果子时，我们知道他并没有在肉体上死去，因为此后他还生了孩子，并活了许多年。那么，如果他在灵魂上没有死，就会得出不敬的结论：预言他犯罪那天必死的那一位自己说了谎。但必须理解死亡有两种方式：要么是停止生活，要么是关乎生活的方式。那么，当说人的灵魂在吃禁物时死了，意思不是停止生活，而是关乎生活方式——他以后要生活在痛苦中，而他原本被造是要幸福快乐地生活。因此，那个把我兄弟\Person{若望}主教寄给我的书中这段标记为异端的人是白拉奇派；因为他的看法显然就是\Person{白拉奇}的看法，宗徒\Person{保禄}在书信中明确驳斥了这种看法。我不需要引用他书信中的具体段落，因为我写信给一个知道的人。但是，在厄弗所大公会议中被谴责的\Person{白拉奇}坚持这一看法，旨在表明我们被基督救赎是不真实的。因为，如果我们没有通过亚当在灵魂上死亡，我们就不是真正被救赎的，这样说是不敬的。此外，在查阅了厄弗所大公会议的法案后，我们丝毫没有发现关于\Person{阿德尔斐乌斯}、\Person{萨瓦}以及据说在那里被谴责的其他人的内容，我们认为，正如加采东大公会议在某处被君士坦丁堡教会篡改一样，对厄弗所大公会议也可能做了类似的事。因此，请你的爱德彻底搜寻这部大公会议法案的古老抄本，从而查看是否能找到任何类似内容，并将你找到的抄本寄给我，我一读完就会归还。因为新的抄本不完全可靠；正是由于这个原因，我一直犹豫不决，至今还不愿就此事回复我上述的兄弟\Person{若望}主教。此外，罗马的抄本比希腊的抄本正确得多，因为，正如我们没有你们的聪明才智，我们也没有任何欺骗。\par

现在关于司铎\Person{若望}，要知道他的案子已在会议中裁决，我由此清楚地确定，他的对手们希望并长期努力想把他定为异端，但完全失败了。\par

以我的名义问候你的朋友们，他们也是我们的朋友；我们中属于你的朋友们，也通过我衷心问候你。愿全能的天主以祂的手在如此众多的荆棘中保护你，使你毫发无损地采摘那些主所选的花朵。\par

\SourceRecord{Translated by James Barmby.；Nicene and Post-Nicene Fathers, Second Series；Vol. 12.；Edited by Philip Schaff and Henry Wace.；Buffalo, NY: Christian Literature Publishing Co.,；1895.；<http://www.newadvent.org/fathers/360206014.htm>.}

% structure: p0001 source=h1 target=section reason=page-title

\section{第六卷，第十五封信}

\Emph{致若望主教。}\par

额我略致君士坦丁堡主教若望。\par

正如异端的邪恶必须以正确信德的热情加以抑制，同样，天主教会信德的纯正也应受到接纳。因为，若有人自称信仰纯正却被轻视，所有人的信仰便陷入怀疑，而因轻率的严苛就会产生致命的错误。因此，不仅迷途的羊不被召回主的羊栈，就连那些在羊栈内的羊，也暴露在野兽的利齿下被残忍撕裂。最亲爱的弟兄，让我们充分思考这一点，绝不容许任何真正宣认天主教会信德的人，在异端借口下受困扰，也不允许（愿天主阻止）以纠正异端为名，反而使异端愈加滋生。\par

然而我们甚感惊奇：你们委派为信仰之事审判加采东教会司铎若望的法官，竟听信传言，不顾真理，当他清楚宣认信仰时仍不相信；尤其当他的控告者被问及他们所说的马尔基翁派异端是什么，并以此试图证明他有罪时，他们明白承认自己不知道。由此明显可见，他们不顾天主，不公正，且损害自己的灵魂，仅仅出于一己私意想要伤害他本人。因此，我们根据已经举行的会议（如面前记录所载），彻底审查并考虑了所有必要事项。由于我们未能发现上述司铎在任何方面有罪，尤其因为他向你们委派的法官提交的辩护词完全符合正确信德的纯正，故此，我们不赞同上述法官的判决，借着我们的天主和救主基督启示的恩宠，以我们的确定判决宣布他为天主教徒，免于一切异端指控。既然我们已经将他送回你们的圣座，你们应以一向的仁慈接纳他，向他施以司铎的爱德，保护他免受一切骚扰，不许任何人给他制造麻烦；正如你们保护他人免受压迫，同样不应拒绝向他提供援助。\par

\SourceRecord{Translated by James Barmby.；Nicene and Post-Nicene Fathers, Second Series；Vol. 12.；Edited by Philip Schaff and Henry Wace.；Buffalo, NY: Christian Literature Publishing Co.,；1895.；<http://www.newadvent.org/fathers/360206015.htm>.}

% structure: p0001 source=h1 target=section reason=page-title

\section{第六卷，第十六封书信}

\Emph{致奥古斯都·\Person{玛里齐乌斯}。}\par

\Person{格里高利}致\Person{玛里齐乌斯}等。\par

既因在您，最基督化的君王，不腐的正信光辉如从天降下的光柱照耀，且众所周知，陛下热切拥抱并全心热爱那纯正的宣认，因天主眷顾您在此宣认上有能力，我们感到十分必要为那些同被同一信仰光照的人恳求，以使君主们的虔诚能以恩宠保护他们，并护卫他们免受一切骚扰。当某些人轻视此类人的信仰宣认时，他们便显出是抵触真信仰。因为，既然宗徒宣告说口宣认是为得救恩，那不肯相信正确宣认的人，就在拒绝他人时控告了自己\BibleRef{罗 10:10}。\par

如今，关于对\Place{加采东}教会司铎\Person{若望}的一切诉讼，已在公会议上依次宣读并审议后，我们发现他遭受了更大的不公，因为当他声明并表明自己是一位天主教徒时，压垮他的不是他的罪过，而是一个长期悬而未决的模糊指控；以至于他的控告者在公开答复中声称，他们不知道他们所提及的\OriginalTerm{马西昂派}{Marcionists}的异端。而且，虽然他们在审判一开始就应当被驳回，但他们却得以作为模糊不清的指控者留在法庭上控告他。但是，以免传闻可能伤害他，他提交了一份书面的信仰宣认，目的是清楚地表明他是正信的宣认者和追随者。然而，由我们最神圣的兄弟和同为主教者\Person{若望}委任的审判官们，不公且无理地忽略了这一点；因此，当他们竭尽全力打压他时，反而显得比他更有罪。因为无人怀疑，对忠信者不信实即是不忠。既然一切已彻底查究并考虑，神圣公会议与我一起，藉天主德能的启示恩宠，已宣布上述司铎\Person{若望}为天主教徒，且在他身上未发现任何异端邪恶的污点，我恳求陛下的虔诚保护，下令保护他免受一切烦扰，不容许一位天主教信仰的宣认者遭受任何骚扰。因为不相信真诚宣认的人，并非铲除异端，而是制造异端。若容忍此事，将产生不信的机会，人们自己将陷入他们不慎想要纠正的罪责。\par

因此，愿最安宁的君王以虔诚的谨慎考虑这些事，并且，如我已恳求，我再次以恳切请求恳求，不要允许一个无辜的人好像有罪者一样再度受难；以使全能天主，祂见陛下热爱并捍卫天主教正直的纯洁，使您在征服敌人后，统治一个和平的共和国，并与祂的圣人们一同在永生中为王。\par

\SourceRecord{Translated by James Barmby.；Nicene and Post-Nicene Fathers, Second Series；Vol. 12.；Edited by Philip Schaff and Henry Wace.；Buffalo, NY: Christian Literature Publishing Co.,；1895.；<http://www.newadvent.org/fathers/360206016.htm>.}

% structure: p0001 source=h1 target=section reason=page-title

\section{第六卷 第十七封信}

\Emph{致\Person{德奥蒂斯图斯}。}\par

\Person{格列高利}致\Person{德奥蒂斯图斯}，皇帝亲属。\par

我们知道，阁下您的基督信仰始终致力于善工，因此我们为您提供收获赏报的机会，您必定乐于见到，以便通过这样的提供，我们也能分享您的功德。\par

因此我们通知您，持此信的\Person{若望}司铎已从那些控告他的人中被宣告无罪。因为应他的请求，我们召开了会议，并对其信仰进行了仔细的审查，发现他并无任何错误宣认。而且，由于他藉天主的仁慈，似乎是正确信仰的宣认者和追随者，我们以明确的判决宣告他无罪；尤其他的控告者声称，他们并不知道他们所谈论的\OriginalTerm{马尔基翁派}{Marcionists}异端是什么。为此，我们以慈父般的情感问候您，请求您以恩惠眷顾保护他。同时，唯恐日后有人再无故骚扰他，或以任何方式在此事中令他烦恼，愿阁下您的庇护如此为他辩护和捍卫——并鉴于您自己的赏报，更加急切地——使他不再遭受任何不义的折磨，并使您以真诚宣认所敬拜的人类造物主与救赎主，在您的诸多善工中，回报您在此事上的行动。十月。第十四小纪。\par

\SourceRecord{Translated by James Barmby.；Nicene and Post-Nicene Fathers, Second Series；Vol. 12.；Edited by Philip Schaff and Henry Wace.；Buffalo, NY: Christian Literature Publishing Co.,；1895.；<http://www.newadvent.org/fathers/360206017.htm>.}

% structure: p0001 source=h1 target=section reason=page-title

\section{第六卷·第十八封信}

\Emph{致若望主教}\par

\Person{额我略}致书于\Place{叙拉古}主教\Person{若望}。\par

受宗座之仁爱及古老惯例之命，我们以为宜授予你的贤弟——已知你接受了治理叙拉古教会之职责——使用\OriginalTerm{披带}{pallium}之权利；即，在你确知你的前任无疑曾使用的时间和方式；然而劝诫你，正如你因从我们接受了此装饰以光荣你的司铎职务而喜乐，同样也要以品德和行为的正直竭力装饰你所接受的职务，以归于我们在基督内的光荣。因为如此，你将以相互呼应的装饰而显耀，如果身体的服饰与心灵的卓越也相和谐。\par

对于所有已知先前授予你教会的特权，我们以我们的权威予以确认，并下令它们应保持不可侵犯。\par

\SourceRecord{Translated by James Barmby.；Nicene and Post-Nicene Fathers, Second Series；Vol. 12.；Edited by Philip Schaff and Henry Wace.；Buffalo, NY: Christian Literature Publishing Co.,；1895.；<http://www.newadvent.org/fathers/360206018.htm>.}

% structure: p0001 source=h1 target=section reason=page-title

\section{第六卷，第二十二封信}

\Emph{致伯多禄主教。}\par

\Person{额我略}致\Place{科西嘉岛}\Place{阿莱里亚}主教\Person{伯多禄}。\par

既然在\Place{科西嘉岛}，在名为\Place{尼杰乌努姆}的地方，在那属于我们依天主的眷顾所服务的圣罗马教会、称为\Place{切拉斯·库皮亚斯}的产业上，我们已下令建造一座大殿，附设一个洗礼堂，以光荣真福宗徒之长\Person{伯多禄}和殉道者\Person{老楞佐}；因此，我们在此敦促您的兄弟情谊立即前往上述地点，并遵守可敬的奉献礼的隆重仪式，庄严地祝圣上述教堂和洗礼堂。还要虔诚地安放您已领受的圣髑（\OriginalTerm{圣髑}{sanctuaria}）。\par

\SourceRecord{Translated by James Barmby.；Nicene and Post-Nicene Fathers, Second Series；Vol. 12.；Edited by Philip Schaff and Henry Wace.；Buffalo, NY: Christian Literature Publishing Co.,；1895.；<http://www.newadvent.org/fathers/360206022.htm>.}

% structure: p0001 source=h1 target=section reason=page-title

\section{第六卷，第二十四封信}

\Emph{致\Person{马里尼亚努斯}主教。}\par

\Person{格里高利}致\Place{拉文纳}主教\Person{马里尼亚努斯}。\par

我们已通过执事维吉利乌斯收到您弟兄的来信，信中您告知我们，某些神职人员和民众喧嚷说，在贵教会与院长克劳狄之间的案件在此审理和裁决是违反法律和教规的。但是，如果他们留意教会秩序以及案件所涉当事人的身份，他们本应完全避免无谓的抱怨；尤其是因为该案件无法在那里进行辩护，因为上述院长在那里抱怨曾受到您前任的不公正对待，且仍遭受此不公。或许，如果他没有向更高权威上诉并寻求让其案件的权利在此得到裁决，这一异议或许可以成立。然而，您自己难道不知道，司铎约翰针对我们的弟兄和同袍主教君士坦丁堡的约翰所提起的案件，曾提交至宗座，并由我们裁决？那么，如果连来自君王所在之城的案件都能由我们审理，您二人之间的事务岂不更应在查明真相后在此了结？至于您，切勿被那里愚人之言所动，切勿相信通过我们会对贵教会造成任何损害。因为，如果您询问您的执事、天主的仆人塞昆迪努斯，以及我们的公证人卡斯托里乌斯，您将从他们那里了解到，您前任曾多么希望解决此案。但您弟兄明智地派人前来处理此事，并未听从虚言。现在我们信赖全能的天主，愿此案能以悦乐天主的方式了结，使抱怨再无余地，双方均不受不公之苦。您那里的前执事和\OriginalTerm{保护人}{defensor}、我们最亲爱的儿子彼得所留下的剑，已交付您前任，请通过天主的仆人塞昆迪努斯和公证人卡斯托里乌斯（持信者）送还我们。\par

\SourceRecord{Translated by James Barmby.；Nicene and Post-Nicene Fathers, Second Series；Vol. 12.；Edited by Philip Schaff and Henry Wace.；Buffalo, NY: Christian Literature Publishing Co.,；1895.；<http://www.newadvent.org/fathers/360206024.htm>.}

% structure: p0001 source=h1 target=section reason=page-title

\section{第六卷，第二十五封信}

\Emph{致萨洛纳的马克西穆斯。}\par

额我略致马克西穆斯，萨洛纳教会的闯入者。\par

当你以各种借口拖延服从我们的书信，当你经多次劝告仍一再推迟前来我们这里查明真相时，你反而使加于你的指控更显可信；而且，即使没有其他对你不利且有害的事，单是你的拖延便已使你成为有罪并遭谴责。最终请你谦卑下来，顺从于服从，并立即前来我们这里，不可有任何借口，好使真相在天主的敬畏中得到调查和确定，从而本着公平和教会法典的规定作出决定。因为你当确信，我们将以正义和法典的规条对待你，并藉着真理的创造者天主的启示，公正地了结你的案件。至于你要求我们派人到你所在的城里，以便在你面前证明所控之事，这在某种程度上情有可原，倘若理性曾要求被告承担举证责任的话。然而，既然这责任不在你身上，而在你的控告者身上，你就不要犹豫，按我们先前所说的，不再拖延地前来我们这里；你的控告者将立即到场，以适当的证据支持关于\OriginalTerm{买卖圣职的异端}{simoniacal heresy}或其他事项的指控；否则，就此事得到妥善解决而言，藉着宗徒之长伯多禄的干预，将有一项公正的处理发生；如此，在我们知晓这些事情之后，不会因任何纵容而在天主面前有罪疚扰乱我们。至于你声称我们最安详的主上们曾命令在你的城里审理此事，我们确实未曾就此收到他们其他的命令，只知你当前来我们这里。但是，即使他们因忙于诸多谋划和忧思，为了他们藉天主恩赐被授予的共和国的福祉，此事曾被提出，且一项命令被悄然获得，然而，既然我们和众人都知道我们最虔诚的主上们喜爱纪律，尊重等级，敬重教会法典，且不干预司铎的案件，我们仍将立即执行有利于他们灵魂及共和国之事，并且我们因着对可怕而可畏的审判的敬畏而不得不这样做。\par

因此，停止所有借口，勿延迟在此出现，好使我们藉真相的调查得以最终了结你的案件。但是，我们获悉，你极其恐惧和战栗，生怕我们因你未经我们同意便不规则地强行进入司铎秩位这一众所周知的事实而报复你——这确是一项不可容忍的不端行为：然而，根据我们最安详的皇帝主上的命令，我们宽恕你这一点，只要你绝不继续坚持你\OriginalTerm{顽固抗命}{contumacy}的错误；我们绝不会因此事而对你动怒。但其他报告给我们的情况，我们不能不加查问而听之任之。\par

既然我们早就致信警告你，在我们查明上述最安详的主上的旨意之前，你绝不可胆敢举行\OriginalTerm{弥撒圣祭}{solemnities of mass}，而你却狡猾地使此信未能落入你手，尽管你仍以某种方式知晓其内容，却拒绝遵从——因此，我们确认先前以书面方式发给你的命令，即在你所受的全部指控得到彻底查究和审查之前，你不得胆敢举行弥撒圣祭。如果你仍以悖逆的胆大妄为擅自举行，当知你并未脱离先前断绝共融的威胁。因为，即使没有其他过犯，单凭这\OriginalTerm{骄傲之罪}{sin of pride}，我们就剥夺你与\OriginalTerm{主的体血}{body and blood of the Lord}的共融。因此，显出你应有的服从，按我们所说，赶紧竭尽全力前来我们这里；但给予你三十天的时间准备行程；如此，放下一切借口，不要延迟在此出现。\par

此外，若有来自法官、军队或民众的任何阻碍你行程的情况发生，我们明白事情是如何巧妙安排的。那么，你自己当思量，你将如何就这一义务向这里的人们或全能的天主在后来的审判中交代，因为你已因蔑视而招致对你不利的严厉判决。\par

再者，我获悉我的\OriginalTerm{兄弟与同侪主教}{brother and fellow bishop}保利诺，以及萨洛纳教会的\OriginalTerm{总执事}{archdeacon}奥诺拉图斯，因拒绝赞同你的擅权而遭受你严重的骚扰，以致被迫提供担保，使他们不得自由离开城市或自己的房屋。若情况属实，你在收到此信后，终于——尽管为时已晚——恢复理智，停止骚扰他们中任何一人，好使他们有自由，或若愿意可前来我这里，或可往他处寻求益处。\par

\SourceRecord{Translated by James Barmby.；Nicene and Post-Nicene Fathers, Second Series；Vol. 12.；Edited by Philip Schaff and Henry Wace.；Buffalo, NY: Christian Literature Publishing Co.,；1895.；<http://www.newadvent.org/fathers/360206025.htm>.}

% structure: p0001 source=h1 target=section reason=page-title

\section{\WorkTitle{第六卷，第二十六封信}}

\Emph{致萨洛纳人。}\par

\Person{额我略}致寄居\Place{萨洛纳}的至爱之子、神职人员和贵族们。\par

我听闻，有些心术不正的人，为毒害你们的心灵（亲爱的），竟试图向你们暗示：我是出于对\Person{玛克西穆}的某种怨恨而行事，我所渴望的不是依规而行，而是愤怒所驱。然而，司铎之心在任何事上绝不应受私情左右，万不可如此！相反，亲爱的，我是为你们着想，并因畏惧全能天主对我自己灵魂的审判，才希望彻底查清这位\Person{玛克西穆}本人是否背负着阻碍圣秩的罪行，是否试图通过\OriginalTerm{西莫尼异端}{simoniacal heresy}——即向某些选举人赠送贿赂——来谋取司铎之职。倘若他未受自身罪孽束缚，他才能在主前自由地作你们的代祷者。\par

然而，他的骄傲之罪已显明：他奉召前来见我，却以种种借口抗拒，逃避前来，害怕前来。那么，倘若他的良心并未就所受指控之事控告他，他又怕什么呢？看，亲爱的，你们已久无牧者，愿全能天主使你们明白，我是多么诚恳、发自内心地同情你们的匮乏。因为我听说主羊群遭受了何等严重的蹂躏。但若没有牧人，谁能防范豺狼呢？因此，请催促上述\Person{玛克西穆}前来我们这里，以便我们若发现他无辜，便可确认他；但若关于他的传闻属实，那么，亲爱的，愿你们不再因他本人的阻挠而继续匮乏。\par

至于我，请相信，我并非出于任何怨恨或私情的敌意而反对他；凡合于规诫与正义之事，我将在天主的助佑下决定。\par

但我极为惊讶的是，在\Place{萨洛纳}教会如此众多的神职人员和信众中，几乎找不到两位圣秩人员——即我们的兄弟和同为主教的\Person{保利诺}，以及我最亲爱的儿子、该教会的总执事\Person{奥诺拉托}——他们拒绝与篡夺司铎职位的\Person{玛克西穆}共融，并记得自己是基督徒。\par

至亲爱的儿子们，你们本应思量自己的圣秩，视被宗座所弃绝者为弃绝，使他先能被洗净（若可能的话），除去对他的指控，然后你们的爱德才能与他共融，而不分担他的罪责。然而，我们以仁爱之心关怀你们的爱德；既然我们得知你们中有人受到压力，被迫接纳他并与他共融，我们恳求全能天主，赦免你们自己一切罪过的罪责，以及你们在他人罪责中所有牵连，赐予你们今生祂保护的恩宠，并使我们得以在永恒家乡为你们喜乐。\par

\SourceRecord{Translated by James Barmby.；Nicene and Post-Nicene Fathers, Second Series；Vol. 12.；Edited by Philip Schaff and Henry Wace.；Buffalo, NY: Christian Literature Publishing Co.,；1895.；<http://www.newadvent.org/fathers/360206026.htm>.}

% structure: p0001 source=h1 target=section reason=page-title

\section{第六卷·第二十七封信}

\Emph{致雅德拉的圣职人员和民众}。\par

\Person{额我略}致居住在\Place{雅德拉}的、与背信者\Person{马克西默斯}共融的司铎、执事、圣职人员、贵族及民众。\par

本座得知，你们中有人因无知受骗或受胁迫，与那些因过错（如你们所知）被宗座剥夺共融的人共融；另有人以健全的分辨力，在主护下持守不参与。我既为那些持守者欣喜，亦为那些迷途者哀痛，因为他们领受了圣体奥迹，而这奥迹本是天主慈爱赐予我们为赦罪之用的，反倒成了他们灵魂的损害。并且，因我（恳求全能天主彰显于你们）真心诚意地同情你们，我以慈父之情恳请并劝诫你们：你们每人当戒绝不法共融，完全避开宗座不接受共融之人，以免有人因可能得救之事，在永恒审判者面前成为有罪。\par

此外，我发现在你们地区有些心术不正的人，试图暗示我因怨恨而反对\Person{马克西默斯}，且我不愿按教会法规行事，而是按愤怒行事。但司铎之心绝不可被私怨所动。至于我，正是为那些地区的民众和我的灵魂着想，并畏惧全能天主的审判，我愿调查此\Person{马克西默斯}之事，并求天主指引，依教会法规裁决。既然我多次写信给他，命他在我获知其案件详情前不得举行弥撒圣祭，他无论如何已被剥夺共融；如今他的骄傲之罪昭然若揭：因他（如我所言）屡次被规劝前来，却以各种借口推辞、回避、害怕前来。若他良心不指证所言之过，他又何所惧？既然你们已知这些事，不能再以无知为借口，我恳求、劝诫、警告你们：完全禁止与禁止的共融有份，且不得有一人，违背自己的灵魂，与任何与上述\Person{马克西默斯}共融的司铎共融。\par

然而我在前已说过，听说你们中有人无知而跌倒，甚至有人被迫共融。我恳求全能的主，以祂永远的护佑保守那些未认同此不法行为的人，并以祂所愿的恩惠应允他们；至于那些被党派精神、无知或任何其他原因引入过错的人，求祂赦免他们一切罪恶的过犯，免除他们因他人牵连的责任，并在现世赐予他们祂护佑的一切恩宠，且使我得以在永恒故乡为他们喜乐。因此，为使这代祷对我们在天主救主前有益，请为你们的灵魂听从我们的劝诫，并从那些你们所知已戒绝且仍戒绝与上述\Person{马克西默斯}共融的人手中，领受圣体。\par

\SourceRecord{Translated by James Barmby.；Nicene and Post-Nicene Fathers, Second Series；Vol. 12.；Edited by Philip Schaff and Henry Wace.；Buffalo, NY: Christian Literature Publishing Co.,；1895.；<http://www.newadvent.org/fathers/360206027.htm>.}

% structure: p0001 source=h1 target=section reason=page-title

\section{第六卷 第二十九封}

\Emph{致\Person{马里尼亚诺斯}主教。}\par

\Person{额我略}致\Place{拉文纳}主教\Person{马里尼亚诺斯}\par

我们感到惊奇，为何阁下的明辨在短时间内发生如此大的变化，以至于不反思自己所求之事。为此我们深感忧虑，因为你提供了明显的证据，表明邪恶谋士的话在你身上所起的作用，超过了神圣学问对你的裨益。当你本应保护隐修院，并全力聚集其中的修道者，从灵魂的聚集中获得益处时，你却相反，如你的信函所证，渴望操练压迫它们；更糟的是，你试图使我们成为你过失的共犯：即希望得到我们同意，以管理其财产和事务为名，压迫你的前任所创立的隐修院。\par

因为你应当回想，在你面前，以及在你的几位司铎、执事和圣职人员的面前，我们曾应他们请求，颁布了一道与你前任遗嘱相悖的谕令。然而，虽然他针对该隐修院所做的处置在其中仍得到确认，你如今却掩饰这一点，并要求我们下令相反的事。我们确实知道这计谋并非出于你本人；但当你拒绝不听那些说乖谬之言的人时，你不仅损害自己的声誉，也损害灵魂。因此，既然我非常爱你，我急切地劝诫你——请仔细考虑此事——不要关心钱财胜过关心灵魂。前者应当附带地看待；但后者应当倾注全部心智，并奋力追求。在这事上辛勤付出你的劳苦和关切吧，因为我们的救赎主从司铎职分中所寻求的不是金子，而是灵魂。\par

此外，我们听闻，在你阁下治下的隐修院受到圣职人员的缠扰和各种烦扰。但愿此事不再发生，你要以严格的禁令加以制止，好使居住其中的隐修士能自由地在我们天主的赞美中欢欣。\par

关于圣职人员\Person{罗马努斯}和\Person{多米尼库斯}，他们鲁莽大胆，未经我们祝福就擅自离开此城，本应受到更严厉的惩罚，然而应当以仁慈之心给予宽免，以劝勉他们返回自己的职责。四月，小纪第十四。\par

\SourceRecord{Translated by James Barmby.；Nicene and Post-Nicene Fathers, Second Series；Vol. 12.；Edited by Philip Schaff and Henry Wace.；Buffalo, NY: Christian Literature Publishing Co.,；1895.；<http://www.newadvent.org/fathers/360206029.htm>.}

% structure: p0001 source=h1 target=section reason=page-title

\section{\WorkTitle{第六卷，第三十封信}}

\Emph{致\Person{塞孔都斯}。}\par

\Person{额我略}致\Place{拉文纳}的天主仆人\Person{塞孔都斯}。\par

如今\Person{卡斯特里乌斯}已返回，并将您与\Person{阿吉卢尔夫}王之间所行的一切都告知了我们。我们已留意尽快将他送回您那里，以免有人以迟延为由对我们提出非难。因此，既从他得知一切应行之事，请您对此事倾注热忱，全力以赴促成这项和约的议定，因为据传闻，有些人正试图阻挠它。为此，请迅速采取行动，以免您的辛劳落空。因为这部分地区和各个岛屿已陷入极大的危险。\par

请用您所能使用的言词激励我们的兄弟\Person{马里尼亚努斯}主教；因为我怀疑他沉睡了。因为某些人来到我这里，其中有几位年老的乞丐，我问他们收到了什么、从谁那里收到的；他们一一告诉我他们在旅途中得到了多少，以及是谁给的。但是，当我询问他们我的上述兄弟给了他们什么时，他们回答说他们曾向他请求，但什么也没从他那里得到；以至于他们连路上的面包也没有得到，而那座教会向来有施舍给所有人的习惯。因为他们说，他回答说：\Quote{我没有什么可以给你们。}我感到惊讶，如果他有衣服、钱和仓库，却没有什么给穷人。\par

那么，告诉他，他应当随着他的职位也改变他的性情。让他不要以为仅仅阅读和祈祷就足够了，以致他想要置身事外，完全不肯亲手结出果实；而要让他有一双慷慨的手；让他救助那些贫乏的人；让他视他人的需要为己任；因为，如果没有这些，他不过是空负主教之名。我确实曾在信中给他一些关于他灵魂的劝诫；但他没有给我任何回音；由此我猜想他甚至不屑于阅读它们。为此，我现在无需再在他给我的信中对他进行任何劝诫；因而我只写了作为他在世俗事务上的顾问能够口授的内容。因为，我为一个人口授，而他却不肯读对他说的话，这并非我的责任。那么，请你的爱德私下与他谈论所有这些事，并劝诫他应当如何自持，以免因眼前的疏忽而失去他以前生活的益处——愿天主不许此事发生。\par

\SourceRecord{Translated by James Barmby.；Nicene and Post-Nicene Fathers, Second Series；Vol. 12.；Edited by Philip Schaff and Henry Wace.；Buffalo, NY: Christian Literature Publishing Co.,；1895.；<http://www.newadvent.org/fathers/360206030.htm>.}

% structure: p0001 source=h1 target=section reason=page-title

\section{第六卷·第三十二封信}

\Emph{致\Person{福图纳图斯}主教}\par

\Person{格列高利}致\Place{拿坡里}主教\Person{福图纳图斯}。\par

我们此前已致信贤弟：倘若有奴隶因天主的感召，愿意从犹太迷信皈依基督信仰，其主人不得将其出售；自他们表明意愿之时起，他们便完全有权获得自由。然而，据我们所知，他们（即犹太主人）既不体贴我们的意愿，也不考虑法律的规定，认为在外教奴隶的情况下他们不受此条件约束。因此，贤弟应关注此类情况，若其奴隶中有人——无论是犹太人或外教人——愿意成为基督徒，在其意愿公开表明之后，任何犹太人不得以任何手段或理由为掩饰有权将其出售；凡愿皈依基督信仰者，你应全力支持其获得自由之权利。不过，唯恐那些因此失去奴隶的人认为其利益受到了不合理的损害，你应审慎留意遵守如下规则：倘若外教人因贸易之故从外邦带来，偶然逃至教会，声称愿意成为基督徒，甚或在教会外宣布此意愿，那么，在三个月期限内（在此期间可为他们寻找买主出售），他们（即犹太主人）可以收取其价格，亦即从基督徒买主处收取。但若在上述三个月之后，此类奴隶中有人宣布其意愿并渴望成为基督徒，则此后任何人不得再购买他，其主人也不得以任何借口出售他；应让他毫无保留地获得自由之益；因为在此情况下，他（即主人）被视为是为自己服役而获得他，并非为了出售。因此，务请贤弟如此警醒地遵守这一切，使任何人的恳求或情面都不能引诱你偏离。\par

\SourceRecord{Translated by James Barmby.；Nicene and Post-Nicene Fathers, Second Series；Vol. 12.；Edited by Philip Schaff and Henry Wace.；Buffalo, NY: Christian Literature Publishing Co.,；1895.；<http://www.newadvent.org/fathers/360206032.htm>.}

% structure: p0001 source=h1 target=section reason=page-title

\section{第六卷·第三十四封书函}

\Emph{致\Person{Castorius}，公证人。}\par

\Person{Gregory}致\Person{Castorius}，我们于\Place{Ravenna}的公证人。\par

当\Place{Ravenna}教会的执事\Person{Florentinus}代表我们最可敬的兄弟兼同道主教\Person{Marinianus}就使用大披肩一事与我们商议时，我们问他古代惯例为何，他回答说，\Place{Ravenna}教会的主教在所有连祷中都使用\OriginalTerm{大披肩}{pallium}。但我们从他人那里得知事实并非如此，而且我们从前任主教\Person{若望}的书信中明显看出这点，我们把这些书信出示给他。但他只是说了奉命要说的话。因为，当这位\Person{若望}被你们禁止，不得擅自且冒失地使用大披肩时，他写信给我们说，古代的惯例是：该城主教应在隆重的连祷中使用大披肩。我们寄给你他的书信副本，以供参考。然而，当上述教会执事\Person{Adeodatus}在此地时，同样就大披肩的使用强烈地向我们施压，我们为了查明真相，同样让人问他惯例如何；他为了说服我们相信他，并成功地从我们这里获得他所求的，便宣誓作证说，他所在城市的主教在四次或五次隆重的连祷中使用大披肩是古代惯例。因此，阁下应勤勉地留心此事，并极其仔细地查问自古以来有多少隆重的连祷。也要注意，查问时称它们为大祷文，而非隆重的连祷；这样一来，当我们通过前述执事\Person{Adeodatus}向我们的作证以及前述主教\Person{若望}的书信所承认的，得知这些隆重的连祷有多少时，我们因知道在连祷中通常佩戴大披肩多少次，便会最乐意地授予这一特权。但不要向由教职人员推举的那些人进行此查问，而要向你认为是公正的其他人士；并且，你在仔细调查后所发现的任何情况，都要准确地传达给我们，这样，如我们所说，在查明真相后，我们便可减轻我们的兄弟兼同道主教、最可敬的\Person{Marinianus}的忧虑。\par

\SourceRecord{Translated by James Barmby.；Nicene and Post-Nicene Fathers, Second Series；Vol. 12.；Edited by Philip Schaff and Henry Wace.；Buffalo, NY: Christian Literature Publishing Co.,；1895.；<http://www.newadvent.org/fathers/360206034.htm>.}

% structure: p0001 source=h1 target=section reason=page-title

\section{第六卷 第三十五封书信}

\Emph{致副执事\Person{安特弥乌斯}}\par

\Person{额我略}致\Person{安特弥乌斯}，我们的\Place{那波利}副执事。\par

我们无法言喻，因\Place{坎帕尼亚}地区所发生之事，我们的悲伤何其之大，心中的痛苦何其之深；但您可从这场灾难的严重程度自行推断。关于此事，我们经由携带本函的显赫的\Person{斯德望}，派遣您的阁下携带金钱，以援助被掳的俘虏，并劝告您要全力关注此事，且要奋力完成；至于那些您知道没有足够财力赎回自身的自由人，您要赶紧赎回他们。但若有奴隶，而您发现他们的主人极其贫穷，无法出面赎回他们，您也毫不犹豫地将他们赎回。同样，您也要留意赎回因您的疏忽而失落的教会奴隶。此外，凡您所赎回的人，您务必努力制作一份名单，包含他们的姓名、每个人所居之地、所从事之事以及来自何处；这份名单当您前来时随身带来。再者，要尽快在此事上表现如此勤勉，以致那些将被赎回的人不会因您的疏忽而陷于危险，否则您日后在我们面前将极为可责。但尤其要努力做到这一点：如果可能，您能够以适中的价格赎回那些俘虏。但要将所有支出的总额清楚而准确地记录下来，并速速将这份书面记录呈交给我们。五月，小纪第十四。\par

\SourceRecord{Translated by James Barmby.；Nicene and Post-Nicene Fathers, Second Series；Vol. 12.；Edited by Philip Schaff and Henry Wace.；Buffalo, NY: Christian Literature Publishing Co.,；1895.；<http://www.newadvent.org/fathers/360206035.htm>.}

% structure: p0001 source=h1 target=section reason=page-title

\section{第六卷，第三十七封信}

\Emph{致哥伦布主教。}\par

额我略致努米底亚主教哥伦布。\par

我们于呈递此信的执事罗加提安手中收到阁下的信函，满溢司铎的甘饴。其友善言辞令我们甚为欣慰，尤其得知我们所长久期盼的阁下安康。然我们对尊座的热忱早已知悉，如今来信所表，亦与我们所持无异。至于阁下对我们的真诚，无需任何证明，因我们心中环绕的爱已足使我们知晓。我们已给予上述你所荐举的听者致西西里教会产业总管的手谕，命其敦促对方行正义之事，俾使一切拖延借口皆被搁置，争议之事得以速决。\par

今告知尊座，有一人前来，自称伯多禄，自称是主教，请求我们为其伸冤。起初其所言似值得怜悯，但经查问，我们发现其所述与事实大相径庭，其行为令我们极为困扰。然因相距甚远，我们无法详查其案由，故犹豫不决，未能裁决。如今，上述执事将返回你处，请求允此人同行，而此人亦自请前往你处。二人皆知尊座按分际对信德怀有热忱，且喜爱公义，故此提议为我们所接纳，并允其所请。既因你亲临现场，能更详尽查明案情，我们劝谕你对该伯多禄秉持公正与教会法规，务使正道之要求由你全然履行，并使其案件看来是依天主敬畏与教会规则而审判。但若有人涉嫌参与或同谋前述伯多禄所被指控之事，须作精确调查，查明真相后，依教会法同样判决。\par

此外，有一事令我们听闻，全然难以容忍，且与正信为敌：即天主教徒（说来可畏）及献身于宗教者（更甚）竟同意其子女、奴隶或受其管辖之人在多纳特派异端中受洗。此事若属实，请阁下竭尽全力纠正，以使信仰之纯洁经你悉心守护而屹立不摇，并使原可藉天主教洗礼得救的无辜灵魂，不致因异端染污而丧亡。凡上述人等中，若有允许其属下在多纳特派中受洗者，请你竭尽全力、以最迫切之态度召其回归天主教信仰。若任何人以任何借口在未来容忍此事发生于其属下，则当被完全断绝与圣职人员的共融。\par

\SourceRecord{Translated by James Barmby.；Nicene and Post-Nicene Fathers, Second Series；Vol. 12.；Edited by Philip Schaff and Henry Wace.；Buffalo, NY: Christian Literature Publishing Co.,；1895.；<http://www.newadvent.org/fathers/360206037.htm>.}

% structure: p0001 source=h1 target=section reason=page-title

\section{第六卷 第四十三封信}

\Emph{致\Person{梵南修}，贵族。}\par

\Person{额我略}致\Person{梵南修}，贵族，前隐修士。\par

你的来信使我们深感忧虑，因为我们得知你与我们的兄弟和同僚主教\Person{若望}之间产生了嫌隙，而原本我们期望能因你们彼此和睦而欢欣。因为，无论起因如何，愤怒本不应爆发到如此程度，以至于你的武装人员（如我们所闻）冲入主教宫，以敌对方式犯下各种恶行，而此事竟使你与他的慈父般的爱德分离。难道这场争执——无论它是什么——不能平静地解决，以便双方都不受损害，且良好关系不受干扰吗？如今，我们并非不知道我们上文提及的兄弟是何等庄重、何等圣洁、何等温和。由此我们推断，若非过度的烦恼逼迫他，他的兄弟情谊绝不会采取你所声称的使你感到委屈的措施。然而，我们一接到他的来信，便立即写信劝告他，要如从前一样接纳你的奉献，不仅允许在你家中举行弥撒，而且，如果你愿意，甚至由他亲自主礼，并且他应当在保持爱德的情况下进行诉讼。既然我们不愿任何人彼此不合或继续不合，我们已注意重申这一劝诫。因此，最亲爱的儿子，你应当如子女一般，向他表示对司铎应有的敬重，不要激怒他的心神。因为，如果你与你的司铎不合（愿天主禁止），你将与谁保持可靠的和睦呢？所以，放下骄傲的心，设法处理你们彼此之间的事务，使爱德不受侵犯，同时你们彼此的利益能和平地达成。\par

\SourceRecord{Translated by James Barmby.；Nicene and Post-Nicene Fathers, Second Series；Vol. 12.；Edited by Philip Schaff and Henry Wace.；Buffalo, NY: Christian Literature Publishing Co.,；1895.；<http://www.newadvent.org/fathers/360206043.htm>.}

% structure: p0001 source=h1 target=section reason=page-title

\section{第六卷，第四十四封}

致\Person{若望}主教\par

\Person{额我略}致\Place{锡拉库萨}主教\Person{若望}。\par

虽然可能有理由激起您手足之情的义怒，使您既不接受\Person{文南休斯}主的献仪，也不允许在他家中举行神圣的弥撒圣祭；然而，既然我们应按此方式处理现世的利益，使任何争执都不能使我们脱离爱德之纽，因此，我们正如先前所写信劝告的那样，劝勉您的圣座：您应以全部的温柔和悦乐天主的真诚接受上述那人的献仪，并允许在他家中举行弥撒奥迹；并且，正如我们所写的，倘若他愿意，您应亲自前往，借同他一起举行弥撒，以恢复你们旧日的情谊。因为您的职责是对子孙们施以司铎般的爱德，尽管如此，在可能发生的案件中，却绝不可因合理的原因而放弃您教会的管辖权。因此，考虑到这一点，有必要让您的手足之情以明智的温和在这些事情上如此行事，既妥善处理事情性质所要求的，又不远离父爱的爱德。\par

\SourceRecord{Translated by James Barmby.；Nicene and Post-Nicene Fathers, Second Series；Vol. 12.；Edited by Philip Schaff and Henry Wace.；Buffalo, NY: Christian Literature Publishing Co.,；1895.；<http://www.newadvent.org/fathers/360206044.htm>.}

% structure: p0001 source=h1 target=section reason=page-title

\section{第六卷，第四十六书}

\Emph{致\Person{斐理克斯}，\Place{比绍伦姆}（\Place{佩萨罗}）主教}\par

\Person{额我略}致\Person{斐理克斯}主教等\par

我们惊奇于你的兄弟之爱，你竟不顾我们可敬的前任交给你的命令的文义，而以与古老惯例不符的方式祝圣由\Person{若望}（即持信人）建造的隐修院。因为，虽然在所述命令中特别规定你应在没有公开弥撒的情况下祝圣该地方，但正如我们所听到的，你的主教座位已安置在那里，弥撒的神圣礼仪也在那里公开举行。如果这是真的，我们在此劝勉你，撇开一切借口，使你的主教座位完全从那里移走，并且从此不得在那里举行公开弥撒。但是，正如习俗和命令的文义所指示的，如果他们愿意在那里为他们举行弥撒，就由你为此指定一位司铎。\par

此外，我们愿意，赖天主的恩宠，在同一隐修院里，一如既往地常有一群天主的仆人团体，正如上述\Person{若望}所请求的，也如现状。至于他告知我们已被你的兄弟之爱取走的圣爵，如果属实，请速归还。那么，让你的圣德努力完成这些事，以使上述持信人无需因同一事再次求助于我们。\par

\SourceRecord{Translated by James Barmby.；Nicene and Post-Nicene Fathers, Second Series；Vol. 12.；Edited by Philip Schaff and Henry Wace.；Buffalo, NY: Christian Literature Publishing Co.,；1895.；<http://www.newadvent.org/fathers/360206046.htm>.}

% structure: p0001 source=h1 target=section reason=page-title

\section{第六卷·第四十八封信}

\Emph{致乌尔比库斯，院长。}\par

\Person{额我略}致\Person{乌尔比库斯}，\Place{圣赫尔梅斯隐修院}院长，该院位于\Place{帕诺尔穆斯}。\par

凡受天主灵感激励，急忙离开今世工作而皈依天主者，应以爱德接待，以仁慈安慰多方慰藉，俾因天主助佑，乐于在其所选生活状态中恒心持守。既然呈函者\Person{阿加托}渴望在您的爱德之隐修院内皈依，我们劝您以一切温柔和爱心接纳他，并以勤勉劝勉激发他对永生的渴望，悉心关切他灵魂的救恩；为使因您的训导，他以虔敬之心恒心事奉我们的天主，这既有助于他离弃世俗，亦使他的皈依成为您自己赏报的增加。然而，惟当他的妻子也愿同样皈依时，方可如此接纳。因为，当两人身体因婚姻之结合而成一体时，部分皈依而部分留于世俗，实为不当。\par

\SourceRecord{Translated by James Barmby.；Nicene and Post-Nicene Fathers, Second Series；Vol. 12.；Edited by Philip Schaff and Henry Wace.；Buffalo, NY: Christian Literature Publishing Co.,；1895.；<http://www.newadvent.org/fathers/360206048.htm>.}

% structure: p0001 source=h1 target=section reason=page-title

\section{第六卷·第四十九书}

\Emph{致\Person{帕拉弟乌斯}主教}\par

\Person{额我略}致\Place{高卢}\Place{桑通}主教\Person{帕拉弟乌斯}（\Emph{桑特}）\par

\Person{额我略}致其司铎、持信者\Person{柳帕里克}前来见我们，告知我们：您的主教阁下为纪念圣宗徒\Person{伯多禄}和\Person{保禄}，以及殉道者\Person{老楞佐}和\Person{庞加爵}，建造了一座圣堂，并在其中安置了十三座祭台，但我们得知其中有四座尚未祝圣，因您渴望将上述诸圣的圣髑安放其中。由于我们已恭敬地提供了圣宗徒\Person{伯多禄}和\Person{保禄}以及殉道者\Person{老楞佐}和\Person{庞加爵}的圣髑，我们劝您以虔诚之心接受，并赖主助佑安放妥当，且首先须确保为在那里服务者提供所需之物，不使缺乏。\par

\SourceRecord{Translated by James Barmby.；Nicene and Post-Nicene Fathers, Second Series；Vol. 12.；Edited by Philip Schaff and Henry Wace.；Buffalo, NY: Christian Literature Publishing Co.,；1895.；<http://www.newadvent.org/fathers/360206049.htm>.}

% structure: p0001 source=h1 target=section reason=page-title

\section{第六卷 第五十函}

\Emph{致\Person{布伦希尔德}王后}\par

额我略致法兰克王后\Person{布伦希尔德}。\par

你的来函显示出宗教精神和虔诚心灵的恳切，这不仅使我们称赞你请求的目的，也使我们乐意允诺你所要求的。因为，拒绝基督徒的虔诚和正直心灵的渴望，尤其是当我们知道你以全心要求并拥抱那既能保护信友的信德又能同样促进灵魂的救恩的事，确实与我们不相宜。因此，我们以适当的荣誉问候陛下，并告知你：通过这些书信的携带着、你所描述的司铎\Person{勒乌帕里克}，我们根据陛下的请求，以应有的尊敬交付了真福宗徒\Person{伯多禄}和\Person{保禄}的一些圣髑。但是，为使这可称赞的宗教虔诚在你们中更加显著，你们必须确保这些圣人的恩惠被以尊敬和适当的荣誉安置，并且那些服侍他们的人不受任何负担或骚扰的烦扰，免得或许在外在需要的压力下，他们在事奉天主时变得无益而迟缓，并且（但愿天主不会允许如此）所赐予的圣髑遭受损伤和被忽视。因此，请陛下照料他们的安宁，以便他们因你的慷慨而免受一切困扰，以平静的心灵向我们的天主献上赞颂，并且为此赏报也在永生中为你积蓄。\par

\SourceRecord{Translated by James Barmby.；Nicene and Post-Nicene Fathers, Second Series；Vol. 12.；Edited by Philip Schaff and Henry Wace.；Buffalo, NY: Christian Literature Publishing Co.,；1895.；<http://www.newadvent.org/fathers/360206050.htm>.}

% structure: p0001 source=h1 target=section reason=page-title

\section{第六卷·第五十一封书信}

\Emph{致前往英格兰（Angliam）的弟兄们}。\par

\Person{额我略}，天主的众仆之仆，致我们的主耶稣基督的仆人们。\par

既然开始行善后中途退缩，不如当初没有开始，所以，最亲爱的儿子们，你们必须完成那藉主的帮助已经开始的好工作。那么，不要让旅途的劳苦，也不要让说坏话之人的舌头阻挠你们；却要以全部的恒心和全部的热忱继续那在天主指引下已开始的工作，要知道巨大的辛劳之后是永远赏报的光荣。要在一切事上谦卑地服从你们的院长\OriginalTerm{院长}{præposito}，他正要回到你们那里，我们也任命他为你们的院长，知道凡通过他的劝诫在你们身上得以成就的事，都会对你们的灵魂有益。愿全能的天主以祂的恩宠保护你们，并赐我能在永恒的家乡看到你们劳苦的果实；这样，即使我不能与你们一同劳苦，也能在赏报的喜乐中与你们在一起；因为我实在渴望劳苦。愿天主保佑你们平安，最亲爱的儿子们。写于八月朔日前的第十日，我们最虔诚的奥古斯都、吾主皇帝莫里奇乌斯·提庇留第十四年，上述陛下执政官第十三年，小纪十四年。\par

\SourceRecord{Translated by James Barmby.；Nicene and Post-Nicene Fathers, Second Series；Vol. 12.；Edited by Philip Schaff and Henry Wace.；Buffalo, NY: Christian Literature Publishing Co.,；1895.；<http://www.newadvent.org/fathers/360206051.htm>.}

% structure: p0001 source=h1 target=section reason=page-title

\section{第六卷，第五十二书}

\Emph{致\Person{佩拉吉乌斯}与\Person{塞雷努斯}二位主教。}\par

格里高利致高卢地区\Place{图尔尼}的\Person{佩拉吉乌斯}与\Place{马西利亚}（马赛）的\Person{塞雷努斯}二位主教。\OriginalTerm{同级}{A paribus}。\par

虽然对于具有悦乐天主的爱德的司铎，虔诚者无需任何推荐，然而，既然写信的适当时机已到，我们仍认为应向您的兄弟之情致函，告知我们已蒙主助，为灵魂的益处，将天主的仆人\Person{奥古斯丁}（我们确信其热忱）与其他天主的仆人派往贵地。您的圣座必须以司铎般的热情援助他，并尽快施以援手。我们也曾吩咐他，为使您更乐意支持他，应将手中事务完全告知您，因为知道一旦您了解此事，必会为天主之故而全心全力按需援助他。\par

此外，我们尽一切方式将我们的共同儿子、司铎\Person{坎迪杜斯}托付于您的爱德，他已被派往管理我们教会的产业。于八月朔日前第十日，小纪14年发出。\par

\SourceRecord{Translated by James Barmby.；Nicene and Post-Nicene Fathers, Second Series；Vol. 12.；Edited by Philip Schaff and Henry Wace.；Buffalo, NY: Christian Literature Publishing Co.,；1895.；<http://www.newadvent.org/fathers/360206052.htm>.}

% structure: p0001 source=h1 target=section reason=page-title

\section{第六卷·书信五十三}

\Emph{致\Person{维吉利乌斯}主教。}\par

\Person{额我略}致\Person{维吉利乌斯}，\Place{阿雷拉特}（\Place{阿尔勒}）主教兼总主教。\par

虽然我们确信你的弟兄之爱致力于善工，并且你自动参与中悦天主的各种事业，但我们仍认为以弟兄般的爱德写信给你是有益的，这样，被我们的书信所激励，你或许会增加你自愿施予的慰藉。因此，我们通知你的圣德，我们已派遣天主的仆人奥斯定——本信的持信人——与其他天主的仆人一同前往，为在他将要到达的地区拯救灵魂，正如他能亲自当面告知你。在此情况下，你必须以祈祷和援助来协助他，并在需要时给予他你帮助的支持，并以父职和司铎的安慰适当地鼓励他，以便当他获得你圣德的帮助后，若他能（如我们所愿）为天主赢得任何收获，你也能因虔诚地以你的丰厚支持助益他的善工而与他一同获得赏报。此外，关于我们的共同之子——司铎坎迪杜斯，以及我们教会的小小产业，让你的弟兄之爱——因与我们同心——注意将两者视为托付给你的；这样，藉你圣德的帮助，从中可以有所增收以养活穷人。既然你的前任多年来掌管这一产业，并将收集的款项留在他自己手中，让你的弟兄之爱考虑这些款项属于谁，应该付给谁，并将它们归还给我们，交给上述我们的儿子——司铎坎迪杜斯。因为，各国国王所保存的，竟被说成被主教们夺走，这是极其可恶的。\par

\SourceRecord{Translated by James Barmby.；Nicene and Post-Nicene Fathers, Second Series；Vol. 12.；Edited by Philip Schaff and Henry Wace.；Buffalo, NY: Christian Literature Publishing Co.,；1895.；<http://www.newadvent.org/fathers/360206053.htm>.}

% structure: p0001 source=h1 target=section reason=page-title

\section{第六卷·第五十四书}

\Emph{致\Person{德西德留}与\Person{塞亚格留}二位主教}\par

\Person{格列高利}致维也纳的\Person{德西德留}与奥古斯托杜努姆（即\Place{欧坦}）的\Person{塞亚格留}，乃高卢地区主教。\Emph{\OriginalTerm{同类}{A paribus}}。\par

鉴于你们的真诚爱德，我们深信你们出于对宗徒之长伯多禄的敬爱，必将慷慨援助我们的属下；尤其是鉴于事情的性质，即使没有请求，你们也应自愿给予援助，更何况见到他们辛勤劳苦。故此，我们通知尊驾，在主安排之下，我们派遣了天主的仆人、持此信者\Person{奥斯定}——他的热忱与勤勉我们深知——连同其他天主的仆人，为拯救那些地区的灵魂前往。当你们从他的陈述中充分了解他所领受的使命后，请你们的兄弟之谊在一切必要之事上予以援助，以便你们能够合宜适当地成为善工的协助者。为此，愿你们的兄弟之谊努力在此事上显示如此的热忱，使你们的行动能向我们证明我们对你们的美好听闻确属实情。我们在此将我们最亲爱的共同之子司铎\Person{坎迪杜斯}全然托付给你们，我们已委托他管理本教会位于那些地区的产业。\par

\SourceRecord{Translated by James Barmby.；Nicene and Post-Nicene Fathers, Second Series；Vol. 12.；Edited by Philip Schaff and Henry Wace.；Buffalo, NY: Christian Literature Publishing Co.,；1895.；<http://www.newadvent.org/fathers/360206054.htm>.}

% structure: p0001 source=h1 target=section reason=page-title

\section{卷六，书信五十五}

\Emph{致\Person{普洛塔西乌斯}主教。}\par

\Person{额我略}致\Person{普洛塔西乌斯}，高卢\Place{阿奎}（\Emph{\Place{艾克斯}}）主教。\par

你对蒙福的宗徒之长\Person{伯多禄}所怀的大爱如何使你出众，是显而易见的，这不仅出自你职务的特权，也出自你对其教会益处所倾注的热忱。我们已从携带此信的天主仆人\Person{奥斯定}的报告中得知此情，故为你内心的深情与对真理的热忱而极其欣喜；我们感谢，虽然身体分离，你仍表明在心意上与我们同在，因你向我们展现了合宜的兄弟之爱德。为使实际行动证实对你的美誉，请告诉我们兄弟和同为主教\Person{维尔吉利乌斯}，将他前任多年收取并扣留的款项交给我们：因为这是穷人的财产。倘若他试图以任何方式推辞——我们不信会如此——你更确切地知道实情，因当时你担任\OriginalTerm{管家}{vicedominus}，请向他说明情况，并劝告他不要扣留圣\Person{伯多禄}及其穷人的财产。不过，尽管我们的人或许不需要此证，但请不要拒绝为此案提供你的证词；如此，为了真理以及你善意的虔诚，蒙福的宗徒\Person{伯多禄}——你因对他的爱而行此事——愿他藉其转祷，在此世和来生报答你。我们衷心将我们的共同儿子，司铎\Person{坎迪杜斯}，托付给您的圣德，我们已将这份产业的管理交付于他。\par

\SourceRecord{Translated by James Barmby.；Nicene and Post-Nicene Fathers, Second Series；Vol. 12.；Edited by Philip Schaff and Henry Wace.；Buffalo, NY: Christian Literature Publishing Co.,；1895.；<http://www.newadvent.org/fathers/360206055.htm>.}

% structure: p0001 source=h1 target=section reason=page-title

\section{第六卷，书信五十六}

\Emph{致院长\Person{斯德望}}\par

\Person{额我略}致\Person{斯德望}等\par

天主的仆人、持信者\Person{奥斯定}给我们带来的报告使我们欢喜，因为他告诉我们，你的爱德正如你应当的那样警醒；他还进一步证实，司铎们、执事们以及全体会众都同心合意地生活。既然首领的良善是属下有益的规范，我们恳求全能天主，藉祂慈爱的恩宠，常激励你们行善，并保护托付给你们的人免受一切魔鬼诡计的诱惑，且恩赐他们与你们在爱德中并与祂中悦的生活方式度日。\par

但是，既然人类的仇敌从不停止谋害我们的行为，以便在某些部分欺骗事奉天主的灵魂，因此，最亲爱的儿子，我们劝告你，要警醒地履行你忧虑的关怀，借着祈祷和警醒如此保守托付给你的人，使游走的狼找不到撕裂羊群的机会；为的是，当你将你所承担了责任的人安然无恙地交付我们的天主时，祂能因祂的恩宠以酬报你的辛劳，并增加你对永生的渴望。\par

我们已经收到了你送来的汤匙和盘子，我们感谢你的爱德，因为你显示了你如何爱穷人，既然你送来了他们需要的东西供他们使用。\par

\SourceRecord{Translated by James Barmby.；Nicene and Post-Nicene Fathers, Second Series；Vol. 12.；Edited by Philip Schaff and Henry Wace.；Buffalo, NY: Christian Literature Publishing Co.,；1895.；<http://www.newadvent.org/fathers/360206056.htm>.}

% structure: p0001 source=h1 target=section reason=page-title

\section{第六卷，第五十七封}

《\Emph{致贵族阿利吉乌斯}》。\par

格列高利致高卢贵族阿利吉乌斯。\par

我们从天主的仆人，这位信件的携带者奥斯定得知，在您身上闪耀着何等伟大的良善，何等伟大的温和，以及中悦基督的爱德；我们感谢全能的天主，祂赐予您这些祂慈爱的恩赐，藉此您得以在人前受尊崇，并在祂眼中——这真正有益——成为荣耀。因此，我们恳求全能的天主，愿祂加倍赐予您这些恩赐，并保护您及您所有的一切在祂的护佑之下，如此安排您尊荣在此世的行为，使它们对您在此世有益，并且——更值得渴望的——在来世也有益。我们以慈父般的温情问候您的尊荣，恳求您，使这封信的携带者以及与他同行的天主仆人们，在需要时能得到您的援助，以便他们在体验您的眷顾时，能更好地完成所托付给他们的事。\par

此外，我们特别向您推荐我们的儿子，司铎堪迪杜斯，我们派遣他管理我们在您区域的教会产业；我们相信，如果您以虔诚的心志援助穷人的事务，您的尊荣将从我们的天主那里获得回报。\par

\SourceRecord{Translated by James Barmby.；Nicene and Post-Nicene Fathers, Second Series；Vol. 12.；Edited by Philip Schaff and Henry Wace.；Buffalo, NY: Christian Literature Publishing Co.,；1895.；<http://www.newadvent.org/fathers/360206057.htm>.}

% structure: p0001 source=h1 target=section reason=page-title

\section{第六卷，第五十八封书信}

\Emph{致\Person{提奥多里克}与\Person{提奥德贝尔特}}。\par

\Person{格列高利}致\Person{提奥多里克}与\Person{提奥德贝尔特}兄弟，\Place{法兰克}人的君王。\Emph{同级书}。\par

既然全能的天主以信德的正直装饰了你们的王国，并以基督宗教的纯洁使它在其他民族中卓然出众，我们对你们寄予厚望，深信你们必会竭尽全力，期望你们的臣民也能归向那使你们得以成为他们君王与主人的信德。既然如此，我们得知，因天主的仁慈，\Place{盎格利}民族渴望皈依基督信仰，但他们邻近的司铎却忽视他们，疏于以自身的劝勉激发其热忱。为此，我们决意派遣天主之仆\Person{奥古斯丁}——此信的携带者，其热忱与勤勉为我们所熟知——连同其他天主之仆前往他们那里。我们也嘱咐他们从邻近地区带一些司铎同行，以便他们能了解\Place{盎格利}人的意向，并尽天主所赐，以规劝助成其愿望。为使他们在这一工作中能展现效力与能力，我们恳请贤明尊驾，以慈父之爱致敬，愿我们所派遣之人堪当蒙受陛下的恩宠。此外，既然这关乎灵魂，愿你们的权力保护并援助他们；如此，那深知你们以虔诚之心和全心关注祂事业的全能天主，将欣然引导你们的事业，并在尘世统治之后，带领你们进入天国。\par

此外，我们恳请陛下将我们至爱的儿子，司铎\Person{坎迪杜斯}，即我们教会产业的管理者，视为受您举荐之人，以使蒙福的宗徒之长\Person{伯多禄}藉其转祷回报您，同时，因您顾念赏报，在他的穷人的事务上施以庇护。\par

\SourceRecord{Translated by James Barmby.；Nicene and Post-Nicene Fathers, Second Series；Vol. 12.；Edited by Philip Schaff and Henry Wace.；Buffalo, NY: Christian Literature Publishing Co.,；1895.；<http://www.newadvent.org/fathers/360206058.htm>.}

% structure: p0001 source=h1 target=section reason=page-title

\section{卷六，第五十九封}

\Emph{致法兰克王后布伦希尔德书}\par

额我略致布伦希尔德等。\par

您的卓越所表现出的基督信仰早已为我们所知，以至于我们丝毫不怀疑您的良善，反而完全确信您必会虔诚而热忱地与我们同心协力于信仰之事，并以您宗教的真诚提供最丰沛的援助。因此，我们满怀确信，并以慈父般的爱德问候您，特此告知：我们获悉，盎格利民族蒙天主恩准渴望成为基督徒，但邻近的司铎们对他们却毫无牧灵关怀。为免他们的灵魂不幸丧亡于永罚，我们已用心派遣此信的携带者——天主的仆人奥斯定（他的热忱与勤勉为我们所深知）——连同其他天主的仆人前往他们那里；以便通过他们，我们得以了解他们的意愿，并尽可能在您亦与我们共同努力下，为他们的归化筹谋。我们也已嘱咐他们，为完成这一计划，应从邻近地区带领司铎同行。因此，敬请您的卓越（您惯行善工）因我们的请求以及出于对天主的敬畏，惠然以各种方式向此人推荐，并热切赐予他您的保护之恩，以您的庇护援助他的辛劳；为使他的工作获得圆满果实，请确保他在您的保护下安全前往上述盎格利民族，好使我们的天主（祂在此世以中悦祂的美善装饰了您）使您既在此世感恩，亦在永恒安息中与祂的圣人们一同感恩。\par

再者，我们将我们亲爱的儿子坎迪杜斯（司铎兼位于贵处之教会产业的\OriginalTerm{总管}{rector}）托付于您的基督信仰，恳请他在一切事上蒙受您保护之恩。\par

\SourceRecord{Translated by James Barmby.；Nicene and Post-Nicene Fathers, Second Series；Vol. 12.；Edited by Philip Schaff and Henry Wace.；Buffalo, NY: Christian Literature Publishing Co.,；1895.；<http://www.newadvent.org/fathers/360206059.htm>.}

% structure: p0001 source=h1 target=section reason=page-title

\section{第六卷·第六十封书函}

\Emph{致欧洛吉主教}\par

额我略致亚历山大的欧洛吉主教。\par

爱德——一切美善之母与守护者——将众多心灵结合为一体，凡在心神眼前所拥有的人，绝不视之为不在眼前。因此，至亲爱的弟兄，我们既因爱德之根而紧密相连，则肉身的别离或地点的远近也无权对我们施加任何主张，因为我们既为一体，必不相隔。我们愿与众弟兄常保此共同爱德。然而，有一特殊原因使我们与亚历山大教会相连，且仿佛因一特别律法，促使我们更倾向于爱她。因为，众所周知，真福圣史马尔谷乃由宗徒之长圣伯多禄派往亚历山大，故此我们因这位师长与其门徒的统一而相连。如此，我似乎因着师长而居于门徒之座，而您则因着门徒而居于师长之座。\par

此外，我们亦因尊座之功德而被这心志统一所系，因为我们知道您有效地追随了您奠基者之规诫，并以全副热忱投身于您师长之怀抱，救恩的宣讲正是从那里传播到您所在的地域。因此，当我们接到尊座的书信时，我们的心灵因您的兄弟般来访而欣喜，却也因您所提及的难以言表的负担而深感忧伤。我们以兄弟般的同情与您一同为您的痛苦叹息。然而，既然各种动荡正在四处蔓延，在共同的需要中，人应少为自己的苦难忧伤，而应通过学习耐心忍受，去克服我们无法完全避免之事。\par

然而，我们自身因伦巴第人的刀剑，在每日对市民的劫掠、伤残与杀戮中所遭受的，我们不愿诉说，以免在叙述自己的悲痛时，因您对我们的同情而增加您的忧愁。\par

此外，不久前我们寄了一封我们自己的信函给萨比尼亚诺——他代表我们的教会驻于皇城——他本应转呈给您的弟兄之尊。若您已收到此信，我们不解为何未见回音。因此，既然必须谨慎，以免任何人的傲慢在教会中引起冒犯，您应当仔细阅读此信，并竭尽全力、全心专注地维护属于您尊严及教会和平之事。\par

愿全能的天主，以祂仁慈之爱的恩宠，赐予您符合司铎身份的性情与爱德，在祂的侍奉中保护您，内外护佑您脱离一切逆境，并仁慈赐予，使迷途者的灵魂因您的宣讲而皈依于祂。\par

我们以当有的爱德接待了本信的携带着，我们的共同儿子——执事依西多禄，他给我们带来了圣史马尔谷的祝福。的确，您因善生的功德而光辉灿烂，给我们送来了芬芳的话，这话接近乐园。然而我们，因我们是罪人，从西方给您送去木材，这木材适于造船，象征我们心灵的动荡，仿佛常在波涛中颠簸；我们原想送更大的木材，但船不够大。八月，第十四小纪。\par

\SourceRecord{Translated by James Barmby.；Nicene and Post-Nicene Fathers, Second Series；Vol. 12.；Edited by Philip Schaff and Henry Wace.；Buffalo, NY: Christian Literature Publishing Co.,；1895.；<http://www.newadvent.org/fathers/360206060.htm>.}

% structure: p0001 source=h1 target=section reason=page-title

\section{Book VI, Letter 61}

\Emph{致卡斯托里乌斯，公证员}。\par

额我略致卡斯托里乌斯等。\par

尊贵的安得烈大人不断催促我恢复在拉文纳教会按照古老习惯使用\OriginalTerm{披肩}{pallium}。你知道若望主教曾写信给我，说该教会的主教们在隆重\OriginalTerm{连祷}{litanies}中使用披肩原是习惯。该教会执事阿德奥达图恳切请求我，并以宣誓使我相信，该地的主教们惯常一年四次在连祷中使用披肩。但前述安得烈大人在信中却说，拉文纳主教习惯在除四旬期外的所有时间在连祷中使用披肩。而且这些连祷，他竟不羞于说是每日的，他断言是隆重的。因此我完全惊讶。但是，阁下，请不要顾及任何人的情面，任何人的话语；只让敬畏天主和正直之心摆在眼前，并向年长之人、该教会的总执事（我想他不会为了别人的荣誉而发假誓），以及在若望主教之前领受圣秩的年长之人，或者还有没有领受圣秩但年龄更大的人查问；让他们来到圣亚玻里纳利的遗体前，触摸他的坟墓，宣誓在若望主教之前是怎样的习俗。因为，如你所知，他是一位非常僭越的人，在骄傲中试图篡取许多事归于己有。凡是忠信和庄重的人按照附后的誓词宣誓的内容，我们愿意在该教会中保留。但要看你不可疏忽，且在这件事上无人败坏你的忠信和热忱；因为我知道你的热忱。勤勉行事，但要使上述教会不因违背正义而受贬低，而是保留若望主教之前存在的惯例。此外，为使自己满意，不要只向两三个人查问，而要尽可能多地找年长且庄重的人，这样我们既不否认该教会的古老习俗，也不让步于它新近贪求并尝试的事情。但要以和蔼温柔的方式做一切事，使你的行为严格而你的言语温和。那把留在拉文纳的剑，正如我们已经写过的，你随身带来；并仔细留意我们的儿子博尼法斯执事和尊贵的茂伦提乌斯\OriginalTerm{档案管理员}{chartularius}写给你的信。\par

我指着圣父、圣子、圣神，不可分离的天主圣三之全能宣誓，并指着真福殉道者亚玻里纳利的遗体，我不偏袒任何人，也无任何个人利益干预，给出我的证词。但我所知并亲身确知的是，在已故若望主教之前，拉文纳主教，在宗座这样或那样的\OriginalTerm{宗座使节}{apocrisiarius}面前，在某些日子，有使用披肩的习惯，并且我不知道他曾在其中暗中僭取，或在宗座使节缺席时僭用。\par

\SourceRecord{Translated by James Barmby.；Nicene and Post-Nicene Fathers, Second Series；Vol. 12.；Edited by Philip Schaff and Henry Wace.；Buffalo, NY: Christian Literature Publishing Co.,；1895.；<http://www.newadvent.org/fathers/360206061.htm>.}

% structure: p0001 source=h1 target=section reason=page-title

\section{第六卷，第六十三封书信}

\Emph{致贵族根纳迪乌斯}。\par

\Person{额我略}致\Place{非洲}贵族\Person{根纳迪乌斯}。\par

我们毫不怀疑，阁下还记得两年前我们曾为我们的兄弟和同僚主教\Person{保禄}写信，请求阁下以尊驾之威望支持他，因据说他正遭受多纳徒派迫害之苦，他渴望来到我们这里；其目的是：既然据报他在那里无法获得任何救助，我们可在查明真相后，以兄弟般的同情给予他建议，并与他就如何以有益的安排应对那毒害性狂妄之疯狂进行磋商。然而，就我们上述兄弟所告知我们的，他不仅未从任何人那里得到援手，反而因种种阻碍而无法安全抵达罗马城。但是，当我们让您来信读给他听后，他回答说，他并非因为镇压多纳徒派而遭受某些人的恶感，而是说他因其对天主教信仰的维护而受到许多人的冷遇；他还告诉了我许多其他事情，但因眼下不是提及这些事情的合适时机，我们认为最好保留不说。\par

既然摆在我们面前的问题不是关于世俗事务，而是关于灵魂的健康，而您的断言和他的说法又不相同，我们便无法特别作出答复，因为我们尚未调查真相；而且，当我们收到阁下信函时，我们正因病卧床。但是，如果全能的天主愿意，待祂使我们恢复先前的健康后，我们将会通过勤勉的查究来审察真相。并根据我们所能够了解的情况，我们将籍由天主的仁慈来处理此案，不但要让那些迷失者的灵魂（您不吝关怀其治疗）恢复健康，而且还要让真信仰的维护者所尚存的一切，籍由我们救赎者保护性的恩宠得以保全。\par

但关于上述主教（您断言他被剥夺了共融），我们大为惊奇的是，为何是阁下的信函，而不是其首席主教的信函，将此事通知了我们。\par

\SourceRecord{Translated by James Barmby.；Nicene and Post-Nicene Fathers, Second Series；Vol. 12.；Edited by Philip Schaff and Henry Wace.；Buffalo, NY: Christian Literature Publishing Co.,；1895.；<http://www.newadvent.org/fathers/360206063.htm>.}

% structure: p0001 source=h1 target=section reason=page-title

\section{第六卷，第六十五封书信}

\Emph{致皇帝莫里斯。}\par

格里高利致莫里斯·奥古斯都。\par

在战争与无数忧虑之中，您以不懈的热忱肩负着基督信仰共和国的治理，令我同普世万民共感欣慰的是，陛下始终以警觉之心守护信仰，正是这信仰使我等主人的帝国光彩夺目。因此，我全然相信，正如您以宗教情怀的爱守护天主的圣业，天主也以其至尊的恩宠守护并助佑您的圣业。至于陛下出于正义与最纯净宗教的热忱，如何因多纳徒派最邪恶的败坏而震动，您所颁发命令的要旨已清楚表明。但来自非洲行省的最可敬的主教们声称，由于不明智的纵容，这些命令竟被如此漠视，以至于在那里既无人敬畏天主的审判，至今也未见帝国命令得到执行；并进而补充说：在上述行省，因多纳徒派的贿赂得逞，天主教信仰竟被公开出售。另一方面，荣耀的根纳狄乌斯也对提出此类控诉者之一表示不满；另有两人随同他证明了此事。然而，由于此事涉及一位世俗法官，我请这些主教前往陛下足前，好让他们亲身向陛下最宁静的耳朵陈述他们声称为了天主教信仰所忍受的一切。\par

为此，我恳请我主上们的基督信仰，为你们灵魂的益处和最虔诚后嗣的生命，下令以严格命令惩罚那些经查明确属所述之人，并以拯救之手遏止丧亡之人的毁灭，为疯狂的心灵施以纠正的药石，治愈他们错误的毒咬；如此，因陛下之筹谋的救治驱逐了致命邪恶的黑暗，真信仰在那些地方散布其宁静的光芒，天上的胜利必在我们救赎主眼前等待你们，因为你们既在外抵御敌人，也于内使他们摆脱邪恶魔计之毒——这乃是一件更为荣耀的事。\par

\SourceRecord{Translated by James Barmby.；Nicene and Post-Nicene Fathers, Second Series；Vol. 12.；Edited by Philip Schaff and Henry Wace.；Buffalo, NY: Christian Literature Publishing Co.,；1895.；<http://www.newadvent.org/fathers/360206065.htm>.}

% structure: p0001 source=h1 target=section reason=page-title

\section{第六卷，第六十六封书信}

\Emph{致亚大纳削司铎}\par

额我略致依索利亚司铎亚大纳削。\par

我们既为那些被异端邪说的错误从教会的合一中割裂的人感到痛苦和悲伤，也同样为那些因宣认天主教信仰而留在教会怀中的人感到欢喜。而且，正如我们有责任以牧灵的关怀去反对前者的不敬，同样也应当对后者的虔诚宣认施予恩惠，并宣布他们的见解是纯正的。\newline{}因此，鉴于有人对你——亚大纳削，里考尼亚省境内名为塔姆纳库斯的圣米肋隐修院的司铎——产生了信仰不纯的怀疑，你为了使信仰宣认的纯正得以显明，选择了诉诸我们治理的宗座，并声称你在身体上受到惩罚后，曾不公正且冲动地做了一些事。尽管在强迫下所做的事绝不受教规的指责，并且被视为无效（因为强迫者本人就使那被迫承认和做的事归于无效），而且那种出于自发意愿的宣认才更应该被接纳和拥抱——就像你在我面前所作的宣认那样——但为了避免任何不确定的可能性，我们仍谨慎地就你的事写信给我们兄弟及同僚、君士坦丁堡城的主教，请他写信告知我们所发生的事。他经我们多次催促后回信说，在你那里发现了一本书卷，其中包含许多异端言论，为此他向你发怒。他为了满足我们的要求，将此书卷借给我们，我们仔细阅读了其前部内容；既然在其中发现了明显的异端邪说毒素，我们便禁止再读下去。但是，既然你向我们保证你是在纯朴中阅读了它，并且为了消除一切不确定怀疑的根据，你亲手交给我们一份文书，在其中阐述你的信仰，最明确地谴责了所有异端以及一切与天主教信仰或宣认的纯正相悖的东西，并声明你一直接纳且仍接纳四大神圣大公会议所接纳的一切，谴责且仍谴责它们所谴责的一切，还承诺接受并坚持那在犹斯定皇帝时代为三章而举行的会议；而且，你被我们禁止阅读那本交织着致命错误毒素的书卷，并拒绝和谴责其中一切违背天主教信仰纯正的言论或隐含之意，并承诺不再读它——我们出于这些理由（你的信仰也因你亲手所写的文书，在天主保佑下，被我们视为纯正天主教信仰），根据你的宣认，判定你免于一切异端邪见的污点，是天主教徒；我们宣布，你借我们的救主耶稣基督的恩宠，证明自己在一切事上是纯正信仰的宣认者和追随者；同时，我们授予你自由许可，不顾一切，返回你的隐修院，恢复你的地位和品位。\par

我们也愿意就此事写信给我们至爱的兄弟、君士坦丁堡城的主教——他是在上述圣若望之后被祝圣的。但是，由于惯例是我们在接到他的教务书信之前不应写信，所以我们延迟了。然而，在接到他的书信后，我们会在方便时将这些事通知他。\par

\SourceRecord{Translated by James Barmby.；Nicene and Post-Nicene Fathers, Second Series；Vol. 12.；Edited by Philip Schaff and Henry Wace.；Buffalo, NY: Christian Literature Publishing Co.,；1895.；<http://www.newadvent.org/fathers/360206066.htm>.}

% structure: p0001 source=h1 target=section reason=page-title

\section{第七卷，第二封信}

\Emph{致哥伦布主教}\par

额我略致努米底亚主教哥伦布。\par

我们经由送信者，您的执事，收到了您兄弟的书信，其中您告知我们在你们那里关于保禄主教本人所做之事。此事办得如此之晚，以致他如今不能亲自出现在这里。因为我们的儿子显贵格纳狄乌斯阁下也派了他的掌玺官到我们这里，涉及同一案件。但当我们派人查问他是否愿意在我们面前控告他（即保禄主教）时，他回答说，他绝非为此意图被派来，而只是从他的教会带来了三个特定的人，他们将提出许多反对他的事。因此，我们既发现他未准备好提起诉讼，也不因那些人的身份而将他们视为适合控告主教的人，我们不能反对或阻止常常提到的那位保禄主教，他请求我们准许他前往皇城；我们立即准许他按自己的请求，和另外两个他将带走的人一起出发。因此，如果曾有任何可以合理反对他的事，相应的做法本应是让他立即前来，并让您的兄弟将一切详情告知我们，正如您现在所做的一样。因为，至于您向我们表示，因我们常以书信访问您而遭受许多人的敌意，那毫无疑问，最可敬的弟兄，善人受恶人的嫉妒，专心于神圣工作的人受乖戾之人的反对所困扰。但是，这些恶事围绕你们越甚，你们就越应当立即专注于所托付的管理工作，看守基督的羊群；不义之人的反对压逼你们越甚，牧灵挂虑的关心就越当激励你们更积极，并且非常确信应许的赏报，好使你们能向大牧者献上受命所做之工的收益。\par

\SourceRecord{Translated by James Barmby.；Nicene and Post-Nicene Fathers, Second Series；Vol. 12.；Edited by Philip Schaff and Henry Wace.；Buffalo, NY: Christian Literature Publishing Co.,；1895.；<http://www.newadvent.org/fathers/360207002.htm>.}

% structure: p0001 source=h1 target=section reason=page-title

\section{\WorkTitle{第七卷，第四封信}}

\Emph{致济利禄主教。}\par

\Person{额我略}致\Person{济利禄}，\Place{君士坦丁堡}主教。\par

我们以应有的爱德接待了我们共同的孩子，司铎\Person{乔治}和你的执事\Person{德奥多禄}；并且我们因你从管理教会事务的职责转向治理灵魂而欣喜，因为按照真理的声音：\BibleQuote{人在小事上忠信，在大事上也忠信}（\BibleRef{路 16:10}）。又对那善于管理的仆人说：\BibleQuote{因为你曾在少许事上忠信，我要派你管理许多事}（\BibleRef{玛 25:23}）；并且随即也对他说到永恒的赏报：\BibleQuote{进来享受你主人的欢乐吧！}你在信中说你极渴望安息。然而，这正表明你恰当地承担了牧职，因为，正如管理之位应被拒绝给那些贪求它的人，也当被给予那些逃避它的人。\BibleQuote{谁也不得自己擅取这尊位，而是蒙天主召选，如同亚郎一样}（\BibleRef{希 5:4}）。那杰出的宣讲者又说：\BibleQuote{若一人替众人死了，那么众人就都死了；基督替众人死了，是为使活着的人不再为自己活，而是为替他们死而复活的那位活}（\BibleRef{格后 5:14}）。又对圣教会的牧者说：\BibleQuote{西满，若望的儿子，你爱我吗？你牧放我的羊}（\BibleRef{若 21:17}）。由此可见，若有人有能力却拒绝牧放全能天主的羊群，便表明他不爱那首席牧者。因为圣父的独生子为了成就众人的益处，从父的隐秘中来到我们中间，若我们为了自己的隐秘而舍弃近人的益处，又该当如何呢？因此，安息本当全心全意地去渴望；然而为了多人的益处，有时也当放下它。因为，正如我们当一心渴望逃避俗务，但若缺乏宣讲之人，我们也必须甘愿肩负起俗务的重担。两位先知的作为教导了我们这一点：一位试图逃避宣讲的职责，另一位却渴望它。因为\Person{耶肋米亚}对派遣他的主回答说：\BibleQuote{唉！吾主上主，我不会说话，因为我太年轻}（\BibleRef{耶 1:6}）。而全能天主寻找宣讲者时说：\BibleQuote{我将派遣谁呢？谁肯为我们去呢？}\Person{依撒意亚}主动献身说：\BibleQuote{我在这里，请派遣我！}（\BibleRef{依 6:8}）。看，两人外表发出了不同的话语，却都源自同一爱泉。\par

因为爱德确实有两条诫命：即爱天主和爱近人。因此，\Person{依撒意亚}渴望以行动生活裨益近人，便渴求宣讲的职务；而\Person{耶肋米亚}渴望以默观生活紧依造物主之爱，便拒绝受派宣讲。那么，前者可赞地渴望，后者可赞地退缩：后者免得因说话而失去静默默观之益；前者免得因沉默而感觉勤奋工作之失。但二者需仔细留意：那位拒绝者最终并未抗拒，那位渴望被派遣者则看见自己先前已被祭坛上的火炭洁净；如此，未经洁净者不敢接近神圣职务，天上恩宠所拣选者也不可假借谦卑而骄傲地拒绝。\par

此外，我在你的信中发现你极其渴求心灵的宁静，渴望没有纷扰的安详思绪。但我不知你的弟兄之尊如何能达到这一点。因为一个承担了船只领航职责的人，越是远离海岸，就越需警醒：时而从迹象预知风暴来临；时而当风暴临到时，若风暴不大，便直线越过；若狂风大作，便侧航避开它们的冲击；当船中所有无责之人都安歇时，他常常独自守夜。再者，既已承担牧职之担，你怎能有安宁的思绪呢？因为经上记载：\BibleQuote{看，巨人在水下呻吟}（\BibleRef{约 26:5}）。按\Person{若望}的话：\BibleQuote{水是指民众}（\BibleRef{默 17:15}）。巨人在水下呻吟，意思是：在此世位高权重者，如同身体庞大，因承担治理民众之责而感受更沉重的磨难。但是，若圣神的德能吹拂受折磨的心灵，那么为以色列子民所行的事，便在我们身上属灵地实现。因为经上记载：\BibleQuote{但以色列子民却在海中干地上走过}（\BibleRef{出 14:29}）。上主藉先知许诺说：\BibleQuote{你从水中经过时，我必与你同在，河川不淹没你}（\BibleRef{依 43:2}）。河川淹没那些人——这世界的俗务以心灵纷扰扰乱他们。但蒙圣神恩宠扶持的人，心灵安稳地渡过河水，且不被河川淹没，因为他在民众之中行走，却不让心灵的头颅沉没于世务之中。\par

我，虽不堪当，却已来到治理之位，也曾决意寻觅一处退隐之所；但见天主的上智安排与我对立，我便将心中的颈项顺服于造物主的轭下。尤其深思这一点：若无天主的恩宠，任何隐密之处都无法拯救灵魂；这一点我们有时可见，即便圣人也会迷失。因为罗特在那腐化的城中曾是义人，却在山上犯了罪\EditorialNote{创世纪 19}。然而，何必提这些例子？我们知道有更重大的事。有什么比乐园更令人愉悦？有什么比天堂更安全？然而人从乐园，天使从天庭，皆因犯罪而坠落。因此，当寻求祂的能力，恳求祂的恩宠；没有祂，我们无论在哪里都免不了过错；有了祂，我们无论在哪里都不会缺少义德。那么，我们当留心，勿让思想的纷扰胜过我们的心智；因为要完全摆脱它绝无可能。但凡处于治理之位者，难免有时需思虑尘世之事，并关心外在事务，使托付给他的羊群能够生存以完成其所当行。但必须极其谨慎，使这种关心不越过适当限度；当它合法地进入心中时，不可任其过度。因此，通过厄则克耳说得对：\BibleQuote{司祭不可剃头，也不可让头发长长，只可修剪头发。}\BibleRef{则 44:20}。因为头上的头发象征什么，岂非心中的思想？它们从脑际无觉地升起，代表现世生活的忧虑；这些忧虑因疏忽的知觉，有时无礼地袭来，仿佛不知不觉地滋长。既然所有管辖他人者，理当有外在的焦虑，却不可过度投身其中，所以司祭被正当地禁止剃头或留长头发，使他们既不完全割舍为属下生活所需的肉性思想，也不让它们过度生长。那里也恰当地说：\BibleQuote{只可修剪头发；}\BibleRef{则 44:20}——意思是：暂时职务的焦虑既应如必需而进行，又当及时剪短，免得长到不合度。因此，通过外在供应的管理，身体的生命得以保全；同时，内心的专注不因同样焦虑的过度而受阻；司祭头上的头发得以覆盖皮肤，又剪短以免遮蔽眼睛。\par

此外，我们以全备的信德收到了你们致我们的书信，并感谢全能的天主，祂藉信友们的彼此宣认，保护那从上到下浑然天成、没有缝线的长衣——即祂的教会——在恩宠的合一中免受一切错谬的撕裂；面对这败亡世界众多罪恶的洪水（可谓），祂用许多木板建造了一只方舟，使全能天主的选民得以在其中保存生命。因为，当我们转寄信仰的宣认给你们，而你们向我们显示爱德时，我们在圣教会中做什么呢？岂非用柏油涂抹方舟，免得任何错谬的波浪进入，杀害所有属神的人，如同人，以及属血肉的人，如同牲畜？\par

但是，你们既然明智地宣认了正确的信仰，无疑地，你们更当谨慎地保持心灵的平安，因为真理说：\BibleQuote{你们当有盐在你们内，彼此和平相处。}\BibleRef{谷 9:50} 宗徒保禄也劝诫说：\BibleQuote{竭力保持圣神的合一，以和平的联系。}\BibleRef{弗 4:3} 他又说：\BibleQuote{务要与众人和平相处，并追求圣德；无圣德，没有人能见主。}\BibleRef{希 12:14} 如果你们远离一个亵圣名称的骄傲，你们便果真与我们有了这平安，正如那位外邦人的导师所说：\BibleQuote{弟茂德啊，要保管好所托付给你的，避免亵圣的新奇言论。}\BibleRef{弟前 6:20} 因为这实在糟糕：那些被立为谦逊宣讲者的人，竟因虚荣名称的夸耀而自傲，而真正的宣讲者说：\BibleQuote{至于我，我决不夸耀，只夸耀我们的主耶稣基督的十字架。}\BibleRef{迦 6:14} 那么，真正可夸耀的人，不是夸耀暂时的权能，而是为基督的名夸耀祂的苦难。因此，我们从心底拥抱你们；我们如此承认你们为司祭，如果你们弃绝言语的虚妄，以圣洁的谦逊占有圣洁的地位。因为看，我们被这亵圣的称呼所困扰，心中怀有并言语表达绝非轻微的怨言。但你们的兄弟情谊知道真理如何说：\BibleQuote{若你在祭台前献礼物时，想起你的弟兄有什么怨你的事，就把礼物留在祭台前，先去与你的弟兄和好，然后再来献礼物。}\BibleRef{玛 5:23} 于此当思：虽然一切过失可由祭献除掉，但引起他人心中怨尤的恶事如此之大，以致主不接受那本应除罪的祭献。那么，要迅速从你们心中除去怨尤之因，使全能天主能悦纳你们奉献的祭品。\par

首先，正确而准确地宣认了正信，我们发现在你所视为受至圣大公会议谴责的人中，你谴责了某一个\Person{欧多克修}；我们在拉丁语中，无论是在会议文件中，还是在诸位真福主教\Person{厄丕法尼}、\Person{奥斯定}、\Person{斐拉斯特}的著作里，都未找到提及他的名字，我们深知他们是反驳异端的主要辩士。如今，若有任何一位公教圣父确实谴责他，我们无疑会随从他们的意见。然而，如果你在教区会议信函中，也欲谴责那些除了圣会议之外而在圣父们的著作中受到谴责的人，那么你的兄弟情谊所列举的人就太少了；但如果你欲谴责大公会议所摒弃的人，那么因著此一人便显得太多了。但在这一切之中，务必记住：为使我们可以自由地宣认真信，并以和平与和谐处理一切当行之事，我们应为最安宁的主上及其后嗣的生命不断祈祷，愿全能的天主将野蛮民族践踏于他们脚下，赐予他们长久而幸福的生命，以便通过基督化的帝国，使在基督内的信德得以统御。\par

\SourceRecord{Translated by James Barmby.；Nicene and Post-Nicene Fathers, Second Series；Vol. 12.；Edited by Philip Schaff and Henry Wace.；Buffalo, NY: Christian Literature Publishing Co.,；1895.；<http://www.newadvent.org/fathers/360207004.htm>.}

% structure: p0001 source=h1 target=section reason=page-title

\section{第七卷，第五封信}

\Emph{致基里亚库斯主教。}\par

\Person{格列高利}致\Person{基里亚库斯}，\Place{君士坦丁堡}主教。\par

当昔日我代表宗座驻跸皇城时，已熟知尊座的善德。今闻灵魂之牧职托付于你，我深为欣慰。我虽不配，仍以全力向全能天主恳祷，愿祂甚至在你内增加祂的恩宠，并使你为永恒家乡聚集灵魂的增益。然而，你自言对托付于你的这项工作力有未逮，我们深知首要之德即是承认软弱；由此我们推断，你能善尽所领受的职务，因我们看出你出于谦卑而承认自身的软弱。因我们众人皆软弱；但更有软弱者，是无力思虑自身软弱之人。然而，至福的弟兄，你之所以强壮，正因不信赖自身之力，而信赖全能天主的权能。\par

然而，我无法以书信言词表达我心与尊位何等紧密相连。但我祈求全能天主藉祂恩宠的恩赐，倍增我们之间同样的爱德，并除去一切冒犯的缘由，免得至圣教会因真信仰的宣认而合一，又因信友们心灵的结合而凝聚，却因司铎间的争执而受损——愿天主禁止。无论如何，我所说的一切，我所言的一切，虽是针对某些人的骄傲行为，但藉全能天主的恩惠，我从未放弃内心爱德的保管；如此，在外执行属于正义之事时，绝不忽视内心属于爱与仁慈之事。也请你常报以我的爱，守护属于和平与仁慈的事；如此，我们保持同心，不容任何不和介入我们之间，更能因心灵的合一而获得向主所求的。\par

此外，我将\Person{若望}，\Place{卡尔西登}司铎，以及\Person{依扫利亚的亚大纳削}托付于尊座，免得有人以暗中诬蔑挑拨你反对他们；因我已彻底审查他们的信仰，发现他们在宣信上纯正，这些宣信也已以书面形式呈交。\par

愿天主圣三以祂的手保护你，使你常在保管灵魂之事上警醒谨慎，终使你在永恒报应中堪当获得冠冕，不仅为你自己的工作，也为你的属下的改善。\par

\SourceRecord{Translated by James Barmby.；Nicene and Post-Nicene Fathers, Second Series；Vol. 12.；Edited by Philip Schaff and Henry Wace.；Buffalo, NY: Christian Literature Publishing Co.,；1895.；<http://www.newadvent.org/fathers/360207005.htm>.}

% structure: p0001 source=h1 target=section reason=page-title

\section{第七卷，第六封信}

\Emph{致摩里士皇帝。}\par

额我略致摩里士皇帝。\par

全能的天主，祂使陛下成为教会和平的守护者，也以同一信德保守您，这信德藉着司铎之间的合一而得以保存；当您谦卑地将心屈服于天上慈爱的轭下时，便藉着天主的恩宠成就了您以勇力践踏仇敌的事。这绝非小可之事：当可敬的若望离世后，陛下长久迟疑，并在敬畏全能天主的顾虑中稍加延宕，为的是使天主的事业能以应有的敬畏妥善安排。因此，我也认为我的兄弟和同僚司铎基亚古斯极适于牧职治理，因为陛下长期的审慎考虑将他提升至此位。我们皆知他在教会事务的管理上长久以来是何等勤勉与干练。因此，我毫不怀疑这是出于天主的安排：那妥善管理小事者，当适当地承担大事，并从事务管理过渡到灵魂的治理。故此，我们在一切祈祷中恳求全能的天主，将这件善工回报给主上们的安宁及其虔诚的后裔，不仅在现世，也以永久的赏报，并恩赐我前述的兄弟和同僚司铎，既被委派管理主的羊群，能在照顾灵魂上显出充分的关怀，使他能无可指摘地纠正属下不当之处，扶持正直，使之增长；以致陛下对他的判断不仅在人前蒙受称许，也在至高威严的眼前蒙受悦纳。\par

可敬的乔治长老和戴奥多勒执事，念及我主们的命令及冬季的临近，我不允许他们在此城中耽搁。\par

\SourceRecord{Translated by James Barmby.；Nicene and Post-Nicene Fathers, Second Series；Vol. 12.；Edited by Philip Schaff and Henry Wace.；Buffalo, NY: Christian Literature Publishing Co.,；1895.；<http://www.newadvent.org/fathers/360207006.htm>.}

% structure: p0001 source=h1 target=section reason=page-title

\section{第七卷，信函七}

\Emph{致\Person{伯多禄}、\Person{多米提安}与\Person{埃尔皮迪乌斯}。}\par

\Person{额我略}致\Person{伯多禄}、\Person{多米提安}与\Person{埃尔皮迪乌斯}主教。\par

我极为喜悦，因为你们欣然接纳了最神圣的\Person{基里亚库斯}的祝圣，他是我的兄弟和同司铎。既然我们从宗徒\Person{保禄}的宣讲中得知：\BibleQuote{若一个肢体喜乐，所有的肢体都一同喜乐}\BibleRef{格前12:26}，所以你们必须明白，我在这事上是怀着何等大的欢欣与你们一同喜乐，因为在这事中，不仅是基督的一个肢体喜乐，而是许多肢体都喜乐了。然而，就我所能粗略阅读的贵弟兄们的信函而言，巨大的喜乐使你们陷入了对这位我兄弟的过度赞扬。因为你们说他出现在教会中如同太阳，以致你们全都呼喊：\BibleQuote{这是主所造的一天；让我们欢欣鼓舞，并在其中喜乐}\BibleRef{咏117:24}。然而，这确实是来世生命的应许，因为经上说：\BibleQuote{义人将如太阳般闪耀}\BibleRef{玛13:43}。因为，无论任何人在任何德行上多么卓越，在现世生活中，他又怎能如太阳般闪耀呢？在这现世生活中，\BibleQuote{可朽坏的身体压得灵魂沉重，地上的居所压低了那思虑万事的头脑}\EditorialNote{\BibleRef{智9:15}}；在这生活中，\BibleQuote{我们看见在我们肢体内另有一条法律，与理智的法律交战，并将我们掳去，顺从那在我们肢体内的罪的法律}\BibleRef{罗7:23}；在这生活中，\BibleQuote{连我们自己身上也带着死亡的判决，为使我们不依靠自己}\BibleRef{格后1:9}；在这生活中，先知也大声呼喊：\BibleQuote{恐惧和战栗降临于我，黑暗覆盖了我}\BibleRef{咏54:6}？因为经上也写着：\BibleQuote{明智的人如太阳般持守；愚笨的人如月亮般变化}\BibleRef{德27:12}；这里太阳的比喻不是用于其光明的辉煌，而是用于行善的恒心。但他的祝圣的良好开端，你们尚不能就其恒心而加以赞扬。至于你们说你们呼喊\BibleQuote{这是主所造的一天}，你们应当想想这话是对谁说的。因为前面是这样说的：\BibleQuote{匠人弃而不用的石头，反而成了屋角的基石。这是主的作为，在我们眼中神妙莫测}\BibleRef{咏117:22}。而关于这同一块石头，随即又加上了\BibleQuote{这是主所造的一天}。因为那因建筑的坚固而被称作石头的，因光照的恩宠而被称作白昼，祂也是被造的，因为祂降生成人。我们奉命在祂内欢欣鼓舞，因为祂以祂美善的光辉在我们内战胜了我们错误的黑暗。那么，在赞美一个受造物时，就不应该使用那唯独适用于造物主的表达。\par

但我为何要指责这些事呢？我知道喜乐如何带走人心。因为你们的爱德在你们内产生了极大的欣喜，而心中的欣喜随之得到了舌头的赞美。既如此，爱德所发现的赞美不能算作过错了。但是，关于我最神圣的兄弟，本应简要地说些我能满意接纳的话，因为我认识他，他长期以来尤其向我证明了此点：尽管他忙于如此多的教会管理事务，他却能在喧嚣的群众中保持一颗平静的心，并始终以温和的态度约束自己。这确实是对一个伟大而坚定的心灵的不少称赞——在没有被事务的纷扰所扰乱。\par

此外，你们的弟兄团体应恒心不断地祈祷，使全能的天主在我们的前述兄弟和同司铎身上守护那已经开始良好的，并一直引领他走向更美善。这应当是你们最神圣的人以及归属于他的百姓常做的祈祷。因为统治者和百姓的功过彼此如此相连，以致臣民的生活常因在上者的过失而变坏，而牧者的生活也常因百姓的恶行而衰败。因为那在上者的恶行对在下者造成的极大伤害，法利塞人就是明证，关于他们经上写道：\BibleQuote{你们封闭了天国，不让人进去。你们自己不进去，也不让那些正要进去的人进去}\BibleRef{玛23:13}。而百姓的过失对牧者生活的严重伤害，我们从达味所做的事可知 \BibleRef{列下2:24}。因为他，受天主见证赞美，深知天上的奥秘，却被隐藏的骄傲之瘤膨胀，在点数百姓时犯了罪；然而刑罚却因达味的罪落在了百姓身上。这是为什么？因为事实上，按照臣服百姓的功过，统治者的心被安置了。现在，公义的审判者通过惩罚那些导致他犯罪的人，来谴责那罪人的过失。但是，因为他自己因骄傲自意志而并非无罪，他也亲自接受了其过失的惩罚。因为那猛烈的愤怒，既打击了百姓的身体，也以内心深处的悲伤击倒了百姓的统治者。因此，你们要彼此考虑这些事；既然那被置于你们和百姓之上的人应为所有人代祷，那么你们所有人都应为他的言行祈祷，使在全能的天主面前，你们能通过效法他而获益，而他也能因你们的功过得到帮助。此外，让我们同心合意，倾尽所能，不断以极大的痛哭为我们最安宁的君主及其虔诚的后代祈祷，使护守的恩宠保佑他们的生命，并使诸国的颈项臣服于基督信仰的帝国。\par

\SourceRecord{Translated by James Barmby.；Nicene and Post-Nicene Fathers, Second Series；Vol. 12.；Edited by Philip Schaff and Henry Wace.；Buffalo, NY: Christian Literature Publishing Co.,；1895.；<http://www.newadvent.org/fathers/360207007.htm>.}

% structure: p0001 source=h1 target=section reason=page-title

\section{第七卷，第十一封}

\Emph{致\Person{鲁菲努斯}，\Place{以弗所}主教。}\par

\Signature{格列高利致\Person{鲁菲努斯}等。}\par

您过去友谊行为中的爱德，促使我们以这封信问候您。因为我们从有关您健康的报告中得知一切安好，深感欣慰，充满喜乐。但既是如此，我们恳求全能的天主，既然在此生——它仿佛是未来生命的影子——祂已恩赐您享受身体暂时的安康，同样在那真正生命的属天家乡，愿祂使我们因您灵魂的圆满救恩而共同欢欣感谢。现在，此信的携带着希望通过我们的信被推荐给您，我们问他是否如一个文书应学的那样识文断字，他回答说他对此一无所知。那么，关于他，我不知还能向您推荐什么；唯愿您关心他的灵魂，以牧者的热忱看顾他，如此，因他不能阅读，您的口舌就是他的书，愿他在您良好的宣讲与工作中看到应当效法的。因为活生生的声音通常比应付差事的阅读更能吸引人心。但是，作为他的师长，您向内供给这属灵教导的同时，也不可缺少对他外在的照顾，为使他藉此帮助渴慕属灵之事，且以免若轻视这种帮助，您将不再有可宣讲的对象。\par

\SourceRecord{Translated by James Barmby.；Nicene and Post-Nicene Fathers, Second Series；Vol. 12.；Edited by Philip Schaff and Henry Wace.；Buffalo, NY: Christian Literature Publishing Co.,；1895.；<http://www.newadvent.org/fathers/360207011.htm>.}

% structure: p0001 source=h1 target=section reason=page-title

\section{\\WorkTitle{第七卷，第十二封书信}}

\\Emph{致女院长雷斯佩克塔}\par

\\Person{额我略}致\\Place{高卢}的\\Place{马西利亚}（\\Place{马赛}）的女院长\\Person{雷斯佩克塔}。\par

虔诚愿望的要求应当由相应的结果来实现，以便所求的恩惠能有效地获得，并且真诚的热忱能值得赞扬地彰显出来。因此，根据我们的子女\\Person{狄纳纽斯}和\\Person{奥勒利亚娜}的请求——他们在宗教热忱中显示，已通过连接建筑物将\\WorkTitle{圣加西亚诺}隐修院与他们在其住处合并——我们已批准允许以下特权：我们规定，在上述隐修院的女院长去世时，不得由外人，而应由修会从自身成员中自行选立一人；此人（只要被认为配得上这一职务）应由同一地方的主教祝圣。此外，关于同一隐修院的财产和管理，我们规定主教或任何圣职人员均不得有任何权力；但下令这些事应完全归属你的关怀，或在你之后同一地方的院长之责。如果在上述隐修院的圣人周年纪念日或奉献日，主教前往那里举行弥撒的神圣礼仪，他的职务仍应如此执行：除非在上述日子他在那里举行弥撒礼仪时，不得放置他的座位。当他离开时，他的座位应同时从同一祈祷所移走。但在所有其他日子，弥撒礼仪应由同一主教所指定的司铎进行。\par

此外，关于天主婢女们的生活和行为，或关于在上述隐修院中被立为女院长者，我们责成主教在敬畏天主中给予密切关注；以便如果其中任何居住者因其过错而应受惩罚，他本人可根据神圣教会法规的力度惩罚该过犯。这些事既由我们下令和批准，你在管理你的修会时，应努力在所有方面如此热切地专注，以致恶毒敌人的恶意在那里找不到任何可污染之处。因此，所有包含在这份训令文件中的这些事，我们命令在基督的保护下，在你的隐修院中永远由所有人和在所有方面遵守，以便所允许的特权的恩惠始终坚定且不可侵犯。十一月，小纪第十五。\par

\SourceRecord{Translated by James Barmby.；Nicene and Post-Nicene Fathers, Second Series；Vol. 12.；Edited by Philip Schaff and Henry Wace.；Buffalo, NY: Christian Literature Publishing Co.,；1895.；<http://www.newadvent.org/fathers/360207012.htm>.}

% structure: p0001 source=h1 target=section reason=page-title

\section{第七卷 第十三封书信}

\Emph{致\Person{福图纳图斯}主教。}\par

\Person{额我略}致\Place{法努姆}主教\Person{福图纳图斯}。\par

如同任何人出售圣器（除非在法律和神圣教规许可的情况下）应受谴责和惩罚一样，若出于赎回俘虏的慈悲目的而出售，则不应受责备或惩罚。因此，我们从你的主教兄弟得知，你为赎回俘虏借了钱，而无力偿还，为此希望凭我们的权威出售一些圣器——在这种情况下，鉴于法律和教规的法令都许可，我们认为宜予批准，并准许你出售这些圣器。但是，为避免出售可能给你带来任何恶感，应在我们的\OriginalTerm{辩护者}{defensor}若望在场的情况下，出售价值足以还清债务的圣器，并将价款支付给债权人，以使此事经过如此处理，既不会使债权人因借钱而受损，也不会使你兄弟现在或将来遭到恶感。\par

\SourceRecord{Translated by James Barmby.；Nicene and Post-Nicene Fathers, Second Series；Vol. 12.；Edited by Philip Schaff and Henry Wace.；Buffalo, NY: Christian Literature Publishing Co.,；1895.；<http://www.newadvent.org/fathers/360207013.htm>.}

% structure: p0001 source=h1 target=section reason=page-title

\section{第七卷，第十五封信}

\Emph{致司铎\Person{乔治}。}\par

\Person{额我略}致\Place{君士坦丁堡}教会司铎\Person{乔治}及执事\Person{德奥多尔}。\par

念及您的良善与爱德，我深感自责，竟让您过早返回；但见您一再恳求准予离去，我思忖您的爱德与我们久留或许会有重大不便。然而，当得知您因冬季而旅途耽搁许久，我承认，我后悔让您如此匆忙离去。因为，若您的爱德未能完成预定行程，那么您与我共处而非远离我，岂不是更好？\par

此外，在您离去后，我从我至爱的儿子执事们所得消息中获悉，您的爱德曾说：我们的全能之主及救主\Person{耶稣基督}，当他降至地狱时，拯救了所有在那里承认他为天主的人，并使他们免于应受的痛苦。关于此事，我愿您的爱德持有截然不同的看法。因为，当他降至地狱时，他仅藉其恩宠拯救了那些既相信他将降临，又在其生活中遵守其诫命的人。显然，在主降生成人之后，即使是持守信仰却无信仰生活的人，也不能得救；因为经上记载：\Emph{\BibleQuote{他们承认认识天主，但在行为上却否认祂}}\BibleRef{铎 1:16}。\Person{若望}也说：\Emph{\BibleQuote{那说认识祂，却不遵守祂诫命的，是撒谎者}}\BibleRef{若一 2:4}。主的兄弟\Person{雅各伯}亦写道：\Emph{\BibleQuote{没有行为的信德是死的}}\BibleRef{雅 2:20}。倘若如今有信仰者不行善功不得救，而无信仰、被弃绝者却因主降至地狱而无善功得救，那么那些从未见过主降生成人者的命运，竟优于那些生在主降生成人奥秘之后的人了。然而，主亲自对门徒们作证，说这样的话或持这样的想法是何等愚蠢，祂说：\Emph{\BibleQuote{许多先知和君王愿意看你们所看见的，却没有看见}}\BibleRef{玛 13:17}。但是，为不使您的爱德纠缠于我个人的争论，请听\Person{菲拉斯特}在其所著关于异端的书中关于此异端的论述。他的原话如下：\Quote{说主降至地狱，死后向所有已在那里的人宣告自己，使他们在那里承认祂而得以获救的人，是异端者；因为这正相反于\Person{达味}先知所说的\BibleQuote{但在阴府里，谁能承认祢？}\BibleRef{咏 6:6}以及使徒所说的\BibleQuote{凡没有法律而犯了罪的，也必没有法律而丧亡}\BibleRef{罗 2:12}。}真福\Person{奥思定}在其所著关于异端的书中，也与他的话一致。\par

因此，思考这一切，请只持守真信仰通过天主教会所教导的：即主降至地狱时，只救拔了那些在肉身生活时，祂藉其恩宠在信德和善行中保全的人。因为祂在福音中所说的\BibleQuote{当我从地上被举起来时，我要吸引众人归向我}\BibleRef{若 12:32}，是指一切被选者。因为一个因恶行而与天主分离的人，死后不可能被吸引到天主那里。愿全能的天主护佑您，使您无论身在何处，都能在灵魂和身体上感受到祂恩宠的助佑。\par

\SourceRecord{Translated by James Barmby.；Nicene and Post-Nicene Fathers, Second Series；Vol. 12.；Edited by Philip Schaff and Henry Wace.；Buffalo, NY: Christian Literature Publishing Co.,；1895.；<http://www.newadvent.org/fathers/360207015.htm>.}

% structure: p0001 source=h1 target=section reason=page-title

\section{第七卷，第十七封信}

\Emph{致\Person{萨比尼安努斯}主教}\par

\Person{额我略}致\Place{雅德拉}主教\Person{萨比尼安努斯}。\par

如果你当初费心慎重考虑教会管理的规则和古代习俗的次序，那么非法的僭越之罪就不会潜入你身，他人也不会因你的罪而招致危险。如今毫无疑问，你知道关于\Person{马克西慕斯}的一些事情传到我们耳中，这些事对他晋升司铎圣秩构成了不小的障碍，我们并未给予同意，而且我们的意愿是，直到关于所说的事情有足够的澄清之前，他不应得到他所追求的地位。但是，当你本应尽一切可能遵守这一点时，反而发生了这样的事：他怀着盲目心灵的贪婪攫取主教职位，使你不经意地倾向于无视我们的禁令而偏袒他。但是，即使如此，为了使报告给我们的事情不至于未经审查，他曾被我们用书信召唤到这里。当他如此乖戾，拖延不来时，我们曾在多封信中告诫他，在禁止领圣体的处罚下，抛开一切借口，赶快到他这里来为自己辩护；但他宁愿接受绝罚也不愿显示服从。因此结果是（说来可怕），他那乖戾性情的邪恶用他自己的沉沦连累了他人。然而现在，既然我们获悉你已不赞同他的邪恶，我们通过本函劝告你（如此，即使为时已晚，你与他断绝关系将有益于你的灵魂），你今后既不可与他共融，也不可在弥撒圣祭的庄严礼仪中提及他的名字；而且你应立即前来见我们，不要延迟，并且还要带上其他你能带来的主教或其他有圣职的人士，以便（在充分审查案情后）如果需要的话，你的赦免可以合宜而体面地得以实现，并且那些陷入类似鲁莽之罪的人，在圣\Person{伯多禄}、宗徒之长的帮助下，通过合乎\Person{基督}心意的安排，被召回救恩之路。再者，任何可能来找我们的主教或宗教人士要知道，他不会遭受任何偏见或不公，一切都会在完全查明真相后安排得合乎我们救赎主的心意；其目的是，即使从我们处理事情的方式本身，在主应允之下，也能向所有人表明，我们并非出于对任何人的私怨，而是出于对天主的热情以及调整教会秩序的热情。\par

\SourceRecord{Translated by James Barmby.；Nicene and Post-Nicene Fathers, Second Series；Vol. 12.；Edited by Philip Schaff and Henry Wace.；Buffalo, NY: Christian Literature Publishing Co.,；1895.；<http://www.newadvent.org/fathers/360207017.htm>.}

% structure: p0001 source=h1 target=section reason=page-title

\section{第七卷 第十九封信}

\Emph{致马里尼亚努斯总主教。}\par

额我略致拉文纳总主教马里尼亚努斯。\par

你的弟兄早已知晓，由于我们祝圣的那位司铎身患众所周知的疾病，阿里米努姆教会至今如何被剥夺了牧灵治理。如今，我们因该地居民的恳求而感动，曾多次劝告他，若能藉主的帮助摆脱那使他滞留的头疾，就应返回他的教会；自准假以来，已等待了四年。当来自该处的圣职人员和市民央求我们并催促我们时，我们急切地劝告他与他们一同返回，主既帮助他，若他能够如此；但他却通过书面请求恳求我们，既然因这缠累他的疾病，他无论如何不能承担该教会的治理，也无法履行他所承担的职务，我们应为此同一教会祝圣一位主教。因此，鉴于托付给我们照顾所有教会的职责，迫使我们务必使信友羊群不再缺少牧灵的护守，且因他们的恳求以及他因自身无能而放弃职务，我们决定应为阿里米努姆同一教会祝圣一位主教；并已按惯例发出我们的谕令，尚未忘记劝告同一教会的圣职人员和人民，以便他们能以一致的选择为自己选出一位主教。因此，我们劝告你的弟兄，凡众人一致同意的人（正如他们自己也请求获准这样做），你应召他到你面前；并谨慎地从各方面考察他。如果，藉主的恩宠，在他身上找不到\WorkTitle{七经}经文中以死惩罚的任何事，且根据可靠人士的报告，他的生活应得到你的认可，那么就将他连同选举的证明文件送来，并附上你的作证信，以便该同一教会的主教，在主的安排下，由我们祝圣。\par

\SourceRecord{Translated by James Barmby.；Nicene and Post-Nicene Fathers, Second Series；Vol. 12.；Edited by Philip Schaff and Henry Wace.；Buffalo, NY: Christian Literature Publishing Co.,；1895.；<http://www.newadvent.org/fathers/360207019.htm>.}

% structure: p0001 source=h1 target=section reason=page-title

\section{第七卷，第二十封书信}

\Emph{致\Place{里米尼}的圣职人员及人民}。\par

\Person{格里高利}致圣职人员等。\par

我们的牧职责任催迫我们以焦虑的关怀扶助任何缺乏司铎治理的教会。因此，既然你们的教会因其司铎如你们所知身患疾病而长期缺乏牧者管理，我们受你们恳求所动，并未停止劝诫所述主教：倘若他感到自己已从该病中康复，则应重新履行他所领受的司铎职务。而他经我们一再警告，如今在同样疾病的压力下，通过书面恳请向我们表示，因该病之故，他完全无法升任所述教会的治理或其所领受的职务。因此，被此人之无望境况所迫，我们认为必须为整顿你们的教会而考虑。我们因而劝勉你们全体一致，无喧哗动乱，在主的助佑下选择一位既不为可敬教规所不容，又堪当如此重大职务的司铎来管理你们。当选后，他应带着由全体签署并经巡视员书面认可的法令前来到我们这里受祝圣，以便你们的教会通过主的安排得以拥有自己的司铎。\par

我们也希望，你们全体一致所选之人须立即带往\Place{拉文纳}我们的兄弟及同主教\Person{马里尼亚诺斯}处，使他接受彻底考核与检验后，在来我们这里时亦能以自己的证词支持他。\par

\SourceRecord{Translated by James Barmby.；Nicene and Post-Nicene Fathers, Second Series；Vol. 12.；Edited by Philip Schaff and Henry Wace.；Buffalo, NY: Christian Literature Publishing Co.,；1895.；<http://www.newadvent.org/fathers/360207020.htm>.}

% structure: p0001 source=h1 target=section reason=page-title

\section{第七卷，第二十三封书信}

\Emph{致\Person{福图纳图斯}与\Person{安特弥乌斯}}\par

\Person{额我略}致\Person{福图纳图斯}主教、\Person{安特弥乌斯}\OriginalTerm{守护者}{defensori}。\par

持信人\Person{卡泰路斯}告知我等：其妹曾与一位\Person{斯德望}订婚，因天主的仁慈感动她，已在\Place{那不勒斯}的一所隐修院内皈依，而同一\Person{斯德望}不正当地扣留了属于她的一所房屋及若干其他物品。再者，鉴于法律法令\EditorialNote{法律引文：Caus. 17, q. 2, c. 28}规定，若订婚女子愿意皈依，不应受任何损失，故请你们，偕同副执事\Person{安特弥乌斯}，勤加查究，探明事实。若你们，如我们所闻，发现上述\Person{斯德望}不正当地占有一所房屋或任何其他财物，则应以你们的劝诫切迫警告他，立即且毫无争论地归还不当扣留之物，不得以任何借口拖延归还非属己有之物。若你们见他忽视你们的劝诫，请将此事告知我们，并详述案件实情，以便在案情明了后，可按照公平，以其他手段强制他做出归还；他因顾全廉耻而轻蔑自行归还之事。我们将这持信人托付于你们，劝请你们不再让他因此事拖延而受苦。\par

\SourceRecord{Translated by James Barmby.；Nicene and Post-Nicene Fathers, Second Series；Vol. 12.；Edited by Philip Schaff and Henry Wace.；Buffalo, NY: Christian Literature Publishing Co.,；1895.；<http://www.newadvent.org/fathers/360207023.htm>.}

% structure: p0001 source=h1 target=section reason=page-title

\section{第七卷·第二十五封书信}

\Emph{致格莱高利亚。}\par

\Person{格列高利}致\Person{格莱高利亚}，\OriginalTerm{内室女官}{cubiculariæ}，属\Person{奥古斯塔}。\par

我收到了你甜美的信，你在信中一直努力指责自己犯了许多罪。但我知道你热切地爱慕全能的主，我坚信祂的仁慈，那关于一位圣妇的判词也适用于你：\BibleQuote{她的许多罪都被赦免了，因为她爱得多}\BibleRef{路 7:47}。随后的事也表明她们如何被赦免：她坐在主的脚前，听祂口中的话\BibleRef{路 10:39}。因为她被提升到默观生活，超越了她的姐妹玛尔大仍从事的积极生活（\BibleRef{路 10:40}）。她也热切地寻求被埋葬的主，俯身于坟墓旁，却没有找到祂的身体。甚至当门徒们离去时，她仍站在坟墓门前，她寻找那视为死者，却得以见到祂活着，并向门徒宣告祂已复活。这是天主仁慈的奇妙安排：生命由妇女的口宣告，因为死亡最初也是在乐园中由妇女的口尝到的。另有一次，她与另一位玛利亚在复活后见到主，并抱住祂的脚。我恳求你，想想那是怎样的手抱住了谁的脚。那位曾是城中的罪妇，那双曾被不洁污染的手，触摸了那坐在圣父右边，超越众天使者的脚。如果我们能够，让我们估量那属天的仁慈之心是怎样的，一个因罪陷入漩涡深处的妇女，竟因恩宠藉着爱之翼被如此高举。亲爱的女儿，这应验了，这应验了先知论及圣教会这个时期所预言的：\BibleQuote{在那一天，达味之家必成为敞开的水泉，为洗涤罪人与不洁者}\BibleRef{匝 13:1}。因为达味之家成为我们罪人的敞开之水泉，我们藉着现在通过达味之子、我们的救主所显示的仁慈，被洗净我们的不洁。\par

至于你的甜美在信中所附加的，即你将不断催促我，直到我写信说已向你启示你的罪被赦免了，你要求了一件困难甚至无益的事：困难是因为我不配得到启示；无益是因为你不应该对自己的罪感到安心，除非在你生命末日不能再哀哭时。但在那一天到来之前，你应当永远怀疑、永远恐惧，害怕过犯，并用每日的眼泪洗净它们。诚然，宗徒保禄已升到三重天，被提到乐园，听见不可言说的隐秘话语\BibleRef{格后 12:2}，但仍恐惧地说：\BibleQuote{我克制自己的身体，使它服从，唯恐我给别人宣讲，自己反而被淘汰}\BibleRef{格前 9:27}。被提到天上的人仍然恐惧，而仍在世上生活的人岂能已不愿恐惧？最亲爱的女儿，请记住，安于现状常常是疏忽的母亲。因此，你今生不应有安于现状，以免变得疏忽。因为经上记载：\BibleQuote{常存敬畏的人是有福的}\BibleRef{箴 28:14}。又说：\BibleQuote{你们当敬畏地事奉上主，战战兢兢地欢乐}\BibleRef{咏 2:11}。简言之，你在今生必须战战兢兢，以便将来在安全的喜乐中永无止境地欢乐。愿全能的天主以祂圣神的恩宠充满你的灵魂，并在你每日祈祷时所流的泪之后，带你进入永恒的喜乐。\par

\SourceRecord{Translated by James Barmby.；Nicene and Post-Nicene Fathers, Second Series；Vol. 12.；Edited by Philip Schaff and Henry Wace.；Buffalo, NY: Christian Literature Publishing Co.,；1895.；<http://www.newadvent.org/fathers/360207025.htm>.}

% structure: p0001 source=h1 target=section reason=page-title

\section{第七卷，第二十六书}

\Emph{致\Person{特奥克蒂斯塔}，贵族}\par

\Person{格列高利}致\Person{特奥克蒂斯塔}等。\par

阁下垂临如此纷繁的事务之中，却满盈圣经的果实，并不断渴慕永恒的喜乐——为此我向全能的天主献上极大的感谢，因为在您身上，我看到了关于蒙选祖先所记载的这句话得以实现：\BibleQuote{以色列子民却在海中干地上行走} \BibleRef{出 15:19}。但另一方面，\BibleQuote{我陷入深渊，洪水卷没了我} \BibleRef{咏 68:3}。而您，如我所见，却在世俗事务的波涛中干足行走，迈向预许之地。那么，让我们感谢那提升所充满之心的圣神；祂在人群的喧嚣中使灵魂孤寂；在祂面前，没有一个因痛悔而感动的灵魂所在不是隐秘之处。因为您呼吸着永恒甘甜的芬芳，并如此炽热地爱慕您灵魂的新郎，以至能与天上的新娘一同说：\BibleQuote{愿你拉着我随你跑！我们因你香膏的芬芳而奔跑} \BibleRef{歌 1:3}。但在阁下的书信中，我发现一个缺陷：您不愿告诉我关于您最宁静的主母，她多么勤奋地阅读，或在阅读中如何为痛悔所感动。因为您的临在应当对她大有裨益，使她在不断受苦并被事务的波涛（不论她愿意与否）吸引到外界之时，能内在地被召回对天上家乡的热爱。您也应当常常探究，每当她为灵魂流泪时，她的痛悔究竟是出于恐惧，还是出于爱。\par

因为痛悔有两种，如您所知：一种畏惧永恒的痛苦，另一种渴慕天上的赏报；因为那渴慕天主的灵魂，先被恐惧所感动而痛悔，而后被爱所感动。首先它因回忆自己的恶行而流泪，害怕因此遭受永恒的惩罚。但当恐惧在长期的痛苦焦虑中消失后，某种安全感便从宽恕的意识中产生，心灵被天上喜乐的爱所点燃。先前因害怕惩罚而哭泣的人，此后因被阻隔于天国之外而极其悲痛地哭泣。因为灵魂默观那些天使的唱诗班、蒙福诸圣的团体、以及天主内在荣光的显现，便更为缺乏永恒之善而哀叹，胜过先前害怕永恒之恶时的哭泣。于是，恐惧的痛悔在完善后，便将心灵引向爱的痛悔。这一切在神圣而真实的历史中以预像的方式描述得很好，经上说：\BibleQuote{加肋布的女儿阿革撒坐在驴上，叹息起来。她父亲问她说：你要什么？她回答说：请给我一件祝福；你给了我一处南方的旱地，请再给我一些有水的土地。她父亲就把上泉和下泉赐给了她。} \BibleRef{苏 15:18}。确实，阿革撒坐在驴上，象征灵魂管理肉体不合理的情欲。她叹息着向父亲求一块有水之地，因为泪水的恩宠应当以极大的渴望向我们的造物主祈求。因为有些人已经白白领受了为正义发言、保护受压迫者、将自己的财物施舍给穷人、拥有热忱信德的恩赐，但尚未获得泪水的恩宠。这些人，也就是说，有南方的旱地，但仍需要水泉；因为当他们忙于善工时（这些善工使他们伟大而热忱），仍深切需要（或因害怕惩罚，或因爱慕天国）为他们在世时所不能没有的罪过而哀哭。但是，正如我所说，既然有两种痛悔，她的父亲就赐给她上泉和下泉。灵魂获得上泉，是指她为渴望天国而以泪水自我折磨；获得下泉，是指她因地狱的惩罚而战栗哭泣。确实，下泉先赐予，上泉后赐予。但因为爱的痛悔远高于另一种低劣的痛悔，所以有必要先提上泉，后提下泉。那么，您这位藉全能之主的作为，亲身体验过两种痛悔的人，应当殷切努力，每日发现您以言语对您最宁静的主母有多少裨益。\par

此外，我恳请您特别留意，以良好的品德教导您所抚养的小主人们，并劝诫指派侍奉他们的光荣的太监们，应当向他们说些能激发他们彼此间相互友爱、并对属下保持温和的话语；免得他们现在若彼此心存芥蒂，日后公开爆发。因为抚养孩子的人的言语，若是善的，就是乳汁；若是恶的，就是毒药。因此，让他们现在对孩子们这样说话，好使孩子们日后能显露出他们从喂养者口中吮吸了什么善言。\par

再者，我亲爱的儿子，执事萨比尼亚努斯带来了三十磅黄金，是陛下送来用于赎回俘虏和分施穷人的；对此我欣喜，但为自己战栗，因为我将要在可畏的审判者面前交账，不仅是关于宗徒之长圣伯多禄的财产，也有关于陛下的财物。但愿全能的天主以天上的回报地上的，以永恒的回报暂时的。我现在要告知您，从意大利境内位于亚得里亚海的克罗托纳城——该城去年被伦巴第人攻陷——许多高贵的男女被掳为俘虏，孩子与父母分离，父母与孩子分离，丈夫与妻子分离，妻子与丈夫分离；其中一些人已被赎回。但由于他们被定价甚高，许多人至今仍留在那些最可憎的伦巴第人手中。我随即用您送来的款项的一半去赎回他们。至于另一半，我安排购买床褥，供给那些你们用希腊语称为\OriginalTerm{隐修女}{monastriæ}的天主婢女；因为在这个严冬的酷寒中，她们在床上极度缺乏被褥；此城中这样的人很多。根据官方名单，她们的人数有三千。她们每年确实从圣伯多禄、宗徒之长的产业中获得八十磅。但这对如此众多的群体，尤其是此城中物价昂贵，又算什么呢？此外，她们的生活是这样，在眼泪和斋戒中如此严格，以致我们相信，若不是她们，我们中没有人能在这地方、在伦巴第人的刀剑下生存这么多年。\par

再者，我送给您一把钥匙，作为来自圣伯多禄宗徒的祝福，取自他的至圣遗体；关于这把钥匙，发生了我现将叙述的奇迹。一个伦巴第人，当他进入波河以外某城时发现了这把钥匙，并不重视它是圣伯多禄的钥匙，反而想为自己牟利，因见它是金的，便取出刀来要切割它。但立刻被魔鬼附身，他将原本要切割钥匙的刀刺入自己的喉咙，当场倒地而死。当伦巴第人的王奥塔里特和他的许多随从来到现场，看到自杀者倒毙在一处，而钥匙在地上另一处时，极大的恐惧降临到所有人身上，以致无人敢从地上拾起这把钥匙。于是，一个被称为热心祈祷和施舍的天主教徒伦巴第人，名叫米努尔夫的，被召来，他亲自从地上拾起了它。奥塔里特鉴于这个奇迹，另做了一把金钥匙，连同这把钥匙一起送给我的前任——神圣记忆者，宣告通过它发生了怎样的奇迹。于是我决定将这把钥匙送给陛下，全能的天主曾通过它砍掉了一个傲慢不信的人，好使您——敬畏并爱慕祂的人——能藉此获得现世和永久的福祉。\par

\SourceRecord{Translated by James Barmby.；Nicene and Post-Nicene Fathers, Second Series；Vol. 12.；Edited by Philip Schaff and Henry Wace.；Buffalo, NY: Christian Literature Publishing Co.,；1895.；<http://www.newadvent.org/fathers/360207026.htm>.}

% structure: p0001 source=h1 target=section reason=page-title

\section{第七卷，第二十七封信}

\Emph{致安纳斯塔斯主教。}\par

额我略致安提约基雅主教安纳斯塔斯。\par

我经由我们共同的儿子萨比尼亚努斯执事之手，收到了你至甘饴的圣座所期盼的信函，其中话语并非来自你的口舌，而是来自你的灵魂。一个生活完美的人说话优美，这并不奇怪。而且，由于你藉着在内心学校里教导你的圣神，学到了生活的规诫——轻视一切地上的事物，奔赴天乡——你越在善上进步，就越认为他人为善。但是，当我听到你的真福的信中说了许多赞扬我的话时，我明白你的用意：你是想描述我应当是什么，而不是我是什么。至于你说我应记住自己的生活方式，绝不给那寻求筛净灵魂的恶灵留地步，我确实记起自己一直是生活方式恶劣的人，我急于克服并结束这种生活方式，如果我能做到的话。然而，如果照你所信，我里面有什么好的东西，我信赖全能天主的助佑，我并没有忘记它。但你的圣座，据我看，通过开头甜言蜜语和后面的话语，希望你的信像一只蜜蜂，既带蜂蜜又带刺，用蜂蜜使我满足，用刺刺伤我。但同时我返回默想所罗门的话：\BibleQuote{朋友加的伤痕，胜过仇敌的亲吻}\BibleRef{箴 27:6}。因此，至于你说我们不应无缘无故地给人绊脚石，这正是你的儿子，我们最虔诚的主上（我们应不断为其生命祈祷）已屡次写信指出的；他出于权力所说的话，我知道你是出于爱而说的。我也不奇怪你在信中使用了帝王式的语言，因为爱和权力之间有非常密切的关系。两者都自以为是地行事，两者都经常以权威说话。\par

诚然，在收到我们的兄弟和同僚基里亚库斯主教的教务信函时，我不值得因那个亵圣的头衔而制造难题，冒扰乱圣教会合一的风险；但我仍然留意告诫他关于这个迷信而骄傲的头衔，说除非他纠正那由第一个叛道者发明的上述说法的傲慢，否则他不能与我们保持和平。然而，你不应说这是无关紧要的事，因为如果我们平心静气地忍受它，我们就是在败坏普世教会的信仰；因为你知道有多少人，不仅是异端者，而且异端首领，出自君士坦丁堡教会。且不说对你的尊严的伤害，如果一位主教被称为普世，那么普世教会就会崩溃，如果那位普世者跌倒的话。但这种轻率离我的耳朵极远，极远。然而我信赖全能天主，祂必将早日实现祂的应许：\BibleQuote{凡高举自己的，必被贬抑}\BibleRef{路 14:11}。\par

在诸多事务中，就这么多。我简要回复了你在信中所说的话：因为我不应在此时以文字表达的事，仍铭记在我的脑海中。我恳求你的真福常在圣洁的祈祷中记念我，以使你的转求能救我脱离现世和永世的邪恶。此外，请热心而热切地为最安详的主上皇帝祈祷；因为他的生命对世界非常必要。我不再多说，因为我相信你知道。\par

\SourceRecord{Translated by James Barmby.；Nicene and Post-Nicene Fathers, Second Series；Vol. 12.；Edited by Philip Schaff and Henry Wace.；Buffalo, NY: Christian Literature Publishing Co.,；1895.；<http://www.newadvent.org/fathers/360207027.htm>.}

% structure: p0001 source=h1 target=section reason=page-title

\section{第七卷·第二十八封书信}

\Emph{致医生\Person{德奥多禄}。}\par

\Person{额我略}致\Place{君士坦丁堡}医生\Person{德奥多禄}。\par

我最亲爱的儿子，执事\Person{撒彼尼亚诺}，在返回我这里时，并未带来阁下您的信件；但他却带来了为穷人和俘虏所寄来的东西；由此我明白了原因。那就是您不愿意以书信向人说话，却以善行向全能天主致辞。因为您这同一善行本身有其声音，它呼喊着进入天主的隐秘之耳，正如经上所载：\BibleQuote{将你的施舍藏在穷人的怀里，它必为你祈求} \BibleRef{德 29:15}。而且，我承认，将不属于自己的东西花费掉，并在我所保管的教会财产账目上增加我最亲爱的儿子、大人\Person{德奥多禄}的财物，这令我难过。然而，我与您的仁慈一同喜乐，因您留心注意并遵守真理所说的话：\BibleQuote{你们施舍，看，一切都为你们洁净了} \BibleRef{路 11:41}；以及所记载的：\BibleQuote{水能灭火，同样，施舍能赎罪} \BibleRef{德 3:33}。\Person{保禄}宗徒也说：\BibleQuote{你们的多余，要补足他们的缺乏，使他们的多余，也能补足你们的缺乏} \BibleRef{格后 8:14}。\Person{多俾亚}告诫他的儿子说：\BibleQuote{如果你有许多，就多多施舍；如果你只有少许，也要心甘情愿地拿出那少许来} \BibleRef{多 4:9}。因此，您遵守了这一切诫命；但我们恳求您为我们祈祷，免得我们不谨慎地、不合需要地分发您劳作的果实；免得我们从那使您减少罪过的事物上，反而积聚罪过。现在，愿全能天主保护您，并赐您在世俗朝廷中获得人的恩宠，以便在长寿之后，引领您到达天上朝廷的永恒喜乐。\par

我们寄给您\OriginalTerm{圣伯多禄宗徒之长的祝福}{as the benediction of Saint Peter, Prince of the apostles}，他是您所深爱的，从他那至圣的遗体上取来一把钥匙，内中封存着他锁链的铁片，但愿那曾束缚他颈项以殉道的，也解开您的一切罪过。\par

\SourceRecord{Translated by James Barmby.；Nicene and Post-Nicene Fathers, Second Series；Vol. 12.；Edited by Philip Schaff and Henry Wace.；Buffalo, NY: Christian Literature Publishing Co.,；1895.；<http://www.newadvent.org/fathers/360207028.htm>.}

% structure: p0001 source=h1 target=section reason=page-title

\section{第七卷，第三十封书信}

致隐修士纳尔塞斯（Narsæ Relegioso）\par

额我略致纳尔塞斯等。\par

当我派遣护卫（defensorem）罗马诺前往皇城时，他长期寻找你的信函，却未能找到；但后来在其他许多人的信函中找到了你的信，你的甘饴在其中向我讲述了你的忧苦和心灵的磨难，并告知恶人对你的反对。但我恳求你，在这一切中，要记起我相信你也永远不会忘记的事：\BibleQuote{凡愿意在基督内虔诚生活的，都必遭受迫害。}\BibleRef{弟后 3:12}。关于这一点，我确信地说，如果你遭受的迫害较少，你或许会生活得不够虔诚。因为让我们听听同一外邦人的导师对他的门徒所说的：\BibleQuote{你们自己知道，弟兄们，我们进到你们那里并非徒然；我们事先受苦并受侮辱。}\BibleRef{得前 2:1}。看哪，最甘饴的儿子，圣善的宣讲者宣告，如果他未曾受侮辱，他的进门将是徒然的；而你的爱德希望说良善的事，却拒绝忍受恶事。因此，你必须在逆境中更加紧紧地束腰，好使逆境本身更增加你对天主之爱的渴望和行善的热诚。同样，庄稼的种子因被霜覆盖而结出更丰硕的果实；同样，火焰被风吹压，却能燃得更旺。我确实知道，从众多恶舌的悖逆言语中，你忍受着猛烈的风暴，心里承载着反驳的波涛。但要记住主通过圣咏作者所说的：\BibleQuote{我在风暴的隐秘处听见了你；我在反驳的水边考验了你。}\BibleRef{咏 80:8}。因为，如果你在反驳你的人中间做属于天主的事，那么你就被证明是真正的工人。\par

此外，你最甘饴的爱德写信给我，要我写些劝告的话，给那些通过你的祈祷和影响，由我们的儿子保罗大人所建立的隐修院。但是，如果它们是天主的器皿，我知道它们借着痛悔的恩宠，内在有智慧的泉源，不应接受我干涸的小水滴。再者，你的完美智慧记得，在乐园中没有雨水，只有一道泉源从乐园中间升起，灌溉地面。因此，那些借着痛悔的恩宠内在有泉源的灵魂，不需要他人舌头的雨水。\par

此外，你在信中告知我厄西基雅夫人的去世；我极其欢欣，因为那善的灵魂，在异乡劳苦之后，已幸福地抵达了它的家乡。此外，请代我问候我光荣的女儿们，多米尼加夫人和欧多基亚夫人。但是，因为我听说上述多米尼加夫人被立为修女院长已有一段时间了，愿你的爱德在这件事上照管她：既然她不再被迫在尘世朝廷的劳苦中服务，愿她完全逃离这个世界的一切喧嚣，全心奉献于天主，不将自己的一部分留于自身之外；而且，愿她也尽可能多地聚集灵魂，为她的造物主服务，使她们的心灵通过她的话语领受痛悔的恩宠，并且她本人也因着通过她的生活和言语，其他人的灵魂也得以挣脱罪恶的束缚，而能更快地摆脱她的一切罪恶。此外，既然世上没有人是无罪的（而犯罪岂不就是逃离天主吗？），我坚信我的这位女儿也有一些罪。因此，为使她能完全满足她的主母，即永恒之智慧，愿那独自逃离的人带着许多人回来。因为，在返回时带来收益的人，逃离的罪责将不被追究。\par

此外，我请求你代我问候亚历山大大人和德奥多禄大人。至于你在信中说，我应该写信给我最杰出的女儿古尔迪雅夫人，以及她至圣的女儿德奥克提斯塔夫人，以及她们尊贵的丈夫马里诺大人和基利斯托多禄大人，并给她们一些关于她们灵魂的劝告，你最甘饴的伟人深知，目前在君士坦丁堡城中，没有人能很好地用拉丁文写成的指示翻译成希腊文。因为他们拘泥于字句，却不留意意义，既不能使人理解字句，又破坏了意义。为此，我简略地写信给我上述女儿古尔迪雅夫人，但没有写信给其他人。此外，我送给你两件长衣（camisiæ）和四条手帕（oraria），并请求你谦恭地献上，连同圣伯多禄的祝福，给上述那些人。另外，有一位临终的人遗赠给我一个小男孩；为照顾他的灵魂，我将他送到你的甘饴那里，使他在此世生活时，服侍那一位，并通过他能够获得天堂的自由。此外，我恳求你最甘饴的爱德，常常探望我最亲爱的儿子，执事阿纳多略，我派遣他去代表教会驻在皇城，使他在世俗事务中忍受劳苦之后，能在你那里藉着天主的话得到安息，并擦干他尘世辛劳的汗水，如同用一块白手巾一样。把他推荐给你所认识的所有人，虽然我相信，若他被完全认识，则无需推荐。但请你对他显示出你对圣宗徒伯多禄和我的爱有多少。愿全能的天主守护你的爱德——对我最甘饴的——免受内外敌人的侵害，并在祂乐意的时候，引领你进入天国。\par

\SourceRecord{Translated by James Barmby.；Nicene and Post-Nicene Fathers, Second Series；Vol. 12.；Edited by Philip Schaff and Henry Wace.；Buffalo, NY: Christian Literature Publishing Co.,；1895.；<http://www.newadvent.org/fathers/360207030.htm>.}

% structure: p0001 source=h1 target=section reason=page-title

\section{\WorkTitle{第七卷·书信三十一}}

\Emph{致基里亚库斯主教}\par

格列高利致君士坦丁堡主教基里亚库斯。\par

我们收到了您的真福的信函，这些信函以心灵而非舌尖的话语向我们说话。因为它们向我敞开了您的心扉，这心扉对我原本并非关闭，因为我自身对同一甘甜保有经验。因此我不断感谢全能的天主，既然众德之母的爱德存于您对我们的心中，您将永不丧失善工之枝，因为您紧握着善之根源。您应当首先以这善工向我和所有弟兄展示这爱德之美——即迅速丢弃那骄傲的言辞，这言辞在教会中引起严重冒犯，从而完全实现所记载的：\BibleQuote{竭力保持圣神所赐的合一，以和平相维系}\BibleRef{弗 4:3}；又说：\BibleQuote{不给敌人毁谤的机会}\BibleRef{弟前 5:14}。因为如果我们因骄傲的榜样而彼此分裂，真正的爱德便不会彰显。至于我，我在灵魂中召唤耶稣作证，我不愿给任何人从最高到最低以冒犯的机会。我愿所有人都伟大而尊贵，但他们的尊贵不可减损全能天主的尊荣。因为凡贪求在天主面前被尊荣的，在我看来并不尊贵。但为使您知道我对您的真福怀有何等善意，我已派我的儿子、执事阿纳托利乌斯到我们最虔敬的君主们脚下，为满足他们的虔敬和您的兄弟情谊：我在此事上不愿伤害任何人，而愿保持悦乐天主的谦卑及圣教会的和谐。由于天主的仇敌敌基督临近，我热切希望他找不到任何属于他自己的东西，不仅在司铎的行为中，甚至在称号中。因此，凡以新方式引入的，应如其被带入那样被移除，我们在主内的平安将不受侵犯地存留。因为如果我们用言语取悦自己，而实际遭受刺痛，我们之间有何愉快、有何爱德？愿您的圣德如此行事，使我们能在内心深处感受到您所说的善事，以致司铎们同心合意，当我们为我们最虔敬的君主们祈求生命时，我们更加堪受垂听，因为平安在神眼前照亮了您的祈祷，没有不和谐的污点使之晦暗。\par

\SourceRecord{Translated by James Barmby.；Nicene and Post-Nicene Fathers, Second Series；Vol. 12.；Edited by Philip Schaff and Henry Wace.；Buffalo, NY: Christian Literature Publishing Co.,；1895.；<http://www.newadvent.org/fathers/360207031.htm>.}

% structure: p0001 source=h1 target=section reason=page-title

\section{第七卷·第三十二书}

\Emph{致安纳斯塔修斯司铎}。\par

额我略致安纳斯塔修斯等。\par

善良的人从心中的善库发出善来 \BibleRef{玛 12:35}，你的爱德已经在你的日常生活中，最近也在你的书信中，表明了这一点。在信中，我发现有两个人就德行展开了争论；也就是说，你自己为爱德而争，另一人为恐惧和谦卑而争。虽然我事务繁忙，虽然我不懂希腊文，但我还是作为你们争论的裁判坐了下来。但说实话，依我看，你通过我在争论中向你所提出的那句话，已经击败了你的对手：\Quote{爱德中没有恐惧，但圆满的爱德驱除恐惧，因为恐惧含着痛苦；那恐惧的，在爱德中还没有得以圆满}。因此，我知道你的兄弟之爱在爱德中得以圆满有多少。既然你深爱全能的天主，你就该对近人大有信心。因为使我们与造物主亲近的，不是地方或地位；而是我们的善功使我们与祂结合，或我们的恶行使我们与祂分离。既然我们仍然不确定每个人的内心如何，你又如何因害怕而不敢写信呢？你不知道我们两人中谁更优越吗？确实，我知道你生活良善，但我自己却意识到自己受许多罪的拖累。虽然你自己也是罪人，但你比我好得多，因为你只背负自己的罪，而我还背负那些托付给我的人的罪。因此，在这事上，我看你是崇高的，我看你是伟大的：在高位且在众人眼前尊贵时，你却没有感到自己有任何提升。因为在那里，虽然人们外表上向你致敬，你的心灵却因被分心的忧虑所重压而沉入深渊。但全能的天主却按经上所记载的对待了你：\BibleQuote{祂在心中设定了上升阶梯，在流泪谷中} \BibleRef{咏 83:6}。然而，在我看来，如果你未曾承担那个被称为尼阿斯隐修院的领导职务，你会显得更为崇高、更为超卓，因为在那隐修院中，据我所知，虽然维持着隐修士的外表，但许多世俗事务却在圣洁的外衣下进行。但即使如此，如果在你领导下，那地方令全能天主不悦的事情得到纠正，我仍会认为天上的恩宠带领你到了那里。\par

但是，由于同一隐修院的院长与耶路撒冷教会的牧者之间常有争执，我相信全能的天主意愿你的爱德和我最圣洁的兄弟、同司铎阿摩斯同时留在耶路撒冷，正是为此目的：使我所提到的争执得以止息。那么，现在请展示你从前有多么爱。因为我知道你们两人都是克己的、博学的、谦卑的；因此，按照圣咏的话，应当在鼓瑟和合唱中赞美我们救主的荣耀 \BibleRef{咏 150:4}。因为在鼓瑟中，皮发出的声音是干燥的；但在合唱中，声音是和谐一致的。那么，鼓瑟所象征的若不是克己，合唱所象征的若不是同心合意，又是什么呢？既然你们通过克己在鼓瑟中赞美主，我恳求你们通过同心合意在合唱中赞美祂。真理本人也说：\BibleQuote{你们应当在自己里面有盐，彼此之间保持平安} \BibleRef{谷 9:50}。盐所象征的若不是智慧，又是什么呢？正如保禄所证实的，他说：\BibleQuote{你们的言谈要时常带着恩宠，用盐调和} \BibleRef{哥 4:6}。既然我们知道你们通过天上圣言的教学有了盐，那么，借着爱德的恩宠，你们务要全心在彼此之间保持平安。亲爱的弟兄，我说这一切，是因为我极其爱你们两人，并且非常担心你们之间的任何分歧会玷污你们祈祷的祭献。\par

你送来的祝福，首先由\Person{埃克斯希拉图斯} \Emph{Secundicerius}，随后由执事\Person{萨比尼亚努斯}转交，我以感恩的心领受了，因为从圣地你该当送出圣物，并通过你的礼物显明你一直事奉的是谁。愿全能的天主用祂的右手保护你，使你免受一切邪恶的伤害。\par

\SourceRecord{Translated by James Barmby.；Nicene and Post-Nicene Fathers, Second Series；Vol. 12.；Edited by Philip Schaff and Henry Wace.；Buffalo, NY: Christian Literature Publishing Co.,；1895.；<http://www.newadvent.org/fathers/360207032.htm>.}

% structure: p0001 source=h1 target=section reason=page-title

\section{第七卷，第三十三封信}

\Emph{致摩里士奥古斯督。}\par

额我略致摩里士奥古斯督。\par

我主上们的眷顾上智，惟恐因司铎们的争执而在圣教会内生出什么绊脚石，已多次屈尊叮嘱我要仁慈地接待我的兄弟和同僚司铎齐里亚古的代表，并准其早日返回。敬爱的陛下，虽然您的一切训令都是合宜且有预见的，但我发觉这样的劝诫使我感到在您眼中被视为不智。然而，即使我的心灵因那骄傲亵圣的称号而受到不小的创伤，我岂能犯下如此大的不智，以致不知道我对信仰的统一和教会的和谐所应尽的义务，并拒绝因任何苦因而接待我兄弟的代表和他所送来的教务会议信函？这绝不可能。这种智慧反倒成了愚拙。因为为维护信仰统一所当行之事是一回事，为抑制高傲所当行之事是另一回事。因此必须区分时宜，以免我前述兄弟的新近之举在任何方面受到扰乱。因此，我满怀深情地接待了他的代表。我向他们显示了所应尽的爱德，并以超出古代惯例的方式尊敬了他们，并让他们与我一同举行弥撒的神圣典礼；因为，正如我的执事不应在展示神圣奥迹时服侍那位犯有高傲之罪或他人犯此罪而自己不加改正的人，同样，在举行弥撒时，他的辅礼者应当服侍我——我，在天主的护佑下，并未陷入骄傲的迷误。\par

不过，我已留意恳切劝诫我前述的兄弟和同僚主教，若他愿与众人保持和平与和谐，就必须戒除那愚妄名号的称谓。关于此事，我主上们的仁爱已在谕令中责成我，说不可因一虚浮的名号在我们中间生出嫌隙。然而，我恳请陛下圣鉴，有些事情虽然虚浮却无大害，另一些则极其有害。当假基督来临并自称天主时，岂非非常虚浮，却极其有害？若就言语的数量而言，不过寥寥几音节；若就错误的严重性而言，则是普世性的灾祸。现在我自信地说：凡自称为普世司铎或渴望被如此称呼的人，其傲慢就是假基督的先驱，因为他骄傲地把自己置于众人之上。而且，他所陷入的迷误并非源于不同的骄傲；因为正如那悖逆者愿意显得高于众人，同样，这位贪求被称为独一司铎的人，也把自己高举于所有其他司铎之上。但由于真理说：\BibleQuote{\Emph{凡高举自己的，必被贬抑}}（\BibleRef{路 14:11}），我知道每一种傲慢，愈是膨胀，就愈快破灭。请陛下责成那些陷入傲慢榜样的，不要因一个虚浮名号的称谓而制造任何嫌隙。因为，我这罪人，在天主助佑下持守谦卑，无需受劝诫要谦卑。愿全能天主长久护佑我们最安宁的主上的生命，以维护圣教会的和平及罗马共和国的利益。因为我们确信，倘若您——敬畏天上之主的人——尚在，您必不容许任何骄傲的行径胜于真理。\par

\SourceRecord{Translated by James Barmby.；Nicene and Post-Nicene Fathers, Second Series；Vol. 12.；Edited by Philip Schaff and Henry Wace.；Buffalo, NY: Christian Literature Publishing Co.,；1895.；<http://www.newadvent.org/fathers/360207033.htm>.}

% structure: p0001 source=h1 target=section reason=page-title

\section{\WorkTitle{第七卷，第三十四封书信}}

\Emph{致主教\Person{厄洛吉}。}\par

\Person{格列高利}致\Place{亚历山大}主教\Person{厄洛吉}和\Place{安提约基}主教\Person{阿纳斯塔修斯}。\par

深深系于你们的爱德使我绝不能缄默，以便圣座知晓我等中间所发生的一切，不受任何虚假传闻的欺骗，并更完美地保持您已完美开始的正义与正直之道。现在，我们的弟兄及同僚主教\Person{基里亚科}的\OriginalTerm{代表}{responsales}来到我这里，带来了他的教务会议公函。诚然，如圣座所知，因一个亵渎名称的称谓，我们与他之间存在着严重的分歧；但我认为，为信仰事宜派来的代表应当受到接待，以免在君士坦丁堡教会中几乎是针对所有司铎而起的傲慢之罪引发信仰动摇和教会合一破裂。我也让这些代表，因为他们非常谦卑地请求，与我一同举行了弥撒圣祭，因为正如我已经向最安详的皇帝陛下所表明的那样，我的弟兄和同僚司铎\Person{基里亚科}的代表应当与我共融，因为靠着天主的帮助，我没有陷入傲慢的错误。但我的执事不应与我们的上述弟兄\Person{基里亚科}一同举行弥撒圣祭，因为他通过一个亵渎的称号犯下了或参与了傲慢之罪；否则，如果他（我的执事）与一个处于如此傲慢地位的人同行，我们可能会（但愿不会）确认那个愚妄名称的虚妄。但我已经告诫我们的这位弟兄改正这种傲慢，因为如果他不改正，他绝不会与我们平安相处。\par

此外，我们的这位弟兄在他的教务会议公函中，靠天主的恩宠，在所有方面都表达得像一位公教徒。但他谴责了一位名叫\Person{欧多克修}的人，我们既没有在教务会议上看到他受到谴责，也没有在他的前任们的教务会议公函中看到他受到弃绝。诚然，君士坦丁堡公会议的法令谴责了\Person{欧多克修}派；但关于他们的创始人\Person{欧多克修}是谁，却只字未提。但罗马教会至今尚未拥有这些法令或那次公会议的法令集，也没有接受它们，尽管它接受了同一公会议针对\Person{马其顿尼}所定义的内容。它确实弃绝了其中提到的其他异端，这些异端早已被其他教父所谴责；但至今对\Person{欧多克修}派一无所知。\Person{索佐门}的\WorkTitle{史书}中确实讲述了一些关于某个\Person{欧多克修}的事，据说他篡夺了君士坦丁堡教会的监督职位。但这部史书本身，宗座拒绝接受，因为它包含许多虚假陈述，过度赞扬\Person{摩普绥提亚的狄奥多勒}，并说他直到去世之日都是教会的一位伟大圣师。因此，如果任何人接受这部史书，他就与在可敬纪念的\Person{优士提尼安}皇帝时代举行的关于\WorkTitle{三章}的公会议相矛盾。但一个不能反驳这次公会议的人必须拒绝这部史书。此外，我们在拉丁语中至今没有找到任何关于这位\Person{欧多克修}的内容，无论是在\Person{斐拉斯特}的著作中，还是在圣\Person{奥思定}关于异端的著作中。所以，请圣座在回信中告知我，在希腊人中是否有哪位公认的教父提到过他。\par

此外，三年前，关于\Place{依索利亚}隐修士的案件——他们被指控为异端——我的弟兄和同僚主教、\Person{若望}大人曾给我送来一封信函以作澄清，他在信中试图证明他们违反了\Place{厄弗所}公会议的定义；并向我转寄了某些章节，据称是同一公会议的定义章节，而他们说这些隐修士反对它们。在这些章节中，除其他事项外，主张\Quote{亚当的灵魂没有因罪而死，因为魔鬼不会进入人的心；凡如此说者当受\OriginalTerm{绝罚}{anathema}。}当我读到此时，深感悲痛。因为如果第一个犯罪的亚当的灵魂没有因罪而死，那关于禁果所说的话又当如何理解：\BibleQuote{你吃的那天必定要死}（\BibleRef{创 2:17}）？诚然，亚当和厄娃吃了禁果，但在肉身中，他们后来又活了九百多年。由此可见，他在肉身中并未死亡。那么，如果他的灵魂没有死，就得出一个亵渎的结论：天主对他所说的话是虚假的，即他吃的当天必定要死。但愿这个错误远离，愿它远离真信仰。因为我们所说的是：第一个人在犯罪当天灵魂死了，并且通过他，全人类都被定罪于这死亡和腐朽的惩罚。但通过第二个人，我们相信可以获得自由：现在从灵魂的死亡中，将来在永恒复活中从肉身的全部腐朽中——正如我们曾对前述代表所说：\Quote{我们说亚当的灵魂因罪而死亡，不是就生活的本质而言，而是就生活的品质而言。因为本质是一回事，品质是另一回事，他的灵魂并没有死到不存在的地步，而是死到不再有福的地步。但这位亚当后来通过忏悔又回到了生命。}\par

然而，若相信福音，就不能否认魔鬼进入人心。因为经上记载：\BibleQuote{耶稣递了饼以后，撒殚便进入他的心。}\BibleRef{若 13:27}。又经上亦说：\BibleQuote{魔鬼已进入犹达斯的心，使犹达斯出卖祂。}（\BibleRef{若 13:2}）。否认这一点的人便陷入\OriginalTerm{白拉奇异端}{Pelagian heresy}。鉴于我们审查了厄弗所大公会议，发现其中并无此类内容，我们又从拉温纳教会取得同一会议的极古\OriginalTerm{抄本}{codex}，发现它与我们所拥有的会议报告完全一致，毫无差异，且在其绝罚与弃绝的法令中，除了弃绝可敬的济利禄的十二章之外，别无他物。但我们曾向他的代表全面且详尽地阐述了这一论点，并使他们完全满意。因此，为避免这些或类似的事情潜入彼处，冒犯圣教会，我们有必要将这些事指示给您的圣座。而且，虽然我们知道我们的弟兄和同为主教的齐亚古为正统，但为了别人，我们应当谨慎，好使错误的种子尚未萌发公开就被践踏。\par

在您的圣座的信件由我们共同的儿子、执事\Person{萨比尼亚诺}送达此地时，我已收到。但由于信使已准备离去，不能耽搁，待我的\OriginalTerm{特使}{responsalis}执事到来时，我将回复。\par

\SourceRecord{Translated by James Barmby.；Nicene and Post-Nicene Fathers, Second Series；Vol. 12.；Edited by Philip Schaff and Henry Wace.；Buffalo, NY: Christian Literature Publishing Co.,；1895.；<http://www.newadvent.org/fathers/360207034.htm>.}

% structure: p0001 source=h1 target=section reason=page-title

\section{第七卷，第三十五封书信}

\Emph{致\Person{多米尼各}主教}\par

\Person{额我略}致\Place{迦太基}主教\Person{多米尼各}。\par

虽然我们相信阁下怀有牧灵之警惕，关注隐修院之事，但仍认为有必要将我们所知的非洲省份某隐修院的情况告知阁下。现有一持信人，即院长\OriginalTerm{昆魁德斯}{Cumquodeus}抱怨说，每当他欲按会规纪律约束所辖隐修士时，他们便立即离开隐修院，且被允许随意游荡。那么，鉴于此事既于彼等自身极为有害，亦为他人树立败亡之榜样，我们恳请阁下，若情况属实，当对彼等施以教会谴责，并以适当惩罚制止如此确实的狂妄；且应以有益的安排引导彼等顺从，使其高傲之心屈服于纪律之轭，以使惩戒能召回那些因彼等榜样而可能犯下类似过犯的人，并教导他们应当服从长上。然而，因他告知我们有些主教庇护迷途的隐修士，请阁下谨慎处理此事，并以警告全面制止此类庇护。七月，\OriginalTerm{小纪}{Indiction}第十五。\par

\SourceRecord{Translated by James Barmby.；Nicene and Post-Nicene Fathers, Second Series；Vol. 12.；Edited by Philip Schaff and Henry Wace.；Buffalo, NY: Christian Literature Publishing Co.,；1895.；<http://www.newadvent.org/fathers/360207035.htm>.}

% structure: p0001 source=h1 target=section reason=page-title

\section{第七卷，第三十八封书信}

\Emph{致多努斯主教。}\par

\Person{额我略}致\Place{墨西拿}主教\Person{多努斯}（\OriginalTerm{墨西拿}{Messene}）。\par

神圣教规与法律的指令均允许出售教会器皿以赎回俘虏。因此，鉴于持信人\Person{福斯蒂努斯}被证实为赎回其女儿们脱离被俘的枷锁而负担了三百三十\OriginalTerm{索利都斯}{solidi}的债务，且其中已偿还三十，显然他没有足够的资金偿还余款，我们通过此信劝告你的弟兄：你务必从你手中属于\Place{梅里恩西安教会}的银子中给他十五磅，并收取其收据；他已知属于该教会。如此，将银子变卖，还清债务，他便可免除其债务的束缚。但你的弟兄亦应注意：倘若上述教会有足够的流通货币，他应从其中收取上述数额；否则，为所述目的，你必须从祝圣过的器皿中向他提供我们所陈述的金额。因为，正如无事出售教会器皿是非常严重的事，另一方面，在这种紧迫的必要下，对于一个极其困苦的教会来说，宁可要其财产而不要其俘虏，或在赎回他们时怠慢，也是错误的。\par

\SourceRecord{Translated by James Barmby.；Nicene and Post-Nicene Fathers, Second Series；Vol. 12.；Edited by Philip Schaff and Henry Wace.；Buffalo, NY: Christian Literature Publishing Co.,；1895.；<http://www.newadvent.org/fathers/360207038.htm>.}

% structure: p0001 source=h1 target=section reason=page-title

\section{第七卷 第三十九封信}

\Emph{致\Person{若望}主教。}\par

\Person{额我略}致\Place{锡拉库萨}主教\Person{若望}。\par

为避免世俗事务的关注使虔敬者的心（但愿天主不许）远离相互的爱德，应尽力使任何已发生争议的事项得以最轻易地了结。鉴于我们从位于\Place{巴亚}地方的\Person{圣伯多禄}隐修院院长\Person{才奇拉留}的信息获悉，在他与设于\Place{锡拉库萨}城的\Person{圣路济亚}隐修院院长\Person{若望}之间，就某些边界问题发生了严重争议，我们，为了避免他们之间争执延长，已考虑由一名土地测量员的裁定来解决他们的纠纷。因此，我们致函\OriginalTerm{护民官}{defensor}\Person{方提努斯}，命他指示那位从\Place{罗马}前往\Place{巴勒莫}的土地测量员\Person{若望}去投奔你的阁下。\par

因此，我们劝勉你与他一同前往有争议的地方，并在双方被召集到一起后，在你的见证下，使争议之地划定边界，但仍保留双方任何一方提出四十年时效的主张。然而，无论决定如何，务请你的阁下尽心留意，务使从此不再引起纷争，也不再有任何投诉到达我们这里。\par

我们相信，你的阁下并非不知道，可敬的院长\Person{才奇拉留}曾是我们旧友；因此，在不失公平的前提下，我们将他在各方面托付给你。而且，鉴于他对于世俗诉讼完全缺乏经验，他需要你的关怀扶助；但是，在此事和一切事上，你都应如理地秉持理性与正义。\par

\SourceRecord{Translated by James Barmby.；Nicene and Post-Nicene Fathers, Second Series；Vol. 12.；Edited by Philip Schaff and Henry Wace.；Buffalo, NY: Christian Literature Publishing Co.,；1895.；<http://www.newadvent.org/fathers/360207039.htm>.}

% structure: p0001 source=h1 target=section reason=page-title

\section{第七卷·第四十封}

\Emph{致尤洛吉主教}\par

格列高利致亚历山大里亚主教尤洛吉。\par

至甘饴的圣座在来信中多次论及宗徒之长圣伯多禄的座位，说现今在继承者身上，伯多禄本人亲自坐在这座位上。我固然自认不配，非但不能与如此主理者相比，甚至不堪立于侍立者之列。然而，既然这位占据伯多禄之座者对我谈论伯多禄之座，我便欣然接受所言。纵使专属于我的尊荣丝毫不令我欢喜，我却因你们至圣者将加于我的恩惠同样加于自身而大为欣慰。因为谁能不知圣教会建立于宗徒之长磐石般的稳固之上？他因心志的坚定而得名，从\OriginalTerm{磐石}{petra}而得称为伯多禄。真理之声曾对他说：\BibleQuote{我要将天国的钥匙交给你}\BibleRef{玛 16:19}。又对他说：\BibleQuote{你回头以后，要坚固你的弟兄}\BibleRef{路 22:32}。同样：\BibleQuote{西满，若望的儿子，你爱我吗？你牧放我的羊}\BibleRef{若 21:17}。因此，虽然宗徒众多，但就首席权而言，唯独宗徒之长之座在权威上坚固壮大，此座虽在三处，却同属一位。因为他亲自尊崇了那蒙他屈尊安息并结束现世生命的座位；他亲自装饰了那派遣其门徒作为福音作者的座位；他亲自奠定了那本要离去却仍坐七年之久的座位。既然同属一席，同为一座，如今三位主教靠天主权威治理，那么凡从你处所闻的善事，我便归于自己。你若以为我有什么善行，也请归于你的功德，因为我们在那一位内原为一体，他曾说：\BibleQuote{愿众人合而为一，父啊，如同你在我内，我在你内，使他们也在我们内合而为一}\BibleRef{若 17:21}。此外，在偿还应向你致候之情时，我向你表明：得知你殷勤劳作，对抗异端者的狂吠，我满心欢悦；我恳求全能的天主以他的护佑扶助你的真福，好借你的口舌从圣教会的怀中拔除一切苦根，免得它再滋生，阻碍多人，且使多人受玷污。因为你既领受了自己的塔冷通，便记起那诫命：\BibleQuote{你们去做生意，直到我回来}\BibleRef{路 19:13}。我虽全无经商之力，却为你的经营所得与你同乐，因我知道：若劳作不使我参与，爱德却使我参与你的辛劳。因为我认为，一个闲立者若能因他人的成就而共同喜乐，邻人的善功便是共有的。\par

再者，我本想送些木材给你，但你的真福并未表明是否需要；我们可以送些更大的，但现今没有船只前来装载，若送小的，我又觉惭愧。请你的真福来信告知我应如何办理。\par

然而，作为爱你者之圣伯多禄教会的小小祝福，我送了你六件阿奎丹式小斗篷（\OriginalTerm{斗篷}{pallia}）和两条手巾（\OriginalTerm{巾帕}{oraria}）。因我情谊深厚，即便微小之物也望蒙悦纳。因为情意本身自有其价值；且无疑，凡因爱而贸然行事者，必无过犯。\par

再者，按照你信末所附清单，我收到了圣史马尔谷的祝福。但既然我不喜饮用\OriginalTerm{滤酒}{colatum}和\OriginalTerm{药酒}{viritheum}，便冒昧索取\OriginalTerm{香料酒}{cognidium}——去年隔了许久之后，你的圣座曾使此酒在本城为人所知。因我们在此地从商贩口中得知其名，却未得实物。现恳请你的圣座以祈祷扶持我，对抗此生所遭受的一切苦楚，并以代祷求全能的天主护佑我。\par

\SourceRecord{Translated by James Barmby.；Nicene and Post-Nicene Fathers, Second Series；Vol. 12.；Edited by Philip Schaff and Henry Wace.；Buffalo, NY: Christian Literature Publishing Co.,；1895.；<http://www.newadvent.org/fathers/360207040.htm>.}

% structure: p0001 source=h1 target=section reason=page-title

\section{第七卷，第42封信}

\Emph{致马里尼亚诺主教。}\par

\Person{额我略}致\Person{马里尼亚诺}，\Place{拉文纳}主教。\par

从贵弟兄来信中获悉，\Place{科尔内留姆}教会的子女不断恳求您为他们祝圣一位主教，以接替他们那位已经失足的前任主教，而您对此事应如何处理感到犹豫，并等待我们明确的指示。既然任何理由都不允许一个犯有罪行的人被召回他已失足的职任，且神圣教规的法令不允许一个教会超过三个月没有主教，以免（但愿天主阻止）古老的仇敌伺机撕裂主的羊群，贵弟兄应当应允他们的请求，在失足者的位置上祝圣一位主教。因为，既然您本应在他们请求之前就以劝勉提醒他们此事，那么当他们要求您时，您便没有理由拒绝，因为天主的教会不应长期没有自己的主教。\par

\SourceRecord{Translated by James Barmby.；Nicene and Post-Nicene Fathers, Second Series；Vol. 12.；Edited by Philip Schaff and Henry Wace.；Buffalo, NY: Christian Literature Publishing Co.,；1895.；<http://www.newadvent.org/fathers/360207042.htm>.}

% structure: p0001 source=h1 target=section reason=page-title

\section{第七卷，第四十三封信}

\Emph{致马里尼亚诺斯主教。}\par

额我略致拉文纳主教马里尼亚诺斯。\par

我们从许多人的报告中得知已有一段时间，在拉文纳地区设立的隐修院普遍受到你方圣职人员的管辖欺压；以致——说来令人痛心——他们以管理的借口将这些隐修院据为己有，如同自己的产业。我们对这些隐修院深感同情，曾致信你的前任，命他纠正这一恶习。然而，眼见不久后他被生命的终结所笼罩，我们记得也曾同样致函贤弟，以免隐修院的负担持续下去。而且，由于我们发现此事至今拖延未改，我们认为宜借此信再次致函于你。因此，我们劝勉你，抛开一切迟延和借口，努力解除这些隐修院所受的这种委屈，使圣职人员或领受圣秩者今后除了为祈祷之目的，或偶而被邀请为庄严举行弥撒圣祭之外，不得有任何进入隐修院的许可。但是，惟恐因某位隐修士或院长的晋升而使隐修院承受负担，你须留意：若某隐修院的院长或隐修士获得任何圣职职务或圣秩，他将如我们所说，不再拥有那里的权力，以免借此机会隐修院被迫承担我们所禁止的负担。那么，在第二次告诫之后，愿你的圣德不要迟延，以警醒的关怀纠正这一切，以免今后我们发现你疏忽懈怠（我们相信不会如此），而不得不另行设法确保隐修院的安宁。须知，我们绝不再容忍天主仆人的团体受此类要求之辖制。然而，为免就隐修士之事提出任何借口，请贤弟毫不迟延地派遣你认为合适的人来，我们将派遣隐修士同他一起到你处；为了供养他们，你必须将他们安置在隐修院内，如果你们那里确实有可以养活他们的地方的话。\par

\SourceRecord{Translated by James Barmby.；Nicene and Post-Nicene Fathers, Second Series；Vol. 12.；Edited by Philip Schaff and Henry Wace.；Buffalo, NY: Christian Literature Publishing Co.,；1895.；<http://www.newadvent.org/fathers/360207043.htm>.}

% structure: p0001 source=h1 target=section reason=page-title

\section{第八卷，第一封信}

\Emph{致彼得主教。}\par

额我略致科西嘉主教彼得。\par

收到你的兄弟的信函后，我们向全能天主献上极大的感谢，因为你以许多灵魂被聚集的消息使我们在安慰中振作。因此，你的兄弟应努力焦虑，藉主的助佑，将你开始的工作达致完美。至于那些曾经信主，但因疏忽或被迫而重投偶像崇拜的人，你要速速将他们带回信德，对他们施加数日的补赎，使他们哀哭自己的罪过，并更坚定地持守他们所回归的，因天主帮助他们，正如他们彻底痛悔所离弃的；至于那些尚未受洗的人，你的兄弟应速速藉着劝诫、恳求、以即将来临的审判警吓他们，同时也给出理由为何不应崇拜木石，将他们聚集到全能天主前；如此，在祂来临时，当严厉的审判之日到来，你的圣德可在圣人行列中被找到。因为还有什么工作比这更有利或更高尚：那就是为灵魂的苏醒与聚集而操心，并将不朽的收益带给你的主，祂已将宣讲的职位赐予你？\par

此外，我们给你的兄弟送去五十\Emph{索利多}，用于购置为受洗者准备的礼服；我们也已将你的兄弟所请求的、位于内格乌格努斯山的教会所属的地产赐给了该教会的司铎，以便其价值从他一贯领取的钱中扣除。\par

此外，你的兄弟请求获准在离同一座山不远的教会内为自己建造一座主教住所；我十分乐意同意这一提议，因为你越接近，就越能为那里的灵魂行善。\par

鉴于你的圣德为他的代祷，我们已将此信的持有者擢升为辅祭，并派他回去服侍你，为使他若在拯救灵魂方面仍有更大用处，则可得到进一步的提升。\par

\SourceRecord{Translated by James Barmby.；Nicene and Post-Nicene Fathers, Second Series；Vol. 12.；Edited by Philip Schaff and Henry Wace.；Buffalo, NY: Christian Literature Publishing Co.,；1895.；<http://www.newadvent.org/fathers/360208001.htm>.}

% structure: p0001 source=h1 target=section reason=page-title

\section{第八卷·书信二}

\Emph{致安提约基雅主教阿纳斯塔西乌斯}\par

额我略致安提约基雅宗主教阿纳斯塔西乌斯\par

我收到了您最甘饴的真福的来信，那信字字含泪。因为我从中看见一片云飞向高空，如同云常做的事；但它带着忧苦的黑暗而来，我难以在开头就察觉它从何而来、往何处去，因为由于我所说的黑暗，我没有完全明白它的起源。然而，至圣者，你应当常记着外邦人的宣讲者所说的话：\BibleQuote{在末期，危险的时代必将来临，那时人要自私、贪财、傲慢}\BibleRef{弟前 4:1}；以及随后的话，我不便详述，您也不需听。看，在您神圣的晚年，您的真福承受许多磨难；但请想想您坐在谁的座位上。岂不就是那位，真理之声曾对他说：\BibleQuote{你年老时，别人要给你束上腰，带你到你不愿意去的地方}\BibleRef{若 21:18}？但我说这话时，想起您的圣座从青年时就曾在许多逆境中劳苦。那么，请与善王一同说：\BibleQuote{在我灵魂的苦楚中，我要追忆我所有的年岁}\BibleRef{依 38:15}。因为有许多人，如您在信中所说，以我们的创伤为快乐；但我们知道谁曾说过：\BibleQuote{你们将要哀哭哭泣，世人却要欢乐；你们将忧愁}\BibleRef{若 16:20}；他又立即加上：\BibleQuote{但你们的忧愁将转变为喜乐}。但既然我们已经承受了所预言的事，我们就该盼望所应许的事。至于您所说的那些人，他们自己把应当减轻的重担加在别人身上，我知道他们是披着羊皮来到的，内心却是贪婪的狼\EditorialNote{玛 7}。但他们更应被容忍，因为他们迫害我们不仅心怀恶意，而且披着宗教的外衣。而且，他们想要为自己拥有超过他人的东西，即使与弟兄共享也不合适，我们对此毫不困扰，因为我们信赖全能的天主：那些贪图别人之物的人，将更早被剥夺自己的所有。因为我们知道谁说过：\BibleQuote{凡自高的，必被降卑}\BibleRef{路 14:11}。经上又记着：\BibleQuote{心骄者必先跌倒}\BibleRef{箴 16:18}。\par

但在这些日子里，据我发现，异端的新战争正在兴起，关于他们我先前已写信给您的真福。他们试图否定先知、福音书和所有教父的言论。但是，只要您的圣座活着，我们信赖我们护佑者的恩宠：他们反对真理之稳固的口将更快被阻止，因为无论所用的刀剑多么锋利，当它们击打岩石时，就会折断而回。此外，由于全能天主的极大恩宠，那些与神圣教会教义分离的人之间没有合一，因为凡分裂的王国必不得站立\EditorialNote{路 11}。而神圣教会总是因异端的质问而更充分地装备自己的教导，以致应验了圣咏作者论天主反对异端所说的话：\BibleQuote{他们因他面容的愤怒而分散，他的心中接近了我们}\BibleRef{咏 54:22}。因为当他们分散在邪恶的错误中时，天主将他的心接近我们，因为通过矛盾，我们更彻底地学会认识他。\par

此外，我们从蛮族的刀剑和法官的邪恶中遭受的祸患，我不愿向您的真福述说，免得增加您的哀叹，而本应通过安慰减轻。但在这一切中，我导师的诫命安慰了我，他说：\BibleQuote{我给你们讲了这些事，为使你们在我内有平安。在世界上你们要受苦难}\BibleRef{若 16:33}。因为我想到他曾对谁说过：\BibleQuote{这是你们的时候和黑暗的权势}\BibleRef{路 22:53}。那么，如果光明之时刻随后到来，因为曾对蒙选者说：\BibleQuote{你们是世上的光}\BibleRef{玛 5:14}，又正如经上所记：\BibleQuote{清晨，义人要管辖他们}\BibleRef{咏 48:15}，那么无论我们在黑暗权势的时刻遭受什么，都不应悲伤。\par

此外，您最甘饴的圣座告诉我，您曾经愿意，如果可能的话，不用纸张和笔墨与我交谈，并因我们之间相隔几乎如东从西的距离而悲伤。但我所感受到的，我坦诚地宣认，这是真实的：在纸上，您的灵魂对我说话，没有纸张，因为在您圣座的话语中，只有爱德在回响，而由于全能天主的恩赐，我们在爱德的联系中结合，并不被地方的距离分开。那么，您为何寻求得到那银羽的鸽子的翅膀，而您早已拥有它们？因为这些翅膀正是对天主和对邻人的爱德。藉着这些，神圣教会高飞，藉着这些，她超越一切尘世的事物；如果您的圣座没有这些，您就不会以如此大的爱德通过书信来到我这里。\par

此外，我恳求您为我的心弱而恳切祈祷，以使全能的天主藉着您的转求，保护我的灵魂脱离一切邪恶，并且尽快将我从此时代如此众多的风暴中夺出，带到永恒安息的海岸。\par

我已收到了所有赠给我的非常丰富的福祉，您作为神贫的天主之人，将它们寄给了我，并论及它们说：一个穷人除了贫穷之物，还能给什么呢？但如果您不是藉着谦逊之神而贫穷，您的福祉就不会丰富。愿全能的天主以祂的保护护卫您远离一切邪恶；并因您的生命对所有善人非常必要，在多年之后，带您抵达天乡的喜乐。\par

\SourceRecord{Translated by James Barmby.；Nicene and Post-Nicene Fathers, Second Series；Vol. 12.；Edited by Philip Schaff and Henry Wace.；Buffalo, NY: Christian Literature Publishing Co.,；1895.；<http://www.newadvent.org/fathers/360208002.htm>.}

% structure: p0001 source=h1 target=section reason=page-title

\section{第八卷 第三封信}

\Emph{致\Place{梅西纳}（\Place{西西里}）主教\Person{多努斯}。}\par

\Person{格列高利}致\Person{多努斯}等。\par

最有口才的人，我们的儿子\Person{福斯蒂诺}来到我们这里，抱怨说他的已故父亲\Person{佩尔特拉修斯}将一些不属于他自己的东西留给您的教会作为他埋葬的代价。他本人知道，我们也听说了，世俗法律在此类案件中的规定；即，如果父亲遗赠了不属于他自己的东西，继承人有义务支付。但是，由于我们知道您的弟兄之爱是依照天主的法律而非世界的法律生活，在我看来，非常不公的是，一个琥珀杯和一个据说属于他产业上某座位于\Place{科森扎}教区内的教堂的男童，竟被您的弟兄之爱扣留。因为，当现任主教、当时是总执事的最可敬的\Person{帕伦布斯}证实事情正如我所说时，您本应当相信他的话，并归还不属于您的东西。此外，在我看来，您本应考虑那枚金胸针——若是有任何东西可以供养他身后遗下的人，这几乎是他的全部财产——并当时就接受它作为他的埋葬费用。尽管如此，您知道我们的训令，即我们已经完全禁止了我们教会中的旧习，也不赞同任何人被允许为获得埋葬人身体的墓地而支付代价。因为，如果舍根人——我们假定他们是外邦人——曾无偿向\Person{亚巴郎}提供埋葬死去的\Person{撒辣}的墓地，使她能葬在她自己的地方，并且\Person{亚巴郎}再三恳求，他们才勉强接受了她墓地的价格；那么我们这些被称为主教的人，难道应该为埋葬信友的遗体而收费吗？因此，我们将此事交付您的弟兄之爱裁决。\par

上述最有口才的人也抱怨此事：您的教会的监护人（\OriginalTerm{监护人}{defensor}）\Person{西西尼乌斯}无理地扣押属于他所有的奴隶；关于这些奴隶，他还断言，依可敬的\Person{马克西米安}主教的判决，扣押者应予交出，但他至今故意拖延交还。因此，我们劝勉您的弟兄之爱，若案件已明确判决，则应按所定执行。否则，应指派某人处理此案，让他到我们的兄弟和同僚主教\Person{塞昆迪努斯}的辖区寻求裁决，以便当他的判决宣告这些奴隶属于谁时，一方面不会使一方似乎受损，另一方面也不会使另一方心怀怨恨。\par

\SourceRecord{Translated by James Barmby.；Nicene and Post-Nicene Fathers, Second Series；Vol. 12.；Edited by Philip Schaff and Henry Wace.；Buffalo, NY: Christian Literature Publishing Co.,；1895.；<http://www.newadvent.org/fathers/360208003.htm>.}

% structure: p0001 source=h1 target=section reason=page-title

\section{第八卷，第五封书信}

\Emph{给各位都主教和主教}。\par

格列高利致得撒洛尼的\Person{欧瑟比}、都拉基翁的\Person{乌尔比提}、\Place{米兰}的\Person{君士坦提}、\Place{尼科波利}的\Person{安德肋}、\Place{格林多}的\Person{若望}、\Place{第一尤斯蒂尼安娜}的\Person{若望}、\Person{若望}（\Place{克利坦西·斯科里塔诺}）、\Place{拉里萨}的\Person{若望}、\Place{拉文纳}的\Person{马里尼亚诺}、\Place{撒丁岛}的\Place{卡利亚里}的\Person{雅奴雅里}，以及\Place{西西里}的所有主教。\par

我曾小心地将最虔诚的皇帝所颁布的法律传达给你们的弟兄，内容是关于那些受军役或公共职责约束之人，不得为了逃避被追究的风险，擅自取得圣职身份或成为隐修士。这一点我特别叮嘱你们：不可草率接纳涉足世俗职分者进入教会的圣职班，因为他们在圣职生活中行为与先前无异，根本不是在试图逃避世俗事务，而是在变换它们。但是，如果有人甚至寻求隐修院，除非他们先被免除公共职责，否则绝不得被接纳。另外，若来自军旅的人急于成为隐修士，不得草率接纳，除非他们的人生已得到充分查考。按照常规，他们应当经受三年的考验，然后，若天主恩准，再穿隐修会服。如果他们这样被证实并接纳，并为了他们灵魂的益处而渴望为他们所犯的罪做补赎，那么，为着他们天上的生命和利益，不应拒绝他们发隐修誓愿。关于此事，请相信我，最宁静且最虔诚的皇帝已经完全安抚，并乐意准许那些他知道未涉入公共职责的人发隐修愿。十二月，第一\OriginalTerm{小纪}{Indiction}。\par

\SourceRecord{Translated by James Barmby.；Nicene and Post-Nicene Fathers, Second Series；Vol. 12.；Edited by Philip Schaff and Henry Wace.；Buffalo, NY: Christian Literature Publishing Co.,；1895.；<http://www.newadvent.org/fathers/360208005.htm>.}

% structure: p0001 source=h1 target=section reason=page-title

\section{第八卷·第六封信}

\Emph{致\Place{耶路撒冷}宗主教\Person{阿摩斯}。}\par

\Person{额我略}致\Place{耶路撒冷}主教\Person{阿摩斯}。\par

确信你的兄弟之爱重视教规的规定与纪律的权威，以免你的一位圣职人员的欺瞒得逞，欺骗了你，从而逃过教会秩序的严格管教；我们以为理应告知你他的过失，好藉你的关怀，使他接受他已逃避的纪律。我们得知，有一位名叫\Person{伯多禄}的辅祭，我们曾命他在我们的儿子、执事\Person{撒比尼雅努斯}（即我们在皇城的教会代表）手下服务，现已逃走，并投奔到你的教会。若此事属实，愿你的兄弟之爱尽力扣留他，并一有机会便将他送回这里。但倘若他因惧怕此事而离开你的教会，潜伏各处以避侦查，则命人细心搜查你所有的堂区，一旦找到，便如前所述送他回来。我们也愿通过你告知他已被剥夺共融：他不得胆敢领受主的体与血的奥迹，直到他返回我们这里，除非他偶然面临死亡的迫近危险。\par

\SourceRecord{Translated by James Barmby.；Nicene and Post-Nicene Fathers, Second Series；Vol. 12.；Edited by Philip Schaff and Henry Wace.；Buffalo, NY: Christian Literature Publishing Co.,；1895.；<http://www.newadvent.org/fathers/360208006.htm>.}

% structure: p0001 source=h1 target=section reason=page-title

\section{第八卷，第十封信}

\Emph{致雅德拉主教撒彼尼亚努斯}\par

额我略致撒彼尼亚努斯等。\par

对于坚持错误的人，理应施以惩罚；但对于回心转意的人，也应予以宽恕。因为，正如在前一种情况下，对犯过者的义怒是理所当然的，所以在后一种情况下，所显示的善意常常能促进和好。因此，既然您的兄弟之谊现在已因忆及司铎职位的崇高而远离了与马克西穆斯的交往和共融（此前轻率曾使您陷入其中）；而且，您如此深切地感到，仅仅与他的分离并不能使您满足，还必须退居隐修院，通过隐修生活来哀悼您过去的过犯，因此，毋庸置疑，您再次被接纳进入我们的恩宠和共融：因为，正如您的过犯先前曾激怒我们一样，您的忏悔也已平息了我们。所以，最亲爱的弟兄，我们劝勉您，要勤于以牧灵的热忱照料主的羊群，并要时刻警醒，使托付给您的羊群获得益处；这样，在永恒的审判者来临时，您所献上的劳苦果实越是丰硕，您所得的丰厚赏报也越加丰盛。因此，努力去拯救那些陷于罪恶中的人吧；努力为迷途者指明回头的道路吧；努力将那些被剥夺了共融的人健康地召回恩宠的共融中吧。您的回归应使您肩负起拯救他人的责任，并成为得救的榜样；如此，当您的殷切关怀引导迷途的羊群回到大牧者的羊栈时，他们自己就不致暴露在狼的利齿之下，而（这是最值得渴望的）相称的赏报也将为您存留在永生中。\par

关于您写信给我们所提的那件事，即请求我们防止有人在皇城暗中对您不利，请不要为这件事烦心。因为我们已经尽一切可能嘱咐我们的\OriginalTerm{代理司铎}{responsalis}，要他保持警觉并加以防范。而且，我们信赖我们天主的德能，事情正在如此进行，以致没有任何人的反对能对抗真理，从而在任何方面困扰您或给您带来重压。\par

此外，厄庇道鲁斯城的居民极其迫切地请求我们将佛洛伦提乌斯归还给他们，他们声称此人是他们的主教，并断言他只是由于纳塔利斯主教的个人意愿而被非法流放。因此，如果您的兄弟之谊对此案有所了解，请写信详细告知我们。但如果您至今尚不了解，请进行调查，然后向我们报告，以便我们在主的助佑下，能全面了解情况后审慎决定如何处理此事。二月，第一小纪。\par

\SourceRecord{Translated by James Barmby.；Nicene and Post-Nicene Fathers, Second Series；Vol. 12.；Edited by Philip Schaff and Henry Wace.；Buffalo, NY: Christian Literature Publishing Co.,；1895.；<http://www.newadvent.org/fathers/360208010.htm>.}

% structure: p0001 source=h1 target=section reason=page-title

\section{第八卷，第十三封信}

致\Person{哥伦布斯}。\par

\Person{额我略}致\Place{努米底亚}主教\Person{哥伦布斯}。\par

我们如何可以推测您的爱德，是从我们自身对您的情感中得出的。我们也不认为您爱宗座会与宗座爱您有所不同。因此，我们必须特别推荐那些我们知道在宗徒之长真福伯多禄的教会中应当献身于您的人，而您的生活既受到司铎职位的尊严的装饰，我们又从过去的经验中已经确知您的真诚。\par

至于我们的兄弟及同为主教的\Person{保禄}，这位持信者，他告诉我们他在您那边遭受的风浪与逆境，对您的圣德而言并非不知。由于他声称，您告诉我们传到您耳中的对他的指控并非真实，而是其对手煽动起来的，并信靠主的帮助能够克服一切，有真理支持他，有您审查，我们恳请您，最亲爱的兄弟，在显然正义在他一边的任何事情上，您都合宜地伸出援手，并以司铎的同情援助他。因此，不要让任何情况或任何人的影响使您偏离对公平的认真关注。要依靠主的诫命，藐视任何违背正义的事物。在捍卫任何一方时，始终坚持正义。不要害怕为真理招致敌意，若有的话；好使您在\OriginalTerm{救赎主}{Redeemer}降临时，获得更大的赏报，因为您没有忽略祂的命令，而是献身于维护和捍卫正义。三月，第一禀政年。\par

\SourceRecord{Translated by James Barmby.；Nicene and Post-Nicene Fathers, Second Series；Vol. 12.；Edited by Philip Schaff and Henry Wace.；Buffalo, NY: Christian Literature Publishing Co.,；1895.；<http://www.newadvent.org/fathers/360208013.htm>.}

% structure: p0001 source=h1 target=section reason=page-title

\section{第八卷，第十四封信}

\Emph{\Person{博尼法斯}，\OriginalTerm{首位辩护人}{Defensorem}}。\par

\Person{额我略}致\Person{博尼法斯}，论\OriginalTerm{辩护人}{Guardians}的特权。\par

那些为教会利益忠心劳苦的人，应该得到适当报酬的恩惠，如此既可显示我们对其服务作出了相称的回报，而他们因蒙受所赐的慰藉之恩，也能显得更有用处。既然如此，那些担任辩护人职务的人，众所周知是勤于教会事业并服务于教宗们，我们认为理应让他们享有以下赏予他们的特权作为酬报——规定：正如在公证人和副助祭的\OriginalTerm{团体}{schola}中，因昔日教宗们的宽厚而设立了\OriginalTerm{教区公证人}{regionarii}，同样，在辩护人中，凡经证明有功绩的七位，应获得教区公证人的尊荣。我们规定，在教宗缺席时，他们有权在任意神职人员的集会中坐在任何位置，并在各方面享有其尊严的特权。此外，若有人获得此优先地位，偶然为自身利益居住在其他省份，他仍须在各方面占据其优先席位，以便成为所有辩护人之首，因为他甚至在获得优先地位之前，就已通过勤勉的个人关注，不断献身于教会利益与教宗之服务。因此，我们制定这些为辩护人的特权和组织而颁布的法令，指定它们必须永久有效且不可违反——无论是我们以书面方式所规定的，还是那些被认为是在我们面前所颁布的；我们也规定，它们不得被任何教宗在任何情况下全部或部分推翻或更改。因为这是一种极其严苛的行为，尤其与司铎的良好品行相悖：任何人以任何借口企图废除已经妥善规定的事，并且以自身榜样教导他人在自己时代之后解散自己的法令。四月，第一\OriginalTerm{小纪}{Indiction}。\par

\SourceRecord{Translated by James Barmby.；Nicene and Post-Nicene Fathers, Second Series；Vol. 12.；Edited by Philip Schaff and Henry Wace.；Buffalo, NY: Christian Literature Publishing Co.,；1895.；<http://www.newadvent.org/fathers/360208014.htm>.}

% structure: p0001 source=h1 target=section reason=page-title

\section{第八卷，第十五封信}

\Emph{致\Person{马里尼亚努斯}，\Place{拉文纳}主教。}\par

额我略致\Person{马里尼亚努斯}等。\par

维护隐修院的安宁并采取措施确保其永久安全的必要性，你阁下因先前担任隐修院管理职务而深为了解。因此，鉴于我们得知，位于\Place{克拉西斯}城的\Person{圣若望}及\Person{圣斯德望}隐修院（由我们共同的儿子、院长\Person{克劳狄}主持）曾因你的前任们而遭受许多损害和不满，理应由你阁下为其未来安宁作出有益安排，为的是修士们在那里事奉天主，并蒙祂的恩宠助佑，得以心神自由地坚持。但为避免因那更应纠正的习俗，任何人任何时候胆敢在那里造成任何烦扰，有必要使你阁下细心注意，确保我们下面列举的几点得到遵守，以免将来可能找到使他们不安的借口。今后，任何人不得再以任何形式的调查，从上述隐修院或任何与其相关的地方的收入或契据中减少任何东西，或尝试任何形式的强占或诡计。但如果\Place{拉文纳}教会与上述隐修院之间发生任何争端，且不能友好解决，则应在由双方选出的敬畏天主的人面前，在至圣福音书上宣誓，无故意拖延地了结。此外，院长去世时，不得任命外人，而应由会众从同一会众中自由选择一人，且该人当选时无任何欺诈或贿赂。但如果他们在自己中间找不到合适的人，他们也应同样明智地从其他隐修院中选择一人来任命。并且，当一位院长到来时，在任何情况下，不得在其隐修院中设立任何人高于他，除非（愿天主禁止）出现圣教法典规定应受惩罚的罪行。同样必须严格遵守这条规定：未经该隐修院院长的同意，不得将修士调离，用于充实其他隐修院、授予圣职或任何神职职务。但如果修士充足，足以颂扬天主并满足隐修院的需要，院长应虔诚地从那些可省下的人中，提供他在天主面前认为合适的人。但如果院长在有足够数量的情况下拒绝提供，那么\Place{拉文纳}主教可从那些可省下的人中取用，以充实其他隐修院。然而，除非该地院长在接到通知后自愿提供，否则不得从该隐修院中挑选任何人担任神职职务。任何从上述隐修院获得任何圣职的人，此后在该处既无任何权力，亦不得居住在那里。\par

还应注意到，如有需要，不得由神职人员制作该隐修院财产和契据的清单；而应由该地院长与其他院长一起制作财产清册。\par

此外，凡院长可能希望为了其隐修院的利益前往或派遣使者到罗马教宗处，他应有完全的自由这样做。\par

此外，虽然隐修院应渴望主教的探访，但鉴于我们听说上述隐修院在你的前任时期曾因接待而负担沉重，所以你阁下理应以适当方式加以规定，使得该城主教可以随时为探访和劝勉之故进入隐修院。但主教应在那里履行爱德职责，以免隐修院遭受任何负担。现在，上述院长不仅不害怕你阁下频繁进入隐修院，反而渴望如此，因为他知道隐修院的财产绝不可能因你而受负担。于四月，第一小纪颁发。\par

\SourceRecord{Translated by James Barmby.；Nicene and Post-Nicene Fathers, Second Series；Vol. 12.；Edited by Philip Schaff and Henry Wace.；Buffalo, NY: Christian Literature Publishing Co.,；1895.；<http://www.newadvent.org/fathers/360208015.htm>.}

% structure: p0001 source=h1 target=section reason=page-title

\section{第八卷，第十七封信}

\Emph{致\Person{马伦提乌斯}}\par

\Person{格列高利}致军团长\Person{马伦提乌斯}。\par

我最亲爱的儿子，执事\Person{西彼廉}，他的归回本令我甚为欢喜，倘若他全人已回到我这里。但现在你仍留在西西里，我确切知道，他身体确实回来了，心神却留在西西里。然而，说此话时，我因你的宁静而与你一同喜乐，如同我为自己忙碌的职务而叹息。为此我诚恳劝诫你：若内心甘饴的美好滋味已触动你心口的味觉，愿你的心神如此沉浸于自身，以至于一切外在的声音、一切外在的愉悦都令你感到乏味。此外，我称赞你回避人群的聚集，因为那渴望借着痛悔的恩宠在天主内更新的心神，常常因邪恶的交谈和言语而重坠旧态。我曾寻找一些人加入你的读经团体，却未寻得一人，我极其痛惜善事之稀缺。虽然我这罪人非常忙碌，但若你愿意来到蒙福宗徒\Person{伯多禄}的门槛，你将能得我为研读圣经的亲密同伴。愿全能的天主将你置于祂的护佑之下，并赐你得以躲避古敌的陷阱。\par

\SourceRecord{Translated by James Barmby.；Nicene and Post-Nicene Fathers, Second Series；Vol. 12.；Edited by Philip Schaff and Henry Wace.；Buffalo, NY: Christian Literature Publishing Co.,；1895.；<http://www.newadvent.org/fathers/360208017.htm>.}

% structure: p0001 source=h1 target=section reason=page-title

\section{第八卷 第十八封信}

\Emph{致泰拉奇纳主教阿格内卢斯。}\par

额我略致阿格内卢斯等。\par

我们听闻——说来令人震惊——有人在你所属的地区敬拜树木，并做出许多其他违背基督信仰的违法之事。我们感到奇怪，为何您的兄弟迟迟没有以严厉的惩罚来纠正此事。因此，我们以这封信劝勉您，要您令人仔细查问，找出这些人，并予以惩戒，使天主得以息怒，而他们的处罚也能成为责备他人的榜样。\par

我们也已致信莫鲁斯子爵，要他在这件事上向您的兄弟提供帮助，这样您就再也找不到借口不抓捕他们。此外，我们发现许多人以守城为借口推脱责任，因此请您的兄弟注意，不得以我们或你们教会的名义，或任何其他理由，允许任何人免除守城之责；而应强制所有人一律参与，以便在全体守城时，赖主的助佑，城市的守卫能得到更好的保障。\par

\SourceRecord{Translated by James Barmby.；Nicene and Post-Nicene Fathers, Second Series；Vol. 12.；Edited by Philip Schaff and Henry Wace.；Buffalo, NY: Christian Literature Publishing Co.,；1895.；<http://www.newadvent.org/fathers/360208018.htm>.}

% structure: p0001 source=h1 target=section reason=page-title

\section{第八卷，第二十封信}

\Emph{致\Place{拉文纳}主教\Person{马里尼亚诺}}\par

\Person{额我略}致\Person{马里尼亚诺}等\par

\Person{约翰}，持此信者，申诉其妻为逃避\Person{乔治}之骚扰，久居圣所之内，迄今未获援助。因其声称身份存有争议，并请求将此事托付于阁下，我等兹劝请阁下庇护此妇，不得任人无理欺凌。若其身份问题仍悬而未决，请留意使案件得以依法公正审理，免受任何压迫；如此，待查明真相、依律裁断后，双方皆不得诉称遭受不公。五月，小纪第一年。\par

\SourceRecord{Translated by James Barmby.；Nicene and Post-Nicene Fathers, Second Series；Vol. 12.；Edited by Philip Schaff and Henry Wace.；Buffalo, NY: Christian Literature Publishing Co.,；1895.；<http://www.newadvent.org/fathers/360208020.htm>.}

% structure: p0001 source=h1 target=section reason=page-title

\section{第八卷，第二十一书}

\Emph{致\Person{若翰}，\Place{西拉库萨}主教}\par

\Person{额我略}致\Person{若翰}等。\par

\Person{斐理克斯}，持此函者，向我们诉苦说，他生于基督徒父母，却被某位基督徒给予了一个\OriginalTerm{撒玛黎雅人}{Samaræan}为奴（即\Emph{为奴}）。这是何等凶残之事，言之可愕。虽然法律秩序和宗教虔敬均不容许此类迷信之人以任何方式拥有基督徒奴隶，但他声称在那人手下服役达十八年之久。然而他说，当你的前任，圣善记忆的\Person{马克西米安}得知此事后，便以司铎应有的热忱将他从该撒玛黎雅人的奴役中释放了。但是，因为该撒玛黎雅人的儿子据说在五年后成了基督徒，而某些人正试图据他的说法将上述\Person{斐理克斯}召回服役，请尊座仔细查问我们被告知的事实。若查实为真，则尽力保护他，绝不允许任何人以任何借口使他受委屈。因为，既然法律明文禁止，在信德上领先于其主人的那种迷信教派的奴隶被召回服役，那么更何况这个生于基督徒父母、自幼即为基督徒的人，岂能在此争议中受到任何牵累？尤其是，他既不能成为那另一人之父的奴隶，而那人显然因其邪恶的妄为而应受法律惩罚。因此，如我们所说，愿尊座的保护如此合理地庇护他，使任何人都不得以任何借口，在任何程度上随意加害于他。\par

\SourceRecord{Translated by James Barmby.；Nicene and Post-Nicene Fathers, Second Series；Vol. 12.；Edited by Philip Schaff and Henry Wace.；Buffalo, NY: Christian Literature Publishing Co.,；1895.；<http://www.newadvent.org/fathers/360208021.htm>.}

% structure: p0001 source=h1 target=section reason=page-title

\section{第八卷·第二十二封信}

\Emph{致贵族夫人\Person{Rusticiana}}.\par

\Person{Gregory}致\Person{Rusticiana}，等。\par

我记得以前曾写信给阁下，并反复敦促您不可迟延地重新访问真福\Person{伯多禄}，宗徒之长的门槛。而我不知您何以如此喜爱\Place{君士坦丁堡}城，却遗忘了\Place{罗马}城。至今我尚未蒙恩从您那里得到任何关于此事的消息。因为这对您的灵魂在获得永生赏报方面有多少益处，以及在各方面如何适合您光荣的女儿贵妇\Person{Eusebia}，我们已充分关注，您也同样可以充分考虑。但是，如果您询问我的儿子\Person{伯多禄}，您的仆人——我发现他年纪虽小却智慧过人，且正在努力学习以求成熟——您将会发现所有住在这里的人对阁下怀有多么大的爱，以及他们多么渴望有朝一日能再次见到您。而且，如果主教导我们，\WorkTitle{圣经}告诫我们应爱仇敌，那么连爱我们的人都不表示爱，那是多么错误，我们应当思考。但是，如果我们偶然被说成是被爱的，我们最确凿地知道，没有人能对他不想见的人有感情。然而，如果您害怕\Place{意大利}的刀剑和战争，您应当留心观察真福\Person{伯多禄}，宗徒之长，在这城市中是何等大的保护；在这城里，没有大量军队，没有军事援助，我们却在刀剑之下蒙天主保全了这么多年。我们这么说，是因为我们爱你们。但愿全能天主赐予凡祂看出永远有益于您灵魂、且现时有益于您家族声望的一切。\par

阁下为赎取俘虏而送来的十磅黄金，我已由我前述之子\Person{伯多禄}手中收到。但我祈求，那曾赐予您为您的灵魂的赏报而施舍它们的天上恩宠，也赐予我以毫无罪染的方式分施它们；免得我们被那借以洗除罪恶的事所玷污。愿全能天主，眷顾您身体的软弱和您的旅途，常以祂的恩宠安慰您，并藉着我最甜蜜的儿子\Person{Strategius}阁下的生命与健康安慰您；如此，祂可以养育他，既为您多年，也为祂自己直到永远，并以现世的福祉充满您和您全家，赐予您从上而来的恩宠。我们亦恳请，代我们问候光荣的阁下\Person{Eudoxius}。\par

\SourceRecord{Translated by James Barmby.；Nicene and Post-Nicene Fathers, Second Series；Vol. 12.；Edited by Philip Schaff and Henry Wace.；Buffalo, NY: Christian Literature Publishing Co.,；1895.；<http://www.newadvent.org/fathers/360208022.htm>.}

% structure: p0001 source=h1 target=section reason=page-title

\section{第八卷，第二十三封书信}

\Emph{致守护者\Person{凡蒂诺}}\par

\Person{额我略}致\Person{凡蒂诺}等。\par

从\WorkTitle{圣斯德望隐修院}（位于\Place{阿格里真托}地区）的女院长处获悉，许多犹太人因天主的恩宠感动，愿意皈依基督徒信仰；但需要有人奉我们的命令前往那里。因此，我们授权并命令你，放下一切借口，速赴上述地点，在天主的眷顾下，以你的劝勉帮助他们实现愿望。然而，若他们认为等待逾越节庆典过于漫长和乏味，而你发现他们此刻急切渴望领受圣洗圣事，那么，为避免拖延可能改变他们的心意（愿天主阻止此事），请与那地方我们的主教弟兄商议，在为他们规定四十天的补赎和斋戒之后，由他在主日或任何可能适逢的重大庆节中，在全能天主的仁慈保护下为他们施洗；因为当前时局因即将来临的灾难也促使我们不应以任何拖延推迟满足他们的愿望。此外，凡你查明其中贫穷且无力为自己购买洗礼服饰的人，我们希望你为他们提供洗礼服饰；并且要知道，你为此所付的款额应记入你的账目。但是，如果他们选择等待神圣的复活节期，请再与主教商议，使他们暂时成为望教者，主教应经常探访他们，悉心照料，并以劝勉的训导激发他们的心灵，使他们越是期待遥远的庆节，就越发预备自己，以炽热的渴望等待它。\par

此外，你当以全部的热忱和勤勉查明上述隐修院（由前述女士管理）是否拥有足够的资源，或是否有所匮乏。无论你确实查知什么，以及对于那些渴望领受圣洗圣事的人所行之事，都请迅速详细告知我们。六月，第一小纪。\par

\SourceRecord{Translated by James Barmby.；Nicene and Post-Nicene Fathers, Second Series；Vol. 12.；Edited by Philip Schaff and Henry Wace.；Buffalo, NY: Christian Literature Publishing Co.,；1895.；<http://www.newadvent.org/fathers/360208023.htm>.}

% structure: p0001 source=h1 target=section reason=page-title

\section{第八卷，第二十四封书信}

\Emph{致\Person{萨比尼亚努斯}，\Place{雅德拉}主教}。\par

格列高利致\Person{萨比尼亚努斯}等。\par

至爱的弟兄，我在你的真诚中深感欣喜，知道你的真诚以谨慎判断的鉴别力，既在应当服从之处服从，也在应当抵抗之处以司铎的热忱抵抗。因为你以何等欣然的虔诚服从了我们因你过去过犯的过错所命令的事，已由你经送信者寄给我们的书信内容向我们披露。事实上，我所爱的弟兄不可能以别的方式接受一位爱他者所命令的事。因此，我信赖全能天主的慈悲，祂的恩宠如此保护你，使你既已如此从其他罪过中获得赦免，得以因健康地服从而喜乐。至于你的爱德所提及的，因被绝罚的背信者\Person{马克西穆斯}的嫉妒而烦恼一事，你不应受困扰；反而你应当以忍耐承受，抵挡那徒然稍许汹涌的波涛，并以恒忍的德能制服浪花的泡沫。因为忍耐知道如何使粗糙变得平滑，恒常知道如何克服凶猛。因此，不要让逆境打击你的精神，反要激励它。让司铎的刚毅使你在一切事上更显勇敢。因为这是真理的真实证明：在艰难环境中表现得更敏捷，在逆境中更勇敢。因此，为使任何打击都不能动摇你正直的坚定从良善的决心，像你已经开始做的那样，将你灵魂的脚步扎根于那磐石的坚固之上——你知道我们的救赎者在那磐石上建立了遍及普世的教会——这样，一颗真诚之心的正确步履就不会在弯曲的道路上绊倒。\par

至于你所写的事，或送信者在我们的面前所解释的事，不要以为我们忽略了它们：我们正在非常仔细地考虑。\par

此外，我们此前和现在都已将一切详情准确地告知我们至爱的儿子、执事\Person{阿纳托利乌斯}，劝他藉我们造物主的助佑，毫不迟延地严格而热诚地行动，处理一切关乎你的爱德及你的儿子们之利益与安宁的事。因此，愿悲伤不使你的弟兄之情忧虑，也不让任何人的敌意使你痛苦。因为，藉神圣恩宠的助佑，我们相信不久之后，那被绝罚的背信者的狂妄将受到更严厉的压制，而你的安宁将如你所愿到来。我们也丝毫没有忽略就他的乖戾写信给我们最卓越的儿子、\Person{总督}，他急于将他推荐给我们。\par

至于那位司铎，你的弟兄之情通过送信者的陈述向我们咨询了关于他的事，要知道，他在失足之后绝不能留在或恢复至其神圣品秩。尽管如此，仍应稍加温和地对待他，因为他据称已欣然承认了自己的过错。此外，这位送信者同时谈到了某些由我们前任授予你的教会的特权。\par

关于你爱德所提及的这些文件，我们希望得到更准确的信息。或者，如果其中一些文件存放在你教会的档案中，有必要将副本发送至此，以便我们能乐意更新一切关乎你尊严之尊重或前述教会之特质的事。\par

如果我们的共同儿子、荣耀的大人\Person{玛尔塞鲁斯}愿意前来此处，请极力劝说他这样做；因为从各方面来说，我都渴望见到他。但如果他选择留在原地，那么你们要以合宜的爱德向他展现自己，以便能恰当地回应他对你们的感情。愿全能天主以祂恩宠的恩赐保护你，并点燃你的心去做祂喜悦的事。\par

\SourceRecord{Translated by James Barmby.；Nicene and Post-Nicene Fathers, Second Series；Vol. 12.；Edited by Philip Schaff and Henry Wace.；Buffalo, NY: Christian Literature Publishing Co.,；1895.；<http://www.newadvent.org/fathers/360208024.htm>.}

% structure: p0001 source=h1 target=section reason=page-title

\section{第八卷，第二十九封信}

\Emph{致亚历山大主教厄娄其乌斯。}\par

\Person{额我略}致\Person{厄娄其乌斯}等。\par

一位博学之人的讲话总是有益的，因为听者要么学到自己已知有所不知的事，要么，更有甚者，开始认识到自己不知自己有所不知。如今后一种听者正是我，您至圣的真福有意写信给我，要求我将虔诚纪念的君士坦丁时代由凯撒勒雅的欧瑟伯收集的全部殉道者行传寄给您。但在收到您的真福的信之前，我并不知道这些行传是否已被收集。因此我感谢您，藉着您至圣的教导，我开始认识我所不知道的事。因为除了同一欧瑟伯的书中关于圣殉道者行传的内容外，我并不知道我们这教会的档案或罗马城的图书馆中有任何收集，除非是少数编成一卷的东西。我们确实有一卷收集了几乎所有殉道者的名字，以及他们分配于特定日子的苦难；我们在那些日子举行弥撒圣祭以纪念他们。但在这卷书中并未指明每人是谁、如何受苦；只记载了名字、地点和苦难的日子。因此可知许多来自不同地区与省份的人戴着殉道的荣冠，正如我所说，藉着他们各自的日子。但我们相信这些您已有。然而，您所希望寄给您的那些，我们已寻找，未曾找到；虽然未曾找到，我们仍将继续搜寻，若能找到，必当寄送。\par

关于您所写的木材长度不足，原因在于运输船只的类型；因为若有更大的船来，我们可以运送更大的木材。至于您说若我们寄送更大块，您将付款，我们确实感谢您的慷慨，但我们不能接受付款，因为福音禁止这件事。因为我们不购买我们运送的木材；我们怎能接受付款，正如经上记载：\BibleQuote{你们白白得来的，也要白白分施}（\BibleRef{玛 10:8}）？因此，我们现在通过船长寄送了短木材，以适应船的大小，兹附上清单。然而，明年若全能的天主愿意，我们将准备更大的木材。\par

我们已怀着它所伴随的善意，收到了圣史玛尔谷的祝福，不，更确切地说，是圣宗徒伯多禄的祝福；并向您问候，恳请您的真福垂念为我们祈祷，好使我们堪当早日摆脱当前的邪恶，并得以不被排除于未来的喜乐之外。\par

\SourceRecord{Translated by James Barmby.；Nicene and Post-Nicene Fathers, Second Series；Vol. 12.；Edited by Philip Schaff and Henry Wace.；Buffalo, NY: Christian Literature Publishing Co.,；1895.；<http://www.newadvent.org/fathers/360208029.htm>.}

% structure: p0001 source=h1 target=section reason=page-title

\section{\WorkTitle{第八卷·第三十书}}

\Emph{致亚历山大城主教厄瓦格里乌斯}\par

\Person{额我略}致\Person{厄瓦格里乌斯}等。\par

我们共同的儿子，持此信者，当他带来你的圣座的书信时，发现我患病，离开时我仍患病；由此导致，我简短书信的涓涓细流几乎无法渗入你的圣座的广阔泉源。但在身体痛苦之际，我收到你的圣座的信函，得以因你训导亚历山大城异端者以及促进信友和谐而喜乐，这实是天上的恩赐；如此，我心灵中的喜乐缓解了病痛的严重性。我们确实以新的欢跃为听到你的善行而喜乐，尽管我们绝不认为你如此完美行事是新事。因为圣教会信徒增加、为天上粮仓的属灵庄稼倍增，我们从未怀疑这是全能天主的恩宠，这恩宠丰沛地流向你——最蒙福者。因此，我们向全能天主谢恩，因我们在你身上看到经上所载的：\BibleQuote{出产丰饶之处，可见牛力之强}\BibleRef{箴 14:4}。因为，若没有一头强壮的牛将舌头的犁拉过听众心田的土壤，如此多的信友增长绝不会出现。\par

但是，既然在你的善行中我知道你也与别人同乐，我回报你的恩惠，并宣布与你所做的并非不同的事；因为身处世界一隅的盎格鲁民族，至今仍不信奉木石崇拜，我决定藉着你为我祈祷的帮助，若天主允准，从我隐修院派遣一名隐修士去宣讲福音。他获得我的许可，由日耳曼地区的主教们祝圣为主教，在他们的协助下前往世界尽头——上述民族处；我们已经收到信件，告知我们他的平安和工作；据说他及与他同去的人在所述民族中闪耀如此伟大的奇迹，以致他们似乎因所显示的标记效法宗徒的能力。此外，在本小纪第一年发生的主诞生庆节中，据报超过一万名盎格鲁人由我们同一的这位兄弟和同为主教者受洗。我将此事告知你，为叫你知道你在亚历山大城人民中藉言语所成就的，以及在世界尽头藉祈祷所成就的。因为你的祈祷在你不在之处，而你的圣善事工显示在你所在之处。\par

其次，关于异端者尤多克修斯其人，我对于他的错误在拉丁语中未曾发现任何信息，我欣喜地看到我已被你的圣座极其充分地满足了。因为你引用了强有力的人物——巴西略、额我略和厄丕法尼的见证；我们承认他显然已被杀死，因我们的英雄们向他掷出了如此多的镖枪。至于那些被证明已在君士坦丁堡教会中出现的这些错误，你在各个要点上作出了极其博学的答复，并且恰如其分地表达了如此伟大宗座的判断。因此，我们感谢全能天主，盟约的法版仍在天主的约柜里。因为司铎之心若非盟约之柜又是什么？既然属灵的教义在其中保持活力，无疑法律的法版安放在其中。\par

你的圣座也小心地宣布，在致某些人的信函中，你现在不使用骄傲的称谓——这些称谓源于虚荣之根，并且你致函我时说：\Quote{如同你所命令的}。我恳求你从此不再让我听见\Emph{命令}这个词，因为我知道我是谁，你是谁。就职位而言，你是我的兄弟；就品德而言，你是我的父亲。我所以并非命令，而是愿意指出似乎有益的事。然而，我发现你的圣座似乎不愿完全记得我提醒你的这一点。因为我曾说，无论对我还是对任何人，你都不应写这类东西；且看，在你致我——那位禁止此举的人——的书信序言中，你竟认为可以使用一个骄傲的称谓，称我为\OriginalTerm{普世教宗}{Universal Pope}。但我恳求你那最为甘饴的圣座不要再这样做，因为给予他人超出理性所要求的，就是从你自己身上减损。至于我，我不寻求藉言语昌盛，而是藉我的行为。我也不把那视为荣誉，因我知道那会使我的兄弟失去荣誉。因为我的荣誉是普世教会的荣誉；我的荣誉是我兄弟们的坚固活力。只有当给予每个人应得的荣誉不被拒绝时，我才真正受到尊敬。因为如果你的圣座称我为普世教宗，你就否认你本人就是你所普遍称呼我的那个。但这事绝不可发生在你我之间。摒弃那些膨胀虚荣、伤害爱德的言语。\par

确实，你的圣座知道，在加采东大公会议中，以及后来由继任的教父，这曾献给我的前任们。然而，他们没有一人曾使用这个称号，为的是在这世上尊重所有司铎的荣誉，同时在全能天主面前保持自己的荣誉。最后，在向你致以应有的问候的同时，我恳求你惠然在神圣的祈祷中纪念我，为使主因你的代祷赦免我罪恶的束缚，因我自身的功德无济于事。\par

\SourceRecord{Translated by James Barmby.；Nicene and Post-Nicene Fathers, Second Series；Vol. 12.；Edited by Philip Schaff and Henry Wace.；Buffalo, NY: Christian Literature Publishing Co.,；1895.；<http://www.newadvent.org/fathers/360208030.htm>.}

% structure: p0001 source=h1 target=section reason=page-title

\section{第八卷，第三十三封信}

\Emph{致\Person{多米尼库斯}}\par

\Person{额我略}致\Place{迦太基}主教\Person{多米尼库斯}。\par

我们经由持此信者手中收到的贵座阁下的信函，如此表达了司铎的谦逊，以至于仿佛以其作者的肉身临在安抚了我们。的确，沟通的稀少并未造成任何损害，只要爱的情意在心内不受阻隔。此外，爱德的力量何其伟大，亲爱的兄弟，它以真诚的链环将彼此的心在互相的情感中联结，不让它们从恩宠的凝聚中松开；恩宠结合离散的，保持联合的，并使素未谋面的人在爱中相识。因此，无论何人将心固定于爱德的枢纽上，任何逆境都无法将他从天国乡居的住所中撕裂，因为无论他转向何方，他都不离开诫命的门槛。因此，卓越的宣讲者也在赞扬同一种爱德时如此说：\BibleQuote{这就是全德的联系}\BibleRef{哥 3:14}。我们因而看出，那不仅能在灵魂中生出成全，更将其束缚的爱德，应受何等赞美。\par

因此，既然你的信函显示你被这德行的火焰所燃烧，我以极大的喜乐在主内欢欣，并希望它在你内越来越多地闪耀，因为牧人的火焰乃是羊群的光。因为主的司铎理应在品行和生活上发光，好使托付给他的子民，仿佛在他生活的镜子里，既能选择效法之事，又能看见应改之处。\par

此外，深知司铎祝圣在非洲地区起始于何处，你以明智的回忆，在对宗座的爱中，回归到你职务的起源，并以堪得称赞的恒心继续持守你对它的情感。因为确实，你以司铎的智慧对它表显的任何尊敬和奉献，你都增添到了自己的荣耀上；因为你借此邀请它以相应的爱来与你束缚。\par

最亲爱的兄弟，我们尚须以不断的祈祷恳求全能的天主，指引我们心灵的行径进入祂真理的道路，并带领我们到达天国，藉祂保护的恩宠，赐予我们能在行动中表显我们所名义上承担的职分。八月，第一小纪。\par

\SourceRecord{Translated by James Barmby.；Nicene and Post-Nicene Fathers, Second Series；Vol. 12.；Edited by Philip Schaff and Henry Wace.；Buffalo, NY: Christian Literature Publishing Co.,；1895.；<http://www.newadvent.org/fathers/360208033.htm>.}

% structure: p0001 source=h1 target=section reason=page-title

\section{第八卷，书信三十四}

\Emph{致斯基拉西乌姆主教若望}。\par

额我略致若望等。\par

显然，这是一件非常严重的事，与司铎所应追求的目标相悖，即希望扰乱先前赐予任何隐修院的特权，并企图使为安宁所安排的一切归于无效。如今，在贤弟城市中的\Place{卡斯蒂利安隐修院}的隐修士们向我们申诉，说你正采取措施，对上述隐修院施加某些与你的前任所允许并经长期习惯认可相抵触的事，并以一种有害的新奇手段扰乱古老的安排。为此，我们在此劝告贤弟，若果真如此，你当停止以任何借口烦扰此隐修院，且不可借任何篡夺之机，推翻它长期保有的权利；反而应不加辩驳地努力保全它的一切特权不受侵犯，并要知道，就上述隐修院而言，你所合法享有的，并不超过你的前任所合法享有的。\par

此外，他们同样申诉说，贤弟以所谓献仪为借口，从隐修院拿走了某些物品；因此，若你记得曾以不当方式收取过任何东西，必须立即归还，以免贪婪之罪严重定你的罪——你身为司铎，本应慷慨对待隐修院。因此，当你保全所有由你的前任所允许并保存的一切时，你当留心密切监督居住在那里的隐修士的行为和生活；若发现有人生活不当，或（愿天主禁止）犯有任何不洁之罪，就当以严格而正规的惩戒予以纠正。因为，我们既希望贤弟避免不当的篡夺，也同样劝告你要在一切有关纪律正直及灵魂守护的事上，悉心关切。\par

上述隐修院的隐修士们还告知我们，那个名为\Place{斯基拉西乌姆}的营地建在属于他们隐修院的土地上；因此，居住在那里的人书面承诺每年支付一笔\OriginalTerm{慰藉金}{solatium}；但后来他们却轻视此事，无故扣留了约定的付款。因此，请贤弟仔细查明真相；若果真如此，务必催促他们不能拖延支付他们所承诺的、以及理性所要求的款项；这样，他们既能安宁地拥有他们所持有的，隐修院的权益也不会遭受损害。\par

此外，上述隐修院的隐修士们向我们申诉说，他们的院长曾以建造教堂为借口，在\Place{斯基拉西乌姆}营地内，以赠与之名，向贤弟让与了六百尺的土地；因此，我们愿意，凡建成后的教堂墙壁所能围绕的土地，均应归属教堂所有；但教堂墙壁之外的任何土地，均应毫无争议地归还给隐修院。因为无论是世俗法律还是神圣教规的典章，都不允许隐修院的产业以任何名义从其所有权中分离。因此，请归还此项违背理性让与的土地赠与。\par

\SourceRecord{Translated by James Barmby.；Nicene and Post-Nicene Fathers, Second Series；Vol. 12.；Edited by Philip Schaff and Henry Wace.；Buffalo, NY: Christian Literature Publishing Co.,；1895.；<http://www.newadvent.org/fathers/360208034.htm>.}

% structure: p0001 source=h1 target=section reason=page-title

\section{卷八，第三十五函}

\Emph{致前执政官\Person{莱昂提乌斯}}\par

\Person{额我略}致\Person{莱昂提乌斯}等\par

既然\BibleQuote{在大户人家，不但有金器银器，也有木器瓦器；有作贵重的，有作卑贱的\BibleRef{弟后 2:20}}，谁能不知在普世教会的怀抱中，有些如卑贱的器皿被派作最低下的用途，而另一些如尊贵的器皿则适于清洁的用途？然而常发生这样的事：巴比伦的居民为耶路撒冷服苦役，而耶路撒冷的居民，即天国之城的居民，却被派往巴比伦服苦役。因为当天主所拣选的人，具备道德卓越、以节制著称、不寻求私利，却被派去处理世俗事务时，这岂非圣城耶路撒冷的居民在巴比伦服役？而当某些人在道德上放纵无度，却占据神圣职位，并在他们似乎做得好的事上寻求自己的荣耀时，这岂非巴比伦的居民在为天上的耶路撒冷服苦役？因为犹达斯混杂在宗徒中，长久宣讲人类的救世主，并与其他宗徒一同行奇迹；但因为他曾是巴比伦的居民，他为天上的耶路撒冷所执行的工作如同苦役。反之，若瑟被带到埃及，侍奉世俗的宫廷，承担了管理世俗事务的责任，尽其所能地处理了暂存王国的一切正义之事；但因为他仍是圣城耶路撒冷的居民，他管理巴比伦的职务，如前所见，仅仅是作为苦役而已。我相信你是他的追随者，善人，因我深知你虽投身于世俗行动，却以温和的精神行事，在各方面持守谦卑的堡垒，并给予每人应得的正义。因为你的荣耀有如此多的善行报告于我，我虽宁愿亲眼见而非听闻，但我仍因这善名而感到饱饫，尽管我不得亲见。然而，那从玉瓶倒香液的女人，代表了圣教会，即所有蒙选者的形象，使满室充满香液\EditorialNote{路7}。同样，每当我们听闻善人的事迹时，我们仿佛通过鼻孔吸入芬芳。当宗徒保禄说：\BibleQuote{我们是基督的馨香，献给天主\BibleRef{格后 2:15}}，这清楚表明他呈现自己为在场的滋味，为不在场的芬芳。因此我们虽不能由你在场的滋味得滋养，却因你不在的芬芳而得滋养。\par

为此我们也非常喜乐，因为你送来的礼物与你的品性相符。我们确实收到了圣十字架油和沉香木：前者以触摸祝福，后者在点燃时散发馨香。因为一个好人理应送来能平息天主对我们愤怒之物。\par

你还为我们的库房送来许多其他物品，因为我们既由灵魂和肉身组成，两方面都需要得滋养。然而在传递这些东西时，你最甘甜的心灵宣告感到极大的羞愧，并以爱德之盾抵挡这羞愧。但我对这些话完全喜乐，因为从这灵魂的见证，我知道一个即使在施予自己之物时也会羞愧的人，决不能夺取他人之物。你所送的礼物，虽你称为微小，却是伟大的；但我认为你的荣耀之谦逊更使其增大。你请求我欣然接纳。但同时请回忆那寡妇的两个小钱\EditorialNote{路27}。因为若她以善意献上微薄之物而中悦天主，为何那以谦卑之心献上丰厚之物的人不应中悦人呢？此外，我们送给你一个从圣伯多禄、宗徒之长最神圣坟墓取得的钥匙，其中嵌有他锁链的祝福，使那为殉道而束缚他颈项的，得以解开你的一切罪过。\par

\SourceRecord{Translated by James Barmby.；Nicene and Post-Nicene Fathers, Second Series；Vol. 12.；Edited by Philip Schaff and Henry Wace.；Buffalo, NY: Christian Literature Publishing Co.,；1895.；<http://www.newadvent.org/fathers/360208035.htm>.}

% structure: p0001 source=h1 target=section reason=page-title

\section{第九卷，第一封信}

\Emph{致\Person{雅努阿里乌斯}，\Place{卡利亚里}（Cagliari）主教。}\par

\Person{格列高利}致\Person{雅努阿里乌斯}等。\par

全能天主的宣讲者、宗徒保禄说：\BibleQuote{不可斥责老年人}（\BibleRef{弟前 5:1}）。但这条规则应在长者之过不通过其榜样将幼小者的心引向毁灭的情况下遵守。然而，当长者树立榜样使幼者走向毁灭时，就应用严厉的斥责来打击他。因为经上记载：\BibleQuote{你们全是幼小者的网罗}（\BibleRef{依 42:22}）。先知又说：\BibleQuote{百岁的罪人受诅咒}（\BibleRef{依 65:20}）。但是，关于你的年高，竟有如此大的恶行上报到我们这里，若非我们心怀仁爱，本当用明确的诅咒打击你。因为有人告诉我，你在主日、举行弥撒圣祭之前，竟出去犁毁这位持信人的庄稼，犁完后才举行弥撒圣祭。而且，在弥撒圣祭之后，你竟不怕拔除那块田产的地界。凡听闻此事者，都知道这样的行为应受何等惩罚。不过，我们曾怀疑你是否有这么大的悖逆；但我们的儿子、院长\Person{西里亚库斯}受我们询问后声明，他在\Place{卡利亚里}时确知此事。看，我们仍然宽恕你的白发，老人，请你终有所悟，自制于如此轻浮的行为和悖逆的事。你离死亡越近，就越该谨慎畏惧。确实，一项惩罚的判决已针对你发出；但因我们知道你的单纯与年迈共存，我们暂且保持沉默。然而，那些你听从其建议而做了这些事的人，我们判定他们被绝罚两个月；但若在这两个月期间按人之常情发生什么，他们不被剥夺\OriginalTerm{天路行粮}{viaticum}的祝福。但今后你要谨慎，远离他们的劝告，免得你在恶事上成为他们的门徒——你本应在善事上作他们的师傅——那时我们不再宽恕你的单纯或你的年迈。\par

\SourceRecord{Translated by James Barmby.；Nicene and Post-Nicene Fathers, Second Series；Vol. 13.；Edited by Philip Schaff and Henry Wace.；Buffalo, NY: Christian Literature Publishing Co.,；1898.；<http://www.newadvent.org/fathers/360209001.htm>.}

% structure: p0001 source=h1 target=section reason=page-title

\section{Book IX, Letter 2}

\Emph{致撒丁岛\OriginalTerm{护守者}{Defensorem}维大利。}\par

额我略致维大利等。\par

持此信者以及我们信函的副本将充分告知您关于我们的弟兄雅努雅略主教所知晓的事；因此，愿您的\OriginalTerm{阁下}{Experience}明智地执行我们所决定对其悖逆顾问宣布的绝罚，使他们因跌倒而学会不再鲁莽行事。\par

此外，我们已由持此信的\OriginalTerm{护守者}{defensorem}雷登普斯送还了以礼物名义寄给我们的麦子。愿您的\OriginalTerm{阁下}{Experience}注意，无论是您还是带来它的人，都不得从中取用任何东西作为赏赐，而应将全部麦子毫无减损地归还给各人，或一起归还给他们，并将他们的收据寄给我；因为，若我发现任何事没有按我所指示的去做，我将以不小的严厉惩罚这一过失。\par

\SourceRecord{Translated by James Barmby.；Nicene and Post-Nicene Fathers, Second Series；Vol. 13.；Edited by Philip Schaff and Henry Wace.；Buffalo, NY: Christian Literature Publishing Co.,；1898.；<http://www.newadvent.org/fathers/360209002.htm>.}

% structure: p0001 source=h1 target=section reason=page-title

\section{卷九·第三封信}

\Emph{致卡拉利斯（卡利亚里）主教雅努亚略。}\par

格列高利致雅努亚略等。\par

最尊贵的女士内雷依达向我们申诉，责备您的弟兄毫不羞愧地向她索取一百个\OriginalTerm{苏勒德斯}{solidi}作为其女儿的丧葬费，并要在她的悲伤呻吟之外，再增加费用的烦恼。如果事实果真如此——因为索要已然腐烂之躯的泥土的代价，并欲从他人的悲痛中牟利，实属严重之事，且与司铎的职分相去甚远——那么，请您的弟兄停止此项索求，不要再打扰她；尤其她告诉我们，那位她声称与之生下此女的霍尔图拉努斯，曾对您的教会慷慨奉献，为数不少。至于这一恶习，我们在天主的准许下就任主教之尊后，便在自己的教会中彻底禁止，绝不允许这不良习俗重新兴起，因为我们记得，当亚巴郎向哈毛尔之子，即厄斐龙，则哈尔之子，请求购买墓地埋葬其妻子遗体时，后者拒绝接受代价，以免显得他从尸体中牟利\EditorialNote{创世纪23}。如果一位外邦人都表现出如此大的体恤，那么我们这些被称为司铎的人，岂不更不应做这样的事？因此我劝诫您，即使是面对陌生人，也不可再冒这出于贪婪的恶习。但是，若有时您允许某人在您的教会中安葬，而该人的父母、亲属或继承人自愿为灯烛献上一些东西，我们不禁止接受。但我们完全禁止任何索求或强取，因为这是极其不虔敬的行为，以免（但愿天主不许）教会被人说成是买卖之地，或您似乎因他人的死亡而欣喜，倘若您以任何方式企图从死尸中牟利。\par

关于上述内雷依达请愿书中包含的其他案件，我们劝您，如有可能，以和解的方式解决；否则，务必不要遗漏派遣一位有知识的人士，由我们委派，前往法庭。为此，我们已派遣我们的\OriginalTerm{辩护人}{defensorem}、此信的携带者雷德姆普图斯前往您所在地区，以便他强制双方出庭受审，并对可能作出的判决立即执行。\par

\SourceRecord{Translated by James Barmby.；Nicene and Post-Nicene Fathers, Second Series；Vol. 13.；Edited by Philip Schaff and Henry Wace.；Buffalo, NY: Christian Literature Publishing Co.,；1898.；<http://www.newadvent.org/fathers/360209003.htm>.}

% structure: p0001 source=h1 target=section reason=page-title

\section{\WorkTitle{第九卷，第四函}}

\Emph{致\Person{雅努阿里乌斯}，\Place{卡拉利斯}（\Place{卡利亚里}）主教。}\par

\Person{额我略}致撒丁岛主教\Person{雅努阿里乌斯}。\par

我们早在您的兄弟的书信到达我们之前，就已知道我们的敌人在撒丁所行的事。而且，我们一段时间以来一直害怕事情会如此，现在我们与你们一同叹息我们所预见的事情已经发生。但是，如果我们写给我们的最杰出的儿子\Person{根纳狄乌斯}以及给您本人的信中所说的事情——告诉你们事情会如此——被注意到了，那么敌人本不会进入你们的地区，或者当他们进入时，他们本会遭到他们所造成的那种危险。即使现在，也要让所发生的事为将来磨砺你们的警觉。因为我们也绝不忽略我们所能为善所做的任何事，主帮助我们。\par

此外，要知道，我们在相当久以前派往\Person{阿吉卢尔夫}那里的院长，已因天主的仁慈与他达成了和平，正如最杰出的总督在书面指示中所指定的。因此，在直到为确认这项和平的协议被拟定之前，以免我们的敌人趁着目前的延迟可能再次倾向于进入那些地区，你要使人保持城墙的警戒，并在各处给予仔细注意。我们信赖我们救赎者的权能，我们的敌人的入侵或阴谋不会再伤害你们。\par

至于你在信中说许多人在我们面前向你提出控诉，这是真的；但在各种事情中，没有什么比我们所最亲爱的儿子、院长\Person{基里亚库斯}向我们报告的更让我们痛心了，即：你在主日、弥撒前，让人在属于\Person{多纳图斯}的田地里翻耕谷物，并且，似乎这还不够，在祭献结束后，你亲自前往该地，挖掘了地界。为此，我劝告你要以焦虑的关注思虑你所担任的职务，完全避免任何可能损害你名誉或你灵魂的事情，并且不要让任何人说服你再做类似的事。因为要知道，你没有承担世俗事物的管理，而是灵魂的领导。因此，你应当将你的心、你的忧虑、你的全部热爱固定在这上面，并勤勉思考赢得灵魂的事，以便当你在主的来临时将祂托付给你的\OriginalTerm{塔冷通}{talents}倍增后交还给祂时，你可以被认为配得从祂那里接受报偿的果实，并在永恒的光荣中，在祂忠信的仆人中受举扬。然而，要知道，我现在所说责备或责难的话，并非出于严厉，而是出于兄弟之爱，因为我希望你被找到时是一位在全能天主面前的司铎，不仅在名义上（那只会导致惩罚），而且在功德上（那是为了报答）。因为我们在我们救赎者的身体内是一个肢体，正如我在你的过失中被撕裂，我也在你的善行中喜乐。\par

此外，关于你希望我们从我们身边（\OriginalTerm{从我们身边}{a nostro latere}）派遣一个人，你可以向他详细传达需要提交给我们的案件，请将你的任何意愿写信给我们最亲爱的儿子\Person{伯多禄}和顾问（\OriginalTerm{顾问}{consiliario}）\Person{狄奥多尔}，以便当通过他们传达给我们后，凡理性所赞成的，便可得到解决，由主启示道路。再者，关于我们的兄弟和同为主教的\Person{马里尼亚努斯}，在和平与前述\Person{阿吉卢尔夫}完全确认后，将会处理，并做任何理性秩序所指示的。\par

\SourceRecord{Translated by James Barmby.；Nicene and Post-Nicene Fathers, Second Series；Vol. 13.；Edited by Philip Schaff and Henry Wace.；Buffalo, NY: Christian Literature Publishing Co.,；1898.；<http://www.newadvent.org/fathers/360209004.htm>.}

% structure: p0001 source=h1 target=section reason=page-title

\section{第九卷·第五函}

\Emph{致达尔马提亚代理总督玛策禄}。\par

格里高利致玛策禄等。\par

我们已经收到阁下的来信，你在信中提到你得罪了我们，并希望直接做补赎以蒙我们悦纳。诚然，我们听说过有关阁下的事，这些事绝不应是一个忠信之人所犯的。因为所有人断言，在玛克西穆的事件中，你是那一切重大祸端的始作俑者，而那座教会的劫掠、如此众多灵魂的丧亡以及那闻所未闻的狂妄胆量，都源自于你。的确，关于你寻求蒙我们悦纳，你理当全心全灵、流着泪，如同你所当行的，为这些事向我们的救赎主做补赎；因为，除非向他做了补赎，我们的宽恕或恩待又能给你什么确定的好处呢？然而，当我们看到你仍牵连在僭越者的毁灭性行为中，或为那些迷失者辩护时，我们看不出你对天主或对人所做的补赎是哪一类的。因为那时你的阁下才会知晓，当你将偏离的带回正直，将狂妄的带回谦逊的规诫时，你便公开且明显地满足了天主和人。如果这事做到了，你便可知道，你便会如此蒙天主和人的悦纳。\par

\SourceRecord{Translated by James Barmby.；Nicene and Post-Nicene Fathers, Second Series；Vol. 13.；Edited by Philip Schaff and Henry Wace.；Buffalo, NY: Christian Literature Publishing Co.,；1898.；<http://www.newadvent.org/fathers/360209005.htm>.}

% structure: p0001 source=h1 target=section reason=page-title

\section{第九卷，第六封信}

\Emph{致\Person{雅努阿里乌}，卡利亚里主教。}\par

\Person{额我略}致\Person{雅努阿里乌}等。\par

那些从你们城市来到这里的犹太人向我们投诉，说伯多禄——他已被天主的旨意从他们的迷信带到了基督信仰的敬拜中——在受洗后的第二天，即复活节庆期的主日，带着一些无秩序之人，严重恶表地、未经你们同意，强占了他们在卡利亚里的犹太会堂，并在那里放置了我们的天主和主之母的像、可敬的十字架，以及他从洗礼池上来时所穿的\OriginalTerm{白衣}{birrum}。关于此事，我们的儿子、光荣的\OriginalTerm{军长}{Magister militum}欧帕特利乌斯和敬虔的、在主内虔诚的尊贵总督的信函也一致证实同样的事。他们还补充说，您已预见到此事，并禁止上述伯多禄尝试去做。得知此事，我们完全称赞您，因为，正如真正善牧当有的，您不愿有任何事发生以致招来合理的指责。但是，由于您完全没有卷入这些错误行为，显示您不悦于所行之事，我们，鉴于您在这事上的心意倾向，尤其您的判断，兹劝勉您，以合宜的敬意移走那里的像和十字架，归还被强夺之物；因为，正如法律不准许犹太人建造新的会堂，同样也允许他们不受干扰地保留旧的。以免上述伯多禄或在此不法混乱行为中给予他帮助或纵容的其他人，回答说他们是出于信仰的热忱而为之，为借此把归化的必要性强加给犹太人；应告诫他们，他们当知道，对他们应使用温和态度；如此才能从他们那里引出愿意服从的意志，而不是违背其意愿将他们带进来。因为经上记载：\BibleQuote{我要向祢自愿献祭} \BibleRef{咏58:8}；并且：\BibleQuote{我要自愿向祂忏悔} \BibleRef{咏27:7}。所以，您的圣座，请带着与您一同不悦这些事的儿子们，努力促进您城中居民的和睦，因为尤其在此时，当有来自敌人的警报时，你们不应有分裂的民众。但是，我们为自己担心同等也为你们担心，我们认为应当让您明白，当目前休战结束后，阿吉卢尔夫王不会与我们和好。因此，您贤弟必须注意加固你们的城或其他地方使之更安全，并认真关注在其中储备粮草，这样，当敌人——天主义怒临于他——到来时，他将无法造成伤害，反会狼狈撤退。但是我们也尽力为你们着想，并敦促那些有责任的人准备抵抗，因为，正如你们视我们的苦难为你们自己的，同样我们也视你们的苦难为我们自己的。\par

\SourceRecord{Translated by James Barmby.；Nicene and Post-Nicene Fathers, Second Series；Vol. 13.；Edited by Philip Schaff and Henry Wace.；Buffalo, NY: Christian Literature Publishing Co.,；1898.；<http://www.newadvent.org/fathers/360209006.htm>.}

% structure: p0001 source=h1 target=section reason=page-title

\section{第九卷，第七书函}

\Emph{致\Place{卡利亚里}（Cagliari）主教\Person{雅努亚里}。}\par

\Person{额我略}致\Person{雅努亚里}等。\par

法律明确规定，为过隐修生活而进入隐修院者，不再有立遗嘱的自由，其财产应归该隐修院所有。此事几乎尽人皆知，但我们十分惊讶地接到圣加维诺与圣路索里隐修院院长加维尼亚的通报，称其隐修院院长西里卡在获得管理职务后，立下遗嘱，遗赠若干财产。当我们向您圣座询问为何容忍属于隐修院的财产被他人扣留时，我们共同的儿子，您的总铎埃皮法尼乌斯亲临我们面前，回答说，该院长直至去世之日拒绝穿着隐修会衣，而一直穿着该地女司铎所穿服装。对此，上述加维尼亚回答说，此做法因习惯而几乎成为合法，声称在上述西里卡之前的院长也曾穿着此类服装。当我们对服装的性质开始颇为怀疑时，似乎有必要与我们的法律顾问及本城其他博学之士一同考虑，在法律上应当如何处理。他们经审议后答复：一位院长经主教隆重祝圣，并多年主持隐修院管理直至离世，其服装性质可能归咎于主教允许如此，但仍不能对隐修院造成损害；她的财产自她进入隐修院并被任命为院长之时起，显明地属于该隐修院。因此，既然她（即院长加维尼亚）声称一所旅客招待所（\OriginalTerm{旅客招待所}{xenodochium}）不正当地占有了非法遗赠的财产，我们在此劝勉你，鉴于隐修院与招待所均位于你的城中，务必以一切谨慎与勤勉作出安排，以便若此占有并非源于先前的契约，而是源于上述西里卡的遗赠，则该财产应毫无争议或推诿地归还给该隐修院。但若偶然有称其来自另一契约，则要么由阁下查明双方实情，按法律秩序所需作出裁定；要么由双方合意选出仲裁人，以裁定双方的诉求。凡他们所作裁决，均应在您的督导下遵守，以免两个可敬场所之间存留任何怨恨；两者本应全然以彼此平安与和睦为重。因此，凡上述西里卡遗嘱下扣留的其他一切财产，既然无一为法律所许可，就必须通过阁下司铎般的悉心照料，审慎归还隐修院占有；因为帝国宪章明确规定，凡违反法律所为，不仅无效，且应被视为从未发生。\par

\SourceRecord{Translated by James Barmby.；Nicene and Post-Nicene Fathers, Second Series；Vol. 13.；Edited by Philip Schaff and Henry Wace.；Buffalo, NY: Christian Literature Publishing Co.,；1898.；<http://www.newadvent.org/fathers/360209007.htm>.}

% structure: p0001 source=h1 target=section reason=page-title

\section{第九卷，第八封信}

\Emph{致\Place{撒丁尼亚}诸位主教。}\par

\Person{额我略}致\Person{味增爵}、\Person{依诺增爵}、\Person{马里尼亚努斯}、\Person{利伯提努斯}、\Person{阿加托}和\Person{维克多}，\Place{撒丁尼亚}诸位主教。\par

我们得知，在你们的岛上，在复活节后，你们习惯前往或派代表到你们的教省主教那里，而他则会以书面通知告知你们关于下一个复活节的事宜，无论你们是否知道时间。而且，据报告，你们中有些人疏忽了按习惯这样做，还使其他人的心也转向不服。还有，你们中有些人在处理涉及他们教会的事而前往海外时，冒险在未告知上述教省主教或未从他那里获得正统典章规定的信件的情况下旅行。因此，我们劝告你们的兄弟友爱，要符合你们教会的习惯，既关于复活节的宣布，又如果有需要迫使你们中任何人因自己的事务前往任何地方，要根据加于你们的规定向你们的教省主教请求许可；除非（我们希望不会发生）你们碰巧与你们的教省主教有讼案，那么那些因此而急于寻求宗座审判的人，可以这样做，正如你们所知，古教父的制度甚至在典章中也是允许的。\par

\SourceRecord{Translated by James Barmby.；Nicene and Post-Nicene Fathers, Second Series；Vol. 13.；Edited by Philip Schaff and Henry Wace.；Buffalo, NY: Christian Literature Publishing Co.,；1898.；<http://www.newadvent.org/fathers/360209008.htm>.}

% structure: p0001 source=h1 target=section reason=page-title

\section{卷九，书信九}

\Emph{致\Person{卡利尼库斯}，意大利总督。}\par

\Person{额我略}致\Person{卡利尼库斯}等。\par

在你向我宣布的你战胜斯拉夫人的消息中，我知道自己因以下事而大感欣慰：这些信件的携带者，从\Place{卡普里塔纳}岛赶来加入圣教会的合一，被阁下您送至真福\Person{伯多禄}，\OriginalTerm{宗徒之长}{Prince of the Apostles}。因为藉此您将更胜过仇敌，若您使那些您所知为天主仇敌的人屈服于真正的主的轭下；并且您以真诚虔敬的心维护天主的事业，将更有效地在人间办理您的事务。\par

现在，关于您希望让我看到一份为分裂者辩护而发给您的命令副本，您这位对我最亲切的阁下应当仔细考虑：尽管那道命令已被争取到，您在其中仍未被命令拒绝那些前来加入教会合一的人，而只是在这动荡时期，不强迫那些不愿意来的人。因此，您必须尽快将这些事禀报我们最虔诚的皇帝，好让他们知道：在他们时代，藉全能天主的助佑和您的努力，分裂者们正自愿赶回归。关于我在\Place{卡里塔纳}岛上事务安排的决定，阁下将通过我们最可敬的兄弟和同僚主教\Person{马里尼亚努斯}得知。但我要您知道，这件事使我感到不小的忧虑：您的总管，他负责那位希望回归的主教的请愿书，声称他丢失了它，而且后来他被教会的反对者抓住了；依我之见，这一行为并非由于他的疏忽，而是由于他的贪婪。因此我感到奇怪，阁下竟完全没有追究他的过错。但我随即责备自己对此感到奇怪，因为哪里大人\Person{尤斯提努斯}提供建议，哪里就不能控告异端者。\par

此外，您告诉我们，您希望在\Place{罗马}城举行\Person{伯多禄}宗徒之长的周年纪念。我们祈求全能天主以祂的仁慈保护您，并成全您的愿望。但我恳求，上述那位极善辞令的人能与您同来，或者，如果他不来，他应退出对您的侍奉。或者，的确，如果阁下因可能发生的事务不能前来，要么让他与圣教会的合一共融，要么我恳求他不要参与您的议事。因为我听说他是一个好人，要不是陷入最有害的错误。至于\Person{马克西姆}的案件，既然我们再也无法抵挡您的恳求，您将从公证人\Person{卡斯托里乌斯}处得知我们的决定。\par

\SourceRecord{Translated by James Barmby.；Nicene and Post-Nicene Fathers, Second Series；Vol. 13.；Edited by Philip Schaff and Henry Wace.；Buffalo, NY: Christian Literature Publishing Co.,；1898.；<http://www.newadvent.org/fathers/360209009.htm>.}

% structure: p0001 source=h1 target=section reason=page-title

\section{第九卷，第十封信}

\Emph{致拉文纳主教马里尼亚努斯}\par

\Person{格列高利}致\Person{马里尼亚努斯}等。\par

持此信者，即最显赫者副主教与辩护者，前来向我们声称：有一位名叫\Person{若望}的主教，来自\Place{潘诺尼亚}，已被安置在一座名为\Place{诺瓦埃}的城堡中，而该城堡已将其所属的\Place{卡普里塔纳}岛作为教区附属于其管辖。他们补充说，这位主教被强行逐出同一城堡后，另有一人在那里被祝圣；然而，关于此人，他们声称已决定他不应居住在上述城堡中，而应居住在他自己的岛上。他们进一步说，当他与他们一同住在那里时，他不愿继续留在裂教错误中，并与全体民众向我们的至杰出之子总督\Person{卡利尼库斯}呈递了一份请愿书，渴望与他所有跟随者一同重新加入天主教会，正如我们已经说过的那样。但他们说，后来他被裂教徒说服而反悔，如今上述岛屿的所有居民都失去了主教的护佑，因为当他们渴望与圣教会合一时，现在却无法接纳那位已转向裂教错误的人；他们希望另有一人被祝圣给他们。但我们在必须严格且彻底审查一切的情况下，已下令采取以下措施：即你的爱德应派遣使者到该主教那里，劝诫他返回天主教会合一并回到自己的百姓中。若劝诫后他仍不屑返回，则天主的羊群不应在其牧者的错误中被欺骗；因此，你的圣座应在该情况下在那里祝圣一位主教，并将上述岛屿作为他的教区，直到\Place{伊斯特里亚}的主教们回归大公信仰为止；如此，每座教会都能保留其自身教区的权利，而失去牧者的百姓也不致缺乏管理与监督的保护。然而，在所有这些事情上，你的爱德务要警惕小心，确保这回归教会的百姓受到十分勤勉的劝诫，以使他们在回归中坚定稳固，免得因摇摆不定的思想而跌回错误的深渊。但请务必恳请至杰出的总督在他的信函中将这些事禀报皇帝陛下的至虔敬之耳，因为尽管已向他发出的命令似乎是由他们那里得到的，但该命令并未禁止他允许那些愿意的人回归教会，而只是目前不得强迫不情愿者。因此，让我们的上述儿子承担起处理此事的管理职责，以便他能如此措辞他的报告，使得他所规定的一切都无疑问。然而，我们也亲自写信给我们共同的儿子\Person{阿纳托利乌斯}，命令他详尽地将这些事禀报给至虔敬的君主们。\par

我多次收到至杰出之子、总督\Person{卡利尼库斯}大人为\Person{马克西穆斯}之事而写的迫切信件。被他的纠缠所压倒，我看除了将\Person{马克西穆斯}的事托付给你的爱德之外，别无他法。因此，若这位\Person{马克西穆斯}来到你的爱德处，也请让他教会的总执事\Person{霍诺拉图斯}到场；你的圣座便可查明他是否被合法祝圣，是否未犯任何西满异端，是否在身体越轨方面无可指责，是否在明知自己被绝罚时仍妄自举行弥撒；凡你在敬畏天主之事上认为合宜的，请予以裁定，好使我们顺应你的安排，在天主之下予以认可。然而，若我们的上述儿子对你的爱德有所猜疑，也请我们至可敬的兄弟、\Place{米兰}主教\Person{康斯坦提乌斯}来\Place{拉文纳}与你一同商议；你们应共同判决此事：凡你们二人看来合宜的，请确信在我也看来合宜。因为正如我们不应向谦卑者固执，同样也应向骄傲者表现严厉。因此，你的爱德，正如你从圣经篇页中所学到的，请在此时按你视为公正的来判决此事。\par

\SourceRecord{Translated by James Barmby.；Nicene and Post-Nicene Fathers, Second Series；Vol. 13.；Edited by Philip Schaff and Henry Wace.；Buffalo, NY: Christian Literature Publishing Co.,；1898.；<http://www.newadvent.org/fathers/360209010.htm>.}

% structure: p0001 source=h1 target=section reason=page-title

\section{卷九 第十一封}

\Emph{致布伦希尔德王后}\par

\Person{格列高利}致\Person{布伦希尔德}，法兰克人的王后。\par

阁下的心如何坚定地定于对全能天主的敬畏，你在所做的其他善事中，也通过你对祂的司铎的爱，以可赞美的方式显示出来；我们因你的基督信仰而大感喜乐，因为你努力以荣誉提拔那些你因他们真正是基督的仆人而爱戴并尊敬的人。因为最优秀的女儿，这适合你，这适合你，使你能够服从一位在你之上的主。因为当你将你心灵的颈项顺服于全能的主的敬畏时，你也确认你对附属民族的统治；并且通过使自己服从于造物主的服务，你更忠诚地将你的臣民与你捆绑在一起。因此，收到你的书信后，我们告知你，阁下的热切愿望大大取悦了我们；并且我们一直渴望将大司牧披肩送给我们的兄弟和同主教\Person{西格里乌斯}，因为我们的最威严的皇帝陛下的心意也是赞成的；并且，据我们的执事（他是我教会在皇帝朝廷的代表）告知，他完全希望此物能被赐予。关于我们前述兄弟，我们收到了许多良好报告，既有来自你的见证，也有来自他人的；尤其我们从区域执事\Person{若望}返回我们这里时得知了他的生活如何。并且听到他在我们兄弟\Person{奥斯定}的事上所做的一切，我们赞美我们的救赎主，因为我们感到他用自己的行为履行了他司铎之名的意义。\par

但是有许多障碍暂时阻止了我们做这件事。首先，那位前来接受这个大司牧披肩的人被牵连在分裂者的错误中；其次，你希望让人以为它是主动送出的，而非应你的请求。此外还有：那个希望使用它的人并未通过专门写给我们的请愿书来请求赐予他；并且未经他的请求就同意如此重大之事对我们来说绝不合理，尤其因为一项古老习俗已经确立，即大司牧披肩的尊严不应被赐予，除非案件的是非曲直要求如此，且是给一个急切请求的人。然而，免得我们似乎想以任何借口推迟阁下的愿望，我们已经安排将大司牧披肩送给我们最亲爱的儿子司铎\Person{坎迪杜斯}，嘱咐他以适当的谨慎代表我们交付。因此，我们上述的兄弟和同主教\Person{西格里乌斯}必须期待它，当他主动与他的几位主教一起起草一份请愿书时；他必须将此交给前述司铎，以便他能够在天主的恩宠下恰当地获得使用同一大司牧披肩的权力。\par

因此，为了使你所承担的责任在我们造物主眼中为你结出果实，让你基督徒的忧勤谨慎地警惕，不允许任何在你统治之下的人通过金钱给予、任何人的庇护或亲属关系而获得圣秩；而应选举这样的人为主教或任何其他神圣职务，即其生活与品行已证明是相称的；免得如果（正如我们所不愿见到的）司铎的尊严成了可买卖的，那么西满异端，这首先在教会中出现的并被教父们的判决所谴责的，会在你们的地区兴起，并且（愿天主禁止）削弱你们王国的能力。因为出卖那救赎万物的圣神，是严重的事，是难以言表的邪恶。\par

但也要注意这一点：既然，如你所知，那位卓越的宣讲者完全禁止一个初学之人登上司铎的统治地位，你就不许任何人从平信徒被祝圣为主教。因为那未曾做过门徒的人会成为怎样的师傅呢？或者，那先前未曾服从于牧人的管教的人，能给主的羊群提供什么样的领导呢？因此，如果任何人的生活表明他配得提升到这一等级，他应该先在教会的职务中服务，以便通过长期实践的经验，他能够看到要效法什么，并学会要教导什么；免得他职位的崭新无法承受治理的重担，并因他晋升的不成熟而导致毁灭的机会。\par

此外，阁下如何对待我们的兄弟和同主教\Person{奥斯定}，以及通过天主的感动，你赐予他多么大的爱德，我们从多位信友的叙述中得知；为此我们表示感谢，并恳求天主大能的慈悲在此保护你，并使你如同在人间统治一样，也在度过许多年岁之后在永生中统治。\par

此外，你们当为了自己的赏报，努力将那些被裂教徒的错误从教会合一中分离的人，召回到和谐的合一之中。因为，他们之所以如此深陷于无知的黑暗，无非是为了逃避教会的纪律，并得以随心所欲地过悖逆的生活——他们既不理解自己所辩护的，也不理解自己所追随的。至于我们，我们尊敬并完全遵循加采东大公会议，而他们却从那里为自己获取了毒害性的借口乌云；若有人胆敢对这会议的信仰加以删减或增添，我们必弃绝他。但他们如此被错误的污点所浸透，以至于信从自己的无知，非理性地、恶意地拒绝了普世教会和四位宗主教；以致由你们尊贵所派遣到我处者，当我们问他为何与普世教会分离时，他承认自己并不知道。他所说的，以及他所听闻的其它事，他都没有能力知道。同样，我们也劝勉你们：你们当以纪律约束其余臣民，禁止他们向偶像献祭、崇拜树木、或行亵渎性的牲畜首级献祭；因为我们听说，许多基督徒既去教堂，（可怖的是！）又不放弃对魔鬼的崇拜。然而，既然这些事完全令我们的天主不悦，祂也不接纳分裂的心思，你们当设法使他们从这些不法行为中得到有益的约束；免得（愿天主不许！）圣洗圣事不仅不能拯救他们，反而成为对他们的惩罚。因此，如果你们知道有暴虐者、有奸淫者、有偷窃者、或从事其它恶行者，务请速速以他们的悔改来平息天主，免得祂将不忠信民族应受的惩罚加于你们身上——就我们所见的，那惩罚已为惩治许多民族而举起；免得（我们不信会发生的事）因恶人的行为而点燃天主义怒的怒火，战争之灾毁灭那些天主的诫命未能召回到正道上的罪人。所以，我们必须以极大的热忱和不断的祈祷，急切地投奔我们救赎主的仁慈，那里是所有人在安全与极大保障中的避难所。因为，凡坚稳地留在那里的人，危险不能压碎他，恐惧也不能惊扰他。\par

我们已遵您来信所愿，将圣卷送至上述我们极爱的儿子、司铎坎迪杜斯，由他呈献于您；我们渴望共享您的善志。愿全能的天主将您置于祂的保护之下，并以祂伸出的手臂保卫您的王国免于不信的民族，并在漫长岁月之后引领您进入永恒的喜乐。\newline{}颁于十月，第一小纪。\par

\SourceRecord{Translated by James Barmby.；Nicene and Post-Nicene Fathers, Second Series；Vol. 13.；Edited by Philip Schaff and Henry Wace.；Buffalo, NY: Christian Literature Publishing Co.,；1898.；<http://www.newadvent.org/fathers/360209011.htm>.}

% structure: p0001 source=h1 target=section reason=page-title

\section{第九卷，第十二封信}

\Emph{致锡拉库萨主教若望}\par

额我略致若望等。\par

有位来自西西里的人告诉我，他的一些朋友——不知是希腊人还是拉丁人——仿佛因对圣罗马教会的热忱而对我关于礼仪的安排有所怨言，\Quote{他怎能如此安排以遏制君士坦丁堡教会，却在各方面追随她的用法？} 我对他说：\Quote{我们追随了她的哪些用法？} 他答道：\Quote{你们在五旬期以外的弥撒中诵念了阿肋路亚；你们命副执事不穿祭衣进行进；你们诵念垂怜经；并立即在正典后诵念天主经。} 我回答他说，在这些事中，我们没有一件是追随了其他教会。\par

因为，我们这里诵念阿肋路亚的习惯，据说是由真福热罗尼莫的传统、在真福教宗达玛苏时代从耶路撒冷教会传来的；因此，在这件事上，我们反而缩减了先前从希腊人那里继承的用法。\par

再者，关于我命副执事不穿祭衣进行进，这是教会的古老用法。但某位教宗——我不知道是哪位——曾命他们穿着亚麻长衣进行进。难道你们的教会曾从希腊人那里接受过任何传统吗？那么，如今副执事穿着亚麻长衣进行进的习惯，若不是从她们的母亲罗马教会领受的，又是从哪里来的呢？\par

再者，我们过去和现在诵念垂怜经的方式都不像希腊人那样：因为希腊人是全体一起念；而在我们这里，是由圣职人员诵念，民众回应；每次诵念时也诵念基督垂怜，这在希腊人那里根本不念。此外，在日常弥撒中，我们省略了一些通常诵念的部分，只念垂怜经、基督垂怜，以便在这些恳求的言辞上稍作停留。至于\OriginalTerm{天主经}{orationem Dominicam}，我们在\OriginalTerm{祈祷后立刻}{mox post precem}诵念，原因是宗徒们的习惯是在同一祈祷中祝圣献祭的祭品。而我认为，由一位学者编撰的祈祷文用在祭品上，却不诵念我们的救赎主为祂的圣体圣血所作的祈祷，这是很不合适的。此外，天主经在希腊人中是由全体民众诵念，而在我们这里则只由司铎独念。那么，我们在哪方面追随了希腊人的习惯呢？我们要么修改了旧的习惯，要么制定了新的有益的习惯，但并未显示我们在模仿他人。因此，请你的爱德找机会前往卡塔尼亚教会；或在锡拉库萨教会中教导那些你相信或认为可能对此事有所怨言的人，在那里召开会议，仿佛为了另一目的，不要停止教导他们。因为关于他们对君士坦丁堡教会的言论，谁能怀疑那教会是隶属于宗座的呢？正如最虔诚的君主皇帝和我们那位身为兄弟的该城主教所不断承认的。然而，如果这个或其他教会有任何好的事物，我准备在好的事上甚至效法我的下属，同时禁止他们做非法的事。因为，那种自以为居首而耻于学习所见一切好事的人是愚蠢的。\par

\SourceRecord{Translated by James Barmby.；Nicene and Post-Nicene Fathers, Second Series；Vol. 13.；Edited by Philip Schaff and Henry Wace.；Buffalo, NY: Christian Literature Publishing Co.,；1898.；<http://www.newadvent.org/fathers/360209012.htm>.}

% structure: p0001 source=h1 target=section reason=page-title

\section{第九卷，第十七封信}

\Emph{致\Person{德梅特里安}与\Person{瓦莱里安}。}\par

\Person{额我略}致\Person{德梅特里安}与\Person{瓦莱里安}，\Place{费尔莫}（Fermo）的文书。\par

神圣教规的法令及法律权威均允许，教会财产可合法用于赎回俘虏。因此，我们获悉，近十八年前，最可敬的\Person{法比乌斯}（Fabius），已故的\Place{费尔莫}教会主教，为赎买你及你的父亲\Person{帕西乌斯}（Passivus）——现为我们的兄弟及同为主教，当时仍是文书——以及你的母亲，曾支付该教会十一磅银子予敌人。你因此有所担忧，唯恐所付之款日后会被向你们追讨。我们认为宜藉此命令的权威消除你的疑虑，规定你及你的继承人从此不得因追讨此款而遭受任何烦扰，亦不得有任何人对你提起诉讼。因为公平原则要求，以虔诚意愿支付之物，不应给被赎回者带来负担或痛苦。\par

\SourceRecord{Translated by James Barmby.；Nicene and Post-Nicene Fathers, Second Series；Vol. 13.；Edited by Philip Schaff and Henry Wace.；Buffalo, NY: Christian Literature Publishing Co.,；1898.；<http://www.newadvent.org/fathers/360209017.htm>.}

% structure: p0001 source=h1 target=section reason=page-title

\section{第九卷，第十八封书信}

\Emph{致\Person{罗马努斯}，护卫者（\OriginalTerm{护卫者}{Defensorem}）。}\par

\Person{格里高利}致\Person{罗马努斯}等。\par

我们为面前的目标所怀的关切，促使我们将教会事务的照管托付给勤勉的人。因此，既然我们发现你，\Person{罗马努斯}，是一位可靠而勤勉的护卫者，我们以为适合从本第二小纪起，将我们因天主的仁慈所服务的圣罗马教会的产业，托付给你管理，这些产业位于\Place{锡拉库萨}、\Place{卡塔纳}、\Place{阿格里真托}和\Place{米莱}地区（\OriginalTerm{米莱地区}{partibus Milensibus}）。因此你必须立即前往那里，要顾及天主的审判，并记住我们的劝诫，努力使自己表现得如此得力而忠信，以致不会因疏忽或欺诈而招致危险——愿天主禁止如此。但你当更加如此行，好因你的忠信和勤勉而蒙受天主恩宠的赞许。此外，我们已照惯例向同一产业的\OriginalTerm{家族}{familia}下达了命令，以免有任何事阻碍你执行所托付给你的任务。\par

\SourceRecord{Translated by James Barmby.；Nicene and Post-Nicene Fathers, Second Series；Vol. 13.；Edited by Philip Schaff and Henry Wace.；Buffalo, NY: Christian Literature Publishing Co.,；1898.；<http://www.newadvent.org/fathers/360209018.htm>.}

% structure: p0001 source=h1 target=section reason=page-title

\section{卷九，第19封信}

\Emph{致锡拉库萨产业上的\OriginalTerm{佃农}{Colonos}}。\par

\Person{额我略}致\OriginalTerm{佃农们}{Coloni}等。\par

我愿你们知道，我们已安排将你们置于我们的\OriginalTerm{护卫者}{defensoris}的照管之下。因此，我们命令你们毫不迟疑地服从他，在他认为适宜并为教会的利益而吩咐你们去做的事上。我们授予他这样的权力，使他能对任何试图不顺从或倔强的人施加严厉惩罚。我们也同样嘱咐他立即注意收回任何躲藏在其辖区之外的奴隶，或任何越界侵入的人，将其重新置于教会管辖之下。因为要知道，他已被警告，如冒险行事，绝不可在任何借口下，对属于他人的东西行不义或掠夺。\par

\SourceRecord{Translated by James Barmby.；Nicene and Post-Nicene Fathers, Second Series；Vol. 13.；Edited by Philip Schaff and Henry Wace.；Buffalo, NY: Christian Literature Publishing Co.,；1898.；<http://www.newadvent.org/fathers/360209019.htm>.}

% structure: p0001 source=h1 target=section reason=page-title

\section{第九卷，第二十三书}

\Emph{致\Person{若望}，\Place{锡拉库萨}主教。}\par

\Person{额我略}致\Person{若望}等。\par

我们的儿子，显赫的\OriginalTerm{前执政官}{exconsul}\Person{莱昂提乌斯}向我们对我们的兄弟及同为主教\Person{良}提出严重控诉；他的控诉使我们完全不安，因为一位主教不应如此轻率和鲁莽行事。此案我们已委托我们的\OriginalTerm{护慰者}{defensoris}\Person{罗马努斯}，当他到你那里时，详细调查。此外，由他（即\Person{莱昂提乌斯}）所派的差役控诉你的爱德，称在辩护杰出的医生\Person{阿基劳}时，我们的兄弟和同为主教，都主教\Person{多米提安}的利益受损。你的爱德确实应当公正地保护你的儿子们，或在此案中保护圣教会的利益，不给仇敌诋毁的借口。然而，我毫不怀疑，即使我这样说，你也注意此事：但我们已吩咐同一位\Person{罗马努斯}，当他到你那里时，与你一起安排妥善处理此案。\par

\SourceRecord{Translated by James Barmby.；Nicene and Post-Nicene Fathers, Second Series；Vol. 13.；Edited by Philip Schaff and Henry Wace.；Buffalo, NY: Christian Literature Publishing Co.,；1898.；<http://www.newadvent.org/fathers/360209023.htm>.}

% structure: p0001 source=h1 target=section reason=page-title

\section{第九卷，书信第二十四}

\Emph{致\Person{罗玛诺斯}，\OriginalTerm{监护人}{Defensorem}。}\par

\Person{额我略}致\Person{罗玛诺斯}等。\par

我们的儿子\Person{德奥多西}，即由已故贵族\Person{利贝里乌斯}在\Place{坎帕尼亚}所建隐修院的院长，已知会我们：已故的显贵妇人\Person{鲁斯提卡}约二十一年前在订立的遗嘱中，首先指定其丈夫\Person{斐理斯}为继承人，并委托他建立一所\Place{西西里}的隐修院；但附有条件——若他未能在一年内付清所有遗赠给其被释放者的款项，或未按她的意愿建立上述隐修院，则圣罗马教会对她所拥有的\Place{库迈}庄园份额拥有无可争议的求偿权，并应协助支付上述遗赠及建造该隐修院。因此，鉴于据称该遗产迄今尚未完全移交给同一隐修院，且部分产业仍被其继承人扣留，兹令阁下详查并审理此案。首先，若在遗嘱条件下有任何可据以向教会提出申诉的继承权，望您调查清楚，并按事务所需为穷人的利益行事；然后立即关注该院落的正当建立及被遗赠财产的追回，以使立遗嘱者的虔诚意愿在两件事上均得实现，并使不义扣留财产者因正当损失而认识到其不当扣留的罪责。因此，我们恳切期望您以最大的热忱调查此案，并赖主的帮助将其了结，以使立遗嘱者的虔诚奉献终能生效。但我们也期望您，在正义允许的范围内，以各种方式扶助该隐修院，以免那些本应提供援助的世俗人士，如其所声称的那样，有权以建立者的名义造成损害。\par

\SourceRecord{Translated by James Barmby.；Nicene and Post-Nicene Fathers, Second Series；Vol. 13.；Edited by Philip Schaff and Henry Wace.；Buffalo, NY: Christian Literature Publishing Co.,；1898.；<http://www.newadvent.org/fathers/360209024.htm>.}

% structure: p0001 source=h1 target=section reason=page-title

\section{第九卷，第二十六封信}

\Emph{致\Person{罗玛诺}，\OriginalTerm{护卫官}{Defensorem}。}\par

\Person{格列高利}致\Person{罗玛诺}等。\par

虽然法律有理由不允许将归属教会之物转让，但有时在规则的严格上应稍加缓和，但凡顾及仁慈诱使如此，尤其当数量巨大，施予者未受重负，而受者的贫困大为减轻之时。因此，由于本信的递交者斯德法尼亚带着她的小儿子卡里克塞努斯（她声称她是与已故丈夫伯多禄所生，并说自己一直处于极度贫困之中）来到此地，以恳求和眼泪要求我们促使将位于卡塔纳城的一所房屋的产权归还给同一卡里克塞努斯——该房屋是已故的婆婆、卡里克塞努斯的祖母阿摩尼亚以赠与的名义献给我们教会的；她断言所述阿摩尼亚无权转让它，且该房屋完全属于前述的儿子卡里克塞努斯；此断言为我们至爱的儿子、执事西彼廉（他熟悉此案）所反驳，他说上述妇女的控诉并无正义依据，她不能合理地以她儿子的名义要求或寻求收回那所房屋；但是，为了避免我们似乎使上述妇人的眼泪落空，并选择严苛之路而非拥抱怜悯之请，我们以此谕令命你将所述房屋归还给上述卡里克塞努斯，连同阿摩尼亚关于同一房屋的赠与契据——该契据据悉在西西里——因为，如我们所说，在存疑案件中，宁可不执行严格，而倾向于仁慈，尤其当因些许让步而使教会不受重负，并仁慈地救助一个贫困孤儿之时。\par

颁于十一月，\OriginalTerm{小纪}{Indiction}二期。\par

\SourceRecord{Translated by James Barmby.；Nicene and Post-Nicene Fathers, Second Series；Vol. 13.；Edited by Philip Schaff and Henry Wace.；Buffalo, NY: Christian Literature Publishing Co.,；1898.；<http://www.newadvent.org/fathers/360209026.htm>.}

% structure: p0001 source=h1 target=section reason=page-title

\section{第九卷 第二十七封信}

\Emph{致\Person{罗马努斯}，\OriginalTerm{辩护者}{Defensorem}。}\par

\Person{额我略}致\Person{罗马努斯}等。\par

我们听说，某些人全然缺乏辨识，渴望我们卷入他们的风险，并希望得到教会人士的如此庇护，以至于教会人士自身也被他们的罪责所束缚。因此，我以本次训令劝诫你，并通过你劝诫我们的兄弟和同为主教的，主\Person{若望}，以及其他有关人士：关于对人们的教会庇护（无论你从我这里收到了信件，还是没有任何信件寄给你），你应当如此适度地施予，以致如果有人卷入公共贪污，他们不可显得被我们不公正地辩护，免得我们因冒然进行不谨慎的辩护，将恶行者的恶名转移到我们自己身上；而是要尽可能按照教会所当行的，通过劝诫和运用代祷的话语，救助你能救助的人；这样你既能给予他们援助，又不玷污圣教会的声誉。\par

\SourceRecord{Translated by James Barmby.；Nicene and Post-Nicene Fathers, Second Series；Vol. 13.；Edited by Philip Schaff and Henry Wace.；Buffalo, NY: Christian Literature Publishing Co.,；1898.；<http://www.newadvent.org/fathers/360209027.htm>.}

% structure: p0001 source=h1 target=section reason=page-title

\section{第九卷 第三十三封书信}

\Emph{致安德肋}。\par

\Person{额我略}致\Person{安德肋}。\par

听说阁下曾为忧愁和疾病所苦，我极表慰问。但随即得知疾病已完全离开你，我立刻转悲为喜，并衷心感谢全能的天主，因为他击打是为了医治，磨难是为了引导人获得真正的喜乐。因为经上记载：\BibleQuote{主所爱的，他必惩戒；他所收纳的每个儿子，他必鞭打。}\BibleRef{希 12:6}因此真理亲自说：\BibleQuote{我的父是园丁；凡在我身上不结实的枝条，他必剪去；凡结实的枝条，他必清理，使它结出更多的果实。}\BibleRef{若 15:1}因为不结果实的枝条被剪去，因为罪人会被彻底连根拔起。但结果实的枝条被称为被清理，因为它被纪律修剪，以便被引向更丰盛的恩宠。因为谷粒的穗子，被打谷器具击打，脱去了芒刺和糠秕。橄榄被压榨器挤压，流出了肥美的油。成串的葡萄被脚踩踏，化成了酒。因此，好人啊，喜乐吧，因为在这对你的鞭打和提升中，你看到你被永恒的审判者所爱。\par

此外，我请求以我的名义问候我的女儿格洛里奥萨，你的妻子。愿全能的天主保护你们在天主的眷顾之下，现在以丰富的恩赐安慰你们，将来以报应的酬报安慰你们。\par

\SourceRecord{Translated by James Barmby.；Nicene and Post-Nicene Fathers, Second Series；Vol. 13.；Edited by Philip Schaff and Henry Wace.；Buffalo, NY: Christian Literature Publishing Co.,；1898.；<http://www.newadvent.org/fathers/360209033.htm>.}

% structure: p0001 source=h1 target=section reason=page-title

\section{第九卷 第三十六封信}

\Emph{致\Person{福图纳图斯}，\Place{那不勒斯}主教。}\par

\Person{额我略}致\Person{福图纳图斯}等。\par

得知您的兄弟为了基督徒奴隶——犹太人在\Place{高卢}领土上购买他们——而被怎样的热忱所燃烧，我们通知您，您的关心如此令我们喜悦，以至这也是我们自己的审慎判断：他们应被禁止从事这种交易。但我们从希伯来人\Person{巴西略}——他与其他犹太人一起来到这里——发现，这种购买是由共和国的多位法官命令给他们的，并且基督徒与外教人都被这样获取。因此，必需通过谨慎的安排来调整此事，以便既不让下这些命令的人受阻，也不让那些声称被迫服从的人不公平地承担任何费用。因此，让您的兄弟以警惕的关心确保此事被遵守和保持；这样，当他们从上述省返回时，可能被他们带来的基督徒奴隶，要么交给那些下命令的人，要么无论如何在四十天内卖给基督徒买家。在这些天数结束后，不得有任何一人以任何方式留在犹太人手中。但是，如果这些奴隶中有人碰巧患上疾病，以至于无法在指定天数内出售，则应当注意，当他们恢复健康后，务必按照上述方式处理。因为任何人都不应当因一项无可指责的交易而遭受损失。但是，由于每当制定新规时，通常如此规定未来，以免在巨额费用上谴责过去，如果任何奴隶从前一年的购买中仍留在他们手中，或者最近被你们从他们那里带走，让他们有机会在你们那里处置他们。这样，他们便不可能因禁令前无知所做的事而遭受损失，正如他们在被禁止之后理应承受的那样。\par

此外，我们听说，上述\Person{巴西略}希望以赠礼的名义将某些奴隶让与他的儿子们——他们因天主的仁慈已成为基督徒。意图是，借此提供的机会，他们可以几乎名义上以他为家主服侍他；并且，如果此后有人碰巧以为他们可以逃往教会避难以成为基督徒，他们不可被要求恢复自由，而应归回先前被给予的那些人的统治之下。此事上，您的兄弟宜保持适当的警惕。如果他想将任何奴隶给予儿子，为除去一切欺诈的机会，务必使他们成为基督徒，并且不让他们留在他家中；但是，当情况可能需要他们为他服务时，应命令他们向他提供从儿子们那里，且为天主的缘故，理应提供给他的任何东西。\par

\SourceRecord{Translated by James Barmby.；Nicene and Post-Nicene Fathers, Second Series；Vol. 13.；Edited by Philip Schaff and Henry Wace.；Buffalo, NY: Christian Literature Publishing Co.,；1898.；<http://www.newadvent.org/fathers/360209036.htm>.}

% structure: p0001 source=h1 target=section reason=page-title

\section{第九卷 第四十一封信}

\Emph{致尤利安努斯书记官}。\par

\Person{额我略}致\Person{尤利安努斯}等。\par

如果在世俗官职中，我们祖先传下来的秩序和\OriginalTerm{纪律}{discipline}被遵守，那么谁又能忍受看到教会的秩序被扰乱，听到这些事情时置之不理，并以不当的宽容推迟纠正呢？你确实做得很好，热爱\OriginalTerm{爱德}{charity}并劝人和谐。但是，由于我们因职位的考量并为天主的缘故，不得不对已知之事进行彻底调查，当\Person{马克西穆斯}到来时，我们将留意向他要求他所说的关于他的事件的严格交代。我们信赖我们的\OriginalTerm{造物主}{Creator}的护佑，我们不会因任何人的好恶而偏离维护\OriginalTerm{教规}{canons}和正直的公义之道，而是乐意遵行合乎理性之事。因为如果我们（但愿天主不许）忽视了教会的关切和热忱，懈怠就会摧毁纪律，而且必定会对信友的灵魂造成伤害，因为他们看到牧者们为他们树立了这样的榜样。至于你在信中说，宫廷的好意和人民的爱戴并未离开他，这一情况并不会使我们从对正义的热忱中退缩，也不会因我们的罪过而导致我们查究真理的决心落空。那么，尊贵的儿子，每个人都应当努力为自己赢得天主的爱。因为没有天主的眷顾，人世间的爱日后能为我们做什么呢？即使在我们中间，它反而更有害于我们。\par

\SourceRecord{Translated by James Barmby.；Nicene and Post-Nicene Fathers, Second Series；Vol. 13.；Edited by Philip Schaff and Henry Wace.；Buffalo, NY: Christian Literature Publishing Co.,；1898.；<http://www.newadvent.org/fathers/360209041.htm>.}

% structure: p0001 source=h1 target=section reason=page-title

\section{第九卷，第四十二封书信}

\Emph{致伦巴第人国王\Person{阿吉卢尔夫}}\par

\Person{格列高利}致\Person{阿吉卢尔夫}等\par

我们感谢陛下，因为您听从了我们的请求，缔结了这样一项对双方都有利的和平，正如我们确信您会做的那样。为此，我们高度赞扬您的智慧与仁善，因为您选择和平，表明您爱天主，祂是和平的缔造者。因为，如果不幸和平未能达成，除了在罪恶与危险中，双方可怜农民的鲜血流淌之外，还能有什么结果呢？他们的劳动本可使双方受益。但是，为了让我们感受到这项和平对我们有利，既然它已经由您缔结，我们以慈父般的爱德向您致意，恳请您，每当有机会时，通过信件命令您各处的地方公爵，尤其是那些驻扎在这些地区的公爵，要他们如所承诺的那样不可侵犯地保持这项和平，且不要为自己寻找任何可能引起争执或恶感的借口，好让我们能更加为您的善意感恩。我们以应有的深情接待了这些信件的携带者，他们确实是您的人民，因为以爱德接待和送别这些明智的人，他们宣告了因天主的恩宠和平已经缔结，这是正确的。\par

\SourceRecord{Translated by James Barmby.；Nicene and Post-Nicene Fathers, Second Series；Vol. 13.；Edited by Philip Schaff and Henry Wace.；Buffalo, NY: Christian Literature Publishing Co.,；1898.；<http://www.newadvent.org/fathers/360209042.htm>.}

% structure: p0001 source=h1 target=section reason=page-title

\section{第九卷，第四十三封信}

\Emph{致\Place{伦巴底}王后\Person{特奥德琳达}}\par

\Person{额我略}致\Person{特奥德琳达}等\par

从我们的儿子、院长\Person{普罗布斯}的报告中，我们得知您的贤明如何一如既往地热心而仁慈地为缔造和平而努力。确实，以您的基督徒精神，我们不会期望别的，只会期望您以各种方式展现您对和平事业的勤勉和良善。为此，我们感谢全能的天主，祂以慈爱治理您的心，使您既拥有正信，又能常做祂眼中喜悦的事。最贤明的女儿，您可以确信，为拯救双方免于流那么多血，您已获得了不小的赏报。因此，我们为您的善意而感谢，并恳求天主的仁慈在现世和来世赐您身心俱获善报。\par

此外，我们以慈父之情问候您，劝勉您如此与您最贤明的配偶相处，使他不要拒绝基督徒共和国的联盟。因为，如我相信您自己所知，他倾向寻求其友谊在多方面是有益的。那么，您当一如既往，常为促进双方的和睦与和解而努力，并在任何获得赏报的机会出现时劳作，使您的善行在天主眼前更为可嘉。\par

\SourceRecord{Translated by James Barmby.；Nicene and Post-Nicene Fathers, Second Series；Vol. 13.；Edited by Philip Schaff and Henry Wace.；Buffalo, NY: Christian Literature Publishing Co.,；1898.；<http://www.newadvent.org/fathers/360209043.htm>.}

% structure: p0001 source=h1 target=section reason=page-title

\section{第九卷，书信四十九}

\WorkTitle{致安提约基雅主教阿纳斯塔希乌斯}。\par

额我略致阿纳斯塔希乌斯等。\par

我收到了贤弟的信函，信中正确地持守了信仰的宣认；我向全能的天主献上极大的感谢，因为祂在更换羊群的牧者时，即使在更换之后，仍守护祂一次交付圣教父们的信仰。杰出的宣讲者说：\Quote{除已奠立的根基，即耶稣基督外，任何人不能再奠立别的根基。}（\BibleRef{格前 3:2}）因此，凡怀着对天主及对近人的爱，坚定持守在基督内的信仰者，便为自己奠定了同一的耶稣基督、天主子与人之子作为根基。可见，凡以基督为根基之处，也应期待善功的建筑随之而来。真理本身也亲自说：\Quote{那不从门进入羊栈，而从别处爬进去的，便是贼，是强盗；但从门进去的，才是羊的牧人。}（\BibleRef{若 10:1}）稍后祂又补充说：\Quote{我就是门。}因此，那从门进入羊栈的，便是经过基督进入的。而经过基督进入的，是那对人性的创造主与救赎主怀有并宣讲真实道理，且持守其所宣讲的；他以承担重担的职务、而非贪图短暂尊荣的虚荣心来承受最高的统治地位。他也明智地看守所托付的羊栈，免得乖戾的人以悖逆之言撕裂天主的羊，或恶毒之灵以诱人的恶乐来蹂躏它们。\par

我们确实记得，可敬的雅各伯曾为他的妻子们长久服事，他说：\Quote{这二十年我同你在一起；你的母绵羊和母山羊从未流产；你羊群中的公羊我从未吃过；被野兽撕裂的，我没有给你看过。我赔偿了一切损失；凡被偷去的，你都向我索还。白天黑夜我受尽干渴和霜冻，睡眠也逃离了我的眼睛。}（\BibleRef{创 31:38}）如果喂养拉班羊群的人尚且如此劳苦警醒，那么喂养天主羊群的人应当何等专注劳苦与警醒呢？但在这一切事上，愿那为我们而成为人、屈尊就卑成为祂所创造的那一位教导我们。愿祂将祂自己爱的神倾注于我的软弱和你的爱德之中，并在一切谨慎与警醒的谨慎中开启我们内心的眼睛。\par

然而，对于正信之人晋升神圣品秩一事，应不断向同一全能的天主谢恩，并常为我们最虔诚、最基督徒的皇帝陛下及其最宁静的配偶、以及他们最温和的子孙祈祷；在他们的时代，异端者的口被缄默；因为尽管他们的心因悖逆思想的疯狂而沸腾，但在公教皇帝的时代，他们不敢说出所想的恶事。\par

再者，贤弟提及对神圣大公会议的信守，声明你遵守了第一次神圣的厄弗所大公会议。但是，由于从皇宫送来的一份异端文件中所作的陈述，我发现按照该文件，某些公教立场与异端立场一同被谴责，因为有些人以为那是在同一城市中曾由异端者召集的第一次厄弗所会议，因此，你务必向亚历山大里亚和安提约基雅的教会索取该会议的法案，查明事实真相。或者，你若愿意，我们将从这里寄给你我们档案馆中自古保存的副本。因为那假托第一次厄弗所会议召开的会议声称，某些提交给它的立场得到了认可，而这些正是塞勒斯提乌和佩拉纠的明确主张。既然塞勒斯提乌和佩拉纠已在那个会议上被谴责，这些立场如何能得到认可呢？它们的作者已被谴责了啊。\par

此外，听说在东方各教会中，无人因行贿授职而能达到神圣品秩，贤弟若发现确有其事，当首先向全能的天主献上初果，即在你辖下的各教会中抑制西满异端的谬误。因为，其他暂且不论，那些不是凭功德、而是凭贿赂被提升到神圣品秩的人，能成为怎样的人呢？愿全能的天主以天上的恩宠护守你的爱德，并恩赐你从所托付的人身上获得倍增的果实与满溢的收获，带往永恒的喜乐。\par

\SourceRecord{Translated by James Barmby.；Nicene and Post-Nicene Fathers, Second Series；Vol. 13.；Edited by Philip Schaff and Henry Wace.；Buffalo, NY: Christian Literature Publishing Co.,；1898.；<http://www.newadvent.org/fathers/360209049.htm>.}

% structure: p0001 source=h1 target=section reason=page-title

\section{第九卷，书信五十五}

致\Person{凡蒂诺}，\Place{帕诺尔穆斯}（\Place{巴勒莫}）的守护者（\OriginalTerm{守护者}{Defensorem}）。\par

\Person{额我略}致\Person{凡蒂诺}等。\par

不久前我们写信给我们的兄弟和同为主教的\Person{维克托}，即——因某些犹太人提交给我们的请愿书中抱怨说，位于\Place{帕诺尔穆斯}的会堂及其客房被他无理占据——他应远离他们的集会，直到查明此事是否公正地执行，以免无辜者似乎因他们自己的意愿而受到伤害。的确，考虑到他的司铎职务，我们很难相信我们上述兄弟曾做出任何不当之事。但是，由于我们从随后在场的我们公证人\Person{萨拉里乌斯}的报告中发现，占据那些会堂没有合理理由，而且它们是被轻率且鲁莽地祝圣的，因此我们命令您的经验，既然一旦祝圣就不能再归还给犹太人，您应负责确保我们上述兄弟和同为主教支付我们的儿子，光荣的贵族\Person{韦南提乌斯}和院长\Person{乌尔比库斯}所估价的会堂本身及其下方或紧邻墙壁的客房，以及相邻花园的价格；这样他所占据的就可归属教会，而他们（犹太人）绝不会受压迫或遭受任何不公。此外，被拿走的书籍或装饰物也应同样查找。如果任何物品被明显带走，我们同样希望它们毫无歧义地归还。因为，正如我们已写过的，他们在自己的会堂内不得有任何超出法律规定的许可，同样，也不应使他们遭受违背正义与公平的损失或任何费用。\par

\SourceRecord{Translated by James Barmby.；Nicene and Post-Nicene Fathers, Second Series；Vol. 13.；Edited by Philip Schaff and Henry Wace.；Buffalo, NY: Christian Literature Publishing Co.,；1898.；<http://www.newadvent.org/fathers/360209055.htm>.}

% structure: p0001 source=h1 target=section reason=page-title

\section{第九卷，第五十八封}

\Quote{致\OriginalTerm{学者}{Scholasticus}\Person{玛尔定}}。\par

\Person{额我略}致\Person{玛尔定}等。\par

鉴于民事事务中产生的问题，如阁下所知，需要非常详尽的调查，那么主教们的案件应以何等谨慎和警觉加以调查，请阁下智慧思量。然而，在您随此函送交我们的信中，关于您受我们弟兄和同道主教\Person{克雷孟提乌斯}派遣前来与我们商议的那些问题，您只作了肤浅的陈述，而对其根源只字未提。但是，倘若这些问题之起源与内在性质为我们所知晓，我们本会知道应如何决定，并以清晰而适当的答复安抚我们前述弟兄的心。然而，令我们完全不悦的是，您向我们表示，有些主教未经其首席主教之文书便去了宫廷，并且他们举行了非法的集会。但既然如上所述，我们对这些问题的起源与性质全然不知，我们无法明确宣判任何事，唯恐因对不完全了解之事作出判决而显得极为可责。因此，非常必要的是，为使我们获得充分信息，阁下本应在您逗留\Place{西西里}期间前来此地答复我们的问题。尽管如此，既然您已见过我们的弟兄和同道主教\Person{若望}，我们相信您在他身上也见到了我们。因此，由于他本人也已费力就同样的问题写信给我们，我们已回信给他，告知我们认为正确之事。既然他是一位经验丰富且判断谨慎的司铎，若您愿意就他所受命处理的问题与他商议，我们确信您会从他那里找到既有利又合理之处。\par

\SourceRecord{Translated by James Barmby.；Nicene and Post-Nicene Fathers, Second Series；Vol. 13.；Edited by Philip Schaff and Henry Wace.；Buffalo, NY: Christian Literature Publishing Co.,；1898.；<http://www.newadvent.org/fathers/360209058.htm>.}

% structure: p0001 source=h1 target=section reason=page-title

\section{第九卷，书信五九}

\Emph{致 \Place{西拉库萨}主教 \Person{若望}}\par

\Person{额我略}致 \Person{若望}等\par

我已收到你的手足之书，其中你告知我，最善辞令的 \Person{马丁}从\Place{阿非利加省}前来，私下向你传达了一些事。确实，你的手足每当找到机会，总是不停地显示你对蒙福的宗徒 \Person{伯多禄}的爱。为此，我们感谢全能的天主，因为你在哪里，我们就不被视为缺席。然而，你的圣座尚未完全明了此事。原来，\Place{比则切纳}省的首席主教因某项指控被控告，最虔诚的皇帝希望我们依教规审理他。但随后，在收到十磅黄金后，\OriginalTerm{军事长官}{magister militum} \Person{德奥多禄}反对此事。然而，最虔诚的皇帝劝告我们委派某人，并做一切合乎教规的事。但鉴于人的矛盾，我们不愿判决此案。现在，这位首席主教又就自己的意图说了一些话。他向我们说这些话是出于真诚，还是因他正受同僚主教攻击，实在难以确定。至于他说自己服从宗座，若主教们被发现有过错，我不知道哪位主教不服从宗座。但若无需特别处理，按谦逊原则，众人皆是平等的。不过，你可视情况与前述最善辞令的 \Person{马丁}谈话。因为该做什么由你考虑；我们已就此事简短回复你，因为我们不应轻易相信素不相识的人。然而，你若亲眼见他，认为应对他说些更明确的话，我们便将此事交托于你的爱德，确信你在全能天主的恩宠中怀有爱。你所视为如何的，无疑被视为是由我们所做。\par

\SourceRecord{Translated by James Barmby.；Nicene and Post-Nicene Fathers, Second Series；Vol. 13.；Edited by Philip Schaff and Henry Wace.；Buffalo, NY: Christian Literature Publishing Co.,；1898.；<http://www.newadvent.org/fathers/360209059.htm>.}

% structure: p0001 source=h1 target=section reason=page-title

\section{卷九，第六十封信}

\Emph{致\Person{罗马努斯}及其他教会产业\OriginalTerm{护卫者}{defensores}。}\par

\Person{额我略}致护卫者\Person{罗马努斯}、护卫者\Person{范提努斯}、副执事\Person{撒比努斯}、护卫者\Person{塞尔吉乌斯}、护卫者\Person{波尼法爵}（\OriginalTerm{同职}{a paribus}）及六位\OriginalTerm{保护人}{patroni}。\par

既然，正如谨慎的远见知道如何阻挡过失的道路，并避免有害之事，同样，疏忽为放纵敞开道路，并常招致本应防范的事；我们应当非常用心地注意，既顾及我们弟兄与司铎的名誉，也顾及他们的安全。如今我们有耳闻，某些主教在帮助的借口下，与妇女同住一屋。因此，为避免由此给嘲笑者提供正当的诽谤口实，或人类古老的仇敌利用这轻易的欺骗机会，我们以此命令之内容，命令你努力表现自己为勤勉而关切者。若在托付给你的产业范围内，有任何主教与妇女同住，你应完全制止此事，并今后绝不允许任何妇女与他们同住，除非神圣法典的审查所允许的，即母亲、姑母、姐妹及此类无可疑恶意者。然而，他们若甚至连与这些人也不同住，则更佳。因我们读到，真福\Person{奥斯定}甚至拒绝与自己的姐妹同住，说：\Quote{与我姐妹同住的人，不是我的姐妹。}\par

有学识之人的谨慎应为我们的重大教训。因为一个较弱的人不畏惧强者所害怕之事，实为不慎的冒失之标志。因他明智地克服了非法之事，他已学会连允许他的事也不使用；的确，在此事上我们不束缚任何人违反其意愿，而是如医生通常所做，为了健康而规定谨慎，即使暂时令人痛苦。因此，我们不施加必要的义务；但若有人选择效法有学识的圣人，我们留待其自愿。那么，让你的\OriginalTerm{经验}{Your Experience}以热诚和关切，致力于遵守我们所命令禁止之事。因为今后若被发现另有情况，要知道你将与我们一同面临不小的风险。此外，你应负责劝诫这些同为弟兄的主教，叫他们训导属下，即那些处于神圣品级者，以他们的榜样，在一切事上遵守他们所遵守的；只添加这一点，即这些人，按照法典权威所规定的，不得离开他们应贞洁管理的妻子。\Signature{于三月，\OriginalTerm{小纪}{Indiction}第二年颁。}\par

\SourceRecord{Translated by James Barmby.；Nicene and Post-Nicene Fathers, Second Series；Vol. 13.；Edited by Philip Schaff and Henry Wace.；Buffalo, NY: Christian Literature Publishing Co.,；1898.；<http://www.newadvent.org/fathers/360209060.htm>.}

% structure: p0001 source=h1 target=section reason=page-title

\section{\WorkTitle{第九卷，第六十一封信}}

以下是哥特王雷卡雷德致真福罗马主教额我略的书信之始。\par

雷卡雷德致可敬的主、至福的教宗、主教额我略。\par

当主以祂的慈悲使我们脱离不虔的亚略异端，圣而公教会将我们聚集于她的怀中，在信德的道路上使我们改善时，那时我们心中渴望欣喜地寻求一位如此可敬的人——您，超越所有其他主教的权威者，愿他竭尽全力为我们人类荐举这样一件堪受天主之事。然而，由于我们忙于诸多治理之事，为各种事务缠身，三年过去，心中的渴望仍未得满足。此后，我们拣选了几位隐修院院长，派他们前往您处，好让他们到您面前，献上我们托圣伯多禄转送的礼物，并更清晰地告知我们您神圣尊驾的健康状况。但当他们匆匆赶路，几乎望见意大利海岸时，不幸却在马赛附近触礁，仅得保全性命。如今我们恳求您荣耀所派至马拉加城（\OriginalTerm{马拉加城}{civitatem Malicitanam}）的一位司铎前来与我们相见。但他因身体病弱，始终未能抵达我王国境内。然而，我们确知他是由您圣座所遣，便送上了一个外饰宝石的金杯，献给您的圣座（我相信您必将俯允），以为堪当献给那位在尊威中首放光芒的宗徒。我也恳求您尊驾，若遇时机，请以您神圣的金字书信寻访我们。因我真心热爱您，我相信，在主启迪下，此事对您胸中的丰硕之心并非隐藏。有时，陆地或海洋分隔的人，基督的恩宠却仿佛有形地将其相连。因为对于那些未能亲见您者，声誉已显明您的善行。\par

此外，我以全心的敬仰，在基督内将希斯帕利斯教会司铎良德托付于您的圣座，因为藉着他，您的慈爱已清楚地显示给我们；当我们与这位主教谈论您的生平，我们自认在善行上不及您。至可敬、至圣者，我欣然听闻您的健康；并恳请您的基督之智，时常在祈祷中向我们的共同主荐举我们及我们的子民——他们在天主之后归我治理，并在您的时代被基督所获——好使向天主的真诚爱德，能使那被世界广袤分离者安居于福祉之中。\par

\SourceRecord{Translated by James Barmby.；Nicene and Post-Nicene Fathers, Second Series；Vol. 13.；Edited by Philip Schaff and Henry Wace.；Buffalo, NY: Christian Literature Publishing Co.,；1898.；<http://www.newadvent.org/fathers/360209061.htm>.}

% structure: p0001 source=h1 target=section reason=page-title

\section{第九卷，第六十二封信}

\Emph{\Person{罗玛努斯}，\OriginalTerm{监护人}{Defensorem}}。\par

\Person{额我略}致\Person{罗玛努斯}等。\par

我们得知，在西西里的\Emph{\OriginalTerm{剃发者}{tonsuratores}}竟敢冒昧地擅自采用\Emph{\OriginalTerm{监护人}{defensores}}的名号，他们不仅对教会利益毫无用处，反而借此机会犯下许多不法行为。因此，我们凭此谕令责成你的经验详查此事。如果你发现任何人，除那些持有信件授权从事此类事务者外，今后僭用此称号，应严加纠正，予以制止。然而，若你发现有人在教会事务中证明自己积极忠实，则必须向我们呈上一份详尽特别的报告，以便我们判断他们是否配得一份信函。\par

此外，我们希望你彻底审查福图纳图斯的账目；当他清偿了所有对他不利的债务后，不再允许他与教产或我们教会的任何事务有关，因为，据我们所闻，他的行为举止已使他今后不应再与我们的人员有任何往来。\par

此外，有人向我们报告，一位自称\Emph{\OriginalTerm{监护人}{defensor}}的玛尔提亚努斯，拒绝服从我们的兄弟和同为主教的若望——我们曾将教产的管理委托给他。因此请查问；若属实，应将他流放，以使他的不顺从——他违背了那一位，从后者（若望）的教会中为自己窃取了一个虚假的荣誉头衔，而后者正在促进同一教会的利益——不致不受惩罚。但是，若还有其他人违背我们上述兄弟的命令，你务必以严厉惩罚对待他们。\par

\SourceRecord{Translated by James Barmby.；Nicene and Post-Nicene Fathers, Second Series；Vol. 13.；Edited by Philip Schaff and Henry Wace.；Buffalo, NY: Christian Literature Publishing Co.,；1898.；<http://www.newadvent.org/fathers/360209062.htm>.}

% structure: p0001 source=h1 target=section reason=page-title

\section{\WorkTitle{第九卷，书信六五}}

\Emph{致卡辣里（卡利亚里）主教雅努雅里乌斯。}\par

额我略致撒丁尼亚主教雅努雅里乌斯。\par

我等闻知，尔之某些圣职人员，受高傲之灵所煽（此事甚为严重），不服从尔兄弟之命，反而忙于他人之事务与劳作，弃其本堂需用之职责。是故，我等大为惊异，为何尔不持守纪律之规，任其放纵游荡，并以严格之缰绳约束之，使其尽职守。又闻，某些顽抗之圣职者，为求支援以抗尔，竟投靠我等之\OriginalTerm{监护人}{defensoris}维塔利斯。因此，我等已致函于彼，告其今后不得再无理支持尔之任何圣职人员以抗尔；但若出现过失，并非严重而可宽恕者，则应以调停者而非支持者之身份接近尔。尔当留意，勿使再有此类属民轻视尔之报告传至我等耳中。\par

我等又闻，有一寡妇将其财产遗留与圣犹利安隐修院，而此财产竟被尔之一圣职人员掠夺，此人在寡妇生前曾掌管其事务，而今逃避归还。故我等劝尔，若所言属实，当以严格程序迫其迅速如数归还隐修院之财产，并责令其即使丧失名誉亦须交出，盖其本不应在保持名誉纯洁时妄取之。然此等可耻之事，竟需我等提醒尔兄弟以纪律之力约束尔之圣职人员，我信尔心中亦自知之。\par

此外，对于拜偶像者、占卜者及算命者，我等恳切劝尔兄弟以牧者之警醒留意，在民众中公开斥责行此等事之人，以劝导之言使其远离如此亵圣之污染、天主审判之诱惑及现世之危险。然而，若尔发现其不愿悔改自新，我等望尔以热忱之火热拿获之；若其为奴隶，则以鞭打及折磨惩治之，使其改过；若其为自由人，则以适当而严格之监禁引导其补赎。如此，凡轻忽救赎之言、拒绝脱离死亡危险者，至少可因身体之苦痛而回归所渴望之心灵健全。我等又闻，尔将教会产业委托于某些平信徒，彼等被发觉侵吞尔之农民财物后逃亡，既拒绝归还其不正当地视为己有之财产，因不受尔管制，又藐视向尔报告其行为。若果如此，尔当严格调查此事，并彻底审理彼等与尔教会农民之间之案件。凡查出之欺诈，应依法处以罚金，迫其赔偿。但尔兄弟此后务必注意，勿将教会产业托付于不受尔管辖之世俗人，而应托付于尔属下之合格圣职人员。若在彼等中发现任何不法行为，尔作为其上司，得以纠正非法之事，因其身份之义务使其受尔传唤而非推诿。\par

\SourceRecord{Translated by James Barmby.；Nicene and Post-Nicene Fathers, Second Series；Vol. 13.；Edited by Philip Schaff and Henry Wace.；Buffalo, NY: Christian Literature Publishing Co.,；1898.；<http://www.newadvent.org/fathers/360209065.htm>.}

% structure: p0001 source=h1 target=section reason=page-title

\section{第九卷，书信第六十七}

\Emph{致\Person{君士坦提乌斯}，\Place{米兰}主教。}\par

额我略致\Person{君士坦提乌斯}等。\par

\Person{玛克西穆斯}，\Place{萨洛纳}教会的背信者，在未能通过世上更大的权势获得任何利益之后，便转而投靠较小的权势；他通过过多的祈祷和对其善行的证明，力图说服我们。既然如此，我认为若一个自称十分敬畏我的人，却完全无法在我这里找到某种程度的宽仁，那将是不人道的。因此，我决定由我们最可敬的兄弟和同为主教的\Person{马里尼亚努斯}在\Place{拉文纳}城审理他的案件。然而，倘若因故他的身份受到怀疑，我们希望你的弟兄之爱（若对你不算过于劳苦）也费心前往同一座城，并与我们上述的兄弟一同参与该审判。无论你们的圣善认为怎样妥当，要知道那在我眼中也必妥当；你们的审判我视为自己的审判；你们二位认为应当宽免的，确信我也宽免；但须审慎留意，以免我们显得有罪地懈怠，或严厉至损害圣教会。我们已将此事委托给掌玺官\Person{卡斯托里乌斯}执行，以便他将所行的一切详尽报告我们。\par

\SourceRecord{Translated by James Barmby.；Nicene and Post-Nicene Fathers, Second Series；Vol. 13.；Edited by Philip Schaff and Henry Wace.；Buffalo, NY: Christian Literature Publishing Co.,；1898.；<http://www.newadvent.org/fathers/360209067.htm>.}

% structure: p0001 source=h1 target=section reason=page-title

\section{第九卷，第六十八封信}

\Emph{致\Place{得撒洛尼}的\Person{尤西比乌斯}。}\par

\Person{额我略}致\Place{得撒洛尼}的\Person{尤西比乌斯}、\Place{都拉齐乌姆}的\Person{乌尔比库斯}、\Place{尼科波利斯}的\Person{安德肋}、\Place{格林多}的\Person{若望}、\Place{第一犹斯提尼亚纳}的\Person{若望}、\Place{克里特}的\Person{若望}、\Place{拉里萨}和\Place{斯库台}的\Person{若望}，以及许多其他主教。\par

我们因所肩负的治理之责，必须警醒地扩展我们职务的关怀，并以劝诫之言教导我们弟兄的心智，以免任何错误的推测足以欺骗无知者，也不让任何掩饰为知情者开脱。务请各位兄弟知晓：前君士坦丁堡主教\Person{若望}，反对天主，反对教会的和平，轻蔑并伤害所有司铎，逾越了谦逊和他自身尺度的界限，并在主教会议中非法篡夺了那骄傲而有害的\OriginalTerm{大公}{œcumenical}，即普世称号。当我们的前任、荣福记忆的\Person{贝拉基}得知此事后，他以完全有效的谴责废除了那次主教会议的全部记录，唯除其中为荣福记忆的安提约基雅主教\Person{额我略}案件所做之事；他以最严厉的斥责责备他，警告他戒绝那个新的、鲁莽的迷信名称；甚至还禁止他的执事与他同行列队，除非他改正如此大的恶行。而我们，完全遵循他正直的热忱，在天主的保护下，坚固地遵守他的命令，因为那永恒审判者的法庭正等待他为同一治理职分交账，所以他理应沿着前任的直路行走而不绊倒。在这件事上，为免我们显得忽略任何关乎教会和平的事，我们曾一再致函那位同样的至圣\Person{若望}，命令他放弃那骄傲的名称，并将他心中的高傲转向我们的师傅与主所教导我们的谦卑。我们发现他毫不理会后，出于对和谐的渴望，我们未曾停止向他的继任者、我们至真福的兄弟与同僚司铎\Person{基里亚库斯}发出同样的劝诫。但是，既然我们看到，如今世界末日临近，人类之敌已借其先驱出现，而通过这些骄傲的头衔，那些本应以善生和谦卑抵挡他的司铎，竟然成为了他的前驱。我劝勉并恳求，你们之中绝不可有一人接受此名，不可有一人赞同此名，不可有一人书写此名，不可有一人在它无论何处被书写时予以承认，也不可有一人加上其签署。反之，正如全能天主的仆人所宜，你们各人应远离这种毒害，不给狡猾的潜伏者任何地步，因为此事正在造成对全教会的伤害与分裂，并且，如我们所说，是对你们所有人的轻蔑。因为如果有一人，如他所自认的，是普世主教，那么你们就不是主教了。\par

此外，我们得知，你们的兄弟团体已被召集至\Place{君士坦丁堡}。虽然我们至虔诚的皇帝不允许在那里做任何非法之事，但是，恐怕不正之人借你们聚会之机，寻求机会以逢迎你们支持此迷信名称，或者想到举行关于另一问题的会议，意在以狡猾的计谋将此名引入，——虽然未经宗座的权威与同意，所通过的任何事都不会有效力——然而，在全能天主面前，我恳求并警告你们，不要让任何阿谀、贿赂或威吓取得你们任何一人的同意。反之，要顾及永恒的审判，健康而同心地表决反对不正当的企图；并凭牧者的恒心与权威，阻挡那闯进来的强盗和豺狼，不给那为撕裂教会而狂怒者任何地步；也不允许因任何花言巧语而就此事举行会议，那会议其实不会是合法的，也不可称为会议。我们也同时劝告你们，假若偶然没有提到这荒谬名称的事，但某个会议或许因另一事情而召集，你们在各方面都要谨慎、周全、警醒、留意，免得有任何违反任何地方或个人的、有偏见的、非法的或违背教规的谕令在此会中通过。但是，如果出现任何有益处理的问题，则所讨论的问题应取一种不致推翻任何古老规条的形式。因此，我们在天主及其诸圣面前再次劝告你们，要以最大的注意和全部的心意遵守所有这些事。因为如果有人，我们并不相信会这样，在任何部分忽略本信，愿他知道，他将被隔绝于真福\Person{伯多禄}、宗徒之长的和平之外。那么，你们的兄弟团体应如此行动，以便当牧者的牧者来审判时，你们不会因所接受的治理之职而被发现有罪。\par

\SourceRecord{Translated by James Barmby.；Nicene and Post-Nicene Fathers, Second Series；Vol. 13.；Edited by Philip Schaff and Henry Wace.；Buffalo, NY: Christian Literature Publishing Co.,；1898.；<http://www.newadvent.org/fathers/360209068.htm>.}

% structure: p0001 source=h1 target=section reason=page-title

\section{第九卷，第七十八封信}

\Emph{致\Place{亚历山大}宗主教\Person{尤洛吉乌斯}}\par

\Person{格列高利}致\Person{尤洛吉乌斯}等。\par

我从送信人手中收到了您至甘饴之圣座的信，信中对我谈及您的事将会迅速了结。然而，他一到达就了解到他从我教会所求的那产业是如何持有的，很快便对此感到满意。他与别人之间的事务也和平解决了。\par

然而，关于那件本应无论如何都要写信告诉我的事，您的圣座却只字未提，您也认为我在其中是迟缓的。实在说来，我担心它会引发分裂的丑闻，因此我不愿成为这种分裂的起因。因为我宁愿任何事情的发生都是通过别人。但将来，蒙天主恩赐，您将会证明：在我渴望讨天主喜悦的事上，我并不畏惧人。关于此事，我早已留心写信给您，甚至在您前往\Place{君士坦丁堡}时也是如此。\par

至于木材，我原本准备了一些更大尺寸的木材，正如您的祝福曾在信中所要求的；但这里派来的船太小，除非将木材砍短，否则无法装载。但我不愿将它们砍短，并保留了您的判断，看应如何处理。如果您不需要它们，我们将把它们用于这里的其他用途。此外，我恳求您的圣座为我恳切祈祷，因为我一直不断地受到痛风、野蛮人的刀剑和令人忧伤的挂虑的压迫。但是，如果您以祈祷的帮助赐予我，我相信您将大力助我抵抗一切逆境。\par

\SourceRecord{Translated by James Barmby.；Nicene and Post-Nicene Fathers, Second Series；Vol. 13.；Edited by Philip Schaff and Henry Wace.；Buffalo, NY: Christian Literature Publishing Co.,；1898.；<http://www.newadvent.org/fathers/360209078.htm>.}

% structure: p0001 source=h1 target=section reason=page-title

\section{第九卷 书信第七十九}

\Emph{致\Person{马里尼亚努斯}，\Place{拉文纳}主教}。\par

\Person{额我略}致\Person{马里尼亚努斯}等。\par

关于\Person{玛克西穆斯}一案应如何处理，您已从我们先前寄给您的信函中得知。然而，由于我们从本信呈递者、我方的档案员\Person{卡斯托里乌斯}的报告中得知，您的兄弟情谊在此事上的意愿，或更确切地说，请求为何，因此，若上述\Person{玛克西穆斯}在您和我方前述档案员面前，就买卖圣职异端宣誓自证清白，并关于其他指控，按我们所写，在\Person{圣亚坡理纳}遗体前，仅在受讯问时回答自己无罪，则我们将他的案件交由您的兄弟情谊裁决，即关于他胆敢在受绝罚期间举行弥撒圣祭一事，应处以何种补赎来洁净此过犯。因此，凡按天主之意认为妥当之事，请您无畏地裁定，且莫对我们有任何疑虑。因为凡由您就此案所定之事，我们皆感欣慰地接受并乐意允准。然而，我们劝诫您应谨慎，并如此调和您所准备施行之事，既若合宜，善待于他，又通过适当安排，如您所当行，保持教会严正精神的特质。我们已嘱咐上述送信者，他在我们这里时，应如何与您共同行事；您从他那里全面了解一切后，请在各方面如此行事，以致我们在您焦虑的关怀中感到我们与您同在。\par

\SourceRecord{Translated by James Barmby.；Nicene and Post-Nicene Fathers, Second Series；Vol. 13.；Edited by Philip Schaff and Henry Wace.；Buffalo, NY: Christian Literature Publishing Co.,；1898.；<http://www.newadvent.org/fathers/360209079.htm>.}

% structure: p0001 source=h1 target=section reason=page-title

\section{第九卷，第八十书}

\Emph{致公证人\Person{卡斯托里乌斯}}\par

\Signature{额我略致\Person{卡斯托里乌斯}等}\par

你越见自己受我们信任，并在需要时受托处理案件，就越当显示出你的勤勉与关切。因此，若\Place{萨洛纳}的\Person{马克西穆}起誓宣称他无罪于买卖圣职的异端，且在其他事上，仅在\Person{圣亚波利纳利}遗体前受询问时答称无辜，并已按我们指示为他的不服从而做了补赎，我们愿你，有经验的阁下，将我们写给他的信交给他，以安慰他。信中我们已表明已恢复了他我们的眷顾与共融。因为，正如对坚持抗命者我们应当严厉，对再次谦卑与忏悔者我们也不应拒绝给予宽恕之地。\par

此外，关于我们的弟兄、\Place{雅德拉}主教\Person{萨比尼亚努}、\Place{萨洛纳}总执事\Person{奥诺拉图}，以及其他曾求助于宗座者，必须非常恳切地与\Person{马克西穆}处理，使他能以合宜的爱德接待他们，心中绝不保留任何怨气，而是以纯正的善意和真诚的情谊与他们相处。\par

\SourceRecord{Translated by James Barmby.；Nicene and Post-Nicene Fathers, Second Series；Vol. 13.；Edited by Philip Schaff and Henry Wace.；Buffalo, NY: Christian Literature Publishing Co.,；1898.；<http://www.newadvent.org/fathers/360209080.htm>.}

% structure: p0001 source=h1 target=section reason=page-title

\section{第九卷，第八十一封书信}

\Emph{致\Person{玛克西穆}，\Place{撒罗纳}主教。}\par

\Person{额我略}致\Person{玛克西穆}等。\par

尽管你在起初的祝圣中已有不当之处，后来又因不顺从之过增加了严重的恶行，然而我们以合宜的节制运用宗座的权柄，从未如案情所要求的那样对你发怒。但由于受托付的责任感深深折磨我们，恐怕我们似乎忽略了所听闻的你的某些非法行为，因此你对我们所激起的不悦持续更久。而且，你若好好思考，就会看到你自己因迟迟不满足我们而证实了这些传闻，从而更加激怒了我们。但现在，你听从了有益的劝告，谦卑地顺服于服从之轭，并且你的爱德借着做补赎，按照我们的指示，通过适当的赔补洁净了自己，要知道弟兄爱德的恩宠已归还于你，并当感恩，因为你重新被接纳进入我们的共融：因为，正如我们应该对坚持错误的人严厉，同样我们应对回心转意的人仁慈宽恕。如今，既然你的兄弟知道他已经恢复了与宗座的共融，就让他按惯例派一个人来我们这里，接受并带给他肩带。因为我们不容许非法之事发生，却也绝不拒绝合乎惯例之事。此外，虽然我们地位的职责本可以要求我们让步，但我们更大大受到我们最甜美、最卓越的儿子，总督大人\Person{加利尼各}的请求所迫，要求我们温和地对待你。我们不能违背他极其亲爱的意愿，也不能使他忧伤。\par

\SourceRecord{Translated by James Barmby.；Nicene and Post-Nicene Fathers, Second Series；Vol. 13.；Edited by Philip Schaff and Henry Wace.；Buffalo, NY: Christian Literature Publishing Co.,；1898.；<http://www.newadvent.org/fathers/360209081.htm>.}

% structure: p0001 source=h1 target=section reason=page-title

\section{卷九，信函八十二}

\Emph{致\Place{君士坦丁堡}执事\Person{阿纳托利}}。\par

\Person{额我略}致\Person{阿纳托利}等。\par

对于善良而虔敬的儿子，我们值得费力如此回应，以致加倍偿还我们所欠的债，这本该由我们主动赐予他们的。鉴于持此信者，我们的儿子、卓越的\Person{玛尔切利诺}，在我们兄弟与同为主教\Person{玛克西穆}以及伊斯特里亚人的事上如此行事，并热心为教会的利益效力，因此，为使他能更充分地不仅在言语上，也在行动上显示其真诚的情感，我们在此恳请你的爱德，在他来到皇城时，以全副热忱与诚意与他合作，并尽力以你所能的一切援助支持他，使他蒙全能天主及你的助佑，在那边少遇困难。你也要注意如此款待他，如同对待真正属于我们自己的人，并将你爱德的效能如此施予他，使他既能认出现世对他过去的回报，也能因他承诺在教会服务中所要表现的虔诚而对未来的报酬抱持极大的希望。然而，就我们所知，最静谧的主上皇帝曾命令我们前述卓越的儿子立即前往觐见，因此，你宜寻找机会告知，滞留他的并非违命的过错，而是我们兄弟与同为主教\Person{玛克西穆}的事。这件事虽然延迟，但经他的努力，终已解决。但我们希望你的爱德注意这一点：切勿让自己牵连到任何压迫穷人的事中，以免迫于权势的压力，做出不利于你灵魂的事。因此，在敬畏天主中处理一切事务时，要特别思念永恒的赏报。\par

\SourceRecord{Translated by James Barmby.；Nicene and Post-Nicene Fathers, Second Series；Vol. 13.；Edited by Philip Schaff and Henry Wace.；Buffalo, NY: Christian Literature Publishing Co.,；1898.；<http://www.newadvent.org/fathers/360209082.htm>.}

% structure: p0001 source=h1 target=section reason=page-title

\section{\WorkTitle{第九卷，第九十一封书信}}

\Emph{致\Place{纳波里}（\Place{拿波里}）主教\Person{福图纳图斯}。}\par

\Person{额我略}致\Person{福图纳图斯}等。\par

既然我所派往\Place{纳波里}城的天主仆人之父，已按天主的安排，如祂所愿，离世而去，我认为应当差遣持信者，隐修士\Person{巴尔巴提亚努斯}，去管理同一批隐修士。目前我们决定他先担任\OriginalTerm{副院}{Prior}，如此，倘他的生活能博得你的\OriginalTerm{兄弟之爱}{Fraternity}的认可，你可稍后祝圣他为他们的\OriginalTerm{院长}{Abbot}。因他有些好品质值得推荐。但他有一个大缺点，就是过于自以为聪明。显然可知，从这根源能生出多少罪的分支。因此，你的圣德要谨慎监督他；若你发现他治理有方且内心谦逊，那么，在得天主允准后，可擢升他至\OriginalTerm{院长}{Abbot}尊位。但若他在谦逊上进步甚微，则应延迟他的祝圣，并向本人报告。\par

\SourceRecord{Translated by James Barmby.；Nicene and Post-Nicene Fathers, Second Series；Vol. 13.；Edited by Philip Schaff and Henry Wace.；Buffalo, NY: Christian Literature Publishing Co.,；1898.；<http://www.newadvent.org/fathers/360209091.htm>.}

% structure: p0001 source=h1 target=section reason=page-title

\section{\Emph{第九卷·第九十三封信}}

\Emph{致古尔法里斯，军事长官}。\par

\Signature{格列高利致古尔法里斯等}\par

这些呈递书信的人从伊斯特利亚地区来到我们这里，报告了关于阁下如此美好的事情，以致我们热切地激起向您致谢。因为我们得知，在受托治理该地区的诸多责任中，您特别渴望赢得灵魂，并且付出如此努力，召回游荡者的心归向教会的合一，以至于按您的愿望，您希望那里没有一个人与宗徒教会分离；而宗徒之长伯多禄的爱如此激发您，以至于您全心渴望重建那由万有之源的主交付钥匙之人的羊栈。光荣的儿子，从如此伟大的工作中，您应当充满信心地期待天主赏报，其中不仅有我们的劝诫，宗徒的话也确认给您，因为那使罪人从错误道路中悔改的，将救他的灵魂免于死亡，并遮盖许多罪过\EditorialNote{雅各伯书第五章}。因为，无论现世的富足或兴盛有多大，它都有终结——死亡的界限。但您所采取的赢得灵魂的追求，保留了其希望的确定性，即永生的赏报。因此，我们以父亲般的爱问候您，劝勉阁下要更热切地实践合一之热忱，这合一之源天主亲自赐予您；并且，凡您所能从分裂的错误中召回怀抱慈母教会的人，要以不断的劝诫爱护他们。还要做到这一点：以您辩护的帮助保护那些主通过您赐予归回祂羊栈的人，使那些仍在错误中的人无处可去控告那些返回正确意见的人。因为，当您在地上维护天主的事业时，祂自己将以其保护的援助顺利引导您的行动，并且，在您所渴望的永生中，将为您如此伟大的善行存留赏报。\par

\SourceRecord{Translated by James Barmby.；Nicene and Post-Nicene Fathers, Second Series；Vol. 13.；Edited by Philip Schaff and Henry Wace.；Buffalo, NY: Christian Literature Publishing Co.,；1898.；<http://www.newadvent.org/fathers/360209093.htm>.}

% structure: p0001 source=h1 target=section reason=page-title

\section{第九卷，第九十四封信}

\Emph{致\Person{罗马努斯} \OriginalTerm{护守者}{Defensorem}}。\par

\Person{额我略}致\Person{罗马努斯}等。\par

持此信者自伊斯特里亚地区来此寻找他们现今居住在西西里地区的主教，他们请求我们助其成行，我们已为他们安排行程。因此，你的经验应接待他们，并安排他们尽快抵达所述主教处；以免如他们所担心的，当地其他分裂者抢先说服他们。因为据他们表示，这位主教本人有意为了信德的合一而前来与我们相会。故应给予他们协助，使之藉主的助佑，成就他们所渴望的善事。但你的经验应亲自——若就近则亲往，否则以书信——劝勉这位主教，切勿迟延，速速藉主恩宠奔赴宗徒们的门槛，他必被我们以全部爱德接待。我们也愿你支付他前来我们这里的旅费。然而，若他认为前来此处负担过重，而安排定居西西里，并同意在提供保证后，留在教会的合一之中，置身于圣经的曲解者之间，此事亦请速速告知我们，以便我们藉主助佑，安排如何供给他那里的费用。此外，也请为持此信者能抵达所述主教处提供合作与援助，使他们离开我们后仍能得到同样的照料。\par

\SourceRecord{Translated by James Barmby.；Nicene and Post-Nicene Fathers, Second Series；Vol. 13.；Edited by Philip Schaff and Henry Wace.；Buffalo, NY: Christian Literature Publishing Co.,；1898.；<http://www.newadvent.org/fathers/360209094.htm>.}

% structure: p0001 source=h1 target=section reason=page-title

\section{第九卷，第九十八封信}

致拉文纳总管特奥多尔\par

额我略致特奥多尔等\par

虽然我们长期从我们\Emph{代表}的报告听到许多关于您的好消息，使我们的心喜乐，但如今我们的儿子、隐修院院长普罗布斯已回到我们这里，他进一步报告了您荣耀的爱德的事，这些事正应被述说，论及一个真正善良且最基督徒的儿子。既然他告诉我们您如此友善，且在安排和平时如此恳切，甚至在我们以前在您们地区的本城公民中也未出现过，我们祈求天上的保护仁慈地在此世和来世，在身体和灵魂上报答您，因为您没有停止为许多人的益处而行事。\par

因此我们告知您，阿里乌尔夫已宣誓遵守和平，不是像他的国王那样宣誓，而是在这样的条件下：任何过犯都不可对自己做出，且无人应进军对抗阿罗吉斯的军队。这全然不公且狡诈，我们视之为他未曾宣誓——因为从某种程度上，他将轻易为自己找到越界的借口，如果我们不防备他，他将更欺骗我们。\par

但是瓦尔尼夫里德，阿里乌尔夫在各方面都听从他的意见，却完全不屑于宣誓。因此，事情变成这样：从我们如此渴望的和平中，我们在这些地区几乎得不到任何补救，因为我们仍须，且将来，像至今一样防备同样的敌人。\par

此外，请贵荣耀知悉，被派往这里的国王的人催促我们签署条约。但我们记起阿吉卢夫据说曾对最显赫的巴西略所说的侮辱，那是通过我们对蒙福伯多禄的伤害，尽管阿吉卢夫本人完全否认了这一点，我们仍认为谨慎起见，暂不签署，免得我们作为他和我们最卓越的儿子即总督大人之间的请愿者和调停人，在任何方面发现自己被欺骗，万一有什么东西暗中被撤除（\Emph{即从条约中}），且他找到不赞同我们请求的借口。因此我们恳求，如同我们也已请求我们前述最卓越的儿子一样，愿贵荣耀，以您与我们相连的爱德，采取措施，使这些人在从阿罗吉斯那里返回之前，国王给他们寄送快信，但却是转交给我们的，命令他们不要要求我们签署。但如果这样有用，我们将使我们光荣的兄弟，或一位主教，或至少一位总执事来签署。\par

关于奥古斯都，我们感谢您，并正留意他依法公平地解决他与对手的案件；我们不愿将出庭对质的麻烦强加于他，但也不愿拒绝给其对手公正。\par

至于其他事项，既然迄今我们未能充分感谢您，将来我们将派我们的\Emph{代表}到您那里，通过他，在天主的仁慈下，我们或许能在彼此相连的爱德中更加紧密。此外，您的荣耀的忧伤深深影响我们；但既然有智慧的人知道所有能说的安慰之道，我们便省略以言语安慰您；但我们以祈祷陪伴您，恳求全能的天主在祂慈爱的保护下守护您和您全家的生命与健康，并在苦难中安慰您的心灵。\par

\SourceRecord{Translated by James Barmby.；Nicene and Post-Nicene Fathers, Second Series；Vol. 13.；Edited by Philip Schaff and Henry Wace.；Buffalo, NY: Christian Literature Publishing Co.,；1898.；<http://www.newadvent.org/fathers/360209098.htm>.}

% structure: p0001 source=h1 target=section reason=page-title

\section{第九卷，第一○五封信}

\Emph{致\Person{塞雷努斯}，\Place{马西利亚}（\Place{马赛}）主教。}\par

\Person{额我略}致\Person{塞雷努斯}等。\par

我们这么久才给您写信，并非由于懒惰，而是事务繁忙。现在我们由衷地向您推荐这封书信的呈送人，我们最亲爱的儿子，我们隐修院院长\Person{基里亚科}，希望您不要让他滞留在\Place{马西利亚}城，而是在您圣座的帮助下，在天主保护下继续前往我们的兄弟和同侪主教\Person{西亚格里乌斯}那里。\par

此外，我们通知您，我们听说您看到某些拜像者，就打碎并扔倒了教堂中的这些图像。我们确实称赞您对于任何手工制品成为朝拜对象的热情；但我们告诉您，您本不该打碎这些图像。因为教堂中使用图像，是为使那些不识字的人至少可以通过观看墙壁来阅读那些他们在书中不能读到的内容。因此，您应当既保存图像，又禁止人民朝拜它们，为的是使那些不识字的人有途径了解历史，并且使人民绝不因朝拜图像而犯罪。\par

\SourceRecord{Translated by James Barmby.；Nicene and Post-Nicene Fathers, Second Series；Vol. 13.；Edited by Philip Schaff and Henry Wace.；Buffalo, NY: Christian Literature Publishing Co.,；1898.；<http://www.newadvent.org/fathers/360209105.htm>.}

% structure: p0001 source=h1 target=section reason=page-title

\section{\WorkTitle{第九卷，第一〇六函}}

\Emph{致\Person{西格留斯}、\Person{埃特里乌斯}、\Person{维吉留斯}、\Person{德西德里乌斯}主教们}。\par

\Person{额我略}致高卢主教：\Person{西格留斯}（\Place{奥古斯托杜努姆}（\Place{Autun}））、\Person{埃特里乌斯}（\Place{鲁格杜努姆}（\Place{Lyons}））、\Person{维吉留斯}（\Place{阿雷塔莱}（\Place{Arles}））、\Person{德西德里乌斯}（\Place{维也纳}（\Place{Vienne}））。\Emph{A paribus}。\par

我们的元首基督，为此意愿我们成为祂的肢体，好藉着爱德与信德的联系，使我们在祂内成为一个身体。我们应当全心依附于祂，因为离了祂我们一无所是，藉着祂我们才能成为我们所蒙召的。任何事物都不得使我们与元首的堡垒分离，免得我们若拒绝作祂的肢体，被抛在祂之外，如同从葡萄树上砍下的枝条而枯萎。因此，为使我们堪当成为救赎主的居所，让我们以全心的热诚常居于祂的爱内。因为祂亲自说：\BibleQuote{爱我的，必遵守我的话，我父也必爱他，我们要到他那里去，与他同住}\BibleRef{若 14:23}。但是，既然我们若不从自身剪除贪欲——一切罪恶的根源——就无法紧靠万善的根源，因此我们藉着这些文字（这些文字使我们如同在渴望的互访交谈中彼此联系），按照宗徒的规诫接近你们的兄弟情谊，好使我们依靠教父的规则和主的命令，从信德的殿中驱逐贪财——即偶像崇拜，以免在主的房屋内容忍任何有害或无序的事。\par

我告知你们，我们早已听闻俗传，在高卢地区，圣职是通过\OriginalTerm{西蒙尼异端}{simoniacal heresy}授予的。我们满怀忧愤，如果金钱在教会职务中有了位置，神圣之事变得世俗。那么，任何人若以支付代价来买得此物，看重的是头衔而非职务，他渴望的不是作司铎，而只是被称作司铎。这算什么？其结果岂非毫无品行考验，毫不关心道德，毫不查问生活，而唯有能出价的人被视为有资格？由此，如果用真正的天平来衡量，当一个人邪恶地急于攫取有利益的职位以图虚名，他恰恰因追求尊位而更加不配。此外，正如一个被邀请却拒绝、被寻找却逃避的人，应该被带到神圣的祭台前；同样，一个主动请求并强行推进的人，无疑应被驱逐。因为凡如此努力攀升高位者，他在增长中岂非减少，在外在上升中岂非内里下沉？所以，最亲爱的弟兄们，在祝圣司铎时，让真诚取胜，让单纯的同意不带买卖，让纯洁的选举被优先，好使晋升到司铎最高位被认为是出于天主的判断，而非出卖者的投票。因为以价格欲图获取或出售天主的恩赐，是严重的罪行，福音权威可以作证\EditorialNote{玛窦福音第21章}。\par

因为，当我们的主救赎主进入圣殿时，祂推倒了卖鸽子之人的座位。卖鸽子不是别的，正是为覆手而收取代价，将全能天主赐予人的圣神出售。而如此行事之人的司祭职在神眼中坠落，由座位的推倒清晰地表明。然而这邪恶的乖戾仍然发挥其力量。因为它驱使那些被欺骗去购买的人去出售。而且，当人们不留意那神圣声音所命令的：\BibleQuote{你们白白得来的，也要白白施舍}\BibleRef{玛 10:8}，此事就增加并加倍成为同一个罪的传染，即购买者和出售者。而且，众所周知，这异端以毒根在教会中比其他一切更早潜入，并在其起源就被宗徒的憎恶所谴责，为何不防备它？为何不考虑，祝福对那为成为异端者而被提升之人，反成了诅咒？\par

因此，在大多数情况下，灵魂的仇敌既然不能将明显的恶事灌输给人，便设法用一个虔敬的假象来取而代之，诱使人或许以为应当接受有钱之人的财物，以便有东西可以分给穷人，只要他能借此将隐藏在施舍外表下的致命毒药注入。因为猎人若公开设网，就不能欺骗野兽；捕鸟者若明目张胆，就不能捕鸟；渔夫若不用鱼饵藏钩，就不能捕鱼。所以，我们必须处处提防敌人的诡计，以免那些不能被公开诱惑所颠覆的人，反而被隐藏的武器更残忍地杀害。因为，若将非法得来的财物分施给穷人，这并不算作施舍；因为怀着良好意图而错误地接受，结果只有更坏，不会更好。中悦我们救主眼目的施舍，不是那些用非法途径和不义手段聚敛来的，而是出于天主赐给我们的正当所得。因此，这一点也是确定的：即使用为圣秩而给的金钱修建隐修院、医院或其他建筑，也无益于获得赏报；因为当一个邪恶的、收买尊位的人被安置在圣地，并像自己一样用金钱任命别人时，他因邪恶的管理所造成的破坏，远超过接受他钱财而受祝圣之人所能建立的。因此，我们不应在施舍的伪装下试图用罪恶获取任何东西。圣经清楚地警告我们：\BibleQuote{\Emph{恶人的祭品，本就可憎，更可憎的是用邪恶来奉献。}}（\BibleRef{箴 21:27}）因为凡在天主的祭品中用邪恶所奉献的，非但不能平息全能天主的义怒，反而激怒祂。因此又记载：\BibleQuote{你要以正义的劳作光荣上主。}（\BibleRef{箴 3:9}）所以，凡以邪恶手段取得，自以为能善加施舍的人，显然并没有光荣上主。为此，撒罗满也说过：\BibleQuote{凡用穷人的财物作祭献的，犹如在他父亲眼前杀害他的儿子。}（\BibleRef{德 34:24}）现在，让我们想想，若一个儿子在父亲眼前被杀，父亲的悲痛有多大；由此我们便容易明白，当天主接受来自掠夺的祭献时，祂是何等悲伤。因此，最亲爱的弟兄们，应当极力避免在施舍的伪装下犯下\OriginalTerm{买卖圣职的异端}{simoniacal heresy}的罪。因为，因罪而施舍是一回事，而因施舍而犯罪则是另一回事。\par

此外，我们听到的另一件事也当以同样不轻的憎恶对待，即有些人因贪图尊位，在主教去世后便剃发，从平信徒骤然成为司铎，无耻地抢夺修道生活的领导权，甚至还未学习如何作兵士。这些人尚未踏入门徒的门槛，就胆敢占据师傅的位置，我们能期望他们为属下带来什么好处呢？在这种情况下，即使某人功德无疑，也应当经过不同的等级在教会职务中得到操练。他应看见自己所要效法的，被塑造成应持守的形态，以便日后被选为迷途者指明生活之路时，不致失误。因此，他应长期在宗教默想中受磨练，才能中悦天主，并像放在灯台上的灯一样发光，使反对的强风虽吹向学问点燃的火焰，却不至于熄灭，反而增强。因为经上记载：\BibleQuote{先受考验，然后任职。}（\BibleRef{弟前 3:10}）更何况被选为百姓代祷者的人，更应先受考验，以免不良的司铎成为百姓丧亡的原因。因此，对此事无可推诿，无可辩护，因为所有人都清楚知道，那位圣善卓越的导师多么关注此事，他禁止新手接受圣秩\EditorialNote{参阅弟茂德前书第三章}。不过，那时称新手，是指刚在圣洁信仰生活中栽植的人；现在，凡突然栽植于宗教习尚，进而谋求圣秩高位的人，也应视为新手。所以，圣秩应当按部就班地升迁；因为藐视阶梯，以陡峭的攀登直奔最高处的人，是自招跌倒。既然同一宗徒在关于圣秩的其它训示中教导他的门徒：\BibleQuote{不可急速给人覆手。}\EditorialNote{参阅弟茂德前书第五章}那么，还有什么比从最高开始，即尚未担任执事就先作主教，更急促、更鲁莽的呢？因此，凡欲获得司铎职位者，不是为了得意的虚荣，而是为了行善，就应先以自己的能力衡量所要承担的担子；若能力不足，就当戒避；即使自认为胜任，也当怀着敬畏去接近。\par

此外，从理性受造物引申论证之后，再引用我们对无理性事物的使用经验，也不为多余。因为适合建筑的木材从森林中砍伐，但建筑的重量并不加在尚未干燥的木材上，而要等待多日的延迟使其青绿干燥，合于必需之用。若偶然忽略了这一预防，木材很快就会被压在其上的重物折断，原本用来支撑的材料反而导致了崩塌。\par

盖以此故，医师之疗身者，于所需之人，不即施新制之药，而姑置之，使浸渍有时。若有人未成熟而予之，则健康之方必成危险之由，无可疑也。故当学习，司铎于其职分内当学习，即受托管理灵魂之治疗者，当观察诸艺之士于理性指导下所留意者，且当自制其野心，若非出于畏惧，亦当出于羞耻矣。\par

然恐或有仍欲以恶习为借口而自辩护者，愿尔等之弟兄团契以理性之缰绳约束之，不容其堕入非法之行；盖凡应受罚者，不当引为效法之例，而当以为纠正之资。\par

再者，亦不能容尔等疏忽另一事，此亦需纠正。盖若守一切而他处任敌乘隙而入，何益之有？故当禁止妇人与居任何圣秩之人同居。关于此事，恐人类之旧敌得志，当由众人同意规定：除圣典所许者外，不得有其他妇人同住。虽此禁令或于一时为某些人所苦，然无疑日后因其对灵魂之裨益而转甘，若敌人于其所能胜人之事上被克服焉。\par

于此关怀之一部分，亦不可忽略教父们为利益所定关于在各教区召开会议之规定。故为避免弟兄间有纷争，或上下之间有嫌隙，司铎应聚集，以讨论所生之事，并于教会惯例作有益之商议；俾往者得正，来者得规，全能之主因弟兄们之一心而处处受赞颂。须知尔等之聚会中，有谁之临在？经载：\BibleQuote{\Emph{凡两三人因我名聚集者，我即在其中}} \BibleRef{玛 18:20}。若于两三人处，主肯临在，则于众多司铎聚集时，岂不更不缺席？且教父规则所定每年两次会议，并非不知。但恐或有必要不能实施此规，我等仍命每年至少一次，不得有借口，必须举行；如此，在期待会议期间，无人敢行邪恶或非法之事。盖通常，虽非出于爱正义，然因惧审查，人常戒绝已知为众判所不悦之事。至爱之弟兄，愿我等守此规矩，传于后人。且当思圣经所载以训我等之一切，尽力劝人遵循。盖确知若全心留意此有益之训，则脱诸恶之污，因倚赖此建树，则无疑杜绝一切欺诈之隙。\par

故为以上目的，愿尔等之弟兄团契，若天主愿意，召集会议，并于其中，藉我至敬之兄弟、同主教\Person{Aregius}，及我至爱之子\Person{Cyriacus}之调停，凡如前所述违反圣典之事，皆当受绝罚之严厉谴责。即：任何人不得擅自以财物获取圣职，或受财物而授予圣职；或任何人不得径由平身骤登管理之位；或除圣典所许者外，不得有其他妇人与司铎同居。关于诸事，我至敬之兄弟\Person{Syagrius}主教当与全体会议，于我至爱之子\Person{Cyriacus}返回我处时，将所行之事告我；使我能确知所定之规、所用之保障与方式，对全能天主为尔等之生活与行止无间感恩。\par

\SourceRecord{Translated by James Barmby.；Nicene and Post-Nicene Fathers, Second Series；Vol. 13.；Edited by Philip Schaff and Henry Wace.；Buffalo, NY: Christian Literature Publishing Co.,；1898.；<http://www.newadvent.org/fathers/360209106.htm>.}

% structure: p0001 source=h1 target=section reason=page-title

\section{卷九，第一〇七书}

致\Person{阿雷吉乌斯}，\Place{瓦平库姆}主教。\par

额我略致高卢主教\Person{阿雷吉乌斯}。\par

你的兄弟情谊所受的苦难，我们得知你因失去子民而遭受的苦难，使我们深感悲痛，以致爱德使我们两人合而为一，我们感到我们的心特别与你的苦难相连。但在此之中，我们因你已使你的心灵懂得如何耐心承受忧伤，并因对另一生命的希望而不为死亡长久悲伤，而得到很大安慰。尽管如此，惟恐某种苦难仍存于你的灵魂中，我劝你从悲伤中安息，停止忧愁。因为为那些我们相信他们藉死亡已获得真生命的人而沉溺于痛苦的烦恼是不合宜的。那些不认识另一生命、不相信从这个世界能转到一个更好世界的人，或许有长久悲伤的合理借口。然而，我们这些知道、相信并教导此事的人，不应为离世者过于忧伤，以免在别人看来是爱德的表现，在我们身上反而成为责备。因为，因悲伤而受折磨，违反每人所宣讲的，似乎是一种不信任，正如宗徒所说：\Quote{\BibleQuote{但我们不愿你们在弟兄们安眠的事上无知，以免你们悲伤，像其他没有望德的人一样。}}\BibleRef{得前 4:12}\par

因此，最亲爱的弟兄，我们面前有了这个理由，就应如同我们所说，不要为死者忧伤，而应将爱德施于生者，对他们同情才有益处，爱德才能结出果实。让我们从此藉着责备、劝勉、说服、安抚、安慰，尽我们所能去帮助所有人。让我们的舌头成为善人的鼓励，恶人的刺激；打击傲慢者，平息愤怒者，激励怠惰者，藉劝勉点燃闲散者，说服退缩者回来，安抚粗暴者，安慰绝望者；这样，既被称为领袖，我们就应为前进者指教救恩之路。让我们在防守上保持警醒，防范敌人诡计的一切途径。如果有时错误将委托给我们的羊群中的一只羊引上歧途，让我们竭尽全力将它带回主的羊栈，这样我们因所背负的牧者之名，得到的不是惩罚，而是赏报。鉴于此一切都需要天主恩宠的帮助，让我们以不断的祈祷恳求全能天主的仁慈，为使祂赐给我们行这些事的意愿和能力，并藉善工的果实，引导我们走上祂自称是牧者之牧者的道路；这样，藉着祂——没有祂我们无法做成任何事——我们才能完成一切。\par

此外，我们共同的儿子，执事\Person{伯多禄}告诉我们，你的兄弟情谊在此地时曾请求我们允许你及你的总执事使用\OriginalTerm{达马提克祭披}{dalmatics}。但由于受你子民疾病的催迫，你如此匆忙地离开，以致你心中的悲伤不容你再坚持此事，尽管这是合宜的且符合你请求的性质；又因为我们有许多事务，而教会礼仪的考虑不允许我们轻率突然地批准一件新事；因此，所请求之事的执行被长期推迟。然而，现在回忆起你爱德的善功，凭此我们权威的文书，我们批准你的请求，并已准许你或你的总执事以使用达马提克祭披作为装饰；并且我们已通过我们最亲爱的儿子，院长\Person{基里亚库斯}之手，送去了同样的达马提克祭披。\par

此外，在会议中——我们已命令通过我们的兄弟及同僚主教\Person{西阿格里乌斯}召集，以反对\OriginalTerm{西满异端}{simoniacal heresy}——我们希望你出席；并且我们已下令，将我们已为上述兄弟送去的\OriginalTerm{披带}{pallium}相应地交给他，条件是他承诺通过会议的定义，将我们所禁止的不合法事务从神圣教会中除去。关于这次会议，我们希望你以信件详尽地告诉我们所有进程，以便你——我们深知你的圣德——能向我们报告一切。\par

\SourceRecord{Translated by James Barmby.；Nicene and Post-Nicene Fathers, Second Series；Vol. 13.；Edited by Philip Schaff and Henry Wace.；Buffalo, NY: Christian Literature Publishing Co.,；1898.；<http://www.newadvent.org/fathers/360209107.htm>.}

% structure: p0001 source=h1 target=section reason=page-title

\section{第九卷 第一〇八封信}

\Emph{致\Person{叙亚格留}主教}\par

\Person{额我略}致\Place{奥古斯托杜努姆}（\Emph{Autun}）主教\Person{叙亚格留}。\par

万善之母是爱德，它不与任何外物、任何粗鲁、任何混乱相干；它如此操练并坚固人心，以致一切都不沉重，不难，凡所作的都成为甘甜。既然，它的特有性质是促进和谐的事物，保全联合的事物，联合分离的事物，纠正错误的事物，并用自身完美的壁垒巩固所有其他德行，凡将自己嫁接到它根上的人，既不失去青翠，也不缺乏果实，因为有效的工作不会失去丰饶的滋润。因此，我非常喜爱你，我至爱的兄弟，并在主内为你喜乐，因为我发现你，根据许多人的见证，如此赋予了这同一爱德，以至于你既得体地表现出司铎应有的风范，又可称赞地为他人树立了效法的榜样。\par

既然，在宣讲工作中（我经过长时间思考，已通过当时我的隐修院的院长（\OriginalTerm{院长}{præpositum}）、现在我们的兄弟和同为主教的\Person{奥斯定}，悉心向盎格鲁民族提供），我发现你理所当然地如此热忱、虔诚，并在各方面协助，以至于在此事上我大大亏欠于你，因此，考虑到如此大的恩情，我不忍放下你的请求，以免在你面前显得无用。因此，按照你请求的内容，我们已在主内提供，让你因使用白羊毛披带（\OriginalTerm{白羊毛披带}{pallium}）而尊贵，佩戴于你自己的教堂内，仅限举行弥撒时。然而，我们决定只在你首先承诺通过教务会议的决定修正我们命令纠正的事项的条件下才赐予你；因为我们确实认为，既知你以天主仁慈在其中杰出的庄重性情，外在服饰的更显赫装饰应归于你；尤其是我们认为你请求它，并非出于无谓的傲慢炫耀，而是为了你教会的特点和尊严。并且，免得在这件服饰上我们似乎只是给予一种空泛的恩赐，我们同时考虑也赐予这一点——即，在总主教各方面保留其地位和尊严的情况下，奥古斯托杜努姆教会应位于\Place{吕格杜努姆}（\Emph{Lyons}）教会之后，并借我们权威的恩准主张这一位置和等级。至于其他主教，我们规定，应按照祝圣日期就座，无论是在会议中、签署文件或任何其他事务中，并应主张各自等级的特权：因为我们觉得合乎情理，在使用白羊毛披带的同时，我们也应如上所述一同赐予一些特权。\par

但是，既然随着尊贵的增加，责任感也应当增加，以使行为的装饰与服饰的装饰相称，你应在一切追求中更热切地操练自己。对你的下属的行为保持警惕；让你的榜样成为他们的教训，你的生活成为他们的导师。通过你口中的劝诫，让他们学习所当畏惧的，并被教导所当热爱的；这样，当你把托付给你的塔冷通连本带利交还时，在报应之日，你将被认为配得听到：\BibleQuote{做得好，善良忠信的仆人：进入你主的喜乐吧。}\BibleRef{玛 25:23}\par

\SourceRecord{Translated by James Barmby.；Nicene and Post-Nicene Fathers, Second Series；Vol. 13.；Edited by Philip Schaff and Henry Wace.；Buffalo, NY: Christian Literature Publishing Co.,；1898.；<http://www.newadvent.org/fathers/360209108.htm>.}

% structure: p0001 source=h1 target=section reason=page-title

\section{第九卷，第一〇九书函}

\Emph{致\Person{布倫希尔德}，法兰克王后}\par

\Person{额我略}致\Person{布倫希尔德}等\par

既然陛下在治国诸事上的王室关怀值得赞扬，为增进您的荣耀，您应当更加警醒，谨慎，不可让您以智谋治理的人在外表之下于内心丧亡。如此，您便能通过敬虔关怀的果实，在占据现世王国的至高之位后，在天主之下获得永恒的王国和喜乐。我们相信您能通过以下方式获得成功：若您在其他善行之外，关注司铎的圣秩授予；据我们所知，在您的地域，这一职务已变得如此受人觊觎，以至于司铎常立即从平信徒中被祝圣。此事极为严重。因为那些贪图被选为主教的人——并非为行善，而是为显赫——能有什么成效？他们能对人民有什么益处？这些人尚未学会必须教导之事，又有什么成效？不过是少数人的非法晋升成为多人的丧亡，并且教会治理的遵守陷入混乱，因为没有遵守常规秩序。因为但凡轻率匆忙地前来掌管治理的人，能用什么样的劝诫来建树他治下的人呢？他的榜样教给他们的不是理性，而是错误。诚为可耻，诚为可耻，命令别人遵守他本人不知如何遵守的事。\par

我们也不略过另一件同样需要改正的事，反而视其为彻底可憎、极其严重的事而加以憎恶；那就是在您的地域，圣秩因\OriginalTerm{买卖圣职的异端}{simoniacal heresy}而授予，这异端最先起来反对教会，并被严厉诅咒定罪。因此，司铎职的尊严受到轻视，神圣的尊荣受到谴责。于是敬畏消亡，纪律被摧毁，因为本该纠正过错的人却犯了过；由于邪恶的野心，可敬的司铎职反受责难和贬损。因为谁还会尊敬被买卖之物？谁不认为买来的东西毫无价值？因此我深感悲痛，并为那地哀叹；既然他们不屑于将圣神当作天赐来接受，反以贿赂获得全能天主愿借覆手赐予人的圣神，我不认为司铎职能在那里长久存续。因为天上的恩宠恩赐被卖掉时，所求的不是为天主服务的生命，而是金钱在相反天主处受崇拜。鉴于如此重大的邪恶不仅对他们构成危险，也对您的王国损害不小，我们以父爱问候陛下，恳求您通过改正这一严重罪行而使天主眷顾您。而且，为消除今后犯此罪的机会，应奉您的命令召开一次主教会议，在我们极可爱的儿子、院长\Person{西里亚库斯}出席下，严格禁止任何人敢从平信徒身份骤然晋升至主教品级，或任何人胆敢为圣职给予或收受任何事物，违者处以绝罚；如此，我们的主和救赎主将如何对待属于您的事物，正如祂看见您以虔诚热忱关心属于祂的事物。但我们特别将这次我们决定召开的主教会议的职责和管理委派给我们的兄弟和同为主教的\Person{西雅格里乌斯}，我们知道他特别属于您；我们恳请您屈尊愿意倾听他的恳求，并以您的援助支持他；以便那会归于您赏报的事，即虔诚且悦乐天主的司铎圣秩授予，因这邪恶毒害被移除，而在您管辖的所有范围内生效。\par

至于这位我们的兄弟，因他在天主之下被派往盎格鲁民族的使命中显明自己极为热诚，我们送去了一件\OriginalTerm{披肩}{pallium}，用于弥撒的隆重礼仪，以使他在属灵之事上给予援助后，能因宗徒之长的恩宠而在属灵秩序上得以提升。\par

此外，我们十分惊异，为何在你的王国中你允许犹太人拥有基督徒奴隶。因为基督徒若非基督的肢体，又是什么？我们都知道你真诚地尊敬这些肢体的头。但请陛下考虑，尊敬头却让肢体被祂的敌人践踏，是何等矛盾。因此我们恳请陛下的法令从你的王国中消除这不义的祸害；这样你便能证明自己更配作全能天主的崇拜者，因为你将祂的信徒从祂的敌人手中释放出来。\par

\SourceRecord{Translated by James Barmby.；Nicene and Post-Nicene Fathers, Second Series；Vol. 13.；Edited by Philip Schaff and Henry Wace.；Buffalo, NY: Christian Literature Publishing Co.,；1898.；<http://www.newadvent.org/fathers/360209109.htm>.}

% structure: p0001 source=h1 target=section reason=page-title

\section{第九卷，第一一〇封信}

\Emph{致法兰克国王狄奥多里克与狄奥德贝尔特}.\newline{}\par

格里高利致狄奥多里克等人。\newline{}\par

既然贵国声名远播，因着基督宗教的恩宠，历久以来超越其余诸国，故当竭力在你们比别国更为荣耀之处，更完美地取悦那赐予君王康泰与财富的全能之主，并使你们所持守的信德在各方面都有助于你们。最杰出的儿子们，我们本只愿以友爱的问候向你们致辞，以尽爱德之职，表达为父之情。然而，因一不法之事令我们极为忧心，宜当兼顾两事：既表问候，亦不可对那需纠正之事缄默不语。若你们留心细察，将发现我们所言全然是为你们福祉的稳妥。\newline{}\par

据闻，在贵国境内，买卖圣职异端（此乃魔鬼最初植入教会内，并自兴起之初即遭惩罚之剑击打与谴责的异端）盛行，虽则司铎应兼具信德与善行。\newline{}\par

若缺乏善行，信德便无价值，正如圣雅各伯所证言：\BibleQuote{没有行为的信德是死的}\BibleRef{雅 2:18}。然而，一位司铎若被证实以贿赂谋取如此伟大圣事的尊位，他又能有何善行？由此导致，连渴望圣秩之人也不努力改善生活或整顿品行，反忙于聚敛财富以换取神圣尊位。于是，无辜与贫穷者因被排斥与轻视，而对圣秩畏缩不前。穷人无辜反遭厌弃，而另一方的贿赂无疑称许了过失；因为凡黄金悦人之处，恶习亦然。因此，这不仅给祝圣者与受祝圣者的灵魂带来致命创伤，而且贵国亦因主教们的过失而受重压——本应借他们的代祷获得助佑。因为，若某人被认为堪当司铎之职，并非因其行为的功绩，而是因其贿赂的丰盈，那么，严肃与勤勉便再也无法对教会尊位提出任何要求，惟不敬的贪金得到一切。当恶习以尊位得报偿时，那本应受惩罚者反被擢升至惩罚者的地位；于是司铎被发现非但无益于人，反而自趋丧亡。因为当牧者受伤，谁能为医治羊群而施药？或那将自己暴露于敌人投枪之下者，如何能以祈祷之盾保护百姓？或那自身根株因重病而感染者，能结出何种果实？故此，在那些地方，当这种代祷者被擢升至治理之位时，更重大的灾祸便可预见，因为他们更激起天主对他们的义怒——本应借他们本人为百姓平息这义怒。\newline{}\par

此外，我们听闻教会所属田产并未纳税；我们对此深感惊异：若对连合法赋税皆获豁免者，仍索取非法款项。\newline{}\par

我们的忧虑也不容我们忽略另一邪恶：有些人受虚荣煽惑，从平信徒身份陡然攫取司铎尊位，并且（说来可耻，虽至关重大不容缄默）那些本应受管辖者不知羞愧亦不畏惧以管辖者自居，本应受教导者以导师自居。他们厚颜无耻地担任灵魂领袖，而自己对领袖应走之路全然不知，甚至不知自己走向何处。此等行径何其恶劣与冒险，即使世俗秩序与纪律亦以显明。因为，既然统军之将若非经辛劳与审慎考验便不得被选，那么，那些以不成熟之急切渴望攀登主教高位者，至少可借此比较思量：灵魂领袖应为何等之人。他们应避免骤然尝试未经考验的事工，免得对尊位的盲目野心既成为自己的惩罚，又向他人播下有害错误的种子——他们自己尚未学会所当教导之事。因此，我们以父亲般的深情致意，恳求最杰出的儿子们，务将如此可憎的恶行从贵国境内驱逐，切勿让任何借口、任何有损你们灵魂的谏言在你们心中立足；因为那有能力纠正却疏于纠正者，无疑分担行恶者的罪责。为使你们能向全能天主献上大礼，请下令召开一次主教会议，在此会议中（正如我们已嘱咐我们的弟兄与同为主教者），在我们最心爱的儿子，院长基里亚库斯面前，以绝罚为约束，制定法令：任何人永远不得为教会圣职给予或收受任何财物，任何平信徒亦不得骤升司铎；如此，我们的救赎主——你们不容其司铎被敌人从中毁灭——必在今世与来生赏报你们的这项服务。\newline{}\par

此外，我们深感惊讶，在贵国竟允许犹太人拥有基督徒奴隶。因为所有基督徒岂非基督的肢体？我们知道你们都忠实地尊崇这肢体的头：但请贵国思量，尊崇头颅却允许祂的肢体被敌人践踏，是何等矛盾。故此，我们恳求，请贵国颁布法令，从贵国境内除去此项不义之恶行，如此你们更能显示自己是全能天主相称的敬拜者，因为你们将祂忠信的仆人从祂的敌人手中释放出来。\par

\SourceRecord{Translated by James Barmby.；Nicene and Post-Nicene Fathers, Second Series；Vol. 13.；Edited by Philip Schaff and Henry Wace.；Buffalo, NY: Christian Literature Publishing Co.,；1898.；<http://www.newadvent.org/fathers/360209110.htm>.}

% structure: p0001 source=h1 target=section reason=page-title

\section{第九卷，书信第一百一十一}

\Emph{致\Place{阿尔勒}主教\Person{维吉利乌斯}（Arles）。}\par

\Person{额我略}致\Person{维吉利乌斯}等。\par

虔诚旨意的渴望与可赞敬虔的倾向，理应由司铎的热切努力时常予以扶助；因此，应当谨慎注意，切勿因懈怠、疏忽或僭妄而扰乱那些为隐修士及修道生活的安宁所制定的任何规定。但，正如理智所要求者理应得到有益的规范，同样，已经规范者不应被违反。如今，荣福的先王、法兰克国王希尔德贝尔特，因受天主教信仰的热爱所激励，在阿尔勒城墙内为自身的报偿建立了一座男隐修院，一如我们在书面记载中所见，他为院内住客的维持生活而赐予了一些财物。并且，唯恐他的目的受挫，以及为隐修士安宁所安排的事被打扰，他在信函中恳求，凡他所让与该隐修院的任何权利，均可藉权威加以确认；同时，他对他的请求另加一项，即某些特权也应同时授予同一隐修院，既涉及事务管理，也涉及院长的任命。他这样做，是因为他知道信友们如此尊敬宗座，以至于凡经由宗座谕令所定之事，此后便不会被任何非法僭越的骚扰所动摇。因此，由于君王的旨意以及所渴望之事迫切要求予以实施，我们的前任罗马主教维吉利乌斯曾致函你的前任奥勒留，其中凡为成全那目的所要求的全部事项，均乐意藉权威的支持得到确认，因为此类请求，既然提出，便不容遭遇困难。但是，为使你的兄弟团体知晓当时所规定的一切，我们已安排将前述前任的书面命令附于此信。这些命令一经阅读，我们劝你以身为司铎所应有的热忱，切实予以全部维护，勿使任何不当或非法之事强加于该隐修院，亦勿使上述命令因任何僭越而受侵犯。因为，凡经宗座权威所认可之事，从来不会缺乏效力；但我们仍额外以我们的权威再次确认我们的前任为这事的安宁所规定的一切。因此，务请你的兄弟团体在遵守这些规定时如此尽职，以杜绝一切扰乱的机会，并劝导他人实行这些事，同时你本人要以应有的谨慎和忠诚，恪守那位已亡者最虔诚的意愿。\par

\SourceRecord{Translated by James Barmby.；Nicene and Post-Nicene Fathers, Second Series；Vol. 13.；Edited by Philip Schaff and Henry Wace.；Buffalo, NY: Christian Literature Publishing Co.,；1898.；<http://www.newadvent.org/fathers/360209111.htm>.}

% structure: p0001 source=h1 target=section reason=page-title

\section{卷九，书信114}

\Emph{致\Person{Virgilius}与\Person{Syagrius}主教。}\par

\Person{Gregory}致\Person{Virgilius}（\Place{Arelate}，即\Place{Arles}的主教）与\Person{Syagrius}（\Place{Augustodunum}，即\Place{Autun}的主教）。\par

我所受委托之职的性质，最亲爱的弟兄们，迫使我发出悲痛的呼喊，并以爱德的焦虑砥砺你们的爱心，因为据说在你们的地区，你们过于疏忽和懈怠，而正义的端正与对贞洁的热忱本应燃起你们的勤勉。如今我们听闻，一位名叫\Person{Syagria}的妇女曾度奉献生活，甚至改变了装束，后来竟被强行嫁给了一个丈夫（此事令人发指），而你们却未曾因忧伤而动身干预以保护她。若果真如此，我为此更加沉重地叹息，唯恐在全能的主面前（愿天主不许）你们仅担雇工之职，而无牧者之功，竟任由一只羊在狼口中被撕咬而不加挣扎。因为你们将向未来的审判者说什么，或献上怎样的交代？你们未曾被强暴的淫邪所触动，未曾因对奉献生活的尊重而奋起护卫，未曾因司铎的职务而激励去保护贞洁的纯真！即使现在，让你们的疏忽回到记忆中；让对这过错的回忆激起你们，让你们职务的考虑推动你们去劝诫上述妇女。而且，恐怕随着时间推移，强迫竟变成了自愿同意，愿你们的舌成为她的疗药，通过你们的劝勉，让她勤于祈祷；愿忏悔的哀叹不离她的记忆；愿她向我们的救赎主展示一颗忏悔的心，并以哭泣补偿贞洁的丧失——这贞洁在她身体内未被允许保存。\par

因此，既然上述妇女据说甚至如今仍愿将她的财产用于敬虔之事，我们劝勉你们，让她在此事上体验你们兄弟之情的眷顾并享有支持，并允许她在为其子女保留适当份额后，自由决定她的财产。因为毫无疑问，你们若帮助那些愿意行善的人，你们自己也在行善。因此，最亲爱的弟兄们，请思量我们所说的话出自何等深厚的爱，并以同样的爱德精神接纳这一切。因为我们在基督内成为一个身体，凡我所感到对你们有损的，我也与你们一同燃烧。我以何等恳切和深情寄此信函，愿真理的创始者向你们的心启示。因此，愿这兄弟般的劝诫不使你们烦恼，因为即使是一杯苦药，为健康而施时也欣然接受。最后，最亲爱的弟兄们，让我们以联合的祈祷恳求我们天主的仁慈，愿祂在祂的敬畏中仁慈地安排我们的生活，使我们在此世能如司铎应做的侍奉祂，并在将来能在祂面前安然无惧地站立。\par

\SourceRecord{Translated by James Barmby.；Nicene and Post-Nicene Fathers, Second Series；Vol. 13.；Edited by Philip Schaff and Henry Wace.；Buffalo, NY: Christian Literature Publishing Co.,；1898.；<http://www.newadvent.org/fathers/360209114.htm>.}

% structure: p0001 source=h1 target=section reason=page-title

\section{第九卷，第一一五函}

\Emph{致\Place{奥古斯托杜努姆}（\Place{奥坦}）主教\Person{瑟雅格略}。}\par

\Person{额我略}致\Person{瑟雅格略}等。\par

如果在世俗事务中，每个人都应保持他的权利和应有的地位，那么在教会的安排中，更不容许有任何混乱；免得和平的恩赐本应由此而生，却让不和在此找到位置。而这样做就能确保：不向权力让步，只向公平让步。\par

据报告，我们最亲爱的兄弟\Place{塔乌里尼}城主教\Person{乌尔西基诺}，在遭受俘虏和掠夺之后，他的教区（据说位于\Place{法兰克人}境内）受到了严重损害，甚至有人不顾教会的法令，在无可指摘的罪行下另立主教。而且，唯恐这一损害之举似乎轻微，还进一步夺取了他本可持有的教会财产。如果事情果真如此，鉴于这是极其残酷的、违背神圣教规的行为——任何人的野心都不应从一个无辜的司铎那里夺走他的祭台，而他因罪孽并不应被取代——就让所有人都视他的案件为己任，竭力反对施加于他人、而他们自己也不愿忍受的事。因为若邪恶之事在长久敞开之前不被堵住入口，它就会因习惯而变得更宽；凡理性所明显禁止的，习俗也会容许。但是，超过其他一切，你的仁爱之关心，鉴于我们的推荐和你自己对天主所负的责任，应热切地致力于他的辩护，不容许他再无理地被逐出他的教区。此外，你既要亲自出面，也要向最卓越的君王们（我们相信他们在任何方面都不会使你忧伤）恳求，使这件已做错的事得以纠正，被暴力夺走的能在真理的庇护下归还；因为经上记载：\BibleQuote{一个兄弟帮助另一个兄弟，必得高举}\BibleRef{箴 18:19}，你的爱德当知，你们帮助弟兄时，越是愉快而坚定地执行天主的诫命，就越能从全能的天主那里得到更多。\par

\SourceRecord{Translated by James Barmby.；Nicene and Post-Nicene Fathers, Second Series；Vol. 13.；Edited by Philip Schaff and Henry Wace.；Buffalo, NY: Christian Literature Publishing Co.,；1898.；<http://www.newadvent.org/fathers/360209115.htm>.}

% structure: p0001 source=h1 target=section reason=page-title

\section{卷九，第一一六函}

\Emph{致\Person{狄奥多里克}与\Person{狄奥德伯特}，法兰克人的国王。}\par

\Person{额我略}致\Person{狄奥多里克}等。\par

国王的首要之善在于培养正义，维护每个人的权利，不让臣民遭受被强权所迫之事，而是合乎公平之事。我们确信您既热爱且完全追求这一点，因此邀请我们向尊座指出需要修正之事，以便我们能用信件援助受压迫者，并为您获得赏报。\par

现在他们说，我们的兄弟与同僚主教，\Place{都灵}城的主教\Person{乌尔西奇努斯}，在他的堂区内（这些堂区位于你们王国的境内）遭受非常严重的损害，以致于，违反教会礼仪，违反司铎的庄重，违反神圣教规的规定，在没有需要这样做的罪行的情况下，竟有另一人不怕在那里被祝圣为主教。而且，似乎以为仅违法还不够，除非在违法之上再加违法，据称，他教会的财产也被夺走了。如果事实如此，那么，一个人因没有罪行伤害他而遭受武力压迫，这是极其不能容忍的。我们恳请你们，首先以父爱般的爱德向你们致意，请求你们：凡是出于对教会的尊敬和对公平的重视，尊座本可自愿赐予的，请你们因我们的代祷而更加仁厚地赐予，并按照我们对您公平的良善的信赖，务必在他的一切方面秉持正义；并且，在查明真相后，请命令纠正非法之事，并将从他那里不当夺走的财产公平地归还于他。而且，不应因他的教会目前被敌人扣留而对他不利；反而，这应更加激发您基督徒的心意去援助他，使他因受到您慷慨恩赐的安慰，不至于感到他所忍受的囚禁造成的损失。因此，为了您灵魂的益处，愿我们的劝诫在您那里得到接纳，以便您能为了自己的赏报，用伸出的正义之手重新振作他的沮丧，从而因您对司铎们秉持公平，您将因他们的祈祷在天主眼前永远昌盛。\par

\SourceRecord{Translated by James Barmby.；Nicene and Post-Nicene Fathers, Second Series；Vol. 13.；Edited by Philip Schaff and Henry Wace.；Buffalo, NY: Christian Literature Publishing Co.,；1898.；<http://www.newadvent.org/fathers/360209116.htm>.}

% structure: p0001 source=h1 target=section reason=page-title

\section{\WorkTitle{第九卷，第一一七封书信}}

\Emph{致法兰克王后布伦希尔德}\par

\Person{额我略}致\Person{布伦希尔德}等：\par

治理王国既有赖于勇毅之与正义相配，亦需权力与公平相辅，二者缺一则不足以胜任。陛下对此二端何等珍视，从您治理众民如此堪受赞誉之事实中，已昭然若揭。有鉴于此，谁还能怀疑陛下之良善，或对获得所求之事疑虑不定呢？他若知道您乐意赏赐臣民之物，自然认为请求是合宜的。今呈此函者，即陛下之仆役\Person{希拉里}，他深信我等向陛下进言能助他一臂之力，故恳请我等赐予推荐书信；他确信，若有我们的代祷为之说情，则陛下赐予他人之恩惠，他亦将更丰沛地获得。为此，我们以慈父般之爱向陛下致以问候之词，恳求：如他声称因某些人之不义而遭受逆境，愿陛下之庇护能保卫他；且为避免他可能遭受不公之压迫，望陛下下令保障其安全。如此，若无人得以恣意无理反对，我们亦将因获得所求而感谢——这所求更多是为陛下自身之赏报——而蒙福的宗徒之长\Person{伯多禄}，因您以基督徒之虔诚尊重我等并应允所请，必将回报陛下。\par

\SourceRecord{Translated by James Barmby.；Nicene and Post-Nicene Fathers, Second Series；Vol. 13.；Edited by Philip Schaff and Henry Wace.；Buffalo, NY: Christian Literature Publishing Co.,；1898.；<http://www.newadvent.org/fathers/360209117.htm>.}

% structure: p0001 source=h1 target=section reason=page-title

\section{卷九，书信一二〇}

\Emph{致 \Place{西班牙}的 \Person{克劳狄}}.\par

\Person{额我略}致\Person{克劳狄}等.\par

善行的美名如同香膏芬芳，您的荣耀之馨香从西方之地一直飘散至此。被这气息的甘甜所浸润，我宣称我深深爱了一位素未谋面之人，并在我的心怀内以爱之手将您攫住；我并非在没有认识他之前就爱他，因为我已得知他的善行。因为对于那位以极深的情感为我所知、却以肉身之眼仍属未知的人，我确能如实说：我认识其人，尽管不识其家居。如今这是您美誉的一大明证：您的荣耀据说常与那位杰出的哥特王亲近；因为善人与恶人常不相悦，既然您讨好了那位善人，就能断定您是善的。为此，我以当得的问候致意您，望您在所行之事上勤勉操练，好使\Person{所罗门}的那句真言在您身上应验：\BibleQuote{义人的途径好似黎明的光，越来越亮，直到日中。}\BibleRef{箴 4:18} 因为，如今真理之光已照耀我们，天国的甘甜已向我们的心灵展现，这确实是白昼，却还不是全然的白昼。它将变成全然的白昼，那时灵魂中再没有罪恶之夜。但您要增长直至全然的白昼，以便在天乡显现之前，善工在此不断增多；好使在将来的报应中，赏报的果实因现在劳苦的热忱而更大。因此，我们将我们至爱的儿子、我方隐修院之父\Person{齐里亚库斯}托付给您的荣耀，使他在完成所托之事后，能无阻碍地早日返回。愿全能的天主以祂天臂的保护守护您，并赐您如今在人间荣耀，在漫长岁月之后在天使中亦得荣耀。\par

\SourceRecord{Translated by James Barmby.；Nicene and Post-Nicene Fathers, Second Series；Vol. 13.；Edited by Philip Schaff and Henry Wace.；Buffalo, NY: Christian Literature Publishing Co.,；1898.；<http://www.newadvent.org/fathers/360209120.htm>.}

% structure: p0001 source=h1 target=section reason=page-title

\section{第九卷，第一二一封信}

\Emph{致\Place{伊斯帕利斯}（\Place{塞维利亚}）主教\Person{莱安德}}\par

额我略致西班牙主教\Person{莱安德}。\par

我收到了您的圣座用爱德之笔写来的信。因为舌头转述到纸上的，已从心灵获得了色彩。当这封信被宣读时，有好而明智的人在场，他们立刻感动得心肠澎湃。每个人都开始用爱的手在内心深处拥抱您，因为在那一封信中，人们不是听到，而是看到了您品格的甘饴。每个人都被点燃，人人都赞赏，而听众的火焰本身也显示了说话者的热忱。因为，火炬若不燃烧自身，就不会点燃别人。于是，我们看到了您的心灵燃着多大的爱德，因为它也这样点燃了别人。您的生命，我始终以极大的敬意纪念，他们并不了解；但您心灵的崇高，却从您语言的谦卑中向他们显明了。至于我的生命，您的信称它值得众人效法：但愿那并非如此的事物，因被说成如此而变成如此，免得一个不惯于说谎的人说谎。对此，我简要引用一位好妇人的话：\Emph{\BibleQuote{不要叫我纳敖米，即美丽，却要叫我玛辣，因为我很痛苦。}}\BibleRef{卢 1:20} 因为，好人啊，我今天已不是您从前所认识的那个人。我承认，在外表上我或许在前进，但内里却跌得很深，我担心自己属于那些被记载的人：\Emph{\BibleQuote{当他们被高举时，您使他们跌倒。}}\BibleRef{咏 72:18} 因为，一个在荣誉上上升、在品行上却堕落的人，就是当他被高举时被摔下。因为我，追随我元首的道路，曾决心成为世人的轻蔑和民众的弃绝，并奔跑在那人的道路上，关于他，圣咏作者又说：\Emph{\BibleQuote{他在心中安排了上升，在流泪谷中。}}\BibleRef{咏 83:7}；也就是说，我越谦卑地外在躺在流泪谷中，我内心就越真实地上升。但现在沉重的荣誉大大压垮了我，无数的挂虑喧嚣着我，当我的心灵为天主收敛自己时，它们像利剑一般用攻击将它劈开。我的心没有安息。它匍匐在最低处，被思虑的重量压抑。默观之翼极少或根本不能将它高举。我迟钝的灵魂昏昏欲睡，被世俗的挂虑围绕着吠叫，几乎已陷入麻木，被迫时而处理世俗之事，时而甚至分配肉性之事；不，有时因厌恶之故，被迫带着罪过处理一些事情。我还需要多说吗？被自己的重量压倒，它流汗如同流血。因为，除非罪被称为血，圣咏作者不会说：\Emph{\BibleQuote{求您救我脱离血债。}}\BibleRef{咏 50:16} 但是，当我们罪上加罪时，我们也就应验了另一位先知所说的：\Emph{\BibleQuote{血接触了血。}}\BibleRef{欧 4:2} 因为，当罪加于罪，以致增加不义的负担时，就是血接触了血。在这一切之中，我恳求您因全能的天主，用您祈祷的手扶持我这个陷于纷扰波涛中的人。因为当我度隐修院的宁静生活时，我仿佛顺风顺水航行；但一场暴风骤然卷起狂风巨浪，使我陷入混乱，失去了航行的顺遂；因为在失去平安时，我遭遇了海难。看呀，现在我在波浪中颠簸，寻求您代祷的木板，好让我纵然不配载着完整船只抵达港口，至少能在损失之后藉木板漂到岸边。\par

您圣座提及遭受痛风之苦，我也因持续患病而疲惫不堪。但是，安慰近在咫尺，只要在所受的鞭打中，我们想起自己所犯的过犯；然后我们就会看出这些不是鞭打，而是恩赐，如果我们藉肉身的痛苦来洁净我们为肉身的快乐所犯的罪过。\par

此外，我们已将白羊毛披肩送给您，并附有圣宗徒之长\Person{伯多禄}的祝福，仅供举行弥撒时使用。在送给您时，我理应就您应当如何生活多多劝诫您；但我抑制言语，因为您的生活方式已先于我的言辞。愿全能的天主护佑您，并因多倍的灵魂果实，带您到天国的赏报。至于我，这封短信见证了我是被何等繁重的职务和何等软弱所压垮，在这信中，我对一个我非常挚爱的人说得很少。\par

\SourceRecord{Translated by James Barmby.；Nicene and Post-Nicene Fathers, Second Series；Vol. 13.；Edited by Philip Schaff and Henry Wace.；Buffalo, NY: Christian Literature Publishing Co.,；1898.；<http://www.newadvent.org/fathers/360209121.htm>.}

% structure: p0001 source=h1 target=section reason=page-title

\section{第九卷，第一二二号信函}

\Emph{致西哥特王雷卡雷德}。\par

额我略致雷卡雷德等。\par

最卓越的儿子，我无法用言语表达我对你的工作和生活有多么喜悦。因为听说在我们时代新奇迹的威力，即整个哥特民族通过你的卓越领导，从亚略异端的错误被带到正信的坚定之中，人便不禁与先知同声宣告：\BibleQuote{这是至高者右手所行的改变}（\BibleRef{咏 76:11}）。因为谁的心胸，即使如同顽石，在听到如此伟大的工作时，不会在赞美全能天主和爱慕你的卓越中软化呢？至于我，我宣称，我常常乐于向来到我这里的神子们讲述这些通过你所成就的事，并常与他们一同赞叹。这些事也多半激励我反躬自省：当君王们为天国收益而劳苦收集灵魂时，我却慵懒、无用、在懈怠的安逸中萎靡。那么，在那可怕的审判中，我将对来临的审判者说什么呢？倘若我那时空手而去，而你的卓越却将带领成群的忠信者紧随其后——这些人是你藉着不断而持续的宣讲，如今已吸引到真正信仰的恩宠之中。但善人，因天主的恩赐，这给我带来极大的安慰：那我自己所没有的圣善工作，我在你身上却喜爱。当我为你的作为而欢欣鼓舞时，你劳苦的成果便因爱德而成为我的。因此，关于哥特人的归化，无论是为你的工作，还是为我们的欢欣，我们大可与众天使同声高呼：\BibleQuote{光荣归于至高天主，地上平安归于主所喜悦的人}（\BibleRef{路 2:14}）。因为我认为，我们更应感谢全能天主，尽管我们与你并未同工，却因与你一同喜乐而分享你的工作。此外，蒙福的宗徒之长伯多禄何等欣然地接纳了你的卓越的礼品，你的生活本身已向众人清楚地见证。因为经上记载：\BibleQuote{义人的心愿是祂的喜悦}（\BibleRef{箴 15:8}）。的确，在全能天主的审判中，所看重的不是所献之物，而是献物者是谁。\par

因此，经上记载：\BibleQuote{上主看顾了亚伯尔和他的祭品，却没有看顾加音和他的祭品}（\BibleRef{创 4:4}）。这就是说，祂正要提及上主看顾祭品之前，先谨慎地声明祂看顾了亚伯尔。如此便清楚表明：献祭者并非因祭品而蒙悦纳，反而祭品是因献祭者而蒙悦纳。因此，你显示你的奉献是多么蒙悦纳：你即将献上金子，却先藉着归化臣服于你的民族，献上了灵魂的礼品。\par

至于你告诉我们，那些被派来将你的奉献带给蒙福的宗徒伯多禄的院长们，因海上的风浪而疲惫不堪，未能完成航行便回到了西班牙，你的礼品并未被扣留，因为它们后来到达了我们这里；然而，被派遣者们的恒心受到了考验，看他们是否懂得以神圣的渴望克服途中的危险，即使身体疲惫，也绝不使心灵疲倦。因为出现在善意目标途中的逆境，是德性的考验，而非弃绝的审判。谁能不知道，蒙福的宗徒保禄前来意大利宣讲，尽管途中遭遇海难，却是何等顺利的事？然而，心灵的船在海浪中屹立无恙。\par

此外，我必须告诉你，我从我最钟爱的儿子、司铎普罗比诺的报告中所知，使我更加为你的工作赞美天主：即你的卓越颁布了一项反对犹太人背信的法令后，相关的人试图通过献上金钱来扭曲你正直的心意；但你的卓越对此不屑一顾，力求满足全能天主的审判，宁要纯洁而不要金子。关于此事，我想起了达味王所做的事：当顺从的兵士将渴求的、被敌人围困的白冷水窖的水带来给他时，他说：\BibleQuote{上主决不让我喝义人的血}（\BibleRef{编上 11:19}）。因为他将水倒出，不肯喝，经上记载：\BibleQuote{他将之作为奠祭献给了上主}。如果武装的君王鄙视水，并将其转变为献给天主的祭品，那么，那位为了爱主而拒绝收受（不是水，而是）金子的君王，献给了全能天主何等的祭品呢？因此，最卓越的儿子，我将确信地说：你把那不愿违背祂而收受的金子，作为奠祭献给了上主。这些事是伟大的，归光荣于全能天主。\par

但在这一切之中，我们必须警醒留意，提防那古老仇敌的陷阱，它看见人间的恩赐越大，就越以巧妙的图谋企图夺去它们。因为强盗也不去拦劫路上空手的旅客，而是那些携带金银器皿的人。的确，现世的生命是一条道路。每个人所携带的恩赐越大，就越需要提防埋伏的恶灵。因此，阁下有责任就您在归化属民时所领受的如此巨大的恩赐，全力守护：首先是内心的谦逊，其次是身体的洁净。因为经上记载：\BibleQuote{凡高举自己的，必被贬抑；凡贬抑自己的，必被高举}\BibleRef{路 14:11}，显然，真正爱慕崇高的人，不会使自己的灵魂与谦逊之根分离。因为那恶灵为摧毁它先前无法反对的善行，常在工成之后进入工人的心中，以自夸的静默思绪搅扰它，使受骗的心灵因所做的伟业而自赏。于是，这心灵因隐藏的肿胀而自高，便失去了那赐予恩惠者的恩宠。因此，先知曾向那骄傲的灵魂说：\BibleQuote{你仗恃你的美貌，因你的名声而行淫}\BibleRef{则 16:15}。的确，灵魂仗恃自己的美貌，就是内心自以为义。它因自己的名声而行淫，就是当它行正事时，不愿传扬造物主的赞美，却寻求自己名声的光荣。因此，先知又写道：\BibleQuote{你既然更美丽，就下去吧}\BibleRef{则 32:19}。灵魂因更美丽而下去，是因为它借着美德的风采——本该借此在天主前被高举——却因骄傲从天主的恩宠中坠落。那么，在这种情况下该做什么呢？当恶灵利用我们已做的善事来抬高我们的心时，我们应当时常回忆自己的恶行，以便承认我们所犯的罪是属于我们自己的，而远离罪恶则是全能天主恩赐的礼物。身体的洁净也应在行善的努力中得到守护，因为按照宗徒宣讲者的声音：\BibleQuote{天主的宫殿是圣的，这宫殿就是你们}\BibleRef{格前 3:17}。他又说：\BibleQuote{天主的旨意就是要你们圣化}\BibleRef{得前 4:3}。关于这圣化，他紧接着说明其含义：\BibleQuote{你们务要戒避邪淫，使你们每一个人知道，怎样以圣洁和尊贵持守自己的器皿，不要放纵邪情私欲。}\par

在处理与臣民的关系时，您的王权治理也应当有所节制，免得权力潜入您的心灵。因为王国治理得好，是在于统治的光荣没有主宰其性情。还要谨慎，免得愤怒潜入，以防任何合法之事被仓促施行。因为愤怒，即使是在追究过失者的过错时，也不应当像女主人那样走在心灵之前，而应当像婢女一样跟在理智的后面，待命而前。若愤怒一旦开始占据心灵，它就会把残酷所行视为正义。因此经上记载：\BibleQuote{人的愤怒并不成全天主的正义}\BibleRef{雅 1:20}。又说：\BibleQuote{各人要敏于听，迟于言，迟于怒}\BibleRef{雅 1:19}。然而，我深信在天主的引导下，您已注意这一切。不过，既然现在有了劝诫的机会，我私下里与您的善行联合，使您虽未受劝诫而行的事，现在有了劝诫者，就不致独自而行。愿全能的天主以祂的圣臂护佑您的一切行事，赐您现世顺遂，并在多年之后赐予永恒的喜乐。\par

我们从真福宗徒伯多禄的至圣遗体上，送您一把传递祝福的小钥匙，内含他锁链的铁片，使那曾为殉道束缚他颈项的，能为您解除一切罪过。我们也给了这封信的呈递者一个十字架，内有一些主十字架的木头，以及真福若翰洗者的头发，您可借此永远得到救主的助佑，借着祂前驱的代祷。\par

此外，我们已送给我们至为可敬的弟兄和同僚主教良达里乌斯一件\OriginalTerm{披带}{pallium}，来自真福宗徒伯多禄的宗座，这既出于古老传统，也出于对您的品格及他的善良与庄重之所应得。\par

很久以前，当一位拿波里的青年来到此地时，您至为甘饴的阁下曾嘱托我致信最虔诚的皇帝，请他在档案库中查找从前与可敬记忆的犹斯提尼安王公就贵国王权问题所订立的条约，以便从中推知他对您应持何种态度。但有两件事严重阻碍我这样做。其一是在上述可敬记忆的犹斯提尼安王公时期，档案库曾因一场突然蔓延的火灾而烧毁，以致他那个时期的文件几乎荡然无存。其二是不言而喻的：您应当在您自己的档案中查找对您不利的文件，并将它们呈上，而非由我来做。因此，我恳请阁下根据您的品格妥善安排，并认真推行一切促进和平的事宜，以便您的统治时代能历经多年，以巨大的赞誉而永志不忘。此外，我们另送您一把来自真福宗徒伯多禄至圣遗体的钥匙，当它以应有的尊荣被保存时，愿它祝福并倍增您所享受的一切。\par

\SourceRecord{Translated by James Barmby.；Nicene and Post-Nicene Fathers, Second Series；Vol. 13.；Edited by Philip Schaff and Henry Wace.；Buffalo, NY: Christian Literature Publishing Co.,；1898.；<http://www.newadvent.org/fathers/360209122.htm>.}

% structure: p0001 source=h1 target=section reason=page-title

\section{卷九，第一二三封书信}

\Emph{致\Person{韦南修斯}与\Person{伊塔莉卡}}。\par

额我略致韦南修斯大人、显贵，及其妻伊塔莉卡。\par

我怀着应有的关切，向来西西里的人询问了您贵体的健康。但他们向我报告说您屡次患病，令人忧伤。说到我自己，我也没什么可告诉您的，只是，因我的罪，看着吧，至今已十一个月，我偶尔能起身下床的情况极为罕见。因为我深受痛风之苦，又为诸多烦扰所困，我的生命对我而言是最痛苦的折磨。每天，我在痛苦中昏厥，叹息着期待死亡带来解脱。事实上，在这座城市的神职人员和民众中，热病蔓延得如此之广，几乎没有任何自由人，也没有任何奴隶，还能胜任任何职务或服务。此外，从邻近的城市我们每天都能听到灾祸和死亡的报告。至于非洲如何因死亡和疾病而荒废，我相信您比我们知道得更清楚，因为您离得更近。至于东方，从那里来的人报告了更令人悲伤的荒凉。因此，在这些事中，既然您看到世界末日临近，普遍的打击降临，您不应为自己的困难过于悲伤。相反，作为明智的贵族，应将全部心思转回对灵魂的关怀，越是临近严苛的审判，就越要畏惧。要虔敬，一如经上记载：\BibleQuote{它对现世和来世的生命都有应许} \BibleRef{弟前 4:8}。但是全能的天主既能长久保全您贵体的生命，也能在多年之后带领您进入永恒的喜乐。我恳请代我问候我最亲爱的女儿们，芭芭拉夫人和安东尼娜夫人；我祈愿天上的恩宠保护她们，使她们万事顺遂。\par

\SourceRecord{Translated by James Barmby.；Nicene and Post-Nicene Fathers, Second Series；Vol. 13.；Edited by Philip Schaff and Henry Wace.；Buffalo, NY: Christian Literature Publishing Co.,；1898.；<http://www.newadvent.org/fathers/360209123.htm>.}

% structure: p0001 source=h1 target=section reason=page-title

\section{第九卷，第一二五封信}

致萨洛纳主教马克西穆。\par

额我略致马克西穆等。\par

收到了我们兄弟和同主教马里尼安的信函，而且我们的\OriginalTerm{档案员}{chartularius}卡斯托里也已返回，我们得知，关于那些一度存疑的事，你们的友爱已做了最充分的澄清；我们由衷感谢全能天主，因为从我们内心深处，一切不良猜忌的怨恨已被根除。为此，我本想尽快打发我们的共同之子、你们的执事斯德望回去。但我频繁的病痛迫使我将他留在我身边数日。然而，我一开始稍有起色，就安排他立即欢喜地返回你们那里。\par

因此，我们按照惯例，将举行神圣弥撒礼仪用的\OriginalTerm{披肩}{pallium}送给你们；我们期望你们在各方面维护其意义。因为这件祭衣的尊贵在于谦卑与正义。所以，愿你们的友爱全心全意地努力在顺境中显出谦卑，在逆境中——若它降临——显出正义；对行善者友善，对悖逆者反对；永不使为真理发言的人受窘；按自己的能力在慈悲善工上热切，并且渴望超出自己的能力而热切；同情软弱者，与善意之人同乐；视他人的苦难如己的苦难，为他人的喜乐而喜乐，如同为自己的喜乐；在纠正恶习时严厉，在培养德行时，安抚听者的心灵；在发怒时，保持无怒的明辨，但在平静时，不放弃严肃的监察。至亲爱的弟兄，这就是你们将领受的披肩的意义；如果你们实践它，你们就将内在拥有你们外在所领受的。\par

此外，我完全将我们的兄弟和同主教撒比尼安托付给你们的友爱；如果你们之间有任何争议之事，应暂时搁置。愿爱德在你们之间稳固，这样，即使将来因外在事物发生争执，也可在爱德不离开内心的情况下予以审查。我们也托付我们的共同之子奥诺拉：关于他，若如我们通过我们的\OriginalTerm{档案员}{chartularius}卡斯托里所了解的那样，他使前三任总执事在五年期满后被迫遵守教会惯例而退职，我们确实希望他能体验到你们圣德的仁爱。因为在他自己审判的案件上，不应寻求判决。但若不是这样，那么，应抑制一切内心的膨胀，抛开一切积怨，接纳他，并决不将他从他目前所在的职位上撤除。此外，逃来我们这里的圣职人员梅西安，我们已放心地托付给我们的共同之子执事斯德望照顾，确信在我们亲自送到你们友爱那里的人身上，你们不会显出任何怨恨，反而会施以权威的庇护。愿全能天主在他的保护中守护你们，并赐予我们如此行事，使我们在经历今生世间的波涛后，能欢喜地抵达永恒的事物。\par

\SourceRecord{Translated by James Barmby.；Nicene and Post-Nicene Fathers, Second Series；Vol. 13.；Edited by Philip Schaff and Henry Wace.；Buffalo, NY: Christian Literature Publishing Co.,；1898.；<http://www.newadvent.org/fathers/360209125.htm>.}

% structure: p0001 source=h1 target=section reason=page-title

\section{\WorkTitle{第九卷，第一二七封信}}

由\Person{圣高隆邦}致\Person{教宗额我略}。\par

致神圣的主，并在基督内的父亲，罗马\Emph{教宗}，教会最美丽的装饰，仿佛整个枯萎的欧洲的一朵最灿烂的花，卓越的\OriginalTerm{观察者}{speculator}，似乎享受对纯洁的神圣默观（？）——我，\Person{Bargoma}，基督内可怜的鸽子，致候。\par

愿恩宠与平安，由天主圣父与我们的\Emph{主}耶稣基督赐与你们。神圣的教宗，我乐意认为，您不会觉得被询问关于逾越节的事有何过分，正如那首诗歌所说：\Quote{你当问你的父亲，他必指示你；问你的长者，他们必告诉你。}\BibleRef{申 32:7}因为，虽然我这确实是个\OriginalTerm{琐碎之人}{micrologo}的人，可能会被贴上某位智者的佳评，据说他看到一位\OriginalTerm{涂脂抹粉}{contupictam}的妇人时说：\InnerQuote{我不欣赏这技艺，但我欣赏这额头。}——即我这卑微之人写信给您这位显赫之人；然而，依托我对羞怯福音谦逊的信赖，我冒昧写信给您，并将我忧伤之事托付于您。因为当需要迫使人写信时，即使写给比自己更尊贵的人，写信也不是徒劳的。\par

那么，关于月亮在第二十一或第二十二日过逾越节，您有何看法？——请和平地说——这被许多计算者证明不是逾越节，而是黑暗之时。因为，我相信，您的效能并非不知，\Person{亚纳多略}（一个学问奇妙的人，如圣\Person{热罗尼莫}所说，\Place{凯撒勒雅}主教\Person{欧瑟伯}将他的作品摘录插入其教会史，圣\Person{热罗尼莫}在其目录中称赞了这部关于逾越节的作品）如何强烈反对这个月龄。因为他反对高卢的\OriginalTerm{Rimarii}{Rimarii}（这些人，如他所说，在逾越节上犯了错误），引入了一个可怕的判词，说：\Quote{当然，如果月出延迟到第二更的末尾，那表示午夜，光不战胜黑暗，而是黑暗战胜光；这在逾越节庆典中是绝对不允许的，即让黑暗的任何部分凌驾于光之上，因为主的复活的庆典是光，光与黑暗没有相通。而且，如果月亮直到第三更才出现，毫无疑问月亮是在其第二十一或第二十二日升起的，在这样的月龄中，不可能献上真正的逾越祭品。因为那些认为在这样的月龄中可以庆祝真正的逾越节的人，不仅无法以圣经的权威肯定这一点，而且还犯了亵渎、顽抗和灵魂危险之罪，同时肯定真正的光——它统治一切黑暗——可以在黑暗仍有统治的时候被献上。}此外，在神圣教义书中我们读到：\Quote{逾越节，即主的复活庆典，不能在春分开始过去之前庆祝，也就是说，它不能在春分之前来临。}而\Person{维克托里}在其周期中显然超越了这一规则，并因此在高卢引入了错误，或者说得不太大胆，确认了一个长期存在的错误。因为这两件事中哪一件能合乎理性呢？要么主的复活在其受难之前庆祝（这种想法荒谬），要么法律中主命令的七天（在此期间只允许合法地吃主的逾越节晚餐，这七天要从月亮的第十四日数到第二十日）违反法律和正义而被超越？因为在第二十一或第二十二日的月亮已经脱离光的统治，因为它是在午夜之后升起的；当黑暗战胜光时，据说庆祝光的庆典是不虔敬的。那么，您这位如此智慧的人——您神圣天才的明亮光芒，如同古时一样，散布在世界中——为什么您要庆祝一个黑暗的逾越节呢？我承认，我惊奇，这个高卢的错误，\OriginalTerm{如Schynteneum}{ac si Schynteneum}，为何没有早被您清除？除非我或许猜想——我几乎不敢相信——既然很明显您没有纠正它，它得到了您的赞同。\par

然而，另一种方式可以更体面地为您的专家身份辩护：如果您害怕被贴上\OriginalTerm{Hermagoric}{Hermagoric}新奇的标记，而满足于您前任的权威，尤其是\Person{教宗良}的权威。\par

我恳求你们，在这样的问题上，不要只信赖谦卑或庄重，它们常常受骗；在这难题中，\BibleQuote{\Emph{活的狗}远\Emph{胜于死的狮子}}（\BibleRef{训 9:4}）。因为一位活着的圣人可以纠正先前之人未曾纠正的错误。你们要知道，我们的师长和爱尔兰古人——他们是哲学家和在构建计算中最有智慧的历算家——并不接纳\Person{维克托里}，反而认为他更应被讥笑或原谅，而非具有权威。因此，请以您判断的支持，赐予我这胆怯的陌生人（而非浅学之人）力量，且勿轻看速速赐予我们您宽仁的赞同，以平息我们周围的风暴；因为我已阅读过如此多的作者，却仍不满意那些主教仅说的那句话：\Quote{我们不应与犹太人同守逾越节}。因为这就是\Person{维克托}主教从前说过的话；但东方人或无人接受他的臆想。但这正是\Person{达贡}麻木（\Emph{麻木}？）的脊骨；这正为错误的昏聩所饮入。我请问，这句话有何价值？如此轻浮、粗野，且毫无圣经的见证为依据：\Quote{我们不应与犹太人同守逾越节？}这与问题有何相干？难道该认为那些被弃绝的犹太人现今仍在守逾越节吗？他们既没有圣殿，又在耶路撒冷之外，而且从前被预表的基督已被他们钉死在十字架上。或者，能正确认为逾越节的月第十四日是出于他们自己的规定吗？岂不更应承认是出于天主，因为唯有祂清楚知道，月第十四日被选为出离埃及的通行，具有何等奥秘的意义。或许对于智者及像您这样的人，这比其他人更为清楚。至于那些提出此异议的人，虽然他们无权柄，但若天主不愿我们与犹太人同守逾越节，则让他们指责天主为何不预先在律法中命令他们守九天的无酵饼，以抵制犹太人的悖逆，好使我们庆节的开端不超越他们庆节的终结。因为，如果复活节应在第二十一日或第二十二日庆祝，那么从第十四日到第二十二日将计为九天，即天主所命定的七日，加上人所添加的两日。但是，如果人可凭己意在天主的法令上添加任何东西，我请问，这是否似乎与申命纪的那句话相悖：\BibleQuote{看，他说，我赐予你们的话语，你们不可添加，也不可删减}（\BibleRef{申 4:2}）。\par

但我写这一切，与其说是谦卑，不如说是鲁莽；我知道我已将自己卷入一个充满大困难的骄傲的欧里普斯，或许无力驶出。按我们的地位与阶层，也不适合以讨论的方式向您莫大的权威进言，且我西方人的信件竟可笑地就复活节之事恳求您，您正合法地坐在宗徒及掌钥者\Person{伯多禄}的座位上。但您不应太多考虑我在此事上的卑微，而应考虑许多已故和现今的师长，他们确认了我所指出的，并请想象您正在与他们交谈：因为你们知道，我尽本分张开我的厚唇之口，尽管可能语无伦次、过分夸张。因此，是您来原谅或谴责\Person{维克托里}，因为您知道，若您赞许他，则在您与前述的\Person{热罗尼莫}之间将成为一个信德问题，因为他赞许\Person{亚纳托利}，而亚纳托利与维克托里对立；所以无论谁追随其一，就不能接受另一位。那么，就让您的警醒留意，在赞许上述两位彼此对立的作者中一方的信德时，当您宣明意见时，在您与\Person{热罗尼莫}之间不应有不谐调，免得我们陷入四面困境，不知应同意您还是同意他。在此事上宽免弱者，免得您显露分歧的丑闻。因为我坦诚向您承认，任何反对\Person{圣热罗尼莫}权威的人，在西方教会中将被斥为异端者：因为他们在神圣圣经的一切方面都毫不犹豫地使自己的信德顺应于他。但关于复活节，说到这里就足够了。\par

但我请问，您对那些您曾写为犯有买卖圣职罪的主教有何判断？这些主教被作者\Person{吉尔塔斯}称为害虫。应与他们共融吗？因为在这省份中已知多有这样的人，故而此事更为严重。至于另一些人，他们曾在执事职上受玷污，之后又被选为主教职？因为我们知道有些人在这些事上良心有疑虑；他们与卑微的我商议，希望能确切知道在犯下这些过犯之后是否仍可安然尽牧职；即或因金钱购买职位，或在执事期间犯奸淫。但我指的是与他们的仆从暗中犯奸淫，按我们的师长所教导，此罪不比他罪为轻。\par

关于第三个探究要点，我恳请你，倘若不添麻烦，请答复：对于那些为更亲密地看见天主，或为更成全的生活的热望所激动，而违背其誓愿，离开他们初蒙归化之地，并违反其院长之意志（隐修士的热情迫使他们如此），或得以自由，或逃往旷野的隐修士，应当如何处理？作者\Person{Vennianus}曾就此询问过\Person{Giltas}，后者给了他非常优美的回答：但对一个渴望学问的人而言，更大的畏惧总是不断增长。这些事，以及更多信札之简短所不容纳的事，本可在当面谈话中更谦卑、更清楚地询问，但因身体的软弱和同路朝圣者的照顾将我束缚在家中，尽管我渴望到你那里去，好从那活井的灵性血脉，从那从天涌流、涌向永生的活水知识中汲取。而且，若我的身体能够跟随我的心灵，罗马将再度面临被轻看的危险；因为——正如我们在博学的\Person{热罗尼莫}的叙述中读到，某些人曾从\Place{Heuline coast}的最远边界来到罗马；然后（说来神奇）在罗马之外寻求别的什么——同样，我（对圣人之骨的崇敬除外）也应热切寻求，不是罗马，而是你：因为我承认自己并非有智慧，而是饥渴，若有时间和机会，我定会同样行事。\par

我已阅读你写的包含\WorkTitle{牧灵规则}的书，文辞简短，教导丰富，充满奥秘；并承认对急需者而言，它比蜜更甘甜。所以，我恳请你，赐予我这渴慕你作品的人，关于\BibleBook{厄则克耳}{先知书}的著作，我听说你用奇妙的天才所精心写成的。我已阅读过\Person{热罗尼莫}关于那位先知的六卷书；但他没有诠释中间部分。但是，你若愿意赐恩，请把你其余的一些著作送到城里给我；也就是一卷书的最后讲解，以及（？即）从引文\BibleQuote{我要往没药山和乳香冈去}开始直到结尾的\BibleBook{}{雅歌}，用简短的注释（或出自他人，或出自你自己）加以处理：我恳求你解释\BibleBook{匝加利亚}{先知书}的全部晦涩之处，显明其隐藏的意义，使西方的盲昧为此感谢你。我提出不合理的要求，请求被告知伟大的事：谁能看不到这点呢？但同样真实的是，你拥有伟大的事，并且深知应当运用那从少许和从大量中而来的。愿爱德促使你写信答复；愿我这封信的粗鲁不阻碍你的解释，因为犯错的是我的表达方式，而我心里存着给你当得尊敬的意愿。是我的本分去激发、去询问、去请求：是你的本分不拒绝你所白受的，运用你的才干，按基督所命，将教导之粮给予求者。愿平安归于你和你的家人；请饶恕我的冒昧，有福的教宗，我写得如此大胆；我恳求你在向共同的主的圣洁祈祷中，为我这极卑微的罪人祈祷。我认为向你推荐我的子民是多余的，因为救主判断他们配得接纳，如同行走在祂的名中；并且，若如我从你圣洁的\Person{坎迪杜斯}所听闻，你会倾向于回复说，由古老惯例确定的事不可更改，那么错误显然是古老的；然而谴责它的真理则更为古老。\par

\SourceRecord{Translated by James Barmby.；Nicene and Post-Nicene Fathers, Second Series；Vol. 13.；Edited by Philip Schaff and Henry Wace.；Buffalo, NY: Christian Literature Publishing Co.,；1898.；<http://www.newadvent.org/fathers/360209127.htm>.}

% structure: p0001 source=h1 target=section reason=page-title

\section{第十卷，第十封信}

\Emph{致\Person{罗玛诺}\OriginalTerm{护院}{Defensorem}。}\par

\Person{格里高利}致西西里护院\Person{罗玛诺}。\par

有人向我们报告，我们最可敬的兄弟\Person{巴息略}主教好像平民百姓一样忙于诉讼，无益地出入法庭。既然此事既使他本人卑贱，又损害了司铎应有的尊敬，愿你的经验一接到此命令，立即以严格执行的方式迫使他返回职责，使他在你的坚持之下，不得以任何借口延迟五天；免得你若以任何方式允许他如此延迟，你与他一起在我们面前成为严重有罪之人。颁布于十二月，小纪第三。\par

\SourceRecord{Translated by James Barmby.；Nicene and Post-Nicene Fathers, Second Series；Vol. 13.；Edited by Philip Schaff and Henry Wace.；Buffalo, NY: Christian Literature Publishing Co.,；1898.；<http://www.newadvent.org/fathers/360210010.htm>.}

% structure: p0001 source=h1 target=section reason=page-title

\section{第十卷 第十五封信}

\Emph{致\Person{克莱孟蒂娜}，贵族}。\par

\Person{额我略}致\Person{克莱孟蒂娜}，等等。\par

据某位院长的报告，我们得知有人向阁下进谗言，说我们对您心存芥蒂。果真如此，则捏造此事之人皆以诚恳为表，实则对您两面三刀，既显己忠，又恶毒地使您怀疑我们。然而，我荣耀的女儿，我素知您的善行，尤其自幼相伴的贞洁，一向对您满怀敬重与关爱。但即便现在，以免阁下怀疑我心意有变，我声明：我对您绝无丝毫恶感或怒气；请确信，我对您怀有慈父之情。唯有一事，既已有人告知，便不应缄默，以免因隐讳当说之话而致爱德衰减。\par

实因有人向我报告，说每当有人冒犯您，您便念念不忘，不肯释怀。若果真如此，则我愈爱您，愈感伤痛，恳请您高贵地摒弃此过，莫让仇敌的种子在您善行的田地中生长。愿您记起天主经的话，切莫让责备胜过宽恕。愿阁下的良善战胜过犯，借着有益的宽恕，使冒犯者归附于您，胜过持续的严厉使其背离。当留给他足以使之羞愧的余地，而非存留使之忧伤的事。因为明智的宽恕往往比严苛的报复更能纠正过错；有时前者使人更为忠信顺从，后者却令人执拗怨毒。我们如此对您说，并非要您削弱对义德的热忱，而是使您在小事上，如同在大事上一般。若过犯的性质需要严厉对待，则当如此处理：使惩罚纠正过错，而后对被纠正者不拒绝恩宠。既然我们出于慈父之情，为您的灵魂益处而劝诫，愿您以我们所怀的爱德接纳这些话，并为了您的尊荣而取用，使您的善行在人前更为显明，在天主前更为纯洁。最亲爱的女儿，您尽可信赖我们，凡事皆然；因我们常愿闻您的平安，请多来信，使我们欣慰。\par

\SourceRecord{Translated by James Barmby.；Nicene and Post-Nicene Fathers, Second Series；Vol. 13.；Edited by Philip Schaff and Henry Wace.；Buffalo, NY: Christian Literature Publishing Co.,；1898.；<http://www.newadvent.org/fathers/360210015.htm>.}

% structure: p0001 source=h1 target=section reason=page-title

\section{第十卷，第十八封书信}

\Emph{致贵族女克莱孟蒂娜}。\par

额我略致克莱孟蒂娜等\par

荣耀的女儿，你要知道，司铎阿曼多已被苏伦托的民众选为主教。我们已写信命他前来，你不应为他的离开而忧伤，因为那在心中与你同在的人，不应被认为离开了你。何况，他既曾令你喜悦，如今又为需要主教的人所接纳，为此你当赞颂全能的天主，并以基督徒的虔诚更加喜乐；并要欣然尽力促成他尽快前来我们这里，以利于他人，因为真诚爱德的本分，就在于所爱之人被召叫而得以成长时，引以为荣。\par

\SourceRecord{Translated by James Barmby.；Nicene and Post-Nicene Fathers, Second Series；Vol. 13.；Edited by Philip Schaff and Henry Wace.；Buffalo, NY: Christian Literature Publishing Co.,；1898.；<http://www.newadvent.org/fathers/360210018.htm>.}

% structure: p0001 source=h1 target=section reason=page-title

\section{第十卷，第十九封信}

\Emph{致副执事安提缪斯。}\par

额我略致坎帕尼亚副执事安提缪斯。\par

在一位曾被选为索伦托城主教的人被我们视为不适任之后，他们便选出了阿曼杜斯，他原是卢库卢斯营地内圣塞味利祈祷所的司铎。因此，我们命令你的阁下，放下各种借口，务必尽快将上述司铎送到我们这里，以便若没有妨碍他前来的事，求愿者的愿望能因基督的帮助得以实现。至于他的生活和行为，既然在他长期生活的地方可以更加了解，请你同我们的弟兄和同道主教福图纳图斯一起仔细查问。如果他的晋升至神圣品秩没有障碍，就应该毫不迟延地将他送到我们这里。但是，为避免我们尊贵的女儿克莱门蒂娜对此不满，你的阁下应当去见她，并在获得她同意的情况下办理此事。然而，如果她拒绝，你的阁下仍不要迟延地将他送到这里，因为我们既应安抚我们子女的心，也不应阻碍灵魂的益处。\par

\SourceRecord{Translated by James Barmby.；Nicene and Post-Nicene Fathers, Second Series；Vol. 13.；Edited by Philip Schaff and Henry Wace.；Buffalo, NY: Christian Literature Publishing Co.,；1898.；<http://www.newadvent.org/fathers/360210019.htm>.}

% structure: p0001 source=h1 target=section reason=page-title

\section{第十卷 第二十三封信}

\Emph{致\Place{西西里}公证人\Person{阿德里安}}。\par

\Signature{\Person{额我略}致\Person{阿德里安}等}\par

一件令我们全然憎恶与不名誉的事传到了我们耳中，我们感到奇怪，若此事属实，你为何没有注意。因为\Person{圣维图斯}隐修院（位于\Place{埃特纳山}上）的一位隐修士\Person{玛尔提安}来到我们这里，呈上一份请愿书，其中控诉许多事，包括此隐修院的隐修士们生活如此乖谬邪恶，竟敢让妇女与他们同住——这是提及起来都令人发指的事。有鉴于此，我们已就此事写信给我们的弟兄与同为主教的\Person{良}，使他查明真相后，若果真如此，务必以最严厉的惩戒予以纠正。因此，你的阁下也必须在一切方面切心于查明真相，并惩罚如此重大的恶行，以免发现任何懈怠或疏忽之处。此外，就同一隐修院其他方面的利益，你应在合乎正义的范围内提供协助，以便若如所言有任何侵犯之事，得以依法纠正，并且今后绝不致发生任何有悖于敬畏天主与法律秩序之事。\par

\SourceRecord{Translated by James Barmby.；Nicene and Post-Nicene Fathers, Second Series；Vol. 13.；Edited by Philip Schaff and Henry Wace.；Buffalo, NY: Christian Literature Publishing Co.,；1898.；<http://www.newadvent.org/fathers/360210023.htm>.}

% structure: p0001 source=h1 target=section reason=page-title

\section{第十卷·第二十四书}

\Emph{致\Person{福尔图纳图斯}\Place{尼阿波利斯}（那不勒斯）主教。}\par

\Person{格里高利}致\Person{福尔图纳图斯}等。\par

当您的弟兄团对您管辖下的隐修院关注不足时，您既使自己受责备，又让我们为您的松懈感到遗憾。如今我们听闻，一位名叫\Person{毛里基乌斯}的人，近日在\Place{巴尔巴奇亚努斯隐修院}成为隐修士，现已带领其他隐修士逃离该院。此事中，上述\Person{巴尔巴奇亚努斯}的轻率在我们看来使他极其有罪，因为他竟未经事先考验就鲁莽地为一位世俗者剃发。我们不是曾写信给您，要您先考验他，若合适，再任命他为院长吗？即使现在，也要好好看管您所选的人。因为如果他已经开始如此行为，以至于显明自己不适合管理弟兄，那么您就因他的过失而有过失。\par

此外，请您的弟兄团更严格地禁止所有隐修院，在领受者于隐修生活中完成两年之前，绝不可为那些可能为隐修誓愿而领受的人剃发。但在这段时间内，应仔细考验他们的生活和品行，以免其中任何人要么不满足于自己所渴望的，要么不坚守自己所选择的。因为未经考验的人服从任何长上已是严重之事，更何况未经考验者应被奉献于天主的服务，那又该何等严重呢？\par

再者，若有士兵愿意成为隐修士，任何人在任何情况下都不得未经我们的同意或未向我们报告之前擅自接纳他。若不切实遵守此规，要知道所有属下的罪责都归于您自己，因为事实本身显明您对他们过于疏于关心。\par

\SourceRecord{Translated by James Barmby.；Nicene and Post-Nicene Fathers, Second Series；Vol. 13.；Edited by Philip Schaff and Henry Wace.；Buffalo, NY: Christian Literature Publishing Co.,；1898.；<http://www.newadvent.org/fathers/360210024.htm>.}

% structure: p0001 source=h1 target=section reason=page-title

\section{卷十·书信三十一}

\Emph{致\Person{利贝蒂努斯}，前任裁判官。}\par

\Person{额我略}致\Person{利贝蒂努斯}等。\par

您在此世财物方面的困境，我们并非不知。但既然对于那些处于极度磨难中的人而言，唯一的安慰便是造物主的仁慈，请将您的希望寄托于祂，并以全心转向祂——祂既公义地容许自己所愿的人受苦，也必仁慈地拯救那信赖祂的人。因此，请向祂感恩，并耐心忍受临于您的一切。因为正心之作为，不仅在于顺境时讚美天主，也在于逆境中同声颂扬祂。所以，在您所受的这些事上，切莫让怨天尤人之心潜入您内，因为我们不知道造物主如此行事的目的何在。尊贵的儿子，或许您曾在顺境时得罪了祂什么，祂便愿以慈爱的苦楚来洁净您。因此，不要让暂时的磨难压垮您，也不要让财产的损失令您分心，因为若您在逆境中以感恩之心和忍耐使天主垂顾于您，则失去的必将倍增，此外，永生的喜乐也已为您预备。我恳求您，勿怪我们委托监护官\Person{罗马努斯}下令从我们处拨付二十套衣服给您仆役；须知，即便微薄之物，若出自蒙福的宗徒\Person{伯多禄}的财物，亦当视为莫大的祝福，因为他既有能力赐予您更大的事物，也必能在全能的天主前为您求得永恒的恩惠。六月，第三财税周期。\par

\SourceRecord{Translated by James Barmby.；Nicene and Post-Nicene Fathers, Second Series；Vol. 13.；Edited by Philip Schaff and Henry Wace.；Buffalo, NY: Christian Literature Publishing Co.,；1898.；<http://www.newadvent.org/fathers/360210031.htm>.}

% structure: p0001 source=h1 target=section reason=page-title

\section{\WorkTitle{第十卷\newline{}书信三十五}}

\Emph{致亚历山大宗主教\Person{Eulogius}。}\par

\Person{Gregory} 致 \Person{Eulogius} 等。\par

在去年，我收到了您极其甘饴的圣座的书信；但因我病势极重，至今未能回复。因为，看，我卧病在床已将近两年，为痛风所苦，几乎不能在庆节上起身三个时辰以举行弥撒。且我很快就被剧痛压迫而躺下，以便在间歇的呻吟中忍受折磨。我的疼痛有时缓和，有时剧烈：但既不缓和到离去，也不剧烈到致死。因此，我虽日日濒死，却日日不得死。而我，身为可悲的罪人，长久被囚禁于如此腐朽的牢狱，这并不奇怪。我被迫呼喊：\BibleQuote{求祢领我的灵魂出离监牢，好让我赞美祢的圣名}\BibleRef{咏 141:8}。然而，因我尚未堪当通过祈祷获得此事，我恳求您的圣座的祈祷能以其转求助我，并将我自罪恶与腐朽的重负中释放，进入您所熟知的那光荣的自由，即天主儿女的自由。\par

您那对我至为甘饴、永当敬仰的福分，在您的信中告知我：我们的共同儿子、\Place{君士坦丁堡}城执事\Person{Anatolius}曾写信给您，说某些来自\Place{耶路撒冷}周边的隐修士来到我这里，就\OriginalTerm{阿格诺伊特派}{Agnoitæ}的错谬进行咨询；您说他恳求您的圣座写信给我，以表达您对此咨询的意见。但是，既没有隐修士从\Place{耶路撒冷}周边来我这里咨询，我也不认为我们的那位共同儿子在信中告诉了您不实之事；但我怀疑是翻译误解了他书信的意思。因为同一位执事，在两年多以前写信给我，说来自前述地区的隐修士曾到\Place{君士坦丁堡}城进行如此咨询，并想问我我的看法。我早在收到您的信之前很久，就对同一异端作了与我在您圣座的书信之后所发现的完全相同的答复：我向全能天主致以极大的感谢，因为关于所有问题，\Place{罗马}和\Place{希腊}的教父们——我们是他们的追随者——都以同一圣神发言。因为我在多处发现，您这封书信仿佛就是我阅读拉丁教父们反对前述异端的著作。请您想想，我多么必须爱戴并赞扬我最神圣的兄弟的卓越，我在他的口中认出了我所深爱的可敬教父们。因此，愿赞美归于祂，愿至高荣耀归于祂；祂的恩赐使\Person{马可}的声音仍在\Person{伯多禄}的宝座上高呼；祂的神倾注，当司祭进入至圣所探究奥迹时，属灵的钟声便因宣讲的话语，在圣教会中如同在会幕中响起。因此，您的宣讲是正确且极应赞扬的。但我们恳求全能之主长久保全您在此世的生命，以致从您这天主之乐器中所发出的真理之声，能在此世传扬得更广。并请为我转求，使这为我已变得过于崎岖的朝圣之路能迅速完成，好让我这无法凭自己功绩的人，能凭您的功绩获得永生之国的许诺，并与天上的子民一同喜乐。\par

\SourceRecord{Translated by James Barmby.；Nicene and Post-Nicene Fathers, Second Series；Vol. 13.；Edited by Philip Schaff and Henry Wace.；Buffalo, NY: Christian Literature Publishing Co.,；1898.；<http://www.newadvent.org/fathers/360210035.htm>.}

% structure: p0001 source=h1 target=section reason=page-title

\section{卷十·第三十六书}

\Emph{致\Place{萨洛纳}主教\Person{玛克西穆}}.\par

\Person{额我略}致\Person{玛克西穆}等.\par

当我们的共同的孩子司铎\Person{维特拉努斯}来到罗马城时，发现我因痛风疼痛而虚弱不堪，以致无法亲自回复贤弟的信函。确实，关于使贤弟处于巨大危险的\OriginalTerm{斯拉夫人}{Sclaves}民族，我深感痛苦和不安。我痛苦，是因为我已因你的痛苦而受苦；我不安，是因为他们已开始经由\Place{伊斯特拉}进入意大利。此外，关于书记官\Person{尤利安}，我能说什么呢？因为我看到处处如何因我们的罪而受到惩罚，以致我们被外邦民族和内部法官所困扰。但不要因此感到丝毫忧伤，因为在我们之后生活的人将看到更糟糕的时代，以至于他们会认为与自身相比，我们度过了幸福的日子。但是，只要贤弟有能力，你应当为穷人、为受压迫者挺身而出。即使你无能为力，全能天主所赏赐的虔诚之心也已足够让他悦纳。因为经上记载：\BibleQuote{解救被拉去赴死的人，不要吝于拯救将被杀害的人。}\BibleRef{箴 24:11}你若说，我的能力不足，那洞察人心的自会明白。因此，在你所做的一切中，要渴望使那洞察人心的悦纳你。但凡能悦纳他的事，不要不做。因为人的恐惧与恩惠如同烟雾，被轻风一吹而散。确知此事：无人能同时取悦天主和恶人。因此，贤弟当以为，你越是自知得罪了悖逆之人，就越是取悦了全能天主。然而，你对穷人的辩护应温和庄重，免得若行事过于严苛，人以为你是出于青年的骄傲。但我们为穷人的辩护必须如此，使谦卑者得蒙护佑，压迫者不易找到可凭恶意指责之处。那么，请留意厄则克耳所说的话：\BibleQuote{人子，不信者和毁灭者与你同在，你住在蝎子中间。}\BibleRef{则 2:6}真福约伯说：\BibleQuote{我成了龙的兄弟，鸮鸟的同伴。}\BibleRef{约 30:29}保禄对自己的门徒说：\BibleQuote{在弯曲悖逆的民族中，你们像世上的光闪耀。}\BibleRef{斐 2:15}因此，我们更当谨慎行事，因为我们知道我们生活在天主的仇敌中间。此外，关于\OriginalTerm{佛提诺派}{Photinianists}，贤弟当极其留意；并如你已开始的，设法将他们召回圣教会的怀抱。但若有人愿来见我，接受解释，须先宣誓，即使得到解释，他们也不允许其追随者固执于错误。然后请阁下向他们承诺，他们不会从我这里受到任何不公正，我将给他们解释。他们若承认真理，就接受；若不肯承认，我也将毫发无损地打发他们走。但若他们中有人欲来我这里反对你，贤弟切勿阻挡；因为他们来了，要么接受解释，要么确实再也不能看见那片土地了。\par

\SourceRecord{Translated by James Barmby.；Nicene and Post-Nicene Fathers, Second Series；Vol. 13.；Edited by Philip Schaff and Henry Wace.；Buffalo, NY: Christian Literature Publishing Co.,；1898.；<http://www.newadvent.org/fathers/360210036.htm>.}

% structure: p0001 source=h1 target=section reason=page-title

\section{第十卷，第三十七函}

\Emph{致阿非利加总督\Person{依诺森}。}\par

\Person{额我略}致\Person{依诺森}等。\par

阁下之清辩，以心中之蜜调和，其味渗入我等肺腑，令我辈爱慕不已；是以尔所书者甘甜，尔所行者馨香；此非无故，盖因德行深造者，在公正而非偏私之眼中为伟大也。再者，我等闻尔已就总督之职，喜中带忧。盖一则为我等极亲爱之子之升迁而欢喜，二则因我等从自身之苦楚实感在混乱时世中升居高位是何等重担。故当尽一切心力，使烦扰之境遇成为赏报之机。诚如尔知，谷粒生于长满荆棘之地，玫瑰出于荆棘。尔既有适播种之时，切勿迟延，播下善工之种，俾于收获之日，可携更大捆之喜乐回家，并因在短暂之尊位上的善工，得臻永恒之光荣。我等既知尔为备造快船所费心力，今以佳音释尔焦虑：承天主仁慈，我等已与伦巴第人的国王议和，直至即将到来之第四小纪年的三月为止。然此和约能否持久，我等不得而知，因闻该王已去世，但至今尚不确定。\par

关于\Person{阿纳蒙达鲁斯}之事，尔来信所请，我等已照办；但愿结果能如我等所愿；盖就我等而言，我等未曾拒绝向受苦者施以代祷之助。\par

至于尔欲我等将\WorkTitle{约伯记释义}一书寄予尔，我等甚喜尔之热忱；因我等见阁下切望获得那能使尔不完全离己、并能使尔之心在分心于世俗之忧后回归自身之物。然若尔欲以美味饱飨，请读尔同乡\Person{真福奥斯定}之著作，切勿以我等之糠秕与彼之精麦相比。\par

此外，我等从我等之\Emph{档案员}\Person{希拉里乌斯}的证言得知，阁下为了爱尔之宗徒之长\Person{圣伯多禄}的穷人，赐予何等眷顾与慈惠。为此，我等致以丰厚谢忱，恳求全能天主的仁慈，愿祂以恩宠护卫尔，不容恶人从外在胜过尔，亦不容恶灵在内；惟愿祂仁慈地在敬畏中安排尔之行事，使祂既使尔在世上有荣光，亦使尔在漫长生命之后，列于其诸圣之数中。\par

\SourceRecord{Translated by James Barmby.；Nicene and Post-Nicene Fathers, Second Series；Vol. 13.；Edited by Philip Schaff and Henry Wace.；Buffalo, NY: Christian Literature Publishing Co.,；1898.；<http://www.newadvent.org/fathers/360210037.htm>.}

% structure: p0001 source=h1 target=section reason=page-title

\section{第十卷·第三十九书}

\Emph{致\Place{亚历山大里亚}宗主教\Person{欧洛吉}}\par

\Person{额我略}致\Person{欧洛吉}等\par

\Emph{\BibleQuote{如凉水润渴魂，远地来的好消息也是如此}} \BibleRef{箴 25:25} 然而，就神圣教会之益处而言，还有什么消息能如我至甜蜜之圣座您安康的消息那样美好呢？您因对真理之光的领悟，既以宣讲的圣言光照同一教会，又以您的风范榜样塑造她走向更美的道路。每当我心中追忆您与我同心合一，并感到我常驻您心中时，我就感谢全能天主，因为爱德不能被地域距离所分割。因为，虽然我们在身体上远隔，但在灵魂上却不可分离。\par

我们的共同儿子辅祭\Person{阿纳托利}在来信中告知我，在帝国之城内，没有任何教会事务曾因世俗原因而受扰。但我相信他此前已向我通报，您的真福在教会事业中如何发言；我欣然认为，凡您适逢其会之处，我都不觉得缺少了我。因为我知道，您作为真理的仆役、\Person{伯多禄}的追随者、神圣教会的宣讲者，必会说出通过来自宗徒\Person{伯多禄}之座的教师之口应当听到的话。\par

此外，在这些日子之前，当\Place{亚历山大里亚}的\Person{阿布拉米}前来见我时，我已回复您的圣座，说明了我对您针对\OriginalTerm{无知派}{Agnoitæ}异端所发表之著作的看法，以及我何以迟迟未作回复。但据说，上述\Person{阿布拉米}因航行困难，在\Place{那不勒斯}城逗留已久；因此我再次来信，内容与先前所写一致，因为您反对所谓无知派的教导中有许多令我们钦佩之处，而无任何令我们不悦之处。我也已就相同意思向我们的儿子辅祭\Person{阿纳托利}详细写信。此外，您的教义在各方面均与拉丁教父如此一致，以至于我毫不惊讶地发现，在不同语言中，圣神并未不同。\par

因为，关于你所说的无花果树，\Person{奥斯定}在同一意义上说得恰当；因为当福音作者接着写道\BibleQuote{因为还不是无花果的时候}\BibleRef{谷 11:13}时，这清楚地表明主所寻找的无花果是会堂里的果实，它们有律法的叶子，却没有行为的果实。因为万物的创造者不可能不知道那棵无花果树没有果实；这是所有人都可能知道的事，因为还不是无花果的季节。但是关于经文所记：\BibleQuote{至于那日子和时辰，圣子和天使都不知道}\BibleRef{谷 13:32}，陛下您正确地察觉到，这肯定不是指作为元首的圣子，而是指祂的奥体，即我们。关于此事，同一位真福\Person{奥斯定}在许多地方采用了这种理解（\WorkTitle{论八十三个不同的问题}第60问；\WorkTitle{论天主圣三}卷一第12章；\WorkTitle{圣咏集}第六篇开篇；\WorkTitle{圣咏集}第三十四篇第二篇讲道）。他也提到另一件可以理解为同一圣子的事，即全能的天主有时以人的方式说话，就如祂对亚巴郎说：\BibleQuote{现在我确知你是敬畏天主的了}\BibleRef{创 22:12}。这并不是说天主那时才知道祂被敬畏，而是说祂那时让亚巴郎知道他是敬畏天主的。因为，就如我们常说快乐的一天，并非指那一天本身快乐，而是指它使我们快乐，同样，全能的圣子说祂不知道那日子，是祂使那日子不被知道；并非祂自己不知道，而是祂不允许它被知道。因此也只有圣父被说成知道那日子，因为与祂同性同体的圣子，祂的知识来自祂的神性本性，天使所不知道的，祂凭此高于天使。因此也可以更细致地这样理解：\OriginalTerm{独生子}{Only-begotten}降生成人，为我们成为完人，祂确实在祂的人性本性中知道审判的日子和时辰，但祂知道这并不来自于祂的人性本性。祂在其中所知道的，却并非从其中知道，因为天主成为人，是藉着祂天主性的能力知道审判的日子和时辰：正如在婚宴上，当童贞母亲说缺酒时，祂回答：\BibleQuote{女人，这关我和你有什么关系？我的时辰尚未来到}\BibleRef{若 2:4}。因为天使的主宰并不受时辰的支配，祂在祂所创造的一切事物中创造了时辰和时期；但因为童贞母亲在缺酒时希望祂行一个奇迹，祂便立即回答她：\Quote{女人，这关我和你有什么关系？}仿佛明白地说：我能行奇迹来自我的父，而非我的母亲。因为祂凭着父的本性行奇迹，而祂之所以能死，则是来自母亲。因此，当祂在十字架上快要死时，祂认出了祂的母亲，并将她托付给门徒，说：\BibleQuote{看，你的母亲}\BibleRef{若 19:27}。祂于是说：\BibleQuote{女人，这关我和你有什么关系？我的时辰尚未来到。}这就是说：\Quote{在奇迹中，我并没有从你的本性得到什么，所以我不认你。当死亡的时辰来到时，我会认你为我的母亲，因为我从你那里获得了死的能力。}因此，那种知识，祂并没有从人性本性中获得，在那本性中祂与天使同为受造物，祂否认祂与那些作为受造物的天使拥有这知识。所以，审判的日子和时辰，祂作为天主也是人知道，但正因如此，因为天主就是人。此外，同样明显的是，凡不是\OriginalTerm{聂斯多略派}{Nestorian}的人，绝不可能成为\OriginalTerm{阿格诺伊特}{Agnoite}。因为一个承认\OriginalTerm{天主的智慧}{Wisdom of God}亲自成为血肉的人，怎能说天主的智慧有任何不知道的事呢？经上记着：\BibleQuote{在起初已有圣言，圣言与天主同在，圣言就是天主。万物是藉着祂而造的}\BibleRef{若 1:1}。如果万物，那么无疑也包括审判的日子和时辰。那么谁还会如此无理智地妄称\OriginalTerm{圣父的圣言}{Word of the Father}创造了祂所不知道的事物呢？经上也记着：\BibleQuote{耶稣知道父已把一切交在祂手中}\BibleRef{若 13:3}。如果一切，那么当然包括审判的日子和时辰。那么谁还会如此愚蠢地说圣子把祂所不知道的事物接在手中呢？\par

但是，关于祂对有关拉匝禄的妇女们所说的话：\BibleQuote{你们把他放在哪里了？}\BibleRef{若 11:34}，我的感觉与您完全相同：如果他们说主不知道拉匝禄被埋在哪里，因而询问，那么他们无疑将被迫承认主不知道亚当和厄娃在犯罪后躲藏的地方，当时祂在乐园中说：\BibleQuote{亚当，你在哪里？}\BibleRef{创 3:9}？或者当祂责斥加音说：\BibleQuote{你的弟弟亚伯在哪里？}\BibleRef{创 4:9}？但如果祂不知道，为什么祂立刻又说：\BibleQuote{你弟弟的血从地上向我喊冤？}然而，关于这段经文，\Person{塞维里安·加巴兰}有不同的见解，他说主这样对妇女们说，如同在斥责，因为祂询问她们把死去的拉匝禄放在哪里；仿佛明确地指涉厄娃的罪，祂是说：我把人安置在乐园，而你们却把他放在坟墓里。\par

但是，我们的上述共同之子\Person{执事阿纳托利}通过提出另一个问题对此作了答复：——如果有人反驳我说，就如那不死者屈尊死去以拯救我们脱离死亡，那在万世之前永恒者甘愿受时间的支配，同样，\OriginalTerm{天主的智慧}{Wisdom of God}也屈尊承担了我们的无知，以拯救我们脱离无知，那又如何？但我迄今尚未给他任何答复，因为一直受重病困扰。现在，因着您的祈祷，我已开始康复；如果我恢复得能够口述，靠主的助佑，我将答复他。至于您，我不必对您说任何有关此事的话，以免我似乎教导您已经知道的事，因为药物若用在健康强健的肢体上，也会失去治疗效果。\par

此外，我们告知你们，在此地我们因缺乏优秀的翻译者而遭受严重困难。因为没有人能表达原意，而所有人都试图逐字逐句地翻译：因此他们混淆了所述的全部意思。由此导致的结果是，我们若不付出艰巨的努力，便完全无法理解所译的内容。\par

我已收到\Person{圣马尔谷}福音作者和您的真福的祝福。我一直希望送您一些木材；但前来的船只太小，无法装载。然而，就连\Place{亚历山大}来人见到的那根也是小尺寸的。因为我为您准备了一些更大的，尚未运往\Place{罗马城}：我等待\Place{亚历山大}船到达时运送；它仍留在砍伐之地。\par

愿全能的天主长久保佑您的生命，以建树圣教会，并感动您为我热切祈祷；使我虽被自己的重罪所压，却能因您的祈祷而在全能天主前被高举。\par

\SourceRecord{Translated by James Barmby.；Nicene and Post-Nicene Fathers, Second Series；Vol. 13.；Edited by Philip Schaff and Henry Wace.；Buffalo, NY: Christian Literature Publishing Co.,；1898.；<http://www.newadvent.org/fathers/360210039.htm>.}

% structure: p0001 source=h1 target=section reason=page-title

\section{第十卷，第四十二书}

致德撒洛尼总主教欧瑟伯\par

额我略致欧瑟伯等\par

最亲爱的弟兄，若我们细心默想和平何其崇高，就会认识到我们应以何等热忱去培育它。因为我们的主、救主将和平作为巨大的恩赐惠然留给祂的门徒，遂使那因坚固的信德而与祂合一的人，在爱中与祂共享共融。经上记载：\BibleQuote{缔造和平的人是有福的，因为他们要称为天主的子女。}（\BibleRef{玛 5:9}）因此，凡渴望成为天父继承人的，就应藉持守和平，不拒绝做祂的子女。那给不和让位的人，势必使自己在如此巨大的恩赐中无分。既然藉天主的仁慈，你的信德之纯全已以公教正统向我们显明，正如所应许的，我们便大为惊异：你竟容许那些你知道信仰纯正、思想正确的人，因某些人的过失而无端受困扰，以致你兄弟之名誉因他人的罪过而蒙上阴影。因为一个容忍犯错之人的，怎能避免被怀疑有错误？若不借着公开的补赎来清除信德的热忱所要求清除的，他又能期待别人如何评价自己？\par

确实，据说你的司铎路加和伯多禄拒绝接受加采东大公会议，因此你的正统子女之心被不小的绊脚石搅动。既然他们的热诚不仅应受赞扬，更应完全珍视，我们劝勉你兄弟之关怀，要以一切积极与审慎行动，毫不犹豫地调查此事。若那些人被查明与那乖谬无关，就应借满足你们子女的心，除去他们心中的绊脚石，并在一切异端中，尤其诅咒塞维鲁和聂斯托利，好使净化在那些因爱信德而对那些异端产生猜疑的人中产生爱德；并使纯洁而单一的公教真理宣认所联合的人，在和谐之感中健康地结合。不要以为怀疑者不配得到满足，因为我们受天主之声的教导：\BibleQuote{这些小子中的一个，你们不要轻看。}（\BibleRef{玛 28:18}）因此，谁不希望那教导我们的人被轻视，就不要拒绝教导者的言语；因为我们的救主曾为祂作证，说祂是祂自己所选的器皿的那一位，也劝勉我们要\Quote{以和平的连系，保持心神的合一}（\EditorialNote{厄弗所书第四章}）。因此，凡不愿被这救恩的连系所束缚的，就应追求和平之事，不给仇敌留地步；好使他在弟兄的激烈纷争中得以进展后，当合一建立时，更坚固地被践踏。\par

然而，若他们（我们并不期待如此）被发现受了这错误之箭的伤，就应以教会劝言的治疗施于他们，好使他们或治愈后仍留在主的羊群中，或从教会身体的合一中被切除；如此，从微小的损失中可获大利，部分的除去可使全身自由。因为审慎的牧人也关心：不拖延地将一只不能治愈的病羊从他健康羊群中赶出，免得它以疾病之污沾染其他羊，他知道唯有赶出这一只才能保全其余的健康。因此，我再次以兄弟之爱德劝告你，要以最大警觉调查此事，并以最大谨慎遵守我们所写的，免得因与其他人交往而使你所持的正确信德变得可疑。因为不纠正应切除之事的人，就是认可那些事。因此，你必须以极大关怀和万全准备深思熟虑：切勿使那些人之身份成为他人之绊脚石，或使舆论对你有害；这样，牧人从托付给你的羊群所得的收益将越发增长，因为真诚的爱德与经认可的关怀使你对看守它们殷切关心。\par

\SourceRecord{Translated by James Barmby.；Nicene and Post-Nicene Fathers, Second Series；Vol. 13.；Edited by Philip Schaff and Henry Wace.；Buffalo, NY: Christian Literature Publishing Co.,；1898.；<http://www.newadvent.org/fathers/360210042.htm>.}

% structure: p0001 source=h1 target=section reason=page-title

\section{第十卷，第六十二封书信}

\Emph{致那不勒斯人。}\par

\Person{额我略}致\Place{那不勒斯}的圣职人员及贵族市民。\par

主教选举中，民众的票数分为两派，这不是新鲜事，也不应受谴责。但严重的是，在此类情况下，选举不是凭判断，而只是凭偏爱。因为在你们的信函到达我们之前，我们已从某些人的报告中得知，被另一派选出的执事若望有一个小女儿。因此，如果他们有意讲求理性，那么别人既不会选他，他也不会同意。他的妄自尊大该是何等严重，他竟然在有小女儿的证据下，证明他长久未能节制自己的身体，还敢觊觎主教职位！此外，执事伯多禄，你们说被你们选出，据说他毫无心机。你们知道，在当今之世，被立于统治理高位的人，应当懂得不仅要谨慎照顾灵魂的得救，还要顾及属下的外在利益和安全。但还要知道，我们听到关于他的传闻，说他曾放高利贷；你们应当彻查此事，若果真如此，就另选他人，并立即远离这样的人。因为我们绝不对喜爱高利贷的人按手。然而，如果经过仔细查证，此事证明是假的（因为他的本人我们并不认识，也不知道关于他纯朴的报告是否真实），他必须带着你们的推举文书来见我们，以便我们详细考察他的生活和品行，同时知晓他的才智；这样，如果他符合考察，我们就能在主助佑下，借他实现你们的愿望。此外，你们也要留意寻找另一位合适的人，万一此人显得不适于授任此职，你们可将选择转向别人。因为，万一此人未获批准，而你们的圣职人员说没有其他适合被选举的人，那将是重大的耻辱。\par

\SourceRecord{Translated by James Barmby.；Nicene and Post-Nicene Fathers, Second Series；Vol. 13.；Edited by Philip Schaff and Henry Wace.；Buffalo, NY: Christian Literature Publishing Co.,；1898.；<http://www.newadvent.org/fathers/360210062.htm>.}

% structure: p0001 source=h1 target=section reason=page-title

\section{卷十，第六十三封书信}

\Emph{致\Person{多米尼库斯}，\Place{迦太基}主教}\par

\Person{格列高利}致\Person{多米尼库斯}等。\par

我们已经得知多么大的瘟疫侵袭了非洲地区；而且，既然意大利也未能免于这样的苦难，我们悲伤的呻吟便加倍了。但是，在这些灾祸和无数的其他灾难中，最亲爱的兄弟，我们的心会因绝望而衰颓，若不是主的声预先坚固了我们的软弱。因为很久以前，福音教训的号角就已向信徒吹响，警告他们：在世界末日临近时，战争和许多其他事，正如你所知，现在正在恐惧中的事，会发生\EditorialNote{见：玛窦福音24章；路加福音21章}。因此，我们不应该在遭受那些我们预先知道的事时过于悲痛，仿佛它们先前不为人知。并且，在我们思考他人的死亡时，死亡的方式常会是一种缓解。因为我们见过多少残害、多少残酷，在那里死亡是唯一的解救，而生命是一种折磨！当达味面前有选择死亡的方式时，他不是拒绝了饥荒或刀剑，而选择让他的子民落在天主的手中吗？由此可知，那些因天主的击打而灭亡的人获得了多大的恩宠，因为他们死于那作为恩赐赐给圣先知的呼求。因此，让我们在一切逆境中感谢我们的造物主，并信赖祂的仁慈，耐心忍受一切，因为我们所受的远少于我们应得的。然而，既然我们暂时受鞭打，好使我们不致失去永生的安慰，那么（既然我们通过这些征兆的宣告，不无知地知道即将来临的审判者已经近了），我们更应当以热心善工和忏悔的哀泣，使我们将来要呈交祂审查的账目稳妥；好使如此重大的击打，藉着祂恩宠的眷顾，对我们不是惩罚的开始，而是有益于我们的净化。\par

然而，既然我们软弱的本性使我们不得不为逝去的人悲伤，愿你宗座的教导成为受苦者的安慰。把应许的美善之事将与他们同在的道理灌输给他们；好使他们因最确实的望德而坚强，学会不为暂时事物的损失而悲伤，与将来的恩赐相比。让你的口舌，正如我们相信它所做的，越来越多地约束他们行恶；让它宣布善人的赏报、恶人的惩罚，好使那些对善事缺乏爱德的人，至少对恶事大为恐惧，并远离那些必受惩罚的事。因为置身于鞭打之中还行该受鞭打的事，就是对鞭打者特别傲慢，更激怒震怒者。有人不肯公义地停止自己的恶行，却希望天主不公义地停止祂的报复，这乃是一种首要的疯狂。但由于这一切都需要天主的帮助，亲爱的兄弟，让我们以合一的祈祷恳求全能天主的仁慈，好使祂既恩赐我们如此相称地尽职，又仁慈地激动民众的心去做这些事；为使我们既在祂的敬畏中合宜地安排我们的行为，又堪当既脱离即将临头的灾祸，并在祂恩宠的引领下（没有这恩宠，我们什么也不能做），达至天上的喜乐。\par

八月，小纪第三年。\par

\SourceRecord{Translated by James Barmby.；Nicene and Post-Nicene Fathers, Second Series；Vol. 13.；Edited by Philip Schaff and Henry Wace.；Buffalo, NY: Christian Literature Publishing Co.,；1898.；<http://www.newadvent.org/fathers/360210063.htm>.}

% structure: p0001 source=h1 target=section reason=page-title

\section{第十一卷，第一封信}

\Emph{致院长\Person{若望}}\par

\Person{额我略}致\Place{西乃山}院长\Person{若望}。\par

你的谦卑的书信证明了你生活的圣洁；因此我们向全能天主献上极大的感恩，因为我们知道仍有人为我们的罪祈祷。因为我们，在教会治理的掩饰下，被抛入这个世界的波涛中，它常常吞没我们。但藉天上恩宠的保护之手，我们从深处被举起。那么，你，在如此伟大的安息宁静中度过平静生活，并如同安全地站在岸上，请伸出你祈祷之手给我们这些航行中的人，或者更确切地说，遭遇船难的人，并尽你所能的恳求帮助我们努力到达生者之地，好使不仅为你自己的生命，也为我们的救援，你获得永远的赏报。愿天主圣三用祂保护之右手保护你的爱德，并赐予你在祂眼中，藉祈祷、劝诫、善行表样，牧养交托给你的羊群，为使你能与所牧养的羊群一同到达永生的牧场。因为经上记载：\BibleQuote{我的羊将来并找到牧场}\BibleRef{若 10:27}。这些牧场我们确实找到了，当我们脱免此生之寒冬，饱飨永生之青翠，如同新春。\par

我们从我们的儿子\Person{辛普利修}的报告中得知，在那由一位名叫\Person{伊索鲁斯}建造的\OriginalTerm{老人院}{Gerontocomium}中，缺少床铺和被褥。因此我们送去了15件斗篷、30件\OriginalTerm{拉哈纳}{rachanæ}和15张床。我们还给了钱用于购买床垫和运输，恳请你的爱德不要轻视，而将它们供给所送之处。于九月初一，第四指控期。\par

\SourceRecord{Translated by James Barmby.；Nicene and Post-Nicene Fathers, Second Series；Vol. 13.；Edited by Philip Schaff and Henry Wace.；Buffalo, NY: Christian Literature Publishing Co.,；1898.；<http://www.newadvent.org/fathers/360211001.htm>.}

% structure: p0001 source=h1 target=section reason=page-title

\section{第十一卷，第十二封信}

\Emph{致\Place{利里努斯}（\Place{莱兰}）隐修院院长\Person{科农}。}\par

\Person{额我略}致\Place{利里努斯}隐修院院长\Person{科农}。\par

掌权者的谨慎乃属下的保障，因凡守护所托付于己者，便避免了仇敌的陷阱。然而，关于您管理弟兄们的娴熟，以及您对守护他们的热切警醒，我们已从最可敬的弟兄和同为主教的\Person{门纳}的报告中获悉。正如我们听闻您前任轻率的疏忽时常令我们忧伤，同样，您远见的谨慎令我们欣喜，因毋庸置疑，您热忱的守护对您而言是获得赏报的益处，对他人而言则是行善的榜样。\par

然而，既然我们的仇敌越知自己在各方面受到防备，就越是寻求从隐秘途径闯入，并企图以狡诈伎俩推翻对手，但愿您的爱德之警醒不断燃起更炽烈的关切，并藉天主的助佑，预先巩固一切，使那四处奔突的豺狼，在主羊群中无隙可入。因此，在我们救赎主恩宠的助佑下，请您竭尽全力，禁止并全面守护那些托付给您的人，免于贪饕、骄傲、贪婪、闲言以及一切不洁，从而使您因所托付的管理而获得的赏报更为丰厚，因为您的属下借您的警醒，必将胜过仇敌的不义。\par

因此，让善人感受到您的温柔，恶人感受到您的惩戒。即使在惩戒中，也要知道应遵守这个秩序：您应当爱护人而追究过失；以免如果您偶然行事相反，惩戒会变成残忍，而您会毁灭那些您希望挽救的人。因为您应当这样割除疮口，而不使健康部分面临溃烂的风险；否则，如果您按压刀刃超过病情之所需，就会伤害那您急于救助的人。因而，愿您的温和本身是谨慎的，而非松弛的；愿您的惩戒是仁爱的，而非严厉的。但两者应当如此调和，以致善人在被爱时也能有所警戒，恶人在惧怕时也能有所爱慕。\par

最亲爱的儿子，请留心这些事，认真遵守；这样，当您借这样的管理，将那些从天主领受的人安全地归还给天主时，您将在永恒报应之日，堪当听见祂说：\BibleQuote{好！善良忠信的仆人，你既然在少许事上忠信，我要委派你管理许多大事：进入你主的福乐吧！}\BibleRef{路 19:17}此外，我们切愿我们的儿子\Person{科隆博斯}司铎，他因自身的功绩而受推荐于您的爱德，亦因我们的推荐而在您的爱德中增进。\par

\SourceRecord{Translated by James Barmby.；Nicene and Post-Nicene Fathers, Second Series；Vol. 13.；Edited by Philip Schaff and Henry Wace.；Buffalo, NY: Christian Literature Publishing Co.,；1898.；<http://www.newadvent.org/fathers/360211012.htm>.}

% structure: p0001 source=h1 target=section reason=page-title

\section{第十一卷，第十三封书信}

\Emph{致马赛主教塞雷努斯}。\par

额我略致塞雷努斯等。\par

你的来信的开端显示出你具有司铎应有的善意，使我们为你的兄弟之爱更加喜乐。但信末的内容与开端如此不一致，以致这样的一封信可被归因于不同的人，而非同一个人。不，就在你对我们寄给你的信表示怀疑这一点上，已显明你是何等轻率。因为，若你曾留意我们在兄弟之爱中给你的劝告，你不仅不会怀疑，反而会明白以司铎的庄重你应做什么。因为从前担任院长的基里亚库斯，他是我们信的携带者，并非如你所想象的那般有训练和学识，以致敢于另造一封信，或使你对此人的品格持有虚假的嫌疑。但，由于你不顾我们有益的劝告，你不仅在行为上，而且在质疑上也变得有罪。因为确实有报告传到我们这里，说你被轻率的热情所激动，打碎了圣人们的像，仿佛以它们不应受敬拜为理由。的确，就你禁止它们受敬拜而言，我们完全称赞你；但我们责备你打碎了它们。说吧，兄弟，曾听说过有哪个司铎做过你所做的事吗？即使没有别的，难道连这个念头也不应约束你，免得你轻视其他弟兄，而只以为自己圣洁和智慧吗？因为敬拜一幅画是一回事，借着画的故事学习什么是应当敬拜的又是另一回事。因为文字给读者呈现的，画给无学问的观看者呈现，因为连无知的人也在其中看到他们应当追随的；文盲也在其中阅读。因此，尤其是对於外邦人，画代替了阅读。这事尤其应当被你所留意，因为你生存在外邦人中间，免得因不谨慎地受正确热情所激动，而在野蛮人的心中引起冒犯。既然古代并非无理由地准许圣徒的历史被画在可敬的地方，你若以明智调和热情，你无疑会获得你所追求的，而不会驱散已聚集的羊群，反而会聚集分散的羊群；如此，牧人应有的名声就会使你卓越，而非散播者的责备落在你身上。但由於你凭情感的冲动行事不谨慎，据说你如此冒犯了你的子女，以致他们中的大部分已停止与你的共融。那么，你将何时把迷路的羊带回主的羊栈，既然不能留住已有的羊？此后我们劝告你，甚至现在也要谨慎，戒除这种僭越，并赶快以父亲般的慈爱，用一切努力，用一切热忱，将那些你发现已与你分离之人的心召回给你自己。\par

因为教会分散的子女必须被召集起来，然后必须用圣经的见证向他们表明，任何手造之物都不可受敬拜，因为经上记载：\BibleQuote{你要朝拜上主，你的天主，惟独事奉他。}\BibleRef{路 4:8}。然后，关於那些为教育无学问的人民而制作的图画——为使虽不识字的人，通过将目光转向故事本身，能知道所发生的事——必须补充说，因为你曾看到这些图画被敬拜，所以你被激怒而命令打碎它们。还必须对他们说：若你们为了古代制作图像的这个教导，而希望它们在教会中，我完全允许它们被制作和拥有。并向他们解释，使你感到不悦的，并非图画所悬挂展示的故事本身的景象，而是对图画不当的敬拜。用这样的话平息他们的心；召回他们与你和好。若有人愿意制作图像，绝不可禁止他，但务必禁止敬拜图像。不过，让你的兄弟之爱谨慎地劝告他们，从所描绘事件的景象中，他们应领悟忏悔的热忱，并俯伏敬拜唯一全能的天主圣三。\par

我们怀着对圣教会和你的兄弟之爱的爱，说了这一切。因此，不要因我的责备而在正直的热忱中动摇，反而要在你虔诚管理的热忱中得到帮助。\par

此外，我们听到你的仁爱乐意接纳坏人进入你的团体；甚至有一位司铎成为你的密友，他堕落之后，据说仍生活在他不洁的污染之中。这我们并不完全相信，因为接纳这样的人的人并不纠正邪恶，反而似乎给予他人行同样之事的许可。但是，恐怕他藉任何教唆或掩饰说服你接纳他并仍保持宠爱，你不但应当将他从你这里赶得更远，而且应当以司铎的热忱全力剪除他的过犯。至于其他被报告为坏人的人，你要努力以父亲般的劝告约束他们远离恶行，召回他们走向正直之道。但是，若（愿天主不许）你似乎丝毫不能以有益的劝告使他们受益，你也当注意将他们远远抛弃，免得因接纳他们，他们的恶行似乎丝毫不令你不悦，且免得不仅他们本人不改过，而且其他人也因你接纳他们而被败坏。请想一想，在人前多么可憎，在天主眼前多么危险，如果恶习似乎通过那个本应惩罚罪行的人而得到滋养。因此，最亲爱的弟兄，要殷勤留意这些事；并努力如此行事，既健全地纠正坏人，又避免因与恶人交往而在你子女的心中产生冒犯。\par

\SourceRecord{Translated by James Barmby.；Nicene and Post-Nicene Fathers, Second Series；Vol. 13.；Edited by Philip Schaff and Henry Wace.；Buffalo, NY: Christian Literature Publishing Co.,；1898.；<http://www.newadvent.org/fathers/360211013.htm>.}

% structure: p0001 source=h1 target=section reason=page-title

\section{第十一卷，第二十五封信}

\Emph{致\Person{雅努阿里乌斯}，\Place{卡拉利斯}（\Place{卡利亚里}）主教}。\par

\Person{额我略}致\Person{雅努阿里乌斯}等。\par

得知贵司铎团（\OriginalTerm{司铎团}{Fraternity}）的关切令我们欣慰，因为您适当地表现出了牧灵上对灵魂守护的警惕。事实上，我们已获报告，您禁止在贵教会读经员已故\Person{埃皮法尼乌斯}的宅邸内按照其遗嘱建立隐修院，原因是：鉴于该宅邸毗邻一座天主的婢女隐修院，恐由此导致灵魂的欺骗。我们高度赞扬您，因为您以应有的方式，通过适当的预见防备了古老仇敌的陷阱。但是，由于我们获悉虔诚的女士\Person{蓬皮亚娜}渴望将同一隐修院中的天主婢女带走，使她们回到各自原有的隐修院（她们是从那里被带来的），并在那里建立一座隐修士团体，因此，如果此事得以实现，则死者的安排应完全遵守。但如果此事未能实现，为使遗嘱人的意愿不致完全落空，我们愿意——鉴于位于\Place{卡拉利斯}城外的已故院长\Person{乌尔班}的隐修院据说已荒废到连一位隐修士也不剩——我再说，我们愿意让\Person{约翰}（即上述\Person{埃皮法尼乌斯}指定为在他宅邸中建立之隐修院院长的人）被祝圣为院长（\Emph{即已故\Person{乌尔班}的隐修院}），只要对他没有障碍。\par

并且，原应存放于上述\Person{埃皮法尼乌斯}宅邸中的圣髑应存放于彼处，且同一\Person{埃皮法尼乌斯}为其宅邸中计划建立的隐修院所捐献的一切，应全部用于另一隐修院；如此，即使如上所述，出于安全考虑，他的意愿在地点上未能实现，但所意图的益处仍可完整保存。确实，请贵司铎团与监护人（\OriginalTerm{监护人}{defensore}）\Person{维塔利斯}一同安排这一切，并努力如此妥善处理，使您既能因值得称赞的禁令，也能因对案件的良好解决而获得赏报。最后，尽管将这座隐修院托付给贵司铎团或许是多余的，但我们仍充分敦促您，按照您的本分，在适当顾及正义的情况下，视其为托付于您。\par

\SourceRecord{Translated by James Barmby.；Nicene and Post-Nicene Fathers, Second Series；Vol. 13.；Edited by Philip Schaff and Henry Wace.；Buffalo, NY: Christian Literature Publishing Co.,；1898.；<http://www.newadvent.org/fathers/360211025.htm>.}

% structure: p0001 source=h1 target=section reason=page-title

\section{第十一卷，第二十八封信}

\Emph{致\Person{奥古斯丁}，盎格鲁人的主教}。\par

\Person{额我略}致\Person{奥古斯丁}等。\par

\BibleQuote{光荣归于至高之天主，与世人之善意和平}\BibleRef{路 2:14}；因为一粒麦子落在地里死了，为的是不独自为王于天堂；甚至祂——我们因祂的死而活，因祂的软弱而坚强，因祂的受苦而脱离苦难，因祂的爱我们在不列颠寻找素不相识的弟兄，因祂的恩赐我们找到那些我们不知而寻找的人。然而谁能描述这里在众信友心中涌起的巨大喜乐，因为盎格鲁民族通过全能天主的恩宠和你的兄弟情谊的劳动，已经抛弃了谬误的黑暗，被圣教信德之光所照耀；以最健全的心智，它现在践踏从前在疯狂恐惧中跪拜的偶像；它用纯洁的心跪伏在全能天主面前；它被神圣讲论的规则约束，远离恶行的过犯；它在心中俯首于神圣的诫命，为在理解中得以高举；它甚至在祈祷中自卑至地，以免在理智和灵魂中躺卧于地。这是谁的作为，若不是祂所说：\BibleQuote{我父工作至今，我也工作}\BibleRef{若 5:17}？祂为显示祂转变世界，不是凭人的智慧，而是凭祂自己的大能，选择了没有学问的人作祂的宣讲者，并派遣他们进入世界。如今祂也照样而行，屈尊通过软弱的人在盎格鲁民族中行大能的事。但在这一属天的恩赐中，最亲爱的弟兄，伴随着巨大的喜乐，也有极严肃的恐惧。因为我知道，全能天主通过你的爱德，在祂愿选定的民族中显示了伟大的奇迹。因此，你必须对这一属天的恩赐既喜乐又恐惧，在喜乐中战栗：——喜乐，是因为盎格鲁人的灵魂被外在的奇迹吸引至内在的恩宠；但恐惧，免得在发生的征兆中，软弱的的心智自高自大，以致骄傲，从外在的尊荣被高举，内里却因虚荣而堕落。因为我们应记起，当门徒们从宣讲归来，喜乐地向他们天上的导师说：\BibleQuote{主，因你的名，连魔鬼也屈服于我们}\BibleRef{路 10:17}，他们随即听到：\BibleQuote{不要因这个欢喜，却要因你们的名称已录在天上而欢喜}\BibleRef{路 10:20}。因为他们既以奇迹为乐，便把心思置于私人的和暂时的喜乐上。但他们被从私人的召回共同的，从暂时的召回到永远的喜乐时，便对他们说：\BibleQuote{你们当因这个欢喜，因为你们的名称已录在天上}。因为并非所有被选者都行奇迹；然而他们所有人的名字都登录在天上。因为真理的门徒不应欢喜，除非为他们与他人共有的善，并且在这善中他们的喜乐没有终结。\par

因此，最亲爱的弟兄，在你藉天主的运作而外在实行的事情中，你应当时常在内心审慎地判断自己，审慎地了解你自己是什么，以及在那同一民族中为使他们悔改而赐予你的恩宠是何等伟大，你甚至领受了行奇事的恩赐。如果你有时记得曾以言语或行为冒犯我们的造主，要常将这些事记在心头，使罪过的记忆抑制心中升起的荣耀。无论你将要或已经领受行奇事的能力，应视这些能力并非赐给你自己，而是为了那些因你而得救的人。此外，当我思想这些事时，我记起一位天主仆人的事，甚至是特选的仆人。诚然，正如你的兄弟情谊所知，梅瑟领天主的人民出埃及时行了奇妙的奇迹。在西乃山禁食四十昼夜，他领受了法律版；在闪电和雷声中，全体人民战栗，他专务于全能天主的服务，甚至与祂亲密交谈\EditorialNote{出谷纪 30 至 31 章}；他开辟了红海的路；有云柱带领他的行程；给饥饿的人民从天降下玛纳；在旷野中奇迹般地为渴望肉食的人供应了肉食，甚至过于饱足\EditorialNote{出谷纪 13, 14, 16 章}。但在干旱之时，他们来到磐石，他心存疑惑，怀疑能否从同一磐石中取水，然而仍按\OriginalTerm{上主}{Lord}的命令毫不失败地涌出丰沛的水流。此后他在旷野中三十八年所行的众多奇迹，谁能计算或查考\EditorialNote{出谷纪 17 章；户籍纪 20 章}？每当心中被疑惑之事困扰，他就进入会幕，私下求问上主，并立即得到指示，天主对他说话\EditorialNote{出谷纪 33 章及以下}。当上主向人民发怒时，他藉祈祷的代求平息了祂；那些骄傲自大、纷争不和的人，他使他们被裂开的地口吞没；他以胜利制服了仇敌，并向自己的人民显示征兆。但当终于到达应许之地时，他被召上山，听到他三十八年前所犯的过错，如我所说，就是怀疑从磐石取水的事。为此，他被告知不得进入应许之地\EditorialNote{户籍纪 27 章}。在此我们应思量全能天主的审判何等可畏，祂藉着祂的那位仆人行了如此多的奇迹，却又将过犯记念如此之久。\par

因此，最亲爱的弟兄，如果我们发现连那位我们知道曾特别被全能天主拣选的人，在行了如此多的奇迹之后，竟因一个过失而死，那么我们这些还不知道自己是否被拣选的人，该怀着多大的恐惧颤抖呢？\par

但是我该如何谈论被弃绝的人所行的异能呢？你的仁爱深知真理在福音中所说的：\BibleQuote{\Emph{在那一天，许多人要对我说：主啊，我们因你的名预言，因你的名驱魔，因你的名行了许多异能。但我要向他们声明：我不认识你们；你们这些作恶的人，离开我吧！}}\BibleRef{玛 7:22}？因此，在神迹与异能中间，心智应大大抑制，免得有人在其中寻求自己的光荣，并因自己的高举而暗自欢欣。因为应借着神迹寻求灵魂的收益，并寻求那行这些神迹之天主的光荣。然而，主赐予我们一个神迹，我们可以因此而格外喜乐，并在我们自身承认被选的光荣，因为祂说：\BibleQuote{\Emph{如果你们彼此相爱，世人因此就可认出你们是我的门徒。}}\BibleRef{若 13:35}这神迹正是先知所求的，当时他说：\BibleQuote{\Emph{主啊，求祢赐给我一个善的记号，使恨我的人看见而羞愧。}}\BibleRef{咏 86:17}\par

我说这些话，是因为我想使听者之心谦卑自抑。但让你的谦卑拥有自信。因为我，一个罪人，怀有最确定的希望：借着我们全能造物主和救主，我们的天主和主耶稣基督的恩宠，你的罪已得赦免，并且你被选为这个目的，即借着你使其他人的罪得赦免。将来你也不会为任何罪疚而忧伤，因为你努力使许多人悔改而给天堂带来喜乐。的确，我们的同一造物主和救主谈论人的悔改时，说：\BibleQuote{\Emph{我告诉你们，对于一个罪人悔改，在天上的欢乐比对于九十九个无须悔改的义人还要大。}}\BibleRef{路 15:7}如果为一个悔改者在天上就有这么大的欢乐，那么为如此众多的人民，从他们的谬误中悔改，因信德而来，以补赎谴责了他们所做的恶事，我们该相信欢乐会是多大呢？因此，在这天堂和天使的欢乐中，让我们重复我们开头所说的天使的话：让我们说，让我们所有人说：\BibleQuote{\Emph{光荣归于至高天主，在地上和平归于善人。}}\par

\SourceRecord{Translated by James Barmby.；Nicene and Post-Nicene Fathers, Second Series；Vol. 13.；Edited by Philip Schaff and Henry Wace.；Buffalo, NY: Christian Literature Publishing Co.,；1898.；<http://www.newadvent.org/fathers/360211028.htm>.}

% structure: p0001 source=h1 target=section reason=page-title

\section{第十一卷 第二十九函}

\Emph{致盎格鲁人的王后伯莎}。\par

\Person{格列高利}致\Person{伯莎}等。\par

凡渴望在尘世统治之后获得天国荣耀的人，应当热切努力为他们的创造主带来收益，以便能藉其行动的阶梯升达所渴慕的事物；正如我们乐于知道你所做的那样。事实上，我们至爱的儿子、司铎\Person{劳伦斯}和隐修士\Person{伯多禄}回来时向我们报告了陛下如何对待我们至可敬的弟兄和同为主教的\Person{奥古斯丁}，以及你赐予他何等巨大的帮助和爱德。我们赞美全能的天主，祂仁慈地乐意将盎格鲁民族的归化保留为你的赏报。因为，正如祂藉着享有盛誉的\Person{海伦娜}——最虔诚的皇帝\Person{君士坦丁}的母亲——点燃罗马人心中的信仰之火，同样，我们相信祂也藉着陛下的热忱在盎格鲁民族中工作。诚然，你身为真正的基督徒，本应早已运用你的睿智，以善加影响我们荣耀的儿子、你的丈夫的心，使他为了他的王国和他自己灵魂的利益，追随你所宣认的信仰，以便因他以及藉着他整个民族的归化，你得以在天堂的喜乐中获得相称的酬报。因为，如我们所说，陛下既以正统信仰坚固，又受学识教养，此事对你应非迟缓或困难。如今，因天主的旨意，时机已经成熟，所以请藉着恩宠的合作继续努力，以便能弥补以往的疏忽并有所增益。因此，以不断的劝勉坚固你荣耀的丈夫对基督信仰的热爱；让你的关切向他注入对天主日益增长的爱，如此点燃他的心，甚至为了他所辖全民族的彻底归化，使他能藉你虔诚的热忱向全能的主献上伟大的祭献，并使关于你的事迹得以增长并在各方面显为真实；因为你的善行不仅为热切为你的生命祈祷的罗马人所知，也传遍各处，甚至传到了君士坦丁堡最宁静的皇帝耳中。因此，正如你的基督信仰之慰藉使我们大得喜乐，同样，愿你完成的工作在天上也带来喜乐。请在归化你的民族一事上，全心全意地协助我们上述至可敬的弟兄和同为主教的\Person{奥古斯丁}以及我们派往你那里的天主仆人们，以便你与我们荣耀的儿子、你的丈夫在此世幸福地统治，并在漫长岁月之后，也获得未来永无止境的生命之乐。现在我们祈求全能的天主，愿祂以恩宠之火点燃陛下的心去完成我们所说之事，并为你那蒙祂悦纳的工作赐予永恒赏报的果实。\par

\SourceRecord{Translated by James Barmby.；Nicene and Post-Nicene Fathers, Second Series；Vol. 13.；Edited by Philip Schaff and Henry Wace.；Buffalo, NY: Christian Literature Publishing Co.,；1898.；<http://www.newadvent.org/fathers/360211029.htm>.}

% structure: p0001 source=h1 target=section reason=page-title

\section{第十一卷 第三十封书信}

\Emph{致前隐修士、叙拉古贵族韦南提乌斯}。\par

格列高利致韦南提乌斯等。\par

在向您致以应有的问候时，我本打算讲述我所受的苦。但我认为无需向您说明您已知晓的事。因为我正受着痛风的折磨，这病痛同样折磨着您和我，在我们身上日益加剧，我们的生命也随之消减。在这些痛苦中，我们还能做什么呢？除了追忆我们的过犯，并感谢全能的天主。因为我们因肉体的纵欲在许多事上犯了罪，如今藉着肉体的磨难得以炼净。我们也当知道，现世的痛苦若能使受痛苦者悔改，便是先前罪过的终结；但若不能使人转向敬畏主，便是后续痛苦的开始。因此我们务必谨慎，以全心归向，用尽全力以眼泪警醒，以免我们从磨难走向磨难。我们还应思量，我们的创造者以何等慈爱的安排对待我们：祂不断击打我们这些该死的人，却仍未杀死我们。祂警告祂将要做的事，却仍未做，为要叫先来的痛苦惊吓我们，当我们归向这位严厉审判者的敬畏时，就能在生命结束后免受祂的惩罚。谁能述说，谁能计数，有多少人沉溺于淫乱、妄言和骄傲中，继续抢劫和不义直至死日，他们在今世如此生活，竟从未头痛过，却因突然一击被交付地狱之火？我们则因不断受鞭打而有不被抛弃的标记，正如圣经所证：\BibleQuote{因为主惩戒他所爱的，鞭打他所接纳的每个儿子}\BibleRef{希 12:6}。因此，在天主的鞭打下，让我们追忆祂的恩赐和我们罪过的亏损。让我们思量，祂何等美善地降福于我们的恶行，而我们在祂的良善下又行了何等恶事。让我们实行主藉先知所说的：\BibleQuote{请你提醒我，让我们一同辩论}\BibleRef{依 43:26}。让我们如今在思想中与天主辩驳，以免日后受天主严厉审判。因为保禄说：\BibleQuote{如果我们先省察自己，就不致受主的审判}\BibleRef{格前 11:31}。因此，凡愿逃避将来审判严厉判决的人，就当藉补赎的苦痛为自己断绝今世一切甘饴。此外，无论这一类恩赐是什么，除了我们的创造者之外，谁的恩赐呢？但那藉本身的喜乐使我们与天主的爱分离的，不应当被视作天主完全给我们的恩赐，免得我们宁愿选择恩赐而轻忽施予者，并在接受善物时（尽管我们自己不善）因那本应使我们敬畏增长的事物而脱离了对祂的敬畏。愿万物的创造者，即全能的天主，藉祂圣神的感动，将我们以书信向您所说的话倾注于您心中，洗净您一切罪恶的污秽，赐予您在此世祂安慰的喜乐，并在日后赐予您与祂同在的永恒赏报。我请求代我问候我最亲爱的女儿们，\Person{芭芭拉}夫人和\Person{安东尼娜}夫人。\par

\SourceRecord{Translated by James Barmby.；Nicene and Post-Nicene Fathers, Second Series；Vol. 13.；Edited by Philip Schaff and Henry Wace.；Buffalo, NY: Christian Literature Publishing Co.,；1898.；<http://www.newadvent.org/fathers/360211030.htm>.}

% structure: p0001 source=h1 target=section reason=page-title

\section{第十一卷 第三十二封书信}

\Emph{致\Place{拉文纳}主教\Person{马里尼亚诺斯}。}\par

\Person{额我略}致\Person{马里尼亚诺斯}，等。\par

当此函的携带着——院长\Person{卡尼迪斯}前来祈求圣髑（并已获准），我因见到您的兄弟之爱所提供的护理帮助而欣喜，由此显明了您对我之关怀，但同时我又因未能如我所愿地享受他的陪伴而痛苦，因为他发现我卧病在床，离开时我仍体弱无力。因为距今已久，我未能起床。因为时而痛风之痛折磨我，时而一种不知何物的火灼之痛弥漫全身，且通常同时有灼痛折磨我，身心俱疲。此外，我还遭受了哪些其他巨大的病痛，我无法尽述。但我可简要言之：一种有害体液的感染如此侵蚀我，以致活着对我来说是痛苦，我焦急地盼望死亡，唯有死亡能解除我的呻吟。因此，至圣的兄弟，请为我祈求神圣仁慈的垂怜，愿它仁慈地减轻对我打击的鞭笞，并赐我忍耐之力，以免（愿天主不许）我的心因过度疲乏而爆发不耐，因抱怨而增加那本可借鞭笞痊愈的罪债。颁于二月，小纪第四年。\par

\SourceRecord{Translated by James Barmby.；Nicene and Post-Nicene Fathers, Second Series；Vol. 13.；Edited by Philip Schaff and Henry Wace.；Buffalo, NY: Christian Literature Publishing Co.,；1898.；<http://www.newadvent.org/fathers/360211032.htm>.}

% structure: p0001 source=h1 target=section reason=page-title

\section{卷十一，第三十三信}

\Emph{致\Person{马里尼亚努斯}，\Place{拉文纳}主教。}\par

\Person{额我略}致\Person{马里尼亚努斯}等。\par

一位\Place{拉文纳}人来到这里，我得知您因吐血患病，极为悲伤。为此，我们已仔细逐一咨询我们所知的博学医师，并将他们的各别意见和处方以书面呈送您的圣座。然而，众医皆首重安静与沉默，我甚怀疑您在您自己的教会中能否做到。因此，在我看来，当\Place{拉文纳}教会已得到妥善安排——无论由能完成弥撒礼仪者，或能负责主教职务、并能施行款待和接待者，或能监督隐修院监护者——之后，您应在夏季之前到我这里来，以便我尽力特别照料您的疾病，并使您不受打扰，因为医师们说夏季对此类疾病极其危险。我甚恐如果您因季节不利而有所忧虑，此症可能带来更大危险。我本人也很虚弱，且无论从哪方面说，您蒙天主眷顾健康返回教会是最有益的；或者，如果您将蒙召，则应在朋友们手中蒙召；而我自知离死甚近，若全能天主乐意在我之先召您，我愿在您手中去世。但如果当前时局妨碍您来，可给\Person{阿戈}少许礼物，以便他派人护送您至\Place{罗马}。那么，如果您感到被此症沉重压制，并安排前来，您须带少数随从，因为当您与我同住主教官邸（\Emph{episcopium}）时，您会得到此教会的每日侍奉。\par

再者，我既不劝勉也不告诫您，而是严格命令您，绝不可擅行禁食，因为医师们说此习对这类疾病极为有害；除非偶尔有重大庆典要求禁食，我允许一年五次。您也必须避免守夜；并让\Place{拉文纳}城习惯在蜡炬上诵念的祈祷，以及司铎在逾越节庆典前后所作的福音诠释，由他人代行。亲爱的，绝不可担任何超乎您能力的事务。我说这些话是为了，如果您感到康复并推迟来访，您知道应遵行我的命令。\par

\SourceRecord{Translated by James Barmby.；Nicene and Post-Nicene Fathers, Second Series；Vol. 13.；Edited by Philip Schaff and Henry Wace.；Buffalo, NY: Christian Literature Publishing Co.,；1898.；<http://www.newadvent.org/fathers/360211033.htm>.}

% structure: p0001 source=h1 target=section reason=page-title

\section{第十一卷，第三十五封书信}

\Emph{致\Person{巴巴拉}与\Person{安东尼纳}}。\par

\Person{额我略}致\Person{巴巴拉}等\par

收到了尊荣的信函，其中言辞含泪，我们最亲爱的女儿们，对于你们父亲的疾病，我们心中的哀伤并不亚于你们。因为我们不能将那因\OriginalTerm{爱德}{charity}而成为我们自己的悲伤视为无关。但是，既然在任何绝望的境况中，都不应怀疑我们\OriginalTerm{救赎主}{Redeemer}的仁慈，请为安慰你们父亲而振作精神，将\OriginalTerm{望德}{hope}寄托于\OriginalTerm{全能天主}{Almighty God}的手中，我们相信，藉祂的保护，祂将护佑你们免受一切灾祸，安慰你们的苦难，并按你们父亲的意愿使你们得到妥善安置。然而，若祂偿还了我们人类命运的债务，即使那时，也不让任何绝望压垮你们，也不让任何人的话语使你们惊慌。因为在天主——孤儿的管理者和保护者之后，我们将如此关心你们最甜美的\OriginalTerm{尊荣}{Glory}，并将在主的助佑下尽力为你们的利益提供一切，使不义之人的粗暴对待不会扰乱你们，并使我们能以各种方式偿还从你们父母的良善中所欠下的债务。愿\OriginalTerm{天上的恩宠}{heavenly grace}以它的眷顾滋养你们，以它的保护护佑你们免受一切邪恶，使你们的平安成为我们的喜乐。\par

\SourceRecord{Translated by James Barmby.；Nicene and Post-Nicene Fathers, Second Series；Vol. 13.；Edited by Philip Schaff and Henry Wace.；Buffalo, NY: Christian Literature Publishing Co.,；1898.；<http://www.newadvent.org/fathers/360211035.htm>.}

% structure: p0001 source=h1 target=section reason=page-title

\section{\WorkTitle{第十一卷·第三十六封书信}}

\Emph{致叙拉古主教若望}。\par

额我略致若望等。\par

我收到贵弟兄的信函，告知我最亲爱的儿子文南提乌斯老爷患病，并述及关于他的一切进展。但我同时得知他病重危急，且不义之人竟染指孤儿财产，我心悲伤，几乎无法自持。然而，其中的安慰在于，泪水缓解了我的呻吟。因此，你的圣德不应忽视那本该是你首要关切之事：关心他的灵魂，通过劝勉、恳求、向他揭示天主的可怕审判，并许诺祂不可言喻的仁慈，以劝诱他即使在最后时刻也能回归他先前的生命状态，免得如此重大过犯的罪债在永恒的审判中与他为敌。然后，你有责任考虑他的女儿们——芭芭拉小姐和安东妮娜小姐——应当如何处置，以免给恶人留下可乘之机。因为他在恳求我焦虑地照顾她们之后，又补充说我要负责她们的婚配，接着在信中提及一事，我思量此事时，无疑认为它可能成为障碍。他说我应再三恳求最虔诚的皇帝陛下，让他亲自为她们的婚配作出安排。你看出这与先前的愿望是何等不同。我担心这会给西西里那些寻找一切机会干涉他事务的人提供恰当的时机。因为一旦此事为人所知，那些据闻已试图查封他财产的人会做些什么？难道理由不会站在他们一边，并为他们提供似乎正当的干预依据吗？倘若他们说：这些女孩已托付给皇帝陛下，我们不能忽视此事，若不干预，后果自负；我们保管财产，直到皇帝陛下命令将她们带往君士坦丁堡为止——我请你告诉我，在这种情况下我能做什么？父亲的托付似乎支持了一个有权势的人。因为他恳求我负责安排她们，使她们或留在罗马城，或不从西西里被带走；而他这样行事，竟不留任何将她们带来或留在那里的余地。但你，要尽你所能反对这些恶人。为了全能的天主，保护她们的财产如同保护你自己的；如果仍有可能，设法让一切不义的机会都从上述文南提乌斯老爷的遗嘱中除去。但是，如果认为应将她们托付给宫廷，他不应把如此重担加于我，希望我以灵魂负责她们婚配之责；关于此事，全能的天主知晓我如何思虑，这就足够了。因此，我已立即致信我最亲爱的儿子、执事阿纳托利乌斯，命他尽力与光荣的贵族妇人鲁斯蒂齐娅谈话，并告诉他如何询问我并告知我关于那些已传给我姓名的人的情况，以便他能迅速告知我们一切，并藉天主的安排，决定应当做什么。\par

此外，在送抵我们的信函中，我们发现贵弟兄因我们不愿你前来此地而忧伤，仿佛这是出于某种不悦；而我们行事纯然出于利益考虑，深知因你所在地区的人，你在那里极其必要。但，免得你因此以为我们对你有任何情绪或不悦（愿天主禁止），若你愿意来我们这里，请适时出现在宗徒们的门槛前。因为，就我们而言，我们如此爱你的爱德，以致渴望常见到你。\par

\SourceRecord{Translated by James Barmby.；Nicene and Post-Nicene Fathers, Second Series；Vol. 13.；Edited by Philip Schaff and Henry Wace.；Buffalo, NY: Christian Literature Publishing Co.,；1898.；<http://www.newadvent.org/fathers/360211036.htm>.}

% structure: p0001 source=h1 target=section reason=page-title

\section{第十一卷　书信三十七}

\Emph{致\OriginalTerm{护民官}{Defensorem}罗玛诺。}\par

\Person{额我略}致\Place{西西里}护民官\Person{罗玛诺}。\par

我们得知，若有人对某位圣职人员提起诉讼，你竟将这些圣职人员带到你面前受审，无视他们的主教。倘若如此，此事显然极不合适，兹凭我们的权柄命令你，不得再如此妄为。但是，若有人对圣职人员提起诉讼，应前往该圣职人员的主教处，由主教亲自审理，或至少由他委派法官；如果是仲裁案件，则由他委派的行政当局强制当事人选择一位仲裁法官。但是，若有圣职人员或平信徒对主教提起诉讼，则你应当介入，或由你亲自审理他们之间的案件，或经你劝告让他们自行选择法官。因为，如果每位主教都不能保留自己的司法权，岂非正是由我们——本应维护教会秩序的人——使教会秩序陷入混乱？\par

此外，有人向我们报告，我们可敬的弟兄\Person{若望}主教曾将若干有过失的圣职人员处以补赎，而你却未告知他，便擅自将他们撤除补赎。倘若属实，须知你所行之事极不恰当，理应受到严厉斥责。因此，即刻将这些圣职人员交还给他们的主教。今后切勿再犯此过：你若疏忽，须知必将在极大程度上招致我们的愤怒。\par

\SourceRecord{Translated by James Barmby.；Nicene and Post-Nicene Fathers, Second Series；Vol. 13.；Edited by Philip Schaff and Henry Wace.；Buffalo, NY: Christian Literature Publishing Co.,；1898.；<http://www.newadvent.org/fathers/360211037.htm>.}

% structure: p0001 source=h1 target=section reason=page-title

\section{第十一卷，书信三十八}

\Emph{致\Person{Vitus}，\OriginalTerm{护卫者}{Defensorem}。}\par

\Person{Gregory}致\Person{Vitus}等。\par

如果你不受任何身体服役的条件或责任的约束，也未曾是任何其他城市的圣职人员，并且对你没有法律上的反对理由，那么为了教会的利益，我们愿意并乐意你接受教会护卫者的职务，以便你能廉洁而勤勉地执行我们为贫民之益所委托给你的一切；善用我们经过深思后授予你的这项特权，忠实地尽力完成我们委托给你的一切，因为你必须在我们的天主的审判台前为你的行为交账。这封书信我们已经口授给\Person{Paterius}，我教会的\Emph{\OriginalTerm{副公证人}{secundicerio notario}}，并已签署。\par

\SourceRecord{Translated by James Barmby.；Nicene and Post-Nicene Fathers, Second Series；Vol. 13.；Edited by Philip Schaff and Henry Wace.；Buffalo, NY: Christian Literature Publishing Co.,；1898.；<http://www.newadvent.org/fathers/360211038.htm>.}

% structure: p0001 source=h1 target=section reason=page-title

\section{第十一卷，第四十封信}

\Emph{致\Place{拉文纳}主教\Person{马里尼亚诺}}\par

\Person{额我略}致\Person{马里尼亚诺}等\par

极大的软弱束缚着我们，至爱的弟兄。如果我们能摆脱它，我们似乎就应受到责备。然而，既然在这脆弱的身体内，我们只有顺从它的软弱才能存活，我们就不应以必须之事为耻。因此，既然医生们都说，对出血患者而言，斋戒是有害的，我们藉本次信函劝勉你的贤弟：你要记住你过去因疾病所受的苦，绝不可将斋戒的劳苦强加于自己。然而，若因天主的仁慈，你自知健康状况大有改善，以致有足够的力量，我们允许你每周斋戒一两次。但在这一切事上，你首先应当注意，绝不可使自己陷入任何恼怒的情绪，免得那被认为现已减轻、暂缓的疾病，因激怒而以后更重地发作。\par

\SourceRecord{Translated by James Barmby.；Nicene and Post-Nicene Fathers, Second Series；Vol. 13.；Edited by Philip Schaff and Henry Wace.；Buffalo, NY: Christian Literature Publishing Co.,；1898.；<http://www.newadvent.org/fathers/360211040.htm>.}

% structure: p0001 source=h1 target=section reason=page-title

\section{第十一卷，第四十四封书信}

\Emph{致鲁斯蒂西亚娜，贵妇}。\par

格列高利致鲁斯蒂西亚娜等。\par

我收到了阁下的书信，在我正患重病时，这些信使我因您的健康、虔诚与温良而完全得安慰。然而有一件事我颇感不悦，即在同一封致我的书信中，反复出现了本可只说一次的话：\Quote{你的婢女}与\Quote{你的婢女}。因为，我因着主教的重担成了众人的仆人，在我接受主教职务之前本是属于她的人，她又有何理由自称是我的婢女呢？因此，我恳求全能的天主，愿我再也看不到您写信给我时使用这个词。此外，您出于最纯洁真挚之心献给真福伯多禄、宗徒之长的礼物，已被收下并悬挂在那里，全体圣职人员均在场。然而，我的儿子，尊贵的君西玛库，见我因痛风而患病且几乎无望，便推迟将您的信交给我，直到帷幔收到很久之后才给；我后来在阁下的信中发现，这些帷幔本应伴随着祷文被送至真福伯多禄的圣殿。因此这事没有做到，因为正如我所说的，我们在收到信之前就收到了帷幔。尽管如此，上述的君西玛库与您的全家一起做了您希望我们与圣职人员一起做的事。但即使缺少人的声音，您的奉献本身在天主面前也有其声音。我信赖祂的慈爱，愿那位您用帷幔遮盖其遗体于地上者的转求，能在天上保护您脱离一切罪恶，并以祂的眷顾治理您的家，以祂的守护护卫它。\par

关于您信中告知已降临到您身上的痛风之苦，我既感到痛苦又极其喜乐：喜乐，因为有害的体液侵蚀了您身体的下部，已完全离开了上部；但痛苦，因为我担心您在如此纤弱的身体中承受着过度的疼痛。因为，哪里缺乏肉体，哪里又有力量抵抗疼痛呢？至于我自身，您知道我从前是怎样的；但如今灵魂的苦涩与不断的恼怒，加之痛风的折磨，使我的身体如同在坟墓中一般干枯。因此，我如今很少能从床上起来。如果痛风的痛苦使我的身体干瘦到如此程度，那么我对您的身体又该如何想象呢——它在疼痛来临前就已经过于干瘦了？\par

至于您赐予真福宗徒安德肋隐修院的施舍，我无需多说，因为经上记载：\BibleQuote{要将你的施舍藏在穷人的怀里，它要为你祈祷}（\BibleRef{德 29:15}）。如果善行本身在天主隐秘的耳中发出声音，那么无论我们高声呼喊还是保持沉默，这件您所行的善事本身就在呼喊。此外，我声明：在这同一所宗徒的隐修院里，有如此多的奇迹、如此多的关怀与对隐修士的守护，以至于仿佛这位宗徒本人就是隐修院的特别院长。因为，从我跟隐修院院长与院牧的叙述中所学到的众多事件中，只说几件：有一天，两位弟兄被派出去为隐修院的用途购买物品，一位是较年轻的，似乎以谨慎著称，另一位是年长的，被派去保护年轻的。两人出发了，从他们收到的作为购买物品价款的银钱中，那位被派保护年轻的弟兄竟背着对方私藏了一些。两人随即返回隐修院，来到祈祷所的门口，那犯偷窃的弟兄被恶魔附身跌倒，开始受折磨。当恶魔离开他后，围上来的隐修士们问他是否从所收的钱中私藏了什么；他否认，于是第二次受折磨。他否认了八次，八次受折磨。但在第八次否认后，他承认私下拿了多少钱。他悔改，伏在地上承认自己犯了罪，当他做了补赎后，恶魔不再来搅扰他。\par

另有一次，在同一宗徒的周年纪念日，弟兄们正在午间休息时，突然一位弟兄睁着眼睛瞎了，开始颤抖，高声喊叫，以此表明他无法忍受所遭受的。弟兄们跑向他，见他睁着眼睛却瞎了，颤抖着，叫喊着，精神恍惚，对外界的事毫无知觉。他们用手扶起他，把他扔在圣安德肋宗徒的祭台前，也伏在地上为他祈祷。他立刻恢复知觉，讲述了所遭受的事：一位老人出现在他面前，放出一条黑狗来撕咬他，说：你为什么想要逃离这座隐修院？而他说，我根本无法逃脱那只狗的撕咬，这时几位隐修士来了，为我在那位老人前求情，他立即命令狗离开，我便恢复了神志。他后来多次承认说：在我遭受这些事的那天，我本有逃离这座隐修院的念头。\par

另一位隐修士也暗中想要离开同一座隐修院。他在心中考虑了此事，正要进入祈祷所时，立刻被交付给恶魔，受到极重的折磨。但当恶魔离开他时，如果他留在祈祷所外面，就不会受伤害；但若他试图进入，便立刻被交给恶魔受折磨。这事多次发生后，他承认了自己的过失，即他正考虑离开隐修院。于是，弟兄们聚集为他祈祷，决定为他连续祈祷三天，他竟如此痊愈，此后恶魔再没有临到他。他还说，在他受折磨时曾见到同一的真福宗徒，并因他想离开隐修院而受到他的斥责。\par

另外两位弟兄也从同一座隐修院逃走了，他们先前在与弟兄们交谈时透露了一些迹象，说他们将经由\Place{亚壁古道}下去，前往耶路撒冷；但出去后，他们偏离了道路。而且，为了使任何追赶他们的人都不可能找到他们，他们在\Place{弗拉米尼亚门}附近找了一些偏僻的地窟，藏身其中。傍晚时分，当他们被寻找却未在隐修院中找到时，几位弟兄骑马追赶他们，从\Place{梅特罗努斯门}出去，沿着\Place{拉丁大道}或\Place{亚壁古道}追赶。但突然间，他们起意沿\Place{萨拉里亚大道}继续寻找他们；于是，在出城后，他们转道\Place{萨拉里亚大道}。然而没有找到他们，他们决定经由\Place{弗拉米尼亚门}返回。当他们返回时，他们骑的马来到那两人藏身的地窟前面，突然停住了，尽管受到鞭打和驱赶，却拒绝移动。隐修士们认为这种情形不可能没有奥秘。他们观察了地窟，看到入口被一堆石头堵住了，但由于他们的马不肯往任何方向走，他们下了马。他们挪开堆积在地窟口的石头，走了进去，发现那两人在这黑暗的地下藏身处惊恐万分。他们被带回隐修院，并因这一奇迹而大为改进，以致于这次短暂逃离隐修院对他们大有裨益。\newline{}\par

我告诉您这些事，是为了让阁下知道您施舍救济的祈祷所是属于谁的。愿全能的天主在灵魂和身体上保护您和您的全家，使您蒙受祂天上的庇佑，并赐您长寿以安慰我们。我恳请我最亲爱的儿子\Person{斯特拉特吉乌斯}大人与您光荣的父母以及您的孩子们代我问候。\par

\SourceRecord{Translated by James Barmby.；Nicene and Post-Nicene Fathers, Second Series；Vol. 13.；Edited by Philip Schaff and Henry Wace.；Buffalo, NY: Christian Literature Publishing Co.,；1898.；<http://www.newadvent.org/fathers/360211044.htm>.}

% structure: p0001 source=h1 target=section reason=page-title

\section{第十一卷·第四十五封信}

\Emph{致贵妇\Person{特奥克蒂丝塔}}。\par

\Person{格列高利}致\Person{特奥克蒂丝塔}等。\par

我们应当向全能天主献上极大的感恩，因为我们的至为虔诚、至为仁慈的皇帝们身边有同族的亲属，他们的生活与言谈给我们大家带来极大的喜乐。因此我们也该不断为这些我们的主上祈祷，愿他们的生命及所有属于他们的人，在天上恩宠的保护下，得以长久而平安地延续。\par

然而我须告知您，我从某些人的报告中得知，因人民的轻浮，有毁谤的骚动起来反对您。我听说阁下因此感到不小的烦恼。如果真是这样，我甚为惊讶，为何地上人的言语能搅扰您，您已将心安置于天上。因为真福\Person{约伯}，当那些前来安慰他的朋友发出斥责时，说：\BibleQuote{因为，请看，我的见证在天上，认识我的在高处}\BibleRef{约 16:20}。凡在天上有其生活见证的人，不应害怕地上人的判断。\Person{保禄}，善人的导师，也说：\BibleQuote{我们的光荣在于良心的见证}\BibleRef{格后 1:12}。他又说：\BibleQuote{各人应省察自己的行为，这样，他只在自身有光荣，而不在别人身上}\BibleRef{迦 6:4}。因为，如果我们因赞美而喜乐，因毁谤而沮丧，我们就是将光荣不放在自己内，而放在别人的口舌上。的确，愚拙的童女没有在器皿里带油，但明智的童女在灯里带了油\EditorialNote{玛 25}。我们的灯就是善行；经上记载：\BibleQuote{照样，你们的光也当在人前照耀，好使他们看见你们的善行，光荣你们在天之父}\BibleRef{玛 5:16}。我们若不是为了从邻人的奉承中寻求善行的光荣辉煌，而是将它保存在良心的见证中，我们就是在器皿里与灯一同带了油。对于外人所说的一切，我们应当回到灵魂的隐秘处。即使所有人都辱骂我们，良心中不责备我们的人仍是自由的；而即使所有人都赞美，若良心责备他，他也不能自由。因此真理论及\Person{若翰}说：\BibleQuote{你们出去到荒野里要看什么？风摇动的芦苇吗？}\BibleRef{玛 11:7}。这实际上是以否定的方式说的，而非肯定，因为接着又说：\BibleQuote{但你们出去要看什么？穿细软衣服的人吗？看，穿细软衣服的人是在王宫里}\BibleRef{玛 11:8}。因为，虽然按照福音的真理，\Person{若翰}穿着粗毛衣服，但其意义是：那些喜爱奉承和赞美的人穿的是细软衣服。而\Person{若翰}被否认是风摇动的芦苇，因为没有任何人的口舌能动摇他心志的刚毅。至于我们，若因赞美而高升，或因辱骂而沮丧，我们就是风摇动的芦苇。但願此事远离您的心。我知道您勤读外邦人的导师，他说：\BibleQuote{我若还讨人的喜欢，就不算是基督的仆人了}\BibleRef{迦 1:10}。\par

然而，若因此事在您心中生起丝毫的忧伤，我相信这是全能天主仁慈地允许的。因为即使在今生，祂也没有应许祂的选民欢乐的喜悦，而是苦难的苦楚；好使如同药物一般，他们能通过苦杯恢复永恒救恩的甘甜。因为祂说什么？\BibleQuote{世界要欢乐，你们要哀哭}\BibleRef{若 16:20}。怀着什么希望？什么应许？稍后又说：\BibleQuote{我要再见你们，你们的心要欢乐，你们的喜乐没有人能从你们夺去}\BibleRef{若 16:22}。故此祂又对门徒说：\BibleQuote{你们常存忍耐，就必保全灵魂}\BibleRef{路 21:19}。\par

我请求您思考，若无可忍受之事，忍耐将在何处？我猜想，没有一个\Person{亚伯尔}而没有他的弟兄\Person{加音}。因为若善人没有恶人相伴，他们就不能完全为善，因为他们不会被炼净；而恶人的团体正是善人的炼净。\Person{诺厄}在方舟中有三个儿子，其中一个是讥笑父亲的人，他本人虽蒙祝福，却在他儿子身上承受了谴责。\Person{亚巴郎}在娶\Person{刻突辣}之前有两个儿子；然而他的按肉性而生的儿子迫害了按应许而生的儿子\EditorialNote{创世纪第9章}。伟大的导师阐释此事说：\BibleQuote{那按肉性而生的迫害了按圣神而生的，现今也是这样}\BibleRef{迦 4:29}。\Person{依撒格}有两个儿子；但一个属神的人，因他肉性弟兄的威胁而逃跑。\Person{雅各伯}有十二个儿子；但一个生活正直的，被十个弟兄卖到埃及。关于\Person{达味}先知，因为他内有需要炼净之处，就允许他遭受儿子的迫害。真福\Person{约伯}论恶人的团体说：\BibleQuote{我成了龙的弟兄，鸱枭的伴侣}\BibleRef{约 30:29}。主对\Person{厄则克耳}说：\BibleQuote{人子，不信的人和毁灭者与你同在，你住在蝎子中间}\BibleRef{则 2:6}。在十二位宗徒中有一个被弃者，好使有一个人通过他的迫害来试炼那十一位。宗徒之长如此对他的门徒说：\BibleQuote{祂拯救了正义的\Person{罗特}，他因恶人的不法行为及言谈而受压迫。因为他在观看和听闻中为义人，住在那些每日以不义行为折磨义人灵魂的人中间}\BibleRef{伯后 2:7}。宗徒\Person{保禄}也写信给他的门徒说：\BibleQuote{在弯曲悖逆的世代中，你们显在这世界中如同明灯，紧握生命之言}\BibleRef{斐 2:15}。\par

既然我们从圣经的见证得知，在此生中善人不能没有恶人，那么阁下绝不应被愚人的声音所扰乱，尤其当我们因善行而在世上遭遇逆境时，更确信全能的天主，为的是在永恒的赏报中为我们保留圆满的酬报。因此，在神圣的福音中，真理也说：\Quote{几时人们为了我的缘故而用各种坏话诬蔑你们，你们是有福的}\BibleRef{玛 5:11}。为了安慰我们，祂屈尊以祂自己的辱骂为例，说：\Quote{如果人们称家主为\InnerQuote{贝尔则步}，更何况他的家人呢？}\BibleRef{玛 10:25}\par

但有许多人或许赞美善人的生活超过应有的程度；为了避免赞美滋生骄傲，全能的天主允许恶人爆发诽谤和斥责，以便若因赞美之口而在心中生出任何罪恶，能因责骂之口而被扼杀。因此，外邦人的导师见证他继续宣讲\Quote{通过恶名和善名}\BibleRef{格后 6:8}，也说：\Quote{像是欺骗者，却是真诚的}。如果曾有这样的人对保禄加以恶名，称他为欺骗者，那么现今哪个基督徒应以为基督忍受辱骂之言为艰难之事呢？此外，我们知道我们的救赎者的前驱具有多么伟大的德能，他在圣经中不仅被称为比先知更大，甚至被称为天使；然而，正如他的死亡历史所见证，在他死后他的身体被迫害者焚烧。但我们为何说圣人们的这些事呢？让我们谈论圣中之圣本身，即为我们在世上成为人的天主，祂在死前听到被诬蔑为附有魔鬼的指控，死后被迫害者称为欺骗者，当有人说：\Quote{我们知道那个欺骗者说过：三天后我要复活}\BibleRef{玛 27:63}。那么，我们这些罪人，在永生的审判者来临时将要受审判，若那位连作为审判者来临者也忍受了如此之多，不仅在死前，也在死后，我们怎能不必须承受恶人的口舌和手呢？\par

最甘饴最杰出的女儿，我简短地说这些事，免得每当您听到愚人诋毁您的话语时，心中甚至受到一丁点忧愁的触动。但是，既然愚人的这种怨言不能通过缄默的理性平息，我认为若忽略可做之事的实行，便是罪过。因为当我们安抚疯狂的心灵，使他们恢复健康时，绝不应引起他们的反感。因为有些冒犯是完全可以轻视的；但有些冒犯，当能避免而无罪时，便不应轻视，以免因纵容它们而犯罪。我们从神圣福音的宣讲中学到这个；因为当真理说\Quote{不是入口的污秽人，而是出口的污秽人}\BibleRef{玛 15:11}，门徒们回应说：\Quote{您知道法利塞人听了这话，就跌倒了吗？}\BibleRef{玛 15:12}，祂立刻回答：\Quote{凡不是我天父所栽种的，必要拔出来。由他们吧！他们是瞎子，且是盲人的向导}\BibleRef{玛 15:13}。然而，当被要求纳税时，祂先说明不必纳税的理由，紧接着说：\Quote{然而，为免得罪他们，你去海边投下钓钩，拿那首先上来的鱼，开它的口，必得一枚斯塔特。拿去给他们，为我和你}\BibleRef{玛 17:26}。为什么祂允许一种冒犯存在，却禁止引起另一种冒犯？为什么？无非是为教导我们：一方面要轻视那些使我们陷入罪的冒犯，另一方面要尽力减轻那些能无罪平息冒犯。\par

因此，阁下在天主保护下，能十分平静地避开恶人的冒犯。因为他们中的首领，您应主动私下召见他们，说明理由，并谴责他们猜想您所持的某些错误观点。而且，如果（如所说的那样）他们怀疑这种谴责不是出于真心，您甚至可用誓言来确认，声明您并不持有、也从未持有那些观点。也不要觉得这样满足他们有失身份；也不要因您皇室的血统而心中对他们有任何轻视之情。因为我们所有人都是在一个君主的权能下受造，并由祂的血所赎的弟兄。因此，我们决不应在任何事上轻视我们的弟兄，无论他们多么贫穷卑微。\par

的确，伯多禄曾领受了天国的权柄，使他在世上所束缚或释放的，在天上也应被束缚或释放；他步行海上，以他的影子治愈病人，以他的言语击杀罪人，以他的祈祷复活死人。又因圣神的指示，他进入外邦人科尔乃略的家，信徒们便质疑他为何进入外邦人中并与他们一同吃饭，为何又收纳他们领受圣洗。然而，这位充满如此恩宠的宗徒之长，虽以如此奇迹的德能为后盾，却不用权柄，而用理由回应信徒们的抱怨，并依次向他们解释实情：他如何看见一个器皿，像一块布，其中有地上的四足兽、野兽、爬虫和天空的飞鸟从天降下，并听见有声音说：\BibleQuote{\Emph{伯多禄，起来，宰了吃！}}\BibleRef{宗 11:5}；三个人如何来叫他到科尔乃略那里去；圣神如何吩咐他与他们同去；那位在犹太地惯常于受洗后降临在受洗者身上的圣神，如何在洗礼之前就降临在外邦人身上。因为，若他在受信徒责备时，顾念他在圣教会中领受的权威，他本可回答说，羊不该胆敢指责托付给他们的牧人。但若他为回应信徒的抱怨而提及自己的权能，他便不能真正成为温良的导师。因此，他谦卑地说明理由，安抚了他们，甚至提出见证人为自己辩护，说：\BibleQuote{同我去的，还有这六位弟兄。}\BibleRef{宗 11:12}。如果教会的牧者、宗徒之长，独一地行过标记和奇迹，在为自己辩护免受责备时，都不屑谦卑地说明理由，那么我们这些罪人，在因某事受责备时，更应当谦卑地说明理由，以安抚责备我们的人！\par

因为，如你所知，当我住在皇城中我主人们的脚下时，许多被控涉及上述问题的人常来找我。但我声明，我的良心为我作证，我从未在他们身上发现任何错误、任何邪恶，或任何他们被指控的事。因此，我轻看传闻，特意亲切地接待他们，并宁愿为他们辩护，免受控告者的攻击。因为人们控告他们说，他们以宗教为借口拆散婚姻；又说他们主张圣洗不能完全除去罪过；又说，若有人为自己的不义做三年补赎，此后便可恣意妄为；又说，若他们被迫说他们诅咒那些被指控的事，他们绝不受诅咒的束缚。若真有人坚持并持守这样的观点，毫无疑问他们不是基督徒。凡是这样的人，我与所有天主教主教及普世教会都予以绝罚，因他们思想违反真理，言语也违反真理。因为，若他们说为了宗教应拆散婚姻，须知虽然人的法律准许了此事，但天主的法律却禁止了。因为真理本身说过：\BibleQuote{\Emph{天主所结合的，人不可拆散。}}\BibleRef{玛 19:6}祂又说：\BibleQuote{\Emph{凡休妻的，除非为了淫乱，是不许可的。}}\BibleRef{玛 19:9}谁能反对这位天主的立法者呢？我们知道经上记载：\BibleQuote{\Emph{两人成为一体。}}\BibleRef{玛 19:5}那么，若丈夫和妻子是一体，而丈夫为宗教的缘故休妻，或妻子为宗教的缘故休夫，而对方仍留在世上，甚至或许转向不法之事，这算什么皈依呢？同一肉身，一部分进入节制，另一部分却留在污秽中？但若双方都愿意过节制的生活，谁敢指责他们呢？因为全能的天主既然赐予了较小的恩赐，并不禁止更大的恩赐。事实上，我们知道许多圣人，先前与配偶一起过节制生活，后来转入了圣教会的规矩。因为圣人们惯常以两种方式甚至从合法的事上禁欲：有时是为了在天主面前增加功劳，有时是为了洗去先前生活的罪过。因为当那三个服从巴比伦王的青年要求吃蔬菜，不愿食用王的食物时，并非因为他们吃天主所造之物会犯罪。他们不愿取用他们可合法取用的东西，是为了让他们的德行通过节制而增长。但达味，他曾夺取他人之妻，并因此罪受到严厉的鞭打，后来渴望喝白冷井里的水；当他最勇敢的兵士将水带给他时，他却不肯喝，而将水奠献给上主。因为若他愿意，他本可以合法地喝；但因他记得曾做了不合法的事，他便值得称赞地甚至戒绝了合法的事。他先前毫不畏惧使战死的兵士流血，如今却认为若是喝水，便是流了活兵的鲜血，说：\BibleQuote{\Emph{我岂能喝这些冒死者的血吗？}}\BibleRef{编上 11:19}因此，当好丈夫和好妻子为了增加功劳或消除先前生活的过失，而愿意约束自己过节制的生活并追求更好的生活时，这是合法的。但是，若妻子不随从丈夫所追求的节制，或丈夫拒绝妻子所追求的节制，则婚姻不可拆散，因为经上记载：\BibleQuote{\Emph{妻子没有自己的主权，而是丈夫有；同样，丈夫也没有自己的主权，而是妻子有。}}\BibleRef{格前 7:4}\par

但是，若有人声称在圣洗圣事中罪过只是被表面地除去，还有什么比这种宣讲更违背信德呢？他们借此竟要取消信德的圣事本身——在这圣事中，灵魂主要是与天上洁净的奥秘相结合，以便从所有罪过中完全获得赦免，能紧紧依附于那一位，关于祂先知说：\BibleQuote{但我亲近天主对我是有益的}\BibleRef{咏 72:28}？因为红海的穿越无疑是神圣圣洗的预像：后面的敌人死了，但在旷野中又遇到了其他的敌人。同样，所有在神圣圣洗中受洗的人，他们过去的一切罪过都被赦免了，因为他们的罪过如同埃及的敌人一样死在身后。但在旷野中我们遇到其他敌人，因为当我们还活在此世、尚未到达应许之地时，许多诱惑困扰我们，并在我们前往永生之地的路上竭力阻挡我们。因此，凡说罪过在圣洗中没有完全被除去的人，就让他说埃及人并没有真正死在红海里吧。但如果他承认埃及人真正死了，他就必须承认罪过在圣洗中完全死了，因为真理在赦免我们方面当然比真理的影儿更有力。在福音中，主说：\BibleQuote{洗过澡的人，不需要再洗，而是完全洁净了}\BibleRef{若 13:10}。因此，如果罪过在圣洗中没有完全被除去，那么洗过澡的人怎能完全洁净呢？因为他若仍有任何罪过，就不能说是完全洁净的。但无人能抗拒真理的声音：\BibleQuote{洗过澡的人，是完全洁净了}\BibleRef{若 13:10}。因此，任何罪的污秽都不再留给他，因为救赎他的那一位亲自宣告他是完全洁净的。\par

但是，若有人声称罪过须在三年内做补赎，三年之后就可以纵情享乐，这些人既不懂真信德的宣讲，也不懂圣经的训导。反对这些人，那位杰出的宣讲者说：\BibleQuote{那随从肉情撒种的，必由肉情收获败坏}\BibleRef{迦 6:8}。反对这些人，他又说：\BibleQuote{随从肉性的人，不能得天主的欢心}\BibleRef{罗 8:8}；接着他对门徒们说：\BibleQuote{你们已不属于肉性，而是属于圣神}\BibleRef{罗 8:9}。\par

现今那些生活在肉体享乐中的人就是属于肉性的人。反对他们，经上说：\BibleQuote{朽坏也不能承受不朽坏的}\BibleRef{格前 15:50}。但是，若他们声称短期的补赎足以对付罪过，以致人可以再次回到罪中，那么第一位牧者的判词正击中他们，他说：\BibleQuote{在他们身上正应验了这句恰如其分的俗语：\InnerQuote{狗呕吐的，它又回来吃；洗净的母猪，又到污泥里打滚。}}\BibleRef{伯后 2:22}。因为补赎对罪过确实有大功效，但只限于人若坚持在这补赎中。因为经上记载：\BibleQuote{那坚持到底的，才可得救}\BibleRef{玛 10:22}。因此又记载：\BibleQuote{人从死尸身上洗净，再触摸它，他的洗净有什么用处？}\BibleRef{德 34:30}。死尸就是指一切邪恶的行为，它把人引向死亡，因为人不在正义的生活中生活。因此，一个人从死尸上洗净，而又触摸它，就是指他哀悼自己记得做过的恶行，但在流泪之后又陷入同样的恶行。因此，从这样的死尸上洗净，对任何灵魂都无益处，如果它重做自己所哀叹的事，并且不通过补赎的哀叹上升到正义的端正。因为真正做补赎，不仅是哀悼所犯的罪，而且要远离所哀悼的事。\par

但是，若有人声称，如果任何人因不得已的情况而诅咒了什么东西，他并不受诅咒之约束，这些人本身就是见证，证明他们不是基督徒。因为他们以为借着虚妄的努力就能解除圣教会的束缚，并且他们也不把圣教会给予信友的赦免视为真实的，如果他们以为她的束缚无效的话。对这样的人不必再辩论，因为他们应当完全被藐视和被诅咒；他们以为能逃避真理之处，就让他们在那里真正被束缚在自己的罪中。\par

那么，若有人以基督徒的名义胆敢宣讲或在自己心中默默持守我们上面所提到的那些错谬观点，这些人无疑曾被我们诅咒过，并且现在也诅咒他们。然而，如同我之前所说的，我在皇城中所遇到的那些来见我的人，我没有发现他们对上述任何一点有错谬，我也认为他们没有任何错谬。因为若有的话，我理应发现。不过，既然有许多信友被不明智的热忱所激动，并且常常在攻击某些人如同异端者时，自己制造了异端，就应当顾念他们的软弱，并且如前所述，应当以理性和温和来安抚他们。因为他们确实像经上所记载的那些人：\BibleQuote{我为他们作证：他们对天主怀有热忱，但不是按照真知}\BibleRef{罗 10:2}。因此，阁下您常在诵读、流泪和施舍中生活，应如我所请求的，以温和的劝告和答复来安抚他们的无知，以便不仅在您自己身上，也在他们身上，找到永恒赏报的荣耀。我极大的爱德促使我对您说这一切，因为我认为您的喜乐就是我的收获，您的忧伤就是我的损失。愿全能的天主以天上的恩宠保护您，并保全我主陛下的虔诚和最虔诚的皇后陛下的安宁，延长您的寿命，以教育小主人们。\par

\SourceRecord{Translated by James Barmby.；Nicene and Post-Nicene Fathers, Second Series；Vol. 13.；Edited by Philip Schaff and Henry Wace.；Buffalo, NY: Christian Literature Publishing Co.,；1898.；<http://www.newadvent.org/fathers/360211045.htm>.}

% structure: p0001 source=h1 target=section reason=page-title

\section{第十一卷 第四十六封书信}

\Emph{致耶路撒冷主教依撒基乌}。\par

额我略致依撒基乌等。\par

与历史的真理相符，当时在洪水之时，方舟之外的人类灭亡，方舟之内的人类得以保存生命，这事实意味着什么？不就是我们现在清楚看见的，即所有不信者都在自己罪恶的波浪下灭亡，而圣教会的统一，如同方舟的严实，将其忠信者保存在信德和爱德之中吗？这方舟确实是用不朽的木材构建的，因为它是由坚强的灵魂、持之以恒于善的灵魂建造的。当某一个人从世俗生活中皈依时，仿佛从山上砍下木材。但按照圣教会的秩序，当一个人被指定看管他人时，就好像方舟是用锯好的木材拼合起来以保存人类生命。而事实上，洪水过后，方舟停在了山上，因为这尘世的腐朽结束后，恶行的洪波过去后，圣教会将在天上之乡安息，如同在高山上。因此，为了建造这方舟，我们阅读了您的兄弟之信后，欣喜地发现您以正确信仰的严谨提供帮助；我们向全能的天主致以极大的感谢，祂虽更换祂羊群的牧者，却使祂曾一次交付给圣教父们的信仰，即使在他们之后，也保持不变。如今这位杰出的宣讲者说：\BibleQuote{除已奠立了的根基，即耶稣基督外，任何人不能再奠立别的根基}\BibleRef{格前 3:11}。因此，凡以爱天主爱近人的心，坚心持守那在基督内的信德者，他便是从圣父那里为自己打下了同一根基，即耶稣基督，天主子和人。因此希望，凡以基督为根基者，善工的建筑也会随之而来。真理本身也亲自说：\BibleQuote{凡不由门进入羊栈，而从别处爬进去的，便是贼，是强盗。但由门进去的，才是羊的牧人}\BibleRef{若 10:1}。稍后祂又加说：\BibleQuote{我是门}。所以，由门进入羊栈的，就是经由基督进入。而经由基督进入的人，是对于同一人类创造者和救赎者思想并宣讲真理，实践他所宣讲的，并以负担沉重的职务为目的而承当治理的最高位置，而非因渴望短暂尊荣的荣耀。他也明智地看管他所承担的羊栈职责，以免要么乖谬的人讲逆耳之言撕裂天主的羊，要么恶毒的灵以诱惑他们纵情于邪恶的享乐来败坏他们。\par

但在这一切事上，愿那为我们的缘故而成为人的那一位教导我们。愿那——祂曾屈尊成为祂自己所造的人——将祂爱之精神注入我的软弱和您的爱德，并在一切谨慎和警醒的注视中开启我们心灵的眼目。\par

但那些信仰纯正的人被晋升至圣职，应当不停地感谢同一位全能的天主，并应常为我们最虔诚、最基督徒的皇帝陛下的寿命祈祷，也为他的最宁静的配偶和他最温和的子孙祈祷。在他们的时代，异端者的口是沉默的，因为尽管他们的心在乖戾意见的疯狂中沸腾，但在正统皇帝的时代，他们不敢说出所持的错误意见；因此我们清楚看到所写的经句应验了：\BibleQuote{他把海水聚集成堆，如同装入瓶内}\BibleRef{咏 32:7}。因为海水被聚集成堆如同装入瓶内，是因为异端者苦涩的知识所拥有的任何错误意见，如今都保留在胸中，不敢公开表达。但您的兄弟灵性上受教，完全阐明了正确的信仰，并彻底说明了应当寻求的事。因此，您的信仰就是我们的信仰。我们持守您所说的，并说出您所持守的。\par

但是，因为在我们耳中听闻，在东方的教会中，没有人获得圣职不是通过贿赂；如果您的兄弟发现确实如此，您应将制止您所辖教会中西满异端的错误作为献给全能天主的首祭。因为，且不说其他，那些靠贿赂而非功绩晋升圣职的人，能成为什么样的人呢？我们知道宗徒之长以何等严厉谴责了这一异端，他第一个对西满宣告了定罪，说：\BibleQuote{愿你的银钱与你一同丧亡！因为你想天主的恩赐可以用银钱买得}\BibleRef{宗 8:20}。我们主天主本身，人类创造者和救赎者，也用细绳做成鞭子，推倒了并赶出圣殿中卖鸽子的座位\EditorialNote{见玛窦福音第21章}。因为在圣殿中卖鸽子，不就是用代价在圣教会中换取那授予圣神的按手礼吗？但卖鸽子的座位被推倒，因为这样的人的司祭职不被视为司祭职。\par

此外，我得知在被称为 \Place{Neas} 的教会中，常与你在 \Place{耶路撒冷} 城的教会发生争执。因此，您的圣德应当仔细考虑一切，温和地纠正一些事情，而对那些无法纠正的事情则要平静地容忍。因为我们清楚地看到，圣教会藉圣咏作者的口所说的：\BibleQuote{\Emph{罪人在我的背上建造}} \BibleRef{咏 128:3}。因为担子是在背上背负的。因此，当我们容忍那些无法纠正的人时，罪人就在我们的背上建造。因为船的舵手，当他考虑到逆风时，通过径直驶过一些波涛来克服它们，但对于那些他预见到无法克服的波涛，则明智地转向避让。所以，因此，愿您的圣德通过遏制一些恶行，容忍另一些恶行来减轻它们，以便在各方面维护居住在 \Place{耶路撒冷} 圣教会中的人们之间的和平。因为经上记载：\BibleQuote{\Emph{你们要追求与众人和睦，并要追求圣德，没有圣德谁也见不到天主}} \BibleRef{希 12:14}。因为在争吵中，灵魂的光明，即善意的光明，被阻塞了。因此圣咏作者说：\BibleQuote{\Emph{我的眼睛因愤怒而昏花}} \BibleRef{咏 6:8}。如果我们从心中失去了和平，而没有和平我们就不能见到主，那么在我们身上还能剩下什么善行呢？因此，你要如此行动，甚至要从那些因争执可能使其丧失的人那里，获取你赏报的利益。愿全能的天主以天上的恩宠护守你的爱德，并赐予你从你所受托的人那里带走丰硕的果实和满溢的尺度，直达永恒的喜乐。\par

\SourceRecord{Translated by James Barmby.；Nicene and Post-Nicene Fathers, Second Series；Vol. 13.；Edited by Philip Schaff and Henry Wace.；Buffalo, NY: Christian Literature Publishing Co.,；1898.；<http://www.newadvent.org/fathers/360211046.htm>.}

% structure: p0001 source=h1 target=section reason=page-title

\section{卷十一，书信四十七}

\Emph{致\Place{君士坦丁堡}的执事\Person{阿纳托利乌斯}。}\par

\Person{额我略}致\Person{阿纳托利乌斯}等。\par

您的爱德写信给我说，我们最虔诚的主上命令，为我在\Place{第一犹斯提尼亚纳}的最可敬的兄弟\Person{若望}主教任命一位继任者，因为他患有头部疾患，以免该城若无主教管辖而遭敌人毁灭——但愿天主不让此事发生。然而，教会法规绝无因疾病而免去主教职务的规定。而且，若身体出现病患，便剥夺病人的职务，这是完全不公正的。因此，此事绝不能通过我们来做，以免因罢免他而使我灵魂犯罪。但可以建议，若执政者患病，可找一位管理者，承担其全部职责，维持并代理其地位，而不必罢免他，在教会的治理和城池守护上；这样既不冒犯全能的天主，也不致使城市被疏忽。然而，若同一位最可敬的\Person{若望}因疾病而请求免除主教职务的尊荣，则在他提交书面请求后应予同意。但若非如此，我们出于对全能天主的敬畏，完全不能做此事。但是，若他不愿如此请求，那么最虔诚的皇帝所喜欢的事，他下令要做的，都在他的权力之内。就按他决定的，让他去安排。只求他不要让我们卷入对这样一位处境的人的罢免。尽管如此，他所做的，若是符合教会法规，我们便遵守。若不符合，我们将在不使自己犯罪的前提下，尽可能忍受。\par

\SourceRecord{Translated by James Barmby.；Nicene and Post-Nicene Fathers, Second Series；Vol. 13.；Edited by Philip Schaff and Henry Wace.；Buffalo, NY: Christian Literature Publishing Co.,；1898.；<http://www.newadvent.org/fathers/360211047.htm>.}

% structure: p0001 source=h1 target=section reason=page-title

\section{卷十一，信函第五十}

\\Emph{致\\Person{阿德里安}，公证人。}\par

\\Person{额我略}致\\Place{帕诺尔穆斯}的公证人\\Person{阿德里安}。\par

呈送此信的\\Person{阿加托萨}控告说，她的丈夫违背她的意愿，在\\Person{乌尔比库斯}院长的隐修院中出家了。由于此事无疑触及该院长的信誉和名声，我们命令你的阁下勤勉调查，查明该男子的出家是否得到了她的同意，或者她本人是否承诺过改变身份。如果查明是这样，就让他留在隐修院，并强迫她按承诺改变身份。但如果两种情况都不成立，而且你发现上述妇女并未犯有任何淫乱之罪（因这罪男人可以合法休妻），那么，为了防止他的出家可能成为被留在世俗的妻子丧亡的缘由，我们希望你，不许有任何借口，将她的丈夫归还给她，即使他已经剪发。因为，尽管世俗法律宣称，为了出家可以不顾一方的意愿而解除婚姻，但神律却不允许这样做。因为，除了因淫乱的缘故，男人绝不允许休妻，因为夫妻通过婚姻的结合已成为一体，不能一部分出家，一部分留在世俗。\par

\SourceRecord{Translated by James Barmby.；Nicene and Post-Nicene Fathers, Second Series；Vol. 13.；Edited by Philip Schaff and Henry Wace.；Buffalo, NY: Christian Literature Publishing Co.,；1898.；<http://www.newadvent.org/fathers/360211050.htm>.}

% structure: p0001 source=h1 target=section reason=page-title

\section{第十一卷，第五十四封信}

\Emph{致高卢主教德西德里乌斯}。\par

\Person{额我略}致\Person{德西德里乌斯}等。\par

关于您的行事，有许多美事传至我们耳中，我们心中充满喜乐，以致不忍拒绝您弟兄所请求赐予之事。但后来我们得知一事——提及此事我们不禁羞愧——即您弟兄惯于向某些人讲授\OriginalTerm{文法}{grammar}。此事我们深为不满，极不赞成，乃至我们将先前所言转为叹息与忧伤，因为对\OriginalTerm{基督}{Christ}的赞颂与对\OriginalTerm{朱庇特}{Jupiter}的赞颂不能同出一口。您自己思量，主教们歌唱甚至一位虔敬的平信徒也不适宜歌唱的东西，这是何等严重而可憎的罪过。尽管我们至爱的儿子坎迪杜斯司铎来到我们这里时，曾就此受严格查问，予以否认并尽力为您辩护，但此念仍未从我们心中消去：因其与司铎相关的行为如此可憎，故当以严格而真实的证据查明是否确有其事。因此，若此后所报之事显然虚假，且确知您并未沉溺于无聊的世俗文辞，我们将感谢我们的天主，因祂未容许您的心沾染那可憎的亵渎赞颂；我们将毫无疑虑、毫无踌躇地处理您所请求赐予之事。\par

我们郑重向您推荐这些隐修士，他们与我们至爱的儿子劳伦提乌斯司铎及梅利图斯院长一同被派往我们最可敬的弟兄及同为主教的\Person{奥斯定}处，愿藉您弟兄之助，他们的行程不遭耽搁。\par

\SourceRecord{Translated by James Barmby.；Nicene and Post-Nicene Fathers, Second Series；Vol. 13.；Edited by Philip Schaff and Henry Wace.；Buffalo, NY: Christian Literature Publishing Co.,；1898.；<http://www.newadvent.org/fathers/360211054.htm>.}

% structure: p0001 source=h1 target=section reason=page-title

\section{第十一卷，第五十五封}

\Emph{致\Place{阿爾勒}主教\Person{維吉留斯}}。\par

\Person{格列高利}致\Person{維吉留斯}等。\par

既然聖經的見證稱貪婪為朝拜偶像，那麼應當以何等熱誠將它從天主的殿中驅逐出去，這是眾所周知的；然而（我們含淚而言）有些司鐸卻不重視此事。因為猛烈的貪慾俘虜了人心，說服人相信它所命令的是合法的，如此行進，以同一把劍同時刺殺給予者與接受者。那麼，如果天主的教會因敗壞的司鐸而向貪婪敞開，此後還有什麼安全的地方能抵擋貪婪呢？他如何能保持羊棧不受侵犯，卻引狼入室呢？可恥啊！他以不合法的賄賂玷污自己的雙手，卻想要以自己的祝福提升他人，而自己卻因自己的不義而俯伏，並被自己的野心所俘虜。既然這掠奪的惡行從未進入你心靈的堡壘，你宣稱在聖職授任之事上你的雙手純潔無瑕，就應感謝全能的天主，並承認你是祂的債務人，因為在祂的保護下，你未被這疾病的傳染所傷害。但是，如果你未曾謹慎地禁止他人也有此事，你身上的這善行對你的益處就會少於它本可成就的。正如你自己已厭惡這邪惡，你也應對你的弟兄懷有熱誠去反對它。因為，既然天主誡命勸告我們要愛人如己，那麼忽視它們，不為他人害怕我們自己所退避的事，就是不小的過錯。因此，最親愛的弟兄，現在就專心修補你因糾正上的疏忽而在他人身上所失去的，並盡你所能阻止任何人陷入這邪惡；堅持召開一次會議以根除這同一異端，好使你的愛德獲得賞報，凡經大眾決議所譴責的，賴天主的恩准，也將由眾人更好地防範。\par

此外，我們聽說我們的弟兄兼同作主教的\Place{馬賽}的\Person{塞雷努斯}接納惡人進入他的親密圈子，以至於最終有一位長老成為他的好友；據說這位長老在墮落之後仍然沉溺於他的不義之中。你應仔細查問此事。如果證實如此，你要注意代替我們糾正此事，好使那位接納這樣一個人的人學會不要因親密而鼓勵他，反而要以懲罰來約束他；而那位被接納的人也要學會以眼淚洗淨自己的罪過，不要因不潔的生活而堆積罪惡。請你的弟兄情誼全面照顧我們派往我們的弟兄兼同作主教的\Person{奧古斯丁}那裡的隱修士們，並努力幫助他們上路，與他們協作，好使他們藉著你的援助，在天主的保護下，能夠迅速到達目的地。\par

\SourceRecord{Translated by James Barmby.；Nicene and Post-Nicene Fathers, Second Series；Vol. 13.；Edited by Philip Schaff and Henry Wace.；Buffalo, NY: Christian Literature Publishing Co.,；1898.；<http://www.newadvent.org/fathers/360211055.htm>.}

% structure: p0001 source=h1 target=section reason=page-title

\section{第十一卷，第五十六封信}

\Emph{致\Person{艾瑟里乌斯}，\Place{里昂}主教}\par

\Person{额我略}致\Person{艾瑟里乌斯}，\Place{高卢}主教\par

您信函的措辞，充满可敬的庄严，如此吸引了我们的内心情感，以至于我们乐于不断彼此交谈，以致若我们不能享受您身体的临在，只要这种交流在我们之间继续，缺席对我们来说就没有差别。因为您身上闪耀着对教会秩序何等大的爱，对纪律何等大的尊重，对遵守健康法规何等大的热忱，这从您谦卑且完全乐意地接受我们的劝诫，并宣布将不可侵犯地遵守它表现出来。既然您怀着一颗乐于纠正他人的心，并以合宜的自由声音谴责一个长期存在的邪恶，而且看到我们其他的弟兄和同僚主教也持同样的态度，您有责任一致起来反对主的仇敌，并通过主教会议的决议将贪婪从天主的殿中驱逐出去。在授予圣职时，不要让对黄金的强烈贪求得到任何满足；不要让谄媚窃取任何利益；不要让偏袒给予任何东西：让一个人的生命获得荣誉的回报，他的谦逊促进他的晋升；这样，当这种遵守获得时，既寻求通过贿赂升迁的人被认为不配，而品行良好见证的人得到应有的荣誉。让这成为您的关注，至爱的弟兄，让这种焦虑时刻守护您的思想，以便您通过行动证明您信函中显示的热忱是您内心的见证。因此，不断且立即推动召开主教会议；并如此热忱地履行您的职责，以符合您职位尊严的方式管理您的职务。\par

关于您请求基于古老惯例授予您教会的事项，我们已命人在我们的档案中查找，但没有找到任何东西。因此，请将您声称拥有的那些信函寄给我们，以便从中我们可以得知应该授予您什么。\par

至于真福\Person{依肋乃}的言行录或著作，我们现已长期搜寻，但至今未能找到任何一部分。\par

此外，请您的兄弟之谊注意，务必将我们派往我们的弟兄和同僚主教\Person{奥斯定}的隐修士们视为完全托付给您，并为了天主的缘故向他们显示您的爱德；并如此热忱地与他们同心协力，以司铎般的热情，并如此迅速以您的帮助协助他们继续旅程，以便在您所辖地区没有任何延宕的原因阻碍他们，他们可以更快地前进，而您也会因您为他们所做的一切获得报酬。于七月十日发出，第四小纪年。\par

\SourceRecord{Translated by James Barmby.；Nicene and Post-Nicene Fathers, Second Series；Vol. 13.；Edited by Philip Schaff and Henry Wace.；Buffalo, NY: Christian Literature Publishing Co.,；1898.；<http://www.newadvent.org/fathers/360211056.htm>.}

% structure: p0001 source=h1 target=section reason=page-title

\section{Book XI, Letter 57}

\Emph{致\Person{阿雷吉乌斯}，\Place{瓦平库姆}主教}。\par

\Person{额我略}致\Place{高卢}主教\Person{阿雷吉乌斯}。\par

在兄弟之爱中既然一心一意，心灵因他人的顺遂而喜乐，也因他人的逆境而忧伤，因为按爱德之律，它必分担两者。因此，您的忧伤曾使我们格外悲伤，唯恐长久的悲痛之苦持续击打您的心，并以呻吟加重您的生命。但收到您爱德的来信后，我们得以所盼望的喜乐得安慰，并感谢全能的天主，因我们现在知道您的平静安然无恙，您的心已恢复安慰。确实，对您而言，我们本不期望别的情况，您必以司铎的忍耐战胜任何逆境。\par

此外，我们清楚记得，您的兄弟之情往昔在铲除\OriginalTerm{买卖圣职的异端}{simoniacal heresy}时曾如何炽燃。因此，我们恳切劝勉您对此认真留意，并且，在我们所写的其他事项中，应以公会议之严格定义予以谴责；如此，藉您关怀之助，我们的愿望得以实现，您即可在纠正恶习上向全能天主献上最中悦的祭献，并为他人的建树显示牧职之光辉在您身上彰显。此外，我们对您生活之体验，已知远超众人，促使我们在此事上预期您的大力协助。故此，请完成您的善举，如同您在天主内已开始的那样，以使您心中以正确目标开始的善行，藉万物造主天主的助佑，得以完满。\par

再者，请您的兄弟之情将惯有的爱德施予我们派往我们最可敬的弟兄及同僚主教\Person{奥古斯丁}处的隐修士们；并尽力帮助他们继续行程，无论是亲自还是藉他人，尽您所能，以致因您的供应，他们在您所辖地区无困难或耽搁，我们即可感到对您的信赖并非徒然，全能的天主也必将因他们被派往拯救之灵魂而赐予您恩宠的赏报。\par

\SourceRecord{Translated by James Barmby.；Nicene and Post-Nicene Fathers, Second Series；Vol. 13.；Edited by Philip Schaff and Henry Wace.；Buffalo, NY: Christian Literature Publishing Co.,；1898.；<http://www.newadvent.org/fathers/360211057.htm>.}

% structure: p0001 source=h1 target=section reason=page-title

\section{第十一卷，第五十八封信}

\Emph{致高卢诸位主教。}\par

\Person{额我略}致\Person{泰洛纳}（\Place{土伦}）的\Person{梅纳斯}、\Person{马西利亚}（\Place{马赛}）的\Person{塞雷努斯}、\Person{卡比洛努姆}（\Place{索恩河畔沙隆}）的\Person{卢普斯}、\Person{梅特}（\Place{梅斯}）的\Person{艾古尔夫}、\Person{巴黎西}（\Place{巴黎}）的\Person{辛普利修斯}、\Person{罗托尼乌斯}（\Place{鲁昂}）的\Person{梅兰提乌斯}，以及\Person{李锡尼}，诸位法兰克主教。\OriginalTerm{同等}{A paribus}\par

虽然你们所肩负的牧职之责任提醒你们的弟兄之爱，应当尽一切努力协助热心人士，尤其是那些为灵魂劳苦的人，然而，我等以信函劝勉亦非多余，以激发你们的勤勉，因为如同火焰因风吹而更旺，善意的意图亦因赞美而进步。由于藉着我们救赎者的恩宠合作，如此众多的盎格鲁民族群众正归向基督信仰的恩宠，以至于我们至为可敬的共同兄弟、同为主教的\Person{奥斯定}声称，与他同在的人不足以在各地完成这项工作，我们已安排派遣几位隐修士，由我们最亲爱的共同儿子、司铎\Person{老楞佐}和院长\Person{梅利托}带领。因此，请你们的弟兄之爱向他们显示相称的仁爱，并迅速在他们所需之处提供援助，使他们因你们的协助而在你们地区不致有所延误，不但他们自己因你们的安慰而欢喜，而且你们因提供援助，将在他们被派遣所促进的事业中成为有份者。\par

\SourceRecord{Translated by James Barmby.；Nicene and Post-Nicene Fathers, Second Series；Vol. 13.；Edited by Philip Schaff and Henry Wace.；Buffalo, NY: Christian Literature Publishing Co.,；1898.；<http://www.newadvent.org/fathers/360211058.htm>.}

% structure: p0001 source=h1 target=section reason=page-title

\section{第十一卷，第五十九封}

致\Place{法兰克}王\Person{狄奥多里克}\par

\Person{额我略}致\Person{狄奥多里克}等。\par

陛下之函，诚为心灵的写照，以清晰的言辞之流，彰显出您与王权相伴的卓越智慧，以致一切关于您的赞扬传闻，其真实性无可置疑。既然您信中以赞扬表明，我们的劝勉如此令陛下心满意足，以至于您愿意凡所知关乎我们天主的敬礼、对教会的尊敬或对司铎的尊荣之事，都能被谨慎确立并全面守护，为此，为求您更大的赏报，我们再次劝勉，请您召集一次主教会议，并如我们先前所写，命司铎的肉体恶习与\OriginalTerm{西满异端}{simoniacal heresy}的邪恶，经由所有主教的一致裁决予以谴责，并在您的王国内予以铲除，决不容许金钱再比主的诫命更有效力。因为一切贪财都是偶像崇拜，凡对此不警醒提防者，尤其在授予教会尊荣时，即使似有信仰，若轻视信仰，则必陷于不信的丧亡。因此，如同对抗外在仇敌，同样也要在你们中间对抗灵魂的仇敌，请你们切实留意，因你们对天主仇敌的这项忠信对抗，你们在此世能在祂保护下兴盛治世，并且来世能在祂恩宠引领下得享永恒喜乐。\par

此外，陛下赠予我们最可敬的弟兄及同为主教的\Person{奥斯定}，在其前往\Place{安格利民族}途中之恩惠，已有从他那裡返回的几位隐修士告知我们。为此，我们深表感谢，并恳求您也愿对此批派往他处的隐修士给予充分支援，并助其前行，使他们所侍奉的全能天主，因您愈慷慨对待他们，而赐予您愈大的回报。\par

\SourceRecord{Translated by James Barmby.；Nicene and Post-Nicene Fathers, Second Series；Vol. 13.；Edited by Philip Schaff and Henry Wace.；Buffalo, NY: Christian Literature Publishing Co.,；1898.；<http://www.newadvent.org/fathers/360211059.htm>.}

% structure: p0001 source=h1 target=section reason=page-title

\section{卷十一，书函六十}

\Emph{致法兰克王泰奥德贝尔特}。\par

格列高利致泰奥德贝尔特等。\par

凡以乐意之心接受父爱的劝言，并将其拥抱在心怀中的，无疑表明自己是愿意改正过错的人。正因如此，陛下绝对的许诺已足以使我们安心。因为我们视一个能守信之人的言语为担保。所以，请陛下惠赐，坚守我们天主的诫命，殷勤关注召集一次主教会议，以使司铎中所有形体的恶习及买卖圣职异端——这异端最早是因邪恶的野心而在教会中兴起的——能因陛下权力的谴责之威胁，通过会议的决议而被铲除，并被连根拔起；否则，若在你们的地区，金子比天主更受爱戴，那么那位现在因祂的戒律被藐视而保持安静的天主，日后必会在复仇中被感受到愤怒。事实上，因为我们说这一切都是为了你们自身的利益，所以我们不断地再三敦促你们，以求能甚至通过恳切请求，为我们最卓越、最亲爱的儿子们谋福利。因为如果那些地区违反天主的事能因陛下的改正而得到纠正，这将完全有益于你们的王国。\par

此外，关于陛下在我们最可敬的兄弟及同为主教的奥斯定往盎格鲁民族途中给予的善意服务，我们从某些从奥斯定那里返回我们这里的隐修士的报告中得知。为此我们向陛下致以最大的感谢，并请求您慷慨地将恩惠施予这些信件的携带者——隐修士们，我们已派遣他们前往我们的上述兄弟处，为的是在他们的旅程中，当他们在您庇护下，在你们的地区不遇困难，而是依靠基督的帮助轻易完成所负行程时，您可以在我们天主面前收获更丰厚的赏报果实。\par

\SourceRecord{Translated by James Barmby.；Nicene and Post-Nicene Fathers, Second Series；Vol. 13.；Edited by Philip Schaff and Henry Wace.；Buffalo, NY: Christian Literature Publishing Co.,；1898.；<http://www.newadvent.org/fathers/360211060.htm>.}

% structure: p0001 source=h1 target=section reason=page-title

\section{第十一卷·第六十一封信}

\Emph{致\Place{法兰克}王\Person{克洛塔尔}}。\par

\Person{格里高利}致\Person{克洛塔尔}等。\par

在您为治理治下民众而承受的如此众多的忧虑与焦虑之中，您堪受极大的赞誉与丰厚的赏报，因为您帮助那些为天主的事业而劳苦的人。既然您已借着已完成的善行表明自己如此，使我们能对您怀有更美好的期望，我们便极其乐意请求您那将为您自己赢得赏报的事。如今，某些与我们至为可敬的弟兄及同为主教的\Person{奥古斯丁}一同前往\Place{安格利}民族的隐修士已返回，并告诉我们，当我们的这位弟兄在您面前时，陛下以何等大的爱德款待了他，他离开时您又以何等支持援助了他。但是，既然那些不从已开始的善行退缩者的工作，在我们天主的眼中是悦纳的，我们以慈父般的深情向您致意，恳请您将这些书信的携带者——我们连同我们最钟爱的儿子，即司铎\Person{劳伦提乌斯}与院长\Person{梅利图斯}一起派往我们前述弟兄处的隐修士们——视为特别托付给您的人。您从前向他显示的任何善意，也请您向他们显示，使您的赞誉更加丰富，好使他们因您的供应而能无延迟地完成他们已开始的旅程，全能的天主便会成为您善行的报答者，在顺境中作您的护卫，在逆境中作您的援助者。\par

再者，我们听说，在您的地区，神圣的品级是以金钱贿赂授予的。我们极其忧虑，天主的恩赐若不是因功绩而得，而是靠贿赂强夺。因这\OriginalTerm{西蒙尼异端}{simoniacal heresy}，是首先在教会中兴起的，曾为宗徒的权威所谴责，我们为您的赏报恳请您召集一次主教会议，好使这异端被所有司铎的决议压制并根除，从此在您的地区不再有能力危害灵魂，也绝不允许将来以任何借口兴起。如此，全能的天主必会按照祂看到您在履行祂的命令时所怀的热忱，以及您对那些因这恶行的利剑而濒临丧亡的灵魂的救恩所怀的关切，而高举您以对抗您的仇敌。\par

\SourceRecord{Translated by James Barmby.；Nicene and Post-Nicene Fathers, Second Series；Vol. 13.；Edited by Philip Schaff and Henry Wace.；Buffalo, NY: Christian Literature Publishing Co.,；1898.；<http://www.newadvent.org/fathers/360211061.htm>.}

% structure: p0001 source=h1 target=section reason=page-title

\section{第十一卷，第六十二封书信}

\Emph{致法兰克王后布吕尼希尔德}。\par

额我略致布吕尼希尔德等。\par

我们感谢全能的天主，祂在赐予陛下的一切其他慈恩中，使您如此充满对基督信仰的热爱，以至于您以虔诚的心和热诚的孝爱，不断将一切您所知为灵魂得益、为信德传播之事付诸实施。至于您以何等大的恩宠和帮助扶持我们最可敬的弟兄和同僚主教奥斯定前往盎格鲁人地区的行程，早已有名声传扬；之后，有些从他那里返回我们这里的隐修士，也向我们详细报告了此事。\par

诚然，那些较少知道您恩惠的人会对您这些基督信仰的表征感到惊奇；但对我们这些亲身体验的人来说，这些不是惊奇，而是喜乐，因为您通过施予他人而使自己喜悦。如今，我们的救赎者在上述民族的皈依中行了何等样、何等大的奇迹，陛下已然知晓。因此，您应当大有喜乐，因为您所给予的援助在此中占有更大的份，正是由于您的帮助，天主之后，宣讲的言语才得以在那些地区广传。因为帮助他人的善工，就是使之成为自己的善工。然而，为使您的赏报之果日益丰盛，我们恳请您慈祥地赐予您的庇护支持，给这些携带着此信的隐修士——我们已派遣他们与我们最亲爱的儿子们，司铎劳伦斯和院长梅里图斯，一同前往我们上述最可敬的弟兄和同僚主教那里，因为他告诉我们他身边的人手不足——并惠然在所有事上扶持他们，以便经由陛下的良好赞助，他们能获得更大的成功，并在您地区内不会遇到任何延迟或困难；这样，您为爱天主而在此类事业中谦卑地行事，便能向我们的天主为你们和你们最亲爱的孙子们求得祂的慈恩。\par

\EditorialNote{在\WorkTitle{保禄·执事文集}中，日期为七月第十日前一日（六月廿二日），第四小纪。}\par

\SourceRecord{Translated by James Barmby.；Nicene and Post-Nicene Fathers, Second Series；Vol. 13.；Edited by Philip Schaff and Henry Wace.；Buffalo, NY: Christian Literature Publishing Co.,；1898.；<http://www.newadvent.org/fathers/360211062.htm>.}

% structure: p0001 source=h1 target=section reason=page-title

\section{卷十一，第六十三封信}

\Emph{致法兰克王后布伦希尔德}\par

格列高利致布伦希尔德等。\par

从天上赐予您的恩赐何其美好，神圣的恩宠又以何等的虔诚充满您，这一切，在您其他德行的种种证明之中，已清楚地向所有人表明：您既以睿智谋略治理野蛮民族的凶悍心灵，且（更值得称颂的是）以智慧装点您的王权。既然您在这两方面超越许多民族，同样也在信仰的纯洁上超越他们，我们对您改正非法之事怀有极大信心。因为您已寄给我们的信函内容，足证陛下如何接纳我们的劝诫，以及您怀有何等热忱渴望完成同样的事。但是，既然那赐予善愿的，也常是善愿的助力，我们相信，祂必以慈爱眷顾引导您的事务，越发顺利，因祂见您勤勉于祂的事业。您行天主之事，天主必行您的事。因此，命人召集一次主教会议，并如我们先前所写，在会议决议中严厉禁止您王国中西摩尼异端之罪。向天主献上祭献，以战胜内在的仇敌，这样，借着祂的助佑，您必能战胜外在的仇敌；并且，根据您对祂仇敌所表现的热忱，您将感受到祂在扶助您时也是如此。此外，请相信我，正如我们从许多经验中得知的，凡与罪一同聚敛的，必在损失中耗散。所以，若您不愿不义地失去什么，就应尽力使任何东西都不是不义得来的。因为在世俗事务中，损失总是源于罪。因此，您若想立于敌对民族之上，若愿在得到天主许可后迅速战胜他们，就当颤抖着接受同一全能天主的诫命，好让祂亲自为您对抗您的仇敌，祂曾在圣经中应许说：\BibleQuote{主必为你们作战，你们应安静。}\BibleRef{出 14:14}\par

\EditorialNote{在保禄·执事的收集本中：日期为某月第十日，小纪4年。在雷米吉亚努斯本中：日期为七月第十日，小纪4年。}\par

\SourceRecord{Translated by James Barmby.；Nicene and Post-Nicene Fathers, Second Series；Vol. 13.；Edited by Philip Schaff and Henry Wace.；Buffalo, NY: Christian Literature Publishing Co.,；1898.；<http://www.newadvent.org/fathers/360211063.htm>.}

% structure: p0001 source=h1 target=section reason=page-title

\section{卷十一，第六十四函}

\Emph{致盎格鲁人的主教\Person{奥古斯丁}}。\par

\Emph{以下开始真福罗马城教宗\Person{额我略}的书信，阐释各种问题，他寄往海外的\Place{萨克森}，致\Person{奥古斯丁}——他亲自派遣其代替自己宣讲。}\par

\Emph{序言}。——我通过至爱的儿子司铎\Person{劳伦爵}和隐修士\Person{伯多禄}，收到了贤弟的书信，其中你勤于就许多问题问询于我。然而，由于我前述诸子发现我苦于痛风，而在他们催促我尽快打发他们离去时，他们获准成行，留下我仍受同样痛苦折磨；我未能如我所应当的，就每一个问题更详尽地答复。\par

\Emph{\Person{奥古斯丁}的第一个问题：}我请问，最真福的父，关于主教们应如何与他们的圣职人员共同生活；以及关于在祭台收取的信友献仪，应分为几份，以及主教在教会内应如何处置它们？\par

\Emph{罗马城教宗圣\Person{额我略}的答复：}圣经——无疑你熟知——作证，尤其是真福\Person{保禄}致\Person{弟茂德}的书信，在其中他悉心教导他应如何在天主的家中处世。现在，宗座有一种习惯，在祝圣主教时交付训令，即所有进款都应分为四份：即一份归主教及其家属，用于款待与宴客；另一份归圣职人员；第三份归穷人；第四份用于修缮教会。然而，由于贤弟曾在隐修院规下受训，在盎格鲁人的教会中不应与你的圣职人员分开生活，这教会最近在天主引领下被带入信仰，因此宜于建立那种生活方式，即初期婴儿教会我们教父们的生活方式，其中无人说他所拥有的任何东西是自己的，但一切都公有。\EditorialNote{宗徒大事录 4}\par

\Emph{\Person{奥古斯丁}的第二个问题：}我愿受教，不能节制的圣职人员是否可以结婚；如果结婚，他们是否应回到世俗？\par

\Emph{真福教宗\Person{额我略}的答复：}然而，如果有任何不属于神品的圣职人员不能自制，他们应娶妻，并分别领取他们的俸禄，因为我们知道，关于我们先前提及的那些教父们记载着，每人按需分配。因此应当考虑并为他们的俸禄作出预备，并使他们受教会的规矩约束，使他们度良好的生活，致力于咏唱圣咏，并借天主的助佑保守他们的心、舌和身体远离一切不法之事。至于那些度团体生活的人，我们还需要说什么关于分配份子、或款待、或施舍呢？因为超出他们所需之余，应花费于虔诚和宗教用途，正如我们共同的主与导师所说：\BibleQuote{把余剩的施舍给人，看，一切对你们都是洁净的}\BibleRef{路 11:41}？\par

\Emph{\Person{奥古斯丁}的第三个问题：}既然只有一个信仰，为什么教会的礼仪如此不同？在罗马教会持守一种弥撒礼仪，在高卢教会则持守另一种？\par

\Emph{真福教宗\Person{额我略}的答复：}贤弟知道罗马教会的礼仪，你受培育于其中。但我赞同你仔细挑选任何你发现可能更令全能天主喜悦之事，无论是在罗马教会或高卢教会，或任何其他教会，并以一项特别规定，引入盎格鲁人的教会——它尚在信仰中为新生者——将你能从许多教会中收集的付诸实践。因为我们不应为地方而喜爱事物，而应为事物而喜爱地方。因此，从各个教会选择那些虔诚、敬神、正义之事，并将它们如同集为一束，植入盎格鲁人的心灵，供其使用。\par

\Emph{\Person{奥古斯丁}的第四个问题：}请教我，如果有人可能以偷窃从教会拿走任何东西，应受什么惩罚？\par

\Emph{真福教宗\Person{额我略}的答复：}在此情形下，贤弟可考虑，关于偷窃者的具体情况，如何最好地纠正他。因为有些人有资源却行窃，另一些人由于匮乏而在这事上犯罪。因此，有必要用罚款纠正某些人，但用鞭打纠正另一些人，有些更严厉，有些较轻微。并且，当有人受较严厉处置时，应以爱德处置，而非愤怒；因为受纠正者本人所受的惩罚是定下的，免得他被交付地狱之火。因为我们应当如此对信徒维持纪律，正如好父亲惯常对待儿子们：他们既为过失以棍棒责打他们，却又设法使那些受疼痛的人成为他们的继承人，并为他们似乎以愤怒攻击的同一批人保存他们所拥有的。那么，这爱德应存于心中，使得任何事情都不超越理性的规则。\par

你还问，他们应如何归还以偷窃从教会拿走的东西。但愿教会绝不至于在收回其所失去的现世之物时还加增，并试图从损失中获利。\EditorialNote{异文：de damnis, de vanis。见\Person{贝德}。}\par

\Emph{\Person{奥古斯丁}的第五个问题：}请问，两个兄弟是否可以娶两个姐妹，她们与他们在血统上相隔甚远？\par

\Emph{真福教宗\Person{额我略}的答复：}这完全可以做。因为在圣经中找不到任何似乎反对之。\par

\Emph{\Person{奥古斯丁}的第六个问题：}关于信友应与亲属结婚到第几代，以及是否允许与继母和兄弟的妻子结婚？\par

\Emph{真福教宗\Person{额我略}的答复：} 在罗马共和国中，有一条世俗法律允许兄弟姊妹的子女、或两兄弟的子女、或两姊妹的子女彼此结婚。但我们从经验得知，这样的婚姻不能产生后裔。神圣法律禁止揭露亲属的下体。因此，只有信徒的第三代或第四代才能合法地结合。至于我们所说的第二代亲属，务必要彼此远离。但与继母交合是严重的罪行，因为法律上也记载着：\BibleQuote{你不可揭露你父亲的下体}\BibleRef{肋 18:7}。当然，儿子不能揭露父亲的下体；但既然法律上写着：\BibleQuote{二人成为一体}\BibleRef{创 2:24}，那敢于揭露继母下体的人——继母曾与父亲成为一体——实际上揭露了父亲的下体。与兄弟的妻子交合也是被禁止的，因为她通过先前的结合已成为兄弟的骨肉。为此，\Person{若翰洗者}被斩首，并荣受神圣的殉道冠冕。他没有被命令否认\Person{基督}，却因宣认\Person{基督}而被杀；因为我们的主\Person{耶稣基督}曾说：\BibleQuote{我是真理}\BibleRef{若 14:6}；而若翰为真理被杀，便是为\Person{基督}倾流了鲜血。\par

\Emph{奥斯定的第七个问题：} 我请求明示，是否应当命令那些如此肮脏结合的人分离，并拒绝给他们提供圣体祭献？\par

\Emph{真福教宗\Person{额我略}的答复：} 但是，既然在盎格鲁人中，有许多人还在不信时据称已缔结了这样不圣洁的婚姻，他们应在接受信仰时受到劝告，要彼此远离，并使他们明白这是严重的罪。他们应畏惧天主可怕的审判，以免为了肉欲的享乐而招致永罚的痛苦。然而不应因此剥夺他们领受主的身体和血的共融，免得我们似乎要惩罚他们因无知而在圣洗圣事之前所束缚的事。因为在此时期，圣教会热忱地纠正一些事，柔和地容忍一些事，深思熟虑地默许和容忍一些事，以便常常通过容忍和默许来抑制其所反对的。但所有来到信仰中的人都应受到警告，不得胆敢再做任何这样的事；如果有人胆敢这样做，就必须被剥夺主的身体和血的共融，因为正如在无知中做了这事的人，其过失应在某种程度上被容忍，同样，对于那些明知故犯、不怕犯罪的人，则应严厉惩罚。\par

\Emph{奥斯定的第八个问题：} 我请问，如果路途遥远，主教们不易聚集，是否可以在没有其他主教在场的情况下祝圣一位主教？\par

\Emph{真福教宗\Person{额我略}的答复：} 在盎格鲁人的教会中，既然到目前为止你是唯一的主教，你只能在没有其他主教的情况下祝圣主教。因为，当主教们从高卢到来时，他们将作为见证人协助你祝圣主教。但我们切愿你的友爱在英格兰如此祝圣主教，使主教们彼此之间不被遥远距离分隔，以致在祝圣任何主教时，他们没有必须聚集的理由。因为其他一些牧者的在场也极为有利；因此他们应当尽可能容易地聚集。因此，当天主恩准，在相距不远的地方祝圣了主教之后，祝圣主教决不可没有三、四位主教聚集而举行。因为在属灵的事物本身中，为了能明智而成熟地安排，我们甚至可以从属血肉的事物中取得范例。的确，当世界上举行婚礼时，会召集一些已婚的人，使那些在婚姻道路上先行的人也能参与随之而来的喜乐。那么，为何在这种属神的祝圣中——人与天主通过神圣奥秘相结合——不应聚集那些既为被祝圣为主教者的晋升而喜乐、又向全能的主献上祈求祂护佑的祈祷的人呢？\par

\Emph{奥斯定的第九个问题：} 我请问，我们应该如何对待高卢和不列颠的主教们？\par

\Emph{真福教宗额我略的答复：}对于高卢的主教们，我们不给予你任何权柄，因为从我先祖的古代起，\Place{阿尔勒}（\Emph{Arles}）的主教已接受披肩，我们绝不该剥夺他已获得的权柄。因此，如果你的兄弟情谊恰好进入高卢各省，你应与同一位\Place{阿尔勒}的主教合作，使主教们的恶习（若有）得以纠正。倘若他偶然在纪律的热忱上冷淡，须由你兄弟情谊的热忱来激发他。我们也写信给他，吩咐他当你的圣德在高卢时，全心协助你，好使你们一起在主教的品行中，抑制一切有违我们造物主命令之事。然而，你自己无权凭自己的权威审判高卢的主教们，但要以劝导、吸引，并展示你自己的善行作为他们的榜样，从而将恶人的心性塑造成对圣德的关怀。因为法律上记载：\Quote{路过别人已熟的庄稼，不可伸镰刀，只可用手搓麦穗来吃} \BibleRef{申 23:25}。因此，你不能将审判的镰刀放入显然已托付给别人的庄稼中；但你可以通过和善的善举，将主的麦粒从其缺陷的糠秕中剥离，并借着劝告与说服，如同咀嚼一般，将其转化为教会的身体。但凡是须以权威行事之事，应与上述\Place{阿尔勒}的主教一同进行，以免有人轻视古代教父的制度所规定的事。至于所有不列颠的主教，我们将其管理托付给你的兄弟情谊，使无学识者得受教导，软弱者因劝告得以坚强，乖戾者因权威得以纠正。\par

\Emph{奥斯定的请求：}我请求将殉道者\Person{圣西斯笃}的圣髑送给我们。\par

\Emph{额我略的恩准：}我们已照你所请求的做了，为的是那些原先声称在某处敬礼殉道者\Person{圣西斯笃}遗体（你兄弟情谊认为那既非真遗体，也非真圣骸）的民众，能从这至圣且获认可的殉道者获得确定的益处，而不致尊敬不确定的事物。然而依我看来，如果民众所相信的某位殉道者的遗体，在他们中间并无奇迹显扬，且也没有年长者声称从祖先那里听过其受难的经过，那么你所请求的圣髑应如此分别存放：将前述遗体所在之地完全封堵，不让民众离弃确定之物而尊敬不确定之物。\par

\Emph{奥斯定的第十个问题：}怀孕的妇女是否应受洗？或者，她生产后，经过多少天才可进入教堂？又，为免其婴儿猝死，婴儿在出生后多少日可领受圣洗圣事？或者，丈夫在产后多少天可与她行房？或者，她若患妇女病，可否进入教堂或领受圣体圣事？或者，男人与妻子同房后，尚未用水洗净，可否进入教堂，甚至前往圣体圣事的服务（\OriginalTerm{服务/职务}{ministerium}：在\Person{伯达}的著作中为\OriginalTerm{奥迹/圣事}{mysterium}）？所有这一切，我们都应当为粗鲁的盎格鲁民族而得知。\par

\Emph{真福教宗额我略的答复：}我不怀疑你的兄弟情谊曾问过这些问题，而且我认为我已经向你提供了回答。但我相信，你希望你自己能说的和能想的，能得到我回应的确认。因为，为什么怀孕的妇女不应受洗？在全能天主的眼中，肉体的多产并非过错。因为，当我们的原祖父母在乐园中犯罪后，因天主的公义审判，他们失去了所领受的不朽。因此，全能的天主不愿因他们的过错而完全灭绝人类，便为了人的罪剥夺了人的不朽，但在祂仁慈的眷顾中，仍为人保留了生育后代的能力。既然如此，有什么理由将全能天主赐予人类的恩赐，从圣洗圣事的恩宠中排除呢？因为，以为恩宠的恩赐可能与那完全消除一切人类罪的奥迹相冲突，这实在是非常愚蠢的事。\par

至于妇女生产后多少天才可进入教堂，你已从旧约的指示得知：若是男婴，她应离开三十三天；若是女婴，则为六十六天。然而应当知道，这是按奥秘意义理解的。因为，如果她在生产后的同一时刻进入教堂，她并不承受任何罪担。因为肉体的快乐（而非痛苦）才是罪过。而在肉体的交合中，快乐所在；而生产后代则有痛苦和呻吟。因此，甚至对万民的第一位母亲也如此说：\Quote{你将在痛苦中生育儿女} \BibleRef{创 3:16}。因此，如果我们禁止妇女在产后进入教堂，就等于将她的惩罚本身算作罪过。此外，绝不应禁止妇女或她所生的婴儿在面临死亡危险时立刻受洗——即使她就在生产后的同一时刻，或婴儿在出生的同一时刻。因为，对于活着且有分辨能力的人，圣事的恩宠应以极大的审慎来对待；但对于那些处于死亡迫近危险中的人，该恩宠应毫无延迟地施予，免得在寻求施行救赎奥迹的时间时，死亡很快介入，而无法弥补失去的时间。\par

此外，丈夫在妻子所生的孩子断奶之前，不应与她同居。但在已婚人士中兴起了一种恶习，即妇女不屑于亲自哺乳自己所生的孩子，而将孩子交给其他妇女哺乳。这种恶习似乎完全是为了不能克制情欲的缘故，因为她们不愿节制自己，所以不屑于哺乳自己的子女。因此，那些依照恶习将孩子交给别人哺乳的妇女，除非她们的洁净期已过，否则不应与丈夫同房；因为即使没有生育的理由，她们在惯常的经期内也被禁止与丈夫同房；神圣的法律甚至将任何与行经的妇人同房的人处以死刑（\BibleRef{肋 20:18}）。然而，一个处在惯常经期中的妇女不应被禁止进入教会，因为本性的多余不能归咎为她的罪过，而且因她所不情愿忍受的而剥夺她进入教会的权利是不公平的。因为我们知道，那位患血漏的妇人谦卑地来到主背后，触摸了他的衣边，她的疾病立刻离开了她（\EditorialNote{路 8}）。那么，如果一位患血漏的人能可赞美地触摸主的衣服，为什么一个患月经流血的人进入主的教会就不合法呢？\par

但你会说，那妇人是被病痛所迫；而这些妇女是被惯常的经期所困。然而，最亲爱的弟兄，请想一想，我们在这可朽的肉身中所忍受的一切，都是本性的软弱，是由天主公义的审判在罪之后所命定的。因为饥饿、干渴、炎热、寒冷、疲倦，都是本性的软弱。寻求食物以抗饥饿，饮料以解渴，凉风以避暑，衣物以御寒，休息以解乏，这不正是为疾病寻找某些治疗的补救办法吗？因为在女性中，月经流血也是一种疾病。因此，如果那位在软弱状态中的妇人触摸主的衣服做得对，那么为何赐予一个软弱之人的恩惠不应赐给所有因本性缺陷而处于软弱中的妇女呢？\par

此外，在这些日子里，她不应被禁止领受圣体圣事的奥迹。然而，如果她出于极大的尊敬，未敢领受，则应受称赞；但如果她领受了，也不应受判断。因为善于善心的人，即使在没有罪的地方，也会以某种方式承认自己的罪；因为常有事情虽无罪，却出自罪。因此，当我们饥饿时，我们吃饭而无罪，尽管我们饥饿是出自原初之人的罪。因为女性的经期并非罪，既为自然发生；然而，那本性本身已被败坏以至即使没有意志的意向也显为污秽，这是一个出自罪的缺陷，人性由此可以认识到自己因审判而成了什么，以至于那自愿犯下罪的人，要非自愿地承担罪的罪责。因此，女性若考虑自己正处于经期，若因尊敬不敢接近主圣体圣血圣事，她们应因正确的思虑而受称赞。但若出于虔敬生活的习惯，她们对同一奥迹生发爱慕之心，则不应禁止她们，正如我们所说。因为正如在\WorkTitle{旧约}中注重外在的行为，在\WorkTitle{新约}中，则不仅是外表的行为，更是内心的思想受到密切注意，以便以严格的审判加以惩罚。因为法律虽然禁止吃许多东西为不洁，然而主在福音中说：\BibleQuote{不是入口的污秽人，而是从心里出来的，这些才污秽人。}（\BibleRef{玛 15:11}）随即祂解释添加说：\BibleQuote{从心里发出恶念。}（同上 19）由此充分表明，全能天主所显为行为上污秽的，正是从污秽思想的根源所生发的。因此宗徒保禄也说：\BibleQuote{在洁净的人，凡物都洁净；但在那污秽不信的人，什么也不洁净。}（\BibleRef{铎 1:15}）并且立即，为说明这污秽的原因，他接着说：\BibleQuote{因为他们的心思和良心都污秽了。}那么，如果食物对一个心思不污秽的人不污秽，为什么一个妇女以纯洁的心思从本性所受的，要算为她的污秽呢？\par

此外，一个男人与自己的妻子同寝后，除非用水洗涤，不应进入教会；即使洗涤了，也不应立即进入。法律曾命令古时的百姓：男人与妇人同房后，既要用水中洗涤，也要在日落前不进入教会。这可以按属灵的意义理解为：当人的心思在思想中与非法情欲的乐者结合时，就是与妇人同房；除非他心中的情欲之火冷却，他不应认为自己配得上弟兄们的集会，因为他看到自己被邪恶欲望的淫荡所重压。因为，在此事上，不同的民族有不同的观念，一些民族遵守一种做法，另一些遵守另一种做法；然而罗马人自古以来的惯例一直是：男人与自己的妻子同寝后，既要寻求沐浴的洁净，也要虔诚地暂时不进入教会。\par

我们这样说，并非认为婚姻是罪。但是，即使夫妻之间的合法交合也无法避免肉身的快感，因此应当避免进入圣地，因为这种快感本身绝不可能无罪。因为那说\BibleQuote{看，我是在罪恶中受孕，我母亲在罪孽中怀了我}\BibleRef{咏 50:7}的人，并非生于奸淫或通奸，而是生于合法的婚姻。因为他知道自己是在罪恶中受孕，便为生在罪中而哀叹，因为树枝从树根吸收了邪恶的汁液，便将其传递到枝条上。然而，他这些话并非称夫妻的交合本身为邪恶，而只是指交合的快感。因为有许多事情是允许且合法的，但在实行时我们却在某种程度上被玷污；正如我们常常因愤怒而攻击过错，扰乱了自己心灵的平静。尽管所做的事是正确的，但心灵在其中受到扰乱却不可称许。例如，那位说\BibleQuote{我因愤怒而眼睛昏花}\BibleRef{咏 6:8}的人，曾对犯法者的恶行发怒。因为心灵若不安宁，便无法提升到默观之光中，他因愤怒而眼睛昏花，遂感悲伤，因为虽然从高处攻击恶行，但仍不免在默观至高之事时感到混乱和烦扰。因此，他对恶行的愤怒是可称赞的，但这仍给他带来困扰，因为他感到自己在受扰中犯了一些过错。因此，合法的肉身交合应当是为了生育后代，而非为了快感；肉身的结合应当是为了生养儿女，而非满足软弱。那么，如果有人使用自己的妻子不是被快感的欲望所驱使，而仅仅是为了生育后代，那么他是否进入教堂或领受主的身体和血的奥迹，应由他自己判断，因为我们不应禁止那些在火中却不感到灼烧的人领受。但是，当交合行为中不是生育后代的爱，而是快感占主导时，已婚者就应当在他们的交合中有所哀叹。因为神圣的训导虽准许他们此事，但在准许的同时又以恐惧动摇心灵。当宗徒保禄说\BibleQuote{不能节制的，就让他有自己的妻子}时，他随即加添说\BibleQuote{但我说这话，是出于宽免，并非出于命令}\BibleRef{格前 7:7}。因为公正和正直的事是不被宽免的；他所称为宽免的事，显明是过错。\par

此外，应当仔细留意：主在西乃山上，当祂要向下民说话时，首先吩咐百姓要远离妇女。如果在那里，主通过受造之物向人说话时，尚且如此谨慎地要求身体的纯洁，使那些要听天主言语的人不得与妇女交合，那么那些要领受全能的主的身体的人，岂不更应当在自身内保持肉身的纯洁，以免被这不可估量的奥迹的伟大所压垮！因此，司祭也曾对达味关于他的仆人们说，如果他们远离妇女是洁净的，就可以吃供饼；而除非达味首先宣告他们远离妇女是洁净的，否则他们根本不能领受。然而，一个人与妻子交合后用水洗净，仍可领受神圣共融的奥迹，因为按照上述意见，他进入教堂是被允许的。\par

\Emph{奥古斯丁的第十一问：}我也问，在梦遗（即常在梦中发生的那种幻觉）之后，是否有人可以领受主的身体，或者，如果他是司祭，是否可以举行神圣奥迹？\par

\Emph{真福教宗国瑞的答覆：}对于这样的人，旧约的法律（正如我们在上一节已经说过的）确实宣告他为不洁，并且不允许他在傍晚之前或未用水洗净时进入教堂。但是，若有人不仅按字面理解那法律，而且按属灵意义理解，他就会以我们前面所说过的同一理智观念来看待它；也就是说，那在梦中产生幻觉的人，就好比是受到不洁的诱惑，被真实的幻象玷污了思想。而用水洗净，意思是用眼泪洗去思想的罪过。而且，除非诱惑的火已经过去，否则他应当感到自己有罪，直到傍晚为止。\par

但在同一幻觉中，分辨至关重要，因为必须仔细考虑它因何缘故出现在睡者的意念中。因为此事有时来自过饱，有时来自本性之过剩或虚弱，有时来自思虑。确实，当它来自本性之过剩或虚弱时，绝不应惊慌，因为意念更应被同情为无意中遭受此事，而非有意为之。但当贪饕欲在进食时超出度量，以致体液容器被充满时，意念便负有某种罪责，然而尚未达到禁止领受神圣奥迹或举行弥撒祭礼的程度——假若恰逢庆节需要，或因当地无其他司铎而必须由该司铎奉献奥迹。因为，若有其他能执行奥迹者在场，则因过饱所致的幻觉不应阻碍领受神圣奥迹，但我认为，应谦逊地避免祭献神圣奥迹，只要睡者的灵魂未被污秽想象所动摇。因为有些人，幻觉多半如此产生，以致其意念虽在睡梦中，却未被污秽想象所玷污。关于此事，有一种情况表明灵魂本身有罪，甚至不能逃脱自己的判断：即当它忆起身体睡眠时未见任何事物，却仍忆起在身体清醒时落入淫荡。然而，若睡者灵魂中的幻觉源于其清醒时的污秽思虑，则意念之罪自明。因为人可看出那污秽从何根源流出，若他无意中遭受了他曾有意思想的事。但应考虑那思虑是由\OriginalTerm{诱惑}{suggestion}、\OriginalTerm{喜悦}{delight}或罪恶的\OriginalTerm{同意}{consent}而来。因为所有罪的完成有三种方式：即诱惑、喜悦和同意。诱惑来自魔鬼，喜悦来自肉体，同意来自精神；因为就原罪而言，蛇诱惑了\Person{厄娃}，\Person{厄娃}作为肉体喜悦了它，而\Person{亚当}作为精神同意了它。需要极大的分辨力，使意念作为自身之法官，区分诱惑与喜悦、喜悦与同意。因为当恶灵在灵魂中诱惑犯罪时，若随后无喜悦于罪，则绝未犯罪。但当肉体开始感到喜悦时，罪便开始。若陷入有意的同意，则罪便算完成。因此，在诱惑中有罪的种子，在喜悦中有其滋养，在同意中有其完成。并且常常发生的是，恶灵在思想中播下的，肉体将其吸入喜悦，然而意念并不赞同这喜悦。并且，虽然肉体不能无灵魂而喜悦，但意念虽与肉体之欢愉搏斗，却仍以某种方式违背己意被束缚于肉欲之喜悦中，以致它凭理性力量抗议而不赞同，但仍被喜悦束缚，且因其被束缚而深深叹息。为此，天军的那位首领也叹息说：\BibleQuote{我发觉在我的肢体内，另有一条法律，与我理智的法律交战，并把我掳去，叫我隶属于那在我肢体内之罪的法律。}\BibleRef{罗 7:23}。然而，若他已被掳，他就不会战斗。但他确实战斗了，因此他并非被掳。因此他凭理智的法律战斗，而肢体内的法律与之交战。若他如此战斗，他就未被掳。看，人可说是既被掳又自由：就他所爱的义德而言是自由的；就他所不愿忍受的喜悦而言是被掳的。\par

\SourceRecord{Translated by James Barmby.；Nicene and Post-Nicene Fathers, Second Series；Vol. 13.；Edited by Philip Schaff and Henry Wace.；Buffalo, NY: Christian Literature Publishing Co.,；1898.；<http://www.newadvent.org/fathers/360211064.htm>.}

% structure: p0001 source=h1 target=section reason=page-title

\section{\WorkTitle{第十一卷·第六十五封书信}}

\Emph{致盎格鲁人主教奥斯定}。\par

\Person{额我略}致\Person{奥斯定}等。\par

虽然为全能天主劳苦的人，确已保留了天国不可言喻的赏报，但我们仍需将尊荣赐予他们，使他们因报酬而更努力地投身于属灵工作。而且，既然新盎格鲁教会藉同一主的恩惠及你们的辛劳，已被领受全能天主的恩宠，我们在此教会中仅准你在举行弥撒时使用\OriginalTerm{披肩}{pallium}，以使你可在十二个不同地点祝圣主教，隶属你的管辖，以便未来伦敦城的主教由其自己的主教会议祝圣，并从这圣宗座（赖天主的恩宠，我在此服务）领受\OriginalTerm{披肩}{pallium}的尊严。此外，我们希望你派遣一位你认为适合祝圣的主教至约克城；如此，若该城及邻近地区接受天主圣言，他也可祝圣十二位主教，享有\OriginalTerm{教省总主教}{metropolitan}的尊严；因为如果我们生命延续，我们也将在他身上，蒙主恩佑，赐予\OriginalTerm{披肩}{pallium}；但我们仍愿他受你爱德弟兄的管辖。但在你死后，他应凌驾于他所祝圣的主教之上，丝毫不受伦敦主教的管辖。此外，未来伦敦与约克主教之间应有如下尊荣区别：先祝圣者被视为首席。但凡因基督热忱所应行之事，他们应以共同会议、和谐行动安排；他们应在正义之事上同心合意，并完成他们所愿之事，不彼此不和。\par

但你爱德弟兄在我们天主之下，不仅应管辖你所祝圣的主教，以及约克主教所祝圣的，还应管辖不列颠的所有司铎，使他们能从你圣德的言谈和生活中学习正确信仰和良好生活的典范，并在信德和行为上善尽职责，待主所悦之时，得以抵达天国。愿天主保佑你平安，最可敬的弟兄。颁于七月历前第十日，在我主莫里斯·提比略皇帝在位第十九年，同帝执政后第十八年，第四小纪。\par

\SourceRecord{Translated by James Barmby.；Nicene and Post-Nicene Fathers, Second Series；Vol. 13.；Edited by Philip Schaff and Henry Wace.；Buffalo, NY: Christian Literature Publishing Co.,；1898.；<http://www.newadvent.org/fathers/360211065.htm>.}

% structure: p0001 source=h1 target=section reason=page-title

\section{\WorkTitle{第十一卷 第六十六封}}

\Emph{致盎格利王艾迪伯特}\par

额我略致艾迪伯特，等。\par

为此，全能天主拣选善人治理万民，为要藉着他们，将祂慈爱的恩惠赐予一切受他们管辖的人。我们在盎格利民族中便见到此事，陛下被置于其上的目的，就是要藉着赐予你的恩惠，将天上的福泽施予你所治理的民族。所以，光荣的儿子，要谨慎守护你从天上领受的恩宠。速速在你治下的民族中扩展基督信仰，加倍你正直的热忱以劝化他们，摧毁偶像崇拜，拆毁他们的庙宇，以劝勉、恐吓、诱导、纠正和善行的榜样，建立你臣民圣洁的生活；这样，你在地上传扬了祂的名和知识，就必在天上得着祂为你的报偿。因为你若寻求并维护祂在万民中的尊荣，祂就要使你的名号越发荣耀于后世。正如昔日至为虔诚的君士坦丁皇帝，使罗马共和国转离悖逆的偶像崇拜，将它与自己一同屈服于我们的全能主天主耶稣基督，并全心与他的臣民转向天主。因此，那人在赞美中超越了古代王侯的名声，在声望和善行上都胜过前人。所以，如今陛下也当速速将天主圣父、圣子、圣神的知识注入你治下的君王和人民，使你既在声望和功德上超越你种族往昔的君王，并且你在臣民中越是洗除了他人的罪过，在全能天主可怕的审判面前，你就越能安心于自己的罪过。\par

此外，你身边有我们最可敬的弟兄奥古斯丁主教，他精通修道规则，满溢圣经知识，蒙天主的恩宠行善。要欣然听从他的劝诫，虔诚追随，勤加记念；因为你若听从他以全能天主之名所说的话，那么当他为你祈祷时，这位全能的天主也必更快地垂听他。因为倘若你（愿天主不许）轻忽他的话，那时全能的天主岂能为你而垂听他呢？你竟为了天主而轻忽听从的人。所以，要全心以信德的热忱与他结合，并以他由上赐给你的权能援助他的努力，好使你使人在你的王国中接受并持守的那位信仰，祂自己使你成为祂国度的有分者。\par

此外，我们愿陛下知道，正如我们从圣经中全能主的话语所学，现今世界的终结已经临近，圣人的国度即将来临，那国度没有终结。如今这世界的终结临近，许多前所未有的事即将发生：即天气的变化，天上的恐怖，时节反乎常序，战争、饥荒、瘟疫、多处地震。然而这些事不会在我们有生之年发生，但在我们之后，这一切都将接踵而至。因此，你若在你的土地上观察到任何这些事，切莫心慌，因为这些世界末日的先兆被预先遣来，为的是要我们关心自己的灵魂，猜疑死亡之刻，并以善行预备迎接即将来临的审判者。光荣的儿子，我们现在已略述这些事，为的是当基督信仰在你的王国扩展后，我们对你所说的话也能延展得更长，并且我们心中因你民族彻底归化而倍增的喜乐，将使我们的谈论愈加充实。\par

\Signature{我已送你一些小礼物，当你以真福宗徒伯多禄的祝福接受时，它们对你将不会微小。因此，愿全能天主守护并完善祂在你内开始的恩宠，在此世延长你的寿命许多年，并在长寿之后接纳你进入天上家乡的聚集中。愿天上的恩宠保佑阁下安康，\OriginalTerm{我主，吾子}{domine fili}。颁于七月朔日前的第十日，我主至为虔诚的皇帝毛利斯·提比略·奥古斯都在位第十九年，同一位我主执政官任期后的第十八年，第四小纪。}\par

\SourceRecord{Translated by James Barmby.；Nicene and Post-Nicene Fathers, Second Series；Vol. 13.；Edited by Philip Schaff and Henry Wace.；Buffalo, NY: Christian Literature Publishing Co.,；1898.；<http://www.newadvent.org/fathers/360211066.htm>.}

% structure: p0001 source=h1 target=section reason=page-title

\section{第十一卷，第六十七封书信}

\Emph{致奎里库斯主教等}\par

\Person{额我略}致\Person{奎里库斯}主教及\Place{伊比利亚}的其他天主教主教。\par

既然对于爱德而言，没有什么是遥远的，就让那些被地方分开的人通过书信联结吧。这封信的携带者来到真福\Person{伯多禄}宗徒之长的教会，断言他从你们的弟兄团体收到了给我们的信函，但在\Place{耶路撒冷}城连同其他物品一起丢失了。据他说，在这些信中，你们渴望询问关于那些在聂斯多略异端错误中迷失的司铎及信众的问题：当他们返回作为所有蒙选者之母的天主教会时，是否应该受洗，或者仅通过宣认唯一真信仰而联入同一慈母教会的怀抱？\par

实际上，我们从教父的古老制度中得知，凡在异端中因圣三之名受洗者，当他们返回圣教会时，或通过傅圣油，或通过按手，或仅通过宣认信仰，即可被召回到慈母教会的怀抱。因此，西方通过按手将亚略派与圣而公教会和好，而东方则通过傅圣油。但一性论者及其他则仅通过真诚的宣认而被接纳，因为他们在异端中所领受的圣洗，其洁净能力在他们身上获得实效：要么前者通过按手领受圣神，要么后者因其对真信仰的宣认而与圣而公教会的怀抱结合。然而，那些未因圣三之名受洗的异端，例如\Emph{\OriginalTerm{波诺西乌斯派}{Bonosiaci}}和\Emph{\OriginalTerm{卡塔弗里吉亚派}{Cataphrygæ}}，因为前者不信基督主，而后者以乖谬的理解相信某个恶人\Person{蒙达努斯}是圣神，类似这样的人还有许多——这些人来到圣教会时，应受洗，因为他们在错误中所领受的，既然不是因圣三之名，便不是洗礼。这也不能称为重洗，因为如前所述，那原非因圣三之名所施。至于聂斯多略派，既然他们因圣三之名受洗——尽管被其异端错误所蒙蔽，以致仿效犹太不信的方式，不信独生子的降生成人——当他们来到圣而公教会时，应通过坚定持守和宣认真信仰，教导他们相信同一天主子并人子、我们的主天主耶稣基督，同一者在万世之前已存在于天主性中，同一者在末期成为人，因为\BibleQuote{圣言成了血肉，居住在我们中间}（\BibleRef{若 1:14})。\par

但我们说：圣言成为血肉，并非藉失其所是，而是藉取其所非是。因为在其降生成人的奥迹中，圣父的独生子增加了属于我们的，而未减少属于他的。因此，圣言与血肉是一个位格，正如他亲自所说：\Emph{\BibleQuote{没有人升过天，除了那从天上降下来的人子，他是在天上的}}\BibleRef{若 3:14}。他在天上原是圣子，在地上说话的人子。因此若望说：\Emph{\BibleQuote{我们知道天主子来了，并赐给了我们理智}}\BibleRef{若一 5:20}。至于他赐给我们理智，他随即加上：\Emph{\BibleQuote{使我们认识真天主。}}在这里，他指的真天主，除了全能圣父之外，还有谁呢？但他对全能圣子也有同样的认识，他接着说：并且我们是在他真实圣子耶稣基督内。看，他说圣父是真天主，耶稣基督是他的真圣子。他更清楚地表明他所认识的真圣子是什么：\Emph{\BibleQuote{这是真天主，也是永生}}。因此，若按照聂斯多略的错误，圣言是一回事，人耶稣基督是另一回事，那么真的人就不会是真天主和永生。但独生子，即万世之前的圣言，成了人。因此，他是真天主和永生。诚然，当圣童贞即将怀他，听见天使对她说话时，她说：\Emph{\BibleQuote{看，主的婢女，愿照你的话成就于我}}\BibleRef{路 1:38}。当她怀了他，前往亲戚依撒伯尔那里时，立刻听见：\Emph{我主的母亲到我这里来，这是从哪里得来的呢？}看，同一位童贞被称为主的婢女和主的母亲。她是主的婢女，因为万世之前、与圣父平等的独生子圣言；但她也是母亲，因为在她腹中，因圣神和她的血肉，他成了人。她不是一个的婢女、另一个的母亲；因为当日，当万世之前存在的天主的独生子从她腹中成了人时，藉着不可测的奥迹，她因神性而成了人的婢女，因血肉而成了圣言的母亲。并非血肉先在童贞女腹中受孕，然后神性才进入血肉；而是圣言一进入腹中，立即保存自己本性的卓越，成了血肉。天主的独生子，通过童贞女的腹，生为一个完美的人，即拥有真实的血肉和理性的灵魂。因此他也被称为\Quote{在同伴中受傅者}，如圣咏作者所说：\Emph{\BibleQuote{天主，你的天主，以喜乐之油傅了你，胜过你的同伴}}\BibleRef{咏 44:8}。因为他受油傅，即受圣神之恩赐。但他受傅胜过同伴，因为我们众人先作为罪人存在，然后通过圣神的傅油而成圣。但他，万世之前存在为天主，在末世因圣神在童贞女腹中受孕为人，就在那里，即在他受孕之处，受同一圣神的傅油。并非先受孕然后受傅；而是因圣神由童贞女的血肉受孕本身就是受圣神的傅油。那么，关于他的诞生这一真理，所有从聂斯多略的邪僻错误中被领回的人，都应在你团体神圣的会众前承认，绝罚同一聂斯多略及其所有跟随者，以及所有其他异端。他们也应承诺接受并敬拜普世教会所接受的可敬会议。且你的尊座应毫无犹豫地在你的集会中接纳他们，准许他们保留自己的品位，以便你同时细察他们内心的隐秘，通过真正的知识教导他们应持守的正确之事，并以温和对待他们，关于他们的品位不制造困难或矛盾，从而将他们从古仇的口中夺回；并且永恒荣耀的报偿在天主面前更多地加增于你，因你聚集许多将在主内永远与你一同夸耀的人。愿天主圣三在你为我们祈祷时保护你，并在其爱中赐予你更丰盛的恩赐。\par

\EditorialNote{在\Emph{Colbert.}与\Emph{Collect. Paul}中，\Quote{于七月朔前第十日，小纪第四年颁发。}}\par

\SourceRecord{Translated by James Barmby.；Nicene and Post-Nicene Fathers, Second Series；Vol. 13.；Edited by Philip Schaff and Henry Wace.；Buffalo, NY: Christian Literature Publishing Co.,；1898.；<http://www.newadvent.org/fathers/360211067.htm>.}

% structure: p0001 source=h1 target=section reason=page-title

\section{第十一卷，第六十八封书信}

\Emph{致\Person{维吉利乌斯}，\Place{阿雷拉特}（\Place{阿尔勒}）主教。}\par

\Person{格列高利}致\Person{维吉利乌斯}等。\par

对于自愿前来投奔我们的弟兄应给予何等情谊，从通常出于爱德邀请他们来访即可明了。因此，若我们的共同兄弟\Person{奥古斯丁}主教偶然来到你处，愿您的爱德恰如其分地如此亲切温柔地接待他，既以您安慰之恩使他振作，也教导他人如何培养兄弟般的爱德。再者，由于远处之人常先由他人得知需要纠正之事，若他偶然向您的兄弟情谊提及司铎或其他人的过失，您当与他协力，以极细致的调查探究此事。并且，你们二人要如此严格而热忱地抵制冒犯天主、激起祂义怒之事，为使他人改正，既让惩罚击中有罪者，也不让诬告折磨无辜者。愿天主保佑您平安，最可敬的兄弟。颁于\OriginalTerm{七月朔日前的第十日}{10th day of the Kalends of July}，我们最虔诚的主上\Person{玛乌里奇乌斯·提贝里乌斯·奥古斯都}皇帝统治第19年，同一位主上执政官后第18年，\OriginalTerm{第4小纪}{Indiction 4}。\par

\SourceRecord{Translated by James Barmby.；Nicene and Post-Nicene Fathers, Second Series；Vol. 13.；Edited by Philip Schaff and Henry Wace.；Buffalo, NY: Christian Literature Publishing Co.,；1898.；<http://www.newadvent.org/fathers/360211068.htm>.}

% structure: p0001 source=h1 target=section reason=page-title

\section{第十一卷·第六十九书}

\Emph{致法兰克王后布伦希尔德}\par

\Signature{额我略致布伦希尔德等}\par

经上记载：\BibleQuote{正义使邦国高升，罪恶令人民蒙羞。}（\BibleRef{箴 14:34}）由此可见，当已知的过失迅速被纠正时，王国便被认为是稳固的。如今，我们因许多人的报告而闻知——此事我们提及时不能不心怀极大忧苦——某些地区的司铎生活如此不端、行为如此邪恶，以至我们听闻感到羞耻，述说更令人悲痛。唯恐此不义的传闻既已远播至此，他人的恶行会以罪恶之箭刺伤我们灵魂或您的王国，我们当以热忱奋起报复这些事，以免少数人的邪恶成为多数人的丧亡。因为不良的司铎是人民败亡的根源。若司铎本应为人民的罪祈祷，却犯下更重的过犯，谁还能为人民的罪献上代祷？然而，那些本应追究这些事的人，既无调查的热心，亦无惩罚的热忱，故请您赐函于我。若蒙您许可，我们将派遣一位获得您权威同意的人，与其它司铎一同彻底调查此事，并按天主的旨意予以纠正。因为我们所谈论之事绝不可视而不见；能改正过失却忽略不做的人，无疑使自己成为共犯。因此，请您为您的灵魂着想，为您希望幸福统治的孙儿们着想，为各省着想；在我们的造物主伸手击打之前，请以最大的热忱思量此恶行的纠正，免得祂因长久仁慈的等待而日后责打更严厉。此外，您要知道，若您从您的领土上迅速剪除如此大罪的毒害，您将为我们的天主献上一份伟大的赎罪祭献。\par

\SourceRecord{Translated by James Barmby.；Nicene and Post-Nicene Fathers, Second Series；Vol. 13.；Edited by Philip Schaff and Henry Wace.；Buffalo, NY: Christian Literature Publishing Co.,；1898.；<http://www.newadvent.org/fathers/360211069.htm>.}

% structure: p0001 source=h1 target=section reason=page-title

\section{第十一卷，第七十六书}

致院长梅里多\par

额我略致在法国的院长梅里多。\par

自从我们与你们同在的会众离开后，我们因一直未听到你们旅程成功的消息而极为焦虑。但当全能的天主将你带到我们最可敬的兄弟奥古斯丁主教那里时，请告诉他：我长久以来一直在考虑关于\OriginalTerm{盎格鲁人}{Angli}的事务；也就是说，该国偶像的庙宇不应被摧毁，但其中的偶像本身应被摧毁。应制备圣水，洒在这些庙宇中，建造祭台，并安放圣髑；因为，如果这些庙宇建筑良好，有必要将它们从对偶像的崇拜转变为对真天主的服务；这样，当民众看到这些庙宇没有被摧毁时，他们可能会从心中除去错误，并在认识和崇拜真天主时，更熟悉地前往他们习惯的地方。而且，因为他们惯常为祭献魔鬼而宰杀许多牛，他们也应该以改变了的形式举行某种类似的庆典，以便在奉献日或在安放圣髑的圣殉道者周年纪念日，他们可以在这些已改为教堂的庙宇周围用树枝搭建帐篷，并以宗教盛宴庆祝庆典。不再让他们将动物祭献给魔鬼，而是宰杀动物以赞美天主供自己食用，并为饱足而感谢万有的施予者；这样，当外在保留了一些喜乐时，他们可能更容易将心灵转向内在的喜乐。因为从刚硬的心中一下子剪除一切无疑是不可能的，因为一个人若想攀登最高处，必须借助阶梯或步伐，而非跳跃。因此，主在埃及向以色列民显明了自己；但祂在祂自己的崇拜中仍保留了他们惯常向魔鬼奉献的牺牲，命令他们将动物祭献给祂；目的是，他们的心被改变，他们在牺牲中省略一些事物而保留其他，这样，虽然动物与他们惯常奉献的相同，但既然他们是将它们祭献给天主而非偶像，它们就不再是同样的牺牲了。因此，你的爱德有必要对我们上述的兄弟说这些话，以便他现在在那地方，能妥善考虑如何安排一切。愿天主保佑你平安，最亲爱的儿子。于七月月望前第十五日，我们最虔诚的主上莫里斯·提比略·奥古斯都皇帝在位的第十九年，同一位主上执政官之后的第十八年，第四小纪。\par

\SourceRecord{Translated by James Barmby.；Nicene and Post-Nicene Fathers, Second Series；Vol. 13.；Edited by Philip Schaff and Henry Wace.；Buffalo, NY: Christian Literature Publishing Co.,；1898.；<http://www.newadvent.org/fathers/360211076.htm>.}

% structure: p0001 source=h1 target=section reason=page-title

\section{卷十一，书信七十七}

\Emph{致在科西嘉的护守者\OriginalTerm{护守者}{Defensorem}波尼法爵。}\par

额我略致波尼法爵等。\par

你的经验难辞其咎，因你明知科西嘉的城镇阿莱里亚和阿雅克肖久无主教，却迟延劝诫他们的圣职人员和百姓为自己选举司铎。但鉴于他们不应再长期没有自己的牧者，你接到此谕后，应即刻分别劝勉这些城镇的圣职人员和百姓，使他们彼此和衷共济，而非内部分歧，使每一城镇同心合意为己选举一位司铎以受祝圣。待他们做出决定后，被选者应前来见我们。若他们意见不一，无法一致通过，而在两人之间分歧，则两人同样按常例做出决定后前来见我们，以便在查考其生活与品德后，视谁最为合适而予以祝圣。再者，据闻当地许多穷人受欺压且遭受不公，你的经验务请注意此事，不可容许他们受不义之冤；你当尽力而为，既不使起诉者无理受阻，也不使被诉者蒙受不白之冤。\par

此外，我们听说当地一些圣职人员，在你本人在场的情况下，竟被平信徒拘押。若果真如此，须知此过将归于你身，因你若果真是男子汉，此事本不应发生。为此，你今后务必留意，不得允许类似事件发生；反之，若有任何人向圣职人员提出控诉，应让其诉诸该圣职人员的主教。若主教本身可疑，则由主教委派一位专员——或若原告对此也有异议，则由你委派——此人应促使双方同意共同选定仲裁人。无论仲裁人如何裁决，在你或主教的照管下，均应依法全力执行，使之无需再为争讼而烦扰。\par

\SourceRecord{Translated by James Barmby.；Nicene and Post-Nicene Fathers, Second Series；Vol. 13.；Edited by Philip Schaff and Henry Wace.；Buffalo, NY: Christian Literature Publishing Co.,；1898.；<http://www.newadvent.org/fathers/360211077.htm>.}

% structure: p0001 source=h1 target=section reason=page-title

\section{卷十一，信函第七十八}

\Emph{致\Person{芭芭拉}和\Person{安东尼娜}}。\par

\Person{额我略}致\Person{芭芭拉}等。\par

接到你们的书信，我万分高兴得知你们安好。我恳求全能的天主，愿祂以祂的护佑，在思想上保护你们脱离恶灵，脱离邪僻之人，脱离一切逆境；并愿祂以敬畏祂的恩宠，使你们定居于相称的婚姻中，使我们大家为你们的定居而喜乐。但是你们，最亲爱的女儿们，当寄望于祂的助佑，常常在祂护佑的荫蔽之下，藉着祈祷和善行，躲避恶人的阴谋。因为，无论人间的安慰或逆境可能存在，除非祂的恩宠保护你们，或祂的不悦困扰你们，否则都算不得什么。所以，不要寄望于任何人，而要将你们的整个灵魂交托于对全能天主的信赖。当我们睡觉时，祂会保护你们，关于祂曾写说：\BibleQuote{看，那保护以色列的，既不盹睡，也不睡眠。}\BibleRef{咏 120:4}。\par

至于你们说，你们急于走近真福伯多禄宗徒之长的门槛，我极其渴望并热切期待，在他的教会中看到你们与堪当你们的丈夫结合，以使你们从我这里得到些许安慰，而我也从你们的到来得到极大的喜乐。我已将你们的事宜托付给你们至可敬的弟兄\Person{若望}主教，以及护卫者\Person{罗马努斯}（\OriginalTerm{护卫者}{defensori}），愿他们在天主眷顾下完成他们已开始的事。\par

你们赠送的两件\OriginalTerm{衣料}{racanæ}，据你们来信说是你们自己的手工，我欣然接受。但要知道，我并不相信你们所说的话。因为你们是在寻求从他人工作中得来的赞美，既然你们或许从未亲手拿过纺锤。然而，这一情况并不使我难过，因为我希望你们热爱阅读圣经，这样，当全能的天主将你们与丈夫结合时，你们就知道该如何生活，该如何管理你们的家室。\par

\SourceRecord{Translated by James Barmby.；Nicene and Post-Nicene Fathers, Second Series；Vol. 13.；Edited by Philip Schaff and Henry Wace.；Buffalo, NY: Christian Literature Publishing Co.,；1898.；<http://www.newadvent.org/fathers/360211078.htm>.}

% structure: p0001 source=h1 target=section reason=page-title

\section{第十二卷，第一封信}

\Emph{致\Place{迦太基}主教\Person{多米尼库斯}}\par

\Person{额我略}致\Person{多米尼库斯}等。\par

你的心肠何等丰沛，你藉其诠释者——你的舌——显示出来，同时以它的甘甜调和你书信的言语，以致你所写的一切都令人愉快欣喜。因此，我们虽不能在身体上，却在爱德的怀抱中拥抱你的弟兄情谊。因为爱德的职分，是为那些和谐一致的灵魂，供应距离所拒绝给予的。既然我们最可爱的弟兄的疾病令我们忧伤，正如他们的健康令我们欣慰，我们感谢全能天主，祂以好消息安慰了我们的忧伤。因为，在收到你的信之前，我们听说你染上了非常严重的疾病，处于极大的痛苦之中。但是，当我们从死亡的危险中被抢救出来时，最亲爱的弟兄，我们尚不确定自己被保留为何事，让我们将喘息的时间转为灵魂的益处，并且，既要向即将到来的审判者交账，就让我们以眼泪和善工在祂面前坚固我们的案件，好使我们能被认为配得，对于我们所做的事获得平安。因为在世俗案件中，仁慈的法官也常为此目的而准许缓期，使先前没有预备好的人，以后可以准备妥当前来受审。如果我们为了灵魂的得救，竟忽略了我们在世俗事务上小心留意之事，那将成何等事！因此，既然按宗徒\Person{若望}的话，没有人是无罪的，让我们忆起思想的诱惑、舌头的失检、违法的行为；并且，趁我们还能的时候，以大力叩门，除去我们不义的污点，好使我们的公义而慈爱的救赎主不按我们的功过施行报复，而是依祂的仁慈屈身赦免。并且，既然仅仅为自己的罪过哭泣不足以尽我们的职分，让我们更热切地致力于照管托付给我们的羊群，并以劝服、以劝诫、以警醒、以宣讲，按天上的仁慈所赐给我们的能力，赶紧以行动履行我们的职分，好使我们透过造物主的恩惠，期待那盼望的赏报。但是，既然没有天主的助佑我们什么善事也不能做，让我们最亲爱的弟兄，同心合意地祈求全能天主，愿祂藉其恩宠的引领，引导我们和托付给我们的羊群，走上祂诫命的道路，并愿祂自己，因祂仁慈的恩赐，愿意我们拥有牧者的名号，赐予我们领悟并实行中悦祂之事。此外，我们以你寄来时的爱德，收到了由你的圣德转交给我们的真福殉道者\Person{阿基莱乌斯}的祝福。九月，小纪第五。\par

\SourceRecord{Translated by James Barmby.；Nicene and Post-Nicene Fathers, Second Series；Vol. 13.；Edited by Philip Schaff and Henry Wace.；Buffalo, NY: Christian Literature Publishing Co.,；1898.；<http://www.newadvent.org/fathers/360212001.htm>.}

% structure: p0001 source=h1 target=section reason=page-title

\section{\WorkTitle{第十二卷 第八封信}}

\Emph{致努米底亚主教哥伦布。}\par

额我略致哥伦布等。\par

持信人多纳德斯自称曾任执事，其所呈诉状（附录于后）之严重性，即令闻之亦不可容忍，此将由贵弟兄明鉴。然闻彼曾因肉体之罪而被罢免，务请详查此事；若果如此，当令其行补赎，俾以泪水自所犯淫乱之束缚中得释。若证实其未犯此等过犯，则其诉状所陈一切，须由尔会同会议首席及其他我等弟兄同为主教者，以勤勉审查加以探究。若其诉状属实，则对其主教维克托（彼竟未醒悟，犯下如此大恶，得罪天主及其自身司铎职守）须施以教会纪律之严厉制裁，使之明悟己行之恶；而该人自身则当恢复其品级。盖凡未因自身过失或罪行而罢免其所任职位之品级者，却因某人之意旨而无效地被剥夺，此诚悖理，且明违教会秩序。\par

\SourceRecord{Translated by James Barmby.；Nicene and Post-Nicene Fathers, Second Series；Vol. 13.；Edited by Philip Schaff and Henry Wace.；Buffalo, NY: Christian Literature Publishing Co.,；1898.；<http://www.newadvent.org/fathers/360212008.htm>.}

% structure: p0001 source=h1 target=section reason=page-title

\section{第十二卷 第二十四封信}

\Emph{致\Person{若望}，\Place{拉文纳}副执事}。\par

\Person{额我略}致\Person{若望}等。\par

有些隐修士从已故院长\Person{克劳狄}的隐修院来到我这里，请求我任命隐修士\Person{君士坦提乌斯}为他们的院长。但我对他们的请求极为反对，因为他们似乎完全怀着世俗的心思，竟要一个十分世俗的人做他们的院长。因为我得知这位\Person{君士坦提乌斯}一心想要拥有私产：这最有力地证明他没有隐修士的心。我还得知，他竟敢独自一人，不带任何一位弟兄，前往位于\Place{皮切诺}省的一所隐修院。从他的这一行径我们可知，无人见证就独自行走的人，生活不正；他自己都不知如何持守规矩，又怎能维持他人的规矩呢？\par

因此，他们放弃了他，转而请求任命一位叫\Person{莫鲁斯}的管库员，他的生活和勤勉有许多见证，已故院长\Person{克劳狄}也同某些人一道称赞过他。所以，请你的经验仔细查问；若他的生活适宜管理之位，就请我们的兄弟和同为主教\Person{马里尼亚努斯}祝圣他为院长。但若他确有明显的反对之处，而他们在自己的会众中找不到合适的人，就让他们从别处选一人，所选的这人就立为院长。此外，务必告知我们上述的兄弟和同为主教，要极力革除那隐修院中四、五个隐修士持有私产的事（此事至今未能纠正），并尽快清除该隐修院的这种祸患；因为若隐修士在那里持有私产，在这同与会众中便不可能保持和睦与\OriginalTerm{爱德}{charity}。隐修士的生活状态不就是轻视世界吗？那么，身在隐修院却寻求黄金，又如何轻视世界呢？所以，请你的经验如此行事，既不要拖延该地的管理，也不要再让我们听到任何关于此事的申诉。\par

此外，因为我最亲爱的已故之子\Person{克劳狄}曾听我讲论过\BibleBook{}{箴言}、\BibleBook{}{雅歌}、\WorkTitle{先知书}，以及\BibleBook{列王}{纪}和\WorkTitle{七经}的一些内容，我因病未能书写下来，而他根据自己的理解口述并让人记录，以免遗忘，并打算在适当的时候带给我，以便更正确地口述（因为他向我朗读他所写的时，我发现我所说的意思被很不利地改变了）。因此，有必要请你的经验，不找任何借口或拖延，前往他的隐修院，召集弟兄们，让他们充分并真实地呈上他所带去的一切关于不同圣经的文稿；请取来这些文稿，并尽快传送给我。\par

此外，关于你的返回，听闻你遭遇了严重麻烦，我们将稍后考虑。此外，我不悦地从某些人那里听说，我们最可敬的兄弟和同为主教\Person{马里尼亚努斯}在守夜礼中公开诵读我对真福\BibleBook{}{约伯}的注释；因为这部作品并非通俗之作，对粗浅的听众而言，与其说是进步，不如说是障碍。但请告诉他，在守夜礼中诵读对\BibleBook{}{圣咏}的注释，这能陶冶世俗之人的良好品行。因为我确实不愿意在我尚在肉身之时，我所说过的话轻易为人所知。因为我很介意，最可怀念的执事\Person{阿纳托利}应主上皇帝的请求和命令，将\WorkTitle{牧灵规则}一书呈给了他——此书由我最神圣的兄弟和同为主教\Place{安提约基雅}的\Person{阿纳斯塔西}译成希腊语；我通过书信得知，此书令他甚为喜悦；但我甚为不悦的是，那些拥有更好事物的人竟被琐事所占据。\par

此外，在论真福\BibleBook{}{约伯}的第三部分，在经文中写道：\BibleQuote{我知道我的救赎者活着}，我怀疑我们的上述兄弟和同为主教\Person{马里尼亚努斯}有一份错误的抄本。因为在我们书橱的抄本中，这段文字与我在他人抄本中所见的不同；因此我已改正了这段文字，以便我们多次提及的兄弟能照我们书橱的抄本拥有它。因为其中有四个词，若缺少了，可能给读者造成不小的困难。请彻底而迅速地执行这一切。若你对最卓越的总督无能为力，也要表明你自己没有忽视做你力所能及的事。\par

关于\Place{阿尔比努斯}的地方，我该说什么呢？他们给我们的答复显然违背正义。不过，你应当仔细考虑此事。此外，不久前我们曾嘱咐你的经验与我们最卓越的儿子长官商议，以便将\OriginalTerm{水道}{formarum}的管理委托给副伯爵\Person{奥古斯都}，因为他各方面都勤勉能干。而迄今为止，你竟拖延此事，甚至没有告知我们你的进展。因此，即使现在，也要赶紧竭尽全力与我们的那位最卓越的儿子商谈，务必将水道完全委托给上述那位最杰出的人，以便他能在某种程度上成功修复它们。因为这些水道如此被轻视和忽视，除非给予更多关注，短期内它们将彻底崩溃。如你所知，此事多么必要，对公众多么有益，你必须尽最大努力，使之如我们所说，委托给上述那人细心管理。签署于一月，第五指称期。\par

\SourceRecord{Translated by James Barmby.；Nicene and Post-Nicene Fathers, Second Series；Vol. 13.；Edited by Philip Schaff and Henry Wace.；Buffalo, NY: Christian Literature Publishing Co.,；1898.；<http://www.newadvent.org/fathers/360212024.htm>.}

% structure: p0001 source=h1 target=section reason=page-title

\section{第十二卷，第二十五封书信}

\Emph{致\Person{罗马努斯}，卫护者（\OriginalTerm{卫护者}{Defensorem}）}。\par

\Person{格列高利}致\Person{罗马努斯}等。\par

众所周知于您的阁下，我们已任命\Person{彼得}为卫护者（\OriginalTerm{卫护者}{defensorem}），他出身于属于我们教会的名为\Place{维特拉}的庄园。因此，既然我们应当善待他，但又不使教会蒙受损失，我们命令您以此谕令严格告诫他，不得以任何借口或托辞，试图让其子女在合法身份所属的庄园之外结婚。在此事上，您的阁下也须非常谨慎，并警告他们，无论如何不得离开他们因出生而隶属的产业。因为，若他们中有人（我们认为不会发生）胆敢离开它，他可确信，我们绝不会同意他们中任何人在其出生的庄园之外居住或结婚，而且他们的土地也应当被登记。那么请注意，如果因您的疏忽，他们中有人试图做我们所禁止的任何事情，您将承担不小的风险。\par

\SourceRecord{Translated by James Barmby.；Nicene and Post-Nicene Fathers, Second Series；Vol. 13.；Edited by Philip Schaff and Henry Wace.；Buffalo, NY: Christian Literature Publishing Co.,；1898.；<http://www.newadvent.org/fathers/360212025.htm>.}

% structure: p0001 source=h1 target=section reason=page-title

\section{第十二卷，信函第28}

\Emph{致\Place{努米底亚}主教\Person{哥伦布}}。\par

\Person{格列高利}致\Person{哥伦布}等。\par

因我等早已知晓，你的兄弟之爱在司铎的庄重和教会的热忱上何其卓著，故我等认为你有充分理由参与那些需要斥责之事的审理，以免因纵容而拖延，致使人人以为凡力所能为者皆可允许。如今，我等之弟兄、\Place{泰格西斯}城主教\Person{保利诺}，被其执事及领受圣秩者指控在体罚方面对他们过于严苛，其事如何，无需赘述，因在此申诉到达我等之前，据他们陈述，此事已为尔所知。并且，既然上司不应有权残暴惩罚属下，我等已致函我兄及同为主教的\Person{维克多}——他在你们中持有首席权——请他与你的兄弟之爱以及其他你认为合适的弟兄及同为主教一同审理并彻底调查我等前述主教弟兄与其神职人员之间的案件。愿你的爱德如此密切而审慎地关注此事，使呈报于我等之事不致未经听证即行放过，以免教会内滋生裂痕——而教会本应竭力消除裂痕。若其神职人员对他的控诉确凿，则请你以我等教会判决之力，审处其罪过——他既不屑自行纠正——使其一方面当下感受到所犯之重罪，另一方面今后学会不能做超乎合法之事。尤其，我等劝勉你，要热切研习，为天主的缘故而操练我们所知你已拥有的热忱。\par

再者，由于我前述弟兄\Person{保利诺}据称通过买卖圣职的异端授予圣秩——此事听闻令人惊骇——请你与前述首席主教或其他人士一起，也以勤勉彻底调查此事。若发现属实（愿天主禁止），则须努力并采取行动，使那既不怕接受贿赂、也不怕给予贿赂的人，都受到教规惩罚的打击，以使对他们的纠正成为对许多人的儆戒。在此致命之根尚未壮大、尚未杀害更多人之前，愿它以全体会议的判决被定罪，从而无人再敢接受或给予任何东西以换取任何圣秩，也无人因私情而被擢升，而应全凭功绩，以免教会秩序混乱，且美善生活被藐视——若有不配者获得功绩的报酬。\par

此外，我等已命令我等的\OriginalTerm{档案员}{Chartularius}\Person{希拉留斯}，若案件有此需要，他不得拒绝参加你们的调查。\par

因此，若有需要，你可致函通知他，表示希望他到你处，以便与他共同处理此事，从而更好地决定应规定之事。三月，五年小纪。\EditorialNote{注意：此日期不见于若干抄本。}\par

\SourceRecord{Translated by James Barmby.；Nicene and Post-Nicene Fathers, Second Series；Vol. 13.；Edited by Philip Schaff and Henry Wace.；Buffalo, NY: Christian Literature Publishing Co.,；1898.；<http://www.newadvent.org/fathers/360212028.htm>.}

% structure: p0001 source=h1 target=section reason=page-title

\section{第十二卷·第二十九封书信}

\Emph{致\Person{维克多}主教}．\par

\Person{格列高利}致\Person{维克多}等．\par

一方面，我们得知弟兄们以父爱般的慈心关怀自己的子女，这令我们喜悦；另一方面，当对其它弟兄的尊重或对自己司铎职位的考虑都不能制止他们做违法之事时，我们也同样感到悲伤。那么，针对我们的弟兄、\Place{特格西斯}城主教\Person{保利诺}，由他的圣职人员和那些领受神圣神品者所提出的控诉是多么严重与严厉，我相信贵弟兄也清楚知道，因为从远处传到我们这里的事，在附近的你们决不会不知晓。而且，由于需要非常谨慎，以免他们所控诉的他在职权范围以外所造成的身体伤害被纵容发生，或因默许而恶化，明显的过犯应始终被教规管制所抑制，以便一项裁决既能惩戒过去，又能作为未来的规范。因此，你应同我们最亲爱的共同弟兄\Person{哥伦布}主教，以及你认为合适召请的其他司铎，通过彻底调查来审查我们上述弟兄与他圣职人员之间的案件。如果申诉人的控诉属实，你们当通过正规的惩戒来纠正此事，使他既明白自己所做的恶事，又学会未来不逾越职权的界限。不要允许他像据说的那样藐视你的职位等级，免得他的藐视危及他自己，并归咎于你。因为凡属下所犯的过失，若不认真纠正，都会影响到居上位者。\par

还有另一件事，即我们的同一位弟兄\Person{保利诺}据说以金钱授予教会神品，你应当全面且极其严格地调查。如果这事被清楚证实（我们希望不会如此），愿你对天主的热情燃烧起来，为这错误报复，使授予者的贪吝转为惩罚，并且，使这非法的圣秩无效，被授予者不得享受他所渴望的野心目标。在此我们劝诫你，并首先提醒你，贵弟兄要如此操心：在一个人犯罪导致西蒙尼异端的不义在你们地区壮大之前，通过一次勤勉召开的公会议，用你判决的修剪刀将它从根部斩断。因为凡是不因自己的职位而热切燃烧去纠正这暴行的人，不要怀疑他将会与那使这特别恶行得以开始者同分一份。因此，如同我们所说，你必须警醒而热切地行动，以使你们的公会议——迄今在天主的护佑下未沾染此类恶名——绝不至于被这邪恶的毒药污染和毁坏。\par

此外，我们已命令我们的书记官\Person{希拉鲁斯}，如果情况需要，他不得推迟加入你们。因此，如有必要，请你写信通知他需要他来你处，以便你与他一起，在天主助佑下，能以有益的方式决定所有这些事。\par

\SourceRecord{Translated by James Barmby.；Nicene and Post-Nicene Fathers, Second Series；Vol. 13.；Edited by Philip Schaff and Henry Wace.；Buffalo, NY: Christian Literature Publishing Co.,；1898.；<http://www.newadvent.org/fathers/360212029.htm>.}

% structure: p0001 source=h1 target=section reason=page-title

\section{第十二卷，第三十二封书信}

\Emph{致比扎西乌姆会议全体主教}。\par

\Person{格里高利}致全体，等等。\par

如同对长上表示应有的敬意与尊崇是值得称赞和明智的，同样，若长上身上有需要纠正之处，出于正直与敬畏天主，也不应因任何纵容而拖延，免得疾病（愿天主不许）开始蔓延全身，而头部的不治之症却未被治愈。不久前，有人向我们报告了关于你们的首席主教、我们的弟兄\Person{克莱门提乌斯}的某些事情，令我们心中感到不小的悲痛。但由于各种苦难的压力，尤其是身边敌人的猖獗，我们没有时间调查此事。鉴于此事十分严重，绝不可不经调查就放过，我们现以万般谨慎与急切的心情敦促你们的弟兄手足，务必以一切方式查明实质真相，以便：若事情属实，可藉教会法律惩罚加以剪除；若为虚假，我们的弟兄的无辜不致长久遭受恶名报告的撕裂。为此，为使调查中不致有懈怠的麻木，我们劝告你们：无论任何人的利益、偏袒、甜言蜜语，或其他任何事物，都不可软化你们任何人，使你们在筛核所报告之事时偏离真理之路；反而你们应如司铎般束腰，以探求真理。因为若有人胆敢在此事上怠惰或表现出疏忽，就让他知道，他在全能天主面前将成为上述罪行的同谋——他竟未因对天主的热情而被推动去彻底查明凶恶罪行的缘由。\par

\SourceRecord{Translated by James Barmby.；Nicene and Post-Nicene Fathers, Second Series；Vol. 13.；Edited by Philip Schaff and Henry Wace.；Buffalo, NY: Christian Literature Publishing Co.,；1898.；<http://www.newadvent.org/fathers/360212032.htm>.}

% structure: p0001 source=h1 target=section reason=page-title

\section{第十二卷，第五十封书信}

\Emph{致\Person{厄娄吉}，亚历山大宗主教。}\par

\Person{额我略}致\Person{厄娄吉}等。\par

这些书信的携带者，来到西西里后，脱离了异端\OriginalTerm{一性论者}{Monophysites}的错误，与至圣公教会联合。他们前往真福\OriginalTerm{宗徒之长}{Prince of the apostles} \Person{圣伯多禄}的圣堂，请求我写信将他们推荐给您的真福，使他们如今不致遭受附近异端者的任何伤害。又因为其中一人说，他所住的隐修院是由其亲属建立的，他渴望得到您的圣座的权柄，好使院内的异端者或回归圣教会的怀抱，或被逐出同一隐修院。我们向您指明此事即可：因为我们知道您的真福，凡关乎全能天主的热诚之事，您必全力以赴。至于我，请你们为我祈祷，因为在我忍受的\Place{伦巴第}人的刀剑之中，我深受痛风之苦。\par

\SourceRecord{Translated by James Barmby.；Nicene and Post-Nicene Fathers, Second Series；Vol. 13.；Edited by Philip Schaff and Henry Wace.；Buffalo, NY: Christian Literature Publishing Co.,；1898.；<http://www.newadvent.org/fathers/360212050.htm>.}

% structure: p0001 source=h1 target=section reason=page-title

\section{第十三卷，书信一}

\Emph{致罗马市民}\par

\Person{额我略}，天主众仆之仆，致他最亲爱的儿子们，\Place{罗马}市民。\par

已传到我耳中，有些心术不正之人在你们中间散布了一些与圣教信德相悖的错误说法，禁止在安息日做任何工作。对此，我只能称他们为\Emph{假基督}的宣讲者；假基督来临时，将既使安息日、也使主日免于一切工作。因为他假装死亡并复活，所以他要人尊敬主日；又因为他强迫人民犹太化，以恢复法律的外在仪式，并使犹太人的背信服从自己，所以他希望遵守安息日。\par

因为先知所说的这句话：\BibleQuote{在安息日，你们不可从你们的城门搬入重担}（\BibleRef{耶 17:24}），在法律仍可按照文字遵守的时候，是可以持守的。但是，自从全能天主的恩宠，我们的主耶稣基督显现以后，法律中那些以预像方式所说的话，就不能再按字面遵守了。因为，若有人说应当遵守这安息日的规定，他必然也主张应当献上血肉的祭献；他也必然主张肉身的割损礼仍需保留。但让他听宗徒保禄反对他所说的：\BibleQuote{如果你们受割损，基督对你们就毫无益处}（\BibleRef{迦 5:2}）。\par

因此，我们以属神的方式接受并持守关于安息日所记载的这一切。因为安息意即休息。但我们拥有真正的安息，就在我们的救赎主，主耶稣基督自身内。凡在祂内承认信德之光的人，若借着双眼将私欲偏情的罪恶带入自己的灵魂，就是在安息日从城门搬入重担。所以，如果我们不借着身体的感官将任何罪的重担带入灵魂，我们就是在安息日不从城门搬入重担。因为我们读到，我们的主和救赎主在安息日行了许多事，以致祂斥责犹太人说：\BibleQuote{你们中谁不在安息日解开他的牛或驴，牵去饮水呢？}（\BibleRef{路 13:15}）如此，如果真理本身亲自命令安息日不应按文字遵守，那么，谁若按照法律的文字遵守安息日的休息，他岂非正是与真理本身矛盾？\par

还有一事也传到了我耳中：那就是有邪恶之人向你们宣讲，说任何人都不应在主日沐浴。诚然，若有人为了奢侈享乐而渴望沐浴，我们在任何其他日子也不允许这样做。但若是为了身体的需要，即便在主日我们也不禁止。因为经上记载：\BibleQuote{从来没有人恨恶自己的肉身，而是保养顾惜它}（\BibleRef{弗 5:29}）。又记载：\BibleQuote{不要为肉体安排，去满足私欲}（\BibleRef{罗 13:14}）。因此，禁止为肉体安排去满足私欲的人，当然允许为身体的需要而安排。因为，如果在主日沐浴身体是罪，那么在那一天也不该洗脸。但如果身体的一部分洗脸是被允许的，为何在需要时整个身体反被禁止呢？然而，在主日应停止世上的劳作，全心专注于祈祷，这样，若六日之内有何疏忽，可在主复活之日以恳祷赔补。\par

最亲爱的儿子们，你们要怀有坚定的恒心和正确的信德，持守这些事；蔑视愚人之言，不要轻易相信你们所听闻他们所说的一切，而要以理智的天平衡量，好使你们在稳固的坚定中抵挡谬误之风，得以获得天国稳固的喜乐。\par

\EditorialNote{两个手抄本中，一为\Emph{Colbert}，一为\Emph{Vatic. F.}，注明\Quote{九月，第六小纪}}\par

\SourceRecord{Translated by James Barmby.；Nicene and Post-Nicene Fathers, Second Series；Vol. 13.；Edited by Philip Schaff and Henry Wace.；Buffalo, NY: Christian Literature Publishing Co.,；1898.；<http://www.newadvent.org/fathers/360213001.htm>.}

% structure: p0001 source=h1 target=section reason=page-title

\section{第十三卷，第五封信}

\Emph{致\Person{埃泰里乌斯}，\Place{里昂}主教}.\par

\Person{额我略}致\Person{埃泰里乌斯}主教。\par

虽然我们所说的话令我们深感痛苦，而兄弟般的同情更令我们哀哭，不容我们对所闻之事妄下论断；然而，我们对所承担治理职责的关怀，以急迫的鞭策刺动我们的心，促使我们以极大的审慎关注教会的利益，并在其可能遭受不可挽回的损失之前，妥善安排应行之事。因此，我们从某些人的报告中得知，一位主教的头部遭受病患，其精神错乱时所行之事，闻之令人悲叹哭泣。为此，为避免牧人生病时，羊群遭受潜伏者獠牙的撕裂（愿天主禁止），或教会自身利益遭受不可挽回的损失，我们必须以谨慎的措施处理此事。所以，既然一位主教在世时，除非他自愿辞职，否则因不可避免的疾病而非罪行使他不能履职时，没有理由在他不辞职的情况下另立他人；那么，若他常有清醒的间歇，应自行提出声明，宣告因疾病导致理智受损，已不再胜任这一职务，并请求另立他人接替。如此行后，由全体选举，将另一位堪当之人隆重祝圣为主教接替其位；但在此前提下，只要这位主教仍在世，其应有之费用应由同一教会供给。然而，若他从未恢复健全心智，则须选立一位经核准、生活堪受称许的可靠之人，堪当管理教会，照顾灵魂的益处，以纪律的约束抑制不安分者，照管教会财产，并在各方面显出成熟与能干。并且，若他能在现任患病主教去世后仍存留，则应在他去世后祝圣为其接替者。\par

您的兄弟之职应知此事保留给您的兄弟之职，因该教会隶属您的兄弟之职的教区，以便您的兄弟之职能查考被选之人的生命、品行与操守。若您的兄弟之职满意，且其身上无违教规之严苛之处，则他必须经由您的兄弟之职的祝圣才能获得其命定之品秩。因此，请您的兄弟之职以警觉的审慎行事并如此安排，使天主的教会不再因任何疏忽而受损，并且您的兄弟之职不仅以言语，更以榜样警告您的兄弟之职的同僚司铎，当勤劳照管神圣场所。\par

\SourceRecord{Translated by James Barmby.；Nicene and Post-Nicene Fathers, Second Series；Vol. 13.；Edited by Philip Schaff and Henry Wace.；Buffalo, NY: Christian Literature Publishing Co.,；1898.；<http://www.newadvent.org/fathers/360213005.htm>.}

% structure: p0001 source=h1 target=section reason=page-title

\section{第十三卷，第六封}

\newline{}\Emph{致法兰克人的王后布伦希尔德}。\par

\newline{}额我略致布伦希尔德等。\par

\newline{}在您其他诸般卓越之处中，这一点超越其余而居首位：在这尘世的风浪中——这些风浪常以烦扰的骚动混乱统治者的心神——您却将心志引回对神圣敬礼的热爱，并为圣洁处所的安宁而筹谋，仿佛没有其他忧虑困扰您。因此，既然权贵们的此种行为常成为臣民极大的保障，我们宣称法兰克民族超乎其他民族而有福，因为它堪当拥有一位如此充满一切美德的王后。\par

\newline{}从您信函所载的信息中，我们得知您已在奥古斯托杜努姆（\Emph{欧坦}）郊外建造了圣玛尔定教堂，以及一座天主婢女的隐修院，并在同一城中建造了一所救济院。我们为此大大欢喜，并向全能的天主谢恩，是祂激励您内心的诚恳从事这些事。为使我们在此事上也被视为与您的善行有某种程度的分享，我们已按照您的意愿，为这些处所及其居住者的安宁与保护颁赐了特权；我们并未延迟接纳陛下您的愿望，即使是在最小程度上。\par

\newline{}此外，我们首先以慈父般的爱德向您致候，并通知您：我们已按照您来信所嘱，接见了我们杰出的儿子们，亦即您的臣仆和使节布尔戈阿尔杜斯与瓦尔马利卡里乌斯，并给予他们私下会谈；他们已向我们详细透露了他们自称受命传达的一切。关于这些事情，将来如何办理，我们将负责通知陛下。至于我们，凡可能的、有益的、有助于您与共和国之间和平安定的事，我们愿在天主之下，以最大的热忱使之成就。\par

\newline{}至于我们至为可敬的弟兄与同僚主教门纳斯，我们在查问了对他的指控后，发现他毫无可责之处，且他已在该受祝福的宗徒伯多禄至圣身体前宣誓并作出满意答复，从而证明了他所受的对其名誉的反对并未影响他。因此，我们允许他洗清罪名、获释回任。因为，若他在任何方面有罪，我们依教会法惩罚其过犯固然正当；但当他有无辜为凭时，我们不应再扣留他或以任何方式困扰他。\par

\newline{}此外，关于某位主教（正如前述卓越人士告知我们的）因头部疾病而无法管理其职务，我们已致函我们的弟兄与同僚主教埃泰里乌斯：若该主教有此疾病的间歇清醒期，他应呈递请愿书，声明自己不能胜任其职位，并请求为其教会另祝圣一位主教。因为，在主教生前——其退出职务管理并非出于自身过失，而是因病——神圣教规绝不允许另有一人代替他被祝圣。但若他从未恢复健全心智的运用，则应寻求一位生活与言行均美善的人，此人既能照顾灵魂，又能以有益的管控处理同一教会的事务与利益；他应是这样的人：若主教去世，他能继任其位。至于若有任何人应被擢升到神圣等级或任何圣职，我们已规定此事应提交并告知我们上述至为可敬的弟兄埃泰里乌斯，前提是该事属于他的教区；如此，经过查询，若这些人并无神圣教规所谴责的过失，则由他亲自祝圣他们。那么，请陛下您的关怀与我们的安排相结合，以便您所极度关切的教会利益不受损害，同时也使您善行的赏报得以增加。\par

\newline{}同样，有人询问一名重婚者可否被接纳进入神圣等级，我们已根据教规禁令完全禁止。因为，天主不许在您的时代——您在此时代做了如此多虔敬与宗教的事——您容许任何违反教会法规的事发生。\par

\newline{}此外，前述卓越人士、我们的儿子们，在向我们递交一份清单后，要求我们——他们说这是奉您的命令——从我们这里派遣一人前往高卢，以便在召开主教会议时，在全能天主的引导下，纠正一切违犯至圣教规的事。于此，我们认出您的光荣所怀的关切：您如何为灵魂的生命和您王国的稳定着想。因为，您敬畏我们的救赎者，并处处遵守祂的诫命，在这件事上也如此行动，好使您王国的治理长久存续，并使您本人经过漫长的岁月后，也能从地上的王国过渡到天上的王国。在适当的时机，若我们所说的蒙天主喜悦，我们将负责实现陛下您的可敬愿望。\par

\newline{}因此，为保护陛下您来信中提及的处所，我们已谨慎地按您的意愿安排了诸事。但为避免我们的法令万一被那些地方的执政者以他们被禁止做某些事为由而压制，此同一法令必须载入公共档案，以便它可保存在您的王室档案中，如同保存在我们的档案中一样。\par

\newline{}愿全能的天主永远使陛下您因敬畏祂而得以保全，并藉着宗徒之长、蒙福的伯多禄的代祷——您将之托付于他——如此成就您与我们诸位杰出的儿子、您的孙儿们的国王的愿望，使您在他们持续的福祉中获得稳定的喜乐，正如您所愿。颁于十一月，小纪第六年。\par

\SourceRecord{Translated by James Barmby.；Nicene and Post-Nicene Fathers, Second Series；Vol. 13.；Edited by Philip Schaff and Henry Wace.；Buffalo, NY: Christian Literature Publishing Co.,；1898.；<http://www.newadvent.org/fathers/360213006.htm>.}

% structure: p0001 source=h1 target=section reason=page-title

\section{第十三卷，第七函}

\Emph{致\Place{法兰克}王\Person{特奥多里克}。}\par

\Person{额我略}致\Person{特奥多里克}等。\par

我们怀着喜悦收到了您表示健康平安的书面致函，由此看出，您的智慧超乎年龄，显然，天上恩宠的眷顾将王权统治托付给阁下，实为法兰克民族的福气。您这一点，尤其在别的事上，足以博得称颂和赞叹：您知道我们女儿，即您卓越的祖母，因着对全能天主的爱而渴望之事，您便极其热心地迅速予以协助，从而既在此世幸福为君，又在来世与天使一同为王。既然这出自天主的恩赐，是伟大明智判断的结果，我们便如此迅速而乐意地满足了阁下所愿，以至我们执行的敏捷显出了您的善行何等令我们喜悦。\par

此外，我们以慈父般的柔情向您致候，并告诉您，您托付您麾下的显贵人物，即我们之子\Person{布尔戈阿尔杜斯}和\Person{瓦尔马里卡里乌斯}，与我们交涉的一切事务，均已在一场私下会晤中向我们披露。我们极口称赞您，既明智地关注当前，如同您所当行，又急于通过您与共和国之间持久的和平，为未来的安全作出安排，使你们合而为一，从而有益地长久扩展您王国的稳固。关于此事，我们将在适当时候告知您天主乐意安排之事。因为就我们而言，凡经证明有利于和平之事，我们都渴望并力求实现。唯有一点：我们的意愿关乎适宜之事，天主的意愿亦当如此，离了祂我们什么也不能作。愿\OriginalTerm{天主圣三}{Holy Trinity}使您常在其敬畏中前进，并如此依祂所悦纳的节制安排您的心怀，以致现今从您那里赐予您的臣民以喜乐，将来从祂自身赐予您以无尽的喜乐。\par

\SourceRecord{Translated by James Barmby.；Nicene and Post-Nicene Fathers, Second Series；Vol. 13.；Edited by Philip Schaff and Henry Wace.；Buffalo, NY: Christian Literature Publishing Co.,；1898.；<http://www.newadvent.org/fathers/360213007.htm>.}

% structure: p0001 source=h1 target=section reason=page-title

\section{第十三卷·第八封信}

致院长\Person{塞纳托尔}\par

\Person{格列高利}致\Person{塞纳托尔}，医院（或客舍，\OriginalTerm{客舍}{xenodochii}）的司铎兼院长。\par

当天主教君王的心，等等。\par

\EditorialNote{参见下一封信（第九封信），此信与之完全一致，第十封致\Person{卢波}的信亦同，只是在所涉及的人名和地名上称谓不同，且在第八和第九封信中，在\Quote{或宣判她（\Emph{他}）无罪}一句之后，增加了以下段落。}\par

按同样的规定，根据创立者的意愿，我们明令：日后被任命为同一客舍及隐修院的院长或司铎者，不得以任何秘密谋划胆敢接受主教职位，除非他首先被解除院长职务，另由他人接替；以避免他因不公正的开支耗费客舍或隐修院的财产，给穷人、旅客或其他仰赖其资源生活者造成严重匮乏。此外，我们禁止主教在没有院长和司铎同意的情况下，有权从该处调离任何隐修士以擢升教会职务，或因任何原因调离，以免这方面的擅权发展到如此程度，以致本应通过吸纳人员来建立的场所反而因人员调离而遭毁坏。\par

\SourceRecord{Translated by James Barmby.；Nicene and Post-Nicene Fathers, Second Series；Vol. 13.；Edited by Philip Schaff and Henry Wace.；Buffalo, NY: Christian Literature Publishing Co.,；1898.；<http://www.newadvent.org/fathers/360213008.htm>.}

% structure: p0001 source=h1 target=section reason=page-title

\section{第十三卷，第九书}

\Emph{致院长\Person{塔拉西雅}}\par

\Person{额我略}致\Person{塔拉西雅}等。\par

当天主教君王的心被天主的恩宠所预先推动，燃炽着热切的渴望，以致自愿要求那些本应由主教劝诫才能引发他们去做的事，那么，这些事自应更以欢欣喜悦的心情予以授予，因为即使他们不愿去做，我们也本应要求他们去做。因此，依据我们最卓越的王室后裔\Person{布伦希尔德}及其孙\Person{特奥德里克}的书信，我们将以下特权凭我们现时权威的法令赐予、授予并确认给圣玛利亚隐修院——该院位于\Place{奥古斯托杜努姆}城，由可敬的\Person{希亚格略}主教所建，其中有一群天主的婢女组成团体，由你管理——我们规定：无论任何君王、任何主教，或任何身具任何尊位的人，或任何其他人等，都不得以任何理由或藉口，减少或剥夺该隐修院所拥有的、由上述君王后裔所赐予的财物，或将来由任何其他人以其自身财产所施予的财物，亦不得将之挪为己用，或假借其他虔敬用途以掩饰其贪婪而转让之。凡已奉献或将来可能奉献于该院的一切，我们愿你以及在你职位和地位上继承你的人，从现在起，不受侵犯、不受干扰地拥有，但你们须将这一切完全用于那些为供养和管理而赐予的人的用途上。\par

我们也规定，在上述隐修院的院长去世后，不得通过任何暗中策划的诡计来任命另一位院长，而只应任命那位同省君王在敬畏天主并征得修女们同意后选定并筹备祝圣的人。\par

在此项下，我们还补充规定，为杜绝一切贪腐之机，任何君王、任何司铎，或其他任何人，无论本人或通过代理人，都不得因该院长的祝圣，或任何与该隐修院相关的事由，接受任何黄金或任何形式的报酬；同时，这位院长本人也不得因其祝圣而给予任何财物，以免藉此机会消耗已奉献或可能奉献给虔敬之地的财物。再者，由于据说在你们地区，不法之徒常常寻找欺骗虔诚妇女的机会，我们规定：除非因犯罪而必须如此，否则不得以任何方式剥夺或撤换该隐修院的院长。因此，若有人提出此类控告，不但应由\Place{奥古斯托杜努姆}城的主教审理案件，还须召请其他六位同僚主教协助调查，彻底查明事实，如此全体意见一致，严格按教会法判决，或在她有罪时予以惩处，或在她无罪时宣告无罪。\par

因此，我们这份谕令和法令文书所载全部内容，我们规定为你以及所有在你同一职衔和地位上继承你的人，并为所有与此相关的人，永久遵守。此外，任何人——无论是君王、司铎、法官或世俗人士——若明知我们此项法令而试图违犯，应丧失其权柄和尊位之尊严，并须知因其所行之不义，在天主审判前已陷于有罪。除非他归还所不正当夺取之物，或以相称的补赎为自己违法之行而哀痛，否则应被禁止领受我们的天主和主、救赎者耶稣基督至圣的圣体宝血，并在永久的审判中遭受严厉的惩罚。愿所有维护该院正义之人都蒙受吾主耶稣基督的平安，使他们在此世获得善行的果实，并在严厉的审判者手中获得永久的平安之赏报。\par

\SourceRecord{Translated by James Barmby.；Nicene and Post-Nicene Fathers, Second Series；Vol. 13.；Edited by Philip Schaff and Henry Wace.；Buffalo, NY: Christian Literature Publishing Co.,；1898.；<http://www.newadvent.org/fathers/360213009.htm>.}

% structure: p0001 source=h1 target=section reason=page-title

\section{第十三卷，第十封信}

\Emph{致卢波院长。}\par

额我略致卢波司铎兼院长。\par

\Emph{当天主教君王们的心，等等。}\par

\SourceRecord{Translated by James Barmby.；Nicene and Post-Nicene Fathers, Second Series；Vol. 13.；Edited by Philip Schaff and Henry Wace.；Buffalo, NY: Christian Literature Publishing Co.,；1898.；<http://www.newadvent.org/fathers/360213010.htm>.}

% structure: p0001 source=h1 target=section reason=page-title

\section{第十三卷，第十二封信}

\Emph{致\Person{帕斯卡修斯}，\Place{内亚波利}（\Place{那不勒斯}）主教。}\par

\Person{格里高利}致\Person{帕斯卡修斯}等。\par

那些怀着纯正意向愿意使基督宗教的外方人皈依真信仰的人，应当讲求友善，而非严厉；以免那些本可因委婉的劝谕而受吸引的人反因对立而被远远驱走。因为凡行事相反、并假借这种意向想使人们中断其持续遵守本族礼仪的人，显然是在谋求自己的事业而非天主的事业。就是说，居住在那不勒斯的犹太人向我们诉苦，声称某些人不合理地试图阻止他们庆祝某些节日的典礼，以致他们不能像自远古以来他们及其祖先一直合法地遵守和举行一样来庆祝他们的节庆。如果这是真的，这些人看来是徒劳无益。因为，如果连这种长期不惯的禁令对他们的信仰和皈依都毫无用处，那又有什么意义呢？我们为何要为犹太人制定他们该如何举行庆典的规则，如果我们不能因此而赢得他们呢？因此，我们应当如此行事：即他们应更受理性和善意的感召而愿意跟随我们，而非逃离我们；并且，通过用他们自己的圣经向他们证明我们所说的话，我们就能在天主的助佑下，使他们皈依慈母教会的怀抱。\par

所以，你的兄弟之爱，只要可能，当藉天主的助佑，激励他们皈依，不再允许他们因自己的典礼而受困扰；但应让他们享有完全的自由，去遵守和庆祝他们所有的节庆和假日，正如至今他们和他们的祖先长期以来一直遵守和举行的那样。\par

\SourceRecord{Translated by James Barmby.；Nicene and Post-Nicene Fathers, Second Series；Vol. 13.；Edited by Philip Schaff and Henry Wace.；Buffalo, NY: Christian Literature Publishing Co.,；1898.；<http://www.newadvent.org/fathers/360213012.htm>.}

% structure: p0001 source=h1 target=section reason=page-title

\section{第十三卷，第十八封信}

\Emph{给西西里的某些主教}\par

\Signature{格里高利致西西里主教莱奥、塞孔迪努斯、若望、多努斯、卢奇杜斯、特拉扬。}\par

正如我们受到宗徒言论的劝诫，彼此分享属灵的助佑——同样，在天主安排下我们因被赋予管理穷人事务的治理职责而必须处理的事宜上，司铎的援助也不应缺席。因此，在派遣持信者、我们的\OriginalTerm{书记官}{Chartularius}阿德里安前往治理我们教会的产业，即锡拉库萨地区时，我们认为有必要将他推荐给各位主教，以便在习俗要求之处，你们能给予他帮助；如此，当他蒙你们以身体的援助支持其工作，并以你们祈祷的属灵助佑使他顺利执行所承担的一切时，在天主与他合作下，他就能圆满完成我们所托付给他的事。至于你们自己，应当在全能天主面前以善行如此行事，以至在你们的作为中找不到任何可被天主的审判击打之处，或可被任何伺机攻击你们的人指控之处。因为我们已命令我们上述的书记官，若他得知我们至可敬的弟兄主教们有任何不当行为，他应先私下以谦逊的劝诫责备他们；若这些事未得改正，他应迅速告知我们。\par

此外，据报告给我们，在我们可敬的前任时代，由当时负责教会产业的执事\OriginalTerm{天主的仆人}{Servusdei}安排，你们各教区的司铎，当你们出去给婴儿施行坚振时，不应受到过重的负担。因为已设定了一笔固定的金额，据我所知，经你们同意，由同一些司铎为\OriginalTerm{神职人员}{clericorum}的服务而支付。而这曾获批准的做法，据说如今全然未被遵守。因此我劝告各位主教，努力不成为属下之人的负担，且若有任何怨言，要予以减轻，况且你们也不应偏离已经作出的决定。因为如果你们使托付给你们的人免受怨言，你们就是在为今生和来世自己的利益着想。\par

\SourceRecord{Translated by James Barmby.；Nicene and Post-Nicene Fathers, Second Series；Vol. 13.；Edited by Philip Schaff and Henry Wace.；Buffalo, NY: Christian Literature Publishing Co.,；1898.；<http://www.newadvent.org/fathers/360213018.htm>.}

% structure: p0001 source=h1 target=section reason=page-title

\section{第十三卷，第二十二封信}

\Emph{致贵妇\Person{鲁斯蒂齐亚纳}。}\par

\Person{额我略}致\Person{鲁斯蒂齐亚纳}等。\par

每当有人从皇城来到我们这里，我们总是留心询问您的身体健康；但由于我的罪过之故，我听到的总是令我难过的事，因为您本就虚弱多病，据说痛风之苦仍在加重。但我恳求全能之主，愿您身体所遭遇的一切都能有益于您灵魂的健康，愿暂时的鞭笞为您预备永恒的安息，愿祂通过有终结的痛苦赐予您无终结的喜乐。至于我，我生活在如此呻吟和繁忙之中，以至于我厌恶自己活到了现在所度过的这些日子，我唯一的安慰就是期待死亡。因此，我恳求您为我祈祷，使我早日脱离这血肉的牢笼，不再受如此巨大的痛苦折磨。\par

此外，我必须告知您，有一个名叫\Person{贝亚托尔}的人来到这里，他自称是\OriginalTerm{私产官}{comes privatarum}，并且正在做许多反对众人的事，但主要是反对您尊贵之人和您最尊贵的孙女们所属之人，仿佛是在调查公共事务。我们固然不会允许他行事不公，但也不能阻碍公共利益。因此，请您尽可能与最虔诚的君主们交涉，使他们能制止他的任何非法行为。因为任何混乱都无益于公共利益，而且他似乎也追讨不到什么大数目的东西。我恳求代我问候我最亲爱的儿子\Person{斯特拉特吉乌斯}大人，愿全能的天主为了祂自己和为了您而养育他，并常以祂的恩宠和这位年轻大人的生命来安慰您。此外，关于您返回此地的事，我该写些什么呢？您知道我是多么渴望此事。但当我看到羁绊您的事务时，我又感到绝望；因此我恳求万有的创造者，无论您身在何处，也无论您将来身处何处，愿祂伸出右手保护您，并保全您脱离一切凶恶。\par

\SourceRecord{Translated by James Barmby.；Nicene and Post-Nicene Fathers, Second Series；Vol. 13.；Edited by Philip Schaff and Henry Wace.；Buffalo, NY: Christian Literature Publishing Co.,；1898.；<http://www.newadvent.org/fathers/360213022.htm>.}

% structure: p0001 source=h1 target=section reason=page-title

\section{Book XIII, Letter 26}

\Emph{致副执事\Person{安提米乌斯}。}\par

\Person{格列高利}致\Place{坎帕尼亚}副执事\Person{安提米乌斯}。\par

我们听说，我们的弟兄和同僚主教\Person{帕斯卡西乌斯}多年来懒惰疏忽，以致在任何方面都不被承认为主教；并且，他的教会、隐修院，以及教会的子女或受压迫的穷人，都感受不到他对他们有任何爱德的热情；他也不向恳求他的人提供任何正义的帮助，而且（更严重的是）他无论如何都不能忍受听取智者和爱慕正义之人的劝告，以便他至少能从别人那里学到他自己不能注意的事；相反，他忽略牧者职责之事，全神贯注于造船，毫无益处。由此，据报告，他已经损失了四百索利多甚至更多。他的过失还加上这一点：据说他每天带着一两个圣职人员下到海边，衣着如此卑贱，以至成为他自己臣民的话题，而且给陌生人留下如此恶劣可鄙的印象，以至于被认为没有丝毫主教应有的品格或尊严。若果真如此，要知道这并非没有你们的过错，因为你们迟迟没有责备和约束他。那么，鉴于这一切不仅使他失去信誉，而且明显给司铎职位带来耻辱，我们希望你为此事召集他，在其他司铎或他的一些尊贵的子女面前，并劝诫他：摆脱懒惰的恶习，不要闲散，而要警醒于照顾他的教会和隐修院，向他的子女展示父亲般的爱德，在正义所称赞的案件中谨慎地挺身而出为穷人辩护，并欣然接受智者的劝告；为的是那座城市能因他的关怀而得到安慰，而他自己也能成功地掩盖他怠惰的过失。然而，如果在我们这次劝诫之后，他竟敢像往常那样疏忽（我们不相信会发生这种情况），那么他必须被送到我们这里，以便在我们面前，他能在敬畏天主之下，学会一个司铎该做什么以及如何去做。\Signature{颁于三月，征收期第六年。}\par

\SourceRecord{Translated by James Barmby.；Nicene and Post-Nicene Fathers, Second Series；Vol. 13.；Edited by Philip Schaff and Henry Wace.；Buffalo, NY: Christian Literature Publishing Co.,；1898.；<http://www.newadvent.org/fathers/360213026.htm>.}

% structure: p0001 source=h1 target=section reason=page-title

\section{第十三卷，第27封信}

\Emph{致副助祭安特弥。}\par

额我略致坎帕尼亚副助祭安特弥。\par

凡我们屡次听闻有关我们的弟兄和同主教们的事，显示他们应受责备并令我们悲伤，需要迫使我们不能不稍加考虑他们的改正。既然有人向我们报告，坎帕尼亚的主教们如此疏忽，以致忘记他们职务的尊严与品格，既不向他们的教会，也不向他们的儿子们显示父爱般的警觉照顾，也不关心隐修院，也不保护受压迫者和穷人，因此我们命令你，并特此授权你召集他们，并严格以我们的名义劝告他们，不可再怠惰，而要表现出他们的司铎热忱和关怀，并警觉地做他们应当公正且合乎天主所做的事，以便不再有任何关于他们的怨言激怒我们。然而，如果你发现他们中任何人在此事之后仍然疏忽，则不得容许任何借口，将他送到我们这里来，以便通过惯常的纪律处分，使他们感受到拒绝改正那些应受责备且极其应受谴责的事情是多么严重。\par

\SourceRecord{Translated by James Barmby.；Nicene and Post-Nicene Fathers, Second Series；Vol. 13.；Edited by Philip Schaff and Henry Wace.；Buffalo, NY: Christian Literature Publishing Co.,；1898.；<http://www.newadvent.org/fathers/360213027.htm>.}

% structure: p0001 source=h1 target=section reason=page-title

\section{第十三卷 第三十一封书函}

\Emph{致佛卡皇帝}。\par

\Person{额我略}致\Person{佛卡}奥古斯都。\par

光荣归于至高之天主，祂按经上所载，改变时序，转移邦国，因为祂使众人明白，祂曾藉其先知恩许所言：\BibleQuote{至高者在人国中掌权，将国赐与谁，就赐与谁}（见：\BibleRef{达 4:17}）。因为在全能天主奥秘的安排中，人间统治时有交替；有时，为击打许多人的罪恶，祂兴起一人，使其严苛，使臣民的颈项在磨难的重轭下弯下，正如我们在苦难中长期所体验的。但有时，当慈悲的天主决意以祂的安慰使许多忧伤之心重获喜悦时，祂高举一人至统治的顶峰，并藉祂怜悯的胸怀，将欢愉的恩宠注入众人心中。我们确信将很快在这丰盈的欢跃中得到坚固，因为我们为您的圣德已登至尊之位而欢喜。\BibleQuote{愿诸天欢乐，愿大地踊跃}（见：\BibleRef{咏 95:11}）；愿全体共和国的人民，迄今深陷苦难，因您的仁政而欢欣鼓舞。愿敌人傲慢之心在您统治的轭下被制服。愿臣民被压碎沮丧的心灵因您的仁慈而复苏。愿天恩之权能使您令敌人畏惧，使您的虔敬对臣民和善。愿整个共和国在您最幸福的年代得到安息，使和平的掠夺在法律的幌子下暴露。愿关于遗嘱的阴谋终止，以及强索的恩惠。愿每个人重新拥有自己的财产，使他们能够高兴地、无所畏惧地拥有他们无欺诈所获得之物。愿每个人的自由现在终于在帝国的轭下归还。因为外邦人的君王与共和国的皇帝有区别：外邦人的君王是奴隶之主，而共和国的皇帝是自由人之主。但我们最好通过祈祷而非提醒您说这些事。愿全能天主的恩宠之手中的您的圣德之心，在每一个思想和行为上得以保守；凡应公正而行之事，凡应仁慈而行之事，愿那住在您胸怀中的圣神引领，使您的仁爱既在暂时的王国中被高举，在多年之后又达到天上之国。颁于六月，第六小纪。\par

\SourceRecord{Translated by James Barmby.；Nicene and Post-Nicene Fathers, Second Series；Vol. 13.；Edited by Philip Schaff and Henry Wace.；Buffalo, NY: Christian Literature Publishing Co.,；1898.；<http://www.newadvent.org/fathers/360213031.htm>.}

% structure: p0001 source=h1 target=section reason=page-title

\section{第十三卷，第三十四封信}

\Emph{致书记官潘塔莱奥。}\par

额我略致潘塔莱奥等。\par

\OriginalTerm{阁下}{Your Experience}还记得，您在真福宗徒伯多禄至圣遗体上发了什么样的誓愿。因此，我们也毫不犹豫地将叙拉古地区教产调查的重任托付给您。那么，您有责任常将您的忠信和对同一位真福宗徒伯多禄的敬畏摆在眼前，行事如此，既在今生在人前，又在最后审判时在天主全能者面前，都不致受责。现在，我们从我们的\Emph{文书官}萨莱里乌斯的报告中得知，您发现迫使教会的佃农 (\Emph{coloni}) 缴纳谷物的\Emph{莫底乌斯}，竟有二十五\Emph{塞克斯塔留斯}之巨。我们对此深恶痛绝，并遗憾您发现此事过迟。因此，我们欣喜地得知您已摧毁了那个不公的\Emph{莫底乌斯}，并设立了公义的度量。然而，既然上述\Emph{文书官}也提到了在您阁下任内，承包人 (\Emph{conductores}) 已通过欺诈从两个地区收取了若干粮谷，因此，正如我们为未来而欣喜于您热心摧毁不义之\Emph{莫底乌斯}，我们也思虑过去的罪愆，唯恐若承包人从农夫手中骗取之物归于我们，我们便牵连在他们的罪中。因此，我们希望您的阁下，以全然的忠诚，全然的正直——敬畏全能的天主，并牢记真福宗徒伯多禄的严格——在每一田庄 (\Emph{massam}) 中列出贫苦困乏的佃农名单，并用查明确系诈骗所得的钱款购买牛羊猪只，分给各贫苦佃农。我们愿意您与最可敬的主教若望阁下、我们的\Emph{文书官}兼\Emph{主管}阿德里安商议此事。此外，若为商议起见，也应请我们的儿子尤利安阁下参加，以便不使他人知晓，一切严守秘密。因此，请你们商议，是应以金钱还是以实物救助这些贫苦佃农。但无论共同基金如何，首先，如我所说，列出名单，然后按各人贫困程度尽力分配。因为，如外邦人的导师所证，我样样都有，且富足有余；我并非求金钱，而是求赏报。\EditorialNote{参阅斐理伯书 4。}所以，你们行事，以便在审判之日，能从我托付给您阁下职务的劳苦中，向我显示成果。若你们纯洁、忠实、奋力地做这事，你们不仅在现世从子女身上得回报，将来在永恒审判者的审查中，也将获得圆满的酬报。\par

\SourceRecord{Translated by James Barmby.；Nicene and Post-Nicene Fathers, Second Series；Vol. 13.；Edited by Philip Schaff and Henry Wace.；Buffalo, NY: Christian Literature Publishing Co.,；1898.；<http://www.newadvent.org/fathers/360213034.htm>.}

% structure: p0001 source=h1 target=section reason=page-title

\section{第十三卷，第三十八封信}

\Emph{致福卡斯皇帝。}\par

额我略致奥古斯都福卡斯。\par

我们欢喜感恩地思量，我们应向全能天主献上何等赞美，因为悲伤的轭已被除去，我们在您仁爱陛下的皇恩下进入了自由的时代。至于您的安宁陛下未能按古例在朝廷找到一位宗座的常驻执事，并非由于我的疏忽，而是由于极其严峻的必要。因为，当本教会所有仆役都因害怕那般压迫与艰难的时期而退缩逃避时，不可能强令他们中的任何人前往皇城留在朝廷任职。但现在他们得知，您的宽仁陛下，因天主恩宠的安排，已达至帝国的顶峰，先前极为恐惧前往那里的人，甚至自己赶往您的脚下，为此欢欣鼓舞。但是，鉴于其中一些人因年老而体弱，几乎不能承受辛劳，另一些人则深陷于教会事务之中，而本信的呈递者，曾是我们所有\OriginalTerm{守护者}{defensores}中的首位，我素知他勤勉，在生活、信德和品性上经受考验，我认为他适合被派往您虔诚陛下的脚下。因此，我蒙天主许可，祝圣他为执事，并尽力迅速将他派往您处，以使他在找到适当时机时，能够将此地正在发生的一切禀明您的宽仁陛下。我恳求您的安宁陛下垂听，请您惠予倾听他虔诚的恳求，以便您越是从他的叙述中真实了解我们的苦难，就能越迅速地怜悯我们。因为伦巴第人的日常兵刃和多少次入侵，使我们在如今长达三十五年间受压迫，其程度我们无法用言语完全描述。但我们信赖全能的上主，祂会为我们完成祂已开始的安慰美事，并且，祂在共和国中兴起了虔诚的君主之后，也必消灭残暴的敌人。愿天主圣三保佑您多年长寿，使我们能在长久的等待之后，更长久地为您的虔诚美德而欣喜。\par

\SourceRecord{Translated by James Barmby.；Nicene and Post-Nicene Fathers, Second Series；Vol. 13.；Edited by Philip Schaff and Henry Wace.；Buffalo, NY: Christian Literature Publishing Co.,；1898.；<http://www.newadvent.org/fathers/360213038.htm>.}

% structure: p0001 source=h1 target=section reason=page-title

\section{第十三卷，第三十九封}

\Emph{致\Person{莱昂蒂娅}皇后。}\par

\Person{额我略}致\Person{莱昂蒂娅·奥古斯塔}。\par

什么舌头能够言说，什么心思能够思想，我们应当全能的天主在你们帝国的安宁中向祂献上何等大的感恩，因为长久以来如此沉重的重担已从我们颈上移去，而至尊统御的柔和之轭已然回归，这是臣民乐于承受的！因此，愿荣耀归于万物的创造者，由天使的歌唱诗班颂扬，愿地上的人们献上感恩，因为整个共和国，经历了诸多悲伤的创伤，如今终于找到了你们安慰的膏油。因此，我们更必须恳求全能天主的仁慈，使祂将你们虔诚的心永远握在祂的右手中，并藉天上恩宠的帮助疏导你们的思念，为使你们的安宁能更正义地统治那些服事你们的人，正如你们更真实地知道如何服事万有的主宰。愿祂使你们成为祂在热爱天主教信仰上的勇士，因为祂以祂的仁慈对待，使你们成为我们的皇帝。愿祂将热忱与温和注入你们的心智，使你们总能以虔诚的热忱，不放过任何得罪天主的过犯而不予惩罚，而在面对任何冒犯你们自身的过失时，又能忍耐与宽恕。愿祂藉着你们的虔诚，赐给我们\Person{普尔赫里娅·奥古斯塔}的宽仁，她因对天主教信仰的热忱，在神圣的大公会议中被称作新\Person{赫肋纳}（\Emph{\WorkTitle{迦克墩大公会议纪要}一卷}）。愿全能天主的仁慈赐予你们更长的岁月，与我们最虔诚的主上共同生活，以使你们的生命越延长，臣民的安慰就越得到巩固。\par

我或许本应请求，请陛下将真福宗徒伯多禄的教会特别托付于你们，这教会迄今一直遭受着针对它的严重阴谋。但是，我深知你们爱慕全能的天主，故我不应请求你们出于虔诚的良善而自动展示的事。因为你们越敬畏万有的创造者，就越能深深爱慕祂的教会，就是那一位的教会，曾对他说：\BibleQuote{\Emph{你是伯多禄，在这磐石上我要建立我的教会，阴间的门决不能战胜她；并且我要将天国的钥匙交给你；凡你在地上所束缚的，在天上也要被束缚；凡你在地上所释放的，在天上也要被释放。}} \BibleRef{玛 16:18}。因此，我们毫不怀疑，你们将以何等强烈的爱将自己与那位你们渴望藉以摆脱一切罪过的人联系在一起。那么，愿他成为你们帝国的守护者，愿他成为你们在地上的保护者，愿他成为你们在天上的转求者：为使你们通过解除臣民的重担，并使他们在你们的帝国中喜乐，在多年之后，你们能在天国中喜乐。\par

\SourceRecord{Translated by James Barmby.；Nicene and Post-Nicene Fathers, Second Series；Vol. 13.；Edited by Philip Schaff and Henry Wace.；Buffalo, NY: Christian Literature Publishing Co.,；1898.；<http://www.newadvent.org/fathers/360213039.htm>.}

% structure: p0001 source=h1 target=section reason=page-title

\section{第十三卷 第四十书}

\Emph{致\Person{齐里亚库斯}，\Place{君士坦丁堡}宗主教}\par

\Person{格列高利}致\Person{齐里亚库斯}等。\par

最亲爱的弟兄，你当勤谨地观察，主的声音说：\BibleQuote{我把我的平安赐给你们}（见：\BibleRef{若 14:27}），和平之德何其伟大！因此我们应当如此持守和平之爱，绝不给纷争留地步。然而，我们若不能在心中和行为上持守和平之根的谦卑——这谦卑正是和平的创始者所教导的——便无法活在和平之中。因此，我们以相称的爱德恳求你，用你心灵之足踏碎那常与灵魂为敌的世俗傲慢，赶紧从教会中除去那乖谬傲慢之称号的绊脚石，免得你或许被发现与我们的和平团体分离。但愿我们同有一个心神，一个意念，一个爱德，一个在\Person{基督}内的联系，他愿意我们成为他的肢体。因为愿你的圣德思量，没有你所向他人宣讲的和平，且因此出于骄傲而不去冒犯你的弟兄，这是多么艰难，多么不体面，多么残忍，多么与司铎的目标相悖！反而要致力于这件事：你如何以谦卑之剑击倒那虚荣无益的傲慢的创始者，好使在此胜利中，圣神的恩宠能认你为他自己的居所，从而使经文在你身上清楚地应验：\BibleQuote{天主的宫殿是圣的，你们就是这宫殿}（见：\BibleRef{格后 6:17}）。\par

我们全力将这封信的呈递者、我们最心爱的共同儿子、执事\Person{博尼法斯}推荐给你，使他在一切需要之事上，能恰当地找到你圣德的援助。\par

\SourceRecord{Translated by James Barmby.；Nicene and Post-Nicene Fathers, Second Series；Vol. 13.；Edited by Philip Schaff and Henry Wace.；Buffalo, NY: Christian Literature Publishing Co.,；1898.；<http://www.newadvent.org/fathers/360213040.htm>.}

% structure: p0001 source=h1 target=section reason=page-title

\section{第十三卷，第四十一封信}

\Emph{致亚历山大宗主教\Person{厄瓦格里乌斯}}\par

\Person{格列高利}致亚历山大主教\Person{厄瓦格里乌斯}。\par

某日，我与亲近的朋友们谈论教会习俗时，一位曾在亚历山大城学习医术的人告诉我们，他有一位同学，同一讲座的学生，极端堕落，他说，这少年突然被祝圣为执事。他补充说，此人是靠贿赂和赠礼获得祝圣的；因为他承认，这种习俗在圣亚历山大教会中盛行。听罢，我十分惊讶，极其诧异于最圣洁可赞颂的\Person{厄瓦格里乌斯}主——他的口舌曾使如此多的异端回归大公信仰——竟未能从圣亚历山大教会中根除西莫尼异端。倘若他那伟大而卓越的教导对此置之不理，那么还有谁的劝诫或纠正能够改善呢？\par

为此，为了你灵魂的赦免，为了你赏报的增加，为使你的工作在威严审判者眼中全然完美，你应当赶紧从你那至圣的宗座——这也是我们的宗座——彻底拔除并根绝西莫尼异端，这异端是教会中首先出现的。\par

因为正因如此，教会圣秩的圣德从许多人身上失落，因为人们晋升这些圣秩，并非基于他们的生活和行为，而是基于贿赂。但若是寻求功德而非贿赂，不配的人就不会前来领受祝圣。而你所得的赏报将会愈发增加，正如那些被晋升至神圣圣秩的善人们将自己奉献于拯救灵魂的关怀。\par

\SourceRecord{Translated by James Barmby.；Nicene and Post-Nicene Fathers, Second Series；Vol. 13.；Edited by Philip Schaff and Henry Wace.；Buffalo, NY: Christian Literature Publishing Co.,；1898.；<http://www.newadvent.org/fathers/360213041.htm>.}

% structure: p0001 source=h1 target=section reason=page-title

\section{第十三卷，第四十二书}

\Emph{致亚历山大宗主教厄洛吉乌斯}\par

额我略致厄洛吉乌斯等。\par

我们向全能的天主致以极大的感谢，因为当经上所记载的得以实现时，心中口里便体验到爱德的甘甜芳香：\Quote{有好消息从远方来，就如拿冷水给口渴的人喝}\BibleRef{箴 25:25}。此前，我因驻节代表、书记官博尼法斯来自皇城的信函而深感不安，信中说您那最甘饴、最可亲的圣座曾遭受眼疾之苦。这封信使我陷入深沉的忧伤。但忽然间，藉着我们的创造者与救赎者之顺利恩宠，我收到了您真福的信函，得知我所听闻的身体不适已经痊愈，便极其喜乐，因为随之而来的心灵欢愉，其程度不亚于先前悲伤的苦涩。因为我们深知，赖全能天主的助佑，您的生命是许多人的健康。当一位训练有素、技术娴熟的舵手掌舵时，水手们便能在波涛中安稳航行。\par

此外，在为您健康而喜悦的同时，我还有另一件欢欣鼓舞的理由：我得知藉着您的口舌，教会之仇敌日益减少，主的羊群日益增多。因为通过您舌头的犁铧，天上的谷物每日增长，并在高处的仓廪中增多；因此在您身上，我们欣喜地看到所记载的实现：\Quote{没有牲口，槽头干净；丰收的粮食，全凭牛力}\BibleRef{箴 14:4}。由此我们清楚地看到，您带领逃散者回归全能天主的侍奉越多，在他面前积攒的功德就越大。而您获得的功德越多，就能越充足地获得所求。因此我恳求您更为恳切地为我这个罪人祈祷，因为身体之痛、心灵之苦以及众多蛮族刀剑下巨大的死亡蹂躏，都极大地折磨着我。在这一切之中，我所需的是永恒的慰藉，而非暂时的，我凭自己无法通过祈祷获得，但我深信将藉您真福的代祷而得到。去年我未收到您圣座的任何信函，甚为忧苦。诚然，您的祝福（虽无信函）既已送出亦已收到。但既然您的言语比您的赠礼更令我喜悦，故所赠之物未能使我如所期望的那般满足。然而，我已嘱我二位的共同儿子——执事埃皮法尼乌斯写信给贵最圣教会的执事亚历山大和依西多尔，告知已收到所寄之物。\par

此外，我曾写信给您说，我已备好制造桅杆和舵的大块木材，但所来的小船无法承载；此后您也未回信答复。故此，若您需要这些木材，请告知我们共同的儿子博尼法斯——我们现正派他作为驻节代表前往皇城——以便他通知我准备，并使之在您真福差人来取时备妥。\par

此外，我们给您送去一个小小的十字架，其中嵌有来自您的爱侣——宗徒圣伯多禄和圣保禄——锁链的祝福；愿此祝福常施于您的双眼，因为藉着这同样的祝福，常有诸多奇事发生。\par

愿全能的天主感动您真福之心，使之挂念为我不断祈祷，并以祂的右手保护您及您的一切，并在多年岁月之后，引领您进入天国。\par

我们收到了您至圣的兄弟会所描述的那些祝福，即圣马尔谷的祝福，由您惠赐而来；我们感谢您的仁爱，因为从这些外在之物，我们得知您内在对我们的情谊。\par

\SourceRecord{Translated by James Barmby.；Nicene and Post-Nicene Fathers, Second Series；Vol. 13.；Edited by Philip Schaff and Henry Wace.；Buffalo, NY: Christian Literature Publishing Co.,；1898.；<http://www.newadvent.org/fathers/360213042.htm>.}

% structure: p0001 source=h1 target=section reason=page-title

\section{\WorkTitle{第十四卷 第二封信}}

\Emph{致撒丁岛的\OriginalTerm{守护者}{Defensorem}\Person{维塔利斯}}\par

\Person{格列高利}致\Person{维塔利斯}等。\par

从蒙贤卿告知我们的信息中，我们发现撒丁岛所建立的\OriginalTerm{医院}{xenodochia}正遭受严重忽视。因此，我们至为可敬的兄弟及同主教\Person{雅努阿里乌斯}本应受到最严厉的责备，若非他的高龄、淳朴以及您告知我们他后来又患上的疾病使我们克制。\par

有鉴于此，他既处于如此状况，无法适宜地安排任何事务，就请您在我们的严格授权下警告该教会的司库和\OriginalTerm{总铎}{archpresbyter}\Person{埃皮法尼乌斯}，要他们自行谨慎而有益地努力整顿那些\OriginalTerm{医院}{xenodochia}。因为，此后若再有任何疏忽，让他们知道，他们绝无可能以任何方式或任何程度在我们面前为自己辩解。\par

此外，既然撒丁岛的业主们向我们请愿，鉴于他们受到多种负担的困扰，我们允许您前往君士坦丁堡为他们申诉。我们也写信给我们至爱的儿子\Person{博尼法斯}，请他在为地方争取救济方面尽力相助。\par

此外，关于您告知我们那些没有司铎的教会，我们已写信给我们上述至为可敬的兄弟及同主教\Person{雅努阿里乌斯}，要他供给他们；但条件是，并非所有人都从他自己的教会中被选为主教。因为他应当如此供给其他教会，以免导致自己教会缺乏可能对其有益的人员。\par

关于您告知我们的，有人在低级修会中曾陷于罪，却被擢升管理某些隐修院一事，他们本不应承担院长之职，除非在彻底改过迁善并经过适当的补赎之后。但既然您说他们已经承担了院长职务，就必须关注他们的生活、品行和对职责的专注。如果他们的行为与其职务并无不符，就让他们继续留在现任职务中。否则，应予以撤职，另任命他人，以利于托付给他们的灵魂。\par

此外，关于圣赫尔马斯隐修院，该院由我们的兄弟在虔诚女士\Person{蓬波尼亚娜}的府邸中建立，鉴于应以温和而非严厉对待，请贤卿设法温言劝慰该女士，使她既不至于忽略创立者之意志而自陷于罪，也不致未能妥善地照顾隐修院的利益。此外，关于上述蓬波尼亚娜先前曾将若干女孩的会衣改换，并使她们在隐修院内改宗之事，您决不可容许她们被从她那里带走或受到惊扰；而应在天主的保护下，让她们继续处于目前的生活方式中。\par

关于您所写的收回教会、隐修院或其他用于虔诚用途的财产之事，必须劝诫利益相关者，他们应尽力在您的支持和协助下以一切方式收回。但如果他们意外地疏忽，或无论如何找不到应当收回之人，那么请贤卿尽数查明并收回，但在发现后收回时，不可显得以高压手段对任何人提起诉讼。关于您告知我们的\Person{霍图拉努斯}与\Person{托马斯}的\OriginalTerm{医院}{xenodochia}之事，我们迄无所知。因此，请贤卿仔细调查迄今给予的皇帝谕令，并按其内容安排一切，并将您所行之事告知我们。\par

关于您所写我方兄弟及同主教\Person{雅努阿里乌斯}在举行祭献时常常遭受极大痛苦，以致要经过长时间才能回到\OriginalTerm{感恩经}{canon}中他中断之处，且很多人怀疑是否应从他祝圣的圣体领受圣体一事，应劝诫他们完全不要惊慌，而以全然的信德与安全感领受圣体，因为人的疾病既不改变也不玷污神圣奥迹的祝福。不过，我们仍应私下极力劝告上述兄弟，每当他感到任何不适即将发作时，不应继续举行，以免使自己遭受轻视，并引起软弱者心灵的冒犯。\par

此外，虔诚女士\Person{蓬波尼亚娜}向我们抱怨说，她已故女婿\Person{埃皮法尼乌斯}的遗产——该\Person{埃皮法尼乌斯}曾指定其妻、上述\Person{蓬波尼亚娜}之女\Person{玛特罗娜}为该遗产的\OriginalTerm{用益权人}{usufructuary}，以利于他指示在其府邸建立的隐修院，并在\OriginalTerm{用益权}{usufruct}消灭后也以一切方式有利于该隐修院——连同其他被证明属于同一\Person{玛特罗娜}的占有之物，已被贤卿和我们至为可敬的兄弟及同主教\Person{雅努阿里乌斯}不义地夺走，因此至今未有任何东西支付给她女儿，或对隐修院有益。现在，如果事实如此，而您意识到做了任何不当之事，请毫不延迟地归还所夺之物；或者，如果您认为情况并非如此，以免对方似乎受到有害的委屈，决不要延误将案件提交给她同意的仲裁员，以通过明确的裁决宣告她的抱怨是否真实和公正。\par

\SourceRecord{Translated by James Barmby.；Nicene and Post-Nicene Fathers, Second Series；Vol. 13.；Edited by Philip Schaff and Henry Wace.；Buffalo, NY: Christian Literature Publishing Co.,；1898.；<http://www.newadvent.org/fathers/360214002.htm>.}

% structure: p0001 source=h1 target=section reason=page-title

\section{第十四卷，第四封信}

\Emph{致\Person{法蒂努斯}，\Place{帕诺尔姆斯}的\OriginalTerm{护卫者}{Defensorem}。}\par

\Person{额我略}致\Person{法蒂努斯}，等等。\par

关于我们的弟兄和同为主教的\Person{厄希拉辣图斯}，正如你自己也知道的那样，有些事传到我们耳中，认为应当以严厉的惩罚来处置。但是，由于此事已由我们最可敬的弟兄和同为主教的\Person{良}平息了，他也宣称他曾在那个案件中担任法官，我们认为应当将他（即\Person{厄希拉辣图斯}）送回他自己的教会，考虑到我们通过将他扣留在此如此之久所施加于他的，或许对他来说已经足够了。因此我们命令你的\OriginalTerm{阁下}{Experience}注意他的品行和行为，并频繁地劝告他，使他能够表现出在向他的\OriginalTerm{神职人员}{clericis}广施仁慈的爱德方面的关切，并在必要时纠正过失。但我们也希望你也劝告他的神职人员，让他们向他表现出谦逊，以及主所命令的服从，并且绝不敢在他面前表现出傲慢。如果他们中的任何人，即无论是主教还是神职人员，无视你的劝告，那么你，凭借我们的这项权威，要么以教规的强制来纠正不顺从的罪，随你判断，要么赶紧向我们报告，使我们能够安排如何以纪律的缰绳，使那些被邪恶倾向的刺棒刺激而犯罪的人不偏离正途。\par

\SourceRecord{Translated by James Barmby.；Nicene and Post-Nicene Fathers, Second Series；Vol. 13.；Edited by Philip Schaff and Henry Wace.；Buffalo, NY: Christian Literature Publishing Co.,；1898.；<http://www.newadvent.org/fathers/360214004.htm>.}

% structure: p0001 source=h1 target=section reason=page-title

\section{第十四卷，书信七}

\Emph{致\Place{科尔丘拉}主教\Person{阿尔基松}。}\par

\Person{额我略}致\Person{阿尔基松}等。\par

一颗高傲之心的野心并非不应被压制，因为当它不顾神圣教规的效力，冒昧地觊觎属于他人的东西时，这种冒失的过分表现不仅证明有害并造成花费，而且也违背了教会的和平。在阅读了您的弟兄之书信后，我们得知了\Place{欧里亚}城主教在过去或最近关于\Place{卡西奥普斯}营地（该营地位于您的教区内）所做的事。我们感到难过的是，那些本应因您教会所施予的爱心而成为欠债者的人，反而成为了教会的敌人，毫无羞耻地约束自己；最终，以一种违反教会安排、违反司铎节制、违反神圣教规条例的方式，他们企图将上述营地从您的管辖下撤走，并置于他们自己的权力之下，从而成为像主人一样的人，而他们先前是被当作客人接待的。关于此事，因我们可敬的弟兄、\Place{尼科波利斯}总主教\Person{安德肋}，在皇帝命令（该命令责成他审理此案）的支持下，已知在他颁布的判决中确定（正如我们已经了解的那样），上述\Place{卡西奥普斯}营地应当像一直以来的那样，留在您教会的管辖之下。我们赞同这一判决的形式，并按照正义的认可，以宗座的权威确认它，并裁定它在各方面保持有效。因为没有任何公平的理由，没有任何教规的秩序，允许一个人以任何方式占据另一个人的堂区。因此，尽管这种争论的罪责似乎需要相当严厉的处理，因为他们以恶报善，但仍应注意，不要以过分压倒仁爱，也不要拒绝给予那些处于困境中的陌生弟兄所应得的，以免爱德被认为在主教的心里没有作用，如果那些应得极大同情的人被剥夺了安慰的补救。因此，\Place{欧里亚}城的司铎和圣职人员不应被驱逐出上述\Place{卡西奥普斯}营地的住所，而且他们还应被允许以应有的敬意，将他们所带来的\Person{真福多纳图}的圣洁可敬遗体，安放在上述地点中他们可能选择的某一座圣堂里。但应以这样的方式，为您（您的爱德）获得保障，即（营地位于您的教区内）通过出具一份保证书，\Place{欧里亚}主教在其中承诺，将来不为自己索要其中的任何权力、任何特权、任何管辖权或任何权威，如同他是首席主教一样；但是，一旦蒙天主恩宠恢复和平，他们无论如何都应返回自己的地方，如果他们愿意，则带走\Person{圣多纳图}的可敬遗体。因此，只要铭记这一承诺，他们既不可借口任何理由为自己进一步索要任何统治权，而应始终承认自己是那里的客人；您的弟兄的教会也不得在任何程度上对其权利和特权造成损害。\par

\SourceRecord{Translated by James Barmby.；Nicene and Post-Nicene Fathers, Second Series；Vol. 13.；Edited by Philip Schaff and Henry Wace.；Buffalo, NY: Christian Literature Publishing Co.,；1898.；<http://www.newadvent.org/fathers/360214007.htm>.}

% structure: p0001 source=h1 target=section reason=page-title

\section{第十四卷　第八封书信}

\Emph{致执事波尼法爵。}\par

\Person{额我略}致书于\Place{君士坦丁堡}的执事\Person{波尼法爵}。\par

每当那些本应传播和平的人之间的不和令我们忧伤时，我们就应极其谨慎地设法消除纷争的根源，使彼此分歧之人重归于好。关于位于\Place{科孚岛}的\Place{卡西奥普斯}军营之事，以及\Place{欧里亚}的主教如何企图将该营地脱离\Place{科孚岛}主教的管辖，并非法将其置于自己治下，说来话长。为使你的爱德能完全明了一切，我们寄去了我们的兄弟、\Place{科孚岛}主教\Person{阿尔基松}的书信，并已派其属下前往你处，以口头详细禀报情况。在此我们简述如下：先皇\Person{茂里西乌}曾以欺骗方式获得一项谕令，该谕令既违背法律与神圣教规，因而无效，双方争论未决；随后他又发出另一项谕令给我们的已故兄弟、当时任\Place{尼科波利斯}都主教的\Person{安德肋}，命他因双方均属其管辖而审理此案，并按教规终断。该都主教审理并作出判决（副本附后），裁定前述\Place{卡西奥普斯}军营应归属\Place{科孚岛}主教的权力与管辖，因其历来属于该教区；我们赞同其判决，并认为应以宗座权威予以确认。唯恐我们的决定过于严苛而毫无慈悯，我们当时特意如此安排（正如我们寄给你的判决书正文所示）：既不让\Place{欧里亚}城的主教或司铎承担居住义务，也不以任何方式扰乱\Place{科孚岛}教会的特权。然而，就在诉讼开始之初，有人从最温和的皇帝陛下处骗取了一项谕令，竟不顾\Place{尼科波利斯}都主教的判决（该判决基于教会礼规与教规之理由），将前述\Place{卡西奥普斯}军营移交给了\Place{欧里亚}主教——此事我们闻之悲伤，述之叹息——并更加严重地损害了\Place{科孚岛}主教及其司铎，以至于（哀哉！）完全剥夺了\Place{科孚岛}教会的管辖权，而将首要管辖权几乎全部交给了\Place{欧里亚}主教。有鉴于此，我们认为不宜将我们的判决书送交任何人，以免显得违背最仁慈的皇帝陛下的谕令，或者（天主不容！）轻视了他。因此，愿你的爱德勤勉地将全部事实禀报陛下，并坚定陈明此事全然非法、全然恶劣、全然不义，且严重违反神圣教规。切莫让他允许在其时代引入此类损害教会的罪行。同时，向他陈明前述已故都主教的判决内容，以及我们如何确认其裁决，并努力促成我们的判决书连同陛下的谕令一同发往该地，从而显示我们既尊重了陛下，又以合理方式纠正了妄行之事。在此事中，务必竭力争取陛下也发出谕令，责令遵守我们所决定的内容。若能如此，将来一切作假之机将被杜绝。因此，请赶紧在天主助佑下，运用你的警醒来消除这些不法行为，既不使企图作恶者借势损害教会自古以来的惯例，也不使邪恶行为日后成为先例。\par

此外，为使你知晓我们上述兄弟及同侪主教\Person{阿尔基松}声称因\Place{得撒洛尼}教会的\OriginalTerm{代理人}{actionariis}所受的冤屈与压迫，我们已将他的来信转给你的爱德。因此，请召见上述教会的\OriginalTerm{负责代表}{responsalis}到你处，在其面前审理此案，并按你的理智所及之处，写信给我们的兄弟及同侪主教\Person{欧瑟比}，责令他禁止其属下行事不公，并警告他们不可压迫下属，反而应在一切正义之事上予以帮助。此外，我们还希望你的爱德写信给可能已被任命为\Place{尼科波利斯}城都主教之人，让他针对我们上述兄弟\Person{阿尔基松}所诉其教会遭受的损害进行审理，并作出公正裁决，因为此事本身据称未被其前任判决，而悬而未决。\par

\SourceRecord{Translated by James Barmby.；Nicene and Post-Nicene Fathers, Second Series；Vol. 13.；Edited by Philip Schaff and Henry Wace.；Buffalo, NY: Christian Literature Publishing Co.,；1898.；<http://www.newadvent.org/fathers/360214008.htm>.}

% structure: p0001 source=h1 target=section reason=page-title

\section{第十四卷，第十二封信}

\Emph{致\Person{泰奥德琳达}，伦巴第王后}\par

\Person{格列高利}致\Person{泰奥德琳达}王后。\par

您不久前从\Place{热那亚}地区寄来的信函使我们分享了您的喜乐，因为我们得知，全能天主的眷顾赐给了您一个儿子，而且，这对尊座极有光荣的是，他已被接纳进入天主教信仰的共融。对您的基督徒精神，我们本无他念，深信您必会用公教的正直来坚固天主所赐给您的这位儿子，好使我们救主既承认您是祂亲近的仆人，也为伦巴第民族繁荣地养育一位敬畏祂的新王。因此，我们祈求全能天主，既使您遵行祂的诫命，又使我们最杰出的儿子\Person{阿杜洛瓦尔德}在祂的爱中长进，为的是，如同他在今世成为人间伟人，同样也在我们天主眼中因善行而荣耀。\par

至于尊座在信中请求之事，即我们应对我们至爱的儿子，\Person{塞昆杜斯}院长所写的内容作全面答复——谁能在明知其请求对许多人有益的情况下，仍有意推迟他的祈请或您的意愿呢？——除非是疾病阻碍了。但痛风所致的如此严重的病弱已紧紧束缚了我们，使我们几乎无法起身，不仅无法口述，甚至无法说话；您派遣的使者，即呈递此信者，也知道此事：他们到达时发现我们虚弱，离开时我们正处在生命的极度危险和危机之中。然而，如果全能天主安排我能康复，我将对他所写的一切作出全面答复。不过，我已通过此信的携带者送去了一份在虔诚的\Person{查斯丁尼}时代举行的会议文件，以便我前述至爱的儿子在阅读时能认识到，他所听到的对宗座或天主教会的一切指控都是虚假的。因为我们决不容纳任何异端者的观点，也绝不丝毫偏离我们前任\Person{良}教宗圣善记忆的卷册；我们接受四次神圣会议所定义的一切，并谴责它们所拒绝的一切。\par

此外，我们考虑给我们的儿子国王\Person{阿杜洛瓦尔德}送去一些圣物：即一个带有主圣十字架木片的十字架，以及一本装在\Place{波斯}匣中的圣福音选读。也给我们的女儿，即他的妹妹，送去三枚戒指，其中两枚镶有风信子石，一枚镶有\OriginalTerm{阿尔布拉}{albula}，我请求通过您将这些交给他们，以便我们对他们所怀的爱德因尊座而倍增。\par

此外，我们以父亲般的爱德向您履行问候之职时，恳请您代我们向最杰出的儿子，您的王夫国王，为已缔结的和平表示感谢，并如您素常所做，从今以后以各种方式引导他的心趋向和平；这样，在您的许多善行中，您在天主眼前，将因一个无辜的民族而获得赏报，这民族本可能因冒犯而灭亡。\par

\SourceRecord{Translated by James Barmby.；Nicene and Post-Nicene Fathers, Second Series；Vol. 13.；Edited by Philip Schaff and Henry Wace.；Buffalo, NY: Christian Literature Publishing Co.,；1898.；<http://www.newadvent.org/fathers/360214012.htm>.}

% structure: p0001 source=h1 target=section reason=page-title

\section{第十四卷，第十三封信}

\Emph{致科孚主教阿尔基松}。\par

额我略致阿尔基松等。\par

对于反省而回归健全劝告的弟兄们，不可拒绝仁慈，以免过错在主教们心中似乎比爱德更重要。因此，我们在你爱德的\OriginalTerm{代理人}{responsales}——欧里亚教会读经员彼得——在场的情况下，接待了他，他带着我们的兄弟和同主教若望的信来到我们这里；所带的信读完后，我们便特意询问他是否要对你的那些代理人的指控有所辩驳。他说未受任何指控，除其主教信函所载外无话可答，我们遂毫不迟延地，在天主之下，按教规作出裁定。然而，过了一段长时间，前述彼得拿出了一份文件，声称是其主教给他的；于是案件被拖延。但由于这份文件中发现上述主教说，他曾希望能获准将圣多纳托的圣洁可敬的遗体安放在圣若翰教堂内，该教堂位于所谓的卡西奥普斯军营，并说他准备，因该地经证实在你教区，给你爱德一份保证书，说明此举不会对你造成任何损害——我们认为他的请求不应毫无结果，既然在需要之时，他期望以这样一种方式得到安置，即确保他在所有方面都承认你教会的管辖权。因此，基于这一理由，我们以此信劝勉你的兄弟之爱：不得有任何拖延或借口，你应当提供机会将上述圣人的可敬遗体安放在前述圣若翰教堂内；条件仅是他需事先以书面保证书向你保证，绝不以任何借口声称对该教堂或军营拥有任何管辖权或特权，如同他是该地主教一般，而是要在该地属于你教区的情况下，丝毫不损地维护你教会的所有权利和权力。同时，正如我们这位兄弟所请求的，你也有责任答复他：无论何时，因天主的仁慈恢复和平后，他可以自由返回自己的地方，他应有权不受任何反对地将上述可敬遗体带走。在此，为了避免所做之事似乎是个人的，且可能为重新挑起争端提供机会，你的继任者也应在一切方面被纳入这一承诺之中，以维持现状；如此，通过这预防性保证，他将来既不得违背公平和神圣教规的谕令，在你教区内对任何事物提出主张，而你教会的权利也绝不因这种让步而受到任何损害。\par

\SourceRecord{Translated by James Barmby.；Nicene and Post-Nicene Fathers, Second Series；Vol. 13.；Edited by Philip Schaff and Henry Wace.；Buffalo, NY: Christian Literature Publishing Co.,；1898.；<http://www.newadvent.org/fathers/360214013.htm>.}

% structure: p0001 source=h1 target=section reason=page-title

\section{\WorkTitle{第十四卷，第十六封信}}

\Emph{来自墨西拿主教费利克斯} \Emph{致圣额我略}\par

致最蒙福、最可敬的主，圣父教宗额我略，热爱您的福乐与圣德的费利克斯。\par

您最蒙福的福乐与圣德在天主内的要求是显而易见的。因为，虽然普世大地因宗徒们的宣讲与训导而充满对真信德的遵守，然而基督的正统教会，由建立而奠基，并由忠信教父们最坚固地确立，又藉神圣论述的教导而进一步被建造，同时受到您劝诫训诲的指导。所有至蒙福的宗徒，同等地享有尊严与权威的参与，使众多民族归信；并且藉有益的法令与劝诫，虔诚而圣洁地将那些在神圣预定之恩宠中所预知的人，从黑暗带向光明，从错误带向真信德，从死亡带向生命。追随这些圣宗徒的功德，并完全效法他们的榜样，您可敬的父职以品行的正直和行为的圣洁，偕同他们装饰了天主的教会；并因圣信德与基督徒品行之坚强，命令应当行何事以悦乐天主，且不断履行并完成教宗职务，如此遵守神圣法律的诫命；因为（如宗徒所说）\BibleQuote{\Emph{法律的听者在天主面前不是义人，但法律的实行者将被成义}} \BibleRef{罗 2:13}。\par

当我们思虑这些事时，从某些从罗马来的人那里传来消息，说您曾写信给我们的同僚奥斯定（后来被祝圣为主教，派往盎格鲁民族，并由您可敬的圣德派遣），以及给盎格鲁人（我们早已知道他们通过您归信了），称若已婚的第四亲等血亲不应分开。然而，无论在我与您自幼一同被抚养和教导的那些地方，还是这些地方，从前并无此惯例；我也未在您前任的任何法令中，或在其他教父的普遍或特殊制度中读到过，亦未得知任何智者迄今允许此事。但我从您的圣前任，以及从其他圣教父们，无论是在尼西亚会议还是其他神圣公会议中聚集的，发现这一禁止（即\Emph{禁止婚配}）应遵守至第七亲等血亲；并且我知道，那些生活正直、敬畏主的人对此谨慎遵循。当这些事在我们中间讨论时，又出现了其他事情，似乎有必要就这些事请教您的权威。因为西拉库萨教会的主教本笃以及其他我们同为弟兄的主教们来到我们这里，哭泣着说，他们因世俗及平信徒的过度行为而内心极为不安和痛苦，结果也有人说了一些不公正的话反对他们。\par

在我们的教省中，有些教会关于祝圣圣堂之事存有疑虑；而且，由于它们古老以及看守者的疏忽，不知它们是否曾被主教们奉献过。关于所有这些事，我们恳请由您的圣德和您圣座的权威予以教导；并请求通过您的信函告知我们，是否如前所述，我们听说您写给我们前述同僚奥斯定及盎格鲁民族的内容，是专门写给他们的，还是写给所有人的；我们渴望就此事以及上述其他事得到充分的了解。\par

因为我们并非以挑剔之意（这离我们很远）向您表明我们所读到的、并知道应由信友遵守的内容；而是寻求知道在此事上我们可以合理而忠信地遵守什么。而且，由于我们中间对此问题已有不小的怨言，我们向您请教，如向元首一般，我们应如何回复我们的弟兄、同为主教者；以免我们在此事上存疑，并避免这怨言在您时代及以后的时代中留存在我们中间，并使您的声誉——一向良好且卓越——因诽谤而受损或遭轻蔑，或您的名字（愿天主禁止）在将来的时代被恶言中伤。因为我们，谦卑地心在天主内观察合理之事，被同一爱德之绳与您相连，并作为忠实的弟子在一切事上捍卫您的宗教，从您那里寻求合理之事的知识。因为我们知道，如同首先的使徒们，即圣座的首长，以及他们后来的继承人一贯所做的那样，您也照顾普世教会，尤其是主教们——他们因默观与洞见而被称为主的眼睛；并且您持续思考我们的宗教与法律，如经上所载：\BibleQuote{\Emph{昼夜默思上主法律的人，是有福的}} \BibleRef{咏 1:2}。您的这种默思不仅通过阅读、通过文字的外在表达而被看见，而且因基督的恩宠在您内洋溢，已知确凿地铭刻在您的良心上；而基督主的最神圣法律决不离开您的心；正如先知在圣咏中所说：\BibleQuote{\Emph{义人的口要默思智慧，他的舌头要谈论审判；天主的法律在他的心中}} \BibleRef{咏 36:30}；不是用墨水写的，而是藉永活天主的圣神秘密地写成的；因此不是写在石版上，而是写在心版上。我们恳求您，愿一切黑暗的阴霾被您最智慧的答复和援助驱散，使晨星藉着您，最圣洁的父，照耀我们，并带来一个令各处众人喜乐的教义定义，因为圣教会的光荣教父们被确知曾宣讲了合宜且最虔诚的教义，以获致永生的稳固继承。\par

\Emph{署名}。愿主永远保守您平安、悦乐天主，圣父中之圣父，同时为我们祈祷。\par

\SourceRecord{Translated by James Barmby.；Nicene and Post-Nicene Fathers, Second Series；Vol. 13.；Edited by Philip Schaff and Henry Wace.；Buffalo, NY: Christian Literature Publishing Co.,；1898.；<http://www.newadvent.org/fathers/360214016.htm>.}

% structure: p0001 source=h1 target=section reason=page-title

\section{\WorkTitle{第十四卷，信函第十七}}

\Emph{致\Person{斐理克斯}，\Place{墨西拿}主教。}\par

致我们最可敬的弟兄、主教\Person{斐理克斯}，\Person{格列高利}，天主众仆之仆。\par

我们的头——基督，为此愿意我们作祂的肢体，为使藉着祂丰厚的爱德和信实，祂能叫我们在祂内成为一体。我们理应如此紧紧依附于祂，因为离开祂我们什么也不能做，藉着祂我们才能成为所蒙召的。愿没有什么能使我们与头的高垒分离，否则，若我们拒绝作祂的肢体，就会被祂舍弃，如同从葡萄树上砍下的枝条而枯萎。那么，为使我们堪当成为我们救赎者的居所，让我们全心渴望地坚守在祂的爱内。因为祂说：\Emph{\BibleQuote{凡爱我的人，必遵守我的话，我父也必爱他，我们要到他那里去，与他同住。}}\BibleRef{若 14:23}。如今，最亲爱的弟兄，您的爱德以宗座权威要求我们答复您的询问。虽然我们急于这样做，不是长篇大论而是简洁地（因为某些紧迫的牵挂，因我们罪的阻碍临于我们），但我们仍将以下内容交托给您的注意，以便更广泛地查考和探讨圣教父们的其他规章制度。因为一个被重担和紧迫的忧虑所磨损和压垮的心灵，不能成就如此多的善事，也不能如此自由地谈论这些事，如同一个欢乐且免于沮丧的心灵所能做的。所以，我们并非因偏爱这些牵挂而拒绝给予您的圣德这些和其他我们认为必要的信息，而是为使这里所欠缺的部分得以更充分地探究。\par

因为，效法您前任的榜样，您认为应该就三点咨询宗座——您就是在宗座培育和受教的——即：关于血亲婚姻、关于属下对主教的侵扰、以及关于祝圣教堂的疑虑。那么要知道，我写给盎格鲁民族的主教\Person{奥斯定}（您记得，他是您的弟子）关于血亲婚姻的话，是特别写给他和最近才接受信德的盎格鲁民族的，以免他们因过于严厉的规定而惊恐，从而从他们良好的开端退缩；但这并非普遍地写给其他人。对此，整个罗马城都是我的证人。我也并非有意在这些信函中如此规定：一旦他们在信德中扎根坚固，若发现低于应有的血亲等级就不应分开，或者低于应有的姻亲等级（即直到第七代）就可结合。但是对于那些尚是新奉教者的人，首先常常应当教导他们，并以言传身教，使他们避免非法之事，然后，合理而信实地，排除他们在此类事情上可能已做的事。因为按照那位宗徒所说：\Emph{\BibleQuote{我曾用奶哺养你们，没有用硬食}}\BibleRef{格前 3:2}，我们仅为他们（而不是如上所述为将来）准许了这些宽容，以免迄今根苗脆弱而栽种下的善事被拔除，反而要使已经开始的事得以坚固，并守护直至达于成全。的确，若在这些事上我们做了任何不当之事，就请知道，这并非出于放纵，而是出于怜悯。为此，我也呼求天主作证，祂洞悉众人的思念，万物在祂面前都是赤裸敞开的。因为，若我要摧毁在我之前的人所建立之物，我必被公正地定罪为拆毁者而非建造者，正如真理之声所证：\Emph{\BibleQuote{凡一国自相纷争，必不得站立}}\BibleRef{路 11:17}；并且一切学问和法律若自相纷争，也必被消灭。为此，我们众人必须同心合意地持守我们圣教父们的规诫，毫不争执，而是在一切善愿中同心一致，藉着主的助佑，服从天主的制度和法令。\par

爱德何其美善，它藉着爱，将不在眼前的事物如临眼前般显现在自己心中，它联合分裂的，整顿混乱的，结合不平等的，成全不完全的！杰出的宣讲者正确地称它为成全的纽带，因为其他德行固然产生成全，但爱德如此束缚它们，使它们再不能从爱者的心中松脱。考虑到这一点，在我已经谈过的事上，我以爱德宽容了；我并没有命令，而是建议；我也没有为后来的人立下必须遵守的规则，而是指出两种危险中哪一种更容易避免。那么，若在世俗事务中每个人都应保有自己应有的权利和适当的等级，那么在教会安排中，岂不更不应引起混乱，以免在理应产生和平祝福的地方出现纷争。而这一点将会得到保证，如果什么都不向权力屈从，而一切都向公平屈从。因此，我们的心因您的伟大而大大欢喜，因为我们看到您在行事上如此热忱，以至于关心我们，并费心通过询问我们来探究这些事，为的是这些事不仅使您在人前获得荣耀，而且也在全能的主面前获得赏报。\par

但我们知道，主教的生涯不应受任何过分的搅扰，因为被称为天主宝座的人，实在不应因君王或臣民的任何行动而受干扰。因为，若最公义的达味都不敢伸手加害于显然已被天主废弃的撒乌耳，那么，对于主的受傅者，对于圣教会的宣讲者，就更应小心，不可将毁谤、责难、轻率或侮辱之手加在他们身上，因为对他们的困扰或毁谤，就是触及了基督，他们是在教会中代替基督担任代表职务的！因此，所有信徒都应当极其谨慎，不可暗中或公开地以毁谤或责难撕裂自己的主教，即主的受傅者，要想到米黎盎（即\Person{Miriam}）的例子，她因对天主的仆人梅瑟说话（因为埃塞俄比亚女子）而受到麻风病不洁的惩罚\EditorialNote{户 13}；以及圣咏作者的话：\BibleQuote{不可触碰我的受傅者，不可伤害我的先知}\BibleRef{咏 104:15}。在神圣的法律中，我们也读到：\BibleQuote{不可亵渎神，也不可诅咒你百姓的首领}\BibleRef{出 22:28}。因此，属下，无论是圣职人员还是平信徒，都应当格外小心，若他们偶见主教或上司做了什么可指责的事，也不可鲁莽地加以责备，以免因着指责恶行的立场，反因骄傲的冲动而陷入更深的深渊。还应当告诫他们，当他们考量上司的过失时，不可因此对他们过于胆大妄为。他们要如此思考恶事：因着对天主的敬畏，他们不可拒绝肩负敬畏的轭，因为主教和上司所行的事，即使表面上似乎值得责备，也不可用口中的剑去击打；因为我们知道，我们的前辈和许多其他圣主教都曾规定，羊群不应轻易责备牧人，也不可妄自控告或指责他们，因为当我们冒犯上司时，就是反对那将他们赐予我们的天主安排。因此，梅瑟在得知百姓抱怨他与亚郎时，说：\BibleQuote{我们算什么？你们的怨言不是针对我们，而是针对天主}\BibleRef{出 16:8}。因此，两级的属下都应当被劝诫：当他们观察主人的行为时，应返回自己的内心，不可擅自在责备他们的事上妄为，因为\BibleQuote{门徒不能高过老师，仆人也不能高过主人}\BibleRef{玛 10:24}。\par

关于对教堂奉献礼的疑问（这也是您希望请教我们的事项之一），您应当恰当地持守我们所接受的前人传授的规矩：即，每当对任何人的圣洗或坚振，以及教堂的祝圣产生疑问，且无法通过文字或见证人给出确切说明（是否某人已受圣洗或坚振，或教堂是否已祝圣）时，就应当为这些人施圣洗和坚振，并依教会法规奉献这些教堂，以免此种疑问使信徒陷入丧亡；因为那些不能以确切证据表明已妥善完成的事，在此情形下并非第二次施行。我们愿在神圣恩宠的支持下如此持守，并吩咐您（如您所请）持守和教导；我们不願轻率地违反，而要忠信地遵守圣父们在我们之前所决定的事。为此，我们恳求救赎主的仁慈，以祂的恩宠助佑您，并赐予您将祂所赐的意愿付诸实行，因为在这事上，劳苦的热忱越大，我们所得的赏报也越丰盛。但我们规定，凡已忠信受教、并以坚定根基不可动摇地植立之人，均应遵守其血统直到第七代。只要他们知道自己彼此有姻亲关系，就不得擅自接近这种结合的关联；任何基督徒也不得（且永远不得合法地）与他自己同族的女子结婚，即曾与她同居为妻或曾以任何非法污秽玷污过的女子；因为这样的结合是乱伦的，对天主和所有善人都是可憎的。但我们读到，圣父们很久以前就已决定，乱伦之人不得被视为任何婚姻名义下的成员。因此，我们不愿在这件事上受到您或任何其他信徒的责备，因为在这件事上我们宽恕了英格兰民族，并非是在定立规则，而是顾虑他们可能放弃已经开始的美善，等等。\par

\SourceRecord{Translated by James Barmby.；Nicene and Post-Nicene Fathers, Second Series；Vol. 13.；Edited by Philip Schaff and Henry Wace.；Buffalo, NY: Christian Literature Publishing Co.,；1898.；<http://www.newadvent.org/fathers/360214017.htm>.}
